\documentclass[10pt, b5paper, openany, includefoot]{book}
%b5 = 182mm x 257mm = 18.2cm x 25.7cm = 7.17in x 10.12in
\usepackage[utf8]{inputenc}
\usepackage{fontspec}
\usepackage{ifpdf}
\usepackage[hmargin={2cm, 2cm}, vmargin={2cm, 1.5cm}]{geometry}
\usepackage{tipa}
%\usepackage[stretch=50,shrink=50]{microtype}


\usepackage{enumitem}
\usepackage{multicol}
\usepackage{longtable}


\usepackage[spanish,es-nolayout,es-tabla]{babel}

\usepackage{array}
\usepackage{caption}
\usepackage[explicit]{titlesec}
\usepackage{amssymb}

\usepackage{graphicx}
\usepackage{tipa}
\usepackage{multirow}
\usepackage{rotating}
%\usepackage{natbib}
\usepackage[authoryear,round]{natbib}

\usepackage{gb4e}
%\let\eachwordone=\it

\usepackage{titlesec}
\titlelabel{\thetitle.\quad}

\usepackage{mathabx}

\usepackage[defaultlines=2, all]{nowidow}

\setlength{\columnsep}{0.15in}

\usepackage{fancyhdr}
\pagestyle{fancy}
\renewcommand{\chaptermark}[1]{ \markboth{#1}{} }
\renewcommand{\sectionmark}[1]{ \markright{#1}{} }

%\titleformat{\section}
%{\vspace{0.5em}\bfseries\sffamily\large}
%{\large{\thesection}}{1em}{{#1}}[\vspace{0.5em}]

\newcommand{\gloss}{}
%\newcommand{\gloss}[1]{\textit{\textsf{#1}}{}}

\newcommand{\variantof}{}
\newcommand{\definition}[1]{\textbf{Definition:} \textit{#1}}
\newcommand{\example}[1]{#1}
\newcommand{\exsen}[1]{\textit{#1}}     % Iquito example sentence
%\newcommand{\extran}[1]{\textenglish{#1}}
\newcommand{\extran}{}
\newcommand{\sense}{}
\newcommand{\grammaticalinfo}{}

\newcommand{\relform}{}
\newcommand{\relforms}{}
\newcommand{\discoursenote}{}

\newcommand{\entry}[2]{\small\noindent\hangafter=1\hangindent=0.5em\markboth{#1}{#1}#2 \vskip .03in}

\newcommand{\reventry}[2]{\small\noindent\hangafter=1\hangindent=0.5em\markboth{#1}{#1}\pagebreak[3]#2 \vskip .0in}

\newcommand{\headword}[1]{\hspace{0pt}\textsf{\normalsize{#1}{}}}

\newcommand{\exampleiqu}[1]{$\bullet$~\textsc{ejemplo} \textsf{#1}{}}
%\newcommand{\exampleiqu}[1]{$\bullet$~\textsc{ejemplo} \textit{\textsf{#1}}{}}

\newcommand{\exampleen}[1]{\textit{#1}{}}
%\newcommand{\exampleen}[1]{{#1}{}}
%\newcommand{\exampleen}[1]{\textsl{#1}{}}





\newcommand{\irregpl}[1]{$\blacktriangleright$~\textsc{plural irreg.} \textsf{#1}{}}

\newcommand{\irregplof}[1]{$\blacktriangleright$~\textsc{plural irreg.~de} \textsf{#1}{}}

\newcommand{\lexeme}[1]{$\blacktriangleright\blacktriangleright$~ \textsc{raíz subyacente} \textsf{#1}{}}

\newcommand{\derivroot}[1]{$\blacktriangleright$~\textsc{raíz deriv.} \textsf{#1}{}}
%\newcommand{\impfrt}[1]{\textit{impf.rt.} \textbf{#1}{}}
%\newcommand{\impfrtlab}{\textit{impf.rt.{ }}}

\newcommand{\scientificname}[1]{$\blacktriangleright$~\textsc{nombre científico} \textit{#1}{}}

\newcommand{\pos}[1]{\textsf{(#1)}}

\newcommand{\litmean}[1]{$\blacktriangleright$~\textsc{literalmente} {#1}}

\newcommand{\grammarnote}[1]{$\blacktriangleright$~\textsc{nota gramatical} {#1.}}

\newcommand{\relformiqu}[1]{$\blacklozenge$ \textsc{palabra relacionada} \textsf{#1}{}}
%\newcommand{\relformiqu}[1]{$\blacktriangleright$ \textsc{Palabra relacionada:} {#1.}{}}

\newcommand{\relformiqurt}[1]{$\blacktriangleright$~\textsc{raíz subyacente} \textsf{#1}{}}

\newcommand{\relformpos}[1]{\textsf{(#1)}}

\newcommand{\relformen}[1]{{#1}}

\newcommand{\relformeu}[1]{{#1}}

\newcommand{\note}[1]{$\blacktriangleright$~\textsc{préstamo} {#1}{}}
%\newcommand{\note}[1]{$\bullet$ {#1.}{}}

\newcommand{\sci}[1]{\textit{#1}}

\newcommand{\activemiddle}[1]{$\blacktriangleright$~\textsc{alternante act./med.} \textsf{#1}{}}

\newcommand{\spa}[1]{\textsf{{#1}}}

\newcommand{\spq}[1]{\textit{{#1}}}

% CMB: \iqt is used one time, to make the formatting the same as it appears in the actual dictionary:
\newcommand{\iqt}[1]{\textsf{#1}}

\newcommand{\iqu}[1]{{\textsf{#1}}}
%\newcommand{\iqu}[1]{\textit{\textsf{#1}}}

\newcommand{\eng}[1]{\textit{#1}}

\newcommand{\vartext}[1]{\textsf{#1}}

\newcommand{\gramnoteentry}[1]{$\blacktriangleright$~\textsc{nota gramatical} {#1.}{}}

\newcommand{\posspref}[1]{\textit{#1}.}

\newcommand{\semnote}[1]{$\blacktriangleright$~\textsc{nota semántica} {#1.}{}}
%\newcommand{\semnote}[1]{$\blacktriangleright$~\textsc{Nota semántica:} {#1.}{}}

\newcommand{\semnoteentry}[1]{$\blacktriangleright$~\textsc{nota semántica} {#1.}{}}

\newcommand{\anthronote}[1]{$\blacktriangleright$~\textsc{nota cultural} {#1.}{}}

\newcommand{\anthnoteentry}[1]{$\blacktriangleright$~\textsc{nota cultural} {#1.}{}}

\newcommand{\altpronuncnote}[1]{$\blacktriangleright$~\textsc{pronunciación alternativa} {#1:}{}}
%\newcommand{\pronnote}[1]{\textit{fst.spch.} \textbf{#1}{}}

\newcommand{\altpronunc}[1]{\textsf{#1}}

\newcommand{\socionote}[1]{$\blacktriangleright$~\textsc{nota sociolingüística} {#1.}{}}

\newcommand{\socionoteentry}[1]{$\blacktriangleright$~\textsc{nota sociolingüística} {#1.}{}}

\newcommand{\irregposs}[1]{$\blacktriangleright$~\textsc{forma poseída irreg.} \textsf{#1}{}}

\newcommand{\variants}[1]{{#1}{}}

\newcommand{\rpos}[1]{\textsf{(#1)}}

\newcommand{\sociovarlab}{$\blacktriangleright$~\textsc{var.~sociolingüística{ }}}

\newcommand{\sociovarof}[1]{$\blacktriangleright$~\textsc{var.~sociolingüística de} \textsf{#1}{}}

\newcommand{\irregpllab}{$\blacktriangleright$~\textsc{plural irreg.}{ }de{ }}

\newcommand{\irregposspl}[1]{$\blacktriangleright$~\textsc{forma poseída plural irreg.} \textsf{#1}{}}
%\newcommand{\irregposspl}{\textsc{Plural poseída irreg.}{ }de{ }}
%\newcommand{\irregposspl}{\textsc{$\blacktriangleright$~plural poseída irreg.}{ }}

\newcommand{\freevarlabs}{$\blacktriangleright$~\textsc{vars.~libres{ }}}

\newcommand{\freevarlab}{$\blacktriangleright$~\textsc{var.~libre{ }}}

\newcommand{\freevarof}[1]{$\blacktriangleright$~\textsc{var.~libre de} \textsf{#1}{}} 

\newcommand{\dialectvarlab}{$\blacktriangleright$~\textsc{var.~dialectal{ }}}

\newcommand{\dialectvarof}[1]{$\blacktriangleright$~\textsc{var.~dialectal de} \textsf{#1}{}} 

\newcommand{\dialectvarlabs}{$\blacktriangleright$~\textsc{vars.~dialectales{ }}}

\newcommand{\dialectlabNanay}{$\blacktriangleright$~\textsc{var.~nanaíno{ }}}

\newcommand{\dialectofNanay}[1]{$\blacktriangleright$~\textsc{var.~nanaíno de} \textsf{#1}{}} 
 
\newcommand{\dialectlabChambira}{$\blacktriangleright$~\textsc{var.~chambirino{ }}}    

\newcommand{\dialectofChambira}[1]{$\blacktriangleright$~\textsc{var.~chambirino de} \textsf{#1}{}}

\newcommand{\dialectlabMaasikuuri}{$\blacktriangleright$~\textsc{var.~maasikuuri{ }}}   \newcommand{\pronnote}[1]{/#1/}

                 % you can customize more


\newcommand{\dialectofMaasikuuri}[1]{$\blacktriangleright$~\textsc{var.~maasikuuri de} \textsf{#1}{}}

%\newcommand{\dialectlabInkawɨɨ́raana}{$\blacktriangleright$~\textsc{var.~inkaw}ɨ́ɨ́\textsc{rààna{ }}}    
%\newcommand{\dialectlabInkawɨɨ́raana}{$\blacktriangleright$~\textsc{var.~inkawɨ́ɨ́rààna{ }}}    

%\newcommand{\dialectofInkawɨɨ́raana}[1]{$\blacktriangleright$~\textsc{var.~inkaw}ɨ́ɨ́\textsc{rààna de} \textsf{#1}{}}
%\newcommand{\dialectofInkawɨɨ́raana}[1]{$\blacktriangleright$~\textsc{var.~inkawɨ́ɨ́rààna de} \textsf{#1}{}}

%\newcommand{\dialectlabMajanakáani}{$\blacktriangleright$~\textsc{var.~máájànàkáànì{ }}} 

%\newcommand{\dialectofMajanakáani}[1]{$\blacktriangleright$~\textsc{var.~máájànàkáànì de} \textsf{#1}{}}
       % or any formatting you want



\newcommand{\persvarlabJPI}{$\blacktriangleright$~\textsc{var.~de~JPI{ }}} 

\newcommand{\persvarofJPI}[1]{$\blacktriangleright$~\textsc{var.~de~JPI~de} \textsf{#1}{}}

\newcommand{\persvarlabELY}{$\blacktriangleright$~\textsc{var.~de~ELY{ }}} 

\newcommand{\persvarofELY}[1]{$\blacktriangleright$~\textsc{var.~de~ELY~de} \textsf{#1}{}}
\newcommand{\Fast}{\textsc{Fast}}
\newcommand{\Unspecified}{\textsc{Unspecified}}

\newcommand{\persvarlabHDC}{$\blacktriangleright$~\textsc{var.~de~HDC{ }}} 

\newcommand{\persvarofHDC}[1]{$\blacktriangleright$~\textsc{var.~de~HDC~de} \textsf{#1}{}}

\newcommand{\persvarlabunk}{$\blacktriangleright$~\textsc{var.~personal{ }}}
 
\newcommand{\persvarofunk}[1]{$\blacktriangleright$~\textsc{var.~personal} \textsf{#1}{}}

\newcommand{\constructvarlab}{$\blacktriangleright$~\textsc{var.~construccional{ }}}

\newcommand{\constructvarof}[1]{$\blacktriangleright$~\textsc{var.~construccional~de} \textsf{#1}{}}

\newcommand{\archvarlab}{$\blacktriangleright$~\textsc{var.~arcaica{ }}}

\newcommand{\archvarof}[1]{$\blacktriangleright$~\textsc{var.~arcaica~de} \textsf{#1}{}}

\newcommand{\prepausallab}{$\blacktriangleright$~\textsc{forma prepausal{ }}}
 
\newcommand{\prepausalof}[1]{$\blacktriangleright$~\textsc{var.~prepausal~de} \textsf{#1}{}}
 
\newcommand{\affectvarlab}{$\blacktriangleright$~\textsc{var.~afectiva{ }}} 

\newcommand{\affectvarof}[1]{$\blacktriangleright$~\textsc{var.~afectiva~de} \textsf{#1}{}}

\newcommand{\euphvarlab}{$\blacktriangleright$~\textsc{var.~eufemística{ }}} 

\newcommand{\euphvarof}[1]{$\blacktriangleright$~\textsc{var.~eufemística~de} \textsf{#1}{}}

\newcommand{\playvarlab}{$\blacktriangleright$~\textsc{var.~bromista{ }}} 
 
\newcommand{\playvarof}[1]{$\blacktriangleright$~\textsc{var.~bromista~de} \textsf{#1}{}}
  
\newcommand{\nicknamelab}{$\blacktriangleright$~\textsc{apodo~}} 

\newcommand{\nicknameof}[1]{$\blacktriangleright$~\textsc{apodo de} \textsf{#1}{}}
 
\newcommand{\quantvarlab}{$\blacktriangleright$~\textsc{var. cuantitativa~}}

\newcommand{\quantvarof}[1]{$\blacktriangleright$~\textsc{var.~cuantitativa~de} \textsf{#1}{}} 
 
\newcommand{\irregthirdposs}[1]{$\blacktriangleright$~\textsc{forma poseída de tercera persona irreg.} \textsf{#1}{}} 
%\newcommand{\irregthirdposs}[1]{\textit{3.poss.} \textbf{#1}{}}
%\newcommand{\irregthirdposs}{\textsc{Forma poseída de tercera persona irreg.}{ }}

 \newcommand{\irregfirstposs}[1]{$\blacktriangleright$~\textit{forma poseída de primera persona irreg.} \textbf{#1}{}}

\newcommand{\allomorphlab}{$\blacktriangleright$~\textsc{allomorfo(s)~}} 

\newcommand{\allomorphof}[1]{$\blacktriangleright$~\textsc{alomorfo~de} \textsf{#1}{}}  
 
\titleformat{\chapter}[display]
{\bfseries\sffamily\LARGE}
{}{-4ex}{#1} 

\titleformat{\section}
{\vspace{0.5em}\bfseries\sffamily\large}
{\large{\thesection}}{1em}{{#1}}[\vspace{0.5em}]


\pagestyle{fancy}
\renewcommand{\headrulewidth}{0pt}
\lhead{\textsf{\rightmark}}
\rhead{\textsf{\leftmark}}
\chead{}
%\lfoot{}
\cfoot{}
%\rfoot{}
\fancyfoot[RO]{{\textit{Dictionary Ende{\textendash}English}} \qquad \textpipe \qquad {\small{\thepage}}\\}
\fancyfoot[LE]{{\small{\thepage}} \qquad \textpipe \qquad  \textit{Dictionary Ende{\textendash}English}\\}

\fancypagestyle{plain}{
\fancyhf{}% clears all header and footer fields
\renewcommand{\headrulewidth}{0pt}
%\lhead{\textbf{\rightmark}}
%\rhead{\textbf{\leftmark}}
%\chead{}
%\lfoot{}
\cfoot{}
%\rfoot{}
\fancyfoot[RO]{{\textit{dictionary Ende–English}} \qquad \textpipe \qquad {\small{\thepage}}\\}
\fancyfoot[LE]{{\small{\thepage}} \qquad \textpipe \qquad  \textit{Ende to English}\\}  
}

%THINGS to find-replace and hand-edit
% p.~ej. > p.~ej.
% conj. > subord.
% ~$\parallel$~ } >>> }
% (medio) >>> } \textsc{(med.)
%} (activo) >>> } \textsc{(act.)
% comment out "extra" k=, kw=, p=, pɨ=
% merge U and UU

\begin{document}

\setcounter{page}{1}

%\newpage
%\thispagestyle{plain}
%\mbox{}
%\newpage

\part*{{dictionary\\ Ende–English}}


\twocolumn[]
\raggedright

\setnoclub
\setnowidow

\entry{1940}{\headword{1940}
  \variantof{Dialectal Variant of \vartext{English loanword}}}

\entry{5}{\headword{5}
  \pos{Numeral}  \sense{
    \definition{5}}}

\entry{00}{\headword{00}
  \pos{Numeral}  \sense{
    \definition{zero}}}

\entry{3}{\headword{3}
  \pos{Numeral}  \sense{
    \definition{3}}}

\entry{300}{\headword{300}
  \pos{Numeral}  \sense{
    \definition{three hundred}}}

\entry{12}{\headword{12}
  \pos{Numeral}  \sense{
    \definition{12}}}

\entry{8}{\headword{8}
  \pos{Numeral}  \sense{
    \definition{8}}}

\entry{1}{\headword{1}
  \pos{Numeral}  \sense{
    \definition{1}}}

\entry{2}{\headword{2}
  \pos{Numeral}  \sense{
    \definition{2}}}

\entry{11}{\headword{11}
  \pos{Numeral}  \sense{
    \definition{11}}}

\entry{10}{\headword{10}
  \pos{Numeral}  \sense{
    \definition{10}}}

\entry{6}{\headword{6}
  \pos{Numeral}  \sense{
    \definition{6}}}

\entry{9}{\headword{9}
  \pos{Numeral}  \sense{
    \definition{9}}}

\entry{4}{\headword{4}
  \pos{Numeral}  \sense{
    \definition{4}}}

\entry{7}{\headword{7}
  \pos{Numeral}  \sense{
    \definition{7}}}

\entry{=0}{\headword{=0}
  \variantof{freevarlab of \vartext{=n}}}

\chapter{A}


\entry{a}{\headword{a}  \pronnote{a}
  \pos{Coordinating connective}  \sense{
    \definition{and}    \example{
      \exsen{Eso a mer toto.}      \extran{Thank you and good evening.}}}}

\entry{a-}{\headword{a-}  \pronnote{a}
  \pos{Verb}  \sense{
    \definition{root.ext.tr}}}

\entry{a-}{\headword{a-}  \pronnote{a}
  \pos{Verb}  \sense{\textbf{1.} 
    \definition{rec.intr}}  \sense{\textbf{2.} 
    \definition{fut.intr.2S}}}

\entry{ä-}{\headword{ä-}  \pronnote{ə}
  \pos{Verb}  \sense{
    \definition{marks a third person non-dual patient}}}

\entry{=a}{\headword{=a}  \pronnote{a}
  \pos{Discourse clitic}  \sense{
    \definition{vocative particle}}}

\entry{ä-}{\headword{ä-}  \pronnote{ə}
  \pos{Intransitive verb}  \sense{
    \definition{root.ext.intr}}}

\entry{aba}{\headword{aba}  \pronnote{aba}
  \pos{Interjection}  \sense{
    \definition{go (command given to animals)}}
  \variants{\Fast Speech Variant \vartext{ba}}
  \variants{\Unspecified Variant \vartext{abwa}}}

\entry{=aba}{\headword{=aba}  \pronnote{aba}
  \pos{Pronominal enclitic}  \sense{
    \definition{clitic form of oba}    \example{
      \exsen{Grace a Kate dag mänmän aba peyang.}      \extran{Grace and Kate are with the girls.}}    \example{
      \exsen{Ende llayaba ttoenttoen dan.}      \extran{This is the Ende people's story.}}    \example{
      \exsen{Do gabmaeyaba ttängäm me dan.}      \extran{She is there, where the white people are.}}    \example{
      \exsen{Ngämo mang ba bin a ada, Eric, Francis, Bewag.}      \extran{These are my brothers' names: Eric, Francis, and Bewag.}}}}

\entry{abade}{\headword{abade}  \pronnote{abade}
  \pos{Modifier}  \sense{
    \definition{future, later, upcoming, impending}    \example{
      \exsen{Yesu obo abade kuddäll e a kame ttam e gomballägän.}      \extran{Jesus predicted his impending death and resurrection.}}}}

\entry{=abaene}{\headword{=abaene}  \pronnote{abaene}
  \pos{Pronominal enclitic}  \sense{
    \definition{clitic form of obaene}    \example{
      \exsen{llayabaene tot}      \extran{people's trash}}}}

\entry{Abagigima}{\headword{Abagigima}
  \pos{Proper noun}  \sense{
    \definition{personal name}}}

\entry{abal}{\headword{abal}  \pronnote{abal}
  \pos{Modifier}  \sense{\textbf{1.} 
    \definition{pure}}  \sense{\textbf{2.} 
    \definition{very}    \example{
      \exsen{mer abal nag}      \extran{very good friends}}}  \sense{\textbf{3.} 
    \definition{exactly, just}    \example{
      \exsen{sisri abal}      \extran{right now}}}
  \variants{\Dialectal Variant \vartext{aball}}}

\entry{aball}{\headword{aball}
  \variantof{Dialectal Variant of \vartext{abal}}  \pronnote{abaɽ}}

\entry{Abam}{\headword{Abam}  \note{Place names}  \pronnote{abam}
  \pos{Proper noun}  \sense{
    \definition{Abam (Wipi-speaking village in Oriomo-Bituri LLG; GPS: -8.926607 143.190246)}    \note{Place names}    \example{
}    \example{
}    \example{
}}}

\entry{abe}{\headword{abe}
  \pos{Coordinating connective}  \sense{
    \definition{but}}}

\entry{Abeam}{\headword{Abeam}  \pronnote{abeam}
  \pos{Proper noun}  \sense{
    \definition{male personal name}}}

\entry{Abere}{\headword{Abere}
  \pos{Proper noun}  \sense{
    \definition{female personal name}}}

\entry{Aberegerem}{\headword{Aberegerem}  \pronnote{abereɡerem}
  \pos{Proper noun}  \sense{
    \definition{Aberagerema (in Kiwai Rural LLG)}    \example{
}    \example{
}    \example{
}}}

\entry{Abigail}{\headword{Abigail}  \pronnote{abeɡail}
  \pos{Proper noun}  \sense{
    \definition{female personal name}}}

\entry{=abira}{\headword{=abira}  \pronnote{abira}
  \pos{Pronominal enclitic}  \sense{
    \definition{clitic form of ubira}    \example{
      \exsen{Bogo llɨg kälekäle abira bandra llɨtɨt e gongkamalle.}      \extran{She started to sing to the little kids.}}}}

\entry{abo}{\headword{abo}  \pronnote{abo}
  \pos{Adverb}  \sense{\textbf{1.} 
    \definition{then, afterwards}    \example{
      \exsen{Ngäna abo gongos ma we.}      \extran{I then returned home.}}}  \sense{\textbf{2.} 
    \definition{must, necessative mood}    \example{
      \exsen{Abo bongo ttongo ma sisor de nogo.}      \extran{You must build a new house.}}}}

\entry{Abom}{\headword{Abom}
  \pos{Proper noun}  \sense{
    \definition{Abom (toponym)}}}

\entry{abor}{\headword{abor}  \pronnote{abor}
  \pos{Noun}  \sense{
    \definition{sago beater}    \example{
      \exsen{Ngäna sana pällkapällke abor alle nuebnegan.}      \extran{I beat the little sago pieces with the sago beater.}}    \example{
}    \example{
}}}

\entry{about}{\headword{about}
  \variantof{Dialectal Variant of \vartext{English loanword}}}

\entry{abwa}{\headword{abwa}
  \variantof{Unspecified Variant of \vartext{aba}}}

\entry{ada}{\headword{ada}  \pronnote{ada}
  \pos{Manner demonstrative}  \sense{\textbf{1.} 
    \definition{like this, thus, so}    \example{
      \exsen{Ada daeya.}      \extran{So it was.}}    \example{
      \exsen{Obo mälläng ik a ada ulleulle dageya.}      \extran{His nostrils were this big.}}}  \sense{\textbf{2.} 
    \definition{quotative particle (for quoted/direct speech)}    \example{
      \exsen{Ngänawa "Ngämlle gulag e" ada däga.}      \extran{I said, "Come with me."}}    \example{
      \exsen{Bogo ada, "Ngämo pätt a ttattllong agan."}      \extran{He said, "My body has already gotten sore."}}}  \sense{\textbf{3.} 
    \definition{that}    \example{
      \exsen{Ngämo umllang dan ada bongo Grace bo nag dan.}      \extran{I know that you are Grace's friend.}}}}

\entry{adade}{\headword{adade}  \pronnote{adade}
  \pos{Adverb}  \sense{
    \definition{~ ada}}}

\entry{adako}{\headword{adako}  \pronnote{adako}
  \pos{Adverb}  \sense{\textbf{1.} 
    \definition{~ ada}}  \sense{\textbf{2.} 
    \definition{~ ako}    \example{
      \exsen{Adako bongo mɨnyi nalle ma we.}      \extran{Agree that you will go to the house.}}}}

\entry{adalla}{\headword{adalla}
  \variantof{Dialectal Variant of \vartext{atka}}  \pronnote{adaɽa}}

\entry{Adam}{\headword{Adam}
  \pos{Proper noun}  \sense{
}}

\entry{adame}{\headword{adame}  \pronnote{adame}
  \pos{Adverb}  \sense{\textbf{1.} 
    \definition{this is why, therefore}    \example{
      \exsen{Adame ddobag a gudae ngällbaenen eran yunu atta bäne pate llɨtaem eran.}      \extran{This is why some wake up early in the morning from sleep and come to you.}}}  \sense{\textbf{2.} 
    \definition{at this moment}    \example{
      \exsen{Angde Yesu eka gognän, Zudas adame gongttägän lla gul aba peyang.}      \extran{Right when Jesus was speaking, Judas arrived with a throng.}}}}

\entry{Adasha}{\headword{Adasha}  \pronnote{adasha}
  \pos{Proper noun}  \sense{
    \definition{male personal name}}}

\entry{adatt}{\headword{adatt}
  \variantof{Fast Speech Variant of \vartext{adawatta}}}

\entry{adawae}{\headword{adawae}  \pronnote{adawaj}
  \pos{Adverb}  \sense{
    \definition{the same way}    \example{
      \exsen{Ngäna däbe täl de deyanykoe. Adawae gog a dädme mäse ada deyanykoe.}      \extran{I pulled that bamboo. I kept doing the same, trying to pull so.}}}}

\entry{adawalle}{\headword{adawalle}  \pronnote{adawaɽe}
  \pos{Adverb}  \sense{
    \definition{then, so}    \example{
      \exsen{Ddia da mängalae abal gurlewän iba ttängäm me. Adawalle de iba ttängäm a ddia peyang gogän.}      \extran{The deer reproduced very quickly in our (incl.) village. So our village came to be populated with deer.}}}}

\entry{adawatta}{\headword{adawatta}  \pronnote{adawaʈ͡ʂa}
  \pos{Subordinating connective}  \sense{\textbf{1.} 
    \definition{because (introduces a cause or reason)}    \example{
}    \example{
      \exsen{Diba toto we ngäma pemli aba peyang ngämi mer duwem de dägaeya, adawatta ddia da däbe ddobae gi daeya.}      \extran{That evening, we (excl.) ate well in our family because that deer was very fatty.}}    \example{
      \exsen{Mamal gänyeriballe yaupe adawatta da mo da bäne peyang bottkamän.}      \extran{Jump quickly to this side: otherwise, the bridge might break with you.}}}  \sense{\textbf{2.} 
    \definition{why, therefore (marks a result or consequence)}    \example{
      \exsen{Bodog ada ka Satade dan. Bogo adawatta bel de ekameny nägagan.}      \extran{Bodog thought it was Saturday. That's why he didn't ring the bell.}}}
  \variants{\Fast Speech Variant \vartext{atta, adatt}}}

\entry{adawede}{\headword{adawede}  \pronnote{adawede}
  \pos{Subordinating connective}  \sense{
    \definition{so that}    \example{
      \exsen{Ngäna Ende eka de tameny anggan adawede ubi ngämo kokok erag Ende eka de umllang bognegnän ade.}      \extran{I teach Ende so that my grandchildren will also learn to speak Ende.}}}}

\entry{ade}{\headword{ade}  \pronnote{ade}
  \pos{Adverb}  \sense{
    \definition{also}    \example{
      \exsen{Ngäna Ingglis ekalle eka allan, a iba eka walle ade täräp me.}      \extran{I speak in English and also in our (excl.) language sometimes.}}}}

\entry{Adi}{\headword{Adi}  \pronnote{adi}
  \pos{Proper noun}  \sense{
    \definition{God}    \example{
      \exsen{Bongo era Adi bo llɨg dan.}      \extran{You are the son of God.}}}}

\entry{adingoll}{\headword{adingoll}  \pronnote{adiŋoɽ}
  \pos{Adverb}  \sense{
    \definition{like this}    \example{
      \exsen{Ende eka de ddob ngämo patme ttättle amallo ada, "Adingoll näpany."}      \extran{They correct some of my Ende words, saying, "Say it like this."}}}}

\entry{adräl}{\headword{adräl}  \note{Trees}  \pronnote{adrəl}
  \pos{Noun}  \sense{
    \definition{type of tree}    \note{Trees}    \example{
}    \example{
}    \example{
}}}

\entry{Adu}{\headword{Adu}  \pronnote{adu}
  \pos{Proper noun}  \sense{
    \definition{female personal name}}}

\entry{adult}{\headword{adult}
  \variantof{Dialectal Variant of \vartext{English loanword}}}

\entry{aduwi}{\headword{aduwi}  \note{Trees}  \pronnote{aduwi}
  \pos{Noun}  \sense{
    \definition{type of tree}    \note{Trees}    \example{
}    \example{
}    \example{
}}}

\entry{=ae}{\headword{=ae}  \pronnote{ae}
  \pos{Clitic}  \sense{
    \definition{adverbializing clitic}    \example{
      \exsen{Ngäna dinduangmeae gogne llättmenyae.}      \extran{I was running nonstop.}}}}

\entry{=ae}{\headword{=ae}  \pronnote{ae}
  \pos{Clitic}  \sense{
    \definition{restrictive clitic; only}    \example{
      \exsen{Ngäna Upiara me ae skul gogne.}      \extran{I was schooling in Upiara only.}}}}

\entry{=ae}{\headword{=ae}  \pronnote{ae}
  \pos{Nominal enclitic}  \sense{
    \definition{distal vocative clitic}    \example{
      \exsen{Llamda ae!}      \extran{Hey old man!}}}}

\entry{ae}{\headword{ae}  \pronnote{ae}
  \pos{Interjection}  \sense{
    \definition{ah}}}

\entry{aeb}{\headword{aeb}  \pronnote{aeb}
  \pos{Verb}  \sense{
    \definition{to deny}}}

\entry{aeb}{\headword{aeb}  \lexeme{aeb}  \pronnote{aeb}
  \pos{Noun}  \sense{
    \definition{black-billed/yellow-legged brushturkey}    \scientificname{Talegalla fuscirostris}    \example{
}}}

\entry{aebaeb}{\headword{aebaeb}  \note{Trees}  \pronnote{aebaeb}
  \pos{Noun}  \sense{
    \definition{type of tree}    \note{Trees}    \example{
}    \example{
}    \example{
}}}

\entry{=aebe}{\headword{=aebe}  \pronnote{aebe}
  \pos{Nominal enclitic}  \sense{
    \definition{restrictive clitic; only}    \example{
      \exsen{Kollba da tämamae dadrowän a be kottllam aebe ttam dagirnän.}      \extran{The fish all died and only the turtle lived on.}}}}

\entry{aenen}{\headword{aenen}  \pronnote{ainen}
  \pos{Copular verb}  \sense{
    \definition{copular form of aya (present singular form)}    \example{
      \exsen{Bäne bin a ainen?}      \extran{What's your name?}}    \example{
      \exsen{Ge Sam ainen.}      \extran{This is Sam.}}}
  \variants{\SpellingVariant \vartext{ainen}}}

\entry{aeneya}{\headword{aeneya}  \pronnote{aineja}
  \pos{Copular verb}  \sense{
    \definition{past singular form of aenen}}
  \variants{\SpellingVariant \vartext{aineya}}}

\entry{aengap}{\headword{aengap}  \lexeme{aengap}  \pronnote{aeŋap}
  \pos{Noun}  \sense{
    \definition{type of big yam with a white interior and no thorns}}}

\entry{aetruaetru}{\headword{aetruaetru}  \note{Trees}  \pronnote{aetruaetru}
  \pos{Noun}  \sense{
    \definition{type of tree with edible blue fruits and white flowers, found in the bush and along big creeks}    \note{Trees}    \example{
}    \example{
}    \example{
}}}

\entry{=aeya}{\headword{=aeya}
  \pos{Verbal enclitic}  \sense{
    \definition{cop.pst.sgS}}
  \variants{\Unspecified Variant \vartext{=ye}}
  \variants{\freevarlabs \vartext{=weya, =eya, =aya, =e}}
  \variants{\Dialectal Variant \vartext{=aeyag}}}

\entry{aeya}{\headword{aeya}
  \variantof{SpellingVariant of \vartext{aya}}  \pronnote{aija}}

\entry{=aeyag}{\headword{=aeyag}
  \variantof{Dialectal Variant of \vartext{=aeya}}}

\entry{ag}{\headword{ag}  \note{Parts of the day}  \pronnote{aɡ}
  \pos{Noun}  \sense{
    \definition{morning (approx. 5 AM–11 AM)}    \note{Parts of the day}    \example{
}    \example{
      \exsen{Ag me inu da era ddobae llokttang dan.}      \extran{Sleeping in the morning is really difficult.}}    \example{
}}}

\entry{ag duwem}{\headword{ag duwem}  \note{Parts of the day}  \pronnote{aɡ duwem}
  \pos{Noun}  \sense{
    \definition{breakfast}    \note{Parts of the day}    \example{
}    \example{
}    \example{
}}}

\entry{Agäb}{\headword{Agäb}  \pronnote{'Agäb}
  \pos{Proper noun}  \sense{
    \definition{Agob language (Pahoturi River language)}    \example{
}    \example{
}    \example{
}}
  \variants{\Unspecified Variant \vartext{Agob}}}

\entry{Agan}{\headword{Agan}
  \pos{Proper noun}  \sense{
    \definition{Agan (toponym)}}}

\entry{agde}{\headword{agde}
  \variantof{freevarlab of \vartext{Agob loanword}}}

\entry{agdea}{\headword{agdea}
  \variantof{freevarlab of \vartext{Agob loanword}}}

\entry{Agob}{\headword{Agob}
  \variantof{Unspecified Variant of \vartext{Agäb}}  \pronnote{aɡəb}}

\entry{Agob loanword}{\headword{Agob loanword}
  \pos{}  \sense{
}
  \variants{\freevarlabs \vartext{agde, agdea}}}

\entry{ai}{\headword{ai}  \pronnote{ai}
  \pos{Modifier}  \sense{\textbf{1.} 
    \definition{good}    \example{
      \exsen{ai abal ttängäm}      \extran{a very good place}}}  \sense{\textbf{2.} 
    \definition{plenty, abundant}    \example{
      \exsen{ai kollba, ai wayati}      \extran{plenty of fish, lots of watermelons}}}  \sense{\textbf{3.} 
    \definition{modal adjective (expresses permission, possibility, obligation)}    \example{
      \exsen{Ngäna ai dan gänyme bina peyang bodämen?}      \extran{May I sit here with you (pl.)?}}    \example{
      \exsen{Wel a era za gagällang dan a ai dan lla de ade bäbäddän.}      \extran{The wind is a bad thing that can also kill people.}}    \example{
      \exsen{Bongo ai da ine de näna.}      \extran{You should drink water.}}}  \sense{\textbf{4.} 
    \definition{very}    \example{
      \exsen{ai mängallang mälla}      \extran{a very strong woman}}}
  \variants{\SpellingVariant \vartext{aii}}}

\entry{ai}{\headword{ai}  \pronnote{ai}
  \pos{Interjection}  \sense{\textbf{1.} 
    \definition{ah}}  \sense{\textbf{2.} 
}}

\entry{ai skul}{\headword{ai skul}  \note{School}  \pronnote{ai skul}
  \pos{Noun}  \sense{
    \definition{secondary school, high school}    \note{School}    \example{
      \exsen{Ngäna ai skul i ddone dall.}      \extran{I didn't go to secondary school.}}    \example{
}    \example{
}}}

\entry{aiai}{\headword{aiai}
  \pos{Adverb}  \sense{
    \definition{properly, well}    \example{
      \exsen{Eka de ddob malla aiai panynen amallo.}      \extran{Some don't speak the language properly.}}}
  \variants{\SpellingVariant \vartext{aiyai}}}

\entry{aida}{\headword{aida}
  \variantof{freevarlab of \vartext{aidan}}  \pronnote{aida}}

\entry{aidan}{\headword{aidan}  \pronnote{aidan}
  \pos{Interjection}  \sense{
    \definition{alright, okay}    \example{
      \exsen{Piasorosoro ada eka gogon, "Ibi goeg tatäräp e." Mälla da ada gogon, "Aida, ibi beyareya."}      \extran{Piasorosoro said, "Let's go to the garden to cut down trees." His wife said, "Okay, let's go."}}}
  \variants{\freevarlab \vartext{aida}}}

\entry{Ailin}{\headword{Ailin}
  \pos{Proper noun}  \sense{
    \definition{female personal name}}}

\entry{Ainduru}{\headword{Ainduru}
  \variantof{SpellingVariant of \vartext{Andrew}}}

\entry{ainen}{\headword{ainen}
  \variantof{SpellingVariant of \vartext{aenen}}  \pronnote{ainen}
  \variants{\SpellingVariant \vartext{ainin}}}

\entry{aineya}{\headword{aineya}
  \variantof{SpellingVariant of \vartext{aeneya}}}

\entry{ainin}{\headword{ainin}
  \variantof{SpellingVariant of \vartext{ainen}}  \pronnote{ainin}}

\entry{Ainor}{\headword{Ainor}  \pronnote{ainor}
  \pos{Proper noun}  \sense{\textbf{1.} 
    \definition{Ainor (toponym)}}  \sense{\textbf{2.} 
    \definition{dog name}}}

\entry{aitar}{\headword{aitar}  \note{Trees}  \pronnote{aitar}
  \pos{Noun}  \sense{
    \definition{type of palm tree that grows along creeks with pools; shoots and soft trunk are edible; similar to dumar}    \note{Trees}    \example{
}    \example{
}    \example{
}}}

\entry{Aituru}{\headword{Aituru}  \pronnote{aituru}
  \pos{Proper noun}  \sense{
    \definition{female personal name}}}

\entry{aiyai}{\headword{aiyai}
  \variantof{SpellingVariant of \vartext{aiai}}}

\entry{aii}{\headword{aii}
  \variantof{SpellingVariant of \vartext{ai}}  \pronnote{aii}}

\entry{äk}{\headword{äk}  \note{Conflict/war}  \pronnote{ək}
  \pos{Verb}  \sense{
    \definition{to beat the wife if she doesn't do what the husband says}    \note{Conflict/war}    \example{
}    \example{
}    \example{
}}}

\entry{Aketa}{\headword{Aketa}
  \pos{Proper noun}  \sense{
    \definition{Aketa (in Gogodala Rural LLG)}}}

\entry{ako}{\headword{ako}  \pronnote{ako}
  \pos{Discourse particle}  \sense{\textbf{1.} 
    \definition{also, too}    \example{
      \exsen{Mägda män de kaptte namättan. Mägda ako kaptte amättan.}      \extran{Mother dressed her daughter. Mother also got dressed.}}}  \sense{\textbf{2.} 
    \definition{again}    \example{
      \exsen{Täre imnealle, ttongo pazi ddägatt auma ma ddaddällɨg e lla da täre de ako mɨnyi dägagallo.}      \extran{One year after the feast and destroying the grave, people will have a feast again.}}    \example{
}}  \sense{\textbf{3.} 
    \definition{then}    \example{
      \exsen{Ge lla da bandra da nällɨtan, ako pomer alle nällɨtan.}      \extran{This man sang a song and then whistled it.}}}  \sense{\textbf{4.} 
    \definition{other}    \example{
      \exsen{ako ttongo llo}      \extran{another tree}}    \example{
      \exsen{Ddobagabira gae kaekep a ako llowamang dan.}      \extran{Some people chew gae but others don't like to.}}}
  \variants{\Fast Speech Variant \vartext{ko}}}

\entry{=ako}{\headword{=ako}
  \pos{Nominal enclitic}  \sense{
}}

\entry{Al}{\headword{Al}  \pronnote{al}
  \pos{Proper noun}  \sense{
    \definition{female personal name}}}

\entry{aläm}{\headword{aläm}  \pronnote{aləm}
  \pos{Intransitive coverb}  \sense{\textbf{1.} 
    \definition{to be wary}    \example{
      \exsen{Abo bongo Daru me ai aläm agne.}      \extran{You must be wary in Daru.}}}  \sense{\textbf{2.} 
    \definition{to warn, caution}    \example{
      \exsen{Yesu ubim ada aläm dägnegän sanawang eka walle.}      \extran{Jesus warned them with a parable.}}}}

\entry{Aleks}{\headword{Aleks}
  \variantof{SpellingVariant of \vartext{Alex}}}

\entry{Alex}{\headword{Alex}  \pronnote{alɛks}
  \pos{Proper noun}  \sense{
    \definition{male personal name}}
  \variants{\SpellingVariant \vartext{Aleks}}}

\entry{ali}{\headword{ali}  \pronnote{ali}
  \pos{Noun}  \sense{
    \definition{conch shell}}}

\entry{Ali}{\headword{Ali}  \pronnote{ali}
  \pos{Proper noun}  \sense{
    \definition{place name}    \example{
}    \example{
}    \example{
}}}

\entry{alla}{\headword{alla}  \pronnote{aɽa}
  \pos{Adverb}  \sense{
    \definition{how, what}    \example{
      \exsen{Alla ulle ine dan?}      \extran{How much water is there?}}    \example{
      \exsen{Ge lla da obo lläg komlla de alla tameny yaran.}      \extran{This man is teaching his two kids how.}}    \example{
      \exsen{Ibi alla täräp me ibi ge nge de ibeny e dan?}      \extran{What time are we (incl.) planting this coconut?}}}}

\entry{=alla}{\headword{=alla}  \pronnote{aɽa}
  \pos{Auxiliary verb}  \sense{
    \definition{aux.prs.1|2duS}}}

\entry{alla ingoll}{\headword{alla ingoll}
  \variantof{freevarlab of \vartext{alla ingollang}}}

\entry{alla ingollang}{\headword{alla ingollang}  \note{Question words}  \pronnote{aɽaiŋoɽaŋ}
  \pos{Adverb}  \sense{
    \definition{how}    \note{Question words}    \example{
      \exsen{Ngäna bam alla ingollang bangmingg?}      \extran{How can I help you?}}    \example{
      \exsen{Ge alla ingollang ngäma lla da gurlewän.}      \extran{This is how our (excl.) population grew.}}}
  \variants{\freevarlab \vartext{alla ingoll}}
  \variants{\Fast Speech Variant \vartext{allingoll, allingollang}}
  \variants{\SpellingVariant \vartext{allaingollang}}}

\entry{allaingollang}{\headword{allaingollang}
  \variantof{SpellingVariant of \vartext{alla ingollang}}}

\entry{Allambun}{\headword{Allambun}  \note{Place names around Limol}  \pronnote{aɽambun}
  \pos{Proper noun}  \sense{
    \definition{Allambun (camping place)}    \note{Place names around Limol}    \example{
}    \example{
}    \example{
}}}

\entry{=allan}{\headword{=allan}  \pronnote{aɽan}
  \pos{Auxiliary verb}  \sense{\textbf{1.} 
    \definition{aux.prs.1|3sgS}}  \sense{\textbf{2.} 
}  \sense{\textbf{3.} 
}  \sense{\textbf{4.} 
}}

\entry{alläp}{\headword{alläp}  \pronnote{aɽəp}
  \pos{Noun}  \sense{\textbf{1.} 
    \definition{kundu drum}    \example{
}    \example{
}    \example{
}}  \sense{\textbf{2.} 
    \definition{guitar}}
  \variants{\Dialectal Variant \vartext{allap}}}

\entry{allap}{\headword{allap}
  \variantof{Dialectal Variant of \vartext{alläp}}  \pronnote{aɽap}}

\entry{Alläpma}{\headword{Alläpma}  \pronnote{aɽəpma}
  \pos{Proper noun}  \sense{
    \definition{Alläpma (toponym)}    \example{
}    \example{
}    \example{
}}}

\entry{=alle}{\headword{=alle}  \pronnote{aɽe}
  \pos{Nominal enclitic}  \sense{
    \definition{ablative case clitic; from}    \example{
      \exsen{Ag alle do yäbäd klaklema.}      \extran{From morning to sunset.}}}}

\entry{=alle}{\headword{=alle}  \pronnote{aɽe}
  \pos{Auxiliary verb}  \sense{
    \definition{aux.prs.2sgS}}}

\entry{=alle}{\headword{=alle}  \pronnote{aɽe}
  \pos{Nominal enclitic}  \sense{\textbf{1.} 
    \definition{instrumental case clitic; with (using)}    \example{
      \exsen{Biye de bongo ai dan ibik alle beni.}      \extran{You can plant taro with the ibik stick.}}}  \sense{\textbf{2.} 
    \definition{comitative case clitic; with (together)}    \example{
      \exsen{Nensi bi nane alle känaebag mɨnyi tudi ma beyareyo.}      \extran{Nancy will go fishing tomorrow with her aunt.}}}
  \variants{\Dialectal Variant \vartext{=ane, =wane}}}

\entry{allingoll}{\headword{allingoll}
  \variantof{Fast Speech Variant of \vartext{alla ingollang}}  \pronnote{aɽiŋoɽ}}

\entry{allingollang}{\headword{allingollang}
  \variantof{Fast Speech Variant of \vartext{alla ingollang}}}

\entry{allka}{\headword{allka}
  \pos{Transitive/Intransitive coverb}  \sense{
    \definition{to shout (at)}    \example{
      \exsen{Ge lla da allka allan, obom allka eran.}      \extran{This man is shouting; he's shouting at her.}}}}

\entry{allko}{\headword{allko}  \lexeme{allko}  \pronnote{aɽko}
  \pos{Noun}  \sense{
    \definition{fly (insect)}    \example{
      \exsen{Allko da yuwog abal dag.}      \extran{There are many flies.}}}}

\entry{allko kukut}{\headword{allko kukut}  \pronnote{aɽko kukut}
  \pos{Noun}  \sense{
    \definition{blue fly}    \example{
}    \example{
}    \example{
}}}

\entry{allko wallägnewallägnen}{\headword{allko wallägnewallägnen}  \note{Sago palms}  \pronnote{aɽko waɽəɡnewaɽəɡnen}
  \pos{Noun}  \sense{
    \definition{type of sago}    \note{Sago palms}}}

\entry{=allo}{\headword{=allo}  \pronnote{aɽo}
  \pos{Auxiliary verb}  \sense{
    \definition{aux.prs.3duS}}}

\entry{Alofa}{\headword{Alofa}
  \variantof{SpellingVariant of \vartext{Alopa}}  \pronnote{alofa}}

\entry{Alopa}{\headword{Alopa}
  \pos{Proper noun}  \sense{
    \definition{female personal name}}
  \variants{\SpellingVariant \vartext{Alofa}}}

\entry{Alphones}{\headword{Alphones}  \pronnote{alfons}
  \pos{Proper noun}  \sense{
    \definition{male personal name}}
  \variants{\SpellingVariant \vartext{Alpons}}}

\entry{Alpons}{\headword{Alpons}
  \variantof{SpellingVariant of \vartext{Alphones}}}

\entry{am}{\headword{am}  \pronnote{am}
  \pos{Noun}  \sense{
    \definition{internode (section of bamboo or sugarcane, separated by nodes)}    \example{
      \exsen{Angde ttongo am pättkäp de nätäräpän.}      \extran{After one internode, cut the node.}}}}

\entry{ama}{\headword{ama}  \pronnote{ama}
  \pos{Interrogative pronoun}  \sense{\textbf{1.} 
    \definition{possessive form of ami}}  \sense{\textbf{2.} 
    \definition{possessive form of ami}}  \sense{\textbf{3.} 
    \definition{possessive form of ami}    \example{
      \exsen{Mer dan ada yuwog abal lla da wa mälla da gänyag, ama moko dan mällayabira ngämiminy i.}      \extran{It is good that there are so many men and women here who want (lit. whose desire it is) to help women.}}}}

\entry{ama}{\headword{ama}
  \pos{Noun}  \sense{
    \definition{hammer}}}

\entry{ama bakmall}{\headword{ama bakmall}
  \pos{Interrogative pronoun}  \sense{\textbf{1.} 
    \definition{comitative of ami}    \example{
}}  \sense{\textbf{2.} 
    \definition{comitative form of ami}}  \sense{\textbf{3.} 
    \definition{comitative of ami}    \example{
      \exsen{Ge lla da ngäna ama bakmall skul att dan.}      \extran{These people are who I went to school with.}}}}

\entry{ama kämall}{\headword{ama kämall}
  \pos{Interrogative pronoun}  \sense{\textbf{1.} 
    \definition{comitative form of ami}    \example{
      \exsen{Bibi amakmall panypeny eralla?}      \extran{With whom do you all speak?}}}  \sense{\textbf{2.} 
    \definition{comitative form of ami}    \example{
      \exsen{Ngämi amakämall skul att dag, gänyme ddone dag.}      \extran{Whomever we (excl.) went to school with, they're not here.}}}  \sense{\textbf{3.} 
    \definition{comitative form of ami}}
  \variants{\Fast Speech Variant \vartext{amakmall}}}

\entry{Amadu}{\headword{Amadu}  \pronnote{amadu}
  \pos{Proper noun}  \sense{
    \definition{male personal name}}}

\entry{amakmall}{\headword{amakmall}
  \variantof{Fast Speech Variant of \vartext{ama kämall}}}

\entry{=amalla}{\headword{=amalla}  \pronnote{amaɽa}
  \pos{Auxiliary verb}  \sense{\textbf{1.} 
    \definition{aux.prs.1|2plS|1|2nsgA>plP}}  \sense{\textbf{2.} 
    \definition{rec.1|2plS}}  \sense{\textbf{3.} 
}}

\entry{amalle}{\headword{amalle}  \pronnote{amaɽe}
  \pos{Interrogative pronoun}  \sense{
    \definition{instrumental-comitative form of aya}    \example{
      \exsen{Bibi amalle erang gusiya?}      \extran{Whom did you do an exchange with?}}}}

\entry{=amallo}{\headword{=amallo}  \pronnote{amaɽo}
  \pos{Auxiliary verb}  \sense{\textbf{1.} 
    \definition{aux.prs.3plS|3nsgA>plP}}  \sense{\textbf{2.} 
}}

\entry{amamär}{\headword{amamär}  \pronnote{amamər}
  \pos{Noun}  \sense{
    \definition{woven rope}    \example{
}    \example{
}    \example{
}}}

\entry{amamärang}{\headword{amamärang}  \pronnote{amaməraŋ}
  \pos{Intransitive coverb}  \sense{
    \definition{to echo}    \example{
      \exsen{Ngämo eka da amamärang agan.}      \extran{My voice echoed.}}}}

\entry{Amanda}{\headword{Amanda}  \pronnote{amanda}
  \pos{Proper noun}  \sense{
    \definition{female personal name}}}

\entry{amäramär}{\headword{amäramär}  \pronnote{aməramər}
  \pos{Noun}  \sense{
    \definition{braid}}}

\entry{amba}{\headword{amba}  \note{Trees}  \pronnote{amba}
  \pos{Noun}  \sense{
    \definition{type of tree found in the grassland}    \note{Trees}    \example{
}    \example{
}    \example{
}}}

\entry{ambag}{\headword{ambag}  \pronnote{ambaɡ}
  \pos{Intransitive coverb}  \sense{
    \definition{to cause trouble}    \example{
      \exsen{Bogo mɨnyi ambag bogon.}      \extran{He will cause trouble.}}}}

\entry{amene}{\headword{amene}  \pronnote{amene}
  \pos{Interrogative pronoun}  \sense{\textbf{1.} 
    \definition{ablative-possessive form of aya}    \example{
      \exsen{Bongo amene män de dɨllɨd?}      \extran{Whose daughter did you marry?}}    \example{
}}  \sense{\textbf{2.} 
    \definition{ablative-possessive form of aya}}  \sense{\textbf{3.} 
    \definition{ablative-possessive form of aya}}}

\entry{Amerika}{\headword{Amerika}
  \pos{Proper noun}  \sense{
    \definition{United States of America; America}}
  \variants{\Fast Speech Variant \vartext{Amrika}}}

\entry{ami}{\headword{ami}  \pronnote{ami}
  \pos{Interrogative pronoun}  \sense{\textbf{1.} 
    \definition{who (nonsingular interrogative pronoun, nominative form)}    \example{
      \exsen{Ubi ami dag?}      \extran{Who are they?}}}  \sense{\textbf{2.} 
    \definition{whoever, some people (nonsingular existential pronoun, nominative form)}    \example{
      \exsen{Da ami botngoenegnän, de ngäna mɨnyi any batngoenen.}      \extran{If some people are playing, I will play.}}}  \sense{\textbf{3.} 
    \definition{who (nonsingular relative pronoun, nominative form)}    \example{
      \exsen{Komlla nag a dadegwaeya ami gudae deyagirneyo walle ddage menae me.}      \extran{There were two friends who once lived on the side of a creek.}}}}

\entry{amig}{\headword{amig}  \pronnote{amiɡ}
  \pos{Copular verb}  \sense{
    \definition{contraction of ami dag}}}

\entry{amim}{\headword{amim}
  \pos{Interrogative pronoun}  \sense{\textbf{1.} 
    \definition{accusative form of ami}}  \sense{\textbf{2.} 
    \definition{accusative form of ami}}  \sense{\textbf{3.} 
    \definition{accusative form of ami}    \example{
      \exsen{kumuddäga pemli de ngäna amim nättaemnegan}      \extran{the three families whom I mentioned}}}}

\entry{amira}{\headword{amira}  \pronnote{amira}
  \pos{Interrogative pronoun}  \sense{\textbf{1.} 
    \definition{dative form of ami}}  \sense{\textbf{2.} 
    \definition{dative form of ami}}  \sense{\textbf{3.} 
    \definition{dative form of ami}}}

\entry{amiyeamiye}{\headword{amiyeamiye}  \pronnote{amijeamije}
  \pos{Adverb}  \sense{
    \definition{upwind, against the wind}}}

\entry{amlle}{\headword{amlle}  \pronnote{amɽe}
  \pos{Interrogative pronoun}  \sense{\textbf{1.} 
    \definition{dative form of aya}}  \sense{\textbf{2.} 
    \definition{dative form of aya}    \example{
      \exsen{Abo mudan abal amlle llätät a!}      \extran{You must not tell anyone!}}}  \sense{\textbf{3.} 
    \definition{dative form of aya}}}

\entry{amne}{\headword{amne}  \pronnote{amne}
  \pos{Locational}  \sense{
    \definition{center, middle}    \example{
      \exsen{Ge ttoen a era bem amne me de gongesän.}      \extran{This story happened in the middle of the sea.}}    \example{
      \exsen{Bogo ttängäm amne we gozenän.}      \extran{He went into the village center.}}    \example{
      \exsen{Bogo sisri amne pazi me dan.}      \extran{She's middle-aged now.}}}}

\entry{Amne kona}{\headword{Amne kona}  \note{Places around Limol}  \pronnote{amne kona}
  \pos{Proper noun}  \sense{
    \definition{Central corner (in Limol)}    \note{Places around Limol}    \example{
}    \example{
}    \example{
}}}

\entry{amo}{\headword{amo}  \pronnote{amo}
  \pos{Interrogative pronoun}  \sense{\textbf{1.} 
    \definition{possessive form of aya}    \example{
      \exsen{Amo sisor gall da däbe?}      \extran{Whose new canoe is that?}}}  \sense{\textbf{2.} 
    \definition{possessive form of aya}    \example{
      \exsen{Amo moko da ulle abal lla we, ede ai dan bogo kälsre abal bogon.}      \extran{Whoever wants (lit. whoever's desire it is) to be a very important person, he should be very trivial.}}}  \sense{\textbf{3.} 
    \definition{possessive form of aya}    \example{
      \exsen{ttongo mälla da amo män a gagäll anyke peyang daeya}      \extran{a woman whose daughter was possessed (lit. with bad spirit)}}}}

\entry{amom}{\headword{amom}  \note{Question words}  \pronnote{amom}
  \pos{Interrogative pronoun}  \sense{\textbf{1.} 
    \definition{accusative form of aya}    \note{Question words}    \example{
      \exsen{Bongo amom ikop nagagalle?}      \extran{Whom did you see?}}    \example{
}    \example{
}}  \sense{\textbf{2.} 
    \definition{accusative form of aya}    \note{Question words}    \example{
      \exsen{Yu da amom dättämän, bogo ma we dallän.}      \extran{Whomever the fire burnt went home.}}}  \sense{\textbf{3.} 
    \definition{accusative form of aya}    \note{Question words}    \example{
      \exsen{Lima bo umllang dan bongo mɨnyi amom känaebag ikop nägag.}      \extran{Lima knows whom you will see tomorrow.}}}}

\entry{Amrika}{\headword{Amrika}
  \variantof{Fast Speech Variant of \vartext{Amerika}}  \pronnote{amrika}}

\entry{amteamte}{\headword{amteamte}  \pronnote{amteamte}
  \pos{Noun}  \sense{\textbf{1.} 
    \definition{body part}}  \sense{\textbf{2.} 
    \definition{bow part}    \example{
}    \example{
}    \example{
}}}

\entry{amtet}{\headword{amtet}  \pronnote{amtet}
  \pos{Noun}  \sense{\textbf{1.} 
    \definition{breath}    \example{
      \exsen{Amtet de nängänymällneg.}      \extran{[You] take (lit. swallow) a breath.}}    \example{
      \exsen{Obo amtet a gottämänän.}      \extran{His breath slowed.}}}  \sense{\textbf{2.} 
    \definition{to breathe}    \example{
      \exsen{Bogo amtet gognän.}      \extran{He was breathing.}}    \example{
}}}

\entry{amtetmeny}{\headword{amtetmeny}  \pronnote{amtetmeɲ}
  \pos{Adverb}  \sense{
    \definition{nonstop, without breathing}}}

\entry{Ana}{\headword{Ana}  \pronnote{ana}
  \pos{Proper noun}  \sense{
    \definition{female personal name}}}

\entry{anben}{\headword{anben}  \pronnote{anben}
  \pos{Verb}  \sense{
    \definition{to guard}}}

\entry{and}{\headword{and}
  \variantof{freevarlab of \vartext{English loanword}}}

\entry{andred}{\headword{andred}
  \variantof{Dialectal Variant of \vartext{English loanword}}}

\entry{ändred}{\headword{ändred}
  \variantof{SpellingVariant of \vartext{andred}}}

\entry{andred}{\headword{andred}
  \pos{Numeral}  \sense{
    \definition{hundred}}
  \variants{\SpellingVariant \vartext{andredd, ändred, hundred}}}

\entry{andredd}{\headword{andredd}
  \variantof{SpellingVariant of \vartext{andred}}  \pronnote{andreɖ͡ʐ}}

\entry{Andrew}{\headword{Andrew}  \pronnote{ændru}
  \pos{Proper noun}  \sense{
    \definition{male personal name}}
  \variants{\SpellingVariant \vartext{Ainduru}}}

\entry{=ane}{\headword{=ane}
  \variantof{Dialectal Variant of \vartext{=alle}}}

\entry{=ang}{\headword{=ang}  \pronnote{aŋ}
  \pos{Clitic}  \sense{
    \definition{multipurpose attributive suffix that functions as an imperfective or resultative participle, agentive nominalizer, and generic adjective-forming suffix}    \example{
      \exsen{Ngäna dämenang (< dämen) dan.}      \extran{I am sitting.}}    \example{
      \exsen{Ud a papekang (< papek) dan.}      \extran{The door is closed.}}    \example{
      \exsen{Ge lla da mamoeang (< mamoe) dan.}      \extran{This man is a hunter.}}    \example{
      \exsen{ttänttäm, ttänttämang}      \extran{heat, hot}}}}

\entry{angäm}{\headword{angäm}
  \pos{Manner adverb}  \sense{
    \definition{quickly}    \example{
      \exsen{Angäm eka muu nageyo!}      \extran{[You all] answer me quickly!}}    \example{
      \exsen{Sisri sisor lla da angämae de pätärang anggan.}      \extran{Now, young people are quickly getting white hairs.}}}}

\entry{angde}{\headword{angde}  \pronnote{aŋde}
  \pos{Adverb}  \sense{
    \definition{when, while, as}    \example{
      \exsen{Angde ma we gongosalle, llɨg kälekäle da mɨnyi gumaemalle oblle ikop e.}      \extran{When she would return home, the children would gather to see her.}}    \example{
      \exsen{Bam itrell a angde dangkamän?}      \extran{When did you get sick (lit. sickness start on you)?}}    \example{
      \exsen{Dade angde wiaǃ}      \extran{Come wheneverǃ}}}}

\entry{=anggalla}{\headword{=anggalla}
  \pos{Auxiliary verb}  \sense{
    \definition{aux.prs.1|2nsgA>plP}}}

\entry{=anggalle}{\headword{=anggalle}  \pronnote{aŋɡaɽe}
  \pos{Auxiliary verb}  \sense{
    \definition{aux.prs.2sgA>plP}}}

\entry{=anggallo}{\headword{=anggallo}
  \pos{Auxiliary verb}  \sense{
    \definition{aux.prs.3nsgA>3plP}}}

\entry{=anggan}{\headword{=anggan}  \pronnote{aŋɡan}
  \pos{Auxiliary verb}  \sense{\textbf{1.} 
    \definition{aux.prs.3plS|1|3sgA>plP}}  \sense{\textbf{2.} 
}  \sense{\textbf{3.} 
}}

\entry{anggog}{\headword{anggog}  \note{Parts of the Body}  \pronnote{aŋɡoɡ}
  \pos{Noun}  \sense{
    \definition{thigh}    \note{Parts of the Body}    \example{
}    \example{
}    \example{
}}}

\entry{anggwuanemeny}{\headword{anggwuanemeny}  \pronnote{aŋɡwuanemeɲ}
  \pos{}  \sense{
    \definition{cook}    \example{
}    \example{
}    \example{
}}}

\entry{angkäpäll}{\headword{angkäpäll}  \note{Trees}  \pronnote{aŋkəpəɽ}
  \pos{Noun}  \sense{
    \definition{type of big tree found in the bush with white flowers and fruit, which cassowary eat when they fall}    \note{Trees}    \example{
}    \example{
}    \example{
}}}

\entry{ankom}{\headword{ankom}  \note{Medicine}  \pronnote{ankom}
  \pos{Noun}  \sense{
    \definition{ant}    \note{Medicine}    \example{
}    \example{
}    \example{
}}}

\entry{Anna}{\headword{Anna}  \pronnote{ana}
  \pos{Proper noun}  \sense{
    \definition{female personal name}}}

\entry{Ansel}{\headword{Ansel}  \pronnote{ansel}
  \pos{Proper noun}  \sense{
    \definition{male personal name}}}

\entry{ansi}{\headword{ansi}  \pronnote{ansi}
  \pos{Intransitive coverb}  \sense{\textbf{1.} 
    \definition{to sneeze}    \example{
      \exsen{Obo mälläng a nongonongorang gogon, käsre ansi gogon.}      \extran{His nose was itchy; then he sneezed.}}}  \sense{\textbf{2.} 
    \definition{achoo}}}

\entry{Anton}{\headword{Anton}  \pronnote{anton}
  \pos{Proper noun}  \sense{
    \definition{male personal name}}}

\entry{anu}{\headword{anu}  \pronnote{anu}
  \pos{Verb}  \sense{
    \definition{sleep (command given to babies)}}}

\entry{any}{\headword{any}  \pronnote{aɲ}
  \pos{Transitive verb}  \sense{
    \definition{to paint}    \example{
      \exsen{Ngäna ngämo mittapa de pällämpälläm palle alle anyan.}      \extran{I painted my face with white clay.}}}}

\entry{any}{\headword{any}  \pronnote{aɲ}
  \pos{Pronoun}  \sense{\textbf{1.} 
    \definition{something}    \example{
      \exsen{Obo any de särämbae erallo.}      \extran{They are preparing something.}}    \example{
}}  \sense{\textbf{2.} 
    \definition{like this, thus, so}    \example{
      \exsen{Ause da gänyme ine ik i any gogon.}      \extran{The old woman went into this water like this.}}}}

\entry{anyke}{\headword{anyke}  \pronnote{aɲke}
  \pos{Noun}  \sense{
    \definition{spirit}    \example{
      \exsen{Mer Anyke}      \extran{the Holy Spirit}}    \example{
      \exsen{gagäll anyke}      \extran{demon, evil spirit}}    \example{
      \exsen{Anykemeny agan.}      \extran{I got scared (lit. spiritless).}}}
  \variants{\Dialectal Variant \vartext{anyki}}}

\entry{anyke plenz}{\headword{anyke plenz}  \lexeme{plenz}  \pronnote{plenz}
  \pos{Intransitive verb}  \sense{
    \definition{to be in shock}    \example{
      \exsen{Angde Sowati obo bin di dänttamän, Bodog arle gogon a bogo anyke gomplezän.}      \extran{When Sowati called out his name, Bodog screamed and his spirit left his body.}}}}

\entry{anykeanyke}{\headword{anykeanyke}  \pronnote{aɲkeaɲke}
  \pos{Noun}  \sense{
    \definition{picture, image, reflection}    \example{
      \exsen{Ngämi up anykeanyke ikop nägagalla walle ik mi.}      \extran{We (excl.) saw the reflection of bananas in the water.}}}}

\entry{anykegulgul}{\headword{anykegulgul}  \pronnote{aɲkeɡulɡul}
  \pos{Transitive coverb}  \sense{
    \definition{to scare away the fish, go near someone who is fishing and spoil it}}}

\entry{anyki}{\headword{anyki}
  \variantof{Dialectal Variant of \vartext{anyke}}  \pronnote{aɲki}}

\entry{anzag}{\headword{anzag}
  \pos{}}

\entry{ao}{\headword{ao}  \pronnote{ao}
  \pos{Interjection}  \sense{
    \definition{yes}    \example{
      \exsen{― Bongo dalle skul i? ― Ao, ngäna skul ibiatt dan.}      \extran{― Did you go to school? ― Yes, I went to school.}}}}

\entry{aoao}{\headword{aoao}  \pronnote{aoao}
  \pos{Transitive coverb}  \sense{
    \definition{to chase away (e.g. humans, dogs, chickens, but not wild animals)}    \example{
      \exsen{Llɨg de aoao nägaebeyo.}      \extran{[You all] chase away the children.}}}}

\entry{aoli}{\headword{aoli}  \pronnote{aoli}
  \pos{Interrogative pronoun}  \sense{\textbf{1.} 
    \definition{how many, how much}    \example{
      \exsen{Aoli dag up sana da?}      \extran{How much banana cake is there?}}    \example{
      \exsen{Bäne pazi da aoli dan?}      \extran{How old are you (lit. how many are your years)?}}}  \sense{\textbf{2.} 
    \definition{multiple, several (i.e. more than one but not necessarily many); few (esp. with the restrictive clitic)}    \example{
      \exsen{aoli eka kuttae}      \extran{only a few words}}    \example{
      \exsen{Ngämo baba ttongdae ebdo me ai dan auli ddäddäg de bägäddnegän angde obo pollo me dagernän.}      \extran{My father could kill multiple animals in a single day when he was young.}}}  \sense{\textbf{3.} 
    \definition{almost, nearly}    \example{
      \exsen{Lama wa Mana ubi kandärmang me gowensegäneyo adawatta Emi aoli kuddäll de ine ik mi ikop dägagän.}      \extran{Lama and Mana were both feeling sorry because they saw that Emi almost died in the water.}}}
  \variants{\SpellingVariant \vartext{auli}}}

\entry{aolidae}{\headword{aolidae}  \note{Verbs}  \pronnote{aolidaj}
  \pos{Quantifier}  \sense{
    \definition{few}    \note{Verbs}    \example{
      \exsen{Adi ubira ttam e zäme ebdo de aolidae dägnegän.}      \extran{God has made their days to live few.}}    \example{
}    \example{
}}}

\entry{ap}{\headword{ap}  \pronnote{ap}
  \pos{Noun}  \sense{
    \definition{grassland, savannah}    \example{
      \exsen{Llimoll a ap me da.}      \extran{Limol is in the savannah.}}}}

\entry{Apang}{\headword{Apang}
  \pos{Proper noun}  \sense{
    \definition{language name}}}

\entry{apapi}{\headword{apapi}  \lexeme{apapi}  \pronnote{apapi}
  \pos{Noun}  \sense{
    \definition{butterfly}    \example{
      \exsen{Apapi da ddob erag käträkäträl dag.}      \extran{Some butterflies are multi-colored.}}}}

\entry{apapi bärät}{\headword{apapi bärät}
  \pos{Noun}  \sense{
    \definition{type of yam}}}

\entry{apapun}{\headword{apapun}  \lexeme{apapun}  \pronnote{apapun}
  \pos{Noun}  \sense{
    \definition{type of short grass that disperses seeds by attaching to animal fur}}}

\entry{Apdo}{\headword{Apdo}
  \pos{Proper noun}  \sense{
    \definition{female personal name}}}

\entry{apgllu}{\headword{apgllu}  \note{Trees}  \pronnote{apɡɽu}
  \pos{Noun}  \sense{
    \definition{small plant with a root that makes a maroon pigment for dyeing grass skirts when mixed with ash (e.g. Acacia ash or coconut ash). When not mixed with ash, it makes a yellow pigment.}    \note{Trees}    \example{
}    \example{
}    \example{
}}}

\entry{Apodo}{\headword{Apodo}
  \pos{Proper noun}  \sense{
    \definition{female personal name}}}

\entry{apte}{\headword{apte}  \pronnote{apte}
  \pos{Quantifier}  \sense{
    \definition{one side, one end, one half, one (of a pair)}    \example{
      \exsen{Gombällämenymom, apte we deyareyo a apte ade deyareyo.}      \extran{They split up: two went to one side and two went to the other side.}}    \example{
      \exsen{Apte pud di bod me nowanseg eka we.}      \extran{Put one end (of the flute) in your mouth to play.}}    \example{
      \exsen{Ngämo yae obo tubu da apte gagäll dan.}      \extran{One of my mother's knees is bad.}}}}

\entry{apte gabän}{\headword{apte gabän}  \pronnote{apte ɡabən}
  \pos{Numeral}  \sense{
    \definition{fourteen (body counting numeral)}    \example{
}    \example{
}    \example{
}}}

\entry{apte kllatolma}{\headword{apte kllatolma}  \note{Numbers}  \pronnote{apte kɽatolma}
  \pos{Numeral}  \sense{
    \definition{seventeen (body counting numeral)}    \note{Numbers}    \example{
}    \example{
}    \example{
}}}

\entry{apte mända}{\headword{apte mända}  \note{Numbers}  \pronnote{apte mənda}
  \pos{Numeral}  \sense{
    \definition{fifteen (body counting numeral)}    \note{Numbers}    \example{
}    \example{
}    \example{
}}}

\entry{apte mätkin}{\headword{apte mätkin}  \note{Numbers}  \pronnote{apte mətkin}
  \pos{Numeral}  \sense{
    \definition{eighteen (body counting numeral)}    \note{Numbers}    \example{
}    \example{
}    \example{
}}}

\entry{apte matta}{\headword{apte matta}  \note{Numbers}  \pronnote{apte maʈ͡ʂa}
  \pos{Numeral}  \sense{
    \definition{twelve (body counting numeral)}    \note{Numbers}    \example{
}    \example{
}    \example{
}}}

\entry{apte ngam}{\headword{apte ngam}  \note{Numbers}  \pronnote{apte ŋam}
  \pos{Numeral}  \sense{
    \definition{eleven (body counting numeral)}    \note{Numbers}    \example{
}    \example{
}    \example{
}}}

\entry{apte tärangesa}{\headword{apte tärangesa}  \note{Numbers}  \pronnote{apte təraŋesa}
  \pos{Numeral}  \sense{
    \definition{nineteen (body counting numeral)}    \note{Numbers}    \example{
}    \example{
}    \example{
}}}

\entry{apte ttang kum}{\headword{apte ttang kum}  \note{Numbers}  \pronnote{apte ʈ͡ʂaŋ kum}
  \pos{Numeral}  \sense{
    \definition{thirteen (body counting numeral)}    \note{Numbers}    \example{
}    \example{
}    \example{
}}}

\entry{apte tupi}{\headword{apte tupi}  \note{Numbers}  \pronnote{apte tupi}
  \pos{Numeral}  \sense{
    \definition{sixteen (body counting numeral)}    \note{Numbers}    \example{
}    \example{
}    \example{
}}}

\entry{arabuni}{\headword{arabuni}  \pronnote{arabuni}
  \pos{Noun}  \sense{
    \definition{red backed buttonquail}    \scientificname{Turnix maculosus}}}

\entry{Arägapetkae}{\headword{Arägapetkae}  \pronnote{arəɡapetkae}
  \pos{Proper noun}  \sense{
    \definition{Arägapetkae (toponym)}    \example{
}    \example{
}    \example{
}}}

\entry{arälle}{\headword{arälle}
  \variantof{Unspecified Variant of \vartext{arle}}  \note{Animal verbs}  \pronnote{arəɽe}}

\entry{aräram}{\headword{aräram}  \pronnote{arəram}
  \pos{Transitive coverb}  \sense{
    \definition{to sand or smooth a canoe (last step before burning)}}}

\entry{aräre}{\headword{aräre}
  \variantof{Unspecified Variant of \vartext{arle}}  \pronnote{arəre}}

\entry{are}{\headword{are}
  \variantof{Dialectal Variant of \vartext{English loanword}}}

\entry{area}{\headword{area}
  \variantof{freevarlab of \vartext{English loanword}}}

\entry{ariga}{\headword{ariga}  \pronnote{ariɡa}
  \pos{Noun}  \sense{
    \definition{hook, in Agob language.}    \example{
}    \example{
}    \example{
}}}

\entry{arle}{\headword{arle}  \pronnote{arle}
  \pos{Noun}  \sense{
    \definition{scream}    \example{
      \exsen{Däräng a ddone ada arlle gogän.}      \extran{The dog screamed really badly.}}    \example{
      \exsen{Ttongo lla da lelang atta arle bogon.}      \extran{A person can scream from fear.}}    \example{
}}
  \variants{\Dialectal Variant \vartext{arlle}}
  \variants{\Unspecified Variant \vartext{arälle, aräre}}}

\entry{arlle}{\headword{arlle}
  \variantof{Dialectal Variant of \vartext{arle}}  \note{Some pronounce with one l, some pronounce with double l.}  \pronnote{arɽe}}

\entry{Arua}{\headword{Arua}  \pronnote{arua}
  \pos{Proper noun}  \sense{
    \definition{male personal name}}}

\entry{arup}{\headword{arup}  \note{Clothing}  \pronnote{arup}
  \pos{Noun}  \sense{
    \definition{clothing type}    \note{Clothing}    \example{
}    \example{
}    \example{
}}}

\entry{Arupi}{\headword{Arupi}  \note{Place names}  \pronnote{arupi}
  \pos{Proper noun}  \sense{
    \definition{Arupi (toponym)}    \note{Place names}    \example{
}    \example{
}    \example{
}}}

\entry{Aruwa}{\headword{Aruwa}
  \pos{Proper noun}  \sense{
    \definition{male personal name}}}

\entry{as}{\headword{as}  \note{Names of bananas}  \pronnote{as}
  \pos{Noun}  \sense{
    \definition{type of introduced banana}    \note{Names of bananas}    \example{
}    \example{
}    \example{
}}}

\entry{asa}{\headword{asa}
  \variantof{Fast Speech Variant of \vartext{ngasekäma}}}

\entry{asa bume}{\headword{asa bume}  \note{Birds}  \pronnote{asa bume}
  \pos{Noun}  \sense{
    \definition{type of bird}    \note{Birds}    \example{
}    \example{
}    \example{
}}}

\entry{ase}{\headword{ase}
  \pos{Modifier}  \sense{\textbf{1.} 
}  \sense{\textbf{2.} 
}  \sense{\textbf{3.} 
}
  \variants{\freevarlab \vartext{se}}}

\entry{asiasi}{\headword{asiasi}  \lexeme{asiasi}  \pronnote{asiasi}
  \pos{Noun}  \sense{
    \definition{type of large tree that grows in the bush and along big creeks, with white flowers and big red fruits eaten by cassowary and children; the trunk is used to make canoes}}}

\entry{Asika}{\headword{Asika}
  \pos{Proper noun}  \sense{
    \definition{female personal name}}}

\entry{asip}{\headword{asip}  \note{Names of bananas}  \pronnote{asip}
  \pos{Noun}  \sense{
    \definition{type of introduced banana}    \note{Names of bananas}    \example{
}    \example{
}    \example{
}}}

\entry{aska}{\headword{aska}
  \variantof{Fast Speech Variant of \vartext{ngasekäma}}}

\entry{atata kottllam}{\headword{atata kottllam}  \pronnote{atata}
  \pos{Noun}  \sense{
    \definition{type of turtle}}}

\entry{atka}{\headword{atka}  \pronnote{atka}
  \pos{}
  \variants{\Dialectal Variant \vartext{adalla}}}

\entry{atrepo}{\headword{atrepo}  \note{Types of Taro}  \pronnote{atrepo}
  \pos{Noun}  \sense{
    \definition{taro type}    \note{Types of Taro}    \example{
}    \example{
}    \example{
}}}

\entry{=att}{\headword{=att}  \pronnote{aʈ͡ʂ}
  \pos{Nominal enclitic}  \sense{\textbf{1.} 
    \definition{ablative case clitic; from (used for location, time, source, or cause)}    \example{
      \exsen{Oba moko da up di bigezenallo walle ik att de.}      \extran{They want to take the bananas out from the water.}}    \example{
      \exsen{Ge mätta da tätäm att dan.}      \extran{This yam is from yesterday.}}    \example{
      \exsen{Ge ngäma masar att eka da.}      \extran{This is the language of our (excl.) ancestors.}}    \example{
      \exsen{Nyongo da yäbäd ttänttäm att ibimeny bogon.}      \extran{The road can be unwalkable from the heat of the sun.}}}  \sense{\textbf{2.} 
    \definition{past participle clitic}    \example{
      \exsen{Ud a papekatt dan.}      \extran{The door is closed.}}    \example{
      \exsen{Ende eka walle dadäräbatt dan.}      \extran{It is written in Ende.}}    \example{
      \exsen{Tabatt ttängäm a ngämenmeny dan.}      \extran{The promised land is not yet reached.}}    \example{
      \exsen{Eka tameny att me, ubi guinttemänggeyo.}      \extran{Having discussed, the two of them parted ways.}}}}

\entry{atta}{\headword{atta}
  \variantof{Fast Speech Variant of \vartext{adawatta}}}

\entry{au}{\headword{au}  \note{Dead things}  \pronnote{au}
  \pos{Noun}  \sense{\textbf{1.} 
    \definition{burial}    \note{Dead things}    \example{
      \exsen{Au ddägattalle lla da tämamae ma we dingismällalle.}      \extran{After the burial, all the people return home.}}}  \sense{\textbf{2.} 
    \definition{to bury}    \note{Dead things}    \example{
      \exsen{Kuddäll pätt de au dägagallo.}      \extran{They will bury the body.}}    \example{
}    \example{
}}}

\entry{aulämän}{\headword{aulämän}  \note{Birth}  \pronnote{auləmən}
  \pos{Noun}  \sense{
    \definition{conception}    \note{Birth}    \example{
}    \example{
}    \example{
}}}

\entry{auli}{\headword{auli}
  \variantof{SpellingVariant of \vartext{aoli}}  \pronnote{auli}}

\entry{auma}{\headword{auma}  \note{Dead things}  \pronnote{auma}
  \pos{Noun}  \sense{
    \definition{grave}    \note{Dead things}    \example{
      \exsen{Mällause bo auma da dan de.}      \extran{The old woman's grave is there.}}    \example{
}    \example{
}}
  \variants{\SpellingVariant \vartext{auuma}}}

\entry{auma~}{\headword{auma~}
  \variantof{Derived Variant of \vartext{RED~}}}

\entry{auri}{\headword{auri}  \pronnote{auri}
  \pos{Noun}  \sense{
    \definition{metal}}}

\entry{ause}{\headword{ause}  \lexeme{ause}  \pronnote{ause}
  \pos{Noun}  \sense{\textbf{1.} 
    \definition{old woman}    \example{
      \exsen{Ause da känyerkänyer däganän.}      \extran{The old woman soothed him.}}}  \sense{\textbf{2.} 
    \definition{old (of a woman)}    \example{
      \exsen{Kak säre zäme ause abal allan.}      \extran{Grandmother is sadly very old already.}}}}

\entry{auseause}{\headword{auseause}
  \pos{Noun}  \sense{\textbf{1.} 
    \definition{nonsingular form of ause}    \example{
      \exsen{Ngäna ddob täräp me Ende eka de ddone mermer panypeny eran, auseause da ttättle amallo.}      \extran{Sometimes I don't speak Ende properly, so old women correct me.}}}  \sense{\textbf{2.} 
    \definition{nonsingular form of ause}    \example{
      \exsen{Ngämi zime auseause gogmam.}      \extran{We women (excl.) have already become old.}}}}

\entry{Australia}{\headword{Australia}
  \pos{Proper noun}  \sense{
    \definition{Australia}}}

\entry{auuma}{\headword{auuma}
  \variantof{SpellingVariant of \vartext{auma}}}

\entry{Awaba}{\headword{Awaba}  \note{Place names}  \pronnote{awaba}
  \pos{Proper noun}  \sense{
    \definition{Awaba (toponym)}    \note{Place names}    \example{
}    \example{
}    \example{
}}}

\entry{awäll}{\headword{awäll}  \lexeme{awäll}  \pronnote{awəɽ}
  \pos{Noun}  \sense{
    \definition{type of tree that grows in the bush (where yam gardens are made) with white and blue flowers and dark, edible fruit}}
  \variants{\Unspecified Variant \vartext{awoll}}}

\entry{Awayang}{\headword{Awayang}
  \pos{Proper noun}  \sense{
    \definition{male personal name}}}

\entry{awe}{\headword{awe}  \pronnote{awe}
  \pos{Noun}  \sense{
    \definition{savannah}}}

\entry{awe}{\headword{awe}  \pronnote{awe}
  \pos{Noun}  \sense{
    \definition{cassowary (used when hunting)}}}

\entry{Awi}{\headword{Awi}  \note{Place names}  \pronnote{awi}
  \pos{Proper noun}  \sense{
    \definition{Awi (toponym)}    \note{Place names}    \example{
}    \example{
}    \example{
}}}

\entry{awi}{\headword{awi}  \note{Parts of the day}  \pronnote{awi}
  \pos{Noun}  \sense{\textbf{1.} 
    \definition{evening (approx. 5 PM till dark)}    \note{Parts of the day}    \example{
      \exsen{awi alle do iddob amnong}      \extran{from evening to midnight}}    \example{
}    \example{
}}  \sense{\textbf{2.} 
    \definition{early}    \note{Parts of the day}    \example{
      \exsen{Ngäna awi meae ttängäm e dalle.}      \extran{I went to the garden early (e.g. 8 AM instead of 9 AM).}}}}

\entry{awoll}{\headword{awoll}
  \variantof{Unspecified Variant of \vartext{awäll}}}

\entry{=aya}{\headword{=aya}
  \variantof{freevarlab of \vartext{=aeya}}}

\entry{aya}{\headword{aya}  \note{Question words}  \pronnote{aja}
  \pos{Interrogative pronoun}  \sense{\textbf{1.} 
    \definition{who (singular interrogative pronoun, nominative form)}    \note{Question words}    \example{
      \exsen{Kollba de aya nokowan?}      \extran{Who cut the fish?}}}  \sense{\textbf{2.} 
    \definition{whoever, someone (singular existential pronoun, nominative form)}    \note{Question words}    \example{
      \exsen{Mulldae dan wadär de bängällbänän, wadär alle bäbäddän, a ge da aeya ambag bogon.}      \extran{If there's someone causing trouble, it is possible to get a cane and beat them with it.}}}  \sense{\textbf{3.} 
    \definition{who (singular relative pronoun, nominative form)}    \note{Question words}    \example{
      \exsen{Da ngänäm aya bangkollmällän bogo mɨnyi ddone säremang dae ballän.}      \extran{He who follows me will not go through darkness.}}}
  \variants{\SpellingVariant \vartext{aeya}}}

\entry{Azaya}{\headword{Azaya}
  \pos{Proper noun}  \sense{
    \definition{male personal name}}
  \variants{\SpellingVariant \vartext{Isaiah}}}

\chapter{B}


\entry{b-}{\headword{b-}  \pronnote{b}
  \pos{Verb}  \sense{\textbf{1.} 
    \definition{irr.1|3.S|A}}  \sense{\textbf{2.} 
    \definition{abil}}  \sense{\textbf{3.} 
}  \sense{\textbf{4.} 
}  \sense{\textbf{5.} 
}  \sense{\textbf{6.} 
}}

\entry{ba}{\headword{ba}
  \variantof{Fast Speech Variant of \vartext{aba}}  \pronnote{ba}}

\entry{bab}{\headword{bab}  \lexeme{bab}  \pronnote{bab}
  \pos{Noun}  \sense{
    \definition{type of small yam with a white interior}}}

\entry{baba}{\headword{baba}  \lexeme{baba}  \note{Family ties}  \pronnote{baba}
  \pos{Kinship noun}  \sense{
    \definition{father}    \note{Family ties}    \example{
      \exsen{Ngämo baba lla ulle deya.}      \extran{My father was a big man.}}    \example{
}    \example{
}}}

\entry{babaem}{\headword{babaem}  \note{Seasons}  \pronnote{babaem}
  \pos{Noun}  \sense{
    \definition{season characterized by wind and going hunting (fourth season; corresponds to late February)}    \note{Seasons}    \example{
      \exsen{Ddäddäg gäz a babaem täräp me eran ttongo mer abal ttoen dan.}      \extran{Killing animals in babaem is a very good thing.}}    \example{
}    \example{
}}}

\entry{bäbälläd}{\headword{bäbälläd}  \pronnote{bəbəɽəd}
  \pos{Transitive verb}  \sense{
    \definition{to drop without warning}    \example{
}}}

\entry{Babaze}{\headword{Babaze}  \pronnote{babaze}
  \pos{Proper noun}  \sense{
    \definition{male personal name}}}

\entry{babdu}{\headword{babdu}  \note{Types of Taro}  \pronnote{babdu}
  \pos{Noun}  \sense{
    \definition{type of taro}    \note{Types of Taro}    \example{
}    \example{
}    \example{
}}}

\entry{Bablela}{\headword{Bablela}  \pronnote{bablela}
  \pos{Proper noun}  \sense{
    \definition{male personal name}}}

\entry{bäblem}{\headword{bäblem}
  \pos{Intransitive verb}  \sense{
    \definition{to shiver}}}

\entry{bablle}{\headword{bablle}  \pronnote{babɽe}
  \pos{Personal pronoun}  \sense{
    \definition{dative form of bongo}    \example{
      \exsen{Ngäna sana de gämäll agan bablle.}      \extran{I stole the sago for you.}}}}

\entry{bäbnge}{\headword{bäbnge}  \lexeme{bäbnge}  \pronnote{bəbŋe}
  \pos{Noun}  \sense{
    \definition{type of palm with coconuts with a light yellow and green exocarp}}}

\entry{Babra}{\headword{Babra}
  \pos{Proper noun}  \sense{
    \definition{female personal name}}}

\entry{bäbrem}{\headword{bäbrem}  \note{Animal verbs}  \pronnote{bəbrem}
  \pos{Ideophone}  \sense{
    \definition{sound made by cassowaries}    \note{Animal verbs}    \example{
      \exsen{Dirom a bäbrem agan.}      \extran{The cassowary rumbled.}}    \example{
}    \example{
}}}

\entry{Babu}{\headword{Babu}
  \pos{Proper noun}  \sense{
    \definition{male personal name}}}

\entry{bäd}{\headword{bäd}  \lexeme{bäd}  \pronnote{bəd}
  \pos{Noun}  \sense{
    \definition{type of large tree that grows in the grassland with wood used for firewood and bark used for building walls}}}

\entry{bädab}{\headword{bädab}  \pronnote{bədab}
  \pos{Intransitive verb}  \sense{\textbf{1.} 
    \definition{to shine brightly (of the moon)}    \example{
      \exsen{Kokta da näbdaban.}      \extran{The moon is shining.}}}  \sense{\textbf{2.} 
    \definition{to dawn, break}    \example{
      \exsen{Ttongo ag a angde däbdabän ge ttoen de ngäna ddobagabira dällätne, ddone ada gotngoenegänän.}      \extran{The next morning when it dawn broke, I told this story to others and they laughed a lot.}}}  \sense{\textbf{3.} 
    \definition{dawn}    \example{
      \exsen{Kikiem a bädab e ekawang pa dan.}      \extran{The kikiem is a bird that sings at dawn.}}}}

\entry{bädab piro}{\headword{bädab piro}  \pronnote{bədab piro}
  \pos{Noun}  \sense{
    \definition{morning star}}}

\entry{Bädämalloang}{\headword{Bädämalloang}  \pronnote{bədəmaɽoaŋ}
  \pos{Proper noun}  \sense{
    \definition{Bädämalloang (sago place along the road from Limol to the canoe place where a large tree has fallen halfway over the path (AZ95))}}}

\entry{badar}{\headword{badar}  \pronnote{badar}
  \pos{Noun}  \sense{
    \definition{type of tree that grows in the grassland with white flowers}    \anthronote{Gälläpang dan kakne ttoe a. The liquid from the shoots is also used to treat small sores.}    \example{
}    \example{
}    \example{
}}}

\entry{badd~}{\headword{badd~}
  \variantof{freevarlab of \vartext{RED~}}}

\entry{baddbedd}{\headword{baddbedd}  \pronnote{baɖ͡ʐbeɖ͡ʐ}
  \pos{Intransitive verb}  \sense{
    \definition{to dry up}    \example{
      \exsen{Ine da däbeddnegän, ada källäm a, ada walle mäg a, orbam a wa karama da ade petapeta abal gogon.}      \extran{Water dried up, and the ponds, the big creeks, pools, and swamps became very shallow.}}}}

\entry{bädde}{\headword{bädde}  \note{Types of Taro}  \pronnote{bəɖ͡ʐe}
  \pos{Noun}  \sense{
    \definition{type of big taro}    \note{Types of Taro}    \example{
}    \example{
}    \example{
}}}

\entry{bädde}{\headword{bädde}  \lexeme{bädde}  \pronnote{bəɖ͡ʐe}
  \pos{Noun}  \sense{
    \definition{type of large tree found by rivers}}}

\entry{bade}{\headword{bade}  \pronnote{bade}
  \pos{Personal pronoun}  \sense{
    \definition{additive form of bongo}    \example{
      \exsen{Senti, bade bablle ikom.}      \extran{Senti, you take some for yourself too.}}}}

\entry{bädma}{\headword{bädma}  \note{Medicine}  \pronnote{bədma}
  \pos{Noun}  \sense{
    \definition{type of medicinal plant}    \note{Medicine}    \example{
}    \example{
}    \example{
}}}

\entry{bädma ibeny}{\headword{bädma ibeny}  \note{Conflict/war}  \pronnote{bədma ibeɲ}
  \pos{Noun}  \sense{
    \definition{planting a bädma plant as a gesture of peace}    \note{Conflict/war}    \example{
}    \example{
}    \example{
}}}

\entry{bädmaol}{\headword{bädmaol}
  \pos{Noun}  \sense{
    \definition{small sago flower}}}

\entry{Badu}{\headword{Badu}  \pronnote{badu}
  \pos{Proper noun}  \sense{
    \definition{male personal name}}}

\entry{baebol}{\headword{baebol}
  \pos{Noun}  \sense{
    \definition{Bible}    \example{
}}
  \variants{\SpellingVariant \vartext{Bible}}}

\entry{baet}{\headword{baet}  \lexeme{baet}  \pronnote{baet}
  \pos{Noun}  \sense{
    \definition{cuscus}    \scientificname{Spilocuscus maculatus}    \example{
      \exsen{Abo baet ttoe de näglle nyäng ngasnges e.}      \extran{Skin the cuscus to make a bag.}}}}

\entry{Baewa}{\headword{Baewa}  \pronnote{baewa}
  \pos{Proper noun}  \sense{
    \definition{male personal name}}}

\entry{bägäbägäl}{\headword{bägäbägäl}  \note{Parts of the body}  \pronnote{bəɡəbəɡəl}
  \pos{Noun}  \sense{
    \definition{Achilles/calcaneal tendon (on the back of ankle)}    \note{Parts of the body}    \example{
}    \example{
}    \example{
}}}

\entry{bägäl}{\headword{bägäl}  \lexeme{bägäl}  \pronnote{bəɡəl}
  \pos{Noun}  \sense{
    \definition{bow}    \example{
      \exsen{Bägäl de popo eran.}      \extran{He shapes the bow.}}    \example{
      \exsen{Masar täräp me lla da bägäl alle gogäddnän.}      \extran{In the old times, people fought with bows.}}}}

\entry{bägäl kuwem}{\headword{bägäl kuwem}  \pronnote{bəɡəl kuwem}
  \pos{Noun}  \sense{
    \definition{gunfire}}}

\entry{bägäl mäträt}{\headword{bägäl mäträt}  \note{Hunting}  \pronnote{bəɡəl mətrət}
  \pos{Noun}  \sense{
    \definition{part of bow}    \note{Hunting}    \example{
}    \example{
}    \example{
}}}

\entry{bägäl odar}{\headword{bägäl odar}  \note{Everything in the house}  \pronnote{bəɡəl odar}
  \pos{Noun}  \sense{
    \definition{place for bows and spears}    \note{Everything in the house}    \example{
}    \example{
}    \example{
}}}

\entry{bägäl tangge}{\headword{bägäl tangge}  \note{Hunting}  \pronnote{bəɡəl taŋɡe}
  \pos{Noun}  \sense{
    \definition{extra bowstring}    \note{Hunting}    \example{
}    \example{
}    \example{
}}}

\entry{bägälbägäl}{\headword{bägälbägäl}  \note{Hunting}  \pronnote{bəɡəlbəɡəl}
  \pos{Noun}  \sense{
    \definition{small bow}    \note{Hunting}    \example{
}    \example{
}    \example{
}}}

\entry{bägallem}{\headword{bägallem}  \note{Trees}  \pronnote{bəɡaɽem}
  \pos{Noun}  \sense{
    \definition{type of tree that grows in the grassland with wood used for firewood and yellow fruit}    \note{Trees}    \example{
}    \example{
}    \example{
}}}

\entry{bägäm}{\headword{bägäm}  \note{Trees}  \pronnote{bəɡəm}
  \pos{Noun}  \sense{
    \definition{type of tree found by creeks with bark used for weaving bags or strong rope}    \note{Trees}    \example{
}    \example{
}    \example{
}}}

\entry{bagama}{\headword{bagama}
  \pos{Noun}  \sense{
}}

\entry{bägem}{\headword{bägem}  \note{Trees}  \pronnote{bəɡem}
  \pos{Noun}  \sense{
    \definition{type of tree with pink flowers and fist-sized fruit with a pit}    \note{Trees}    \example{
}    \example{
}    \example{
}}}

\entry{bagen}{\headword{bagen}  \note{Types of Taro}  \pronnote{baɡen}
  \pos{Noun}  \sense{
    \definition{type of big taro}    \note{Types of Taro}    \example{
}    \example{
}    \example{
}}}

\entry{Bäglle}{\headword{Bäglle}
  \pos{Proper noun}  \sense{
    \definition{male personal name}}}

\entry{Baiduwa}{\headword{Baiduwa}  \note{Place names}  \pronnote{baiduwa}
  \pos{Proper noun}  \sense{
    \definition{Baiduwa (toponym)}    \note{Place names}    \example{
}    \example{
}    \example{
}}}

\entry{Baim}{\headword{Baim}  \note{Place names around Limol}  \pronnote{baim}
  \pos{Proper noun}  \sense{
    \definition{Baim (toponym)}    \note{Place names around Limol}    \example{
}    \example{
}    \example{
}}}

\entry{bäkän}{\headword{bäkän}  \note{Trees}  \pronnote{bəkən}
  \pos{Noun}  \sense{
    \definition{type of cultivated plant with leaves eaten with sago}    \note{Trees}    \example{
}    \example{
}    \example{
}}}

\entry{=bakmall}{\headword{=bakmall}  \pronnote{bakmaɽ}
  \pos{Nominal enclitic}  \sense{
    \definition{comitative case clitic; with (only used after oba, ama, and nouns followed by =aba)}    \example{
      \exsen{Ngämi oba bakmall eka tameny amalla.}      \extran{We (excl.) discussed with them.}}    \example{
      \exsen{Ge lla da ngäna ama bakmall skul att dan.}      \extran{These people are who I went to school with.}}    \example{
      \exsen{Inggoemeny e bog tongoenen e bägne, bogmam nagnag aba bakmall.}      \extran{We will tell jokes with our friends.}}}}

\entry{bako}{\headword{bako}
  \pos{Personal pronoun}  \sense{
    \definition{additive form of bongo}    \example{
      \exsen{Bako ikop nas.}      \extran{You close your eyes too.}}}}

\entry{bälämbäl}{\headword{bälämbäl}  \pronnote{bələmbəl}
  \pos{Transitive verb}  \sense{\textbf{1.} 
    \definition{to miss, long for}    \example{
      \exsen{Ngäna ttongo lla de bälämbäl eran.}      \extran{I miss someone.}}}  \sense{\textbf{2.} 
    \definition{to remember, think of}    \example{
      \exsen{Angde lla da borale bo eka de bäddareyo, ubi mɨnyi ngattongatt masarmasar de bämbälaebeyo.}      \extran{When people hear the song of the flute, they will remember their ancestors from before.}}    \example{
      \exsen{Ngäma ddob gagäll gudae att gagäll de dämbälaebnalla.}      \extran{Some of us (excl.) were remembering bad things from the past.}}}}

\entry{Balimo}{\headword{Balimo}  \note{Place names}  \pronnote{balimo}
  \pos{Proper noun}  \sense{
    \definition{Balimo (toponym)}    \note{Place names}    \example{
}    \example{
}    \example{
}}}

\entry{bäll}{\headword{bäll}  \note{Parts of the Body}  \pronnote{bəɽ}
  \pos{Noun}  \sense{
    \definition{thigh}    \note{Parts of the Body}    \example{
}    \example{
}    \example{
}}}

\entry{bäll kutt}{\headword{bäll kutt}  \note{Parts of a reptile}  \pronnote{bəɽ kuʈ͡ʂ}
  \pos{Noun}  \sense{
    \definition{femur}    \note{Parts of a reptile}    \example{
}    \example{
}    \example{
}}}

\entry{bälla~}{\headword{bälla~}
  \variantof{Infinitival reduplicant of \vartext{RED~}}}

\entry{bällabälle}{\headword{bällabälle}  \pronnote{bəɽabəɽe}
  \pos{Transitive verb}  \sense{
    \definition{to find}    \example{
      \exsen{Ngäna kumuddäga ttägäll käp tärpitärpi de nabllenegan.}      \extran{I found three smooth stones.}}}
  \variants{\Fast Speech Variant \vartext{bllablle}}}

\entry{bällae}{\headword{bällae}  \pronnote{bəɽaj}
  \pos{Modifier}  \sense{
    \definition{loud}}}

\entry{bälläg}{\headword{bälläg}  \note{Names of bananas}  \pronnote{bəɽəɡ}
  \pos{Noun}  \sense{
    \definition{type of introduced banana}    \note{Names of bananas}    \example{
}    \example{
}    \example{
}}}

\entry{bälläg}{\headword{bälläg}  \note{Types of palms}  \pronnote{bəɽəɡ}
  \pos{Noun}  \sense{
    \definition{Areca palm}    \scientificname{Areca catechu}    \note{Types of palms}    \example{
}    \example{
}    \example{
}}}

\entry{bällam}{\headword{bällam}  \pronnote{bəɽam}
  \pos{Quantifier}  \sense{
    \definition{every}    \example{
      \exsen{täräp bällam}      \extran{always}}    \example{
      \exsen{Ebdo bällam, yogoll a manen eran.}      \extran{Every day, it rains.}}    \example{
      \exsen{Bongo tämamae ttoen bällam e zizag dan.}      \extran{You are the owner of everything.}}}
  \variants{\Fast Speech Variant \vartext{bllam}}}

\entry{bällamb}{\headword{bällamb}
  \pos{Transitive verb}  \sense{
    \definition{to ambush}}}

\entry{bällämbäll}{\headword{bällämbäll}  \pronnote{bəɽambəɽ}
  \pos{Noun}  \sense{
    \definition{thought}}}

\entry{bällämbällämmeny}{\headword{bällämbällämmeny}  \pronnote{bəɽəmbəɽəmeɲ}
  \pos{Modifier}  \sense{
    \definition{thoughtless, apathetic, uncaring}    \example{
      \exsen{Däbe lla da eran lla bällämbällämmeny dan.}      \extran{That man doesn't care about others.}}}}

\entry{bällängg}{\headword{bällängg}  \pronnote{mbəɽə}
  \pos{Intransitive verb}  \sense{
    \definition{to split up, separate}    \example{
      \exsen{Ubi gombällägmenyamän.}      \extran{They were splitting up.}}}}

\entry{bällangg}{\headword{bällangg}  \pronnote{bombɽaɡ}
  \pos{Transitive verb}  \sense{
    \definition{to find}    \example{
      \exsen{Ngäna mɨnyi obom bombllag.}      \extran{I will find him.}}    \example{
      \exsen{Bogo mɨnyi obom bombllagän.}      \extran{He will find him.}}    \example{
}}}

\entry{ballängg}{\headword{ballängg}
  \variantof{Unspecified Variant of \vartext{ballɨngg}}  \pronnote{baɽəŋɡ}}

\entry{ballbell}{\headword{ballbell}  \pronnote{baɽbeɽ}
  \pos{Intransitive coverb}  \sense{
    \definition{to cook improperly}    \example{
      \exsen{Sana da yu me ballbell allan.}      \extran{Sago is cooking improperly on the fire.}}}}

\entry{balldae}{\headword{balldae}  \pronnote{baɽdaj}
  \pos{Verb}  \sense{
    \definition{to announce}}}

\entry{=bälle}{\headword{=bälle}  \pronnote{bəɽe}
  \pos{Pronominal enclitic}  \sense{
    \definition{clitic form of oblle}    \example{
      \exsen{Obo nane bälle ikop e dallän.}      \extran{She went to see her aunt.}}}}

\entry{bällgab}{\headword{bällgab}  \pronnote{bəɽɡab}
  \pos{Transitive verb}  \sense{
    \definition{to open (e.g. eyes, flowers)}    \example{
      \exsen{Däräng täräpang dagirnän ikop saeyang käsre ikop däbällgabän.}      \extran{Dog stayed there a long time with eyes closed, then opened his eyes.}}}}

\entry{bällgae}{\headword{bällgae}  \pronnote{bəɽɡe}
  \pos{Intransitive verb}  \sense{
    \definition{to spread, scatter, move away}    \example{
      \exsen{Ngämo mosemosen a ddob mälla peyang gognegän, ede ubi gobällgewän.}      \extran{My older brothers got married, so they moved away (from home).}}}}

\entry{ballɨngg}{\headword{ballɨngg}  \pronnote{baɽɪŋɡ}
  \pos{Transitive verb}  \sense{\textbf{1.} 
    \definition{to welcome, greet}    \example{
      \exsen{Ngänaeka alle yamballägallo.}      \extran{In tears, they greeted him.}}    \example{
      \exsen{Sisor pazi di bamballmenyneyo.}      \extran{They will be celebrating the new year.}}}  \sense{\textbf{2.} 
    \definition{to predict}    \example{
      \exsen{Bogo abade ngasnen e ttoen de damballmenynegän.}      \extran{He predicted things that will happen in the future.}}}
  \variants{\Unspecified Variant \vartext{ballängg}}}

\entry{bällkäp}{\headword{bällkäp}  \pronnote{bəɽkəp}
  \pos{Noun}  \sense{
    \definition{ant pupae}}}

\entry{ballma}{\headword{ballma}  \note{Bees}  \pronnote{baɽma}
  \pos{Noun}  \sense{
    \definition{type of biting bee found in trees}    \note{Bees}    \example{
}    \example{
}    \example{
}}}

\entry{bällma}{\headword{bällma}  \pronnote{bəɽma}
  \pos{Noun}  \sense{
    \definition{spit, saliva}}}

\entry{bällma käkan}{\headword{bällma käkan}  \pronnote{bəɽma kəkan}
  \pos{Noun}  \sense{
    \definition{spit, saliva}}}

\entry{bällma ontog}{\headword{bällma ontog}  \pronnote{bəɽma ontoɡ}
  \pos{Transitive compound verb}  \sense{
    \definition{to spit at}    \example{
      \exsen{Ubi obom bällma dawäntmenyneyo.}      \extran{They spat at him.}}}}

\entry{ballme}{\headword{ballme}  \pronnote{baɽme}
  \pos{Noun}  \sense{
    \definition{dawn, daybreak}}}

\entry{ballo bällabällott}{\headword{ballo bällabällott}  \note{Types of Taro}  \pronnote{baɽo bəɽabəɽoʈ͡ʂ}
  \pos{Noun}  \sense{
    \definition{type of big taro}    \note{Types of Taro}    \example{
}    \example{
}    \example{
}}}

\entry{bälltoe}{\headword{bälltoe}  \note{Trees}  \pronnote{bəɽtoe}
  \pos{Noun}  \sense{
    \definition{type of tree}    \note{Trees}    \example{
}    \example{
}    \example{
}}}

\entry{bälwäd}{\headword{bälwäd}
  \variantof{Unspecified Variant of \vartext{bolod}}  \pronnote{bəlwəd}}

\entry{bälwod}{\headword{bälwod}
  \variantof{Unspecified Variant of \vartext{bolod}}  \pronnote{bəlwod}}

\entry{bam}{\headword{bam}  \pronnote{bam}
  \pos{Personal pronoun}  \sense{
    \definition{accusative form of bongo}    \example{
      \exsen{Ngäna bam trungg allan.}      \extran{I am calling you.}}}}

\entry{bam bam doros}{\headword{bam bam doros}  \note{Clothing}  \pronnote{bam bam doros}
  \pos{Noun}  \sense{
    \definition{bum bum pants}    \note{Clothing}    \example{
}    \example{
}    \example{
}}}

\entry{Bämäg}{\headword{Bämäg}  \note{Places around Limol}  \pronnote{bəməɡ}
  \pos{Proper noun}  \sense{
    \definition{Bämäg (toponym)}    \note{Places around Limol}    \example{
}    \example{
}    \example{
}}}

\entry{bämäng}{\headword{bämäng}  \note{Trees}  \pronnote{bəməŋ}
  \pos{Noun}  \sense{
    \definition{type of tree}    \note{Trees}    \example{
}    \example{
}    \example{
}}}

\entry{bämblläd}{\headword{bämblläd}  \note{is this bämblläd or bämbllädd? see example}  \pronnote{bəmbɽəd}
  \pos{Transitive verb}  \sense{
    \definition{to expand, rise}    \note{is this bämblläd or bämbllädd? see example}}}

\entry{bamearoro}{\headword{bamearoro}  \note{Mushrooms}  \pronnote{bamearoro}
  \pos{Noun}  \sense{
    \definition{type of mushroom}    \note{Mushrooms}    \example{
}    \example{
}    \example{
}}}

\entry{Bana}{\headword{Bana}
  \pos{Proper noun}  \sense{
    \definition{female personal name}}}

\entry{bänäm}{\headword{bänäm}  \pronnote{bənəm}
  \pos{Noun}  \sense{\textbf{1.} 
    \definition{type of very small insect that lives on bandicoots}    \example{
      \exsen{Maigag a bänäm peyang dag.}      \extran{Bandicoots have tiny insects.}}}  \sense{\textbf{2.} 
    \definition{thorns on sago leaves}}}

\entry{bänamb}{\headword{bänamb}  \pronnote{bənamb}
  \pos{Transitive verb}  \sense{
    \definition{to open (something folded, e.g. book, mouth)}    \example{
      \exsen{Buk di bämbnaebmam.}      \extran{We will open the book.}}    \example{
      \exsen{Obo kili peyang ttang de imbinaebnegan.}      \extran{He happily arrived with his arms open.}}}}

\entry{bänäne}{\headword{bänäne}
  \variantof{SpellingVariant of \vartext{bänene}}}

\entry{bänbän}{\headword{bänbän}  \pronnote{bənbən}
  \pos{Noun}  \sense{
    \definition{the fifth stage of coconut growth in which the fruits are beginning to form}    \example{
      \exsen{Bänbän de mudag ddädd a, däbe källkae e bazernän.}}}}

\entry{band}{\headword{band}  \note{Trees}  \pronnote{band}
  \pos{Noun}  \sense{
    \definition{type of tree}    \note{Trees}    \example{
}    \example{
}    \example{
}}}

\entry{bändaeg}{\headword{bändaeg}  \pronnote{bəndaeɡ}
  \pos{Transitive verb}  \sense{
    \definition{to pull an all-nighter (stay awake)}    \example{
      \exsen{Ibi iddob mɨnyi beyambadägeya.}      \extran{We (incl.) will pull an all-nighter.}}}}

\entry{bändam}{\headword{bändam}  \note{Trees}  \pronnote{bəndam}
  \pos{Noun}  \sense{
    \definition{type of tree}    \note{Trees}    \example{
}    \example{
}    \example{
}}}

\entry{bandra}{\headword{bandra}  \pronnote{bandra}
  \pos{Noun}  \sense{
    \definition{song}    \example{
      \exsen{Wagiba bandra de nɨllɨtnan sisri ag me.}      \extran{Wagiba was singing this morning.}}}}

\entry{=bäne}{\headword{=bäne}  \pronnote{bəne}
  \pos{Pronominal enclitic}  \sense{
    \definition{clitic form of obene}    \example{
      \exsen{Pip bäne ddäddäg a ttällanenang dan.}      \extran{The sting of the red bee is painful.}}}}

\entry{bäne}{\headword{bäne}  \pronnote{bəne}
  \pos{Personal pronoun}  \sense{
    \definition{possessive form of bongo}    \example{
      \exsen{Bäne däräng a daden.}      \extran{You have a dog (lit. your dog exists).}}}}

\entry{bänene}{\headword{bänene}  \pronnote{bənene}
  \pos{Personal pronoun}  \sense{
    \definition{ablative-possessive form of bongo}    \example{
      \exsen{Ngäna sana de gämäll agan bänene.}      \extran{I stole the sago from you.}}}
  \variants{\SpellingVariant \vartext{bänäne}}}

\entry{bänezaga}{\headword{bänezaga}
  \pos{Personal pronoun}  \sense{
    \definition{yourself (reflexive form of bongo)}    \example{
      \exsen{Ddone dan otät yu watt a, abo bänezaga yu näga.}      \extran{There isn't any cooked food; cook it yourself.}}}}

\entry{bäng}{\headword{bäng}  \pronnote{bəŋ}
  \pos{Noun}  \sense{
    \definition{firestick (to start a fire)}    \example{
}}}

\entry{bänga}{\headword{bänga}  \pronnote{bəŋa}
  \pos{Subordinating connective}  \sense{\textbf{1.} 
    \definition{though}    \example{
      \exsen{Bänga ngäna llamda dan, be ngäna mängallang dan.}      \extran{Though I am an old man, I am strong.}}}  \sense{\textbf{2.} 
    \definition{should}    \example{
      \exsen{Bänga buwensegenyallo.}      \extran{We should leave it alone.}}}}

\entry{banggo}{\headword{banggo}
  \pos{Noun}  \sense{
    \definition{type of long yam with a white interior, thorns, and no hair}}}

\entry{banggu}{\headword{banggu}  \note{Clothing}  \pronnote{baŋɡu}
  \pos{Noun}  \sense{
    \definition{headdress}    \note{Clothing}    \example{
      \exsen{Däbe bäne banggu di aeya diwän?}      \extran{Who wove that headdress for you?}}    \example{
}    \example{
}}}

\entry{bäny~}{\headword{bäny~}
  \variantof{Inflected Form of \vartext{RED~}}}

\entry{bänybäny}{\headword{bänybäny}  \pronnote{bəɲbəɲ}
  \pos{Transitive verb}  \sense{
    \definition{to cut, slice (flesh)}    \example{
      \exsen{Ngämi ge ddia de däbänya.}      \extran{We (excl.) cut this deer.}}}}

\entry{bänz}{\headword{bänz}  \note{Bees}  \pronnote{bənz}
  \pos{Noun}  \sense{
    \definition{type of biting bee found in trees}    \note{Bees}    \example{
}    \example{
}    \example{
}}}

\entry{bänz}{\headword{bänz}  \lexeme{bänz}  \pronnote{bənz}
  \pos{Noun}  \sense{
    \definition{mosquito}}}

\entry{bänzibänzi}{\headword{bänzibänzi}  \note{Sago palms}  \pronnote{bənzibənzi}
  \pos{Noun}  \sense{
    \definition{type of sago}    \note{Sago palms}    \example{
}    \example{
}    \example{
}}}

\entry{baob}{\headword{baob}  \lexeme{baob}  \pronnote{baob}
  \pos{Noun}  \sense{
    \definition{water lily}}}

\entry{Baogab}{\headword{Baogab}  \note{Places around Limol}  \pronnote{baoɡab}
  \pos{Proper noun}  \sense{
    \definition{Baogab (an island in Karama swamp used for camping; has coconuts and bananas)}    \note{Places around Limol}    \example{
}    \example{
}    \example{
}}}

\entry{baottbaott}{\headword{baottbaott}  \pronnote{baoʈ͡ʂbaoʈ͡ʂ}
  \pos{Intransitive coverb}  \sense{
    \definition{to walk around aimlessly}}}

\entry{bäräbäräl}{\headword{bäräbäräl}  \note{Birds}  \pronnote{bərəbərəl}
  \pos{Noun}  \sense{
    \definition{type of bird}    \note{Birds}    \example{
}    \example{
}    \example{
}}}

\entry{bärät}{\headword{bärät}  \lexeme{bärät}  \pronnote{bərət}
  \pos{Noun}  \sense{
    \definition{type of small yam}}}

\entry{Barekam}{\headword{Barekam}  \pronnote{barekam}
  \pos{Proper noun}  \sense{
    \definition{male personal name}}}

\entry{bargae}{\headword{bargae}
  \pos{Noun}  \sense{
    \definition{type of fish}}}

\entry{bärke}{\headword{bärke}  \note{Birds}  \pronnote{bərke}
  \pos{Noun}  \sense{
    \definition{Papuan eclectus}    \scientificname{Eclectus polychloros}    \note{Birds}    \example{
}    \example{
}    \example{
}}}

\entry{bärkebärke}{\headword{bärkebärke}  \lexeme{bärkebärke}  \pronnote{bərkebərke}
  \pos{Noun}  \sense{
    \definition{type of algae that can be red or green like a parrot}}}

\entry{Basido}{\headword{Basido}  \note{Place names around Limol}  \pronnote{basido}
  \pos{Proper noun}  \sense{
    \definition{Basido (toponym)}    \note{Place names around Limol}    \example{
}    \example{
}    \example{
}}}

\entry{Basido kona}{\headword{Basido kona}  \note{Places around Limol}  \pronnote{basido kona}
  \pos{Proper noun}  \sense{
    \definition{Pastor's corner (in Limol)}    \note{Places around Limol}    \example{
}    \example{
}    \example{
}}}

\entry{bät}{\headword{bät}
  \variantof{Unspecified Variant of \vartext{bɨt}}}

\entry{bät}{\headword{bät}  \note{Trees}  \pronnote{bət}
  \pos{Noun}  \sense{
    \definition{type of tree}    \note{Trees}    \example{
}    \example{
}    \example{
}}}

\entry{bätäny}{\headword{bätäny}  \note{Trees}  \pronnote{bətəɲ}
  \pos{Noun}  \sense{
    \definition{type of tree}    \note{Trees}    \example{
}    \example{
}    \example{
}}}

\entry{bätbät}{\headword{bätbät}
  \variantof{SpellingVariant of \vartext{bɨtbɨt}}}

\entry{Bati}{\headword{Bati}  \pronnote{bati}
  \pos{Proper noun}  \sense{
    \definition{female personal name}}}

\entry{batri}{\headword{batri}  \pronnote{batri}
  \pos{Noun}  \sense{
    \definition{battery}}}

\entry{batri käp}{\headword{batri käp}
  \pos{Noun}  \sense{
    \definition{battery}}}

\entry{batt}{\headword{batt}  \note{Stages of life}  \pronnote{baʈ͡ʂ}
  \pos{Modifier}  \sense{
    \definition{adult, mature}    \note{Stages of life}    \example{
      \exsen{lla batt, mälla batt}      \extran{adult man, adult woman}}    \example{
      \exsen{Bogo batt däm de ibeb eran.}      \extran{She plants a mature plant.}}    \example{
}}}

\entry{batt}{\headword{batt}  \pronnote{baʈ͡ʂ}
  \pos{Noun}  \sense{
    \definition{central lateral beam of a house}}}

\entry{bätte}{\headword{bätte}  \note{Snakes}  \pronnote{bəʈ͡ʂe}
  \pos{Noun}  \sense{
    \definition{type of snake}    \note{Snakes}    \example{
}    \example{
}    \example{
}}
  \variants{\Unspecified Variant \vartext{bättekuibiag}}}

\entry{bättekuibiag}{\headword{bättekuibiag}
  \variantof{Unspecified Variant of \vartext{bätte}}}

\entry{baur}{\headword{baur}  \note{Traditional Arrows}  \pronnote{baur}
  \pos{Noun}  \sense{
    \definition{type of spear}    \note{Traditional Arrows}    \example{
}    \example{
}    \example{
}}
  \variants{\SpellingVariant \vartext{bawur}}}

\entry{bawa}{\headword{bawa}  \note{Seasons}  \pronnote{bawa}
  \pos{Noun}  \sense{\textbf{1.} 
    \definition{season characterized by hunting and fishing in heavy rain (ninth season; corresponds to June)}    \note{Seasons}    \example{
      \exsen{Bawa täräp me nyongo da era ddobae tärpi dan.}      \extran{During the shower times, the roads can be very slippery.}}    \example{
}    \example{
}}  \sense{\textbf{2.} 
    \definition{rain shower}    \note{Seasons}}}

\entry{bawa minyminy}{\headword{bawa minyminy}  \note{Seasons}  \pronnote{bawa miɲmiɲ}
  \pos{Noun}  \sense{
    \definition{season characterized by hunting and fishing in light rain (tenth season; corresponds to July)}    \note{Seasons}    \example{
}    \example{
}    \example{
}}}

\entry{bawa pokallmäll}{\headword{bawa pokallmäll}  \note{Water}  \pronnote{bawa pokaɽməɽ}
  \pos{Noun}  \sense{
    \definition{wave}    \note{Water}    \example{
}    \example{
}    \example{
}}}

\entry{bawa sarasaram}{\headword{bawa sarasaram}  \pronnote{bawa sarasaram}
  \pos{Noun}  \sense{
    \definition{drizzle}    \example{
}    \example{
}    \example{
}}}

\entry{bawur}{\headword{bawur}
  \variantof{SpellingVariant of \vartext{baur}}  \note{Hunting}  \pronnote{bawur}}

\entry{bazere}{\headword{bazere}  \lexeme{bazere}  \pronnote{bazere}
  \pos{Noun}  \sense{
    \definition{type of purple yam with hairs and no thorns}}}

\entry{be}{\headword{be}  \pronnote{be}
  \pos{Coordinating connective}  \sense{
    \definition{but}    \example{
      \exsen{Pamker a käpang gogon, be tomato da ddone käpang gogon.}      \extran{The pumpkin plant is bearing fruit, but the tomato plant is not bearing fruit.}}}}

\entry{be~}{\headword{be~}
  \variantof{freevarlab of \vartext{RED~}}}

\entry{beatururang}{\headword{beatururang}  \note{Seasons}  \pronnote{beatururaŋ}
  \pos{Noun}  \sense{
    \definition{season characterized by thunderstorms and flooding (second season; corresponds to early February)}    \note{Seasons}    \example{
}    \example{
}    \example{
}}}

\entry{beawa}{\headword{beawa}
  \variantof{SpellingVariant of \vartext{beyawa}}  \pronnote{beawa}}

\entry{beawaenen}{\headword{beawaenen}
  \variantof{SpellingVariant of \vartext{beyawaenen}}}

\entry{bebe}{\headword{bebe}  \note{Type of pandanus}  \pronnote{bebe}
  \pos{Noun}  \sense{
    \definition{type of pandanus with long fruit (~2 feet)}    \note{Type of pandanus}    \example{
}    \example{
}    \example{
}}}

\entry{bebe}{\headword{bebe}  \pronnote{bebe}
  \pos{Intransitive verb}  \sense{
    \definition{to leak, excrete}    \example{
      \exsen{Källa bebe we nalle.}      \extran{[You] go and pass excrement.}}    \example{
      \exsen{De däräng a källa benanbenan dan.}      \extran{This dog is always pooping.}}    \example{
      \exsen{Sana sisor a nyukukum me bebe allan.}      \extran{The new sago is leaking in the sago bag.}}}}

\entry{Bebelin}{\headword{Bebelin}
  \pos{Proper noun}  \sense{
    \definition{female personal name}}}

\entry{bebi}{\headword{bebi}  \pronnote{bebi}
  \pos{Noun}  \sense{
    \definition{baby}    \example{
      \exsen{Bebi da ngänaeka allan.}      \extran{The baby is crying.}}}
  \variants{\SpellingVariant \vartext{beibi}}}

\entry{begere}{\headword{begere}  \lexeme{begere}  \pronnote{beɡere}
  \pos{Noun}  \sense{
    \definition{type of long purple yam}}}

\entry{beibi}{\headword{beibi}
  \variantof{SpellingVariant of \vartext{bebi}}}

\entry{bel}{\headword{bel}
  \pos{Noun}  \sense{
    \definition{bell}}}

\entry{bem}{\headword{bem}  \lexeme{bem}  \pronnote{bem}
  \pos{Noun}  \sense{
    \definition{sea, ocean}    \example{
      \exsen{Kila a Wala gall alle tudi ittaenen e deyareyo bem e.}      \extran{Kila and Wala went to the sea with the canoe to fish.}}}}

\entry{Ben}{\headword{Ben}  \pronnote{ben}
  \pos{Proper noun}  \sense{
    \definition{male personal name}}}

\entry{Benaeya}{\headword{Benaeya}
  \pos{Proper noun}  \sense{
    \definition{male personal name}}}

\entry{benanbenan}{\headword{benanbenan}  \note{Traditional Arrows}  \pronnote{benanbenan}
  \pos{Noun}  \sense{
    \definition{type of spear}    \note{Traditional Arrows}    \example{
}    \example{
}    \example{
}}}

\entry{bendoe}{\headword{bendoe}  \pronnote{bendoe}
  \pos{Transitive verb}  \sense{
    \definition{to confuse, mix up}    \example{
      \exsen{Ngäna ubim komlla angde ikop deyagän, ngäna yambedowaeyan.}      \extran{When I saw them two, I confused them for each other.}}    \example{
}    \example{
}}}

\entry{bengae}{\headword{bengae}  \pronnote{beŋae}
  \pos{Transitive verb}  \sense{\textbf{1.} 
    \definition{to roof, cover}    \example{
      \exsen{Ubi ttägäll de nyeny alle daibengaemallo.}      \extran{They covered the mumus with nyeny wood.}}}  \sense{\textbf{2.} 
    \definition{roof, roofing}}}

\entry{benmäll}{\headword{benmäll}  \pronnote{benməɽ}
  \pos{Intransitive verb}  \sense{\textbf{1.} 
    \definition{to shine, flash}    \example{
      \exsen{Yesu bo pätt a wa iddpo da gombenmällnegnän.}      \extran{Jesus' body and clothes shined.}}}  \sense{\textbf{2.} 
    \definition{shine, flash}    \example{
      \exsen{Gälas ulle alle mullae dan utale we benmäll a ballän.}      \extran{With a big mirror, it is possible for a flash of light to travel far away.}}}}

\entry{Bensi}{\headword{Bensi}
  \pos{Proper noun}  \sense{
    \definition{female personal name}}}

\entry{Benson}{\headword{Benson}  \pronnote{benson}
  \pos{Proper noun}  \sense{
    \definition{male personal name}}}

\entry{Benta}{\headword{Benta}  \pronnote{benta}
  \pos{Proper noun}  \sense{
    \definition{female personal name}}}

\entry{Ber}{\headword{Ber}  \note{Place names}  \pronnote{ber}
  \pos{Proper noun}  \sense{
    \definition{Ber (toponym)}    \note{Place names}    \example{
}    \example{
}    \example{
}}}

\entry{beräberäl}{\headword{beräberäl}  \pronnote{berəberəl}
  \pos{Transitive coverb}  \sense{
    \definition{to swing}    \example{
      \exsen{Ingnenang aba pite da mermerae abal beräberäl agnegnan.}      \extran{The grass skirts of the dancers swung very nicely.}}}}

\entry{Beradi}{\headword{Beradi}  \note{Places around Limol}  \pronnote{beradi}
  \pos{Proper noun}  \sense{
    \definition{Beradi (toponym)}    \note{Places around Limol}    \example{
}    \example{
}    \example{
}}}

\entry{Bes}{\headword{Bes}
  \pos{Proper noun}  \sense{
    \definition{female personal name}}}

\entry{Bessie}{\headword{Bessie}  \pronnote{bessie}
  \pos{Proper noun}  \sense{
    \definition{female personal name}}}

\entry{Betliem}{\headword{Betliem}  \pronnote{betliem}
  \pos{Proper noun}  \sense{
    \definition{Bethlehem}}}

\entry{bette}{\headword{bette}  \pronnote{beʈ͡ʂe}
  \pos{Noun}  \sense{
    \definition{crimson finch}    \scientificname{Neochmia phaeton}    \example{
}    \example{
}    \example{
}}}

\entry{Bewag}{\headword{Bewag}  \pronnote{bewaɡ}
  \pos{Proper noun}  \sense{
    \definition{male personal name}}}

\entry{beya}{\headword{beya}
  \variantof{Dialectal Variant of \vartext{biye}}  \pronnote{beja}}

\entry{beyambäg}{\headword{beyambäg}  \pronnote{bejambəɡ}
  \pos{Transitive verb}  \sense{
    \definition{to chase}    \example{
      \exsen{Ngäna sɨmell beyambägag dan.}      \extran{I am a pig-chaser.}}    \example{
      \exsen{Yäbäd imneimne gumbiebägän ddone dingmenän.}      \extran{Sun pursued from behind, but did not reach him.}}}}

\entry{beyat}{\headword{beyat}  \note{Trees}  \pronnote{bejat}
  \pos{Noun}  \sense{
    \definition{type of tree that grows in the bush with wood used for house sticks}    \note{Trees}    \example{
}    \example{
}    \example{
}}}

\entry{beyawa}{\headword{beyawa}  \pronnote{bejawa}
  \pos{Personal pronoun}  \sense{
    \definition{he, she (emphatic third person singular pronoun)}    \example{
      \exsen{Beyawa dattkoeyän turik alle.}      \extran{He was the one who chopped it with an ax.}}}
  \variants{\SpellingVariant \vartext{beawa}}}

\entry{beyawaebe}{\headword{beyawaebe}
  \pos{Personal pronoun}  \sense{
    \definition{resctrictive form of beyawa}    \example{
      \exsen{Beyawaebe sana dägdene.}      \extran{Only he beat the sago.}}}}

\entry{beyawaeben}{\headword{beyawaeben}  \pronnote{bejawaibɛn}
  \pos{Copular verb}  \sense{
    \definition{copular form of beyawaebe}    \example{
      \exsen{Adi da ttongdae dan a ttongo ako ddone dan, be beyawaeben.}      \extran{God is the one and there is no other; only He is.}}}}

\entry{beyawaenen}{\headword{beyawaenen}  \pronnote{beawaenen}
  \pos{Copular verb}  \sense{
    \definition{copular form of beyawa (present form)}    \example{
      \exsen{Sisri ngäna erame giddollnen allan? Gänyme, ge ttängäm a beyawaenen.}      \extran{Where am I living now? Here, this village is where.}}}
  \variants{\SpellingVariant \vartext{beyawainin, beyawainen, beawaenen}}}

\entry{beyawaeneya}{\headword{beyawaeneya}
  \pos{Copular verb}  \sense{
    \definition{past form of beyawaenen}    \example{
      \exsen{Ngämo mäg bo bin a beyawaeneya Sara.}      \extran{Sara, it was my mother's name.}}}}

\entry{beyawainen}{\headword{beyawainen}
  \variantof{SpellingVariant of \vartext{beyawaenen}}}

\entry{beyawainin}{\headword{beyawainin}
  \variantof{SpellingVariant of \vartext{beyawaenen}}}

\entry{beyawän}{\headword{beyawän}
  \variantof{Unspecified Variant of \vartext{beyawan}}}

\entry{beyawan}{\headword{beyawan}
  \pos{Copular verb}  \sense{
    \definition{present copular form of beyawa}    \example{
      \exsen{Ngämo bin a beyawan Merol.}      \extran{Merol, it is my name.}}}
  \variants{\Unspecified Variant \vartext{beyawän}}}

\entry{=bi}{\headword{=bi}  \pronnote{bi}
  \pos{Pronominal enclitic}  \sense{
    \definition{clitic form of ubi}    \example{
      \exsen{Nensi bi nane alle känaebag mɨnyi tudi ma beyareyo.}      \extran{Nancy will go fishing tomorrow with her aunt.}}    \example{
      \exsen{Medäda bi ddone kili dägaebeyo.}      \extran{Their fathers were not happy.}}}}

\entry{bib}{\headword{bib}  \note{Gardens}  \pronnote{bib}
  \pos{Intransitive coverb}  \sense{
    \definition{to break ground, surface}    \note{Gardens}    \example{
      \exsen{Mätta da zime bib gognegän.}      \extran{The yams have already surfaced.}}    \example{
}    \example{
}}}

\entry{bib}{\headword{bib}  \pronnote{bib}
  \pos{Noun}  \sense{
    \definition{spring water}}}

\entry{bibi}{\headword{bibi}  \pronnote{bibi}
  \pos{Personal pronoun}  \sense{
    \definition{you all, you (second person nonsingular pronoun, nominative form)}    \example{
      \exsen{Bibi utamom ngattong.}      \extran{You all go first.}}}}

\entry{Bibiae}{\headword{Bibiae}  \pronnote{bibiae}
  \pos{Proper noun}  \sense{
    \definition{female personal name}}
  \variants{\SpellingVariant \vartext{Bibiai}}}

\entry{Bibiai}{\headword{Bibiai}
  \variantof{SpellingVariant of \vartext{Bibiae}}  \pronnote{bibiai}}

\entry{bibim}{\headword{bibim}  \pronnote{bibim}
  \pos{Personal pronoun}  \sense{
    \definition{accusative form of bibi}    \example{
      \exsen{Ngäna bibim truminy anggan.}      \extran{I am calling you all.}}}}

\entry{bibira}{\headword{bibira}
  \variantof{Dialectal Variant of \vartext{bibra}}}

\entry{bible}{\headword{bible}  \note{Types of Taro}  \pronnote{bible}
  \pos{Noun}  \sense{
    \definition{type of big taro}    \note{Types of Taro}    \example{
}    \example{
}    \example{
}}}

\entry{Bible}{\headword{Bible}
  \variantof{SpellingVariant of \vartext{baebol}}}

\entry{bibol}{\headword{bibol}  \note{Birds}  \pronnote{bibol}
  \pos{Noun}  \sense{
    \definition{type of bird}    \note{Birds}    \example{
}    \example{
}    \example{
}}}

\entry{biboz}{\headword{biboz}  \pronnote{biboz}
  \pos{Noun}  \sense{
    \definition{fairywren (emperor, white-shouldered)}    \scientificname{Malurus cyanocephalus, M. alboscapulatus}}}

\entry{bibra}{\headword{bibra}  \pronnote{bibra}
  \pos{Personal pronoun}  \sense{
    \definition{dative form of bibi}    \example{
      \exsen{Känazbag, ngäna mɨnyi bibra sisor ttoenttoen de bɨllɨt.}      \extran{Tomorrow, I will tell you all a new story.}}}
  \variants{\Dialectal Variant \vartext{bibira}}}

\entry{Bidog}{\headword{Bidog}  \pronnote{bidoɡ}
  \pos{Proper noun}  \sense{
    \definition{male personal name}}}

\entry{big}{\headword{big}  \note{Trees}  \pronnote{biɡ}
  \pos{Noun}  \sense{
    \definition{type of very large tree that grows in the bush with wood used for firewood, especially when camping.}    \note{Trees}    \anthronote{Used in a special ceremony involving a poisoned shirt, a long stick, and a sick person.}    \example{
}    \example{
}    \example{
}}}

\entry{Big}{\headword{Big}
  \pos{Proper noun}  \sense{
    \definition{Big (toponym)}}}

\entry{big budar}{\headword{big budar}  \note{Grubs}  \pronnote{biɡ budar}
  \pos{Noun}  \sense{
    \definition{type of grub}    \note{Grubs}    \example{
}    \example{
}    \example{
}}}

\entry{Bigag}{\headword{Bigag}
  \pos{Proper noun}  \sense{
    \definition{male personal name}}}

\entry{Bigia}{\headword{Bigia}
  \variantof{SpellingVariant of \vartext{Bigiya}}}

\entry{Bigiya}{\headword{Bigiya}
  \pos{Proper noun}  \sense{
    \definition{female personal name}}
  \variants{\SpellingVariant \vartext{Bigia}}}

\entry{Bigjay}{\headword{Bigjay}  \pronnote{biɡd͡ʒe}
  \pos{Proper noun}  \sense{
    \definition{male personal name}}}

\entry{bigma}{\headword{bigma}  \pronnote{biɡma}
  \pos{Noun}  \sense{
    \definition{enclosure, pen, sty}    \example{
      \exsen{Bong angde bigma we dallän, sɨmell a dindugän.}      \extran{When Bong went to the pen, the pig had escaped.}}}}

\entry{bikme}{\headword{bikme}  \lexeme{bikme}  \pronnote{bikme}
  \pos{Noun}  \sense{
    \definition{type of palm tree with hanging, poisonous yellow and green fruits that can be eaten after being buried by the creek for up to 2 years and then cooked on the fire}}}

\entry{bikme laelem}{\headword{bikme laelem}  \pronnote{bikme laelem}
  \pos{Noun}  \sense{
    \definition{type of bikme palm with very chewy fruit}}}

\entry{bikme pipllo}{\headword{bikme pipllo}  \pronnote{bikme pipɽo}
  \pos{Noun}  \sense{
    \definition{type of bikme palm with yellow fruit that is very hard}}}

\entry{bikme sanasana}{\headword{bikme sanasana}  \pronnote{bikme sanasana}
  \pos{Noun}  \sense{
    \definition{type of bikme palm with fruit that is soft like sago}}}

\entry{bikme tutu}{\headword{bikme tutu}  \note{Games}  \pronnote{bikme tutu}
  \pos{Noun}  \sense{
    \definition{type of string game}    \note{Games}    \example{
}    \example{
}    \example{
}}}

\entry{bikme yobeg}{\headword{bikme yobeg}  \pronnote{bikme jobeɡ}
  \pos{Noun}  \sense{
    \definition{type of bikme palm with fruit that is medium-firm}}}

\entry{Biks}{\headword{Biks}
  \pos{Proper noun}  \sense{
    \definition{personal name}}}

\entry{Biku}{\headword{Biku}  \pronnote{biku}
  \pos{Proper noun}  \sense{
    \definition{male personal name}}}

\entry{bikwem}{\headword{bikwem}  \note{Everything in the house}  \pronnote{bikwem}
  \pos{Noun}  \sense{
    \definition{fireplace}    \note{Everything in the house}    \example{
      \exsen{Ibi ibra mɨnyi ttongo sisor bikwem de ako bangeseya.}      \extran{We will also make a new fireplace for us.}}    \example{
}    \example{
}}}

\entry{bile}{\headword{bile}
  \pos{Noun}  \sense{
    \definition{salt}}}

\entry{bilod}{\headword{bilod}  \pronnote{bilod}
  \pos{Noun}  \sense{
    \definition{type of spear}}}

\entry{=bim}{\headword{=bim}  \pronnote{bim}
  \pos{Pronominal enclitic}  \sense{
    \definition{clitic form of ubim}    \example{
      \exsen{Ngäna ngämo masar bim ddone ikop dägnegne.}      \extran{I did not see my grandparents.}}}}

\entry{Bimadbn}{\headword{Bimadbn}  \pronnote{bimadbn}
  \pos{Proper noun}  \sense{
    \definition{Bimadbn (toponym)}    \example{
}    \example{
}    \example{
}}}

\entry{bin}{\headword{bin}  \pronnote{bin}
  \pos{Noun}  \sense{
    \definition{name}    \example{
      \exsen{Ngämo mälla bo bin a Kristina.}      \extran{My wife's name is Kristina.}}}}

\entry{bina}{\headword{bina}  \pronnote{bina}
  \pos{Personal pronoun}  \sense{
    \definition{possessive form of bibi}    \example{
      \exsen{Ge bina ma daeya?}      \extran{Is this your (pl.) house?}}}}

\entry{binaene}{\headword{binaene}  \pronnote{binaene}
  \pos{Personal pronoun}  \sense{
    \definition{ablative-possessive form of bibi}    \example{
      \exsen{Lla da binaene eka de mɨnyi bondärmällnegnän.}      \extran{People will be listening to your (pl.) words.}}}}

\entry{binang}{\headword{binang}  \pronnote{binaŋ}
  \pos{Noun}  \sense{
    \definition{namesake}}}

\entry{binang}{\headword{binang}
  \pos{Modifier}  \sense{
    \definition{serious}    \example{
      \exsen{ulle binang itrell peyang lla}      \extran{seriously ill person}}}}

\entry{binbäddbädd}{\headword{binbäddbädd}  \pronnote{binbəɖ͡ʐbəɖ͡ʐ}
  \pos{Adverb}  \sense{
    \definition{fully, completely}}}

\entry{Bine}{\headword{Bine}  \note{Language Names}  \pronnote{bine}
  \pos{Proper noun}  \sense{
    \definition{Bine language (spoken to the east)}    \note{Language Names}    \example{
}    \example{
}    \example{
}}}

\entry{Binyomoll Källäm}{\headword{Binyomoll Källäm}  \note{Places around Limol}  \pronnote{biɲomoɽ kəɽəm}
  \pos{Proper noun}  \sense{
    \definition{Binyomoll Pond (on the road to Kinkin)}    \note{Places around Limol}    \example{
}    \example{
}    \example{
}}}

\entry{binzeg}{\headword{binzeg}  \pronnote{binzeɡ}
  \pos{Transitive verb}  \sense{
    \definition{to heat, warm}}}

\entry{Bipi}{\headword{Bipi}
  \pos{Proper noun}  \sense{
    \definition{Bipi (toponym)}}}

\entry{bir}{\headword{bir}  \pronnote{bir}
  \pos{Noun}  \sense{\textbf{1.} 
    \definition{spit, skewer}}  \sense{\textbf{2.} 
    \definition{to roast}    \example{
      \exsen{Ngatengate da biratt a ddobae moko dan.}      \extran{Roasted possum is very tasty.}}    \example{
}    \example{
}}}

\entry{Birigi}{\headword{Birigi}
  \pos{Proper noun}  \sense{
    \definition{male personal name}}}

\entry{biro}{\headword{biro}
  \pos{Noun}  \sense{
    \definition{pen (writing implement)}}}

\entry{bisbis}{\headword{bisbis}  \note{Bees}  \pronnote{bisbis}
  \pos{Noun}  \sense{
    \definition{type of stingless bee}    \note{Bees}    \example{
}    \example{
}    \example{
}}}

\entry{bisel}{\headword{bisel}  \note{Sago palms}  \pronnote{bisel}
  \pos{Noun}  \sense{
    \definition{type of sago that grows tall and wide}    \note{Sago palms}    \example{
}    \example{
}    \example{
}}}

\entry{Bisenmo}{\headword{Bisenmo}
  \pos{Proper noun}  \sense{
    \definition{Bisenmo (toponym)}}}

\entry{bisnis}{\headword{bisnis}
  \pos{Noun}  \sense{
    \definition{business}}}

\entry{Bisuaka}{\headword{Bisuaka}  \pronnote{bisuaka}
  \pos{Proper noun}  \sense{
    \definition{Bisuaka (Bituri-speaking village in Oriomo-Bituri Rural LLG; has a primary school but no aid post)}    \example{
}    \example{
}    \example{
}}
  \variants{\SpellingVariant \vartext{Biswaka}}}

\entry{Biswaka}{\headword{Biswaka}
  \variantof{SpellingVariant of \vartext{Bisuaka}}}

\entry{bitän}{\headword{bitän}  \note{Animal names}  \pronnote{bitən}
  \pos{Noun}  \sense{
    \definition{type of animal}    \note{Animal names}    \example{
}    \example{
}    \example{
}}}

\entry{bittang}{\headword{bittang}
  \pos{Transitive coverb}  \sense{
    \definition{to litter, dirty}}}

\entry{bittott}{\headword{bittott}  \note{Animal names}  \pronnote{biʈ͡ʂoʈ͡ʂ}
  \pos{Noun}  \sense{
    \definition{grey squirrel}    \note{Animal names}    \example{
}    \example{
}    \example{
}}}

\entry{Bitur}{\headword{Bitur}  \note{Language Names}  \pronnote{bitur}
  \pos{Proper noun}  \sense{
    \definition{Bitur language (spoken to the north)}    \note{Language Names}    \example{
}    \example{
}    \example{
}}}

\entry{biwiz}{\headword{biwiz}  \note{Trees}  \pronnote{biwiz}
  \pos{Noun}  \sense{
    \definition{type of large tree that grows in the bush with purple flowers and wood used for kundu drums and canoes}    \note{Trees}    \example{
}    \example{
}    \example{
}}}

\entry{biye}{\headword{biye}  \pronnote{bije}
  \pos{Noun}  \sense{
    \definition{taro}    \scientificname{Colocasia}    \example{
      \exsen{Biye de bongo ai dan ibik alle nibe.}      \extran{You can plant taro with the ibik.}}}
  \variants{\Dialectal Variant \vartext{beya}}}

\entry{Biyewolatt}{\headword{Biyewolatt}  \note{Places around Limol}  \pronnote{bijewolaʈ͡ʂ}
  \pos{Proper noun}  \sense{
    \definition{Biyewolatt (toponym)}    \note{Places around Limol}    \example{
}    \example{
}    \example{
}}}

\entry{bɨd}{\headword{bɨd}  \pronnote{bɪd}
  \pos{Noun}  \sense{
    \definition{gum tree}}}

\entry{bɨk}{\headword{bɨk}  \pronnote{bɪk}
  \pos{Noun}  \sense{
    \definition{poisoned creek}}}

\entry{bɨt}{\headword{bɨt}  \pronnote{bɪt}
  \pos{Modifier}  \sense{
    \definition{dark}}
  \variants{\Unspecified Variant \vartext{bät}}}

\entry{bɨtbɨt}{\headword{bɨtbɨt}  \note{Colors}  \pronnote{bɪtbɪt}
  \pos{Color term}  \sense{
    \definition{black}    \note{Colors}    \example{
      \exsen{Obo ttoe a bɨtbɨt dan.}      \extran{His skin is black.}}    \example{
}    \example{
}}
  \variants{\SpellingVariant \vartext{bätbät}}}

\entry{bɨtbɨtbɨtbɨt}{\headword{bɨtbɨtbɨtbɨt}  \pronnote{bɪtbɪtbɪtbɪt}
  \pos{Color term}  \sense{
    \definition{purple (color of sawis, purple yam)}}}

\entry{blab}{\headword{blab}  \note{Stages of life}  \pronnote{blab}
  \pos{Intransitive verb}  \sense{
    \definition{to mature, reach puberty}    \note{Stages of life}    \example{
      \exsen{Ge nge da ablaban.}      \extran{This coconut has matured.}}    \example{
      \exsen{Män a zäme ablaban.}      \extran{The girl has already reached puberty.}}    \example{
}}}

\entry{blame}{\headword{blame}
  \variantof{freevarlab of \vartext{English loanword}}}

\entry{blengud}{\headword{blengud}  \pronnote{blɛŋɡud}
  \pos{Noun}  \sense{
    \definition{blanket}}}

\entry{bllablla}{\headword{bllablla}  \lexeme{bllablla}  \pronnote{bɽabɽa}
  \pos{Noun}  \sense{
    \definition{type of cordyline with big leaves, red fruit, and leaves that are used to fan fire and tied around the waist and chest for dancing}    \scientificname{Cordyline fruticosa}}}

\entry{bllablle}{\headword{bllablle}
  \variantof{Fast Speech Variant of \vartext{bällabälle}}}

\entry{blläg}{\headword{blläg}  \pronnote{bɽəɡ}
  \pos{Transitive verb}  \sense{
    \definition{to serve}    \example{
      \exsen{Pilat de bäbllägaebeyo.}      \extran{They will serve plates.}}}}

\entry{bllam}{\headword{bllam}
  \variantof{Fast Speech Variant of \vartext{bällam}}}

\entry{bllolla}{\headword{bllolla}  \note{Trees}  \pronnote{bɽoɽa}
  \pos{Noun}  \sense{
    \definition{type of tree}    \note{Trees}    \example{
}    \example{
}    \example{
}}}

\entry{=bo}{\headword{=bo}  \pronnote{bo}
  \pos{Pronominal enclitic}  \sense{
    \definition{clitic form of obo}    \example{
      \exsen{Ngämo yae bo bin a Nowar.}      \extran{My mother's name is Nowar.}}}}

\entry{bob}{\headword{bob}  \lexeme{bob}  \pronnote{bob}
  \pos{Noun}  \sense{
    \definition{flood}    \example{
      \exsen{Tätäm ttäle bun parga de bob a era dällekän.}      \extran{Yesterday, the flood destroyed the Ttäle Bun bridge.}}}}

\entry{bob}{\headword{bob}  \pronnote{bob}
  \pos{Noun}  \sense{
    \definition{type of tree}}}

\entry{bob lläkäm}{\headword{bob lläkäm}  \note{Mushrooms}  \pronnote{bob ɽəkəm}
  \pos{Noun}  \sense{
    \definition{type of mushroom}    \note{Mushrooms}    \example{
}    \example{
}}}

\entry{bobag}{\headword{bobag}  \pronnote{bobaɡ}
  \pos{Noun}  \sense{
    \definition{flooded}    \example{
      \exsen{Tatuma da bobag daeya.}      \extran{The bathing area was flooded.}}}}

\entry{Bobe}{\headword{Bobe}
  \pos{Proper noun}  \sense{
    \definition{Bobe (toponym)}}}

\entry{bobllem}{\headword{bobllem}
  \pos{Intransitive verb}  \sense{
    \definition{to flap, shake (of skin)}    \example{
      \exsen{Mälla bo pätt a bobllem allan.}      \extran{The skin of the woman's body shakes.}}}}

\entry{bobngätt}{\headword{bobngätt}  \note{Types of Taro}  \pronnote{bobŋəʈ͡ʂ}
  \pos{Noun}  \sense{
    \definition{place name}    \note{Types of Taro}}}

\entry{Bobngätt}{\headword{Bobngätt}  \pronnote{bobŋəʈ͡ʂ}
  \pos{Proper noun}  \sense{
    \definition{Bobngätt (toponym)}    \example{
}    \example{
}    \example{
}}}

\entry{Bobzag}{\headword{Bobzag}  \pronnote{bobzaɡ}
  \pos{Proper noun}  \sense{\textbf{1.} 
    \definition{personal name}}  \sense{\textbf{2.} 
    \definition{dog name}}}

\entry{Bobze}{\headword{Bobze}
  \pos{Proper noun}  \sense{
    \definition{Bobze (toponym)}}}

\entry{bod}{\headword{bod}  \pronnote{bod}
  \pos{Noun}  \sense{\textbf{1.} 
    \definition{mouth}    \example{
      \exsen{Ngäna mɨnyi mälläng ik dae beyazen a bod ik dae bigezän.}      \extran{I will put it in through your nostrils and take it out through your mouth.}}    \example{
      \exsen{Mällaeyaba bod me dan.}      \extran{It's on the women's lips (i.e. the subject of gossip).}}}  \sense{\textbf{2.} 
    \definition{beak}    \example{
      \exsen{Bod a malla ddäddäg ma dan.}      \extran{The beak is not edible.}}    \example{
}    \example{
}}  \sense{\textbf{3.} 
    \definition{lip of bag (braid that goes along the rim of the bag to finish it)}}}

\entry{bod ttoe}{\headword{bod ttoe}
  \pos{Noun}  \sense{
    \definition{lip}}}

\entry{boddo}{\headword{boddo}  \note{Trees}  \pronnote{boɖ͡ʐo}
  \pos{Noun}  \sense{
    \definition{type of large tree with a big trunk, fruit that deer eat, and aerial prop roots}    \note{Trees}    \example{
}    \example{
}    \example{
}}}

\entry{bode}{\headword{bode}  \pronnote{bode}
  \pos{Personal pronoun}  \sense{
    \definition{additive form of bogo}    \example{
      \exsen{Bode ngämlle pate Ende eka walle eka gognän.}      \extran{He also spoke to me in Ende.}}}}

\entry{bodo}{\headword{bodo}  \note{Adjectives}  \pronnote{bodo}
  \pos{Modifier}  \sense{
    \definition{full}    \note{Adjectives}    \example{
      \exsen{Kap a bodo dan.}      \extran{The cup is full.}}    \example{
      \exsen{Ngäna ine de zämaeangae nägawan do angde kap a bodo agan.}      \extran{I poured water until the cup was full.}}    \example{
}}}

\entry{bodobodo}{\headword{bodobodo}
  \pos{Modifier}  \sense{
    \definition{nonsingular form of bodo}}}

\entry{bodobodom}{\headword{bodobodom}  \lexeme{bodobodom}  \pronnote{bodobodom}
  \pos{Noun}  \sense{
    \definition{type of biting ant}}}

\entry{Bodog}{\headword{Bodog}  \pronnote{bodoɡ}
  \pos{Proper noun}  \sense{
    \definition{male personal name}}}

\entry{boe}{\headword{boe}  \lexeme{boe}  \pronnote{boe}
  \pos{Noun}  \sense{
    \definition{type of cultivated tree with edible indigo fruit and white flowers that attract birds and butterflies}}}

\entry{boeb}{\headword{boeb}  \pronnote{boeb}
  \pos{Verb}  \sense{
    \definition{to deny}}}

\entry{bog}{\headword{bog}  \note{Types of Taro}  \pronnote{boɡ}
  \pos{Noun}  \sense{
    \definition{type of taro}    \note{Types of Taro}    \example{
}    \example{
}    \example{
}}}

\entry{boga}{\headword{boga}  \note{Birds}  \pronnote{boɡa}
  \pos{Noun}  \sense{
    \definition{type of bird}    \note{Birds}    \example{
}    \example{
}    \example{
}}}

\entry{boge}{\headword{boge}  \lexeme{boge}  \pronnote{boɡe}
  \pos{Noun}  \sense{
    \definition{mudfish}    \example{
      \exsen{Boge kollba da mer mokowang ddäddäg dan.}      \extran{Mudfish is a very delicious fish.}}}}

\entry{bogel}{\headword{bogel}  \lexeme{bogel}  \pronnote{boɡel}
  \pos{Noun}  \sense{
    \definition{seaweed}}}

\entry{bogo}{\headword{bogo}  \pronnote{boɡo}
  \pos{Personal pronoun}  \sense{
    \definition{he, she (third person singular animate pronoun, nominative form)}    \example{
      \exsen{Bogo daden ma me?}      \extran{Is he home?}}}}

\entry{bogobogo}{\headword{bogobogo}  \lexeme{bogobogo}  \pronnote{boɡoboɡo}
  \pos{Noun}  \sense{
    \definition{type of small yam with a pure white interior}}}

\entry{Bok}{\headword{Bok}  \pronnote{bok}
  \pos{Proper noun}  \sense{
    \definition{Buk (Taeme-speaking village in Morehead Rural LLG; from Limol, one must pass through Kinkin and Kondobol)}    \example{
}    \example{
}    \example{
}}}

\entry{boko}{\headword{boko}  \lexeme{boko}  \pronnote{boko}
  \pos{Noun}  \sense{
    \definition{type of lizard}    \example{
      \exsen{Llɨg kälekäle da boko de näbäddallo sisri.}      \extran{Small children kill boko lizards.}}}}

\entry{boko}{\headword{boko}
  \pos{Personal pronoun}  \sense{
}}

\entry{bol}{\headword{bol}  \pronnote{bol}
  \pos{Noun}  \sense{
    \definition{ball}}
  \variants{\Unspecified Variant \vartext{boll}}}

\entry{boll}{\headword{boll}
  \variantof{Unspecified Variant of \vartext{bol}}}

\entry{bollboll}{\headword{bollboll}  \pronnote{boɽboɽ}
  \pos{Transitive verb}  \sense{
    \definition{to open sago}    \example{
      \exsen{Bogo ako ttongo pättkäp däbollän gädagäde we.}      \extran{He opened another node of the sago trunk to beat.}}}}

\entry{bollga}{\headword{bollga}  \note{Sago palms}  \pronnote{boɽɡa}
  \pos{Noun}  \sense{
    \definition{type of sago}    \note{Sago palms}    \example{
}    \example{
}    \example{
}}}

\entry{Bollga}{\headword{Bollga}
  \pos{Proper noun}  \sense{
    \definition{female personal name}}}

\entry{Bolloll}{\headword{Bolloll}  \note{Places around Limol}  \pronnote{boɽoɽ}
  \pos{Proper noun}  \sense{
    \definition{Bolloll (toponym, on the southward road to Malam)}    \note{Places around Limol}    \example{
}    \example{
}    \example{
}}}

\entry{bolod}{\headword{bolod}  \pronnote{bolod}
  \pos{Noun}  \sense{
    \definition{sugarcane}    \scientificname{Saccharum edule}    \anthronote{Used to make pakos (spear). Dry on fire, straighten it, push arrow on, and tie with tulip bark (~2m).}}
  \variants{\Unspecified Variant \vartext{bälwäd, bälwod}}
  \variants{\SpellingVariant \vartext{bolwod}}}

\entry{bolwod}{\headword{bolwod}
  \variantof{SpellingVariant of \vartext{bolod}}  \pronnote{bolwod}}

\entry{=bom}{\headword{=bom}  \pronnote{bom}
  \pos{Pronominal enclitic}  \sense{
    \definition{clitic form of obom}    \example{
      \exsen{Ngäna Wagiba bom dangmingg.}      \extran{I helped Wagiba.}}}}

\entry{boma}{\headword{boma}
  \pos{Noun}  \sense{
}}

\entry{bomall}{\headword{bomall}  \note{Trees}  \pronnote{bomaɽ}
  \pos{Noun}  \sense{
    \definition{type of tree that grows in the bush with bark used to make sago baskets}    \note{Trees}    \example{
}    \example{
}    \example{
}}}

\entry{bombllo}{\headword{bombllo}  \pronnote{bombɽo}
  \pos{Transitive/Intransitive coverb}  \sense{\textbf{1.} 
    \definition{to increase, proliferate, multiply}    \example{
      \exsen{Ngäma lla da ada ingollang gombllowän.}      \extran{Our (excl.) population grew like this.}}}  \sense{\textbf{2.} 
    \definition{to increase, proliferate, multiply}    \example{
      \exsen{Mätta da bombllo allan.}      \extran{The yams proliferate.}}    \example{
      \exsen{Ngäna komlla lla de bombllo yaran.}      \extran{My two daughters each married into a family (lit. I multiplied two people; i.e. each daughter will now have a family of her own).}}}}

\entry{bombom}{\headword{bombom}  \note{Trees}  \pronnote{bombom}
  \pos{Noun}  \sense{
    \definition{type of tree with big leaves and soft wood that floats and is carved by children}    \note{Trees}    \example{
}    \example{
}    \example{
}}}

\entry{bomo}{\headword{bomo}
  \pos{Noun}  \sense{
    \definition{aerial root (e.g. of pandanus)}}}

\entry{Bomso}{\headword{Bomso}  \pronnote{bomso}
  \pos{Proper noun}  \sense{
    \definition{male personal name}}}

\entry{Bong}{\headword{Bong}
  \pos{Proper noun}  \sense{
    \definition{male personal name}}}

\entry{bongo}{\headword{bongo}  \pronnote{boŋo}
  \pos{Personal pronoun}  \sense{
    \definition{you (second person singular pronoun, nominative form)}    \example{
      \exsen{Bongo erowe ibi alle?}      \extran{Where are you going?}}}}

\entry{Bonibi}{\headword{Bonibi}  \pronnote{bonibi}
  \pos{Proper noun}  \sense{
    \definition{female personal name}}}

\entry{Bonybony}{\headword{Bonybony}  \note{Places around Limol}  \pronnote{boɲboɲ}
  \pos{Proper noun}  \sense{
    \definition{Bonybony (camping, garden, and sago place (AX94))}    \note{Places around Limol}    \example{
}    \example{
}    \example{
}}}

\entry{bonydre}{\headword{bonydre}  \pronnote{boɲdre}
  \pos{Noun}  \sense{
    \definition{goshawk (grey-headed, brown); collared sparrowhawk}    \scientificname{Tachyspiza poliocephala, T. fasciata dogwa; T. cirrocephala}}
  \variants{\Unspecified Variant \vartext{buindre}}}

\entry{bonzro}{\headword{bonzro}  \note{Dancing}  \pronnote{bonzro}
  \pos{Noun}  \sense{
    \definition{dancing on the side playfully}    \note{Dancing}    \example{
}    \example{
}    \example{
}}}

\entry{bor}{\headword{bor}
  \pos{Noun}  \sense{
}}

\entry{borale}{\headword{borale}  \pronnote{borale}
  \pos{Noun}  \sense{
    \definition{traditional bamboo flute used to scare wallabies}    \example{
      \exsen{Borale da eran ttongo ekaeka ma za dan.}      \extran{A flute is a thing that makes noise.}}}
  \variants{\Dialectal Variant \vartext{boralle}}}

\entry{boralle}{\headword{boralle}
  \variantof{Dialectal Variant of \vartext{borale}}  \pronnote{boraɽe}}

\entry{bore}{\headword{bore}
  \pos{Noun}  \sense{
    \definition{traditional bamboo pipe for smoking tobacco}}}

\entry{bormop}{\headword{bormop}  \note{Traditional Arrows}  \pronnote{bormop}
  \pos{Noun}  \sense{
    \definition{type of spear}    \note{Traditional Arrows}    \example{
}    \example{
}    \example{
}}}

\entry{boser}{\headword{boser}  \lexeme{boser}  \pronnote{boser}
  \pos{Noun}  \sense{
    \definition{rock found in creek}}}

\entry{botol}{\headword{botol}
  \variantof{Dialectal Variant of \vartext{English loanword}}}

\entry{bott}{\headword{bott}
  \pos{Noun}  \sense{
    \definition{boat}}}

\entry{botta}{\headword{botta}  \pronnote{boʈ͡ʂa}
  \pos{Noun}  \sense{
    \definition{lateral beam placed directly on the house post}    \example{
      \exsen{Ma da botta käme dämenang da.}      \extran{The house already has the lateral beam.}}}}

\entry{bowe}{\headword{bowe}  \pronnote{bowe}
  \pos{Noun}  \sense{
    \definition{side}}}

\entry{Boze}{\headword{Boze}  \pronnote{boze}
  \pos{Proper noun}  \sense{
    \definition{Boze (Agob- and Bine-speaking village in Oriomo-Bituri Rural LLG; from Limol, one must pass through Malam, Kurunti, and Kibuli)}    \example{
}    \example{
}    \example{
}}}

\entry{Bozorob}{\headword{Bozorob}  \pronnote{bozorob}
  \pos{Proper noun}  \sense{
    \definition{Bozorob (camping place on the road to Kurunti; from Limol, one must pass through Malam)}    \example{
}    \example{
}    \example{
}}
  \variants{\Fast Speech Variant \vartext{Bozrob}}}

\entry{Bozrob}{\headword{Bozrob}
  \variantof{Fast Speech Variant of \vartext{Bozorob}}}

\entry{Breton}{\headword{Breton}  \pronnote{brɛton}
  \pos{Proper noun}  \sense{
    \definition{male personal name}}}

\entry{buata}{\headword{buata}  \pronnote{buata}
  \pos{Noun}  \sense{
    \definition{betel nut, areca nut (fruit of Areca catechu)}    \example{
      \exsen{buata däg}      \extran{betel nut bunch}}    \example{
}    \example{
}}}

\entry{budar}{\headword{budar}  \note{Grubs}  \pronnote{budar}
  \pos{Noun}  \sense{
    \definition{grub, larva}    \note{Grubs}    \example{
}    \example{
}    \example{
}}}

\entry{buddo}{\headword{buddo}  \pronnote{buɖ͡ʐo}
  \pos{Noun}  \sense{\textbf{1.} 
    \definition{problem, issue}    \example{
      \exsen{Mälla papa da buddo ulle dan Papua Niyu Gini mi.}      \extran{Beating one's wife is a big problem in Papua New Guinea.}}}  \sense{\textbf{2.} 
    \definition{weight}    \example{
      \exsen{Buddo da tumang da.}      \extran{It's very heavy (lit. the weight is a lot).}}}}

\entry{buddobuddo}{\headword{buddobuddo}  \note{Adjectives}  \pronnote{buɖ͡ʐobuɖ͡ʐo}
  \pos{Modifier}  \sense{
    \definition{nonsingular form of buddo}    \note{Adjectives}}}

\entry{buddobuddog}{\headword{buddobuddog}  \pronnote{buɖ͡ʐobuɖ͡ʐoɡ}
  \pos{Adverb}  \sense{
    \definition{carrying a load}    \example{
      \exsen{Dagwaeya deyarneyo buddobuddog.}      \extran{They (du.) were going with their loads.}}}}

\entry{buddog}{\headword{buddog}  \note{Adjectives}  \pronnote{buɖ͡ʐoɡ}
  \pos{Modifier}  \sense{\textbf{1.} 
    \definition{heavy}    \note{Adjectives}    \example{
      \exsen{Mälla da ine de buddog de wony eran.}      \extran{The woman is carrying the heavy water.}}    \example{
}    \example{
}}  \sense{\textbf{2.} 
    \definition{deep (of a voice)}    \note{Adjectives}}  \sense{\textbf{3.} 
    \definition{to trouble, bother}    \note{Adjectives}    \example{
      \exsen{Ewatta bibi obom buddog eralla?}      \extran{Why are you all bothering her?}}}}

\entry{budombudom}{\headword{budombudom}  \note{Poison}  \pronnote{budombudom}
  \pos{Noun}  \sense{
    \definition{red ants}    \note{Poison}    \example{
}    \example{
}    \example{
}}}

\entry{bugu}{\headword{bugu}  \pronnote{buɡu}
  \pos{Noun}  \sense{
    \definition{sheath, base, midrib (of a palm leaf)}    \example{
      \exsen{Nge da ddagemeny dan, obo bugu daeben.}      \extran{The coconut palm doesn't have branches, there are only its leaf sheaths.}}    \example{
      \exsen{Bugu dädär da aspunan.}      \extran{Dry palm leaf sheaths fall.}}}}

\entry{Buib}{\headword{Buib}  \pronnote{buib}
  \pos{Proper noun}  \sense{
    \definition{Buib (toponym)}}}

\entry{buidde}{\headword{buidde}  \note{Tools}  \pronnote{buiɖ͡ʐe}
  \pos{Noun}  \sense{
    \definition{club (weapon)}    \note{Tools}    \example{
}    \example{
}    \example{
}}}

\entry{Buiddobuiddog}{\headword{Buiddobuiddog}  \note{Places around Limol}  \pronnote{buiɖ͡ʐobuiɖ͡ʐoɡ}
  \pos{Proper noun}  \sense{
    \definition{Buiddobuiddog (sago place and creek; also a road (Top L AZ96))}    \note{Places around Limol}    \example{
}    \example{
}    \example{
}}}

\entry{buindre}{\headword{buindre}
  \variantof{Unspecified Variant of \vartext{bonydre}}  \pronnote{buindre}}

\entry{buitu}{\headword{buitu}  \note{Tools}  \pronnote{buitu}
  \pos{Noun}  \sense{
    \definition{stick with a round base}    \note{Tools}    \example{
}    \example{
}    \example{
}}}

\entry{buk}{\headword{buk}  \pronnote{buk}
  \pos{Noun}  \sense{
    \definition{book}    \example{
      \exsen{buk ulle}      \extran{big book}}}}

\entry{bullalla}{\headword{bullalla}
  \pos{Noun}  \sense{
    \definition{type of flower}}}

\entry{bullull}{\headword{bullull}  \pronnote{buɽuɽ}
  \pos{Noun}  \sense{\textbf{1.} 
    \definition{Papuan frogmouth}    \scientificname{Podargus papuensis}    \example{
}    \example{
}    \example{
}}  \sense{\textbf{2.} 
    \definition{rufous owl}    \scientificname{Ninox rufa}}}

\entry{bulwem}{\headword{bulwem}  \lexeme{bulwem}  \pronnote{bulwem}
  \pos{Noun}  \sense{
    \definition{type of big yam with a white and light purple interior and no hairs}}}

\entry{bumo}{\headword{bumo}
  \pos{Noun}  \sense{
    \definition{tucker bag}    \example{
      \exsen{Bogo gopällttänän obo wätät bumo peyang.}      \extran{She set off with her filled tucker bag.}}}}

\entry{bumrel}{\headword{bumrel}  \note{Poison}  \pronnote{bumrel}
  \pos{Noun}  \sense{
    \definition{beetle}    \note{Poison}    \example{
}    \example{
}    \example{
}}}

\entry{bun}{\headword{bun}  \lexeme{bun}  \pronnote{bun}
  \pos{Noun}  \sense{\textbf{1.} 
    \definition{head}    \example{
      \exsen{Pa bo bun a ddäddäg ma dan.}      \extran{The bird's head is not edible.}}}  \sense{\textbf{2.} 
    \definition{part of a bow}}  \sense{\textbf{3.} 
    \definition{mouth (of a river)}    \example{
      \exsen{Ngäna walle bun i ibi allan.}      \extran{I am going towards the river mouth.}}}  \sense{\textbf{4.} 
    \definition{owner}    \example{
      \exsen{Ginarang däbaolle eka kutt e bun dan.}      \extran{Ginarang is the owner of that word.}}}}

\entry{bun dänräp}{\headword{bun dänräp}
  \pos{Noun}  \sense{
}}

\entry{bun käp}{\headword{bun käp}
  \pos{Noun}  \sense{
    \definition{head}}}

\entry{bun kom}{\headword{bun kom}  \note{Parts of the Body}  \pronnote{bun kom}
  \pos{Noun}  \sense{
    \definition{hair (on head)}    \note{Parts of the Body}    \example{
}    \example{
}    \example{
}}
  \variants{\SpellingVariant \vartext{bunkom}}}

\entry{bun kutt}{\headword{bun kutt}  \note{Parts of the body}  \pronnote{bun kuʈ͡ʂ}
  \pos{Noun}  \sense{
    \definition{skull}    \note{Parts of the body}    \example{
}    \example{
}    \example{
}}
  \variants{\SpellingVariant \vartext{bunkutt}}}

\entry{bun lla}{\headword{bun lla}  \pronnote{bun ɽa}
  \pos{Noun}  \sense{
    \definition{instructor, leader}    \example{
      \exsen{Lla da mer bun lla dag.}      \extran{The men are good leaders.}}}}

\entry{bun pallkamatt}{\headword{bun pallkamatt}  \pronnote{bun paɽkamaʈ͡ʂ}
  \pos{Modifier}  \sense{
    \definition{smart}}}

\entry{bunbun}{\headword{bunbun}
  \pos{Noun}  \sense{
    \definition{type of plant}}}

\entry{Bundae}{\headword{Bundae}  \pronnote{bundae}
  \pos{Proper noun}  \sense{
    \definition{male personal name}}}

\entry{bungg}{\headword{bungg}
  \pos{Transitive verb}  \sense{
    \definition{to ambush}}}

\entry{bunkälle bunkälle}{\headword{bunkälle bunkälle}  \note{Games}  \pronnote{bunkəɽe bunkəɽe}
  \pos{Noun}  \sense{
    \definition{type of game where the players tie hair and hide}    \note{Games}    \example{
}    \example{
}    \example{
}}}

\entry{bunkom}{\headword{bunkom}
  \variantof{SpellingVariant of \vartext{bun kom}}}

\entry{bunkom tätäp}{\headword{bunkom tätäp}  \pronnote{bunkom tətəp}
  \pos{Noun}  \sense{
    \definition{hair tied with string}}}

\entry{bunkombunkom}{\headword{bunkombunkom}  \note{Trees}  \pronnote{bunkombunkom}
  \pos{Noun}  \sense{
    \definition{type of tall, big tree that grows along creek with wood used for canoes}    \note{Trees}    \example{
}    \example{
}    \example{
}}}

\entry{bunkutt}{\headword{bunkutt}
  \variantof{SpellingVariant of \vartext{bun kutt}}}

\entry{bunkuttang}{\headword{bunkuttang}  \pronnote{bunkuʈ͡ʂaŋ}
  \pos{Noun}  \sense{
    \definition{catfish}}}

\entry{Bunkuttangmälla}{\headword{Bunkuttangmälla}
  \pos{Proper noun}  \sense{
    \definition{female personal name}}}

\entry{bunmat}{\headword{bunmat}
  \pos{Noun}  \sense{
    \definition{center of a garden}}}

\entry{buntok}{\headword{buntok}
  \pos{Transitive coverb}  \sense{
    \definition{to carry on head}}}

\entry{bur}{\headword{bur}  \note{Birds}  \pronnote{bur}
  \pos{Noun}  \sense{
    \definition{type of bird}    \note{Birds}    \example{
}    \example{
}    \example{
}}}

\entry{burag}{\headword{burag}  \note{Marriage}  \pronnote{buraɡ}
  \pos{Noun}  \sense{
    \definition{bride price (given to the bride's family by the groom)}    \note{Marriage}    \example{
}    \example{
}    \example{
}}}

\entry{burara}{\headword{burara}  \lexeme{burara}  \pronnote{burara}
  \pos{Noun}  \sense{
    \definition{water lily}}}

\entry{but}{\headword{but}
  \variantof{freevarlab of \vartext{English loanword}}}

\entry{buwo}{\headword{buwo}  \note{Names of bananas}  \pronnote{buwo}
  \pos{Noun}  \sense{
    \definition{type of native banana}    \note{Names of bananas}    \example{
}    \example{
}    \example{
}}}

\entry{Buyubun}{\headword{Buyubun}  \note{Places around Limol}  \pronnote{bujubun}
  \pos{Proper noun}  \sense{
    \definition{Buyubun (toponym)}    \note{Places around Limol}    \example{
}    \example{
}    \example{
}}}

\entry{buz}{\headword{buz}  \note{Trees}  \pronnote{buz}
  \pos{Noun}  \sense{
    \definition{type of cultivated tree with fist-sized, edible green fruit and light purple flowers}    \note{Trees}    \example{
}    \example{
}    \example{
}}}

\entry{Buzi}{\headword{Buzi}  \pronnote{buzi}
  \pos{Proper noun}  \sense{
    \definition{Buzi (in Kiwai Rural LLG)}    \example{
}    \example{
}    \example{
}}}

\chapter{C}


\entry{Caso}{\headword{Caso}  \pronnote{kaso}
  \pos{Proper noun}  \sense{
    \definition{male personal name}}}

\entry{Cathy}{\headword{Cathy}  \pronnote{kaθi}
  \pos{Proper noun}  \sense{
    \definition{female personal name}}}

\entry{central}{\headword{central}
  \variantof{Dialectal Variant of \vartext{English loanword}}}

\entry{chairman}{\headword{chairman}
  \variantof{Dialectal Variant of \vartext{English loanword}}}

\entry{Charles}{\headword{Charles}
  \pos{Proper noun}  \sense{
    \definition{male personal name}}}

\entry{Christina}{\headword{Christina}  \pronnote{kristina}
  \pos{Proper noun}  \sense{
    \definition{female personal name}}}

\entry{church}{\headword{church}
  \variantof{freevarlab of \vartext{English loanword}}}

\entry{Cynthia}{\headword{Cynthia}
  \variantof{SpellingVariant of \vartext{Sintia}}}

\chapter{D}


\entry{d-}{\headword{d-}  \pronnote{d}
  \pos{Verb}  \sense{
    \definition{rem.tr}}}

\entry{da}{\headword{da}  \pronnote{da}
  \pos{Subordinating connective}  \sense{\textbf{1.} 
    \definition{if, when (introduces a condition)}    \example{
      \exsen{Bongo ai dan da däbe melem de nanges, sisri nanges.}      \extran{If you can do that job, do it today.}}    \example{
      \exsen{Da Joshua sɨmell de bäbäddän, ibi mɨnyi toto duwem e bäddägeya.}      \extran{If Joshua shoots a pig, we will eat it for dinner.}}    \example{
      \exsen{Da Llimoll me lla da kuddäll bogon, komlla o kumuddäga lla da mɨnyi utt alle eka de bängaeyo.}      \extran{When a person dies in Limol, two or three people will deliver the news with a conch shell.}}}  \sense{\textbf{2.} 
    \definition{might, may, could (marks potential mood)}    \example{
      \exsen{Kollokolloe eka panynen a mudan. Da iba eka de bawengameya.}      \extran{Speaking a mixed language (i.e. code-switching) is no good. We might forget our (incl.) language.}}}
  \variants{\Unspecified Variant \vartext{de}}}

\entry{da}{\headword{da}
  \pos{Nominal demonstrative}  \sense{\textbf{1.} 
    \definition{there}}  \sense{\textbf{2.} 
    \definition{he, she, it (third person singular animate/inanimate pronoun, nominative only)}    \example{
      \exsen{Mozaya yu nägagalla. Da därko dägagan.}      \extran{We grilled the mozaya fish. It got hard.}}    \example{
      \exsen{Da näkäp buddo me adämenan.}      \extran{He's sat with a heavy mind.}}}  \sense{\textbf{3.} 
    \definition{last}    \example{
      \exsen{da pazi me}      \extran{last year}}}}

\entry{da}{\headword{da}
  \pos{Nominal demonstrative}  \sense{
    \definition{that (medial demonstrative and pronoun)}    \example{
      \exsen{da ma me}      \extran{in that house}}    \example{
      \exsen{Da dibeya.}      \extran{That's that.}}}}

\entry{da~}{\headword{da~}
  \variantof{Infinitival reduplicant of \vartext{RED~}}}

\entry{=da}{\headword{=da}  \pronnote{da}
  \pos{Nominal enclitic}  \sense{\textbf{1.} 
    \definition{nominative clitic (marks the subject or agent of the verb)}    \example{
      \exsen{Ddia da wa kottllam a dindugmällneyo.}      \extran{The deer and the turtle were going to have a race.}}}  \sense{\textbf{2.} 
    \definition{accusative clitic on conjoined objects (marks the object of the verb in conjoined noun phrases)}    \example{
      \exsen{Ngämi ngäma toboll a bägäl a danserbeaebeya.}      \extran{We (excl.) prepared our bows and arrows.}}}}

\entry{=da}{\headword{=da}  \pronnote{da}
  \pos{Nominal enclitic}  \sense{\textbf{1.} 
    \definition{close possessive clitic}    \example{
      \exsen{Ine kutt da bo lid a popang dan.}      \extran{The water container's lid has a hole.}}}  \sense{\textbf{2.} 
    \definition{close possessive kinship clitic (marks third person possession on the nominal phrase; can only attach to phrases with kinship nominal heads)}    \example{
      \exsen{Däräng käsre nagda bom yagnen dängkamän.}      \extran{Dog then started to look for his friend.}}}}

\entry{da}{\headword{da}
  \variantof{freevarlab of \vartext{dan}}  \pronnote{da}}

\entry{däba}{\headword{däba}  \pronnote{dəba}
  \pos{Nominal demonstrative}  \sense{
    \definition{that (medial demonstrative)}    \example{
      \exsen{Däba ngättäma we gotarän.}      \extran{He slept in that place.}}}
  \variants{\Unspecified Variant \vartext{däbe}}}

\entry{däba}{\headword{däba}  \pronnote{dəba}
  \pos{Transitive verb}  \sense{
    \definition{to place, put}    \example{
      \exsen{Ngämo toboll de dädäbaneg llo mit mi.}      \extran{I placed my spears at the base of the tree.}}}}

\entry{däba}{\headword{däba}  \note{Trees}  \pronnote{dəba}
  \pos{Noun}  \sense{
    \definition{type of tree that grows in the grassland with leaves used to wrap sago and durable wood used for kundu drums, house posts, and formerly, bridges}    \note{Trees}    \example{
}    \example{
}    \example{
}}}

\entry{däbaballe}{\headword{däbaballe}
  \pos{Adverbial demonstrative}  \sense{
}}

\entry{däbae}{\headword{däbae}  \pronnote{dəbae}
  \pos{Intransitive verb}  \sense{\textbf{1.} 
    \definition{to drizzle}    \example{
      \exsen{Sisri abal yogoll a adäbaeyan.}      \extran{It just started to drizzle.}}}  \sense{\textbf{2.} 
    \definition{to knock fruit continuously}}}

\entry{däbaeya}{\headword{däbaeya}  \pronnote{dəbaeja}
  \pos{Copular verb}  \sense{
    \definition{past singular form of däban}    \example{
      \exsen{Däbaeya ngämo ttoenttoen a.}      \extran{That was my little story.}}    \example{
      \exsen{Däbe.}      \extran{That was it. / The end.}}}
  \variants{\freevarlabs \vartext{däbeya, däbaya, däbe, dɨbeya}}}

\entry{däbag}{\headword{däbag}
  \pos{Copular verb}  \sense{
    \definition{present plural form of däban}    \example{
      \exsen{Maten bo llɨg a ngämo kokok a däbag.}      \extran{Martin's children, they are my grandchildren.}}}}

\entry{däbagaeya}{\headword{däbagaeya}  \pronnote{dəbaɡaeja}
  \pos{Copular verb}  \sense{
    \definition{past plural form of däban}    \example{
      \exsen{Däbagaeya ngämo bin a.}      \extran{Those were my names.}}}
  \variants{\freevarlab \vartext{däbageya}}}

\entry{däbageya}{\headword{däbageya}
  \variantof{freevarlab of \vartext{däbagaeya}}  \pronnote{dəbaɡeja}}

\entry{däbageyo}{\headword{däbageyo}
  \pos{Copular verb}  \sense{
}}

\entry{däbagwaeya}{\headword{däbagwaeya}  \pronnote{dəbaɡwaeja}
  \pos{Copular verb}  \sense{
    \definition{past dual form of däban}    \example{
      \exsen{Kurupel a Täm, ngämo baba bi dibagweya.}      \extran{Kurupel and Tam, they were my parents.}}}}

\entry{däbako}{\headword{däbako}
  \pos{Nominal demonstrative}  \sense{
}}

\entry{däbäll}{\headword{däbäll}  \pronnote{dəbəɽ}
  \pos{Transitive verb}  \sense{
    \definition{to touch}    \example{
      \exsen{Nädbällnegan.}      \extran{He touched him.}}    \example{
      \exsen{Nädbällnegan}      \extran{He touched three people at different times.}}    \example{
      \exsen{Lla da ubim yadbällan.}      \extran{He touched three people at the same time.}}}}

\entry{däballe}{\headword{däballe}
  \pos{Demonstrative}  \sense{
}}

\entry{daballe}{\headword{daballe}
  \pos{Adverbial demonstrative}  \sense{
    \definition{ablative form of da; after that}}}

\entry{däbamasäm}{\headword{däbamasäm}  \pronnote{dəbamasem}
  \pos{Adverbial demonstrative}  \sense{
    \definition{ablative form of däba}    \example{
      \exsen{Däbamasem a, ako ttongo källäm me gobäll.}      \extran{From there, we went to another pond.}}    \example{
      \exsen{Däbamatt a ngämi gall deyagllaenalla.}      \extran{After that, we paddled the canoe over.}}}
  \variants{\Unspecified Variant \vartext{däbamattäm, däbamasem}}}

\entry{däbamasem}{\headword{däbamasem}
  \variantof{Unspecified Variant of \vartext{däbamasäm}}}

\entry{däbamattäm}{\headword{däbamattäm}
  \variantof{Unspecified Variant of \vartext{däbamasäm}}}

\entry{däbame}{\headword{däbame}
  \pos{}  \sense{
}}

\entry{däban}{\headword{däban}
  \pos{Copular verb}  \sense{
    \definition{copular form of däba (present singular form)}    \example{
      \exsen{Sisri ngäna pasta dan. Ngämo melem a sisri däban.}      \extran{Now I am a pastor. That is my job.}}}
  \variants{\Dialectal Variant \vartext{däben}}}

\entry{dabän}{\headword{dabän}  \pronnote{dabən}
  \pos{Transitive verb}  \sense{
    \definition{to knock a leaf, e.g. a tobacco leaf}}}

\entry{däbaolle}{\headword{däbaolle}  \pronnote{dəbaoɽe}
  \pos{Adverb}  \sense{
    \definition{allative form of däba}    \example{
      \exsen{Däbaolle ikop e abällan.}      \extran{We went there to see.}}}}

\entry{däbaya}{\headword{däbaya}
  \variantof{freevarlab of \vartext{däbaeya}}  \pronnote{dəbaja}}

\entry{däbe}{\headword{däbe}
  \variantof{Unspecified Variant of \vartext{däba}}  \pronnote{dəbe}}

\entry{däbe}{\headword{däbe}
  \variantof{freevarlab of \vartext{däbaeya}}  \pronnote{dəbe}}

\entry{Dabe}{\headword{Dabe}  \note{Places around Limol}  \pronnote{dabe}
  \pos{Proper noun}  \sense{
    \definition{Dabe (toponym)}    \note{Places around Limol}    \example{
}    \example{
}    \example{
}}}

\entry{däbe}{\headword{däbe}  \pronnote{dəbe}
  \pos{Demonstrative}  \sense{
    \definition{demonstrative, 'that,' could also translate as 'these' or 'those', depending on the context}    \example{
      \exsen{Bogo dindugän do kukiny mama de ikop dägagän dibaeya däbe ik e gotärakän do gotägolän.}      \extran{He ran to a pile of tall grass that he saw, and went inside and hid there.}}    \example{
}}
  \variants{\Unspecified Variant \vartext{dɨbe}}}

\entry{däbem}{\headword{däbem}
  \pos{Nominal demonstrative}  \sense{
    \definition{accusative form of däba}    \example{
      \exsen{Tämamae llɨg klekle da däbem llo de kälnan erallo.}      \extran{All the children are climbing that tree.}}}}

\entry{däbemasem}{\headword{däbemasem}  \pronnote{dəbemasem}
  \pos{Adverb}  \sense{
    \definition{ablative form of däbe}}}

\entry{däben}{\headword{däben}
  \variantof{Dialectal Variant of \vartext{däban}}}

\entry{däbeya}{\headword{däbeya}
  \variantof{freevarlab of \vartext{däbaeya}}  \pronnote{dəbeja}}

\entry{däbi}{\headword{däbi}  \pronnote{dəbi}
  \pos{Noun}  \sense{
    \definition{green-backed honeyeater}    \scientificname{Glycichaera fallax}}}

\entry{Dabi}{\headword{Dabi}  \pronnote{dabi}
  \pos{Proper noun}  \sense{
    \definition{female personal name}}}

\entry{dabit}{\headword{dabit}
  \pos{Noun}  \sense{
    \definition{palm spear}}}

\entry{däd}{\headword{däd}  \pronnote{dəd}
  \pos{Demonstrative}  \sense{
    \definition{there}}}

\entry{dada}{\headword{dada}  \pronnote{dada}
  \pos{Kinship noun}  \sense{
    \definition{older sibling of the same sex (man's older brother or woman's older sister)}    \example{
}    \example{
}    \example{
}}}

\entry{dadaeya}{\headword{dadaeya}  \pronnote{dadaeja}
  \pos{Copular verb}  \sense{
    \definition{past singular form of daden}    \example{
      \exsen{Bogo tätäm dadaeya ma me?}      \extran{Was he home yesterday?}}}
  \variants{\freevarlabs \vartext{dadiweya, dadeya}}}

\entry{dädäk}{\headword{dädäk}  \pronnote{dədək}
  \pos{Noun}  \sense{
    \definition{wall}}}

\entry{dadäp}{\headword{dadäp}  \pronnote{dadəp}
  \pos{Noun}  \sense{
    \definition{plant part}    \example{
      \exsen{Masar bo nge dadäp peyang gognegän.}}}}

\entry{dädär}{\headword{dädär}  \pronnote{dədər}
  \pos{Noun}  \sense{
    \definition{type of big taro}    \example{
}    \example{
}    \example{
}}
  \variants{\Unspecified Variant \vartext{dɨdɨr}}}

\entry{dädär}{\headword{dädär}  \pronnote{dədər}
  \pos{Noun}  \sense{\textbf{1.} 
    \definition{stone, rock}    \example{
      \exsen{Ngängälatt dädär a mänyi tutu atta bolldaeyän kopek e.}      \extran{A round rock will roll from the hill to the valley.}}}  \sense{\textbf{2.} 
    \definition{dry}    \example{
      \exsen{Ttängäm ulle da tämamae dallän dädär abal gogon.}      \extran{All the big villages became very dry.}}}  \sense{\textbf{3.} 
    \definition{hard}    \example{
      \exsen{Bogo dɨdɨr bun peyang dan.}      \extran{He has a hard head (i.e. he is stubborn).}}}
  \variants{\Unspecified Variant \vartext{dɨdɨr}}}

\entry{dadär}{\headword{dadär}  \note{Traditional fishing ways and nets}  \pronnote{dadər}
  \pos{Noun}  \sense{
    \definition{type of net suspended in the water by sticks}    \note{Traditional fishing ways and nets}    \example{
}    \example{
}    \example{
}}}

\entry{dädär käp}{\headword{dädär käp}
  \pos{Noun}  \sense{
    \definition{stone}    \example{
      \exsen{dädär käp turik}      \extran{stone ax}}}}

\entry{dadäräb}{\headword{dadäräb}  \pronnote{dadərəb}
  \pos{Transitive verb}  \sense{\textbf{1.} 
    \definition{to write}    \example{
      \exsen{Ngämo kame dan dadräb a.}      \extran{I can't write.}}    \example{
      \exsen{Pizin eka de ai dan badräb.}      \extran{I can write in Tok Pisin.}}}  \sense{\textbf{2.} 
    \definition{to dress}    \example{
      \exsen{Bogo obom dadäräb eran.}      \extran{She is dressing her.}}}  \sense{\textbf{3.} 
    \definition{to decorate}    \example{
      \exsen{Bogo mɨnyi llo ttam alle godräballe ingong de ngkamalle bandra peyang.}      \extran{She would decorate herself with leaves and start to dance to the song.}}}
  \variants{\Fast Speech Variant \vartext{dadräb}}}

\entry{dädäräb}{\headword{dädäräb}  \pronnote{dədərəb}
  \pos{Transitive verb}  \sense{
    \definition{to cut grass, flowers}}}

\entry{dadäräd}{\headword{dadäräd}  \pronnote{dadərəd}
  \pos{Intransitive verb}  \sense{
    \definition{to clear (e.g. of the mind, sky)}    \anthronote{Traditional way of clearing your mind is by cutting with shell.}    \example{
      \exsen{Näkäp a odardan.}      \extran{The mind is cleared.}}}}

\entry{dadargu}{\headword{dadargu}  \note{Things a person can do}  \pronnote{dadarɡu}
  \pos{Noun}  \sense{
    \definition{disturbance}    \note{Things a person can do}    \example{
}    \example{
}    \example{
}}}

\entry{dade}{\headword{dade}  \lexeme{dade}  \pronnote{dade}
  \pos{Manner adverb}  \sense{\textbf{1.} 
    \definition{maybe}    \example{
      \exsen{Bogo täbädd lla bälle panya wo, wup wo, o dade wayati käp de däddänalle.}      \extran{He used to pick a ripe pineapple, ripe banana, or maybe a watermelon for the strange man.}}}  \sense{\textbf{2.} 
    \definition{ever}    \example{
      \exsen{dade aya}      \extran{whoever}}    \example{
      \exsen{dade erame}      \extran{wherever}}    \example{
      \exsen{dade eremde}      \extran{whichever}}    \example{
      \exsen{dade angde}      \extran{whenever}}}}

\entry{dade}{\headword{dade}  \lexeme{dade}  \pronnote{dade}
  \pos{Noun}  \sense{
    \definition{yam stick}    \example{
      \exsen{Zugu obo mätta dade de era topotopoll dade dae dägawän.}      \extran{Zugu only planted topotopoll yam sticks for his yams.}}}}

\entry{dade}{\headword{dade}
  \variantof{freevarlab of \vartext{daden}}}

\entry{dadeg}{\headword{dadeg}  \pronnote{dadeɡ}
  \pos{Copular verb}  \sense{
    \definition{present plural form of daden}    \example{
      \exsen{Kemu bo llabun a gänyme a Malläm me dadeg.}      \extran{Kemu has relatives here and in Malam (lit. Kemu's relatives exist here and in Malam).}}}
  \variants{\freevarlab \vartext{dadegän}}}

\entry{dadegaeya}{\headword{dadegaeya}  \pronnote{dadeɡaeja}
  \pos{Copular verb}  \sense{
    \definition{past plural form of daden}    \example{
      \exsen{Tätäm ubi kumuddäga dadegaeya ma me?}      \extran{Were they three home yesterday?}}}
  \variants{\freevarlab \vartext{dadegeya}}}

\entry{dadegaeyo}{\headword{dadegaeyo}  \pronnote{dadeɡaejo}
  \pos{Copular verb}  \sense{
    \definition{present dual form of daden}    \example{
      \exsen{Buk a komllaebe dadegeyo.}      \extran{There are only two books.}}}
  \variants{\freevarlab \vartext{dadegeyo}}}

\entry{dadegän}{\headword{dadegän}
  \variantof{freevarlab of \vartext{dadeg}}
  \variants{\Dialectal Variant \vartext{gan}}}

\entry{dadegeya}{\headword{dadegeya}
  \variantof{freevarlab of \vartext{dadegaeya}}}

\entry{dadegeyo}{\headword{dadegeyo}
  \variantof{freevarlab of \vartext{dadegaeyo}}  \pronnote{dadeɡejo}}

\entry{dadegwaeya}{\headword{dadegwaeya}  \pronnote{dadeɡwaeja}
  \pos{Copular verb}  \sense{
    \definition{past dual form of daden}    \example{
      \exsen{Tätäm ubi komlla dadegwaeya ma me}      \extran{Were they two home yesterday?}}}
  \variants{\freevarlab \vartext{dadegweya}}}

\entry{dadegweya}{\headword{dadegweya}
  \variantof{freevarlab of \vartext{dadegwaeya}}}

\entry{dadel}{\headword{dadel}  \note{Gardens}  \pronnote{dadel}
  \pos{Noun}  \sense{\textbf{1.} 
    \definition{harvest}    \note{Gardens}    \example{
      \exsen{dadel ebdo}      \extran{harvest day}}    \example{
}    \example{
}}  \sense{\textbf{2.} 
    \definition{season of harvesting young gardens (seventh season; corresponds to early May)}    \note{Gardens}    \example{
}}}

\entry{daden}{\headword{daden}  \note{Verbs}  \pronnote{daden}
  \pos{Copular verb}  \sense{
    \definition{to exist, have, there is (present singular form)}    \note{Verbs}    \anthronote{The hand sign for the message daden or "s/he is here" is to pat the air with the hand.}    \example{
      \exsen{Bäne ine da daden?}      \extran{Do you have water (lit. does your water exist)?}}}
  \variants{\freevarlabs \vartext{dade, dadi}}
  \variants{\Dialectal Variant \vartext{dadewän, dadewon}}}

\entry{dadewän}{\headword{dadewän}
  \variantof{Dialectal Variant of \vartext{daden}}
  \variants{\SpellingVariant \vartext{dadewon}}}

\entry{dadewon}{\headword{dadewon}
  \variantof{Dialectal Variant of \vartext{daden}}}

\entry{dadeya}{\headword{dadeya}
  \variantof{freevarlab of \vartext{dadaeya}}  \pronnote{dadeja}}

\entry{dadi}{\headword{dadi}
  \variantof{freevarlab of \vartext{daden}}}

\entry{Dadi}{\headword{Dadi}
  \pos{Proper noun}  \sense{
    \definition{male personal name}}}

\entry{dadi}{\headword{dadi}
  \variantof{SpellingVariant of \vartext{dedi}}}

\entry{dädik}{\headword{dädik}  \pronnote{dədik}
  \pos{Transitive verb}  \sense{\textbf{1.} 
    \definition{to calm someone down with cool words from anger}}  \sense{\textbf{2.} 
    \definition{to bump}}}

\entry{dadiweya}{\headword{dadiweya}
  \variantof{freevarlab of \vartext{dadaeya}}  \pronnote{dadiweja}}

\entry{dädme}{\headword{dädme}  \pronnote{dədme}
  \pos{Adverbial demonstrative}  \sense{
    \definition{there}}}

\entry{dadräb}{\headword{dadräb}
  \variantof{Fast Speech Variant of \vartext{dadäräb}}}

\entry{dädri}{\headword{dädri}  \pronnote{dədri}
  \pos{Adverbial demonstrative}  \sense{
    \definition{there, that way, that side}    \example{
      \exsen{Ttongo män a ngäma dädri dan Wasua me.}      \extran{One of our (excl.) daughters is there, in Wasua.}}}}

\entry{dädribattäm}{\headword{dädribattäm}
  \pos{}}

\entry{=dae}{\headword{=dae}
  \variantof{Unspecified Variant of \vartext{=daebe}}}

\entry{=dae}{\headword{=dae}  \pronnote{dae}
  \pos{Nominal enclitic}  \sense{
    \definition{perlative case clitic; along, through}    \example{
      \exsen{Ngäna wälläng ik dae ibi allan.}      \extran{I am going through the bush.}}    \example{
      \exsen{Ngäna walle dae ibi allan.}      \extran{I am going along the river.}}}}

\entry{=daebe}{\headword{=daebe}
  \pos{Nominal enclitic}  \sense{
    \definition{restrictive clitic; only}}
  \variants{\Unspecified Variant \vartext{=dae}}}

\entry{daebeg}{\headword{daebeg}  \pronnote{daebeɡ}
  \pos{Copular verb}  \sense{
    \definition{present plural form of daeben}    \example{
      \exsen{Dedme ddone dan be kumuddäga nyäng daebeg.}      \extran{There's nothing there but three bags.}}}}

\entry{daebegaeya}{\headword{daebegaeya}  \pronnote{daebeɡaeja}
  \pos{Copular verb}  \sense{
    \definition{past plural form of daeben}    \example{
      \exsen{Misdae up daebegaeya ibebatta.}      \extran{The only things planted were bananas.}}}}

\entry{daebegeyo}{\headword{daebegeyo}  \pronnote{daebeɡejo}
  \pos{Copular verb}  \sense{
    \definition{present dual form of daeben}    \example{
      \exsen{Ende ttängäm a komlla daebegeyo, Malläm a Llimoll.}      \extran{There are only two Ende villages: Malam and Limol.}}}}

\entry{daeben}{\headword{daeben}  \pronnote{daeben}
  \pos{Copular verb}  \sense{
    \definition{copular form of daebe (present singular form)}    \example{
      \exsen{Däbe mälla bo melem daeben.}      \extran{That is just women's work.}}}}

\entry{daebeya}{\headword{daebeya}  \pronnote{daebeja}
  \pos{Copular verb}  \sense{
    \definition{past singular form of daeben}    \example{
      \exsen{Diba kui mi ttongdae ddia gullbe daebeya.}      \extran{On that island, there was only one stag.}}}}

\entry{daem}{\headword{daem}
  \pos{Intransitive coverb}  \sense{
    \definition{to blink}    \example{
      \exsen{Ngäna ikop daem agan.}      \extran{I blinked.}}}}

\entry{däen}{\headword{däen}  \note{seems like it's missing a medial consonant}  \pronnote{dəen}
  \pos{Noun}  \sense{
    \definition{type of snake}    \note{seems like it's missing a medial consonant}    \example{
}    \example{
}    \example{
}}}

\entry{Daena}{\headword{Daena}
  \pos{Proper noun}  \sense{
    \definition{female personal name}}}

\entry{daendae}{\headword{daendae}  \lexeme{daendae}  \pronnote{daendae}
  \pos{Noun}  \sense{
    \definition{flowering plant that is said to be ancestral to the area, with many types; flowers used as adornment when dancing}}
  \variants{\SpellingVariant \vartext{daindai}}}

\entry{daes}{\headword{daes}
  \variantof{freevarlab of \vartext{English loanword}}}

\entry{daeya}{\headword{daeya}  \pronnote{daeja}
  \pos{Copular verb}  \sense{
    \definition{past singular form of dan}    \example{
      \exsen{Ngämo baba ttongdae mälla peyang daeya.}      \extran{My father had (lit. was with) one wife.}}}
  \variants{\Fast Speech Variant \vartext{de}}
  \variants{\freevarlabs \vartext{deya, daeyag}}}

\entry{daeyag}{\headword{daeyag}
  \variantof{freevarlab of \vartext{daeya}}}

\entry{Daeyagmälla}{\headword{Daeyagmälla}
  \pos{Proper noun}  \sense{
    \definition{female personal name}}}

\entry{Daeyna}{\headword{Daeyna}
  \pos{Proper noun}  \sense{
    \definition{female personal name}}}

\entry{dag}{\headword{dag}  \pronnote{daɡ}
  \pos{Copular verb}  \sense{
    \definition{present plural form of dan}    \example{
      \exsen{Ubi dag.}      \extran{They are there.}}    \example{
      \exsen{Ikop a ulleulle dag.}      \extran{The eyes are big.}}}}

\entry{däg}{\headword{däg}  \lexeme{däg}  \pronnote{dəɡ}
  \pos{Noun}  \sense{
    \definition{hand, group, bunch, set}    \example{
      \exsen{buata däg}      \extran{betelnut bunch}}    \example{
      \exsen{manggo däg}      \extran{bunch of mango}}    \example{
      \exsen{nge däg}      \extran{coconut bunch}}    \example{
      \exsen{tämamae däg}      \extran{an entire set (e.g. of teeth)}}}}

\entry{daga}{\headword{daga}  \note{Trees}  \pronnote{daɡa}
  \pos{Noun}  \sense{
    \definition{betel}    \scientificname{Piper betle}    \note{Trees}    \example{
}    \example{
}    \example{
}}}

\entry{dägadäga}{\headword{dägadäga}  \pronnote{dəɡadəɡa}
  \pos{Adverb}  \sense{
    \definition{completely}    \example{
      \exsen{Ngäna ngämo sana de ot dägadäga dättemän.}      \extran{I finished eating all my sago.}}}}

\entry{dägädägäl}{\headword{dägädägäl}
  \variantof{SpellingVariant of \vartext{dägäldägäl}}}

\entry{dagaeya}{\headword{dagaeya}  \pronnote{daɡaeja}
  \pos{Copular verb}  \sense{
    \definition{past plural form of dan}    \example{
      \exsen{Obo mälläng ik a ulleulle abal dagaeya.}      \extran{His nostrils were very big.}}}
  \variants{\freevarlab \vartext{dageya}}
  \variants{\Dialectal Variant \vartext{dagaya}}}

\entry{dagal}{\headword{dagal}  \pronnote{daɡal}
  \pos{Transitive verb}  \sense{\textbf{1.} 
    \definition{to board, get on}    \example{
      \exsen{Ngämi deyareya gall e godagaleya.}      \extran{We went to the boat and got on.}}    \example{
      \exsen{Ngäna gongäs gall e godagal.}      \extran{I returned to the boat and got on.}}}  \sense{\textbf{2.} 
    \definition{to board, load, put on}    \example{
      \exsen{Ngämo män kälsre de gall e dädagal.}      \extran{I put my daughter onto the boat.}}}}

\entry{dägäldägäl}{\headword{dägäldägäl}  \note{Parts of a bird}  \pronnote{dəɡəldəɡəl}
  \pos{Noun}  \sense{
    \definition{intestines}    \note{Parts of a bird}    \example{
      \exsen{Dägäldägäl a ddäddäg ma dan.}      \extran{The intenstines are edible.}}    \example{
}    \example{
}}
  \variants{\SpellingVariant \vartext{dägädägäl}}}

\entry{dagaya}{\headword{dagaya}
  \variantof{Dialectal Variant of \vartext{dagaeya}}  \pronnote{daɡaja}}

\entry{dageya}{\headword{dageya}
  \variantof{freevarlab of \vartext{dagaeya}}  \pronnote{daɡeja}}

\entry{dageyo}{\headword{dageyo}  \pronnote{daɡejo}
  \pos{Copular verb}  \sense{
    \definition{present dual form of dan}    \example{
      \exsen{Ubi komlla era nag dageyo.}      \extran{The two of them are friends.}}}
  \variants{\Dialectal Variant \vartext{dagwän}}}

\entry{dägmar}{\headword{dägmar}  \pronnote{dəɡmar}
  \pos{Noun}  \sense{\textbf{1.} 
    \definition{tongue}    \example{
      \exsen{Dägmar a malla ddäddäg ma dan.}      \extran{The tongue is not edible.}}}  \sense{\textbf{2.} 
    \definition{spoon}    \example{
      \exsen{Dägmar alle wätät de notnan kakoll att.}      \extran{He was eating food from a plate with a spoon.}}}}

\entry{dagwaeya}{\headword{dagwaeya}  \pronnote{daɡwaeja}
  \pos{Copular verb}  \sense{
    \definition{past dual form of dan}    \example{
      \exsen{Däräng a wa Baet a mer abal nag dagwaeya.}      \extran{Dog and Cuscus were very good friends.}}}
  \variants{\freevarlabs \vartext{dagwe, dagweya}}}

\entry{dagwän}{\headword{dagwän}
  \variantof{Dialectal Variant of \vartext{dageyo}}}

\entry{dagwe}{\headword{dagwe}
  \variantof{freevarlab of \vartext{dagwaeya}}  \pronnote{daɡwe}}

\entry{dagweya}{\headword{dagweya}
  \variantof{freevarlab of \vartext{dagwaeya}}  \pronnote{daɡweja}}

\entry{Daiba}{\headword{Daiba}  \note{Places around Limol}  \pronnote{daiba}
  \pos{Proper noun}  \sense{
    \definition{garden and sago place of Jerry Dareda (along the road to Bisuaka)}    \note{Places around Limol}    \example{
}    \example{
}    \example{
}}}

\entry{daindai}{\headword{daindai}
  \variantof{SpellingVariant of \vartext{daendae}}  \pronnote{daindai}}

\entry{daindaim}{\headword{daindaim}
  \pos{Noun}  \sense{
    \definition{drizzle}}}

\entry{däkäna}{\headword{däkäna}
  \variantof{SpellingVariant of \vartext{däkna}}  \pronnote{dəkəna}}

\entry{däkna}{\headword{däkna}  \pronnote{dəkna}
  \pos{Noun}  \sense{
    \definition{small black termite mound that burns for a long time; after a woman gives birth, it is heated in the fire, wrapped in bark and cloth, placed under a mat, and used to warm the woman's stomach}    \example{
}    \example{
}    \example{
}}
  \variants{\SpellingVariant \vartext{däkäna}}}

\entry{däl}{\headword{däl}
  \variantof{SpellingVariant of \vartext{dɨl}}}

\entry{dalab}{\headword{dalab}  \pronnote{'dalab}
  \pos{Transitive verb}  \sense{
    \definition{to open, pierce, make a hole}    \example{
      \exsen{Obo käp de pop nedelaban.}      \extran{He made a hole in the fruit.}}}
  \variants{\Dialectal Variant \vartext{darab}}}

\entry{dale}{\headword{dale}  \note{Parts of a fire}  \pronnote{dale}
  \pos{Noun}  \sense{
    \definition{ash}    \note{Parts of a fire}    \example{
}    \example{
}    \example{
}}}

\entry{däm}{\headword{däm}  \pronnote{dəm}
  \pos{Noun}  \sense{
    \definition{plant}    \example{
      \exsen{Bogo gärep kabkab däm de dibewän diba ttängäm me.}      \extran{He planted grape plants in that place.}}}}

\entry{dam}{\headword{dam}  \note{Question words}  \pronnote{dam}
  \pos{Adverb}  \sense{
    \definition{then}    \note{Question words}    \example{
      \exsen{Angde ngäna otaran, dam yogoll agan.}      \extran{I was sleeping, and then it rained.}}}}

\entry{däm ibenen}{\headword{däm ibenen}  \note{Seasons}  \pronnote{dəm ibenen}
  \pos{Noun}  \sense{
    \definition{season when crops are planted (fifteenth season; corresponds to late November)}    \note{Seasons}    \example{
}    \example{
}    \example{
}}}

\entry{dämädämäll}{\headword{dämädämäll}  \pronnote{dəmədəməɽ}
  \pos{Modifier}  \sense{
    \definition{numb, paralyzed}    \example{
      \exsen{apte pätt dämädämällang lla}      \extran{man with one half of his body paralyzed}}}}

\entry{dämar}{\headword{dämar}  \lexeme{dämar}  \pronnote{dəmar}
  \pos{Noun}  \sense{
    \definition{type of palm tree (~2 m) that used to be cooked and eaten; also fed to pigs}}}

\entry{damärärmae}{\headword{damärärmae}  \pronnote{damərərmae}
  \pos{Adverb}  \sense{
    \definition{simultaneously}    \example{
      \exsen{Ibi mɨnyi damärärmae toboll daspunigalla.}      \extran{We will shoot arrows at the same time.}}}}

\entry{damasäm}{\headword{damasäm}
  \pos{Adverbial demonstrative}  \sense{
    \definition{ablative form of da; therefore}    \example{
      \exsen{Ddob Idi lla da dadeg, damasäm Ende eka de klloklloe panypeny erallo.}      \extran{There are also Idi speakers, so people are speaking mixed Ende.}}}
  \variants{\Unspecified Variant \vartext{damasema}}}

\entry{damasema}{\headword{damasema}
  \variantof{Unspecified Variant of \vartext{damasäm}}  \pronnote{damasema}}

\entry{dämbaemeny}{\headword{dämbaemeny}  \pronnote{dəmbaemeɲ}
  \pos{Intransitive verb}  \sense{
    \definition{to doze off}    \example{
      \exsen{Dagirnän e: yunu gondäbaemenynän.}      \extran{and stayed there for a long, long time until he dozed off.}}}}

\entry{dämbag}{\headword{dämbag}  \pronnote{dəmbaɡ}
  \pos{Noun}  \sense{
    \definition{lazy person, weak person}}}

\entry{dämen}{\headword{dämen}  \note{Verbs of Motion}  \pronnote{dəmen}
  \pos{Intransitive verb}  \sense{
    \definition{to sit}    \note{Verbs of Motion}    \example{
      \exsen{Ngäna dämenang dan.}      \extran{I am sitting.}}    \example{
      \exsen{Adämeneyo!}      \extran{(You two) sit down!}}    \example{
      \exsen{Ttongo lla da dämenma toko me adämenan.}      \extran{A person sat on top of the chair.}}}}

\entry{dämenangg}{\headword{dämenangg}  \pronnote{dəmenaŋɡ}
  \pos{Transitive verb}  \sense{\textbf{1.} 
    \definition{causative-applicative form of dämen}    \example{
      \exsen{Nändämenanggan.}      \extran{He sat him down.}}}  \sense{\textbf{2.} 
    \definition{applicative form of dämen}    \example{
      \exsen{Ngäna Kate bom täräb ma me dändämenang.}      \extran{I sat down with Kate at the funeral.}}}}

\entry{dämoe}{\headword{dämoe}  \pronnote{dəmoe}
  \pos{Ditransitive verb}  \sense{\textbf{1.} 
    \definition{to push}    \example{
      \exsen{Ngäna abo däbe sɨmell de misdae ada dandämoe.}      \extran{Then I just pushed the pig like that.}}    \example{
      \exsen{Angde obom ttu wi dandämoeyän, bogo ddone dängllawän, be bogo gonddämän.}      \extran{When she pushed him into the deep, he didn't swim: rather, he drowned.}}}  \sense{\textbf{2.} 
    \definition{to send}    \example{
      \exsen{Mälla de dandmoeaebeya wayati kongkom e.}      \extran{We sent the women to fetch the watermelons.}}    \example{
      \exsen{Giniya män kälsre de nandämoeyan Bodog pate.}      \extran{Giniya sent the girl to Bodog.}}}
  \variants{\Dialectal Variant \vartext{domoe}}}

\entry{dämoengg}{\headword{dämoengg}
  \pos{Transitive verb}  \sense{
    \definition{causative-applicative form of dämoe}    \example{
      \exsen{Ngäna bam sɨmell ikop e dandämoengg.}      \extran{I sent you to see the pig.}}}}

\entry{damona}{\headword{damona}  \note{Numbers}  \pronnote{damona}
  \pos{Numeral}  \sense{
    \definition{1296 (yam counting numeral; 6^4)}    \note{Numbers}    \example{
}    \example{
}    \example{
}}}

\entry{damong}{\headword{damong}
  \pos{Locational}  \sense{
    \definition{against}}}

\entry{damong}{\headword{damong}  \pronnote{'damong}
  \pos{Modifier}  \sense{
    \definition{healthy, well}    \example{
      \exsen{damong mätta}      \extran{a plump yam}}    \example{
      \exsen{Bongo ngänäm damong nagalle, a sisri ngäna damong dan.}      \extran{You comforted me and now I'm okay.}}}}

\entry{dan}{\headword{dan}  \pronnote{dan}
  \pos{Copular verb}  \sense{
    \definition{to be, exist, is, there is (present singular form)}    \example{
      \exsen{Ge nyäng a kälsre dan.}      \extran{This bag is small.}}    \example{
      \exsen{Warama era mer lla dan.}      \extran{Warama is a good man.}}}
  \variants{\Dialectal Variant \vartext{danän}}
  \variants{\freevarlab \vartext{da}}}

\entry{dana}{\headword{dana}  \pronnote{dana}
  \pos{Transitive coverb}  \sense{
    \definition{to sprout}}}

\entry{dänäk}{\headword{dänäk}  \pronnote{dənək}
  \pos{Noun}  \sense{
    \definition{type of small bush with reddish fruit that is black when ripe}}}

\entry{danän}{\headword{danän}
  \variantof{Dialectal Variant of \vartext{dan}}}

\entry{dändak}{\headword{dändak}  \note{Food}  \pronnote{dəndak}
  \pos{Noun}  \sense{
    \definition{type of purple yam with purple skin and no thorns or hairs}    \note{Food}    \example{
}    \example{
}    \example{
}}}

\entry{dändär}{\headword{dändär}  \pronnote{dəndər}
  \pos{Transitive verb}  \sense{\textbf{1.} 
    \definition{to hear, listen}    \example{
      \exsen{Ngäna obom dandär.}      \extran{I heard him.}}    \example{
      \exsen{Ngämne eka de andär!}      \extran{[You] listen to my words!}}}  \sense{\textbf{2.} 
    \definition{to sense, feel}    \example{
      \exsen{Ngäna otät molle de dändär eran.}      \extran{I smell food (lit. sense the smell of food).}}    \example{
      \exsen{Ngämo matta me tätäp dändär eran.}      \extran{I feel pain in my shoulder.}}}  \sense{\textbf{3.} 
}}

\entry{dändär}{\headword{dändär}  \pronnote{dəndər}
  \pos{Transitive verb}  \sense{
    \definition{to stuff}    \example{
      \exsen{Mänmän a sana de zazaba we dändäraebeyo.}      \extran{The girls stuffed the sago into the bag.}}}}

\entry{dändäräm}{\headword{dändäräm}  \note{Trees}  \pronnote{dəndərəm}
  \pos{Noun}  \sense{
    \definition{type of small tree with white and purple flowers and many hard, marble-sized, green fruit that children play with}    \note{Trees}    \example{
}    \example{
}    \example{
}}}

\entry{dändäräm käp}{\headword{dändäräm käp}  \note{Games}  \pronnote{dəndərəm kəp}
  \pos{Noun}  \sense{
    \definition{type of marble game}    \note{Games}    \example{
}    \example{
}    \example{
}}}

\entry{dändäräp}{\headword{dändäräp}
  \variantof{Dialectal Variant of \vartext{dänräp}}  \pronnote{dəndərəp}}

\entry{dändäratt}{\headword{dändäratt}  \pronnote{dəndəraʈ͡ʂ}
  \pos{Noun}  \sense{
    \definition{bagful of sago}}}

\entry{dändärek}{\headword{dändärek}
  \pos{Transitive verb}  \sense{
    \definition{to control, influence, rule, govern}    \example{
      \exsen{Bogo mɨnyi llɨg de gagäll ttoen ngasnges e bändärekän.}      \extran{He will pressure the boy to do bad things.}}}
  \variants{\SpellingVariant \vartext{dändrek}}}

\entry{dändrek}{\headword{dändrek}
  \variantof{SpellingVariant of \vartext{dändärek}}}

\entry{dang}{\headword{dang}  \pronnote{daŋ}
  \pos{Noun}  \sense{
    \definition{length (of a house)}    \example{
      \exsen{dang papek}      \extran{wall spanning the length of a building}}}}

\entry{dang toto}{\headword{dang toto}  \pronnote{daŋ toto}
  \pos{Noun}  \sense{
    \definition{corner post}    \example{
}}}

\entry{dängam}{\headword{dängam}  \note{Birds}  \pronnote{dəŋam}
  \pos{Noun}  \sense{
    \definition{Blyth's hornbil}    \scientificname{Rhyticeros plicatus}    \note{Birds}}}

\entry{dangg}{\headword{dangg}  \pronnote{daŋɡ}
  \pos{Intransitive verb}  \sense{
    \definition{to burn}    \example{
      \exsen{Nyongo da andegan.}      \extran{The road is burning.}}}}

\entry{dangkälmang}{\headword{dangkälmang}  \note{Food}  \pronnote{daŋkəlmaŋ}
  \pos{Noun}  \sense{
    \definition{type of medium-sized, long yam with a white interior and thorns}    \note{Food}    \example{
}    \example{
}    \example{
}}}

\entry{dangkam}{\headword{dangkam}  \pronnote{daŋkam}
  \pos{Intransitive verb}  \sense{\textbf{1.} 
    \definition{to lean}    \example{
      \exsen{Lla da ttongo lla da bo patme dangkam allan.}      \extran{The man is leaning on the other man.}}    \example{
}    \example{
}}  \sense{\textbf{2.} 
    \definition{to rely on}    \example{
      \exsen{Ngäna ngämo mosen patme dangkaemmäll allan ngämingg i.}      \extran{I rely on my brother for help.}}}}

\entry{dangne}{\headword{dangne}  \lexeme{dangne}  \pronnote{daŋne}
  \pos{Noun}  \sense{
    \definition{crawling vine that grows in the grassland with purple flowers; used to make rope}    \anthronote{Old story: when Satan was thrown from heaven, he was tied with this rope.}}}

\entry{dangnebudar}{\headword{dangnebudar}  \note{Grubs}  \pronnote{daŋnebudar}
  \pos{Noun}  \sense{
    \definition{type of grub}    \note{Grubs}    \example{
}    \example{
}    \example{
}}}

\entry{Daniel}{\headword{Daniel}  \pronnote{daniel}
  \pos{Proper noun}  \sense{
    \definition{male personal name}}}

\entry{Danipa}{\headword{Danipa}  \pronnote{danipa}
  \pos{Proper noun}  \sense{
    \definition{male personal name}}}

\entry{dänräp}{\headword{dänräp}  \note{Parts of a fish}  \pronnote{dənrəp}
  \pos{Noun}  \sense{\textbf{1.} 
    \definition{fish scale}    \note{Parts of a fish}    \example{
}}  \sense{\textbf{2.} 
    \definition{scab}    \note{Parts of a fish}}
  \variants{\Dialectal Variant \vartext{dändäräp}}}

\entry{dänyäk}{\headword{dänyäk}  \note{Trees}  \pronnote{dəɲək}
  \pos{Noun}  \sense{
    \definition{small plant that grows in the grassland with purple and white flowers and blue fruit that children like to eat}    \note{Trees}    \example{
}    \example{
}    \example{
}}}

\entry{daolle}{\headword{daolle}  \pronnote{daoɽe}
  \pos{Adverbial demonstrative}  \sense{
    \definition{allative form of da}    \example{
      \exsen{Ibi daolle beyareya.}      \extran{Let's (du.) go over there.}}}}

\entry{daollemae}{\headword{daollemae}
  \pos{Adverb}  \sense{
    \definition{allative form of da with perlative clitic}    \example{
      \exsen{Yu da daollemae obo pite we gogon.}      \extran{The fire got closer to her grass skirt.}}}}

\entry{dape}{\headword{dape}  \note{Musical instruments}  \pronnote{dape}
  \pos{Noun}  \sense{
    \definition{drum head}    \note{Musical instruments}    \example{
}    \example{
}    \example{
}}}

\entry{där}{\headword{där}  \pronnote{dər}
  \pos{Noun}  \sense{\textbf{1.} 
    \definition{pair}}  \sense{\textbf{2.} 
    \definition{to match, balance, agree}    \example{
      \exsen{Mätta de där yag.}      \extran{Match these yams (i.e. bring the same number).}}    \example{
      \exsen{Yuwog abal lla da Yesu bo pallall kuki ikopatt eka de dällätaemneyo, be oba kuki eka da ddone där gognegnän.}      \extran{Many people were making false testimonies about Jesus, but their lies didn't add up.}}}}

\entry{dara}{\headword{dara}
  \pos{Noun}  \sense{
}}

\entry{Dara}{\headword{Dara}  \pronnote{dara}
  \pos{Proper noun}  \sense{
    \definition{male personal name}}}

\entry{dara}{\headword{dara}  \pronnote{dara}
  \pos{Noun}  \sense{
    \definition{type of traditional medicine}    \example{
}    \example{
}    \example{
}}}

\entry{darab}{\headword{darab}
  \variantof{Dialectal Variant of \vartext{dalab}}  \pronnote{darab}}

\entry{daradara}{\headword{daradara}  \pronnote{daradara}
  \pos{Noun}  \sense{
    \definition{type of vine with yellow fruit that are red when ripe}    \example{
}    \example{
}    \example{
}}}

\entry{däradära}{\headword{däradära}
  \pos{Transitive verb}  \sense{
    \definition{to decorate}}}

\entry{Därall}{\headword{Därall}
  \pos{Proper noun}  \sense{
    \definition{Därall (toponym)}}}

\entry{daram~}{\headword{daram~}
  \variantof{freevarlab of \vartext{RED~}}}

\entry{daram}{\headword{daram}
  \pos{Verb}  \sense{
}}

\entry{daramdaram}{\headword{daramdaram}  \pronnote{daramdaram}
  \pos{Adverb}  \sense{
    \definition{shining brightly}    \example{
      \exsen{Yäbäd a daramdaramang gogon.}      \extran{The sun was shining brightly.}}}}

\entry{däräng}{\headword{däräng}  \pronnote{dərəŋ}
  \pos{Noun}  \sense{
    \definition{dog}    \example{
      \exsen{Däräng bo llan a tubutubu dag.}      \extran{The dog's ears are short.}}}
  \variants{\Baby Talk Variant \vartext{deng}}}

\entry{däräng olleolle}{\headword{däräng olleolle}  \note{Animal names}  \pronnote{dərəŋ oɽeoɽe}
  \pos{Noun}  \sense{
    \definition{type of possum-like animal with spotted skin}    \note{Animal names}    \example{
}    \example{
}    \example{
}}}

\entry{därängbun}{\headword{därängbun}  \note{Type of pandanus}  \pronnote{dərəŋbun}
  \pos{Noun}  \sense{
    \definition{type of pandanus with a curved fruit shaped like a dog's head}    \note{Type of pandanus}    \example{
}    \example{
}    \example{
}}}

\entry{därängg}{\headword{därängg}
  \pos{Transitive verb}  \sense{
    \definition{to lead}    \example{
      \exsen{Mälla da obom ttängäm att de ngattong dändrägän menae e.}      \extran{The woman led him outside the village.}}}}

\entry{därängge}{\headword{därängge}  \note{Palms}  \pronnote{dərəŋɡe}
  \pos{Noun}  \sense{
    \definition{small orchid with blue, yellow, white, purple flowers}    \note{Palms}    \example{
}    \example{
}    \example{
}}}

\entry{Därängge}{\headword{Därängge}
  \pos{Proper noun}  \sense{
    \definition{male personal name}}}

\entry{däränggedärängge}{\headword{däränggedärängge}  \pronnote{dərəŋɡedərəŋɡe}
  \pos{Noun}  \sense{
    \definition{large wild orchid}}}

\entry{därängkänan}{\headword{därängkänan}  \pronnote{dərəŋkənan}
  \pos{Noun}  \sense{
    \definition{leader}}}

\entry{därba}{\headword{därba}  \note{Snakes}  \pronnote{dərba}
  \pos{Noun}  \sense{
    \definition{type of snake}    \note{Snakes}    \example{
}    \example{
}    \example{
}}}

\entry{därdärag}{\headword{därdärag}  \note{Dancing}  \pronnote{dərdəraɡ}
  \pos{Adverb}  \sense{
    \definition{in pairs}    \note{Dancing}    \example{
      \exsen{Ubi mänmän de därdäragae dällädnegeyo.}      \extran{They get married with women in pairs.}}    \example{
}    \example{
}}}

\entry{Dareda}{\headword{Dareda}  \pronnote{dareda}
  \pos{Proper noun}  \sense{
    \definition{male personal name}}}

\entry{darkukiny}{\headword{darkukiny}  \note{Water}  \pronnote{darkukiɲ}
  \pos{Noun}  \sense{
    \definition{type of grass}    \note{Water}    \example{
}    \example{
}    \example{
}}}

\entry{därmae}{\headword{därmae}  \pronnote{dərmae}
  \pos{Intransitive verb}  \sense{
    \definition{to wipe out}    \example{
      \exsen{Bogenvil me lla da mäse därmae e gongkaemnegän.}      \extran{In Bougainville, people (i.e. the PNG Defense Force) were trying to wipe them (i.e. the secessionists) out.}}}}

\entry{därmir}{\headword{därmir}  \note{Trees}  \pronnote{dərmir}
  \pos{Noun}  \sense{
    \definition{type of tree that is used to treat sores}    \note{Trees}    \example{
}    \example{
}    \example{
}}}

\entry{därmir}{\headword{därmir}  \note{Names of Bananas}  \pronnote{dərmir}
  \pos{Noun}  \sense{
    \definition{type of introduced banana}    \note{Names of Bananas}    \example{
}    \example{
}    \example{
}}}

\entry{däroledärole}{\headword{däroledärole}
  \pos{Modifier}  \sense{
    \definition{dry}    \example{
      \exsen{Da bongo sana däroledärole de därunggu alle yu nägag, mɨnyi mer abal bogon.}      \extran{If you cook the dry sago with the bamboo tubes, it will be very good.}}}
  \variants{\freevarlab \vartext{droledrole}}}

\entry{därollog}{\headword{därollog}  \pronnote{dəroɽoɡ}
  \pos{Noun}  \sense{
    \definition{brolga}    \scientificname{Antigone rubicunda rubicunda}    \example{
}    \example{
}    \example{
}}}

\entry{darombe}{\headword{darombe}  \note{Musical instruments}  \pronnote{darombe}
  \pos{Noun}  \sense{
    \definition{mouth harp. quarter-moon-shaped bamboo flute with honey inside; takes a day to make}    \note{Musical instruments}    \example{
}    \example{
}    \example{
}}}

\entry{Darren}{\headword{Darren}  \pronnote{darren}
  \pos{Proper noun}  \sense{
    \definition{male personal name}}}

\entry{Daru}{\headword{Daru}
  \pos{Proper noun}  \sense{\textbf{1.} 
    \definition{Daru (capital city of Western Province, located on Daru Island)}}  \sense{\textbf{2.} 
    \definition{Daru Island}}}

\entry{därunggu}{\headword{därunggu}  \pronnote{dəruŋɡu}
  \pos{Noun}  \sense{
    \definition{bamboo tube}    \example{
      \exsen{Da bongo sana däroledärole de därunggu alle yu nägag, mɨnyi mer abal bogon.}      \extran{If you cook the dry sago in the new bamboo, it will be very good.}}}}

\entry{daudau}{\headword{daudau}  \pronnote{daudau}
  \pos{Intransitive coverb}  \sense{
    \definition{to nod}}}

\entry{dauma}{\headword{dauma}  \note{Names of Bananas}  \pronnote{dauma}
  \pos{Noun}  \sense{
    \definition{type of introduced banana}    \note{Names of Bananas}    \example{
}    \example{
}    \example{
}}}

\entry{David}{\headword{David}
  \variantof{SpellingVariant of \vartext{Deibid}}}

\entry{ddä~}{\headword{ddä~}
  \variantof{Infinitival reduplicant of \vartext{RED~}}}

\entry{dda~}{\headword{dda~}
  \variantof{Inflected Form of \vartext{RED~}}}

\entry{ddäb}{\headword{ddäb}  \note{Parts of the body}  \pronnote{ɖ͡ʐəb}
  \pos{Noun}  \sense{
    \definition{anus}    \note{Parts of the body}    \example{
}    \example{
}    \example{
}}}

\entry{ddadd}{\headword{ddadd}  \lexeme{ddadd}  \pronnote{ɖ͡ʐaɖ͡ʐ}
  \pos{Noun}  \sense{
    \definition{type of large tree with white flowers and blue fruit; wood used for making canoes}}}

\entry{ddäddäbeabag}{\headword{ddäddäbeabag}  \pronnote{ɖ͡ʐəɖ͡ʐəbeabaɡ}
  \pos{Modifier}  \sense{
    \definition{independent, resourceful}}}

\entry{ddäddäg}{\headword{ddäddäg}  \lexeme{ddäddäg}  \pronnote{ɖ͡ʐəɖ͡ʐəɡ}
  \pos{Transitive verb}  \sense{
    \definition{to peel, remove}    \example{
      \exsen{Ngäna ma papek de näddägaban.}      \extran{I removed the walls.}}}}

\entry{ddäddäg}{\headword{ddäddäg}  \lexeme{ddäddäg}  \pronnote{ɖ͡ʐəɖ͡ʐəɡ}
  \pos{Noun}  \sense{\textbf{1.} 
    \definition{edible animal, game, meat}    \example{
      \exsen{Bogo ddäddäg de dokonggän oblle.}      \extran{He cut the meat for her.}}}  \sense{\textbf{2.} 
    \definition{to eat meat; bite}    \example{
      \exsen{Ddob ddägnan ma gullem a dadeg.}      \extran{There are some edible snakes.}}    \example{
      \exsen{Ge ddia de däbänya wa däddägaebeya.}      \extran{We cut this deer and ate it.}}    \example{
      \exsen{Llɨg a ddäddäg de näddägan.}      \extran{The boy ate the meat.}}    \example{
      \exsen{Käza da ge ttang me daddägän ngänäm.}      \extran{The crocodile bit me in this hand.}}    \example{
      \exsen{Da bam kungge da naddägän, bongo mɨnyi ttattlle ampug.}      \extran{If a spider bites you, you will become very ill.}}}  \sense{\textbf{3.} 
    \definition{bite (of an animal)}    \example{
      \exsen{Ddob llayabira komo ddäddäg a malla ttällanenang dan.}      \extran{For some people, the centiepede's bite isn't painful.}}}  \sense{\textbf{4.} 
    \definition{hunger for meat}    \example{
      \exsen{Ngämim ddäddäg abal da deyagnän.}      \extran{We (excl.) were very hungry for meat (lit. hunger got us).}}}
  \variants{\Baby Talk Variant \vartext{tata}}}

\entry{ddäddäg}{\headword{ddäddäg}  \lexeme{ddäddäg}  \pronnote{ɖ͡ʐəɖ͡ʐəɡ}
  \pos{Transitive verb}  \sense{
    \definition{to bind}    \example{
      \exsen{Sana dor däpittaemeya a däddägaemeya.}      \extran{We wove and bound the sago stalks.}}}}

\entry{ddäddäg}{\headword{ddäddäg}  \lexeme{ddäddäg}  \pronnote{ɖ͡ʐəɖ͡ʐəɡ}
  \pos{Transitive verb}  \sense{\textbf{1.} 
    \definition{to look, watch}}  \sense{\textbf{2.} 
    \definition{to look, watch}}}

\entry{ddäddäg}{\headword{ddäddäg}  \lexeme{ddäddäg}  \pronnote{ɖ͡ʐəɖ͡ʐəɡ}
  \pos{Transitive verb}  \sense{
    \definition{to pain, ache, hurt}    \example{
      \exsen{Ge ute da naddägnan ngänäm.}      \extran{This sore is hurting me.}}}}

\entry{ddäddäg ttoe}{\headword{ddäddäg ttoe}
  \pos{Noun}  \sense{
    \definition{leather, hide, animal skin}    \example{
      \exsen{ddäddäg ttoe nyäng}      \extran{leather bag}}}}

\entry{ddäddäl}{\headword{ddäddäl}
  \pos{Transitive verb}  \sense{
    \definition{to shove}}}

\entry{ddäddäll}{\headword{ddäddäll}  \note{Weather}  \pronnote{ɖ͡ʐəɖ͡ʐəɽ}
  \pos{Noun}  \sense{
    \definition{thunder}    \note{Weather}    \example{
}    \example{
}    \example{
}}}

\entry{ddaddällɨg}{\headword{ddaddällɨg}  \pronnote{ɖ͡ʐaɖ͡ʐəɽɪɡ}
  \pos{Transitive verb}  \sense{
    \definition{to destroy}    \example{
      \exsen{Kuddäll lla da bäne ma de däddellɨgallo.}      \extran{They destroy the house of the dead man.}}}
  \variants{\Unspecified Variant \vartext{ddaddellɨg}}
  \variants{\SpellingVariant \vartext{ddaddlläg}}}

\entry{ddaddell}{\headword{ddaddell}
  \pos{Transitive verb}  \sense{
    \definition{to destroy}}}

\entry{ddaddellɨg}{\headword{ddaddellɨg}
  \variantof{Unspecified Variant of \vartext{ddaddällɨg}}}

\entry{ddaddlläg}{\headword{ddaddlläg}
  \variantof{SpellingVariant of \vartext{ddaddällɨg}}}

\entry{ddaddu}{\headword{ddaddu}  \pronnote{ɖ͡ʐaɖ͡ʐu}
  \pos{Transitive verb}  \sense{
    \definition{to remove a shoot}}}

\entry{ddaebän}{\headword{ddaebän}  \pronnote{ɖ͡ʐaebən}
  \pos{Transitive verb}  \sense{\textbf{1.} 
    \definition{to divorce}    \example{
      \exsen{Ngäna mɨnyi mälla de bäddebän.}      \extran{I will divorce my wife.}}}  \sense{\textbf{2.} 
    \definition{to adopt}    \example{
      \exsen{Ngäna Jepet bom ngämlle llɨg de däddebän.}      \extran{I adopted Jepet as my son.}}}}

\entry{ddäg}{\headword{ddäg}  \pronnote{ɖ͡ʐəɡ}
  \pos{Noun}  \sense{\textbf{1.} 
    \definition{back}    \example{
      \exsen{Ddia ddäg de nällɨt.}      \extran{[You] cut the back of the deer.}}}  \sense{\textbf{2.} 
    \definition{outside}    \example{
      \exsen{Bongo polle ddäg me nagrine.}      \extran{[You] stay outside the fence.}}}}

\entry{ddäg kutt}{\headword{ddäg kutt}
  \pos{Noun}  \sense{
    \definition{backbone, spine}}}

\entry{ddägaddäge}{\headword{ddägaddäge}
  \pos{Intransitive verb}  \sense{\textbf{1.} 
    \definition{to write}    \example{
      \exsen{Amo ingoll anykeanyke da bin peyang goddgenegän gänya mani me?}      \extran{Whose likeness and name are inscribed on this coin?}}}  \sense{\textbf{2.} 
    \definition{to branch out}}}

\entry{ddägatt}{\headword{ddägatt}
  \variantof{Fast Speech Variant of \vartext{ddägattalle}}}

\entry{ddägattalle}{\headword{ddägattalle}  \pronnote{ɖ͡ʐəɡaʈ͡ʂaɽe}
  \pos{Adverb}  \sense{
    \definition{later, after}    \example{
      \exsen{ttongo pazi ddägatt}      \extran{one year later}}    \example{
      \exsen{au ddägattalle}      \extran{after the burial}}}
  \variants{\Fast Speech Variant \vartext{ddägatt}}}

\entry{ddägddäg}{\headword{ddägddäg}
  \pos{Adverb}  \sense{
    \definition{from behind}    \example{
      \exsen{Bogo ddägddäg deyarän.}      \extran{He came from behind.}}}}

\entry{ddage}{\headword{ddage}  \lexeme{ddage}  \pronnote{ɖ͡ʐaɡe}
  \pos{Noun}  \sense{\textbf{1.} 
    \definition{branch}    \example{
      \exsen{Kollko ddage me godämenän.}      \extran{He sat on the breadfruit tree branch.}}}  \sense{\textbf{2.} 
    \definition{stream, tributary}    \example{
      \exsen{Ngämi däbe ddage de danttogeya.}      \extran{We (excl.) followed that stream.}}}}

\entry{ddage bun}{\headword{ddage bun}  \note{Water}  \pronnote{ɖ͡ʐaɡe bun}
  \pos{Noun}  \sense{
    \definition{river mouth}    \note{Water}    \example{
      \exsen{ddage bun e}      \extran{downstream}}}}

\entry{ddage llätt}{\headword{ddage llätt}  \note{Water}  \pronnote{ɖ͡ʐaɡe ɽəʈ͡ʂ}
  \pos{Noun}  \sense{
    \definition{river source}    \note{Water}    \example{
      \exsen{ddage llätt e}      \extran{upstream}}    \example{
}    \example{
}}}

\entry{ddageddage}{\headword{ddageddage}  \pronnote{ɖ͡ʐaɡeɖ͡ʐaɡe}
  \pos{Noun}  \sense{\textbf{1.} 
    \definition{the fifth stage of sago growth in which the inflorescence has emerged}    \example{
      \exsen{Bäne sana da zäme ddageddage gogon.}      \extran{Your sago palm is already forming flowers.}}}  \sense{\textbf{2.} 
    \definition{with many branches}    \example{
      \exsen{Ge manggo da ddageddage dan.}      \extran{That mango tree has too many branches.}}}}

\entry{ddägnan ma}{\headword{ddägnan ma}  \pronnote{ɖ͡ʐəɡnan ma}
  \pos{Modifier}  \sense{
    \definition{edible}}}

\entry{ddäkop}{\headword{ddäkop}
  \variantof{Unspecified Variant of \vartext{ddokop}}  \note{Parts of the body}  \pronnote{ɖ͡ʐəkop}}

\entry{ddäll}{\headword{ddäll}
  \pos{Intransitive verb}  \sense{
    \definition{to arrive}    \example{
      \exsen{Ubi ddob ttängäm atta guddällmamän.}      \extran{They arrived from another village.}}}}

\entry{ddäll}{\headword{ddäll}  \note{Parts of a bird}  \pronnote{ɖ͡ʐəɽ}
  \pos{Noun}  \sense{\textbf{1.} 
    \definition{chest}    \note{Parts of a bird}    \example{
      \exsen{Llɨg kälsre da mälla bo ddäll me dan.}      \extran{The child is on the woman's chest.}}}  \sense{\textbf{2.} 
    \definition{part of the sago trunk closest to the leaves before the shoot}    \note{Parts of a bird}}  \sense{\textbf{3.} 
    \definition{ten (lit. chest; body counting numeral)}    \note{Parts of a bird}}}

\entry{ddäll kom}{\headword{ddäll kom}  \note{Parts of the body}  \pronnote{ɖ͡ʐəɽ kom}
  \pos{Noun}  \sense{
    \definition{chest hair}    \note{Parts of the body}    \example{
}    \example{
}    \example{
}}}

\entry{ddäll kutt}{\headword{ddäll kutt}  \note{Parts of a reptile}  \pronnote{ɖ͡ʐəɽ kuʈ͡ʂ}
  \pos{Noun}  \sense{
    \definition{sternum, breastbone}    \note{Parts of a reptile}    \example{
}    \example{
}    \example{
}}}

\entry{ddälläb}{\headword{ddälläb}  \pronnote{ɖ͡ʐəɽəb}
  \pos{Intransitive verb}  \sense{
    \definition{to fall over (of a tree)}    \example{
      \exsen{Llo käkäp a goddälläbnalle.}      \extran{Half of the tree fell over.}}}}

\entry{ddällbän}{\headword{ddällbän}
  \pos{Transitive verb}  \sense{
}}

\entry{ddällgoe}{\headword{ddällgoe}  \pronnote{ɖ͡ʐəɽɡoe}
  \pos{Transitive verb}  \sense{\textbf{1.} 
    \definition{to go through thick bush in search of animals}    \example{
      \exsen{Lla da däräng peyang tawa de daddällgoeneyo.}      \extran{The man perused the swamp together with the dog.}}}  \sense{\textbf{2.} 
    \definition{to disturb a beehive}}}

\entry{ddällombog}{\headword{ddällombog}  \note{Dancing}  \pronnote{ɖ͡ʐəɽomboɡ}
  \pos{Transitive verb}  \sense{
    \definition{to miss}    \note{Dancing}    \example{
      \exsen{Bogo nyongo de nanddällbogan.}      \extran{He missed the road.}}    \example{
      \exsen{Bogo angde tuk i nallnan mo de nanddällbogan ttongo.}      \extran{She missed a step going up the ladder.}}    \example{
}}}

\entry{ddällpoyampoyam}{\headword{ddällpoyampoyam}  \note{Mushrooms}  \pronnote{ddällpo'yampo'yam}
  \pos{Noun}  \sense{
    \definition{type of mushroom}    \note{Mushrooms}    \example{
}    \example{
}    \example{
}}}

\entry{ddallwe}{\headword{ddallwe}  \note{Trees}  \pronnote{ɖ͡ʐaɽwe}
  \pos{Noun}  \sense{
    \definition{type of tree}    \note{Trees}    \example{
}    \example{
}    \example{
}}}

\entry{ddäma}{\headword{ddäma}  \pronnote{ɖ͡ʐəma}
  \pos{Noun}  \sense{\textbf{1.} 
    \definition{basket}    \example{
      \exsen{Naomi bo llɨg a ddäma me dan angällbänan.}      \extran{Naomi's baby woke up in the baby basket.}}}  \sense{\textbf{2.} 
    \definition{pouch of a marsupial}}  \sense{\textbf{3.} 
    \definition{uterus}}}

\entry{ddäma mu}{\headword{ddäma mu}  \note{Marriage}  \pronnote{ɖ͡ʐəma mu}
  \pos{Noun}  \sense{
    \definition{birth payment made to the maternal uncle}    \note{Marriage}    \example{
      \exsen{Ngämi ngäma pope aba pate zime ddäma mu de dangeseya.}      \extran{We (excl.) already paid our uncles the basket payment.}}}}

\entry{ddamba}{\headword{ddamba}  \note{Parts of a fish}  \pronnote{ɖ͡ʐamba}
  \pos{Noun}  \sense{\textbf{1.} 
    \definition{wing}    \note{Parts of a fish}}  \sense{\textbf{2.} 
    \definition{pectoral fin}    \note{Parts of a fish}    \example{
}    \example{
}    \example{
}}}

\entry{ddamba mit}{\headword{ddamba mit}  \note{Parts of an insect}  \pronnote{ɖ͡ʐamba mit}
  \pos{Noun}  \sense{
    \definition{metathorax}    \note{Parts of an insect}    \example{
}    \example{
}    \example{
}}}

\entry{Ddämir}{\headword{Ddämir}
  \pos{Proper noun}  \sense{
    \definition{Ddamir (toponym)}}}

\entry{ddämoemkäp kuibiag}{\headword{ddämoemkäp kuibiag}  \note{Snakes}  \pronnote{ɖ͡ʐəmoemkəp kuibiaɡ}
  \pos{Noun}  \sense{
    \definition{type of python}    \note{Snakes}    \example{
}    \example{
}    \example{
}}}

\entry{ddän}{\headword{ddän}
  \pos{Transitive verb}  \sense{
    \definition{to pick, gather, harvest}    \example{
      \exsen{Dädme ddob wayati de däddaebeya.}      \extran{We harvested some watermelons there.}}}}

\entry{ddänddäl}{\headword{ddänddäl}  \pronnote{ɖ͡ʐənɖ͡ʐəl}
  \pos{Intransitive verb}  \sense{
    \definition{to climb}    \example{
      \exsen{Baet a llo ddage me gonddälän.}      \extran{The cuscus climbed on the tree branch.}}}}

\entry{ddänddäm}{\headword{ddänddäm}  \pronnote{ɖ͡ʐənɖ͡ʐəm}
  \pos{Intransitive verb}  \sense{\textbf{1.} 
    \definition{to drown}    \example{
      \exsen{Obom dandämoyän ttu we, angde obom ttu we dandämoyän bogo ddone dängllawän, be bogo gonddämän.}      \extran{She pushed him into the deep, and when she pushed him into the deep, he didn't swim; he drowned.}}    \example{
      \exsen{Kauga angde Matthew bom ikop dägagän ddänddäm me, bogo mängalae gogbänän walle we, Matthew bom dirängänän.}      \extran{When Kauga saw Matthew drowning, he quickly jumped into the water and lifted Matthew out.}}}  \sense{\textbf{2.} 
    \definition{to set (of the moon)}    \example{
      \exsen{Kok a de ddänddäm allan.}      \extran{The moon is setting.}}}  \sense{\textbf{3.} 
    \definition{to drown, sink}    \example{
      \exsen{Kollba da ddone dan, ede enda nett de nänddäman?}      \extran{There are no fish, so what sank the net?}}}  \sense{\textbf{4.} 
    \definition{to worry}    \example{
      \exsen{Bogo ereya gagäll de bällnan eran. Diba ttoen da obo näkäp de ddänddäm eran.}      \extran{She's remembering bad news. This thing makes her worry.}}}  \sense{\textbf{5.} 
}}

\entry{ddänddängeny}{\headword{ddänddängeny}  \pronnote{ɖ͡ʐənɖ͡ʐəŋeɲ}
  \pos{Adverb}  \sense{
    \definition{immediately}    \example{
      \exsen{Ngäna ddänddängeny me gogäbän gall atta.}      \extran{I jumped out of the canoe immediately.}}}}

\entry{ddängall}{\headword{ddängall}  \note{Bees}  \pronnote{'ddängall}
  \pos{Noun}  \sense{
    \definition{type of stinging bee found in trees}    \note{Bees}    \example{
}    \example{
}    \example{
}}}

\entry{ddänggab}{\headword{ddänggab}
  \pos{Transitive verb}  \sense{
    \definition{to hold, grab (with one's teeth)}    \example{
      \exsen{Sɨmell a ngoi alle ere ada dänddgabän.}      \extran{The pig held it with its teeth.}}}}

\entry{ddänggaddängge}{\headword{ddänggaddängge}  \pronnote{ddäng'gaddäng'ge}
  \pos{Transitive verb}  \sense{\textbf{1.} 
    \definition{to crucify}    \example{
      \exsen{Yesu bom llo tärpamatt toko me danddgeyo.}      \extran{They crucified Jesus on a wooden cross.}}}  \sense{\textbf{2.} 
    \definition{to catch, trap}    \example{
      \exsen{Martin käza de danddgewän.}      \extran{Martin trapped the crocodile.}}}}

\entry{ddangoe}{\headword{ddangoe}  \pronnote{ɖ͡ʐaŋoe}
  \pos{Transitive verb}  \sense{
    \definition{to force to do}    \example{
      \exsen{Andrew Yuga bom era tongoe e daddengoenän.}      \extran{Andrew was forcing Yuga to play a game.}}}}

\entry{ddangol}{\headword{ddangol}  \pronnote{ɖ͡ʐaŋol}
  \pos{Noun}  \sense{
    \definition{type of spear}}}

\entry{ddänmäll}{\headword{ddänmäll}  \pronnote{ɖ͡ʐənməɽ}
  \pos{Transitive verb}  \sense{
    \definition{to struggle}    \example{
      \exsen{Ngäna gall bllablle we däddänmällneg.}      \extran{I struggled to find the canoe.}}}}

\entry{ddapall}{\headword{ddapall}  \lexeme{ddapall}  \pronnote{'ddapall}
  \pos{Noun}  \sense{\textbf{1.} 
    \definition{sky}    \example{
      \exsen{Ddapall a siremang gogon.}      \extran{The sky got dark.}}}  \sense{\textbf{2.} 
    \definition{heaven}    \example{
      \exsen{Ddapall e ibi allan.}      \extran{I am going to heaven.}}}}

\entry{ddapall käkan}{\headword{ddapall käkan}  \note{Nature}  \pronnote{ɖ͡ʐapaɽ kəkan}
  \pos{Noun}  \sense{
    \definition{cloud}    \note{Nature}    \example{
      \exsen{Yäbäd de ddapall käkan da dakonewän.}      \extran{Clouds covered the sun.}}    \example{
}    \example{
}}}

\entry{ddapall källa}{\headword{ddapall källa}
  \pos{Noun}  \sense{
    \definition{cloud}}}

\entry{ddapall ma}{\headword{ddapall ma}  \lexeme{ddapall ma}  \pronnote{ɖ͡ʐapaɽ ma}
  \pos{Noun}  \sense{
    \definition{heaven}    \example{
      \exsen{Ddapall ma da kili ma ttängäm dan.}      \extran{Heaven is a place of happiness.}}}}

\entry{ddel}{\headword{ddel}
  \pos{Transitive verb}  \sense{
    \definition{to explain}    \example{
      \exsen{Obo kollmällang aba pate bogo tämamae eka midd de däddelnegnän.}      \extran{He was explaining all the meanings to his disciples.}}}}

\entry{Ddele}{\headword{Ddele}
  \pos{Proper noun}  \sense{
    \definition{Ddele (toponym)}}}

\entry{Ddeleag}{\headword{Ddeleag}  \note{Dialect names}  \pronnote{ɖ͡ʐeleaɡ}
  \pos{Proper noun}  \sense{
    \definition{Ende dialect spoken in Ddele}    \note{Dialect names}    \example{
}    \example{
}    \example{
}}}

\entry{Ddelema}{\headword{Ddelema}  \pronnote{ɖ͡ʐelema}
  \pos{Proper noun}  \sense{
    \definition{male personal name}}}

\entry{ddia}{\headword{ddia}  \lexeme{ddia}  \pronnote{ɖ͡ʐia}
  \pos{Noun}  \sense{
    \definition{deer}    \example{
      \exsen{Ddia da dindu allan.}      \extran{The deer is running.}}}}

\entry{ddo~}{\headword{ddo~}
  \variantof{Infinitival reduplicant of \vartext{RED~}}}

\entry{ddob}{\headword{ddob}  \pronnote{ɖ͡ʐob}
  \pos{Quantifier}  \sense{\textbf{1.} 
    \definition{some}    \example{
      \exsen{Ngämira ddob kollba de deyasiyu.}      \extran{They gave us (excl.) some fish.}}}  \sense{\textbf{2.} 
    \definition{other}    \example{
      \exsen{Ako ddob kemibi ttoen a dadeg ngasnen ma da.}      \extran{There are also many other things to do.}}}}

\entry{ddobae}{\headword{ddobae}  \pronnote{ɖ͡ʐobae}
  \pos{Adverb}  \sense{
    \definition{very}    \example{
      \exsen{ddobae mangallang lla}      \extran{very strong person}}    \example{
      \exsen{ddobae melemang lla}      \extran{very hardworking person}}}}

\entry{ddobae~}{\headword{ddobae~}
  \variantof{freevarlab of \vartext{RED~}}}

\entry{ddobaeddobae}{\headword{ddobaeddobae}  \pronnote{ɖ͡ʐobaeɖ͡ʐobae}
  \pos{Adverb}  \sense{
    \definition{very, extremely}    \example{
      \exsen{Bundae ddobaeddobae mikutt gogon.}      \extran{Bundae was very, very angry.}}}}

\entry{ddobag}{\headword{ddobag}  \pronnote{ɖ͡ʐobaɡ}
  \pos{Pronoun}  \sense{
    \definition{some, others}    \example{
      \exsen{Ddobag a Ende eka de panypeny erallo, ddobag a kllokloe panypeny erallo.}      \extran{Some are speaking Ende; others are speaking a mix.}}}}

\entry{ddoddllem}{\headword{ddoddllem}
  \variantof{Fast Speech Variant of \vartext{ddoddollem}}}

\entry{ddoddollem}{\headword{ddoddollem}
  \pos{Transitive/Intransitive coverb}  \sense{
    \definition{to make noise}    \example{
      \exsen{Bogo auma me do wattällang a endagaeya däbem ddoddollem dägnegän.}      \extran{He used the those things left at the grave to make noise.}}}
  \variants{\Fast Speech Variant \vartext{ddoddllem}}}

\entry{ddoga}{\headword{ddoga}  \note{Trees}  \pronnote{ɖ͡ʐoɡa}
  \pos{Noun}  \sense{
    \definition{type of tree}    \note{Trees}    \example{
}    \example{
}    \example{
}}}

\entry{ddogllop}{\headword{ddogllop}
  \variantof{Fast Speech Variant of \vartext{ddogollop}}  \pronnote{ɖ͡ʐoɡɽop}}

\entry{ddogoll}{\headword{ddogoll}  \pronnote{ɖ͡ʐoɡoɽ}
  \pos{Transitive verb}  \sense{\textbf{1.} 
    \definition{to put together}    \example{
      \exsen{Angde tine pittnen a gottamänän, dibaballe ddäganen de dängkam.}      \extran{After the sago leaf weaving finished, I started putting together the roof.}}    \example{
      \exsen{Lla da päre abo gobällan kottllam bo gollob pallkepallke de dänglläbeyo a kottllam bo patme däddäganeyo.}      \extran{Then the people sadly went to the turtle, collected all the pieces of his shell, and put him back together.}}}  \sense{\textbf{2.} 
    \definition{to stick}    \example{
      \exsen{Tongle da addgollan.}      \extran{The leech is stuck.}}}  \sense{\textbf{3.} 
    \definition{to join}    \example{
      \exsen{Ngäna goddägol.}      \extran{I joined.}}}  \sense{\textbf{4.} 
    \definition{to lean}    \example{
      \exsen{Llo pätt me ddogollag dan.}      \extran{It's leaning on the tree trunk.}}}}

\entry{ddogollop}{\headword{ddogollop}  \note{Parts of a reptile}  \pronnote{ɖ͡ʐoɡoɽop}
  \pos{Noun}  \sense{
    \definition{reptile scale}    \note{Parts of a reptile}    \example{
}    \example{
}    \example{
}}
  \variants{\Fast Speech Variant \vartext{ddogllop}}}

\entry{ddokddok}{\headword{ddokddok}  \note{Traditional Arrows}  \pronnote{'ddok'ddok}
  \pos{Modifier}  \sense{
    \definition{blunt}    \note{Traditional Arrows}    \example{
      \exsen{Tätäm ibik a ddokddok gogon.}      \extran{Yesterday, the digging stick got blunt.}}}}

\entry{ddokddok}{\headword{ddokddok}  \note{Traditional Arrows}  \pronnote{'ddok'ddok}
  \pos{Noun}  \sense{
    \definition{type of spear}    \note{Traditional Arrows}    \example{
}    \example{
}    \example{
}}}

\entry{ddokop}{\headword{ddokop}
  \pos{Noun}  \sense{
    \definition{kidney}}
  \variants{\Unspecified Variant \vartext{ddäkop}}}

\entry{ddol}{\headword{ddol}  \pronnote{ɖ͡ʐol}
  \pos{Noun}  \sense{
    \definition{foam, bubbles, gas}    \example{
      \exsen{Ttongo lla da ddone mullae dan sisor gärep kae ine de gudne nyäng e banzämaeyän, adawatta gärep kae ine da mɨnyi ddol bogon a nyäng de bungolltän.}      \extran{A man cannot pour new wine into an old bag, because the wine will become gassy and it will make the bag explode.}}}}

\entry{ddolag}{\headword{ddolag}
  \pos{Modifier}  \sense{
    \definition{foamy, bubbly, gassy}}}

\entry{ddollombog}{\headword{ddollombog}
  \pos{Transitive verb}  \sense{
    \definition{to misspeak, speak with mistakes}    \example{
      \exsen{Tame eka de da ngäna bäpany, da erame ngäna banddollbog, ede mäzi ngänäm ttättle bageyo.}      \extran{When I speak Taeme, whenever I make mistakes, they will correct me.}}}}

\entry{ddonddo}{\headword{ddonddo}  \pronnote{ɖ͡ʐonɖ͡ʐo}
  \pos{Modifier}  \sense{\textbf{1.} 
    \definition{proud}    \example{
      \exsen{Ngämi bam ge ddobae ddonddo nalla.}      \extran{We (excl.) are very proud of you for this.}}}  \sense{\textbf{2.} 
    \definition{to boast, brag}    \example{
      \exsen{Ubi oba mängall de gonddoneyo.}      \extran{They were boasting about their strength.}}}}

\entry{ddone}{\headword{ddone}  \pronnote{ɖ͡ʐone}
  \pos{Interjection}  \sense{\textbf{1.} 
    \definition{no}    \example{
      \exsen{― Bäne kokok a daden? ― Ddone.}      \extran{― Do you have grandchildren? ― No.}}    \example{
      \exsen{Be bogo ddone allan.}      \extran{But he said no.}}}  \sense{\textbf{2.} 
    \definition{not}    \example{
      \exsen{Da bongo ddone kandärmang ekag, bäne ikop a säremang bognegnän.}      \extran{If you don't say sorry, you will go blind.}}}  \sense{\textbf{3.} 
}}

\entry{ddone ada}{\headword{ddone ada}  \pronnote{ɖ͡ʐone ada}
  \pos{Adverb}  \sense{
    \definition{very, a lot (antiphrasis)}    \example{
      \exsen{Llamda da ddone ada mikutt gogon obo pate.}      \extran{The old man got very angry at them.}}}}

\entry{ddone amom}{\headword{ddone amom}  \pronnote{ɖ͡ʐone amom}
  \pos{Personal pronoun}  \sense{
    \definition{accusative form of ddone aya}    \example{
      \exsen{Ubi ddone amom ikop dägaeyo.}      \extran{They didn't see anyone.}}}}

\entry{ddone any}{\headword{ddone any}
  \pos{Pronoun}  \sense{
    \definition{nothing}    \example{
      \exsen{Ddone any a lel gogalle.}      \extran{She wasn't afraid of anything.}}}}

\entry{ddone aya}{\headword{ddone aya}  \pronnote{ɖ͡ʐone aja}
  \pos{Personal pronoun}  \sense{
    \definition{no one, anyone (nominative form)}    \example{
      \exsen{Ddone aya mɨnyi ngänäm batramän ttongo ttängäm e.}      \extran{No one will take me away to another village.}}}}

\entry{ddongddong}{\headword{ddongddong}  \lexeme{ddongddong}  \pronnote{ɖ͡ʐoŋɖ͡ʐoŋ}
  \pos{Noun}  \sense{
    \definition{thick cluster of short grass}}}

\entry{ddugwemddugwem}{\headword{ddugwemddugwem}  \pronnote{ɖ͡ʐuɡwemɖ͡ʐuɡwem}
  \pos{Adverb}  \sense{
    \definition{stomping}}}

\entry{ddumbi}{\headword{ddumbi}  \note{Traditional Arrows}  \pronnote{'ddumbi}
  \pos{Noun}  \sense{
    \definition{type of spear topped with the claw of a cassowary}    \note{Traditional Arrows}    \example{
}    \example{
}    \example{
}}}

\entry{ddungg}{\headword{ddungg}  \pronnote{ɖ͡ʐuŋɡ}
  \pos{Transitive verb}  \sense{
    \definition{to decapitate}    \example{
      \exsen{Zon bom inkätt danddugän.}      \extran{He decapitated John.}}}}

\entry{de}{\headword{de}
  \variantof{Fast Speech Variant of \vartext{daeya}}}

\entry{de}{\headword{de}  \pronnote{de}
  \pos{Demonstrative}  \sense{\textbf{1.} 
    \definition{that (mesial determiner)}    \example{
      \exsen{de llɨg kälsre}      \extran{that small boy}}}  \sense{\textbf{2.} 
    \definition{there (mesial)}    \example{
      \exsen{Käsre de amne me dantämonän, llo gäba me.}      \extran{Then, he waited there in the middle, in the shade of the tree.}}}  \sense{\textbf{3.} 
    \definition{next}    \example{
      \exsen{de sände me}      \extran{next week}}}}

\entry{de}{\headword{de}
  \variantof{Unspecified Variant of \vartext{da}}  \pronnote{de}}

\entry{=de}{\headword{=de}  \pronnote{de}
  \pos{Nominal enclitic}  \sense{\textbf{1.} 
    \definition{accusative clitic}    \example{
      \exsen{Ngämi nge de dibenyeya.}      \extran{We planted the coconut.}}}  \sense{\textbf{2.} 
    \definition{argument focus marker}    \example{
      \exsen{Ddia da angde yinu kuddäll gogon, kottllam a de deyanzigän.}      \extran{He had been really dead asleep, and turtle had passed him!}}}}

\entry{deba}{\headword{deba}
  \variantof{SpellingVariant of \vartext{diba}}}

\entry{debag}{\headword{debag}
  \variantof{SpellingVariant of \vartext{dibag}}}

\entry{Deboa}{\headword{Deboa}  \pronnote{deboa}
  \pos{Proper noun}  \sense{
    \definition{male personal name}}}

\entry{ded}{\headword{ded}
  \pos{Adverbial demonstrative}  \sense{
    \definition{there}}}

\entry{dedam}{\headword{dedam}  \pronnote{dedam}
  \pos{Adverbial demonstrative}  \sense{
    \definition{then, at that time}    \example{
      \exsen{Skul täräp a dedam ddone deya.}      \extran{There was no schooltime back then.}}    \example{
      \exsen{Angde nyongo amne we gogaebne, yogoll a dedam disamän.}      \extran{When we were halfway through the journey, that's when the rain stopped.}}}}

\entry{dedi}{\headword{dedi}
  \pos{Kinship noun}  \sense{
    \definition{father}}
  \variants{\SpellingVariant \vartext{dadi}}}

\entry{dedme}{\headword{dedme}  \pronnote{dedme}
  \pos{Adverbial demonstrative}  \sense{
    \definition{there}    \example{
      \exsen{Dedme ddone dan be kumuddäga nyäng daebeg.}      \extran{There's nothing there except for three bags.}}    \example{
      \exsen{Ge llɨg a dedme gunziwän.}      \extran{This boy settled there.}}}}

\entry{dedo}{\headword{dedo}
  \pos{Demonstrative}  \sense{
    \definition{there}    \example{
      \exsen{Dedo otät a ddone dan?}      \extran{Isn't there any food?}}    \example{
      \exsen{Ngäna dedo melem gogne.}      \extran{I was working there.}}}}

\entry{dedre}{\headword{dedre}  \pronnote{dedre}
  \pos{Intransitive verb}  \sense{
    \definition{to descend, go down}    \example{
      \exsen{Ubi tutu atta walle we godrenegnän.}      \extran{They were descending from the mountain to the river.}}}}

\entry{Deibid}{\headword{Deibid}  \pronnote{deibid}
  \pos{Proper noun}  \sense{
    \definition{male personal name}}
  \variants{\SpellingVariant \vartext{David}}}

\entry{dektta}{\headword{dektta}
  \pos{Noun}  \sense{
    \definition{doctor}    \example{
      \exsen{Ubi mɨnyi dektta we moko bognegän.}      \extran{They will want a doctor.}}}}

\entry{del}{\headword{del}  \pronnote{del}
  \pos{Noun}  \sense{
    \definition{coconut lorikeet}    \scientificname{Trichoglossus haematodus nigrogularis}    \example{
}    \example{
}    \example{
}}}

\entry{Delema}{\headword{Delema}
  \pos{Proper noun}  \sense{
    \definition{Delema (toponym)}}}

\entry{dem}{\headword{dem}  \note{Sago palms}  \pronnote{dem}
  \pos{Noun}  \sense{
    \definition{type of sago used for paints}    \note{Sago palms}    \example{
}    \example{
}    \example{
}}}

\entry{dem}{\headword{dem}
  \pos{Nominal demonstrative}  \sense{
}}

\entry{deng}{\headword{deng}
  \variantof{Baby Talk Variant of \vartext{däräng}}  \pronnote{deŋ}}

\entry{deodeo}{\headword{deodeo}  \note{Poison}  \pronnote{deodeo}
  \pos{Noun}  \sense{
    \definition{termite}    \note{Poison}    \example{
}    \example{
}    \example{
}}}

\entry{derägmäll}{\headword{derägmäll}
  \pos{Transitive verb}  \sense{
    \definition{to rebuke, scold}    \example{
      \exsen{Yesu ubim dädergmällnegän ubim adawatta oba tikop a llokollokott dagaeya.}      \extran{Jesus rebuked them (i.e. his disciples) for the hardness of their hearts.}}}}

\entry{Derideri}{\headword{Derideri}  \note{Place names}  \pronnote{derideri}
  \pos{Proper noun}  \sense{
    \definition{Derideri (Nambo-speaking village in Morehead Rural LLG; east of Morehead)}    \note{Place names}    \example{
}    \example{
}    \example{
}}}

\entry{Dettall}{\headword{Dettall}
  \pos{Proper noun}  \sense{
    \definition{Dettall (toponym)}}}

\entry{Dewara}{\headword{Dewara}  \pronnote{dewara}
  \pos{Proper noun}  \sense{
    \definition{Dewara (Were/Kiunum-speaking village in Gogodala Rural LLG, along the Fly River; from Limol, one must pass through Upiara and Kondobol)}    \example{
}    \example{
}    \example{
}}}

\entry{deya}{\headword{deya}
  \variantof{freevarlab of \vartext{daeya}}  \pronnote{deja}}

\entry{diaba}{\headword{diaba}  \note{Traditional Arrows}  \pronnote{diaba}
  \pos{Noun}  \sense{
    \definition{type of spear}    \note{Traditional Arrows}    \example{
}    \example{
}    \example{
}}}

\entry{Diandra}{\headword{Diandra}
  \pos{Proper noun}  \sense{
    \definition{female personal name}}}

\entry{diba}{\headword{diba}  \pronnote{diba}
  \pos{Nominal demonstrative}  \sense{\textbf{1.} 
    \definition{that (medial demonstrative)}    \example{
      \exsen{Bongo Saken bo amamär de ewede diba llo me nganae eralle?}      \extran{Why are you coiling Saken's rope around that tree?}}}  \sense{\textbf{2.} 
}
  \variants{\SpellingVariant \vartext{deba}}}

\entry{dibaballe}{\headword{dibaballe}  \pronnote{dibabaɽe}
  \pos{Adverbial demonstrative}  \sense{
    \definition{ablative form of diba; then, thereupon}    \example{
      \exsen{Mällause da dagirnän, dibaballe goeg a dättämän.}      \extran{The old woman was staying; thereafter, the garden burned.}}}}

\entry{dibaballemae}{\headword{dibaballemae}
  \pos{Nominal demonstrative}  \sense{
}}

\entry{dibaeya}{\headword{dibaeya}
  \pos{Copular verb}  \sense{
    \definition{cop.that.pst.sgS}}
  \variants{\freevarlabs \vartext{dibeya, dibaya}}}

\entry{dibag}{\headword{dibag}
  \pos{Copular verb}  \sense{
    \definition{plural present form of diban}    \example{
      \exsen{Ngämo kokok a dibag.}      \extran{Those are my grandchildren.}}}
  \variants{\SpellingVariant \vartext{debag}}}

\entry{dibagaeya}{\headword{dibagaeya}
  \pos{Copular verb}  \sense{
    \definition{past plural form of diban}    \example{
      \exsen{Llɨg a dibagaeya.}      \extran{Those were the boys.}}}}

\entry{dibagaeyo}{\headword{dibagaeyo}
  \pos{Copular verb}  \sense{
}
  \variants{\SpellingVariant \vartext{dibageyo}}}

\entry{dibageyo}{\headword{dibageyo}
  \variantof{SpellingVariant of \vartext{dibagaeyo}}}

\entry{dibagwaeya}{\headword{dibagwaeya}
  \pos{Copular verb}  \sense{
    \definition{past dual form of diban}    \example{
      \exsen{Ngämo baba bi dibagwaeya.}      \extran{Those were my parents.}}}}

\entry{dibamasäm}{\headword{dibamasäm}
  \pos{Adverbial demonstrative}  \sense{
    \definition{ablative form of diba; because}    \example{
      \exsen{Dibamasäma, dallän do Yop e.}      \extran{From there, she went to Yop.}}    \example{
      \exsen{Kollokolloe panypeny erallo, dibamasäma ngämi erag kollokolloe dag gänyme.}      \extran{They are speaking mixed (i.e. code-switching), because we (incl.) are mixed here.}}}
  \variants{\Unspecified Variant \vartext{dibamattäm}}}

\entry{dibamattäm}{\headword{dibamattäm}
  \variantof{Unspecified Variant of \vartext{dibamasäm}}}

\entry{dibame}{\headword{dibame}
  \pos{Nominal demonstrative}  \sense{
}}

\entry{diban}{\headword{diban}  \pronnote{diban}
  \pos{Copular verb}  \sense{
    \definition{copular form of diba (present singular form)}    \example{
      \exsen{Kak Wagiba diban.}      \extran{There's Grandmother Wagiba.}}}}

\entry{dibaolle}{\headword{dibaolle}
  \pos{Adverbial demonstrative}  \sense{
    \definition{allative form of diba}    \example{
      \exsen{Ngäna dibaolle melem e gozen.}      \extran{I entered that job.}}}}

\entry{dibawalle}{\headword{dibawalle}
  \pos{Adverbial demonstrative}  \sense{
}}

\entry{dibaya}{\headword{dibaya}
  \variantof{freevarlab of \vartext{dibaeya}}  \pronnote{dibaja}}

\entry{dibem}{\headword{dibem}
  \pos{Nominal demonstrative}  \sense{
    \definition{accusative form of diba}    \example{
      \exsen{Bogo dibem llo de kälnan eran.}      \extran{She climbs that tree.}}}}

\entry{dibeya}{\headword{dibeya}
  \variantof{freevarlab of \vartext{dibaeya}}  \pronnote{dibeja}}

\entry{dibie}{\headword{dibie}  \pronnote{dibie}
  \pos{Noun}  \sense{
    \definition{spectacled longbill}    \scientificname{Oedistoma iliolophus}}}

\entry{Dibllag}{\headword{Dibllag}
  \pos{Proper noun}  \sense{
    \definition{personal name}}}

\entry{Dibllagmälla}{\headword{Dibllagmälla}
  \pos{Proper noun}  \sense{
    \definition{female personal name}}}

\entry{Dibor}{\headword{Dibor}  \pronnote{dibor}
  \pos{Proper noun}  \sense{
    \definition{female personal name}}}

\entry{diboz}{\headword{diboz}
  \pos{Noun}  \sense{
    \definition{type of bird}}}

\entry{didiri}{\headword{didiri}
  \variantof{SpellingVariant of \vartext{didri}}}

\entry{didri}{\headword{didri}  \pronnote{didri}
  \pos{Adverbial demonstrative}  \sense{
    \definition{there}    \example{
      \exsen{Da ngäna sawe nyongo de bongkollmäll, ngaska didri lla da ddone dan.}      \extran{If I follow the left road, maybe there won't be any people there.}}}
  \variants{\SpellingVariant \vartext{didiri}}}

\entry{didribatt}{\headword{didribatt}
  \pos{Adverbial demonstrative}  \sense{
    \definition{ablative form of didri}    \example{
      \exsen{Ngämi didribatt dag Boze att dag.}      \extran{We (excl.) are from there, from Boze.}}}}

\entry{Didroe}{\headword{Didroe}
  \pos{Proper noun}  \sense{
    \definition{male personal name}}}

\entry{Dieb}{\headword{Dieb}  \pronnote{dieb}
  \pos{Proper noun}  \sense{
    \definition{male personal name}}}

\entry{Diendra}{\headword{Diendra}  \pronnote{diendra}
  \pos{Proper noun}  \sense{
    \definition{female personal name}}}

\entry{Digabo Källäm}{\headword{Digabo Källäm}  \note{Place names around Limol}  \pronnote{diɡabo kəɽəm}
  \pos{Proper noun}  \sense{
    \definition{Digabo Pond (canoe and fishing place in Taolang; road to Taolang is in AX95)}    \note{Place names around Limol}    \example{
}    \example{
}    \example{
}}}

\entry{digodigol}{\headword{digodigol}  \note{Trees}  \pronnote{diɡodiɡol}
  \pos{Noun}  \sense{
    \definition{type of tree}    \note{Trees}    \example{
}    \example{
}    \example{
}}}

\entry{digol}{\headword{digol}  \note{Trees}  \pronnote{diɡol}
  \pos{Noun}  \sense{
    \definition{type of big, tall tree that grows near gardens in the bush with white flowers, green fruits, and special, valuable red wood that is used for ax handles}    \note{Trees}    \example{
}    \example{
}    \example{
}}}

\entry{Dikae}{\headword{Dikae}
  \pos{Proper noun}  \sense{
    \definition{male personal name}}
  \variants{\SpellingVariant \vartext{Dikai}}}

\entry{Dikai}{\headword{Dikai}
  \variantof{SpellingVariant of \vartext{Dikae}}  \pronnote{dikai}}

\entry{Dikiboe}{\headword{Dikiboe}
  \pos{Proper noun}  \sense{
    \definition{personal name}}}

\entry{Dikullowang}{\headword{Dikullowang}  \note{Places around Limol}  \pronnote{dikuɽowaŋ}
  \pos{Proper noun}  \sense{
    \definition{Dikullowang (small island and hunting place in Taolang)}    \note{Places around Limol}    \example{
}    \example{
}    \example{
}}}

\entry{Dikullwang}{\headword{Dikullwang}
  \pos{Proper noun}  \sense{
}}

\entry{dikun}{\headword{dikun}
  \pos{Noun}  \sense{
    \definition{deacon}}}

\entry{dimes}{\headword{dimes}  \note{Trees}  \pronnote{dimes}
  \pos{Noun}  \sense{
    \definition{type of cultivated tree with sour, mango-like fruit}    \note{Trees}    \example{
}    \example{
}    \example{
}}}

\entry{Dimgi}{\headword{Dimgi}
  \pos{Proper noun}  \sense{
    \definition{Dimgi (toponym)}}}

\entry{Dimiri}{\headword{Dimiri}  \pronnote{dimiri}
  \pos{Proper noun}  \sense{
    \definition{Dimiri/Demeri (Idi-speaking village in Morehead Rural LLG; from Limol, one must pass through Kuiwang)}    \example{
}    \example{
}    \example{
}}}

\entry{Dimisisi}{\headword{Dimisisi}  \pronnote{dimisisi}
  \pos{Proper noun}  \sense{
    \definition{Dimisisi (Idi-speaking village in Morehead Rural LLG; from Limol, one must pass through Kinkin and Bok)}    \example{
}    \example{
}    \example{
}}
  \variants{\Unspecified Variant \vartext{Dimsisi, Dimsis}}}

\entry{Dimoe}{\headword{Dimoe}
  \pos{Proper noun}  \sense{
    \definition{female personal name}}
  \variants{\Unspecified Variant \vartext{Dimoi}}}

\entry{Dimoi}{\headword{Dimoi}
  \variantof{Unspecified Variant of \vartext{Dimoe}}  \pronnote{dimoi}}

\entry{Dimsis}{\headword{Dimsis}
  \variantof{Unspecified Variant of \vartext{Dimisisi}}}

\entry{Dimsisi}{\headword{Dimsisi}
  \variantof{Unspecified Variant of \vartext{Dimisisi}}}

\entry{Dimson}{\headword{Dimson}  \pronnote{dimson}
  \pos{Proper noun}  \sense{
    \definition{male personal name}}}

\entry{dindu}{\headword{dindu}  \note{Verbs of Motion}  \pronnote{'dindu}
  \pos{Intransitive verb}  \sense{\textbf{1.} 
    \definition{to run, flee, escape}    \note{Verbs of Motion}    \example{
      \exsen{Ngäna dindu allan.}      \extran{I am running.}}    \example{
      \exsen{Ddia da walle we nindugan.}      \extran{The deer ran to the water.}}    \example{
      \exsen{Bong angde bigma we dallän, sɨmell a dindugän.}      \extran{When Bong went to the pen, the pig had fled.}}    \example{
      \exsen{Obo pate bizmollne.}      \extran{We will be running to her.}}}  \sense{\textbf{2.} 
    \definition{race}    \note{Verbs of Motion}    \example{
      \exsen{Oba abo dindu da dädme llätt gogon.}      \extran{Then, their race stopped there.}}}}

\entry{dinggel}{\headword{dinggel}  \note{Animal names}  \pronnote{'dinggel}
  \pos{Noun}  \sense{
    \definition{sugar glider}    \scientificname{Petaurus breviceps}    \note{Animal names}    \example{
}    \example{
}    \example{
}}
  \variants{\Unspecified Variant \vartext{dinggoll}}}

\entry{dinggi}{\headword{dinggi}
  \pos{Noun}  \sense{
    \definition{dinghy}}}

\entry{dinggoll}{\headword{dinggoll}
  \variantof{Unspecified Variant of \vartext{dinggel}}}

\entry{dini}{\headword{dini}  \note{Trees}  \pronnote{dini}
  \pos{Noun}  \sense{
    \definition{type of small tree that grows in the bush with white flowers and red fruit that cassowaries like to eat}    \note{Trees}    \example{
}    \example{
}    \example{
}}}

\entry{dinidini}{\headword{dinidini}  \note{Trees}  \pronnote{dinidini}
  \pos{Noun}  \sense{
    \definition{type of small tree that grows in the bush with red fruit}    \note{Trees}    \example{
}    \example{
}    \example{
}}}

\entry{Dipa}{\headword{Dipa}  \pronnote{dipa}
  \pos{Proper noun}  \sense{
    \definition{male personal name}}}

\entry{dirindi}{\headword{dirindi}  \note{Food}  \pronnote{dirindi}
  \pos{Noun}  \sense{
    \definition{type of large yam with a white interior and thorns; with or without hairs}    \note{Food}    \example{
}    \example{
}    \example{
}}}

\entry{dirindi}{\headword{dirindi}  \pronnote{dirindi}
  \pos{Noun}  \sense{
    \definition{type of tree}}}

\entry{dirom}{\headword{dirom}  \lexeme{dirom}  \pronnote{di'rom}
  \pos{Noun}  \sense{
    \definition{southern cassowary}    \scientificname{Casuarius casuarius}    \example{
      \exsen{Dirom a ttongo mer mokowang ddäddäg dan.}      \extran{Cassowary is a very delicious animal.}}}}

\entry{dirom gäz}{\headword{dirom gäz}
  \pos{Transitive compound verb}  \sense{
    \definition{to menstruate for the first time}}}

\entry{dirom käp}{\headword{dirom käp}
  \pos{Noun}  \sense{
    \definition{type of small taro}}}

\entry{dirom käp}{\headword{dirom käp}  \pronnote{dirom kəp}
  \pos{Noun}  \sense{
    \definition{cassowary egg}}}

\entry{dirom mas}{\headword{dirom mas}
  \pos{Noun}  \sense{
    \definition{cassowary fibula}}}

\entry{diromdirom}{\headword{diromdirom}  \note{Trees}  \pronnote{diromdirom}
  \pos{Noun}  \sense{
    \definition{type of small tree with round, flat, edible fruit that are green and red when ripe}    \note{Trees}    \example{
}    \example{
}    \example{
}}}

\entry{dis}{\headword{dis}
  \variantof{Dialectal Variant of \vartext{English loanword}}}

\entry{district}{\headword{district}
  \variantof{SpellingVariant of \vartext{distrik}}}

\entry{distrik}{\headword{distrik}
  \pos{Noun}  \sense{
    \definition{district}}
  \variants{\SpellingVariant \vartext{district}}}

\entry{dit}{\headword{dit}  \lexeme{dit}  \pronnote{dit}
  \pos{Noun}  \sense{
    \definition{type of cane used for building houses, bows, and canoes}}}

\entry{Diwa}{\headword{Diwa}  \pronnote{diwa}
  \pos{Proper noun}  \sense{
    \definition{male personal name}}}

\entry{dɨbe}{\headword{dɨbe}
  \variantof{Unspecified Variant of \vartext{däbe}}  \pronnote{dɪbe}}

\entry{dɨbeya}{\headword{dɨbeya}
  \variantof{freevarlab of \vartext{däbaeya}}}

\entry{dɨdɨr}{\headword{dɨdɨr}
  \variantof{Unspecified Variant of \vartext{dädär}}  \note{Colors}  \pronnote{dɪdɪr}}

\entry{dɨl}{\headword{dɨl}  \pronnote{dɪl}
  \pos{Modifier}  \sense{
    \definition{bitter; sour}}
  \variants{\SpellingVariant \vartext{däl}}}

\entry{dɨp}{\headword{dɨp}  \pronnote{dɪp}
  \pos{Noun}  \sense{
    \definition{shoot (of a plant)}}}

\entry{do}{\headword{do}  \pronnote{do}
  \pos{Adverbial demonstrative}  \sense{\textbf{1.} 
    \definition{over there (distal)}    \example{
      \exsen{Ngämo za da dag do Pauma bo patme.}      \extran{My things are there at Pauma's.}}}  \sense{\textbf{2.} 
    \definition{to, until}    \example{
      \exsen{Ngämi däpllätt abo do Taolang gall tapma.}      \extran{Then we (excl.) started walking to Taolang canoe place.}}    \example{
      \exsen{Ako polle kättnan gongkaemallo do dättemänaeballo.}      \extran{They started to build the fences until they were finished.}}    \example{
      \exsen{iddob amnong alle do bädab}      \extran{from midnight to dawn}}}}

\entry{do}{\headword{do}  \pronnote{do}
  \pos{Noun}  \sense{\textbf{1.} 
    \definition{handle}    \example{
      \exsen{Dägmar a do peyang dan.}      \extran{The spoon has a handle.}}}  \sense{\textbf{2.} 
    \definition{femur}}}

\entry{Dobola}{\headword{Dobola}  \pronnote{dobola}
  \pos{Proper noun}  \sense{
    \definition{male personal name}}}

\entry{dodo}{\headword{dodo}
  \pos{Adverbial demonstrative}  \sense{
    \definition{there}    \example{
      \exsen{Ngämi dodo dagirni.}      \extran{We (excl.) were living there.}}}}

\entry{dodo}{\headword{dodo}  \pronnote{dodo}
  \pos{Transitive verb}  \sense{\textbf{1.} 
    \definition{to blow over}}  \sense{\textbf{2.} 
    \definition{to crumble sago into frying pan}}  \sense{\textbf{3.} 
    \definition{to remove thorn}}}

\entry{dodro}{\headword{dodro}  \pronnote{dodro}
  \pos{Intransitive verb}  \sense{\textbf{1.} 
    \definition{to slip}    \example{
      \exsen{Nikki bo kaptte da dronendronen dag.}      \extran{Nikki's clothes are slipping off.}}}  \sense{\textbf{2.} 
    \definition{to die}    \example{
      \exsen{Kollba da tämamae dadrowän a be kottllam aebe ttam dagirnän.}      \extran{The fish all died but the turtle lived on.}}}}

\entry{dodro}{\headword{dodro}  \pronnote{dodro}
  \pos{Transitive verb}  \sense{
    \definition{to clean}    \example{
      \exsen{Bogo llan de komlla deyadurowän mermerae.}      \extran{She washed her two ears well.}}}}

\entry{dogma}{\headword{dogma}  \note{Trees}  \pronnote{doɡma}
  \pos{Noun}  \sense{
    \definition{type of tree that grows in the bush with wood that smells like matches and is good for house posts}    \note{Trees}    \example{
}    \example{
}    \example{
}}}

\entry{Dokoe}{\headword{Dokoe}
  \pos{Proper noun}  \sense{
    \definition{male personal name}}}

\entry{Doli}{\headword{Doli}
  \pos{Proper noun}  \sense{
    \definition{male personal name}}}

\entry{dom}{\headword{dom}
  \pos{Transitive coverb}  \sense{
    \definition{to clench (one's fist)}    \example{
      \exsen{Dom nägagan ttang de.}      \extran{He clenched his fist.}}}}

\entry{domäll}{\headword{domäll}  \note{Type of pandanus}  \pronnote{doməɽ}
  \pos{Noun}  \sense{\textbf{1.} 
    \definition{type of pandanus}    \note{Type of pandanus}    \example{
}    \example{
}    \example{
}}  \sense{\textbf{2.} 
    \definition{old-style sewn mat made of domäll pandanus}    \note{Type of pandanus}}}

\entry{dombak}{\headword{dombak}
  \variantof{Unspecified Variant of \vartext{dompak}}}

\entry{domoe}{\headword{domoe}
  \variantof{Dialectal Variant of \vartext{dämoe}}  \note{Verbs of Motion}  \pronnote{domoe}}

\entry{dompa}{\headword{dompa}  \note{Traditional Arrows}  \pronnote{dompa}
  \pos{Noun}  \sense{\textbf{1.} 
    \definition{type of blunt arrow}    \note{Traditional Arrows}    \example{
      \exsen{Ge baba bäne digol dompa popoatt dan.}      \extran{Father made this dompa out of digol tree.}}    \example{
}    \example{
}}  \sense{\textbf{2.} 
    \definition{penis (slang)}    \note{Traditional Arrows}}}

\entry{dompadompa}{\headword{dompadompa}  \note{Traditional Arrows}  \pronnote{dompadompa}
  \pos{Noun}  \sense{
    \definition{type of spear}    \note{Traditional Arrows}    \example{
}    \example{
}    \example{
}}}

\entry{dompak}{\headword{dompak}  \lexeme{dompak}  \pronnote{dompak}
  \pos{Noun}  \sense{
    \definition{eel}}
  \variants{\Unspecified Variant \vartext{dombak}}}

\entry{Donae}{\headword{Donae}  \pronnote{donae}
  \pos{Proper noun}  \sense{
    \definition{female personal name}}
  \variants{\SpellingVariant \vartext{Donai}}}

\entry{Donai}{\headword{Donai}
  \variantof{SpellingVariant of \vartext{Donae}}  \pronnote{donai}}

\entry{dongkal}{\headword{dongkal}  \pronnote{doŋkal}
  \pos{Intransitive verb}  \sense{
    \definition{to stop, end}    \example{
      \exsen{Wällän a dädär me gondokalmällnegän.}      \extran{The roots ended on the rock.}}}}

\entry{dongkäral}{\headword{dongkäral}  \lexeme{dongkäral}  \pronnote{doŋkəral}
  \pos{Noun}  \sense{
    \definition{type of lizard}}}

\entry{dongki}{\headword{dongki}
  \pos{Noun}  \sense{
    \definition{donkey}}}

\entry{Donsi}{\headword{Donsi}
  \pos{Proper noun}  \sense{
    \definition{female personal name}}}

\entry{dor}{\headword{dor}
  \pos{Noun}  \sense{
    \definition{stalk}    \example{
      \exsen{Sana dor de adingoll däpittaemeya.}      \extran{We weaved the sago stalks like this.}}}}

\entry{Dore}{\headword{Dore}  \pronnote{dore}
  \pos{Proper noun}  \sense{
    \definition{female personal name}}}

\entry{Dorin}{\headword{Dorin}  \pronnote{dorin}
  \pos{Proper noun}  \sense{
    \definition{female personal name}}}

\entry{dorko}{\headword{dorko}  \pronnote{dorko}
  \pos{Modifier}  \sense{
    \definition{dry}    \example{
      \exsen{sana dorko}      \extran{dry sago}}    \example{
      \exsen{Ekaklle da dorko dan.}      \extran{The ground is dry.}}}}

\entry{dorllog}{\headword{dorllog}  \pronnote{dorɽoɡ}
  \pos{Noun}  \sense{
    \definition{rufous-bellied kookaburra}    \scientificname{Dacelo gaudichaud}}}

\entry{doros}{\headword{doros}  \note{Clothing}  \pronnote{doros}
  \pos{Noun}  \sense{
    \definition{pants}    \note{Clothing}    \example{
}    \example{
}    \example{
}}}

\entry{Doumori}{\headword{Doumori}  \pronnote{doumori}
  \pos{Proper noun}  \sense{
    \definition{Doumori (in Kiwai Rural LLG)}    \example{
}    \example{
}    \example{
}}}

\entry{dowa}{\headword{dowa}  \note{Trees}  \pronnote{dowa}
  \pos{Noun}  \sense{
    \definition{type of tree that grows in the grassland along creeks with wood used for firewood}    \note{Trees}    \example{
}    \example{
}    \example{
}}}

\entry{dowa~}{\headword{dowa~}
  \variantof{Derived Variant of \vartext{RED~}}}

\entry{Dowabunang}{\headword{Dowabunang}  \note{Places around Limol}  \pronnote{dowabunaŋ}
  \pos{Proper noun}  \sense{
    \definition{camping, sago, hunting place, and garden of Kaoga Dobola (on the road to Kinkin AZ94)}    \note{Places around Limol}    \example{
}    \example{
}    \example{
}}}

\entry{dowadowae}{\headword{dowadowae}
  \pos{Adverb}  \sense{
    \definition{close together, next to each other, neighboring, adjacent}    \example{
      \exsen{ekaklle me dämadäme dowadowae}      \extran{sitting on the ground next to each other}}}}

\entry{dowae}{\headword{dowae}  \pronnote{dowae}
  \pos{Locational}  \sense{
    \definition{vicinity, proximity}    \example{
      \exsen{Ngäna ibi allan wup dowae e.}      \extran{I am walking towards the banana.}}    \example{
      \exsen{Grace obo zegatt ebdo dowae me dan.}      \extran{Grace's birthday is soon.}}}}

\entry{Dowan}{\headword{Dowan}  \pronnote{dowan}
  \pos{Proper noun}  \sense{
    \definition{name of a mountain}}}

\entry{dowattäm}{\headword{dowattäm}
  \pos{Modifier}  \sense{
    \definition{ablative form of do}    \example{
      \exsen{Obo masar a dowattäm da Kulläntti.}      \extran{Her ancestor is from there, Kurunti.}}}}

\entry{doweae}{\headword{doweae}  \pronnote{doweae}
  \pos{Adverb}  \sense{
    \definition{straight}    \example{
      \exsen{Käsre ada gognegän doweae Bundae bälle yagyagag.}      \extran{Then they went like this straight to find Bundae.}}}}

\entry{dradre}{\headword{dradre}  \note{Trees}  \pronnote{dradre}
  \pos{Noun}  \sense{
    \definition{type of tree that grows in swamp with edible, currant-sized blue fruit}    \note{Trees}    \example{
}    \example{
}    \example{
}}}

\entry{dradre}{\headword{dradre}
  \pos{Transitive verb}  \sense{
    \definition{to dress}    \example{
      \exsen{Ngäna obom dädre.}      \extran{I dressed him.}}}}

\entry{droledrole}{\headword{droledrole}
  \variantof{freevarlab of \vartext{däroledärole}}  \pronnote{droledrole}}

\entry{druwem}{\headword{druwem}  \pronnote{druwem}
  \pos{Transitive coverb}  \sense{
    \definition{to knock on}    \example{
      \exsen{Ngäna ud de druwem nägawan.}      \extran{I knocked on the door.}}}}

\entry{du}{\headword{du}  \pronnote{du}
  \pos{Modifier}  \sense{
    \definition{wild (of plants)}}}

\entry{du~}{\headword{du~}
  \variantof{Derived Variant of \vartext{RED~}}}

\entry{du kyakya}{\headword{du kyakya}  \pronnote{du kjakja}
  \pos{Noun}  \sense{
    \definition{hook-billed kingfisher}    \scientificname{Melidora macrorrhina}}}

\entry{duab}{\headword{duab}  \pronnote{duab}
  \pos{Transitive verb}  \sense{
    \definition{to knock over; blow down}    \example{
      \exsen{Däräng a ine de naduaban.}      \extran{The dog knocked the water over.}}    \example{
      \exsen{Llo de duduaibnegnän a dattkaemnegän.}      \extran{Wind blew down the trees and broke them.}}}}

\entry{Duaba}{\headword{Duaba}  \pronnote{duaba}
  \pos{Proper noun}  \sense{
    \definition{Duaba (Gogodala-speaking village in Gogodala Rural LLG)}    \example{
}    \example{
}    \example{
}}}

\entry{duar}{\headword{duar}
  \pos{}}

\entry{dubllodubllom}{\headword{dubllodubllom}  \note{Trees}  \pronnote{dubɽodubɽom}
  \pos{Noun}  \sense{
    \definition{type of tree that grows in the bush with yellow fruit that have a hard seed, in which there is an edible nut}    \note{Trees}    \example{
}    \example{
}    \example{
}}}

\entry{Duboläpläp}{\headword{Duboläpläp}
  \variantof{Unspecified Variant of \vartext{Dubolläplläp}}}

\entry{Dubolläplläp}{\headword{Dubolläplläp}  \pronnote{duboɽəpɽəp}
  \pos{Proper noun}  \sense{
    \definition{Dubolläplläp (big hill and camping place on the road to Karama swamp; had houses on top in 2015)}}
  \variants{\Unspecified Variant \vartext{Dubollopllop, Duboläpläp}}}

\entry{Dubollopllop}{\headword{Dubollopllop}
  \variantof{Unspecified Variant of \vartext{Dubolläplläp}}}

\entry{dudli}{\headword{dudli}
  \variantof{freevarlab of \vartext{duduli}}}

\entry{duduli}{\headword{duduli}
  \pos{Adverb}  \sense{
    \definition{that way}    \example{
      \exsen{Llɨg a duduli ibi allan.}      \extran{The boy is going that way.}}}
  \variants{\Unspecified Variant \vartext{duliduli}}
  \variants{\freevarlab \vartext{dudli}}}

\entry{Dugal}{\headword{Dugal}  \pronnote{duɡal}
  \pos{Proper noun}  \sense{
    \definition{male personal name}}}

\entry{Dugi}{\headword{Dugi}
  \pos{Proper noun}  \sense{
    \definition{male personal name}}}

\entry{dugo}{\headword{dugo}  \note{Birds}  \pronnote{duɡo}
  \pos{Noun}  \sense{
    \definition{type of bird}    \note{Birds}    \example{
}    \example{
}    \example{
}}}

\entry{Duiya}{\headword{Duiya}  \pronnote{duija}
  \pos{Proper noun}  \sense{
    \definition{male personal name}}
  \variants{\SpellingVariant \vartext{Duya}}}

\entry{Duks}{\headword{Duks}
  \pos{Proper noun}  \sense{
    \definition{male personal name}}}

\entry{Dukumiang}{\headword{Dukumiang}  \note{Places around Limol}  \pronnote{dukumiaŋ}
  \pos{Proper noun}  \sense{
    \definition{Dukumiang (fishing, sago, hunting place, and garden; camping place of Bewag Bewag; R side of AY96, take southward road; northward road goes to Kapal)}    \note{Places around Limol}    \example{
}    \example{
}    \example{
}}}

\entry{duli}{\headword{duli}  \pronnote{duli}
  \pos{Adverbial demonstrative}  \sense{\textbf{1.} 
    \definition{over there (distal)}    \example{
      \exsen{Obo ma da duli dan kona me dan.}      \extran{His house is over there on the corner.}}    \example{
      \exsen{Duli dan ddob ttängäm me dan.}      \extran{She is over there, somewhere else.}}}  \sense{\textbf{2.} 
    \definition{away (from a place towards another direction)}    \example{
      \exsen{Ine da däbe llɨg kälsre de dätramän duli.}      \extran{The water carried the little boy away.}}    \example{
      \exsen{Duli abäll bibi.}      \extran{You all, go away.}}}}

\entry{duliballe}{\headword{duliballe}  \pronnote{dulibaɽe}
  \pos{Adverb}  \sense{
    \definition{ablative form of duli}    \example{
      \exsen{Ubi duliduli gobällän a ako duliballe ngäsmäll eran.}      \extran{They had migrated far away and are now returning from there.}}}
  \variants{\Dialectal Variant \vartext{dulibane}}}

\entry{dulibane}{\headword{dulibane}
  \variantof{Dialectal Variant of \vartext{duliballe}}}

\entry{dulibatta}{\headword{dulibatta}
  \variantof{Unspecified Variant of \vartext{dulibattäm}}}

\entry{dulibattäm}{\headword{dulibattäm}
  \pos{Adverbial demonstrative}  \sense{
    \definition{ablative form of duli}    \example{
      \exsen{Obo masar a dulibatta deyarän gänyaolle.}      \extran{His grandfather came here from over there.}}}
  \variants{\Unspecified Variant \vartext{dulibatta}}}

\entry{duliduli}{\headword{duliduli}
  \variantof{Unspecified Variant of \vartext{duduli}}}

\entry{dum}{\headword{dum}  \lexeme{dum}  \pronnote{dum}
  \pos{Noun}  \sense{
    \definition{width (of a house)}    \example{
      \exsen{dum papek}      \extran{wall spanning the width of a building}}}}

\entry{dum}{\headword{dum}  \lexeme{dum}  \pronnote{dum}
  \pos{Noun}  \sense{
    \definition{type of big tree that grows in bush with strong wood used for making canoes and paddles, sap used to paint bowstrings or to patch holes, yellow flowers, and green fruit; planted near house for shade}    \scientificname{Pterocarpus indicus}}}

\entry{Dum}{\headword{Dum}
  \pos{Proper noun}  \sense{
    \definition{Dum (big hill on the road to Malam; Bamboo and Zarma creek)}}}

\entry{dum}{\headword{dum}  \lexeme{dum}  \pronnote{dum}
  \pos{Noun}  \sense{
    \definition{placenta}}}

\entry{Dum Tutu}{\headword{Dum Tutu}  \note{Places around Limol}  \pronnote{dum tutu}
  \pos{Proper noun}  \sense{
    \definition{Dum Mountain}    \note{Places around Limol}    \example{
}    \example{
}    \example{
}}}

\entry{dumbeg}{\headword{dumbeg}
  \pos{Transitive verb}  \sense{
}}

\entry{dumbi}{\headword{dumbi}  \note{Medicine}  \pronnote{dumbi}
  \pos{Noun}  \sense{
    \definition{type of red tree}    \note{Medicine}    \example{
}    \example{
}    \example{
}}}

\entry{dumdum}{\headword{dumdum}  \pronnote{dumdum}
  \pos{Transitive verb}  \sense{
    \definition{to surround}    \example{
      \exsen{Lla gul ulle da obom dandumeyo.}      \extran{The large crowd surrounded him.}}}}

\entry{Dumoll}{\headword{Dumoll}  \note{Places around Limol}  \pronnote{dumoɽ}
  \pos{Proper noun}  \sense{
    \definition{Dumoll (Taolang side garden and camping place of Gidu Jerry, Kols Baewa, and Wareya Giniya)}    \note{Places around Limol}    \example{
}    \example{
}    \example{
}}}

\entry{dun}{\headword{dun}  \pronnote{dun}
  \pos{Ideophone}  \sense{
    \definition{sound of a drum}}}

\entry{dundu kllamen}{\headword{dundu kllamen}  \note{Games}  \pronnote{dundu kɽamen}
  \pos{Noun}  \sense{
    \definition{type of game involving a race}    \note{Games}    \example{
}    \example{
}    \example{
}}}

\entry{dune~}{\headword{dune~}
  \variantof{Derived Variant of \vartext{RED~}}}

\entry{dune}{\headword{dune}
  \pos{Verb}  \sense{
}}

\entry{duny}{\headword{duny}  \lexeme{duny}  \pronnote{duɲ}
  \pos{Noun}  \sense{
    \definition{beetle}}}

\entry{dupi}{\headword{dupi}  \note{Parts of the body}  \pronnote{dupi}
  \pos{Noun}  \sense{
    \definition{stomach}    \note{Parts of the body}    \example{
}    \example{
}    \example{
}}}

\entry{dur}{\headword{dur}  \lexeme{dur}  \pronnote{dur}
  \pos{Noun}  \sense{
    \definition{type of medium-sized bamboo that grows along creeks; used for dancing; young plants used for cooking sago}}}

\entry{durgu}{\headword{durgu}
  \pos{Noun}  \sense{
    \definition{cliff}}}

\entry{Duwaba}{\headword{Duwaba}  \pronnote{duwaba}
  \pos{Proper noun}  \sense{
    \definition{Duaba (in Gogodala Rural LLG)}    \example{
}    \example{
}    \example{
}}}

\entry{duwar}{\headword{duwar}
  \pos{Modifier}  \sense{
}}

\entry{duwel}{\headword{duwel}  \lexeme{duwel}  \pronnote{duwel}
  \pos{Noun}  \sense{
    \definition{type of tall tree that grows in the bush near yam gardens}}}

\entry{duwel sära}{\headword{duwel sära}  \pronnote{duwel səra}
  \pos{Noun}  \sense{\textbf{1.} 
    \definition{type of sago bundle wrapped in sago leaves}    \example{
}}  \sense{\textbf{2.} 
    \definition{the third stage of sago growth in which the leaves are shorter and the pith is almost ready to be harvested}    \example{
      \exsen{Sana da duwelsärawang gogon.}}}}

\entry{duwem}{\headword{duwem}  \pronnote{duwem}
  \pos{Transitive/Intransitive coverb}  \sense{\textbf{1.} 
    \definition{to eat}    \example{
      \exsen{Ubi tämamae duwem gognegän.}      \extran{They all ate.}}    \example{
      \exsen{Duwem bägaeya.}      \extran{We ate it.}}}  \sense{\textbf{2.} 
    \definition{food, meal}    \example{
      \exsen{Ada gogon, "Ddone, ge ade ngämo da yäbdo duwem e a toto duwem e."}      \extran{‎He said, "No, this is for my lunch and for my dinner."}}    \example{
      \exsen{yäbdo duwem}      \extran{lunch}}    \example{
      \exsen{toto duwem}      \extran{dinner}}}}

\entry{duwemduwem}{\headword{duwemduwem}  \pronnote{duwemduwem}
  \pos{Noun}  \sense{
    \definition{feast, fellowship meal}    \example{
      \exsen{Bangeseya ibi duwemduwem de.}      \extran{We (incl.) will make a feast.}}}}

\entry{duwie ku}{\headword{duwie ku}  \lexeme{duwie ku}  \pronnote{duwie ku}
  \pos{Noun}  \sense{
    \definition{type of big purple yam}}}

\entry{Duya}{\headword{Duya}
  \variantof{SpellingVariant of \vartext{Duiya}}}

\chapter{E}


\entry{e}{\headword{e}  \pronnote{e}
  \pos{Interjection}  \sense{\textbf{1.} 
    \definition{whoa}}  \sense{\textbf{2.} 
    \definition{let's go}}}

\entry{e}{\headword{e}
  \pos{Relative pronoun}  \sense{\textbf{1.} 
    \definition{which, what, that}    \example{
      \exsen{Iba kame da e kollba we gopänaeyän.}      \extran{We (incl.) don't know which fish she turned into.}}}  \sense{\textbf{2.} 
    \definition{which, what}    \example{
      \exsen{E pazi me ka?}      \extran{In which year?}}    \example{
      \exsen{Bongo e bin dan ma me?}      \extran{What is his position in the community?}}}}

\entry{e}{\headword{e}
  \pos{Coordinating connective}  \sense{
}}

\entry{=e}{\headword{=e}
  \variantof{freevarlab of \vartext{=aeya}}}

\entry{=e}{\headword{=e}  \pronnote{e}
  \pos{Discourse clitic}  \sense{
    \definition{vocative}}}

\entry{=e}{\headword{=e}  \pronnote{e}
  \pos{Nominal enclitic}  \sense{\textbf{1.} 
    \definition{allative case clitic; to, into, towards}    \example{
      \exsen{Mälla da maket e nallan.}      \extran{The woman goes to the market.}}    \example{
      \exsen{Ngämo näkäp e gongttagän ada ngaska zime iddob amänong agan.}      \extran{It came to my mind that maybe it was already midnight.}}    \example{
      \exsen{Ngäna mɨnyi bäne eka de bäpänaeneg Ingglis e.}      \extran{I'm going to translate your words into English.}}}  \sense{\textbf{2.} 
    \definition{for, to (purposive use of the allative case; also follows nonfinite verbs to express desire, intention, or the start of an action)}    \example{
      \exsen{Ge mätta da känaebag e dan.}      \extran{This yam is for tomorrow.}}    \example{
      \exsen{Obo moko da gäbän e dan.}      \extran{He wants to jump (lit. his desire is to jump).}}    \example{
      \exsen{Alla täräp me ibi ge nge de ibeny e dan?}      \extran{At what time are we going to plant this coconut?}}    \example{
      \exsen{Gäddgädd e nängkaman.}      \extran{He started to fight.}}}}

\entry{ebagal}{\headword{ebagal}  \note{Traditional Arrows}  \pronnote{ebaɡal}
  \pos{Noun}  \sense{
    \definition{type of spear}    \note{Traditional Arrows}    \example{
}    \example{
}    \example{
}}}

\entry{ebdo}{\headword{ebdo}  \pronnote{ebdo}
  \pos{Noun}  \sense{\textbf{1.} 
    \definition{day}    \example{
      \exsen{ebdo ulle}      \extran{all day}}}  \sense{\textbf{2.} 
    \definition{noon (approx. 11 AM –1 PM)}}
  \variants{\SpellingVariant \vartext{yebdo}}
  \variants{\freevarlab \vartext{yäbdo}}}

\entry{ebdo duwem}{\headword{ebdo duwem}
  \pos{Noun}  \sense{
    \definition{lunch}}}

\entry{Eda}{\headword{Eda}
  \variantof{SpellingVariant of \vartext{Edna}}}

\entry{eddom}{\headword{eddom}  \pronnote{eɖ͡ʐom}
  \pos{Temporal adverb}  \sense{
    \definition{the day before yesterday}}}

\entry{eddrugoe}{\headword{eddrugoe}
  \pos{Transitive verb}  \sense{
    \definition{to shout}}}

\entry{ede}{\headword{ede}  \pronnote{ede}
  \pos{Subordinating connective}  \sense{\textbf{1.} 
    \definition{so, therefore}    \example{
      \exsen{Be ngämo ttɨle da attkaman, ede ngämlle ddone mullae dan, ngäna tongoe e beyar.}      \extran{But my leg broke, so I can't go play.}}}  \sense{\textbf{2.} 
    \definition{then}    \example{
      \exsen{Da obo ngoe a mamam gognegän, ede Yokon a bäne mätta de notan.}      \extran{If her teeth are red, then Yokon has eaten your yam.}}}}

\entry{Edeb}{\headword{Edeb}  \note{Places around Limol}  \pronnote{edeb}
  \pos{Proper noun}  \sense{\textbf{1.} 
    \definition{garden, camping, and hunting place of Tewa (on the other side of Karama swamp)}    \note{Places around Limol}    \example{
}    \example{
}    \example{
}}  \sense{\textbf{2.} 
    \definition{personal name}    \note{Places around Limol}}}

\entry{edeya}{\headword{edeya}
  \pos{Connective}  \sense{
    \definition{but}}}

\entry{Edna}{\headword{Edna}  \pronnote{edna}
  \pos{Proper noun}  \sense{
    \definition{female personal name}}
  \variants{\SpellingVariant \vartext{Eda}}}

\entry{Edward}{\headword{Edward}  \pronnote{edward}
  \pos{Proper noun}  \sense{
    \definition{male personal name}}}

\entry{Egapo}{\headword{Egapo}  \note{Places around Limol}  \pronnote{eɡapo}
  \pos{Proper noun}  \sense{
    \definition{camping place and large community garden (in Limol; road to garden is visible on L of AZ96)}    \note{Places around Limol}    \example{
}    \example{
}    \example{
}}}

\entry{ege}{\headword{ege}  \pronnote{eɡe}
  \pos{}  \sense{
    \definition{now}}}

\entry{ei}{\headword{ei}
  \pos{Interjection}  \sense{
    \definition{hey}}}

\entry{eighty}{\headword{eighty}
  \variantof{SpellingVariant of \vartext{eiti}}}

\entry{eit}{\headword{eit}  \note{Numbers}  \pronnote{eit}
  \pos{Numeral}  \sense{
    \definition{eight (English numeral; also general)}    \note{Numbers}    \example{
}    \example{
}    \example{
}}}

\entry{eiti}{\headword{eiti}
  \pos{Numeral}  \sense{
    \definition{eighty}}
  \variants{\SpellingVariant \vartext{eighty}}}

\entry{eitin}{\headword{eitin}
  \pos{Numeral}  \sense{
    \definition{eighteen (English numeral)}}}

\entry{eiz}{\headword{eiz}  \pronnote{eiz}
  \pos{Noun}  \sense{
    \definition{HIV; AIDS}}}

\entry{eka~}{\headword{eka~}
  \variantof{Infinitival reduplicant of \vartext{RED~}}}

\entry{eka}{\headword{eka}  \pronnote{eka}
  \pos{Noun}  \sense{\textbf{1.} 
    \definition{language}    \example{
      \exsen{Motu eka, pällämpälläm eka}      \extran{Hiri Motu, English}}    \example{
      \exsen{Ngämo Ende eka da umllang dan, be bäne ddone mäzi umllang dan.}      \extran{My Ende is fluent, but yours isn't.}}}  \sense{\textbf{2.} 
    \definition{word, message, news}    \example{
      \exsen{Ngäna mɨnyi bäne eka de bäpänaeneg Ingglis e.}      \extran{I'm going to translate your words into English.}}    \example{
      \exsen{Komlla o kumuddäga lla da mɨnyi utt alle eka de bängaeyo.}      \extran{Two or three men will send word with the conch shell.}}}  \sense{\textbf{3.} 
    \definition{story}    \example{
      \exsen{Ge eka da gongesän ngattong.}      \extran{This story happened first.}}    \example{
      \exsen{Ngämo eka kälsre da däbe.}      \extran{That was my little story.}}}  \sense{\textbf{4.} 
    \definition{sound, song, call}    \example{
      \exsen{Gontomoneyo wutt eka we.}      \extran{They were waiting for the sound of the conch shell.}}    \example{
      \exsen{Alläp de eka eran.}      \extran{He plays the drum.}}}  \sense{\textbf{5.} 
    \definition{to say, speak, talk, tell}    \example{
      \exsen{Ngäna mɨnyi ada eka bog.}      \extran{I will say so.}}    \example{
      \exsen{Män a llɨg de ada eka dägagän, "Ma we nalle."}      \extran{The girl told the boy, "Go home."}}}}

\entry{eka kutt}{\headword{eka kutt}  \pronnote{eka kuʈ͡ʂ}
  \pos{Noun}  \sense{
    \definition{word (single unit)}    \example{
      \exsen{Gudae masamasar aba eka kutt panypeny da kälakälae sapasapang gognegän.}      \extran{The words of our ancestors were slightly different.}}}}

\entry{eka laem}{\headword{eka laem}
  \pos{Intransitive compound verb}  \sense{
    \definition{to argue}    \example{
      \exsen{Eka golaemeyo.}      \extran{They argued.}}}}

\entry{eka llɨtɨtang}{\headword{eka llɨtɨtang}
  \pos{Noun}  \sense{
    \definition{messenger, prophet}    \example{
      \exsen{Bogo Adi bäne ttongo llɨtɨtang dan.}      \extran{He is one of God's messengers. / He is a prophet.}}}}

\entry{eka malam}{\headword{eka malam}  \pronnote{eka malam}
  \pos{Transitive compound verb}  \sense{
    \definition{obey}    \example{
}    \example{
}    \example{
}}}

\entry{eka midd}{\headword{eka midd}
  \pos{Noun}  \sense{
    \definition{meaning}    \example{
      \exsen{Oba pate bogo tämamae eka midd de däddelnegnän}      \extran{He was explaining all the meanings to them.}}}}

\entry{eka mu}{\headword{eka mu}  \pronnote{eka mu}
  \pos{Noun}  \sense{
    \definition{answer, reply}    \example{
      \exsen{Ngäna ubim eka mu anggan Mägayam eka walle.}      \extran{I answer them in Makayam.}}    \example{
      \exsen{Kottllam a ada eka mu gogän, "Kandärmang."}      \extran{Turtle replied, "Sorry."}}}}

\entry{eka pänaenenang}{\headword{eka pänaenenang}  \pronnote{eka pənaenenaŋ}
  \pos{Noun}  \sense{
    \definition{linguist}}}

\entry{eka panypeny}{\headword{eka panypeny}  \pronnote{eka paɲpeɲ}
  \pos{Transitive compound verb}  \sense{
    \definition{to speak}    \example{
      \exsen{Ddone ngänäm mulldae allan eka panypeny e.}      \extran{I am unable to speak.}}    \example{
      \exsen{Ge Llimollang a alla ingollang eka däpanyneyo.}      \extran{This is how Limol people spoke it.}}}}

\entry{eka panypenyang}{\headword{eka panypenyang}  \note{Government}  \pronnote{eka paɲpeɲaŋ}
  \pos{Noun}  \sense{\textbf{1.} 
    \definition{speaker}    \note{Government}    \example{
      \exsen{Ge Ende eka panypenyang lla dan.}      \extran{This is an Ende speaker.}}}  \sense{\textbf{2.} 
    \definition{spokesperson}    \note{Government}    \example{
      \exsen{Ngäna eka panypenyang dan kominiti mi.}      \extran{I am a spokesperson in the community.}}    \example{
}    \example{
}}}

\entry{eka tameny}{\headword{eka tameny}  \note{Polite words and conversation}  \pronnote{eka tameɲ}
  \pos{Intransitive compound verb}  \sense{\textbf{1.} 
    \definition{to discuss, converse}    \note{Polite words and conversation}    \example{
      \exsen{Ngämo nagnag aba kämall eka gontemenyne.}      \extran{I was conversing with my friends.}}    \example{
}    \example{
}}  \sense{\textbf{2.} 
    \definition{to tell; converse with, speak to}    \note{Polite words and conversation}    \example{
      \exsen{Ubi ngämim gänya eka walle mae eka deyantemenynegnän.}      \extran{They were only speaking to us in this language.}}}}

\entry{eka täp}{\headword{eka täp}  \pronnote{eka təp}
  \pos{Noun}  \sense{
    \definition{disagreement, debate, argument}    \example{
      \exsen{Angde Yesu a sedusi lla da eka täp me gognegnän, ttongo zuu llayaba sabi tamenyang a dallän oba pate.}      \extran{When Jesus and the Sadducee were debating, a Pharisee (lit. teacher of the Jewish law) came towards them.}}}}

\entry{eka ttängkɨg}{\headword{eka ttängkɨg}  \pronnote{eka ʈ͡ʂəŋkɪɡ}
  \pos{Transitive compound verb}  \sense{
    \definition{to answer someone with force, to talk against someone}}}

\entry{ekaeka}{\headword{ekaeka}
  \pos{Intransitive coverb}  \sense{
    \definition{to make noise}    \example{
      \exsen{Borale da eran ttongo ekaeka ma za dan.}      \extran{The flute is a thing that makes noise.}}}}

\entry{ekaklle}{\headword{ekaklle}  \pronnote{ekakɽe}
  \pos{Noun}  \sense{\textbf{1.} 
    \definition{land}    \example{
      \exsen{Ddia da ge ulle abal ddäddäg dan iba ekaklle me.}      \extran{Deer is a major animal in our (excl.) land.}}}  \sense{\textbf{2.} 
    \definition{ground}    \example{
      \exsen{Ekaklle da llokttang dan.}      \extran{The ground is hard.}}}  \sense{\textbf{3.} 
    \definition{Earth}    \example{
      \exsen{Kok ekaklle de nganaenen eran.}      \extran{The moon rotates around Earth.}}}  \sense{\textbf{4.} 
    \definition{low}    \example{
      \exsen{De nyonga da ekaklle abal me paplläg eran.}      \extran{That cockatoo is flying very low.}}}}

\entry{ekaklle ulle}{\headword{ekaklle ulle}
  \pos{Noun}  \sense{
    \definition{world}    \example{
      \exsen{Mälla papa da buddo ulle dan Papua Niugini mi a be kemibi ngättma me ge ekaklle ulle me ako.}      \extran{Beating one's wife is a big issue in Papua New Guinea but also in many places in the world.}}}}

\entry{ekakukiang}{\headword{ekakukiang}  \note{Bad characteristics of people}  \pronnote{ekakukiaŋ}
  \pos{Noun}  \sense{
    \definition{liar}    \note{Bad characteristics of people}    \example{
}    \example{
}    \example{
}}}

\entry{ekamatär}{\headword{ekamatär}  \note{Parts of the Body}  \pronnote{ekamatər}
  \pos{Noun}  \sense{
    \definition{throat}    \note{Parts of the Body}    \example{
}    \example{
}    \example{
}}}

\entry{ekateptep}{\headword{ekateptep}  \note{Bad characteristics of people}  \pronnote{ekateptep}
  \pos{Noun}  \sense{
    \definition{liar, conman}    \note{Bad characteristics of people}    \example{
}    \example{
}    \example{
}}}

\entry{eke}{\headword{eke}
  \variantof{Unspecified Variant of \vartext{ke}}}

\entry{eleben}{\headword{eleben}
  \pos{Numeral}  \sense{
    \definition{eleven (English numeral)}}
  \variants{\SpellingVariant \vartext{ileben}}}

\entry{elementary}{\headword{elementary}
  \variantof{Dialectal Variant of \vartext{English loanword}}}

\entry{elementeri}{\headword{elementeri}
  \pos{}}

\entry{elementeri skul}{\headword{elementeri skul}
  \variantof{SpellingVariant of \vartext{elementri skul}}}

\entry{elementri skul}{\headword{elementri skul}  \note{School}  \pronnote{elementri skul}
  \pos{Noun}  \sense{
    \definition{elementary school}    \note{School}    \example{
}    \example{
}    \example{
}}
  \variants{\SpellingVariant \vartext{elementeri skul}}}

\entry{Elisa}{\headword{Elisa}  \pronnote{ɛlisa}
  \pos{Proper noun}  \sense{
    \definition{female personal name}}}

\entry{Elizabeth}{\headword{Elizabeth}
  \pos{Proper noun}  \sense{
    \definition{female personal name}}}

\entry{ellollo}{\headword{ellollo}  \pronnote{eɽoɽo}
  \pos{Intransitive verb}  \sense{
    \definition{to put down load and rest}}}

\entry{Elsie}{\headword{Elsie}  \pronnote{ɛlsi}
  \pos{Proper noun}  \sense{
    \definition{female personal name}}}

\entry{Em}{\headword{Em}  \pronnote{em}
  \pos{Noun}  \sense{
    \definition{Em language (Pahoturi River language spoken in Kurunti, Kibuli, Beyambod)}}}

\entry{em}{\headword{em}  \pronnote{em}
  \pos{Interrogative pronoun}  \sense{
    \definition{accusative form of e}    \example{
      \exsen{Bongo em skul de dättemän?}      \extran{Which level of school did you finish?}}}}

\entry{emaemae}{\headword{emaemae}  \pronnote{emaemae}
  \pos{Modifier}  \sense{
    \definition{various, many different kinds}    \example{
      \exsen{Dageya, emaemae pa de dägäddaebeyo.}      \extran{They went and killed many different types of birds.}}}}

\entry{ematta}{\headword{ematta}  \note{Question words}  \pronnote{emaʈ͡ʂa}
  \pos{Adverb}  \sense{\textbf{1.} 
    \definition{ablative form of e; why (relative)}    \note{Question words}    \example{
      \exsen{Ngämo kame dan ematta bongo mikutt alle.}      \extran{I don't know why you're mad.}}    \example{
      \exsen{Bogo gongosalle ag me ma we a llabunda bim, umllang dägnegalle ewatta kuddäll gogon lla da.}      \extran{He would return to the house of the relatives in the morning and inform them why the person died.}}}  \sense{\textbf{2.} 
    \definition{ablative form of e; why (interrogative)}    \note{Question words}    \example{
      \exsen{Bongo ematta däräng de nazuwalle?}      \extran{Why did you shoot the dog?}}    \example{
      \exsen{Bibi komlla ewatta käptang abal agalla?}      \extran{Why are you two so wet?}}}
  \variants{\Unspecified Variant \vartext{ewatta}}}

\entry{Emi}{\headword{Emi}  \pronnote{emi}
  \pos{Proper noun}  \sense{
    \definition{female personal name}}}

\entry{enaenae}{\headword{enaenae}
  \variantof{Unspecified Variant of \vartext{enanae}}}

\entry{enanae}{\headword{enanae}  \pronnote{enanae}
  \pos{Adverb}  \sense{\textbf{1.} 
    \definition{directly, straight}    \example{
      \exsen{Llɨg kälsre de ine da enanae dae dätramän do ddob llayaba pate dängttägän.}      \extran{The water carried the small boy straight through and he arrived near some people.}}}  \sense{\textbf{2.} 
    \definition{forever, for good}    \example{
      \exsen{Obo ttängäm de enanae dowansegän.}      \extran{He left his village forever.}}}  \sense{\textbf{3.} 
    \definition{final; finally, at last}    \example{
      \exsen{God bo enanane ttoen pallängkmeny}      \extran{God's final judgment}}    \example{
      \exsen{Enanae abal ako ngäna gongkäbämenyne.}      \extran{Then, at long last, I was diving.}}}
  \variants{\Unspecified Variant \vartext{enaenae}}}

\entry{enanaeenanae}{\headword{enanaeenanae}  \pronnote{enanaeenanae}
  \pos{Modifier}  \sense{
    \definition{eternal}    \example{
      \exsen{Era lla da bongnonggän ddone mɨnyi kuddäll bogän, be enanaeenanae ttam de bälldän.}      \extran{He who believes will not die, but will receive eternal life.}}}}

\entry{enda}{\headword{enda}  \note{Question words}  \pronnote{enda}
  \pos{Interrogative pronoun}  \sense{
    \definition{what}    \note{Question words}    \example{
      \exsen{Ge enda?}      \extran{What is this?}}    \example{
}    \example{
}}}

\entry{endaeya}{\headword{endaeya}
  \pos{Copular verb}  \sense{
    \definition{past singular form of endan}    \example{
      \exsen{Diba eka bo bin a endaeya?}      \extran{What was the name of that language?}}}
  \variants{\freevarlabs \vartext{endeya, endaya}}}

\entry{endaeyag}{\headword{endaeyag}  \note{Question words}  \pronnote{endaejaɡ}
  \pos{Copular verb}  \sense{
    \definition{past plural form of endan}    \note{Question words}    \example{
}    \example{
}    \example{
}}}

\entry{endag}{\headword{endag}
  \pos{Copular verb}  \sense{
    \definition{present plural form of endan}    \example{
      \exsen{Bäne eka muu da endag ubira?}      \extran{What answers do you have for them?}}}}

\entry{endagaeya}{\headword{endagaeya}  \pronnote{endaɡaeja}
  \pos{Copular verb}  \sense{
    \definition{past plural form of endan}    \example{
      \exsen{Bogo tämamae wäte endagaeya mermerae dänggllämnegän.}      \extran{Whatever wounds there were, she cleaned them all properly.}}}
  \variants{\freevarlab \vartext{endageya}}}

\entry{endageya}{\headword{endageya}
  \variantof{freevarlab of \vartext{endagaeya}}}

\entry{endan}{\headword{endan}
  \pos{Copular verb}  \sense{
    \definition{copular form of enda (present singular form)}    \example{
      \exsen{Bäne mokowang melem a endan?}      \extran{What is your favorite work?}}}}

\entry{endaya}{\headword{endaya}
  \variantof{freevarlab of \vartext{endaeya}}  \note{Question words}  \pronnote{endaja}}

\entry{enddäna}{\headword{enddäna}  \pronnote{enɖ͡ʐəna}
  \pos{Noun}  \sense{
    \definition{clearing}    \example{
      \exsen{Tatuma da enddäna me dan.}      \extran{The washing place is in the clearing.}}}
  \variants{\Unspecified Variant \vartext{yänddäna}}
  \variants{\Fast Speech Variant \vartext{enddna}}}

\entry{enddänaenddäna}{\headword{enddänaenddäna}
  \pos{Adverb}  \sense{
    \definition{in the open, openly, freely}    \example{
      \exsen{Bogo llameny känyärtto ttängäm me dagirnän, ddone enddnaenddna taon e dallnän.}      \extran{He was staying in quiet places without people; he wasn't going openly into town.}}}}

\entry{enddna}{\headword{enddna}
  \variantof{Fast Speech Variant of \vartext{enddäna}}}

\entry{Ende}{\headword{Ende}  \pronnote{ende}
  \pos{Proper noun}  \sense{
    \definition{Ende language (Pahoturi River language spoken in Limol, Malam, and Kinkin)}    \example{
}    \example{
}    \example{
}}}

\entry{ende}{\headword{ende}  \pronnote{ende}
  \pos{Interrogative pronoun}  \sense{
    \definition{accusative form of enda}    \example{
      \exsen{Sisri iddob e, bongo ende ngasnges e dan?}      \extran{What are you doing tonight?}}}}

\entry{endeya}{\headword{endeya}
  \variantof{freevarlab of \vartext{endaeya}}}

\entry{Endo}{\headword{Endo}
  \pos{Proper noun}  \sense{
    \definition{female personal name}}}

\entry{Englis}{\headword{Englis}
  \variantof{SpellingVariant of \vartext{Ingglis}}}

\entry{English}{\headword{English}
  \variantof{SpellingVariant of \vartext{Ingglis}}}

\entry{English loanword}{\headword{English loanword}
  \pos{Pro-form}  \sense{
}
  \variants{\Dialectal Variant \vartext{isipi, year, very, 1940, ospitol, officer, no, or, teacher, feel, Papua, from, vocational, dis, fortiz, language, resourceful, adult, second, rizen, student, transfer, Guinea, elementary, years, it, chairman, namba, last, old, andred, sosa, training, now, there, got, until, uniform, of, toilet, New, so, group, first, shake, nothing, tunnel, still, secretary, this, government, in, central, are, about, land, many, fellowship, family, grade, independence, fellows, botol, forti}}
  \variants{\freevarlabs \vartext{you, blame, when, food, open, must, linguistic, me, nett, ran, daes, Teachers, example, go, kapi, the, area, greasy, but, and, not, 1940s, church, stael}}}

\entry{enma}{\headword{enma}  \pronnote{enma}
  \pos{Noun}  \sense{
    \definition{nature}}}

\entry{Enza}{\headword{Enza}  \note{Place names}  \pronnote{enza}
  \pos{Proper noun}  \sense{
    \definition{Enza (Bitur-speaking village in Oriomo-Bituri Rural LLG; from Limol, one must pass through Bisuaka)}    \note{Place names}    \example{
}    \example{
}    \example{
}}}

\entry{enzul}{\headword{enzul}
  \pos{Noun}  \sense{
    \definition{angel}}}

\entry{era~}{\headword{era~}
  \variantof{freevarlab of \vartext{RED~}}}

\entry{era}{\headword{era}  \pronnote{era}
  \pos{Relative pronoun}  \sense{\textbf{1.} 
    \definition{which, that, who}    \example{
      \exsen{Bam era mälla da papa nagan ngämo nag dan.}      \extran{The woman that hit you is my friend.}}    \example{
      \exsen{Däräng a bam era da naddägan, ngämo dan.}      \extran{The dog that bit you is mine.}}    \example{
      \exsen{Era lla da bongnonggän ddone mɨnyi kuddäll bogän.}      \extran{The man who believes will not die.}}}  \sense{\textbf{2.} 
    \definition{which, what}    \example{
      \exsen{Ibi ngasekäma era nyongo ibi wi dag?}      \extran{Which road might we (incl.) take?}}    \example{
      \exsen{Bäne era eka dag umllang a?}      \extran{What languages do you speak?}}}  \sense{\textbf{3.} 
    \definition{focus marker}    \example{
      \exsen{Bongo era Adi bo llɨg dan.}      \extran{YOU are the son of God.}}}}

\entry{Erabal}{\headword{Erabal}
  \pos{Proper noun}  \sense{
    \definition{female personal name}}}

\entry{eraballe}{\headword{eraballe}
  \pos{Relative pronoun}  \sense{
    \definition{instrumental form of era}    \example{
      \exsen{nyäng iatt mälla da eraballe za kemibi bokomän}      \extran{bag with which women carry many things}}}}

\entry{eraeya}{\headword{eraeya}  \pronnote{eraeja}
  \pos{Copular verb}  \sense{
    \definition{past singular form of eran}    \example{
      \exsen{Kilikiliangae dinduag a ada gogeyo do oba mokowang abal tatuma da eraeya.}      \extran{They ran happily like this to where their favorite washing place was.}}}
  \variants{\freevarlabs \vartext{ere, eraya, ereya}}}

\entry{eraeyag}{\headword{eraeyag}
  \pos{Copular verb}  \sense{
    \definition{past plural form of eran}}}

\entry{erag}{\headword{erag}
  \pos{Copular verb}  \sense{
    \definition{present plural form of eran}    \example{
      \exsen{Gagäll rabis wattällnen ma da erag duli utale mae.}      \extran{The places that are for discarding rubbish are far away.}}    \example{
      \exsen{Ngattong ddone ada deya, ddäddäg a erag?}      \extran{Before there were lots of them; where are the animals?}}}
  \variants{\Dialectal Variant \vartext{eragän}}}

\entry{eragaeya}{\headword{eragaeya}
  \pos{Copular verb}  \sense{
    \definition{past plural form of eran}    \example{
      \exsen{Ddob kollba da ngäma eragaeya, ddob llayabira dasiaemeya.}      \extran{Some of the fish that were ours (excl.), we gave them away to other people.}}}
  \variants{\freevarlab \vartext{erageya}}}

\entry{eragän}{\headword{eragän}
  \variantof{Dialectal Variant of \vartext{erag}}}

\entry{erageya}{\headword{erageya}
  \variantof{freevarlab of \vartext{eragaeya}}}

\entry{eragwaeya}{\headword{eragwaeya}  \pronnote{eraɡwaeja}
  \pos{Copular verb}  \sense{
    \definition{past dual form of eran}    \example{
      \exsen{Oba tämamae ttoen bällam ngasnen a eragwaeya, llame dagwaeya.}      \extran{Everything that the two of them did, they did together.}}}
  \variants{\freevarlabs \vartext{eragwe, eragweya}}}

\entry{eragwe}{\headword{eragwe}
  \variantof{freevarlab of \vartext{eragwaeya}}}

\entry{eragweya}{\headword{eragweya}
  \variantof{freevarlab of \vartext{eragwaeya}}}

\entry{eräl}{\headword{eräl}
  \pos{Transitive verb}  \sense{
    \definition{to name}}}

\entry{=eralla}{\headword{=eralla}  \pronnote{eraɽa}
  \pos{Auxiliary verb}  \sense{\textbf{1.} 
    \definition{aux.prs.1|2nsgA>3sgP}}  \sense{\textbf{2.} 
}}

\entry{=eralle}{\headword{=eralle}  \pronnote{eraɽe}
  \pos{Auxiliary verb}  \sense{
    \definition{aux.prs.2sgA>3sgP}}}

\entry{=erallo}{\headword{=erallo}  \pronnote{eraɽo}
  \pos{Auxiliary verb}  \sense{
    \definition{aux.prs.3nsgA>3sgP}}}

\entry{Eramang}{\headword{Eramang}  \note{Places around Limol}  \pronnote{ɛramaŋ}
  \pos{Proper noun}  \sense{
    \definition{Eramang (the main canoe place in Limol; for camping, fishing, hunting)}    \note{Places around Limol}    \example{
}    \example{
}    \example{
}}}

\entry{erame}{\headword{erame}  \note{Question words}  \pronnote{erame}
  \pos{Relative pronoun}  \sense{\textbf{1.} 
    \definition{locative form of era; where, in which}    \note{Question words}    \example{
      \exsen{Llimoll a ttängäm dan ngäna erame gozeg.}      \extran{Limol is the place where I was born.}}    \example{
      \exsen{Ngämo kame dan däbe pazi da ngäna erame gozeg.}      \extran{I do not know in which year I was born.}}}  \sense{\textbf{2.} 
    \definition{where, in which (locative form of era)}    \note{Question words}    \example{
      \exsen{Erame bibi sana de nattkoealla?}      \extran{Where did you (pl.) cut the sago?}}    \example{
}}}

\entry{eran}{\headword{eran}
  \pos{Copular verb}  \sense{
    \definition{copular form of era (present singular form)}    \example{
      \exsen{Däbe abo igi mi eran, däbe bamllamän.}      \extran{Then she will hold that one that's on the bottom.}}    \example{
      \exsen{Obo sana da eran?}      \extran{Where is his sago? / Which is his sago?}}    \example{
      \exsen{Bäne eka da eran?}      \extran{What is your language?}}}
  \variants{\Dialectal Variant \vartext{eranän}}}

\entry{=eran}{\headword{=eran}  \pronnote{eran}
  \pos{Auxiliary verb}  \sense{\textbf{1.} 
    \definition{auxiliary verb in the present tense that marks a first or third person singular actor acting on a third person singular patient}}  \sense{\textbf{2.} 
}  \sense{\textbf{3.} 
}}

\entry{eranän}{\headword{eranän}
  \variantof{Dialectal Variant of \vartext{eran}}}

\entry{erang}{\headword{erang}  \lexeme{erang}  \note{Family ties}  \pronnote{eraŋ}
  \pos{Kinship noun}  \sense{
    \definition{exchange sibling; exchange brother (man's wife's brother who marries the man's sister; reciprocal); exchange sister (woman's husband's sister who marries the woman's brother; reciprocal)}    \note{Family ties}    \example{
      \exsen{Ngämlle erang a ddone dan.}      \extran{I don't have an exchange sibling.}}}}

\entry{erängazmäg}{\headword{erängazmäg}
  \variantof{SpellingVariant of \vartext{erngazmäg}}}

\entry{erängg}{\headword{erängg}  \pronnote{erəŋɡ}
  \pos{Transitive verb}  \sense{
    \definition{to test, try; taste}    \example{
      \exsen{Tos de dägazen nyäng ik att de a batri käp de dazer. Angde daerängg mermer gogllayän.}      \extran{I took out the torch from inside my bag and put in the batteries. When I tested it, it properly shone.}}    \example{
      \exsen{Be ami ddäddägeyo, mokwang abal daeränggeyo.}      \extran{But those who tried eating it found that it tasted delicious.}}}}

\entry{erany}{\headword{erany}  \note{Things a person can do}  \pronnote{eraɲ}
  \pos{Noun}  \sense{\textbf{1.} 
    \definition{scream}    \note{Things a person can do}    \example{
      \exsen{Ngäna erany de dandär.}      \extran{I heard a scream.}}}  \sense{\textbf{2.} 
    \definition{to scream, shout, yell (at)}    \note{Things a person can do}    \example{
      \exsen{Bogo erany allan.}      \extran{He is screaming.}}    \example{
      \exsen{Lla da män kälsre de erany nägawan.}      \extran{The man shouted at the little girl.}}}
  \variants{\Unspecified Variant \vartext{yärany}}}

\entry{erär}{\headword{erär}
  \pos{Transitive verb}  \sense{\textbf{1.} 
    \definition{to name; pass down a name}    \example{
      \exsen{Ngäma pemli makäp me gänyan ge ernan erallo.}      \extran{Here in our (excl.) family, we are passing on this name.}}    \example{
      \exsen{Ngämo kame dan aeya ngämo bin di ge daerän.}      \extran{I don't know who gave me this name.}}}  \sense{\textbf{2.} 
    \definition{to measure}    \example{
      \exsen{Ge za da llo ernan ma dan.}      \extran{This thing is for measuring trees.}}}}

\entry{eraya}{\headword{eraya}
  \variantof{freevarlab of \vartext{eraeya}}  \pronnote{eraja}}

\entry{ere}{\headword{ere}
  \variantof{freevarlab of \vartext{eraeya}}  \pronnote{ere}}

\entry{erem}{\headword{erem}  \note{Question words}  \pronnote{erem}
  \pos{Relative pronoun}  \sense{\textbf{1.} 
    \definition{accusative form of era}    \note{Question words}    \example{
      \exsen{Angde ubi gongoseyo, ubi ikop dägageyo mamam mätta da Meri erem de dibeyän audaban.}      \extran{When they returned, they saw that the red yam that Mary had planted disappeared.}}    \example{
      \exsen{Ngäna gem eka de ttättle anggan, ge ddob lla da erem gomoenen amallo.}      \extran{I correct these words that some people mispronounce.}}}  \sense{\textbf{2.} 
    \definition{accusative form of era}    \note{Question words}    \example{
      \exsen{Erem mätta de bongo moko eralle?}      \extran{Which yam do you like?}}    \example{
      \exsen{dade eremde}      \extran{whichever}}    \example{
}}}

\entry{ereya}{\headword{ereya}
  \variantof{freevarlab of \vartext{eraeya}}  \pronnote{ereja}}

\entry{Erga}{\headword{Erga}  \pronnote{ɛrɡa}
  \pos{Proper noun}  \sense{
    \definition{female personal name}}}

\entry{ergod}{\headword{ergod}  \note{Verbs of Motion}  \pronnote{erɡod}
  \pos{Intransitive verb}  \sense{
    \definition{to crawl}    \note{Verbs of Motion}    \example{
      \exsen{Llɨg kälsre da ergodnen allan.}      \extran{The child is crawling.}}    \example{
}    \example{
}}}

\entry{ergodag}{\headword{ergodag}  \note{Stages of life}  \pronnote{erɡodaɡ}
  \pos{Noun}  \sense{
    \definition{toddler}    \note{Stages of life}}}

\entry{eria}{\headword{eria}
  \pos{Noun}  \sense{
    \definition{area}    \example{
      \exsen{Kawam eria me dan oba ttängäm a duli.}      \extran{Their land is there, in the Kawam area.}}}}

\entry{Eric}{\headword{Eric}  \pronnote{ɛrik}
  \pos{Proper noun}  \sense{
    \definition{male personal name}}}

\entry{erkoll}{\headword{erkoll}
  \pos{Noun}  \sense{
    \definition{dirty water}    \example{
      \exsen{Bongo ngänäm dangelbne, ddone ikop dagne, erkoll a dangttawalle.}      \extran{You were looking for me and didn't see me; the dirty water hid me.}}}}

\entry{Erme}{\headword{Erme}
  \pos{Proper noun}  \sense{
    \definition{male personal name}}}

\entry{erngazäg}{\headword{erngazäg}
  \variantof{SpellingVariant of \vartext{erngazeg}}}

\entry{erngazeg}{\headword{erngazeg}  \lexeme{erngazeg}  \pronnote{erŋazeɡ}
  \pos{Kinship noun}  \sense{
    \definition{exchange cousin (one's parent's exchange sibling's child)}}
  \variants{\SpellingVariant \vartext{erngazäg}}}

\entry{erngazenda}{\headword{erngazenda}  \lexeme{erngazenda}  \pronnote{erŋazenda}
  \pos{Kinship noun}  \sense{
    \definition{exchange uncle (one's parent's exchange brother)}}}

\entry{erngazmäg}{\headword{erngazmäg}  \lexeme{erngazmäg}  \pronnote{erŋazməɡ}
  \pos{Kinship noun}  \sense{
    \definition{exchange aunt (one's parent's exchange sister)}}
  \variants{\SpellingVariant \vartext{erängazmäg}}}

\entry{ero}{\headword{ero}  \pronnote{ero}
  \pos{Relative pronoun}  \sense{\textbf{1.} 
    \definition{where}    \example{
      \exsen{Dirom de ngäna däbe darälgoe do llo da ulle da ero daeya,}      \extran{I dragged that cassowary to where that big tree was.}}}  \sense{\textbf{2.} 
    \definition{where}    \example{
      \exsen{Ero dan ma da?}      \extran{Where is the house?}}}}

\entry{Erodias}{\headword{Erodias}
  \pos{Proper noun}  \sense{
}}

\entry{eroe}{\headword{eroe}
  \pos{Locational}  \sense{
    \definition{side}}}

\entry{erongg}{\headword{erongg}  \pronnote{eroŋɡ}
  \pos{Verb}  \sense{
    \definition{to start weaving (crossing strings)}    \example{
      \exsen{Essie nga wiya-o, bongo ngämo naerongg tater de.}      \extran{Essie, you come and start weaving my mat.}}}}

\entry{erowattäm}{\headword{erowattäm}  \pronnote{erowaʈ͡ʂəm}
  \pos{Relative pronoun}  \sense{\textbf{1.} 
    \definition{ablative form of ero}    \example{
      \exsen{Ddone aya umllang gogon ada erowattäm da ge za da gogezän.}      \extran{No one knew from where these things came out.}}}  \sense{\textbf{2.} 
    \definition{ablative form of ero}    \example{
      \exsen{Ge lla da erowattäm da ngänttägmäll allan?}      \extran{From where is this man arriving?}}}}

\entry{erowe}{\headword{erowe}
  \pos{Relative pronoun}  \sense{\textbf{1.} 
    \definition{allative form of ero}    \example{
      \exsen{Ge mälla da iba kame da erowede ibi allan.}      \extran{We (incl.) don't know where this woman is going.}}}  \sense{\textbf{2.} 
    \definition{allative form of ero}    \example{
      \exsen{Bongo erowe ibi alle?}      \extran{Where are you going?}}}}

\entry{es}{\headword{es}  \pronnote{es}
  \pos{Interjection}  \sense{
    \definition{come (command given to animals)}}}

\entry{esam}{\headword{esam}  \note{Medicine}  \pronnote{esam}
  \pos{Noun}  \sense{
    \definition{lemongrass}    \scientificname{Cymbopogon}    \note{Medicine}    \example{
}    \example{
}    \example{
}}}

\entry{Ese}{\headword{Ese}
  \pos{Proper noun}  \sense{
    \definition{male personal name}}}

\entry{eso}{\headword{eso}  \pronnote{eso}
  \pos{Interjection}  \sense{
    \definition{thank you}    \example{
      \exsen{Eso eka nägagan.}      \extran{He thanked him.}}    \example{
      \exsen{Eso ulle.}      \extran{Thank you very much.}}    \example{
      \exsen{Eso ulle bänene kandärmang e.}      \extran{Thank you for the gift.}}    \example{
      \exsen{Eso ulle bänene ngämingg e.}      \extran{Thank you for the help.}}}}

\entry{Essie}{\headword{Essie}  \pronnote{essie}
  \pos{Proper noun}  \sense{
    \definition{female personal name}}}

\entry{Esta}{\headword{Esta}
  \pos{Proper noun}  \sense{
    \definition{female personal name}}}

\entry{Evelyn}{\headword{Evelyn}  \pronnote{ɛvɛlɪn}
  \pos{Proper noun}  \sense{
    \definition{female personal name}}}

\entry{ewatta}{\headword{ewatta}
  \variantof{Unspecified Variant of \vartext{ematta}}}

\entry{ewede}{\headword{ewede}  \pronnote{ewede}
  \pos{Relative pronoun}  \sense{\textbf{1.} 
    \definition{why}    \example{
      \exsen{Census, ddob ako ttoen a dadegaeya, ewede deyaralle patrol a.}      \extran{There was the census, and also other things: that's why the patrol came.}}}  \sense{\textbf{2.} 
    \definition{why}    \example{
      \exsen{Ewede bongo bäne känyer yununin alle?}      \extran{Why are you sleeping by yourself?}}}
  \variants{\freevarlab \vartext{yewede}}
  \variants{\Unspecified Variant \vartext{yowede}}}

\entry{ewembe}{\headword{ewembe}  \lexeme{ewembe}  \pronnote{ewembe}
  \pos{Noun}  \sense{
    \definition{type of very big yam with a white interior, thorns, and hairs}}}

\entry{example}{\headword{example}
  \variantof{freevarlab of \vartext{English loanword}}}

\entry{=eya}{\headword{=eya}
  \variantof{freevarlab of \vartext{=aeya}}}

\entry{ezi}{\headword{ezi}  \note{Tools}  \pronnote{ezi}
  \pos{Noun}  \sense{
    \definition{sharp edge}    \note{Tools}    \example{
}    \example{
}    \example{
}}}

\entry{Ezra}{\headword{Ezra}  \pronnote{ezra}
  \pos{Proper noun}  \sense{
    \definition{male personal name}}}

\chapter{F}


\entry{faeb}{\headword{faeb}
  \variantof{SpellingVariant of \vartext{paeb}}}

\entry{fama}{\headword{fama}
  \variantof{SpellingVariant of \vartext{pama}}}

\entry{family}{\headword{family}
  \variantof{SpellingVariant of \vartext{pemli}}}

\entry{Feliks}{\headword{Feliks}
  \variantof{SpellingVariant of \vartext{Felix}}}

\entry{Felix}{\headword{Felix}  \pronnote{feliks}
  \pos{Proper noun}  \sense{
    \definition{male personal name}}
  \variants{\SpellingVariant \vartext{Feliks}}}

\entry{fellows}{\headword{fellows}
  \variantof{Dialectal Variant of \vartext{English loanword}}}

\entry{fellowship}{\headword{fellowship}
  \variantof{Dialectal Variant of \vartext{English loanword}}}

\entry{femli}{\headword{femli}
  \variantof{SpellingVariant of \vartext{pemli}}}

\entry{feel}{\headword{feel}
  \variantof{Dialectal Variant of \vartext{English loanword}}}

\entry{fifti}{\headword{fifti}
  \variantof{SpellingVariant of \vartext{pipti}}}

\entry{fiftin}{\headword{fiftin}
  \variantof{SpellingVariant of \vartext{piptin}}}

\entry{first}{\headword{first}
  \variantof{Dialectal Variant of \vartext{English loanword}}}

\entry{five}{\headword{five}
  \variantof{SpellingVariant of \vartext{paeb}}}

\entry{Flora}{\headword{Flora}  \pronnote{flora}
  \pos{Proper noun}  \sense{
    \definition{female personal name}}}

\entry{fo}{\headword{fo}
  \variantof{SpellingVariant of \vartext{po}}}

\entry{forti}{\headword{forti}
  \variantof{SpellingVariant of \vartext{poti}}}

\entry{fortiz}{\headword{fortiz}
  \variantof{Dialectal Variant of \vartext{English loanword}}}

\entry{forty}{\headword{forty}
  \variantof{SpellingVariant of \vartext{poti}}}

\entry{foti}{\headword{foti}
  \variantof{SpellingVariant of \vartext{poti}}}

\entry{four}{\headword{four}
  \variantof{SpellingVariant of \vartext{po}}}

\entry{food}{\headword{food}
  \variantof{freevarlab of \vartext{English loanword}}}

\entry{Francis}{\headword{Francis}  \pronnote{fransɪs}
  \pos{Proper noun}  \sense{
    \definition{male personal name}}}

\entry{Frank}{\headword{Frank}  \pronnote{freŋk}
  \pos{Proper noun}  \sense{
    \definition{male personal name}}}

\entry{from}{\headword{from}
  \variantof{Dialectal Variant of \vartext{English loanword}}}

\chapter{G}


\entry{=g}{\headword{=g}
  \pos{Verbal enclitic}  \sense{
    \definition{cop.prs.plS}}
  \variants{\Dialectal Variant \vartext{=gän}}}

\entry{g-}{\headword{g-}  \pronnote{ɡ}
  \pos{Verb}  \sense{
    \definition{rem.intr}}}

\entry{g}{\headword{g}
  \variantof{Dialectal Variant of \vartext{ga}}  \pronnote{ɡ}}

\entry{ga}{\headword{ga}
  \pos{Auxiliary verb}  \sense{
}
  \variants{\Dialectal Variant \vartext{g}}}

\entry{gäbä~}{\headword{gäbä~}
  \variantof{freevarlab of \vartext{RED~}}}

\entry{gäba}{\headword{gäba}  \pronnote{ɡəba}
  \pos{Noun}  \sense{
    \definition{shade}    \example{
      \exsen{Kottllam a bogo ada gongnomenynän adaka ddia da zäme ngattong agan, be bogo de amne me gotarnän llo gäba me.}      \extran{Turtle thought that Deer was already first, but he was in the middle sleeping in the shade of a tree.}}    \example{
      \exsen{Ubi ttoenttoen gogneyo llo gäba me.}      \extran{They were telling stories in the shade.}}}}

\entry{Gabadi}{\headword{Gabadi}
  \pos{Proper noun}  \sense{
    \definition{female personal name}}}

\entry{gäbaeb}{\headword{gäbaeb}
  \pos{Intransitive verb}  \sense{
    \definition{to jump}    \example{
      \exsen{Män a kädkädag walle we mɨnyi bogbaebän.}      \extran{The girl will jump into the cold water.}}}}

\entry{Gäbag}{\headword{Gäbag}
  \pos{Proper noun}  \sense{
    \definition{male personal name}}}

\entry{gäbagäba}{\headword{gäbagäba}  \pronnote{ɡəbaɡəba}
  \pos{Modifier}  \sense{
    \definition{shady}    \example{
      \exsen{Llo gäbagäba de ikop dägagän.}      \extran{He saw a shady tree.}}}}

\entry{gabän}{\headword{gabän}  \pronnote{ɡabən}
  \pos{Transitive verb}  \sense{
    \definition{to cut in half or to split, e.g. to split a deer thigh with someone.}}}

\entry{gabän}{\headword{gabän}  \note{Parts of the Body}  \pronnote{ɡabən}
  \pos{Noun}  \sense{\textbf{1.} 
    \definition{wrist}    \note{Parts of the Body}    \example{
}    \example{
}    \example{
}}  \sense{\textbf{2.} 
    \definition{six (lit. wrist; body counting numeral)}    \note{Parts of the Body}}}

\entry{gäbän}{\headword{gäbän}  \pronnote{ɡəbən}
  \pos{Intransitive verb}  \sense{
    \definition{to jump, hop}    \example{
      \exsen{Gäbänmällang dan.}      \extran{He is hopping.}}    \example{
      \exsen{Män a gall att agäbänan.}      \extran{The girl jumped out of the canoe.}}}}

\entry{gabana}{\headword{gabana}  \note{Government}  \pronnote{ɡabana}
  \pos{Noun}  \sense{
    \definition{governor}    \note{Government}    \example{
}    \example{
}    \example{
}}}

\entry{gäbgäb}{\headword{gäbgäb}  \note{Trees}  \pronnote{ɡəbɡəb}
  \pos{Noun}  \sense{
    \definition{type of tree}    \note{Trees}    \example{
}    \example{
}    \example{
}}}

\entry{gablle}{\headword{gablle}  \note{Parts of the Body}  \pronnote{ɡabɽe}
  \pos{Noun}  \sense{
    \definition{hip; waist}    \note{Parts of the Body}    \example{
      \exsen{Obo gablle mättnan ma källän a ddäddäg ttoe dädär alle ngasngesatt daeya.}      \extran{The belt he wore around his waist was made from animal hide.}}    \example{
}    \example{
}}}

\entry{gabma}{\headword{gabma}  \pronnote{ɡabma}
  \pos{Noun}  \sense{\textbf{1.} 
    \definition{white person}    \example{
      \exsen{ttongo gabma bo bin}      \extran{a white man's name}}}  \sense{\textbf{2.} 
    \definition{white, Western; foreign}}}

\entry{gabma bägäl}{\headword{gabma bägäl}
  \pos{Noun}  \sense{
    \definition{gun}}}

\entry{gäbmäll}{\headword{gäbmäll}  \note{Verbs of Motion}  \pronnote{ɡəbməɽ}
  \pos{Verb}  \sense{
    \definition{to skip}    \note{Verbs of Motion}    \example{
      \exsen{Bogo gäbmäll allan.}      \extran{He is skipping.}}    \example{
      \exsen{Ubi gäbmäll anggan.}      \extran{They are skipping.}}    \example{
}}}

\entry{gäbmällgäbmäll}{\headword{gäbmällgäbmäll}  \pronnote{ɡəbməɽɡəbməɽ}
  \pos{Adverb}  \sense{
    \definition{skipping}    \example{
      \exsen{Bogo ttäle gäbmällgäbmäll ibi allan.}      \extran{He goes, skipping.}}}}

\entry{gabman}{\headword{gabman}
  \variantof{Unspecified Variant of \vartext{gabmantt}}}

\entry{gabmantt}{\headword{gabmantt}  \pronnote{ɡabmanʈ}
  \pos{Noun}  \sense{
    \definition{government}}
  \variants{\Unspecified Variant \vartext{gabman}}}

\entry{gäboll}{\headword{gäboll}  \lexeme{gäboll}  \pronnote{ɡəboɽ}
  \pos{Noun}  \sense{
    \definition{magpie-lark}    \scientificname{Grallina cyanoleuca neglecta}}}

\entry{gädagäde}{\headword{gädagäde}  \pronnote{ɡədaɡəde}
  \pos{Transitive verb}  \sense{\textbf{1.} 
    \definition{to beat sago, pound sago}    \example{
      \exsen{Beyawaebe dägdene.}      \extran{Alone, he beat sago.}}}  \sense{\textbf{2.} 
    \definition{to sharpen}    \example{
      \exsen{Ngäna tri pos dae dägdeneg.}      \extran{I only sharpened three posts.}}}  \sense{\textbf{3.} 
    \definition{to throw a tantrum}}}

\entry{gädän}{\headword{gädän}
  \pos{Intransitive verb}  \sense{
    \definition{to land}}}

\entry{gädd~}{\headword{gädd~}
  \variantof{Infinitival reduplicant of \vartext{RED~}}}

\entry{gäddgädd}{\headword{gäddgädd}  \pronnote{ɡəɖ͡ʐɡəɖ͡ʐ}
  \pos{Transitive verb}  \sense{\textbf{1.} 
    \definition{to fight}    \example{
      \exsen{Lla da gäddgädd e bobllabän.}      \extran{The men go to fight.}}    \example{
      \exsen{Bongo sämell aba peyang gäddgädd alle.}      \extran{You are fighting with the pig.}}    \example{
      \exsen{Ubi gogäddän.}      \extran{They fought.}}}  \sense{\textbf{2.} 
    \definition{to catch}    \example{
      \exsen{Mälla da sisri ag me kok de nägäddaeballo kollba tokong e.}      \extran{This morning, the women caught grasshopers for bait.}}}  \sense{\textbf{3.} 
    \definition{to cut (leaves)}    \example{
      \exsen{Ngäma Dowabunang me gällall de ke ami dägäddaebeyo?}      \extran{Who cut our (excl.) pandanus leaves in Dowabunang?}}}}

\entry{gae nge}{\headword{gae nge}  \pronnote{ɡae ŋe}
  \pos{Noun}  \sense{
    \definition{type of palm with coconuts with a red or green exocarp; the husk is chewed and the young coconut water is drunk}}}

\entry{Gaem}{\headword{Gaem}  \pronnote{ɡaem}
  \pos{Proper noun}  \sense{
    \definition{female personal name}}}

\entry{=gaeya}{\headword{=gaeya}
  \pos{Verbal enclitic}  \sense{
    \definition{cop.pst.plS}}
  \variants{\freevarlabs \vartext{=gaya, =geya}}}

\entry{=gaeyo}{\headword{=gaeyo}
  \pos{Verbal enclitic}  \sense{
    \definition{cop.prs.duS}}
  \variants{\Dialectal Variant \vartext{=gwän}}
  \variants{\Unspecified Variant \vartext{=giyu}}
  \variants{\freevarlab \vartext{=geyo}}}

\entry{gag}{\headword{gag}
  \pos{Transitive verb}  \sense{
    \definition{to store}    \example{
      \exsen{Mätta de gaguma we dägagaeballo.}      \extran{They stored the yams in the yamhouse.}}}}

\entry{gägäb ine}{\headword{gägäb ine}
  \pos{Noun}  \sense{
    \definition{dew}    \example{
      \exsen{Apapi da llo ttam me gägäb ine nanong dan.}      \extran{The butterfly drinks the dew on leaves.}}}}

\entry{gägabäll}{\headword{gägabäll}  \pronnote{ɡəɡabəɽ}
  \pos{Intransitive verb}  \sense{
    \definition{to stand}    \example{
      \exsen{Tutu we dirngäneyo a ubi dägabllabän.}      \extran{She pulled her to land, and they stood.}}    \example{
      \exsen{Dibaya kälepalle danyroene, angde dowae mäse diga, ddia da dam dindugän a dägabällän.}      \extran{I crept slowly towards it, and when I got close, the deer ran away and stood.}}    \example{
      \exsen{Oba ulle binang a dägabälleyo ada ddia bo sawe ttang alle a kottllam bo ttätt ttang alle.}      \extran{Their big bosses stood like this: the deer's on the left hand side and the turtle's on the right hand side.}}}}

\entry{gagägagäll}{\headword{gagägagäll}
  \pos{Adverb}  \sense{
    \definition{badly, poorly}    \example{
      \exsen{Ubi Ende eka de gagägagäll däpanyallo.}      \extran{They spoke the Ende language poorly.}}}
  \variants{\Fast Speech Variant \vartext{gagagäll}}}

\entry{gagagäll}{\headword{gagagäll}
  \variantof{Fast Speech Variant of \vartext{gagägagäll}}}

\entry{gagäll}{\headword{gagäll}  \note{Adjectives}  \pronnote{ɡaɡəɽ}
  \pos{Modifier}  \sense{\textbf{1.} 
    \definition{bad, rotten}    \note{Adjectives}    \example{
      \exsen{Ge kollba da gagäll dan.}      \extran{This fish is bad.}}    \example{
      \exsen{Da ngäna bol de ttang alle bäbädd, ngämo tupi di mɨnyi gagäll bäga.}      \extran{If I hit the ball with my hand, I will hurt my index finger.}}    \example{
      \exsen{Ngämo giddoll a ngattong ai daeyag a be ngäna sisri gagäll ttam e azenan.}      \extran{My life was good at first, but now, I have entered a bad life.}}    \example{
      \exsen{Panya dae ddone gagäll gogon.}      \extran{Only the pineapple did not go bad.}}}  \sense{\textbf{2.} 
    \definition{dead}    \note{Adjectives}    \example{
      \exsen{Ngämo kameny me de gagäll gogon.}      \extran{He died during my absence.}}}  \sense{\textbf{3.} 
    \definition{sin}    \note{Adjectives}    \example{
      \exsen{Ibim deyanglläbnegän gagäll kup att de.}      \extran{He brought us out of the depths of sin.}}}}

\entry{gagäll ine}{\headword{gagäll ine}
  \pos{Noun}  \sense{
    \definition{beer}}}

\entry{gagälläd}{\headword{gagälläd}  \pronnote{ɡaɡəɽəd}
  \pos{Transitive verb}  \sense{
    \definition{to cut leaf}}}

\entry{gageb}{\headword{gageb}
  \pos{Noun}  \sense{
    \definition{hind leg}}}

\entry{gäglib}{\headword{gäglib}  \pronnote{ɡəɡlib}
  \pos{Transitive verb}  \sense{\textbf{1.} 
    \definition{to chase, scare away/off (wild animals)}    \example{
      \exsen{Däräng a ddia de diglibiyu Madura pate.}      \extran{The dogs chased the deer, steering it towards Madura.}}    \example{
      \exsen{Ngäna käza de nägliban.}      \extran{I scared away the crocodile.}}}  \sense{\textbf{2.} 
    \definition{to pluck}    \example{
      \exsen{Pa kom de dägläbaebeyo, tot do dangesaemeyo.}      \extran{They plucked the bird's feathers and made a mess there.}}}}

\entry{gäglläb}{\headword{gäglläb}  \pronnote{ɡəɡɽəb}
  \pos{Transitive verb}  \sense{
    \definition{to weed for the first time (hard work to remove the large plants)}}}

\entry{gaguma}{\headword{gaguma}  \note{Gardens}  \pronnote{ɡaɡuma}
  \pos{Noun}  \sense{
    \definition{yamhouse}    \note{Gardens}    \example{
      \exsen{Mäkat llɨgllɨgag a dan bäne gaguma me.}      \extran{There is a rat with babies in your yamhouse.}}    \example{
}    \example{
}}}

\entry{Gaima}{\headword{Gaima}  \pronnote{ɡaima}
  \pos{Proper noun}  \sense{
    \definition{Gaima (toponym)}    \example{
}    \example{
}    \example{
}}}

\entry{gaimbi}{\headword{gaimbi}  \note{Trees}  \pronnote{ɡaimbi}
  \pos{Noun}  \sense{
    \definition{type of fruit tree}    \note{Trees}    \example{
}    \example{
}    \example{
}}}

\entry{gajen}{\headword{gajen}
  \variantof{SpellingVariant of \vartext{gazen}}  \pronnote{ɡazen}}

\entry{gal}{\headword{gal}  \note{Marriage}  \pronnote{ɡal}
  \pos{Noun}  \sense{
    \definition{food offering}    \note{Marriage}    \example{
}    \example{
}    \example{
}}}

\entry{gäl}{\headword{gäl}  \note{Trees}  \pronnote{ɡəl}
  \pos{Noun}  \sense{
    \definition{type of tree that grows in the bush with white flowers, brown seeds, and fruit with a yellow pericarp and green exocarp}    \note{Trees}    \example{
}    \example{
}    \example{
}}}

\entry{Gäläb}{\headword{Gäläb}  \note{Places around Limol}  \pronnote{ɡələb}
  \pos{Proper noun}  \sense{
    \definition{Gälab (sago place and Geoff Rowak's camping place in Limol)}    \note{Places around Limol}    \example{
}    \example{
}    \example{
}}}

\entry{Gälabi}{\headword{Gälabi}  \pronnote{ɡəlabi}
  \pos{Proper noun}  \sense{
    \definition{Gälabi (Wipi-speaking village in Oriomo-Bituri Rural LLG; from Limol, one must pass through Wipim)}    \example{
}    \example{
}    \example{
}}}

\entry{gälas}{\headword{gälas}  \note{Weaving}  \pronnote{ɡəlas}
  \pos{Noun}  \sense{
    \definition{weaving pattern with concentric squares}    \note{Weaving}    \example{
}    \example{
}    \example{
}}}

\entry{gälbän}{\headword{gälbän}  \pronnote{ɡəlbən}
  \pos{Transitive verb}  \sense{\textbf{1.} 
    \definition{to knock (a fruit)}    \example{
      \exsen{Polle ma ik i angde gozenän, up wo gälbänanatt de ikop dägagän.}      \extran{When he went inside the fence, he saw the fallen ripe banana.}}}  \sense{\textbf{2.} 
    \definition{to defeather, depilate, remove hair; remove teeth}}}

\entry{galbe}{\headword{galbe}  \lexeme{galbe}  \pronnote{ɡalbe}
  \pos{Noun}  \sense{
    \definition{purple/greater yam}    \scientificname{Dioscorea alata}}
  \variants{\SpellingVariant \vartext{galbi}}}

\entry{galbi}{\headword{galbi}
  \variantof{SpellingVariant of \vartext{galbe}}}

\entry{gäleb}{\headword{gäleb}  \note{Trees}  \pronnote{ɡəleb}
  \pos{Noun}  \sense{
    \definition{type of tree}    \note{Trees}    \example{
}    \example{
}    \example{
}}}

\entry{gälgäl}{\headword{gälgäl}
  \variantof{Unspecified Variant of \vartext{gallgall}}}

\entry{galib}{\headword{galib}  \note{Traditional Arrows}  \pronnote{ɡalib}
  \pos{Noun}  \sense{
    \definition{type of spear}    \note{Traditional Arrows}    \example{
}    \example{
}    \example{
}}}

\entry{Galibma}{\headword{Galibma}  \note{Places around Limol}  \pronnote{ɡalibma}
  \pos{Proper noun}  \sense{
    \definition{sacred place of Biku (Madura) Kangge (on the road to Bisuaka, after Limol ma kuddäll; bush area, creek)}    \note{Places around Limol}    \example{
}    \example{
}    \example{
}}}

\entry{galigali}{\headword{galigali}  \note{Types of Taro}  \pronnote{ɡaliɡali}
  \pos{Noun}  \sense{
    \definition{type of small taro}    \note{Types of Taro}    \example{
}    \example{
}    \example{
}}}

\entry{gall}{\headword{gall}  \note{Water}  \pronnote{ɡaɽ}
  \pos{Noun}  \sense{
    \definition{canoe, boat}    \note{Water}    \example{
      \exsen{Amo gall sisor da däbe?}      \extran{Whose new canoe is this?}}    \example{
}}
  \variants{\SpellingVariant \vartext{gaall}}}

\entry{gall gullem}{\headword{gall gullem}  \note{Water}  \pronnote{ɡaɽ ɡuɽem}
  \pos{Noun}  \sense{
    \definition{canoe outrigger}    \note{Water}    \example{
}    \example{
}    \example{
}}}

\entry{gall ngattongag}{\headword{gall ngattongag}  \pronnote{ɡaɽ ŋaʈ͡ʂoŋaɡ}
  \pos{Noun}  \sense{
    \definition{navigator (of a canoe)}}}

\entry{gall onyang}{\headword{gall onyang}
  \pos{Noun}  \sense{
    \definition{operator (of a canoe)}}}

\entry{gallab}{\headword{gallab}  \pronnote{ɡaɽab}
  \pos{Transitive verb}  \sense{
    \definition{to open (one's mouth)}    \example{
      \exsen{Eka tameny a bongo abo mudan, kälae ttatt gallab a mudan.}      \extran{You are not allowed to speak or open your jaw even a little.}}    \example{
      \exsen{Abo ttatt de gallab a mudan.}      \extran{Don't you say a word.}}}}

\entry{gällall}{\headword{gällall}  \note{Type of pandanus}  \pronnote{ɡəɽaɽ}
  \pos{Noun}  \sense{
    \definition{type of pandanus}    \note{Type of pandanus}    \example{
      \exsen{Ngäma Dowabunang me gällall de ke ami dägäddaebeyo?}      \extran{Who cut our pandanus leaves at Dowabunang?}}    \example{
}    \example{
}}}

\entry{gällanggälla}{\headword{gällanggälla}
  \variantof{SpellingVariant of \vartext{gllangglla}}  \pronnote{ɡəɽaŋɡəɽa}}

\entry{gällatater}{\headword{gällatater}  \note{Trees}  \pronnote{ɡəɽatater}
  \pos{Noun}  \sense{
    \definition{type of tree}    \note{Trees}    \example{
}    \example{
}    \example{
}}}

\entry{gällbän}{\headword{gällbän}  \pronnote{ɡəɽbən}
  \pos{Intransitive verb}  \sense{
    \definition{to break}    \example{
      \exsen{Ngämo llimba da agällbänan.}      \extran{My fingernail broke.}}}}

\entry{gälle}{\headword{gälle}  \note{Trees}  \pronnote{ɡəɽe}
  \pos{Noun}  \sense{
    \definition{type of big tree that grows in the bush near big creeks with wood used for canoes, white and blue flowers, and edible red fruit}    \note{Trees}    \example{
}    \example{
}    \example{
}}}

\entry{gallgall}{\headword{gallgall}  \pronnote{ɡaɽɡaɽ}
  \pos{Noun}  \sense{
    \definition{bank, coast, shore}    \example{
      \exsen{Ubi komlla apte walle gallgall e dängkällneneyo.}      \extran{The two of them climbed onto the other shore.}}    \example{
      \exsen{Sämell a dizmollän gälgälag Geleli walle we guspullän a tämamae gomänddnegän.}      \extran{The pigs ran to the banks of the Sea of Galilee and fell in, and all of them drowned.}}}
  \variants{\Unspecified Variant \vartext{gälgäl, gällgäll}}}

\entry{gällgäll}{\headword{gällgäll}
  \variantof{Unspecified Variant of \vartext{gallgall}}}

\entry{gallgell}{\headword{gallgell}  \pronnote{ɡaɽɡeɽ}
  \pos{Intransitive verb}  \sense{\textbf{1.} 
    \definition{to remove leaves, take leaves off; remove skin, skin}}  \sense{\textbf{2.} 
    \definition{to open mouth wide when screaming}}}

\entry{gällnen}{\headword{gällnen}
  \pos{Transitive verb}  \sense{
}}

\entry{gallwab}{\headword{gallwab}  \pronnote{ɡaɽwab}
  \pos{Transitive verb}  \sense{
    \definition{to prune}}}

\entry{Galo}{\headword{Galo}  \pronnote{ɡalo}
  \pos{Proper noun}  \sense{
    \definition{male personal name}}}

\entry{Galwe}{\headword{Galwe}  \pronnote{ɡalwe}
  \pos{Proper noun}  \sense{
    \definition{male personal name}}}

\entry{Gamaewe}{\headword{Gamaewe}  \pronnote{ɡamaewe}
  \pos{Proper noun}  \sense{
    \definition{Gamaewe (toponym)}}}

\entry{gämäll}{\headword{gämäll}  \lexeme{gämäll}  \note{Market}  \pronnote{ɡəməɽ}
  \pos{Transitive coverb}  \sense{
    \definition{to steal}    \note{Market}    \example{
      \exsen{gämällang lla}      \extran{thief}}    \example{
      \exsen{mälla gämäll}      \extran{adultery (lit. 'wife-stealing')}}    \example{
      \exsen{Ngäna sana de gämäll agan bänene.}      \extran{I stole the sago from you.}}    \example{
      \exsen{Ngäna sana de gämäll agan bablle.}      \extran{I stole the sago for you.}}}}

\entry{gämäll tänggag}{\headword{gämäll tänggag}
  \pos{Transitive compound verb}  \sense{
    \definition{to steal}    \example{
      \exsen{Da bongo ddone otät de nonttog, bam mɨnyi gämäll nantägegän.}      \extran{If you do not give him food, he will steal from you.}}}}

\entry{gambäg}{\headword{gambäg}  \pronnote{ɡambəɡ}
  \pos{Intransitive verb}  \sense{
    \definition{to stand closely}}}

\entry{gambän}{\headword{gambän}
  \pos{Intransitive verb}  \sense{
    \definition{to stand}}}

\entry{gämbanengg}{\headword{gämbanengg}
  \pos{Transitive verb}  \sense{
    \definition{causative-applicative form of gambän}    \example{
      \exsen{Ngäna bam danggebänängg.}      \extran{I made you stand.}}}
  \variants{\Unspecified Variant \vartext{gambenängg}}}

\entry{gambenängg}{\headword{gambenängg}
  \variantof{Unspecified Variant of \vartext{gämbanengg}}}

\entry{gameny}{\headword{gameny}  \pronnote{ɡameɲ}
  \pos{Intransitive verb}  \sense{
    \definition{to become cloudy, when dark rain clouds are forming}}}

\entry{gamo}{\headword{gamo}  \pronnote{ɡamo}
  \pos{Noun}  \sense{
    \definition{type of large sea turtle}}}

\entry{gamo}{\headword{gamo}
  \pos{Noun}  \sense{
}}

\entry{gämoe}{\headword{gämoe}  \note{Animal verbs}  \pronnote{ɡəmoe}
  \pos{Transitive verb}  \sense{\textbf{1.} 
    \definition{to miss}    \note{Animal verbs}    \example{
      \exsen{Ddia de ddone dagmoe bazu.}      \extran{Shooting the deer, I didn't miss.}}    \example{
      \exsen{Ttongo eka nagmoeyan.}      \extran{I forgot to say one thing.}}}  \sense{\textbf{2.} 
    \definition{to miss, feel longing for}    \note{Animal verbs}    \example{
      \exsen{Ngäna bam gämoe nallan.}      \extran{I am missing you.}}}}

\entry{gamu}{\headword{gamu}  \pronnote{ɡamu}
  \pos{Noun}  \sense{
    \definition{type of ginger with flat leaves; used as medicine for centipede bites and as bait for catching flying foxes; chew it first and the flying fox will eat it and become lethargic}    \example{
}    \example{
}    \example{
}}}

\entry{gän}{\headword{gän}
  \pos{Noun}  \sense{
    \definition{gun}}}

\entry{=gän}{\headword{=gän}
  \variantof{Dialectal Variant of \vartext{=g}}
  \variants{\Dialectal Variant \vartext{=gan}}}

\entry{=gan}{\headword{=gan}
  \variantof{Dialectal Variant of \vartext{=gän}}}

\entry{=gan}{\headword{=gan}
  \pos{Verbal enclitic}  \sense{
    \definition{cop.prs.plS}}}

\entry{gan}{\headword{gan}
  \variantof{Dialectal Variant of \vartext{dadegän}}}

\entry{gangan}{\headword{gangan}  \note{Trees}  \pronnote{'gangan}
  \pos{Noun}  \sense{
    \definition{type of tree}    \note{Trees}    \example{
}    \example{
}    \example{
}}}

\entry{gänggälläm}{\headword{gänggälläm}  \pronnote{'gänggälläm}
  \pos{Transitive verb}  \sense{
    \definition{to wash (a body part or inanimate object)}    \example{
      \exsen{Ddob mälla da pane moepang de ikrol alle gällämnan amallo.}      \extran{Some women wash dirty pots with ash.}}    \example{
      \exsen{Bogo pilatt de dängllämnegan.}      \extran{He washed the dishes.}}    \example{
      \exsen{Da bongo källama atta angos, ttang anggälläm.}      \extran{If you come back from the toilet, wash your hands.}}}
  \variants{\Fast Speech Variant \vartext{gängglläm}}}

\entry{gängglläd}{\headword{gängglläd}  \pronnote{ɡəŋɡɽəd}
  \pos{Transitive verb}  \sense{\textbf{1.} 
    \definition{to shove, push}    \example{
      \exsen{Ngäna sɨmell de duduli dänggllädalle.}      \extran{I shoved the pig that way.}}}  \sense{\textbf{2.} 
    \definition{to extend (e.g. a garden)}}}

\entry{gängglläm}{\headword{gängglläm}
  \variantof{Fast Speech Variant of \vartext{gänggälläm}}}

\entry{gänggllam}{\headword{gänggllam}
  \pos{Intransitive coverb}  \sense{
    \definition{to snore}    \example{
      \exsen{Ge lla da dämen ma toko me inunin allan a mälläng gänggllam allan.}      \extran{This man is sleeping on a chair and snoring.}}}}

\entry{gany}{\headword{gany}  \pronnote{ɡaɲ}
  \pos{Transitive verb}  \sense{\textbf{1.} 
    \definition{to plant, place in the ground, put in}    \example{
      \exsen{Kup de däkällneg a pos de däganeg.}}    \example{
      \exsen{Angde pos ganen a gottamänän ngäna sipel gog.}      \extran{When the post planting finished, I took a rest.}}}  \sense{\textbf{2.} 
    \definition{to focus, tune}    \example{
      \exsen{Llan de ngämi däganyeya eka we.}      \extran{We tuned our ears for the sound.}}}}

\entry{gänya}{\headword{gänya}  \pronnote{ɡəɲa}
  \pos{Adverbial demonstrative}  \sense{\textbf{1.} 
    \definition{here (proximal)}    \example{
      \exsen{Bogo gänya angttägan.}      \extran{He is coming here.}}    \example{
      \exsen{Ngäna sisri gänya bäne peyang eka tameny allan.}      \extran{Now I am here talking with you.}}}  \sense{\textbf{2.} 
    \definition{this (proximal determiner)}    \example{
      \exsen{Bogo gänya nyongo me ibiag da.}      \extran{He often walks on this road.}}    \example{
      \exsen{Da ngäna ddone yu bipellknegalle bablle, bongo ddone gänya mer yu mi kallkäll bongllaegalle.}      \extran{If I hadn't chopped this wood for you, you wouldn't have this nice fire to warm up by.}}}
  \variants{\Dialectal Variant \vartext{gonya, gonyo}}}

\entry{gänyaballe}{\headword{gänyaballe}
  \pos{Adverbial demonstrative}  \sense{
    \definition{ablative form of gänya}    \example{
      \exsen{Gänyaballe dallän do we.}      \extran{From here, he went there.}}    \example{
      \exsen{Iba ge kollba da mullae dag, ibi gänyaballe ma we bängäsmäll.}      \extran{We (incl.) have enough fish; from here we shall return home.}}}}

\entry{gänyaeben}{\headword{gänyaeben}
  \pos{Copular verb}  \sense{
    \definition{restrictive copular form of gänya (present singular form)}    \example{
      \exsen{Ngämo eka da gänyaeben Ende eka da ge ngäna ere eka allan.}      \extran{The only language that I speak is this Ende language.}}}}

\entry{gänyaeya}{\headword{gänyaeya}  \pronnote{ɡəɲaeja}
  \pos{Copular verb}  \sense{
    \definition{past singular form of gänyan}    \example{
      \exsen{Ede ngämo ttoenttoen llätt a gänyaeya.}      \extran{So here was the end of my story.}}}
  \variants{\freevarlabs \vartext{gänyeya, gonyeya, gänyaya}}}

\entry{gänyag}{\headword{gänyag}
  \pos{Copular verb}  \sense{
    \definition{present plural form of gänyan}    \example{
      \exsen{Oba za da gänyag.}      \extran{Here are their things.}}}}

\entry{gänyagaeya}{\headword{gänyagaeya}
  \pos{Copular verb}  \sense{
    \definition{past plural form of gänyan}    \example{
      \exsen{Oba bin a gänyagaeya ada, Samson, Wawase, a Pedro.}      \extran{Here were their names: Samson, Wawase, and Pedro.}}}
  \variants{\freevarlabs \vartext{gänyagaya, gänyageya}}}

\entry{gänyagaya}{\headword{gänyagaya}
  \variantof{freevarlab of \vartext{gänyagaeya}}  \pronnote{ɡəɲaɡaja}}

\entry{gänyageya}{\headword{gänyageya}
  \variantof{freevarlab of \vartext{gänyagaeya}}  \pronnote{ɡəɲaɡeja}}

\entry{gänyageyo}{\headword{gänyageyo}  \pronnote{ɡəɲaɡejo}
  \pos{Copular verb}  \sense{
    \definition{present dual form of gänyan}    \example{
      \exsen{Ge lla komlla da gänyageyo, ibnin allo.}      \extran{These two people here, they are going.}}}}

\entry{gänyagwaeya}{\headword{gänyagwaeya}
  \pos{Copular verb}  \sense{
    \definition{past dual form of gänyan}    \example{
      \exsen{Tätäm ubi komlla gänyagwaeya gonzämeyo.}      \extran{Yesterday, the two of them were here passing through.}}}
  \variants{\freevarlab \vartext{gänyagweya}}
  \variants{\Fast Speech Variant \vartext{gänyaweya}}}

\entry{gänyagweya}{\headword{gänyagweya}
  \variantof{freevarlab of \vartext{gänyagwaeya}}  \pronnote{ɡəɲaɡweja}}

\entry{gänyamasäm}{\headword{gänyamasäm}
  \pos{Adverbial demonstrative}  \sense{
    \definition{ablative form of gänya}    \example{
      \exsen{Mälla gänyamasäm dan.}      \extran{The woman is from here.}}    \example{
      \exsen{Gänyamasäm a ddone mermer Ende eka de panypeny erallo.}      \extran{Here's why they aren't speaking Ende well.}}}
  \variants{\SpellingVariant \vartext{gänyamasen, gänyamasem}}}

\entry{gänyamasem}{\headword{gänyamasem}
  \variantof{SpellingVariant of \vartext{gänyamasäm}}}

\entry{gänyamasen}{\headword{gänyamasen}
  \variantof{SpellingVariant of \vartext{gänyamasäm}}}

\entry{gänyame}{\headword{gänyame}
  \variantof{freevarlab of \vartext{gänyme}}  \pronnote{ɡəɲame}}

\entry{gänyäme}{\headword{gänyäme}
  \variantof{freevarlab of \vartext{gänyme}}}

\entry{gänyan}{\headword{gänyan}  \note{Verbs}  \pronnote{ɡəɲan}
  \pos{Copular verb}  \sense{
    \definition{copular form of gänya (present singular form)}    \note{Verbs}    \example{
      \exsen{Gänyan obo sana da.}      \extran{Here is his sago.}}}
  \variants{\Dialectal Variant \vartext{gänyanän}}}

\entry{gänyanän}{\headword{gänyanän}
  \variantof{Dialectal Variant of \vartext{gänyan}}}

\entry{gänyaolle}{\headword{gänyaolle}  \pronnote{ɡəɲaoɽe}
  \pos{Adverbial demonstrative}  \sense{
    \definition{allative form of gänya}    \example{
      \exsen{Bäne za de gänyaolle ikom.}      \extran{Bring your things over here.}}    \example{
      \exsen{Eka sasapang a gänyaolle llɨtaem eran mälla walle.}      \extran{Various languages are arriving here along with women.}}}}

\entry{gänyaollemae}{\headword{gänyaollemae}
  \pos{Adverbial demonstrative}  \sense{
    \definition{allative form of gänya with perlative clitic}    \example{
      \exsen{Bongo gänyaollemae wiya.}      \extran{You come over here.}}}}

\entry{gänyawe}{\headword{gänyawe}
  \pos{}}

\entry{gänyaweya}{\headword{gänyaweya}
  \variantof{Fast Speech Variant of \vartext{gänyagwaeya}}}

\entry{gänyaya}{\headword{gänyaya}
  \variantof{freevarlab of \vartext{gänyaeya}}  \pronnote{ɡəɲaja}}

\entry{gänyerballe}{\headword{gänyerballe}
  \variantof{Fast Speech Variant of \vartext{gänyeriballe}}  \pronnote{ɡəɲerbaɽe}}

\entry{gänyeri}{\headword{gänyeri}
  \pos{Adverbial demonstrative}  \sense{
    \definition{this way}    \example{
      \exsen{Ngämo mäda gänyeri guddällän, gänyeri mälla de dällädän.}      \extran{My father came over here and got a local wife.}}}}

\entry{gänyeriballe}{\headword{gänyeriballe}  \pronnote{ɡəɲeribaɽe}
  \pos{Adverbial demonstrative}  \sense{
    \definition{allative form of gänyeri}    \example{
      \exsen{Wendy, mamal gänyeriballe yaupe adawatta da mo da bäne peyang bottkamän.}      \extran{Wendy, jump quickly to this side in case the bridge breaks with you.}}}
  \variants{\Fast Speech Variant \vartext{gänyerballe}}}

\entry{gänyerimae}{\headword{gänyerimae}
  \pos{Adverbial demonstrative}  \sense{
    \definition{gänyeri with perlative clitic}    \example{
      \exsen{Bongo wiya gänyerimae.}      \extran{You come over here.}}}}

\entry{gänyeya}{\headword{gänyeya}
  \variantof{freevarlab of \vartext{gänyaeya}}  \pronnote{ɡəɲeja}}

\entry{gänyme}{\headword{gänyme}  \pronnote{ɡəɲme}
  \pos{Adverbial demonstrative}  \sense{
    \definition{here}    \example{
      \exsen{Kemu bo llabun a gänyme a Malläm me dadeg.}      \extran{Kemu has relatives here and in Malam.}}    \example{
      \exsen{Tämamae lla da Ende eka de gänyme Llimoll me panypeny erallo.}      \extran{All people here in Limol speak Ende.}}}
  \variants{\freevarlabs \vartext{gänyame, gänyäme}}}

\entry{gänyo}{\headword{gänyo}
  \variantof{SpellingVariant of \vartext{gonyo}}  \pronnote{ɡəɲo}}

\entry{Gao}{\headword{Gao}
  \pos{Proper noun}  \sense{
    \definition{male personal name}}}

\entry{gaopi}{\headword{gaopi}  \pronnote{ɡaopi}
  \pos{Noun}  \sense{
    \definition{Australian pelican}    \scientificname{Pelecanus conspicillatus}}}

\entry{gaora}{\headword{gaora}  \note{Sago palms}  \pronnote{ɡaora}
  \pos{Noun}  \sense{
    \definition{type of sago}    \note{Sago palms}    \example{
}    \example{
}    \example{
}}}

\entry{gara}{\headword{gara}  \lexeme{gara}  \pronnote{ɡara}
  \pos{Noun}  \sense{
    \definition{aquatic leech}}}

\entry{Garai}{\headword{Garai}
  \variantof{SpellingVariant of \vartext{Garayi}}}

\entry{Garayi}{\headword{Garayi}  \pronnote{ɡaraji}
  \pos{Proper noun}  \sense{
    \definition{male personal name}}
  \variants{\SpellingVariant \vartext{Garai}}}

\entry{Garaz}{\headword{Garaz}  \pronnote{ɡaraz}
  \pos{Proper noun}  \sense{
    \definition{male personal name}}}

\entry{gärep}{\headword{gärep}  \pronnote{ɡərep}
  \pos{Noun}  \sense{
    \definition{grape}}}

\entry{gastol}{\headword{gastol}
  \pos{Noun}  \sense{
    \definition{type of fish}}}

\entry{gaugau}{\headword{gaugau}  \lexeme{gaugau}  \pronnote{ɡauɡau}
  \pos{Noun}  \sense{
    \definition{type of big tree that grows in the bush along creeks with wood used for canoes}}}

\entry{=gaya}{\headword{=gaya}
  \variantof{freevarlab of \vartext{=gaeya}}}

\entry{gäz}{\headword{gäz}  \pronnote{ɡəz}
  \pos{Transitive verb}  \sense{\textbf{1.} 
    \definition{to kill}    \example{
      \exsen{Emaemae pa de dägäddaebeyo.}      \extran{They killed many different types of birds.}}    \example{
      \exsen{Ge lla da sisri gäz e dan.}      \extran{This man is to be killed now.}}}  \sense{\textbf{2.} 
    \definition{to hit, beat}    \example{
      \exsen{Ge lla da mälla da bom näbäddan.}      \extran{This man beat his wife.}}    \example{
      \exsen{Bibi komllaebe gäddnan a mudageyo.}      \extran{You two, don't fight.}}}}

\entry{gazen}{\headword{gazen}  \pronnote{ɡazen}
  \pos{Transitive verb}  \sense{\textbf{1.} 
    \definition{to come out, get out, exit, escape}    \example{
      \exsen{Nag wiya agezän.}      \extran{Friend, come out.}}    \example{
      \exsen{Bogo ma ik atta gogezänän wälläng e dindugän.}      \extran{He got out from inside and ran into the bush.}}    \example{
      \exsen{Kandärmang, yu da daollemae llamda da bälle gazenma da ddone mullae gogon.}      \extran{Sadly, the fire got closer to the old man and he didn't have a way he could escape.}}}  \sense{\textbf{2.} 
    \definition{to rise, come out; shine}    \example{
      \exsen{Yäbäd a mermer abal gogezänän.}      \extran{The sun came out, shining nicely.}}    \example{
      \exsen{Kok a gazen allan.}      \extran{The moon is rising.}}}  \sense{\textbf{3.} 
    \definition{to take out}    \example{
      \exsen{Käsre tos de dägazen nyäng ik att de a batri käp de dazer.}      \extran{Then I took the torch out from the bag and put in the batteries.}}    \example{
      \exsen{Aeya obom bigezänän?}      \extran{Who will come take him out?}}}
  \variants{\SpellingVariant \vartext{gajen}}}

\entry{gazibra}{\headword{gazibra}  \lexeme{gazibra}  \pronnote{ɡazibra}
  \pos{Noun}  \sense{
    \definition{type of edible water snake with skin like sandpaper}    \example{
      \exsen{Alläp de era gazibra ttoe alle ttoe amallo.}      \extran{Drums are skinned with gazibra snake skin.}}}}

\entry{gaall}{\headword{gaall}
  \variantof{SpellingVariant of \vartext{gall}}}

\entry{ge}{\headword{ge}  \pronnote{ɡe}
  \pos{Nominal demonstrative}  \sense{\textbf{1.} 
    \definition{this (proximal demonstrative and pronoun)}    \example{
      \exsen{Ge wätät a ddobae moko dan.}      \extran{This food is delicious.}}    \example{
      \exsen{Ge obo mälla dan.}      \extran{This is his wife.}}}  \sense{\textbf{2.} 
    \definition{like this; for this long}    \example{
      \exsen{Dagirne ge-e, dirom a daeya känyärdae gongttägän.}      \extran{I stayed there for a long time, the cassowary came up quietly.}}}}

\entry{geagell}{\headword{geagell}  \pronnote{ɡeaɡeɽ}
  \pos{Noun}  \sense{
    \definition{Lewin's rail}    \scientificname{Lewinia pectoralis}}}

\entry{geawe}{\headword{geawe}  \note{Trees}  \pronnote{ɡeawe}
  \pos{Noun}  \sense{
    \definition{type of big tree that grows along creeks with bright purple flowers and wood used for canoes}    \note{Trees}    \example{
}    \example{
}    \example{
}}}

\entry{geddäll}{\headword{geddäll}
  \variantof{Dialectal Variant of \vartext{giddoll}}}

\entry{Geleli}{\headword{Geleli}
  \pos{Proper noun}  \sense{
    \definition{Galilee}}}

\entry{gem}{\headword{gem}
  \pos{Nominal demonstrative}  \sense{
    \definition{accusative form of ge}    \example{
      \exsen{Gem de nɨpnae.}      \extran{You will turn this around.}}    \example{
      \exsen{Ngäna gem eka de ttättle anggan.}      \extran{I correct these words.}}}}

\entry{gemae}{\headword{gemae}
  \pos{Nominal demonstrative}  \sense{
    \definition{perlative form of ge}}}

\entry{geme}{\headword{geme}
  \pos{Nominal demonstrative}  \sense{
    \definition{locative form of ge}}}

\entry{Gene}{\headword{Gene}  \pronnote{ɡene}
  \pos{Proper noun}  \sense{
    \definition{male personal name}}}

\entry{Geoff}{\headword{Geoff}
  \pos{Proper noun}  \sense{
    \definition{male personal name}}}

\entry{Georgina}{\headword{Georgina}  \pronnote{ɡeorɡina}
  \pos{Proper noun}  \sense{
    \definition{female personal name}}
  \variants{\SpellingVariant \vartext{Gogina}}}

\entry{ger}{\headword{ger}  \note{Trees}  \pronnote{ɡer}
  \pos{Noun}  \sense{
    \definition{type of big tree that grows in the bush}    \note{Trees}    \example{
}    \example{
}    \example{
}}}

\entry{Geser}{\headword{Geser}  \pronnote{ɡeser}
  \pos{Proper noun}  \sense{
    \definition{male personal name}}}

\entry{=geya}{\headword{=geya}
  \variantof{freevarlab of \vartext{=gaeya}}}

\entry{geyako}{\headword{geyako}
  \pos{Nominal demonstrative}  \sense{
}}

\entry{=geyo}{\headword{=geyo}
  \variantof{freevarlab of \vartext{=gaeyo}}}

\entry{gi}{\headword{gi}  \pronnote{ɡi}
  \pos{Noun}  \sense{
    \definition{grease, fat}    \example{
      \exsen{Ddia da däbe ddobae gi daeya.}      \extran{That deer was fatty.}}}}

\entry{Gibson}{\headword{Gibson}  \pronnote{ɡibson}
  \pos{Proper noun}  \sense{
    \definition{male personal name}}}

\entry{giddoll}{\headword{giddoll}  \pronnote{ɡiɖ͡ʐoɽ}
  \pos{Noun}  \sense{\textbf{1.} 
    \definition{life}    \example{
      \exsen{Oba giddoll makäp me ubi ttoen de mermerangae dangesnegneyo, wätät yagnen a,}      \extran{In their life, they did all things properly.}}    \example{
      \exsen{Bäne giddoll a mɨnyi ada bogon ddobae llayaba za gämällangag.}      \extran{Your life will be like this, stealing things from people.}}}  \sense{\textbf{2.} 
    \definition{to live, reside}    \example{
      \exsen{Ttongo ause bo bin a Madima, Kinkin ttängäm me giddollag daeya.}      \extran{An old woman, her name is Madima, was living in Kinkin village.}}    \example{
      \exsen{Män a Daru me dagirnän.}      \extran{The girl lived in Daru.}}}  \sense{\textbf{3.} 
    \definition{to stay, remain}    \example{
      \exsen{Bongo polle ddäg me nagirne, ngäna do amne we balle.}      \extran{You stay outside the fence; I will go in the middle.}}    \example{
      \exsen{Ende eka da ddone budabän be bagirnän enanae enanae.}      \extran{The Ende language will not be lost, but will remain forever.}}}
  \variants{\Dialectal Variant \vartext{geddäll}}}

\entry{giddollag}{\headword{giddollag}  \note{Birth}  \pronnote{ɡiɖ͡ʐoɽaɡ}
  \pos{Modifier}  \sense{
    \definition{menstruating, on one's period}    \note{Birth}    \example{
}    \example{
}}}

\entry{giddollmeny}{\headword{giddollmeny}
  \pos{Modifier}  \sense{
    \definition{in menopause}}}

\entry{Gidra}{\headword{Gidra}  \note{Language Names}  \pronnote{ɡidra}
  \pos{Proper noun}  \sense{
    \definition{(possibly derogatory) Wipi language}    \note{Language Names}    \example{
}    \example{
}    \example{
}}}

\entry{gidre}{\headword{gidre}
  \pos{Noun}  \sense{
    \definition{enemy}}}

\entry{Gidu}{\headword{Gidu}  \pronnote{ɡidu}
  \pos{Proper noun}  \sense{
    \definition{male personal name}}}

\entry{giegier}{\headword{giegier}  \pronnote{ɡieɡier}
  \pos{Noun}  \sense{
    \definition{white-browed crake}    \scientificname{Poliolimnas cinereus}}}

\entry{Gilbet}{\headword{Gilbet}
  \pos{Proper noun}  \sense{
    \definition{male personal name}}}

\entry{gilib}{\headword{gilib}  \note{Birds}  \pronnote{ɡilib}
  \pos{Noun}  \sense{
    \definition{type of bird}    \note{Birds}    \example{
}    \example{
}    \example{
}}}

\entry{Gimaga}{\headword{Gimaga}
  \pos{Proper noun}  \sense{
    \definition{personal name}}}

\entry{Gina}{\headword{Gina}
  \variantof{SpellingVariant of \vartext{Zina}}  \pronnote{ɡina}}

\entry{Ginarang}{\headword{Ginarang}
  \pos{Proper noun}  \sense{
    \definition{male personal name (the original Ende man [see SE_PN022])}}}

\entry{Gini}{\headword{Gini}
  \pos{Proper noun}  \sense{
    \definition{female personal name}}}

\entry{Ginia}{\headword{Ginia}  \pronnote{ɡinia}
  \pos{Proper noun}  \sense{
    \definition{male personal name}}}

\entry{Giniya}{\headword{Giniya}  \pronnote{ɡinija}
  \pos{Proper noun}  \sense{
    \definition{male personal name}}}

\entry{girag}{\headword{girag}  \note{Animal names}  \pronnote{ɡiraɡ}
  \pos{Noun}  \sense{
    \definition{long-nosed echymipera}    \scientificname{Echymipera rufescens}    \note{Animal names}    \example{
}    \example{
}    \example{
}}}

\entry{girag dirindi}{\headword{girag dirindi}
  \pos{Noun}  \sense{
    \definition{type of yam}}}

\entry{giragirag}{\headword{giragirag}  \note{Trees}  \pronnote{ɡiraɡiraɡ}
  \pos{Noun}  \sense{
    \definition{type of tree}    \note{Trees}    \example{
}    \example{
}    \example{
}}}

\entry{giri}{\headword{giri}  \pronnote{ɡiri}
  \pos{Noun}  \sense{
    \definition{knife}    \example{
      \exsen{Ge lla da nyäng de giri alle papoe eran.}      \extran{This man is piercing the bag with a knife.}}    \example{
      \exsen{Ngäna giri de dipirngän.}      \extran{I took out the knife.}}}}

\entry{giri busa}{\headword{giri busa}
  \pos{Noun}  \sense{
    \definition{sword}}}

\entry{giritai}{\headword{giritai}  \note{Food}  \pronnote{ɡiritai}
  \pos{Noun}  \sense{
    \definition{type of long yam with a white interior and hairs}    \note{Food}    \example{
}    \example{
}    \example{
}}}

\entry{giriwak}{\headword{giriwak}  \note{Tools}  \pronnote{ɡiriwak}
  \pos{Noun}  \sense{
    \definition{type of tool}    \note{Tools}    \example{
}    \example{
}    \example{
}}}

\entry{giro}{\headword{giro}  \pronnote{ɡiro}
  \pos{Noun}  \sense{
    \definition{chicken pox}}}

\entry{gita}{\headword{gita}
  \pos{Noun}  \sense{
    \definition{guitar}}}

\entry{giwi}{\headword{giwi}  \note{Birds}  \pronnote{ɡiwi}
  \pos{Noun}  \sense{
    \definition{fruit dove (coroneted, orange-bellied, pink-spotted, orange-fronted)}    \scientificname{Ptilinopus coronulatus, P. iozonus, P. perlatus; P. aurantiifrons}    \note{Birds}    \example{
}    \example{
}    \example{
}}}

\entry{Giwo}{\headword{Giwo}  \pronnote{ɡiwo}
  \pos{Proper noun}  \sense{
    \definition{male personal name}}}

\entry{=giyu}{\headword{=giyu}
  \variantof{Unspecified Variant of \vartext{=gaeyo}}}

\entry{gɨg}{\headword{gɨg}  \pronnote{ɡɪɡ}
  \pos{Transitive verb}  \sense{
    \definition{to collect ants}}}

\entry{gɨlgɨl}{\headword{gɨlgɨl}
  \pos{Verb}  \sense{
    \definition{descend}}}

\entry{gɨngg}{\headword{gɨngg}  \pronnote{ɡɪŋɡ}
  \pos{Transitive verb}  \sense{
    \definition{to ambush, gang up on, attack in a large group}    \example{
      \exsen{Nogat bom ubony a ere daeränggeyo, gudae enanae dɨnggɨgeyo.}      \extran{When bees started biting Nogat, they really attacked him.}}}}

\entry{Gladis}{\headword{Gladis}
  \pos{Proper noun}  \sense{
    \definition{female personal name}}
  \variants{\SpellingVariant \vartext{Gledis}}}

\entry{Glan}{\headword{Glan}
  \pos{Proper noun}  \sense{
    \definition{female personal name}}}

\entry{glas}{\headword{glas}
  \pos{Noun}  \sense{
    \definition{glass}}}

\entry{gleb}{\headword{gleb}  \pronnote{ɡleb}
  \pos{Transitive verb}  \sense{
    \definition{to take, steal}    \example{
      \exsen{Aya ngämo up näglebän?}      \extran{Who took my banana?}}}}

\entry{Gledis}{\headword{Gledis}
  \variantof{SpellingVariant of \vartext{Gladis}}}

\entry{Glendis}{\headword{Glendis}
  \pos{Proper noun}  \sense{
    \definition{female personal name}}}

\entry{gllae}{\headword{gllae}  \pronnote{ɡɽae}
  \pos{Transitive verb}  \sense{\textbf{1.} 
    \definition{to float, drift; paddle}    \example{
      \exsen{Net a gllaenen allan.}      \extran{The net is floating.}}    \example{
      \exsen{Ngämi (excl.) angde gogllaenalla, däbe za da dibeya ka ili därän.}      \extran{When we paddled, that thing almost disappeared.}}}  \sense{\textbf{2.} 
    \definition{to paddle; pedal}    \example{
      \exsen{Ngäna mɨnyi gall de bogllae.}      \extran{I will paddle the canoe.}}    \example{
      \exsen{Ngämi kili gogman, kili peyang gall de dagllaeya do Upiara dälltam.}      \extran{We (excl.) were so happy, with joy we paddled the canoe until we arrived at Upiara.}}    \example{
      \exsen{Ngäna gall de mängamängall dagllae ddia koenmäll e, gall alle gumbiebmeny.}      \extran{I paddled the canoe quickly to chase the deer with the boat.}}}  \sense{\textbf{3.} 
    \definition{to dig, spade}    \example{
      \exsen{Gällae de dängkamalle, dagllaealle dagllaealle.}      \extran{They started digging; they dug and dug.}}}}

\entry{gllae}{\headword{gllae}  \note{This is an inflected form.}  \pronnote{ɡɽae}
  \pos{Intransitive verb}  \sense{
    \definition{to shine}    \note{This is an inflected form.}    \example{
      \exsen{Yäbäd a mermerae gogllaewän.}      \extran{The sun shined nicely.}}    \example{
      \exsen{Ngäna angde gongllae ddäg e, ge indrang a ngämo dowae me gogllayän.}}}}

\entry{gllaengg}{\headword{gllaengg}
  \pos{}  \sense{
}}

\entry{gllaglle}{\headword{gllaglle}  \pronnote{ɡɽaɡɽe}
  \pos{Transitive verb}  \sense{
    \definition{to skin, remove skin}    \example{
      \exsen{Ngäna ttall ttoe de nägllewan.}      \extran{I removed the wallaby skin.}}}}

\entry{gllangglla}{\headword{gllangglla}  \note{Things a person can do}  \pronnote{ɡɽaŋɡɽa}
  \pos{Intransitive verb}  \sense{
    \definition{to swim}    \note{Things a person can do}    \example{
      \exsen{Gällanggälla eran.}      \extran{He is swimming.}}    \example{
      \exsen{Tätam ngämi bun i dängallaebeya, ako mutt i dingllaebeya.}      \extran{Yesterday, we (excl.) swam down the river and then back up the river.}}}
  \variants{\SpellingVariant \vartext{gällanggälla}}}

\entry{glle}{\headword{glle}  \pronnote{ɡɽe}
  \pos{Noun}  \sense{
    \definition{type of tree with small edible fruit and wood used to make canoes}}}

\entry{gllo}{\headword{gllo}
  \pos{Transitive verb}  \sense{
    \definition{to take out, remove}    \example{
      \exsen{Däräng källa de ma ngätt att de baglloaebnalla.}      \extran{We will take dog poop out from the yards.}}}}

\entry{gllogllo}{\headword{gllogllo}  \pronnote{ɡɽoɡɽo}
  \pos{Noun}  \sense{
    \definition{marbled frogmouth}    \scientificname{Podargus ocellatu}    \example{
}    \example{
}    \example{
}}}

\entry{Gllu}{\headword{Gllu}
  \pos{Proper noun}  \sense{
    \definition{Gllu (toponym)}}}

\entry{gllu}{\headword{gllu}  \pronnote{ɡɽu}
  \pos{Ideophone}  \sense{
    \definition{splash}}}

\entry{Gloria}{\headword{Gloria}
  \pos{Proper noun}  \sense{
    \definition{female personal name}}}

\entry{go}{\headword{go}
  \variantof{freevarlab of \vartext{English loanword}}}

\entry{go~}{\headword{go~}
  \variantof{Infinitival reduplicant of \vartext{RED~}}}

\entry{go}{\headword{go}  \pronnote{ɡo}
  \pos{Noun}  \sense{
    \definition{drain}    \example{
      \exsen{Mälla da go bangesnegnän ine da eramae bizmollnän.}      \extran{The women will make drains so that the water will run.}}    \example{
}    \example{
}}}

\entry{Goballwang}{\headword{Goballwang}
  \pos{Proper noun}  \sense{
    \definition{Goballwang (camping, hunting, and garden place in Karama swamp between Upiara and Limol)}}}

\entry{god}{\headword{god}  \note{Trees}  \pronnote{ɡod}
  \pos{Noun}  \sense{
    \definition{type of cultivated fruit tree similar to dimes tree}    \note{Trees}    \example{
}    \example{
}    \example{
}}}

\entry{Godd}{\headword{Godd}  \pronnote{ɡoɖ͡ʐ}
  \pos{Proper noun}  \sense{
    \definition{God}    \example{
      \exsen{Godd ibim deyangesnegän.}      \extran{God created us (incl.).}}}}

\entry{goeg}{\headword{goeg}  \pronnote{ɡoeɡ}
  \pos{Noun}  \sense{
    \definition{old garden that is ready to be cleared again}    \example{
      \exsen{goeg kuddäll}      \extran{old garden}}    \example{
      \exsen{Lla da goeg melem anggan.}      \extran{People do garden work.}}}}

\entry{Goeg wälläng}{\headword{Goeg wälläng}  \note{Places around Limol}  \pronnote{ɡoeɡ wəɽəŋ}
  \pos{Proper noun}  \sense{
    \definition{garden place in the bush on the west side of the road to Taolang}    \note{Places around Limol}    \example{
}    \example{
}    \example{
}}}

\entry{gogäle}{\headword{gogäle}  \pronnote{ɡoɡəle}
  \pos{Noun}  \sense{
    \definition{noise}    \example{
      \exsen{Lla llɨg a dandärmom mänmän abaene gogäle de.}      \extran{The boys heard the noise from the girls.}}}}

\entry{Goge}{\headword{Goge}  \pronnote{ɡoɡe}
  \pos{Proper noun}  \sense{
    \definition{male personal name}}}

\entry{Gogina}{\headword{Gogina}
  \variantof{SpellingVariant of \vartext{Georgina}}}

\entry{gogo}{\headword{gogo}  \pronnote{ɡoɡo}
  \pos{Noun}  \sense{
    \definition{varieties of palms with coconuts with a dark green exocarp}}}

\entry{Gogo}{\headword{Gogo}
  \variantof{Fast Speech Variant of \vartext{Gogodala}}}

\entry{gogo}{\headword{gogo}  \pronnote{ɡoɡo}
  \pos{Transitive verb}  \sense{
    \definition{to build}    \example{
      \exsen{Ma de gogo eran.}      \extran{He is building the house.}}    \example{
      \exsen{Nogowallo.}      \extran{They built it.}}    \example{
      \exsen{Ttongo täräp me ngäna ngämo ma gogo we gotäbane.}      \extran{One time, I planned to build my house.}}    \example{
      \exsen{Abo bongo ttongo ma sisor de nogo.}      \extran{You must build a new house.}}}}

\entry{gogo kottllam}{\headword{gogo kottllam}  \pronnote{ɡoɡo koʈ͡ʂɽam}
  \pos{Noun}  \sense{
    \definition{type of turtle with yellow scales on neck}}}

\entry{Gogodala}{\headword{Gogodala}  \note{Language Names}  \pronnote{ɡoɡodala}
  \pos{Proper noun}  \sense{
    \definition{Gogodala language}    \note{Language Names}    \example{
}    \example{
}    \example{
}}
  \variants{\Fast Speech Variant \vartext{Gogo}}
  \variants{\SpellingVariant \vartext{Gogodara}}}

\entry{Gogodara}{\headword{Gogodara}
  \variantof{SpellingVariant of \vartext{Gogodala}}}

\entry{gogodd}{\headword{gogodd}  \note{Parts of a bird}  \pronnote{ɡoɡoɖ͡ʐ}
  \pos{Noun}  \sense{\textbf{1.} 
    \definition{spleen}    \note{Parts of a bird}    \example{
      \exsen{Pa bo gogodd a malla ddäddäg ma dan.}      \extran{The bird's spleen is not edible.}}    \example{
}    \example{
}}  \sense{\textbf{2.} 
    \definition{bladder}    \note{Parts of a bird}}}

\entry{gogodd pänddäg}{\headword{gogodd pänddäg}  \note{Birth}  \pronnote{ɡoɡoɖ͡ʐ pənɖ͡ʐəɡ}
  \pos{}  \sense{
    \definition{water breaking}    \note{Birth}    \example{
}    \example{
}    \example{
}}}

\entry{gol}{\headword{gol}  \pronnote{ɡol}
  \pos{Noun}  \sense{
    \definition{goal}}}

\entry{golgol}{\headword{golgol}  \note{Trees}  \pronnote{ɡolɡol}
  \pos{Noun}  \sense{
    \definition{type of tree that grows in the bush with a straight trunk and wood used for house sticks}    \note{Trees}    \example{
}    \example{
}    \example{
}}}

\entry{gollab}{\headword{gollab}  \pronnote{ɡoɽab}
  \pos{Intransitive verb}  \sense{\textbf{1.} 
    \definition{to pour}    \example{
      \exsen{Ibi kumuddäga ine kutt de dogollaebaebeya.}      \extran{We (incl.) poured out three water buckets.}}    \example{
      \exsen{Nogollaeb.}      \extran{[You] pour it.}}}  \sense{\textbf{2.} 
}}

\entry{gollaeb}{\headword{gollaeb}  \pronnote{ɡoɽaeb}
  \pos{Quantifier}  \sense{
    \definition{many}    \example{
      \exsen{Gaem ai tater gollaeb de inen anggan.}      \extran{Gaem is weaving many good mats.}}}}

\entry{gollaembäg}{\headword{gollaembäg}  \pronnote{ɡoɽaembəɡ}
  \pos{Transitive verb}  \sense{
    \definition{applicative form of gollab}    \example{
      \exsen{Mänmän mängalae ine walle danggollaebägeyo.}      \extran{The girls quickly poured water on someone.}}}}

\entry{gollob}{\headword{gollob}  \pronnote{ɡoɽob}
  \pos{Noun}  \sense{
    \definition{outer layer, hull, shell (e.g. of a turtle, egg)}    \example{
      \exsen{wit gollob}      \extran{wheat hull}}    \example{
      \exsen{Lla da päre abo gobällan kottllam bo gollob pallkepallke de dänglläbeyo.}      \extran{Then the people sadly went to the turtle and collected all the pieces of his shell.}}}}

\entry{golloll}{\headword{golloll}  \pronnote{ɡoɽoɽ}
  \pos{Locational}  \sense{
    \definition{back, behind}    \example{
      \exsen{Llamda da gudae dallän polle we gozenän a up golloll me gotogolän.}      \extran{The old man went through the fence early in the morning and hid behind the banana tree.}}}}

\entry{gollolla}{\headword{gollolla}  \note{Palms}  \pronnote{ɡoɽoɽa}
  \pos{Noun}  \sense{
    \definition{type of palm}    \note{Palms}    \example{
}    \example{
}    \example{
}}}

\entry{gomoe}{\headword{gomoe}
  \pos{Transitive verb}  \sense{
    \definition{to make a mistake}    \example{
      \exsen{Ngäna gem eka de ttättle anggan, ge ddob lla da erem gomoenen amallo.}      \extran{I correct these words, which other people mispronounce.}}}}

\entry{gonagone}{\headword{gonagone}  \pronnote{ɡonaɡone}
  \pos{Transitive verb}  \sense{\textbf{1.} 
    \definition{to cook}    \example{
      \exsen{Bongo nalle, otät e nognen.}      \extran{You, go and cook the food.}}}  \sense{\textbf{2.} 
    \definition{to burn}    \example{
      \exsen{Llo mit de dugonaeballo.}      \extran{They burned the tree bases.}}}  \sense{\textbf{3.} 
    \definition{to heat}    \example{
      \exsen{Ngäna ine ttämttämang yu dägag a ngämo ttäle de dogonaeya.}      \extran{I boiled water and heated my leg.}}}}

\entry{gonangg}{\headword{gonangg}
  \pos{Transitive verb}  \sense{
    \definition{to cook}}}

\entry{gone}{\headword{gone}
  \pos{Verb}  \sense{
    \definition{come}}}

\entry{gongg}{\headword{gongg}  \pronnote{ɡoŋɡ}
  \pos{Intransitive verb}  \sense{
    \definition{to disturb a bees nest and then to feel the bites}}}

\entry{gongglem}{\headword{gongglem}  \pronnote{ɡoŋɡlem}
  \pos{Noun}  \sense{\textbf{1.} 
    \definition{immature coconut that is partially solid inside}    \example{
      \exsen{Gongglem daeben yäbäd ttänttämang me nane ma da.}      \extran{Immature coconuts are the best when drunk in the hot sun.}}}  \sense{\textbf{2.} 
    \definition{the eighth stage of coconut growth during which the endosperm of the fruit is solidifying}}}

\entry{gonggo}{\headword{gonggo}  \pronnote{ɡoŋɡo}
  \pos{Noun}  \sense{
    \definition{piping bellbird/crested pitohui}    \scientificname{Ornorectes cristatus}    \example{
}    \example{
}    \example{
}}}

\entry{gonya}{\headword{gonya}
  \variantof{Dialectal Variant of \vartext{gänya}}}

\entry{gonyeya}{\headword{gonyeya}
  \variantof{freevarlab of \vartext{gänyaeya}}}

\entry{gonyo}{\headword{gonyo}
  \variantof{Dialectal Variant of \vartext{gänya}}  \pronnote{ɡoɲo}
  \variants{\SpellingVariant \vartext{gänyo}}}

\entry{gonz}{\headword{gonz}  \note{Water}  \pronnote{ɡonz}
  \pos{Noun}  \sense{
    \definition{reed}    \note{Water}    \example{
      \exsen{Buk nyäng a gonz alle inenatt a mer abal dan.}      \extran{Gonz reed is very good to weave book bags with.}}    \example{
}    \example{
}}}

\entry{gonzagonzar}{\headword{gonzagonzar}  \note{Trees}  \pronnote{ɡonzaɡonzar}
  \pos{Noun}  \sense{
    \definition{type of tree}    \note{Trees}    \example{
}    \example{
}    \example{
}}}

\entry{gonzagonzar budar}{\headword{gonzagonzar budar}  \note{Grubs}  \pronnote{ɡonzaɡonzar budar}
  \pos{Noun}  \sense{
    \definition{type of grub}    \note{Grubs}    \example{
}    \example{
}    \example{
}}}

\entry{Gonzer}{\headword{Gonzer}
  \pos{Proper noun}  \sense{
    \definition{female personal name}}}

\entry{gora}{\headword{gora}  \note{Musical instruments}  \pronnote{ɡora}
  \pos{Noun}  \sense{
    \definition{rattle}    \note{Musical instruments}    \example{
}    \example{
}    \example{
}}}

\entry{goral}{\headword{goral}  \note{Trees}  \pronnote{ɡoral}
  \pos{Noun}  \sense{
    \definition{type of tree}    \note{Trees}    \example{
}    \example{
}    \example{
}}}

\entry{goro}{\headword{goro}
  \pos{Modifier}  \sense{
    \definition{lush, overgrown, wild}    \example{
}}}

\entry{goro wälläng}{\headword{goro wälläng}
  \pos{Noun}  \sense{
    \definition{jungle}}}

\entry{got}{\headword{got}
  \variantof{Dialectal Variant of \vartext{English loanword}}}

\entry{government}{\headword{government}
  \variantof{Dialectal Variant of \vartext{English loanword}}}

\entry{Grace}{\headword{Grace}  \pronnote{ɡres}
  \pos{Proper noun}  \sense{
    \definition{female personal name}}
  \variants{\SpellingVariant \vartext{Greis}}}

\entry{grade}{\headword{grade}
  \variantof{Dialectal Variant of \vartext{English loanword}}}

\entry{grawa}{\headword{grawa}  \pronnote{ɡrawa}
  \pos{Noun}  \sense{\textbf{1.} 
    \definition{Australasian darter}    \scientificname{Anhinga novaehollandiae}    \example{
}    \example{
}    \example{
}}  \sense{\textbf{2.} 
    \definition{little corella}    \scientificname{Cacatua sanguinea}}}

\entry{greasy}{\headword{greasy}
  \variantof{freevarlab of \vartext{English loanword}}}

\entry{greid}{\headword{greid}
  \pos{Noun}  \sense{
    \definition{grade}    \example{
      \exsen{greid siks me}      \extran{in grade six}}}}

\entry{Greis}{\headword{Greis}
  \variantof{SpellingVariant of \vartext{Grace}}}

\entry{group}{\headword{group}
  \variantof{Dialectal Variant of \vartext{English loanword}}}

\entry{grup}{\headword{grup}
  \pos{Noun}  \sense{
    \definition{group}}}

\entry{gu~}{\headword{gu~}
  \variantof{Inflected Form of \vartext{RED~}}}

\entry{Guar}{\headword{Guar}  \pronnote{ɡuar}
  \pos{Proper noun}  \sense{
    \definition{male personal name}}}

\entry{Gubam}{\headword{Gubam}
  \pos{Proper noun}  \sense{
    \definition{Gubam (in Morehead Rural LLG)}}}

\entry{gubare}{\headword{gubare}  \pronnote{ɡubare}
  \pos{Noun}  \sense{
    \definition{crossbeam}}}

\entry{guboll}{\headword{guboll}  \note{Birds}  \pronnote{ɡuboɽ}
  \pos{Noun}  \sense{
    \definition{New Guinean magpie}    \scientificname{Gymnorhina tibicen papuana}    \note{Birds}    \example{
}    \example{
}    \example{
}}}

\entry{gudae}{\headword{gudae}  \pronnote{ɡudae}
  \pos{Quantifier}  \sense{
    \definition{plenty}    \example{
      \exsen{Wayati da duli ge ttängäm gudae ada enddäna gogon wayati dae gogon.}      \extran{There were so many watermelon there, this place was like a clearing of only watermelon.}}}}

\entry{gudae}{\headword{gudae}  \pronnote{ɡudae}
  \pos{Modifier}  \sense{\textbf{1.} 
    \definition{early morning}    \example{
      \exsen{Gudae ag särem abal me ngäma ngämi otat de mängalae yu dägaebeya.}      \extran{In the darkness of early morning, we (excl.) quickly cooked our food.}}    \example{
      \exsen{Bogo gudae abal angällbänan.}      \extran{He woke up very early in the morning.}}}  \sense{\textbf{2.} 
    \definition{(the) past, before}    \example{
      \exsen{Ngäma ddob gudae att gagäll de dämbälaebnalla.}      \extran{Some of us (excl.) were remembering bad things from the past.}}}  \sense{\textbf{3.} 
    \definition{earlier, previously, in the past, before}    \example{
      \exsen{gudae gudae täräp me}      \extran{long, long ago}}    \example{
      \exsen{Komlla nag a dadegwaeya ami gudae deyagirneyo walle ddage menae me.}      \extran{There were two friends who used to live on the side of the stream.}}    \example{
      \exsen{Dedme mɨnyi bibi obom ikop nägaeyo ada ingollang bogo alla umllang deyagnegän bibim gudae.}      \extran{You (pl.) will see him there, just like how he told you before.}}}}

\entry{gudne}{\headword{gudne}  \note{Adjectives}  \pronnote{ɡudne}
  \pos{Modifier}  \sense{\textbf{1.} 
    \definition{old}    \note{Adjectives}    \example{
      \exsen{gudne lla}      \extran{old man}}    \example{
      \exsen{Ge ma da gudne dan.}      \extran{This house is old.}}    \example{
}}  \sense{\textbf{2.} 
    \definition{long ago}    \note{Adjectives}    \example{
      \exsen{Däräng a gudne ulle gänyaolle ngäsangngäsang allan.}      \extran{The dog grew up here a long time ago and is always returning.}}}}

\entry{guem}{\headword{guem}  \note{Water}  \pronnote{ɡuem}
  \pos{Noun}  \sense{
    \definition{channel; deep part of river}    \note{Water}    \example{
}    \example{
}    \example{
}}}

\entry{gugall}{\headword{gugall}  \pronnote{gu'gall}
  \pos{Noun}  \sense{
    \definition{type of plant with red, yellow, and white flowers and fruit with small, round seeds; children use the fruit in bamboo blow guns}}}

\entry{gugi}{\headword{gugi}  \pronnote{ɡuɡi}
  \pos{Intransitive verb}  \sense{
    \definition{to stand}    \example{
      \exsen{Angde Kwalde amne we gogän, bogo obo däräng de dangermällnegän a dugiwän.}}    \example{
      \exsen{Dedme bugimällämän lla da.}      \extran{People will be standing there.}}}}

\entry{gugu}{\headword{gugu}  \note{Types of Taro}  \pronnote{ɡuɡu}
  \pos{Noun}  \sense{
    \definition{type of big taro}    \note{Types of Taro}    \example{
}    \example{
}    \example{
}}}

\entry{gugu}{\headword{gugu}  \note{Types of Taro}  \pronnote{ɡuɡu}
  \pos{Noun}  \sense{
    \definition{row of leaves going up the roof}    \note{Types of Taro}}}

\entry{Gugu}{\headword{Gugu}
  \pos{Proper noun}  \sense{
    \definition{Gugu (toponym)}}}

\entry{Gugu Gel}{\headword{Gugu Gel}  \pronnote{ɡuɡu ɡel}
  \pos{Proper noun}  \sense{
    \definition{female personal name}}}

\entry{Guim}{\headword{Guim}
  \pos{Proper noun}  \sense{
    \definition{personal name}}}

\entry{Guinea}{\headword{Guinea}
  \variantof{Dialectal Variant of \vartext{English loanword}}}

\entry{gul}{\headword{gul}  \pronnote{ɡul}
  \pos{Noun}  \sense{
    \definition{crowd, group, mob; school (of fish)}    \example{
      \exsen{Bogo lla gul ulle de ikop dägagän.}      \extran{He saw the big mob of people.}}    \example{
      \exsen{Ngämo yae a ngämo baba gogezäneyo Ende gul att.}      \extran{My mother and my father came from the Ende crowd.}}}}

\entry{gulag}{\headword{gulag}  \pronnote{ɡulaɡ}
  \pos{Transitive coverb}  \sense{\textbf{1.} 
    \definition{to accompany}    \example{
      \exsen{Ngäna Senti bom umllang däga ngämlle gulag e.}      \extran{I told Senti to accompany me.}}}  \sense{\textbf{2.} 
}}

\entry{guli}{\headword{guli}  \note{Trees}  \pronnote{ɡuli}
  \pos{Noun}  \sense{
    \definition{type of tree}    \note{Trees}    \example{
}    \example{
}    \example{
}}}

\entry{gulin}{\headword{gulin}
  \pos{Noun}  \sense{
}}

\entry{gull}{\headword{gull}  \lexeme{gull}  \pronnote{ɡuɽ}
  \pos{Noun}  \sense{
    \definition{net}    \example{
      \exsen{Mälla da kemibi kollba de nägäddallo gull alle.}      \extran{The women caught many fish with the net.}}    \example{
      \exsen{Angde ngäna gull de ik i däga, käza da gull ik i gozenän.}      \extran{When I lifted the net up, the crocodile went inside it.}}}}

\entry{gulla}{\headword{gulla}
  \pos{}  \sense{
    \definition{push}}}

\entry{gullabgullab}{\headword{gullabgullab}
  \pos{Modifier}  \sense{
    \definition{fat}    \example{
      \exsen{Ngäna sɨmell de yagnen anggan, käm gullabgullab a.}      \extran{I'm looking for a pig, one with a fat belly.}}}}

\entry{gullba}{\headword{gullba}  \pronnote{ɡuɽba}
  \pos{Noun}  \sense{
    \definition{bundle of sago}}}

\entry{gullbe}{\headword{gullbe}  \pronnote{ɡuɽbe}
  \pos{Kinship noun}  \sense{\textbf{1.} 
    \definition{husband}    \example{
      \exsen{gullbe da walle}      \extran{couple}}    \example{
      \exsen{Bäne gullbe da daden?}      \extran{Do you have a husband (lit. does your husband exist)?}}    \example{
      \exsen{Gullbeda däban.}      \extran{That's her husband.}}}  \sense{\textbf{2.} 
    \definition{male animal}    \example{
      \exsen{ddia gullbe, sɨmell gullbe}      \extran{buck, boar}}}  \sense{\textbf{3.} 
    \definition{huge, big}    \example{
      \exsen{Wel gullbe da gongttägän a mägda bo pite de däzinän.}      \extran{A big wind blew and lifted his mother's skirt.}}    \example{
      \exsen{Ngämi yu dägaebnalla angde yogoll gullbe da dipliwän abo.}      \extran{We (excl.) were cooking them when a rainstorm started.}}}
  \variants{\SpellingVariant \vartext{gullbi}}}

\entry{Gullbe Bikme Auma}{\headword{Gullbe Bikme Auma}  \note{Places around Limol}  \pronnote{ɡuɽbe bikme auma}
  \pos{Proper noun}  \sense{
    \definition{sacred place of Dareda (near Karama swamp)}    \note{Places around Limol}    \example{
}    \example{
}    \example{
}}}

\entry{Gullbe bo llädayatt}{\headword{Gullbe bo llädayatt}  \note{Places around Limol}  \pronnote{ɡuɽbe bo ɽədajaʈ͡ʂ}
  \pos{Proper noun}  \sense{
    \definition{sacred place of Dareda (large hill)}    \note{Places around Limol}    \example{
}    \example{
}    \example{
}}}

\entry{Gullbe bo makollamatt}{\headword{Gullbe bo makollamatt}  \note{Places around Limol}  \pronnote{ɡuɽbe bo makoɽamaʈ͡ʂ}
  \pos{Proper noun}  \sense{
    \definition{Dareda's sacred place (hill on the road to Kinkin)}    \note{Places around Limol}    \example{
}    \example{
}    \example{
}}}

\entry{gullbe peyang}{\headword{gullbe peyang}  \pronnote{ɡuɽbe pejaŋ}
  \pos{Modifier}  \sense{
    \definition{married (of a female)}    \example{
      \exsen{Bäne yae auli gullbe peyang daeya?}      \extran{How many husbands does your mother have?}}}}

\entry{gullbi}{\headword{gullbi}
  \variantof{SpellingVariant of \vartext{gullbe}}}

\entry{gullbog}{\headword{gullbog}  \pronnote{ɡuɽboɡ}
  \pos{Modifier}  \sense{
    \definition{married (of a female)}}}

\entry{gullem}{\headword{gullem}  \lexeme{gullem}  \pronnote{ɡuɽem}
  \pos{Noun}  \sense{
    \definition{snake}    \example{
      \exsen{Ge ekaklle ulle me, gullem sapasapang a dadeg.}      \extran{In this world, there are many kinds of snakes.}}    \example{
      \exsen{Gullem llɨg a däbe agezänan tatuma me.}      \extran{A baby snake came out in that washing place.}}}}

\entry{Gullem suwe}{\headword{Gullem suwe}  \note{Places around Limol}  \pronnote{ɡuɽem suwe}
  \pos{Proper noun}  \sense{
    \definition{gardening place; Galo's sacred place (AZ97)}    \note{Places around Limol}    \example{
}    \example{
}    \example{
}}}

\entry{gullem suwetar}{\headword{gullem suwetar}  \pronnote{ɡuɽem suwetar}
  \pos{Noun}  \sense{
    \definition{type of tree that is used as medicine for snake bites and to repel snakes}}}

\entry{gullme}{\headword{gullme}  \lexeme{gullme}  \pronnote{ɡuɽme}
  \pos{Noun}  \sense{
    \definition{type of flowering tree that grows on the riverside in swamps with wood used for firewood and long, hanging red flowers that smell nice}}}

\entry{gullme käpang}{\headword{gullme käpang}  \note{Traditional Arrows}  \pronnote{ɡuɽme kəpaŋ}
  \pos{Noun}  \sense{
    \definition{type of spear}    \note{Traditional Arrows}    \example{
}    \example{
}    \example{
}}}

\entry{gullun}{\headword{gullun}
  \pos{Modifier}  \sense{
    \definition{old}    \example{
      \exsen{Mälla da gänyan, kaptte gullun alle obo bun di omllawan.}      \extran{The wife is here, her head is wrapped with old cloth.}}}}

\entry{gumae}{\headword{gumae}
  \pos{Transitive verb}  \sense{
    \definition{to go around}}}

\entry{gungg}{\headword{gungg}  \pronnote{ɡuŋɡ}
  \pos{Transitive verb}  \sense{
    \definition{to marry a widow}    \example{
      \exsen{Angde Kwalde kuddäll agan, Kidarga ako mosen da bäne mik di nanggugan.}      \extran{When Kwalde died, Kidarga married his older brother's (i.e. Kwalde's) widow.}}}}

\entry{Gurel}{\headword{Gurel}
  \pos{Proper noun}  \sense{
    \definition{male personal name}}}

\entry{guwaba}{\headword{guwaba}  \note{Medicine}  \pronnote{ɡuwaba}
  \pos{Noun}  \sense{
    \definition{guava tree; water steeped with its leaves is used to wash sores}    \scientificname{Psidium guajava}    \note{Medicine}    \example{
}    \example{
}    \example{
}}}

\entry{guwo}{\headword{guwo}  \pronnote{ɡuwo}
  \pos{Noun}  \sense{\textbf{1.} 
    \definition{heart}    \example{
      \exsen{Yesu lla yaba guwo watt gagäll anyke yuwog de däkällttanegnän.}      \extran{Jesus cast many evil spirits out from the hearts of men.}}    \example{
      \exsen{Sali bo guwo watt dirom llɨg de ekaekong digezänän.}      \extran{Out of Sali's chest came a chirping cassowary chick.}}}  \sense{\textbf{2.} 
    \definition{inside}    \example{
      \exsen{Bogo dadiweya gall guwo me dämenang dagirnän.}      \extran{He was sitting in the canoe.}}}}

\entry{guzi}{\headword{guzi}  \lexeme{guzi}  \pronnote{ɡuzi}
  \pos{Noun}  \sense{
    \definition{yabby (type of crayfish)}    \scientificname{Cherax}    \example{
      \exsen{Bogo guzi de daittän.}      \extran{She caught a crayfish.}}}}

\entry{guziguzi}{\headword{guziguzi}  \note{Trees}  \pronnote{ɡuziɡuzi}
  \pos{Noun}  \sense{
    \definition{type of tree}    \note{Trees}    \example{
}    \example{
}    \example{
}}}

\entry{=gwaeya}{\headword{=gwaeya}
  \pos{Verbal enclitic}  \sense{
    \definition{cop.pst.duS}}
  \variants{\freevarlabs \vartext{=gweya, =gwe}}
  \variants{\Fast Speech Variant \vartext{=weya}}}

\entry{gwaga}{\headword{gwaga}  \note{Trees}  \pronnote{ɡwaɡa}
  \pos{Noun}  \sense{\textbf{1.} 
    \definition{type of big tree that grows in the bush with big, edible fruit that are yellow outside, red inside, and stain lips brown; when ripe, it will open and drop the heart-shaped seed}    \note{Trees}    \example{
}    \example{
}    \example{
}}  \sense{\textbf{2.} 
    \definition{the sixth stage of sago growth in which the fruits have matured and the pith is drier}    \note{Trees}    \example{
}}}

\entry{gwaga käp}{\headword{gwaga käp}  \pronnote{ɡwaɡa kəp}
  \pos{Noun}  \sense{
    \definition{fruit of sago palm}    \example{
      \exsen{Baba bo sana da gwaga käpang gogon kikib me.}      \extran{Dad's sago palm is already producing fruit.}}    \example{
      \exsen{Llɨg kälekäle da sana gwaga käp alle gotongoenegnän.}      \extran{The children were playing with the sago fruit.}}}}

\entry{gwälläd}{\headword{gwälläd}  \note{Trees}  \pronnote{ɡwəɽəd}
  \pos{Noun}  \sense{
    \definition{type of tree}    \note{Trees}    \example{
}    \example{
}    \example{
}}}

\entry{=gwän}{\headword{=gwän}
  \variantof{Dialectal Variant of \vartext{=gaeyo}}}

\entry{gwängäm}{\headword{gwängäm}
  \pos{Noun}  \sense{
    \definition{one of the secrets that was practiced by men only, an animal that was killed and eaten only by men inside the festival of the bullroarer, not women or youths}}}

\entry{gwara}{\headword{gwara}  \note{Weather}  \pronnote{ɡwara}
  \pos{Noun}  \sense{
    \definition{lightning}    \note{Weather}    \example{
}    \example{
}    \example{
}}}

\entry{gwazi}{\headword{gwazi}  \note{Types of Taro}  \pronnote{ɡwazi}
  \pos{Noun}  \sense{
    \definition{type of big taro}    \note{Types of Taro}    \example{
}    \example{
}    \example{
}}}

\entry{=gwe}{\headword{=gwe}
  \variantof{freevarlab of \vartext{=gwaeya}}}

\entry{gwell}{\headword{gwell}  \note{School}  \pronnote{ɡweɽ}
  \pos{Noun}  \sense{\textbf{1.} 
    \definition{rule, law}    \note{School}    \example{
      \exsen{Moses adawatta ge gwell de dadräbnegän bibra, bina tikop a llokollokott dag Adi bäne Mer Eka malaem e.}      \extran{Moses wrote these laws for you all, because your hearts are stubborn to obey God's gospel.}}    \example{
}    \example{
}}  \sense{\textbf{2.} 
    \definition{secret}    \note{School}    \example{
      \exsen{Gwell eka de ngäna llɨtɨt eran.}      \extran{I'm telling a secret story.}}}}

\entry{Gwen}{\headword{Gwen}
  \pos{Proper noun}  \sense{
    \definition{female personal name}}}

\entry{=gweya}{\headword{=gweya}
  \variantof{freevarlab of \vartext{=gwaeya}}}

\chapter{H}


\entry{Hannah}{\headword{Hannah}  \pronnote{hannah}
  \pos{Proper noun}  \sense{
    \definition{female personal name}}}

\entry{hedmasta}{\headword{hedmasta}
  \pos{Noun}  \sense{
    \definition{headmaster}}}

\entry{Helen}{\headword{Helen}
  \pos{Proper noun}  \sense{
    \definition{female personal name}}}

\entry{Hiden}{\headword{Hiden}  \pronnote{hiden}
  \pos{Proper noun}  \sense{
    \definition{female personal name}}}

\entry{hundred}{\headword{hundred}
  \variantof{SpellingVariant of \vartext{andred}}}

\chapter{I}


\entry{i-}{\headword{i-}
  \pos{Verb}  \sense{
    \definition{irrealis aspect marker, marks events that have or will not be realized, this includes meanings such as conditional, should, "try to", "don't want to", it co-occurs with the prefix b- and the suffix -alle/-allo.}}}

\entry{i-}{\headword{i-}  \pronnote{i}
  \pos{Verb}  \sense{
    \definition{venitive prefix}    \example{
      \exsen{A atka gogon, [obo obo känae] obo känaebag ingäsmäll.}      \extran{And he said, "Come back tomorrow."}}    \example{
      \exsen{Llɨg a ddone ako dingäsmällän Bundae bälle yagyag.}      \extran{The boys didn't return to Bundae.}}    \example{
      \exsen{Däräng käsre dallän do ma we gogon, a täl pallkoll de diwenyän, käsre [adaka] adaka gänyaolle nagda Baet bo llan de täräpnan anggan.}      \extran{Dog then went to his house, and brought a piece of bamboo. Then, he cut his friend Cuscus's ears.}}    \example{
      \exsen{Dindugän, do täl pallkoll dingällbällnän.}      \extran{He ran and got a piece of bamboo.}}}}

\entry{i}{\headword{i}  \pronnote{i}
  \pos{Transitive verb}  \sense{
    \definition{to weave; interlock}    \example{
      \exsen{Bogo nyäng de i eran.}      \extran{She is weaving the basket.}}    \example{
      \exsen{Tater de aya diwän?}      \extran{Who wove the mat?}}    \example{
      \exsen{Ngämo yae era lla ngänäm polle dinignän.}      \extran{My mother, she used to build sturdy fences like a man.}}}}

\entry{i abal}{\headword{i abal}  \note{Weaving}  \pronnote{i abal}
  \pos{Noun}  \sense{
    \definition{plain weaving pattern}    \note{Weaving}    \example{
}    \example{
}    \example{
}}}

\entry{Iamega}{\headword{Iamega}
  \variantof{SpellingVariant of \vartext{Yamega}}}

\entry{Ib}{\headword{Ib}  \pronnote{ib}
  \pos{Proper noun}  \sense{
    \definition{female personal name}}}

\entry{ib~}{\headword{ib~}
  \variantof{freevarlab of \vartext{RED~}}}

\entry{iba}{\headword{iba}  \pronnote{iba}
  \pos{Personal pronoun}  \sense{
    \definition{genitive form of ibi}    \example{
      \exsen{Iba nag ttoen a alla ingoll utale ballän?}      \extran{How far will our friendship go?}}}}

\entry{ibae}{\headword{ibae}
  \pos{Transitive verb}  \sense{
    \definition{to dig using one's nose}}}

\entry{ibe}{\headword{ibe}  \lexeme{ibe}  \pronnote{ibe}
  \pos{Noun}  \sense{
    \definition{mist; fog}}}

\entry{ibeny}{\headword{ibeny}  \note{Gardens}  \pronnote{ibeɲ}
  \pos{Transitive verb}  \sense{
    \definition{to plant}    \note{Gardens}    \example{
      \exsen{Bogo ibeny eran.}      \extran{He is planting.}}    \example{
      \exsen{Biye de bongo ai dan ibik alle nibe.}      \extran{You can plant taro with the ibik stick.}}    \example{
      \exsen{Ttängäm e gongttägeyo, misdae up daebegaeya ibebatta.}      \extran{They arrived at the garden; only bananas were planted there.}}}}

\entry{Ibeti}{\headword{Ibeti}
  \variantof{SpellingVariant of \vartext{Ibetty}}}

\entry{Ibetty}{\headword{Ibetty}  \pronnote{ibɛti}
  \pos{Proper noun}  \sense{
    \definition{female personal name}}
  \variants{\SpellingVariant \vartext{Ibeti}}}

\entry{ibi}{\headword{ibi}  \pronnote{ibi}
  \pos{Personal pronoun}  \sense{
    \definition{we (first person nonsingular inclusive pronoun, nominative form)}    \example{
      \exsen{Koe, ai dan ibi gänyme beyagernalla.}      \extran{Koe, it's good the two of us will be living here.}}}}

\entry{ibi}{\headword{ibi}
  \variantof{Dialectal Variant of \vartext{RED~}}  \pronnote{ibi}}

\entry{ibiatt}{\headword{ibiatt}  \pronnote{ibiaʈ͡ʂ}
  \pos{Noun}  \sense{
    \definition{footprint}    \example{
      \exsen{Däräng aba mäda Kwalde bäne ibiatt de dongkollmälleyo.}      \extran{The dogs followed their owner Kwalde's footprints.}}}}

\entry{ibik}{\headword{ibik}  \note{Tools}  \pronnote{ibik}
  \pos{Noun}  \sense{
    \definition{digging stick}    \note{Tools}    \example{
      \exsen{Tätäm ibik a ddokddok gogon.}      \extran{Yesterday, the digging stick became blunt.}}    \example{
      \exsen{Biye de bongo ai dan ibik alle nibe.}      \extran{You can plant taro with the ibik stick.}}    \example{
}    \example{
}}}

\entry{Ibikang}{\headword{Ibikang}  \note{Place names}  \pronnote{ibikaŋ}
  \pos{Proper noun}  \sense{
    \definition{Ibikang (Kawam-speaking settlement of Wim; not far from Limol (AY96))}    \note{Place names}    \example{
}    \example{
}    \example{
}}}

\entry{ibim}{\headword{ibim}  \pronnote{ibim}
  \pos{Personal pronoun}  \sense{
    \definition{accusative form of ibi}    \example{
      \exsen{Ai dan, käza da ddone ibim beyaddägän.}      \extran{It's alright, the crocodile won't bite us.}}}}

\entry{ibira}{\headword{ibira}
  \variantof{Unspecified Variant of \vartext{ibra}}}

\entry{ibra}{\headword{ibra}  \pronnote{ibra}
  \pos{Personal pronoun}  \sense{
    \definition{dative form of ibi}    \example{
      \exsen{Bongo nalle sana ma, ngäna tudi ma balle ibra kollba ma.}      \extran{You go for sago, and I will go fishing to get fish for us.}}}
  \variants{\Unspecified Variant \vartext{ibira}}}

\entry{Ibru}{\headword{Ibru}  \pronnote{ibru}
  \pos{Noun}  \sense{
    \definition{Hebrew}}}

\entry{id}{\headword{id}
  \variantof{SpellingVariant of \vartext{yid}}}

\entry{idab}{\headword{idab}
  \pos{Transitive verb}  \sense{
    \definition{to dig out}}}

\entry{idaida}{\headword{idaida}  \note{Games}  \pronnote{idaida}
  \pos{Noun}  \sense{
    \definition{type of game played outside in open space}    \note{Games}    \example{
}    \example{
}    \example{
}}}

\entry{Idan}{\headword{Idan}  \pronnote{idan}
  \pos{Proper noun}  \sense{
    \definition{male personal name}}}

\entry{idän}{\headword{idän}  \pronnote{idən}
  \pos{Transitive verb}  \sense{
    \definition{to pick, harvest; dig up, uproot}    \example{
      \exsen{Ubi mɨnyi mätta de bidaebneyo.}      \extran{They will dig up the yams and bring them over.}}    \example{
      \exsen{Ada wayati da moko alle dädeyo.}      \extran{[You] pick watermelons as you please.}}}}

\entry{idd}{\headword{idd}  \note{Dead things}  \pronnote{iɖ͡ʐ}
  \pos{Noun}  \sense{
    \definition{ghost}    \note{Dead things}    \example{
      \exsen{Ubi ada ka bogo idd daeya.}      \extran{They thought she was a ghost.}}    \example{
}    \example{
}}}

\entry{idd ma}{\headword{idd ma}  \note{Dead things}  \pronnote{iɖ͡ʐ ma}
  \pos{Noun}  \sense{
    \definition{afterlife}    \note{Dead things}    \example{
}    \example{
}    \example{
}}}

\entry{iddi}{\headword{iddi}
  \pos{Noun}  \sense{
}}

\entry{iddnge}{\headword{iddnge}  \pronnote{iɖ͡ʐŋe}
  \pos{Noun}  \sense{
    \definition{type of coconut palm}}}

\entry{iddob}{\headword{iddob}  \note{Parts of the day}  \pronnote{iɖ͡ʐob}
  \pos{Noun}  \sense{
    \definition{night}    \note{Parts of the day}    \example{
      \exsen{Angde iddob gogon, ngäna inu wi dalle.}      \extran{When night came, I went to sleep.}}    \example{
      \exsen{Bäne moko daden ibi kollba ma sisiri beyareya, ako iddob abal e bonttägmällnalla.}      \extran{You want for us to go fishing now, so we will arrive late at night.}}    \example{
      \exsen{Däbe ttall a iddob me guddällän.}      \extran{That wallaby arrived at night.}}}}

\entry{iddob amänong}{\headword{iddob amänong}
  \variantof{SpellingVariant of \vartext{iddob amnong}}  \pronnote{iɖ͡ʐob amənoŋ}}

\entry{iddob amnong}{\headword{iddob amnong}  \pronnote{iɖ͡ʐob amnoŋ}
  \pos{Noun}  \sense{
    \definition{midnight}    \example{
      \exsen{Ttongdae mɨnyi awi alle do iddob amnong a ttongo ako iddob amnong alle do bädab bagirnän.}      \extran{One will stay from evening to midnight, and another will stay from midnight until dawn.}}}
  \variants{\SpellingVariant \vartext{iddob amänong}}}

\entry{iddoiddob}{\headword{iddoiddob}  \note{Trees}  \pronnote{iɖ͡ʐoiɖ͡ʐob}
  \pos{Noun}  \sense{
    \definition{type of tree}    \note{Trees}    \example{
}    \example{
}    \example{
}}}

\entry{iddpo}{\headword{iddpo}  \note{Clothing}  \pronnote{iɖ͡ʐpo}
  \pos{Noun}  \sense{
    \definition{clothing, clothes}    \note{Clothing}    \example{
      \exsen{Ttongo mälla da iddpo de nällpäganegnan.}      \extran{One woman was hanging the clothes.}}    \example{
      \exsen{Obo iddpo da pällämpälläm a gombenmällnegnän.}      \extran{His clothes were white and shining.}}    \example{
}}
  \variants{\SpellingVariant \vartext{ittpo}}}

\entry{ide}{\headword{ide}  \pronnote{ide}
  \pos{Noun}  \sense{
    \definition{side}}}

\entry{Iden}{\headword{Iden}
  \pos{Proper noun}  \sense{
    \definition{Eden}}}

\entry{Idi}{\headword{Idi}  \note{Language Names}  \pronnote{idi}
  \pos{Proper noun}  \sense{
    \definition{Idi language (Pahoturi River language spoken in Dimsisi, Sibidiri, Dimiri, Iblamand, and Biram)}    \note{Language Names}    \example{
}    \example{
}    \example{
}}}

\entry{Idi loanword}{\headword{Idi loanword}
  \pos{}  \sense{
}}

\entry{idoidog}{\headword{idoidog}  \pronnote{idoidoɡ}
  \pos{Noun}  \sense{
    \definition{harpoon}}}

\entry{Idugoe}{\headword{Idugoe}
  \pos{Proper noun}  \sense{
    \definition{male personal name}}}

\entry{igi}{\headword{igi}  \pronnote{iɡi}
  \pos{Locational}  \sense{\textbf{1.} 
    \definition{bottom}    \example{
      \exsen{Tuk me gogäbaebän dowede ada gogon, ine ik gongkäbägän do igi abal.}      \extran{On top he jumped like this into the water, down to the very bottom.}}    \example{
      \exsen{Igi me era däbe namllam.}      \extran{[You] grab the one on the bottom.}}}  \sense{\textbf{2.} 
    \definition{underneath}    \example{
      \exsen{Ge dädär de igi me era ekaklle da dakonenegän.}      \extran{These rocks underneath, the soil covered them.}}}}

\entry{igi pite}{\headword{igi pite}  \pronnote{iɡi pite}
  \pos{Noun}  \sense{
    \definition{underwear}    \example{
      \exsen{Bongo bäne igi pite de zäme amättalle?}      \extran{Did you already put on your underpants?}}}}

\entry{igiigi}{\headword{igiigi}
  \pos{Locational}  \sense{
    \definition{under, beneath}    \example{
      \exsen{Ge lla da aeya ngämo imnealle beyarän, bogo mängallang abal dan, ngäna obo igiigi me dan.}      \extran{The man that will come after me, he is very powerful; I am beneath him.}}}}

\entry{ik}{\headword{ik}  \pronnote{ik}
  \pos{Locational}  \sense{
    \definition{inside}    \example{
      \exsen{Bogo dindugän do kukiny mama de ikop dägagän dibaeya däbe ik e gotärakän do gotägolän.}      \extran{He ran to a pile of tall grass that he saw, went inside, and hid there.}}    \example{
      \exsen{Lama wa Mana ubi kandärmang me goensegäneyo adawatta Emi aoli kuddäll de ine ik me ikop dägagän.}      \extran{Lama and Mana were both feeling sorry because they saw that Emi almost died in the water.}}    \example{
      \exsen{Bogo poper gogon, ma ik atta gogezänän wälläng e dindugän.}      \extran{He was shocked, got out from inside, and ran into the bush.}}}}

\entry{ikikib}{\headword{ikikib}  \pronnote{ikikib}
  \pos{Noun}  \sense{
    \definition{dizziness}    \example{
      \exsen{Ngänäm ikikib da nallan.}      \extran{My head is spinning (lit. dizziness has gotten me).}}}}

\entry{ikllo}{\headword{ikllo}
  \pos{Noun}  \sense{
    \definition{smoke}    \example{
      \exsen{Yu ikllo da ttongo eka ngangema ttoen dan.}      \extran{Smoke is one way of sending a message.}}}}

\entry{iklloikllowang}{\headword{iklloikllowang}  \pronnote{ikɽoikɽowaŋ}
  \pos{Color term}  \sense{
    \definition{grey; the color of smoke}}}

\entry{ikoikop}{\headword{ikoikop}
  \variantof{Fast Speech Variant of \vartext{ikopikop}}  \pronnote{ikoikop}}

\entry{ikoll}{\headword{ikoll}  \pronnote{ikoɽ}
  \pos{Noun}  \sense{
    \definition{incident, problem, trouble}    \example{
      \exsen{Ubi abo diba ikoll me ma we dängäsmällän.}      \extran{They then returned home from the incident.}}    \example{
      \exsen{Bongo adawede eka dändärmeny agnalle, ikoll a nagan.}      \extran{You weren't listening, so it came back to bite you (lit. 'the trouble got you').}}}}

\entry{ikop}{\headword{ikop}  \pronnote{ikop}
  \pos{Noun}  \sense{\textbf{1.} 
    \definition{eye}    \example{
      \exsen{Llamäg o mällause täräp me iba ikop a mɨnyi säremang bognegän.}      \extran{During old age, our (incl.) eyes will go dim.}}}  \sense{\textbf{2.} 
    \definition{to see}    \example{
      \exsen{Lla da lla de ikop eran.}      \extran{A person sees a person.}}    \example{
      \exsen{Ngäna ddia de ikop nägagan a nazuan.}      \extran{I saw a deer and shot it.}}}  \sense{\textbf{3.} 
    \definition{to look}    \example{
      \exsen{Ikopag imne we.}      \extran{(You) look back.}}    \example{
      \exsen{Omägag a obo nyäng e ikop gogon.}      \extran{The fortune-teller looked into her bag.}}    \example{
      \exsen{Lla da däbe ikop dägneyo käza de.}      \extran{People were looking at that crocodile.}}    \example{
      \exsen{Ge lla de ikop nägayaebeyo!}      \extran{(You all) look at these people!}}}}

\entry{ikop ddäddäg}{\headword{ikop ddäddäg}
  \pos{Intransitive compound verb}  \sense{\textbf{1.} 
    \definition{to look, watch}    \example{
      \exsen{Llɨg a wa män a ubi ikop goddägneyo.}      \extran{The boy and the girl looked at each other.}}}  \sense{\textbf{2.} 
    \definition{to look, watch}    \example{
      \exsen{Masar era ngämim ikop deyaddägnän.}      \extran{Grandfather, he was watching us.}}}}

\entry{ikop glas}{\headword{ikop glas}  \note{Clothing}  \pronnote{ikop ɡlas}
  \pos{Noun}  \sense{
    \definition{eyeglasses, spectacles; goggles}    \note{Clothing}    \example{
      \exsen{kämbmeny ma ikop glas}      \extran{diving goggles}}    \example{
}    \example{
}}}

\entry{ikop käp}{\headword{ikop käp}  \note{Parts of a reptile}  \pronnote{ikop kəp}
  \pos{Noun}  \sense{
    \definition{eyeball}    \note{Parts of a reptile}    \example{
}    \example{
}    \example{
}}}

\entry{ikop kom}{\headword{ikop kom}  \pronnote{ikop kom}
  \pos{Noun}  \sense{
    \definition{eyelash}}}

\entry{ikop ku}{\headword{ikop ku}  \pronnote{ikop ku}
  \pos{Noun}  \sense{
    \definition{pupil}    \example{
      \exsen{Steven a Kange, däräng de dällädeyo, pakos de dängkäneyo, ikop ku att de, däräng a ttam gogän.}      \extran{Steven and Kange grabbed the dog, pulled out the arrow from his pupil, and the dog lived.}}    \example{
}    \example{
}}}

\entry{ikop kutt nängazmeny}{\headword{ikop kutt nängazmeny}  \pronnote{ikop kuʈ͡ʂ nəŋazmeɲ}
  \pos{Phrase}  \sense{
    \definition{stare at someone intently when angry with them.}}}

\entry{ikop sära}{\headword{ikop sära}  \pronnote{ikop səra}
  \pos{Noun}  \sense{
    \definition{wink; eyebrow raise}}}

\entry{ikop särem}{\headword{ikop särem}  \pronnote{ikop sərem}
  \pos{Noun}  \sense{
    \definition{fainting, losing consciousness}    \example{
      \exsen{Ubi mɨnyi ikop särem atta buspullnän nyongo me, adawatta ma we ibima da utale abal dan.}      \extran{They will faint and fall on the road because going home is a very long journey.}}}}

\entry{ikop songgorag}{\headword{ikop songgorag}  \pronnote{ikop soŋɡoraɡ}
  \pos{Modifier}  \sense{
    \definition{blind}    \example{
      \exsen{Däbe ikop songgorag lla da ada eka dägagän, "Ngämo moko da mermer ikopang e dan."}      \extran{That blind man said, "I wish to see properly."}}}}

\entry{ikop su}{\headword{ikop su}  \pronnote{ikop su}
  \pos{Intransitive coverb}  \sense{
    \definition{to peek}    \example{
      \exsen{Llamda da up golloll me gotogolän a ikop su gognän känyertto ubira ikop e ami obo up de gämäll alle däglebneyo.}      \extran{The old man hid behind the banana tree and peeked quietly to see those who were stealing his banana.}}}}

\entry{ikop täbab}{\headword{ikop täbab}  \lexeme{täbab}  \pronnote{təbab}
  \pos{Transitive verb}  \sense{
    \definition{to watch, look after, patrol}    \example{
      \exsen{sip ikop täbaebag}      \extran{shepherd}}}}

\entry{ikop ttoe}{\headword{ikop ttoe}  \pronnote{ikop ʈ͡ʂoe}
  \pos{Noun}  \sense{
    \definition{eyelid}    \example{
}    \example{
}    \example{
}}}

\entry{ikopikop}{\headword{ikopikop}  \pronnote{ikopikop}
  \pos{Transitive/Intransitive coverb}  \sense{
    \definition{to watch}    \example{
      \exsen{Ubi ikopikop gogaebne.}      \extran{They were watching.}}    \example{
      \exsen{Ge lla da obom ikoikop eran.}      \extran{This person is watching him.}}    \example{
}}
  \variants{\Fast Speech Variant \vartext{ikoikop}}}

\entry{ikopmeny}{\headword{ikopmeny}  \pronnote{ikopmeɲ}
  \pos{Modifier}  \sense{
    \definition{unseen, invisible}    \example{
      \exsen{Eiz itrell a ikopmeny dan, lla bo mam me dan.}      \extran{HIV is invisible; it's in a person's blood.}}    \example{
      \exsen{Ause da kollba we gopnaeyän a ine ik mi ikopmeny gogän.}      \extran{The woman turned into a fish and disappeared into the water.}}}}

\entry{ikopse}{\headword{ikopse}  \pronnote{ikopse}
  \pos{Noun}  \sense{
    \definition{prayer}    \example{
      \exsen{Känaebag, ngäna mɨnyi ikopse we beyar, ikopse ttamänatt me mɨnyi Upiara we bongos.}      \extran{Tomorrow, I am going to go pray, and after I finish praying, I will return to Upiara.}}}
  \variants{\SpellingVariant \vartext{ikopsi}}}

\entry{ikopse ma}{\headword{ikopse ma}
  \pos{Noun}  \sense{
    \definition{church, temple}    \example{
      \exsen{Bongo mɨnyi mikutt mi iba ikopse ma de nälläk?}      \extran{Will you burn our (incl.) church in anger?}}}
  \variants{\Unspecified Variant \vartext{ikopse ma ma}}}

\entry{ikopse ma ma}{\headword{ikopse ma ma}
  \variantof{Unspecified Variant of \vartext{ikopse ma}}}

\entry{ikopsi}{\headword{ikopsi}
  \variantof{SpellingVariant of \vartext{ikopse}}}

\entry{ikrol}{\headword{ikrol}  \note{Parts of a fire}  \pronnote{ikrol}
  \pos{Noun}  \sense{
    \definition{ash}    \note{Parts of a fire}    \example{
}    \example{
}    \example{
      \exsen{Ikrol da duwamänän.}      \extran{The ash dissipated.}}}}

\entry{Ilaeza}{\headword{Ilaeza}  \pronnote{ilaeza}
  \pos{Proper noun}  \sense{
    \definition{Elijah}}}

\entry{ileben}{\headword{ileben}
  \variantof{SpellingVariant of \vartext{eleben}}}

\entry{ili}{\headword{ili}  \pronnote{ili}
  \pos{Relative pronoun}  \sense{\textbf{1.} 
    \definition{where}    \example{
      \exsen{Ngämi ili ibi eran mondre ma dan.}      \extran{It's where we (excl.) go gardening.}}    \example{
      \exsen{Da duli melem e ballän, ili Mospi o ili ballän, ede Ingglis eka de mɨnyi bäpanyän.}      \extran{If she goes away for work, in Port Moresby or wherever she goes, then she will speak English.}}}  \sense{\textbf{2.} 
    \definition{where}    \example{
      \exsen{Bongo ili kanyekanye gogäne?}      \extran{Where have you travelled to?}}}}

\entry{ilibattäm}{\headword{ilibattäm}
  \pos{Pronoun}  \sense{
    \definition{ablative form of ili}    \example{
      \exsen{Da ilibattäm lla da bogon, da ada Upiara watt, ngäna obom era ede mɨnyi pällämpälläm eka walle eka bäntameny.}      \extran{If the person is from elsewhere (i.e. non-Ende villages), then I would speak with him using English.}}}}

\entry{imam}{\headword{imam}  \pronnote{imam}
  \pos{Intransitive verb}  \sense{
    \definition{to leave behind; leave without saying goodbye}}}

\entry{imamalle}{\headword{imamalle}  \pronnote{imamaɽe}
  \pos{Adverb}  \sense{
    \definition{intentionally}}}

\entry{Imanuel}{\headword{Imanuel}  \pronnote{imanuel}
  \pos{Proper noun}  \sense{
    \definition{Imanuel}}}

\entry{imne}{\headword{imne}  \pronnote{imne}
  \pos{Locational}  \sense{\textbf{1.} 
    \definition{rear, behind, back}    \example{
      \exsen{Gall me bogo ngattong e, ngäna imne we, käza da amne me daeya.}      \extran{She was in front of the canoe, I was in the rear, and the crocodile was in the middle.}}    \example{
      \exsen{Ddia da säre imne gogon, kottllam a ngattong gongttägän ttamän ma ngättma we.}      \extran{Sadly, Deer was behind; Turtle had arrived to the finish line first.}}    \example{
      \exsen{Ddob lla da ami imne dongkollmällaemneyo lel gognegnän}      \extran{The others who were following in the back were scared.}}    \example{
      \exsen{Ngäna bäne imne ballne.}      \extran{I will walk behind you.}}}  \sense{\textbf{2.} 
    \definition{afterwards, after, later}    \example{
      \exsen{Ngäma mäda gagäll gogon, ngäma mäg ako imne gagäll gogon.}      \extran{Our (excl.) father passed away; later, our mother passed away too.}}    \example{
      \exsen{Imne duwem ag, ngattong ttang de anggälläm.}      \extran{Eat after you wash your hands first.}}}
  \variants{\SpellingVariant \vartext{yimne}}}

\entry{imne~}{\headword{imne~}
  \variantof{Derivational reduplicant of \vartext{RED~}}}

\entry{imnealle}{\headword{imnealle}  \pronnote{imneaɽe}
  \pos{Postposition}  \sense{
    \definition{after, later}    \example{
      \exsen{däbe imnealle}      \extran{after that}}    \example{
      \exsen{kumuddäga pazi imnealle}      \extran{three years later}}}}

\entry{imneimne}{\headword{imneimne}  \pronnote{imneimne}
  \pos{Adverb}  \sense{\textbf{1.} 
    \definition{from behind}    \example{
      \exsen{Yäbäd imneimne gongbeabogän ddone dingmenän.}      \extran{Sun chased him from behind, but did not reach him.}}}  \sense{\textbf{2.} 
    \definition{behind, late; later, after}    \example{
      \exsen{Ngämo nag zäme angttägan be ngäna mɨnyi imneimne bongttäg.}      \extran{My friend already arrived, but I will arrive later.}}}}

\entry{imnong}{\headword{imnong}
  \pos{Modifier}  \sense{
    \definition{after}    \example{
      \exsen{Imnong e ubi dädmawän sipel e.}      \extran{After a while, they sat down to rest.}}}}

\entry{imomdae}{\headword{imomdae}  \note{Adjectives (Swadesh)}  \pronnote{imomdae}
  \pos{Noun}  \sense{\textbf{1.} 
    \definition{truth}    \note{Adjectives (Swadesh)}    \example{
      \exsen{Ngäna bibim imomdae de umllang anggan.}      \extran{I am telling you (pl.) truths.}}    \example{
      \exsen{Bongo imomdae de ekalle.}      \extran{You are speaking the truth.}}}  \sense{\textbf{2.} 
    \definition{true, real, actual}    \note{Adjectives (Swadesh)}    \example{
      \exsen{Ge gänyan Ende imomdae eka dan.}      \extran{This here is the true Ende language.}}    \example{
      \exsen{Ubi anykeanyke dae ikop dägageyo, ubi ddone imomdae up pätt de ikop dägageyo.}      \extran{They saw only the reflection; they did not see the actual banana plant.}}    \example{
      \exsen{Sos me ada eka anggan, "Iba pätt a God bo imomdae giddoll ma dan."}      \extran{They say in church, "Our (incl.) bodies is where God truly resides."}}}  \sense{\textbf{3.} 
    \definition{correct, right}    \note{Adjectives (Swadesh)}    \example{
      \exsen{Bongo imomdae abal de ekalle ada, Adi da ttongdae dan.}      \extran{You are very right in saying that God is one.}}    \example{
}    \example{
}}  \sense{\textbf{4.} 
    \definition{to believe}    \note{Adjectives (Swadesh)}    \example{
      \exsen{Bogo ddone komlla näkäpang bogon obo tikop me, be imomdae bägagän.}      \extran{He will not be doubtful (lit. two-minded) in his heart, but will believe it.}}}}

\entry{imomdae ngasngesang}{\headword{imomdae ngasngesang}  \note{good characteristics of people}  \pronnote{imomdaj ŋasŋesaŋ}
  \pos{Modifier}  \sense{
    \definition{honest}    \note{good characteristics of people}    \example{
}    \example{
}    \example{
}}}

\entry{imomdae ttättlle}{\headword{imomdae ttättlle}  \pronnote{imomdae ʈ͡ʂəʈ͡ʂɽe}
  \pos{Modifier}  \sense{
    \definition{honest}    \example{
      \exsen{Ngäma umllang dan bongo imomdae ttättlle panypenyang dan.}      \extran{We (excl.) know that you speak honestly.}}    \example{
}    \example{
}}}

\entry{imomdae ttoen}{\headword{imomdae ttoen}
  \pos{Noun}  \sense{
    \definition{faith, religion}    \example{
      \exsen{Imomdae ttoen peyang ikopse dae ada ingollang gagäll anyke de bäkällttanegän.}      \extran{Only with faith and prayer can the evil spirits be cast out.}}}}

\entry{imomdae ttoenang}{\headword{imomdae ttoenang}
  \pos{Noun}  \sense{
    \definition{believer, faithful}    \example{
      \exsen{Tämamae ttoen a kuddukuddull dag oblle aeya imomdae ttoenang bognän.}      \extran{Everything is easy for he who may be a believer.}}}}

\entry{imonz}{\headword{imonz}
  \pos{Transitive verb}  \sense{
    \definition{to touch}}}

\entry{imonzimonz}{\headword{imonzimonz}  \pronnote{imonzimonz}
  \pos{Noun}  \sense{
    \definition{tag (game)}    \example{
      \exsen{Mosenda kaptte de dängllämnegnän a mänyanda bi komllaebe tupi imonzimonz gotngoeneyo walle ik mi.}      \extran{The oldest washed the clothes and the younger two played tag in the water.}}}}

\entry{imullgoe}{\headword{imullgoe}
  \pos{Transitive verb}  \sense{
    \definition{to drop, push, make fall}    \example{
      \exsen{Bogo obom daimullgoeyän.}      \extran{She pushed him over.}}}}

\entry{in}{\headword{in}
  \variantof{Dialectal Variant of \vartext{English loanword}}}

\entry{Ina}{\headword{Ina}  \pronnote{ina}
  \pos{Proper noun}  \sense{
    \definition{female personal name}}}

\entry{inam}{\headword{inam}  \pronnote{inam}
  \pos{Transitive verb}  \sense{\textbf{1.} 
    \definition{to cover}    \example{
      \exsen{Ekaklle kok de mɨnyi enanae binamän.}      \extran{Earth will completely cover the moon (i.e. during a lunar eclipse).}}}  \sense{\textbf{2.} 
    \definition{to weigh down, press down}    \example{
      \exsen{Ngäna sɨmell de dinam ekaklle we.}      \extran{I weighed the pig down to the ground.}}}}

\entry{Inapa}{\headword{Inapa}  \pronnote{inapa}
  \pos{Proper noun}  \sense{
    \definition{male personal name}}}

\entry{Inawa}{\headword{Inawa}  \pronnote{inawa}
  \pos{Proper noun}  \sense{
    \definition{male personal name}}}

\entry{inbo}{\headword{inbo}  \lexeme{inbo}  \note{Family ties}  \pronnote{inbo}
  \pos{Kinship noun}  \sense{\textbf{1.} 
    \definition{brother-in-law (woman's husband's brother)}    \note{Family ties}}  \sense{\textbf{2.} 
    \definition{sister-in-law (man's elder brother's wife or woman's husband's younger sister; reciprocal)}    \note{Family ties}}
  \variants{\SpellingVariant \vartext{yinbo}}}

\entry{inbunatt}{\headword{inbunatt}  \pronnote{inbunaʈ͡ʂ}
  \pos{Noun}  \sense{
    \definition{mallet fish (big, lives in creek, lean)}}}

\entry{independence}{\headword{independence}
  \variantof{Dialectal Variant of \vartext{English loanword}}}

\entry{Indonesia}{\headword{Indonesia}  \pronnote{indonesia}
  \pos{Proper noun}  \sense{
    \definition{Indonesia}}}

\entry{indraindrang}{\headword{indraindrang}
  \pos{Adverb}  \sense{
    \definition{clearly}    \example{
      \exsen{Mak ibim indraindrang umllang anggan ada Yesu era Adi bo llɨg aenen.}      \extran{Mark told us (incl.) clearly that Jesus is the Son of God.}}}}

\entry{indrang}{\headword{indrang}  \pronnote{indraŋ}
  \pos{Noun}  \sense{\textbf{1.} 
    \definition{light}    \example{
      \exsen{Ngäna ngämo tos indrang de dɨs.}      \extran{I put out my flashlight.}}    \example{
      \exsen{Indrang a ngämo dowae me gogllaeyän.}      \extran{The light shined close to me.}}}  \sense{\textbf{2.} 
    \definition{luminous, bright}    \example{
      \exsen{ttongo mer indrang ag}      \extran{one very bright morning}}    \example{
      \exsen{Yäbäd a mɨnyi säremang bogon ebdo me a kok a ade mɨnyi ddone indrang bognän iddob me.}      \extran{The sun will be dark during the day and also, the moon will not be bright at night.}}}  \sense{\textbf{3.} 
    \definition{to show, reveal, illuminate, bring to light}    \example{
      \exsen{Ngäna Adi lla bo mer de indrang eran.}      \extran{I am showing God's goodness.}}    \example{
      \exsen{Aeya ge eka de indrang eran?}      \extran{Who is announcing this news?}}}
  \variants{\SpellingVariant \vartext{yindrang}}}

\entry{indre}{\headword{indre}  \note{Trees}  \pronnote{indre}
  \pos{Noun}  \sense{
    \definition{type of tree that grows in the grassland with edible brown-yellow fruit and good, long-burning firewood}    \note{Trees}    \example{
}    \example{
}    \example{
}}}

\entry{indre gullem}{\headword{indre gullem}  \pronnote{indre}
  \pos{Noun}  \sense{
    \definition{type of snake}}}

\entry{indugoeg}{\headword{indugoeg}  \pronnote{induɡoeɡ}
  \pos{Transitive verb}  \sense{
    \definition{to command}    \example{
      \exsen{Bongo gagäll anyke, ngäna bam indugoeg nallan, agezän obo guwo watta.}      \extran{You evil spirit, I command you to come out of his heart.}}    \example{
      \exsen{Mägda män de naindugoegan kaptte mättmätt e.}      \extran{The girl's mother commanded her to put on clothes.}}}}

\entry{indurgoeg}{\headword{indurgoeg}
  \pos{Transitive verb}  \sense{
    \definition{to yell}}}

\entry{ine}{\headword{ine}  \lexeme{ine}  \pronnote{ine}
  \pos{Noun}  \sense{\textbf{1.} 
    \definition{water; liquid}    \example{
      \exsen{nge ine}      \extran{coconut water}}    \example{
      \exsen{Ge ine da nanott dan.}      \extran{This water has been drunk.}}    \example{
      \exsen{Bogo mängalae gogbaebän ine ik e oblle ngämingg e.}      \extran{She quickly jumped into the water to help him.}}}  \sense{\textbf{2.} 
    \definition{alcoholic beverage}    \example{
      \exsen{Ubira mer näkäpngon a ddone dan adawatta ubi ine de nanen amallo.}      \extran{They don't have good judgment because they are drinking alcohol.}}}}

\entry{ine bib}{\headword{ine bib}  \note{Water}  \pronnote{ine bib}
  \pos{Noun}  \sense{
    \definition{spring (water source)}    \note{Water}    \example{
}    \example{
}    \example{
}}}

\entry{ine källa}{\headword{ine källa}  \lexeme{ine källa}  \pronnote{ine kəɽa}
  \pos{Noun}  \sense{
    \definition{mud, clay}    \grammarnote{deer and wallabies come to eat it}}}

\entry{ine konkonang}{\headword{ine konkonang}
  \pos{Noun}  \sense{
}}

\entry{ine kup}{\headword{ine kup}  \lexeme{ine kup}  \pronnote{ine kup}
  \pos{Noun}  \sense{
    \definition{well (water source)}}}

\entry{ine kutt}{\headword{ine kutt}  \note{Everything in the house}  \pronnote{ine kuʈ͡ʂ}
  \pos{Noun}  \sense{
    \definition{bucket, water container}    \note{Everything in the house}    \example{
      \exsen{Ngäna bänene ine kutt de yuwetyuwet yuwenyan.}      \extran{I borrowed your bucket.}}    \example{
}    \example{
}}}

\entry{ine ma}{\headword{ine ma}  \pronnote{ine ma}
  \pos{Noun}  \sense{
    \definition{well (water source)}    \example{
      \exsen{Bongo ine ma nalle oba bongo ine neyatan.}      \extran{You, go to the well and fetch water.}}    \example{
}    \example{
}}}

\entry{ine mallmell}{\headword{ine mallmell}  \note{Snakes}  \pronnote{ine maɽmeɽ}
  \pos{Noun}  \sense{
    \definition{type of snake}    \note{Snakes}    \example{
}    \example{
}    \example{
}}}

\entry{ine takmäll}{\headword{ine takmäll}  \lexeme{ine takmäll}  \pronnote{ine takməɽ}
  \pos{Modifier}  \sense{
    \definition{seasick}}}

\entry{ingellab}{\headword{ingellab}  \note{Animal verbs}  \pronnote{iŋeɽab}
  \pos{Verb}  \sense{
    \definition{to rise up, said of an animal}    \note{Animal verbs}    \example{
      \exsen{Bogo ingellabang dan.}      \extran{He rises up in attention}}    \example{
}    \example{
}}}

\entry{Ingglis}{\headword{Ingglis}
  \pos{Proper noun}  \sense{
    \definition{English language}}
  \variants{\SpellingVariant \vartext{Englis, English, ingglis}}}

\entry{ingglis}{\headword{ingglis}
  \variantof{SpellingVariant of \vartext{Ingglis}}}

\entry{inggoe}{\headword{inggoe}
  \variantof{Dialectal Variant of \vartext{inggol}}  \pronnote{iŋɡoe}}

\entry{inggoemeny}{\headword{inggoemeny}  \note{Oral Traditions}  \pronnote{iŋɡoemeɲ}
  \pos{Intransitive verb}  \sense{\textbf{1.} 
    \definition{to joke}    \note{Oral Traditions}    \example{
      \exsen{Ende walle inggoemeny amalla.}      \extran{We are making jokes in Ende.}}    \example{
      \exsen{Ngäna Ingglis eka walle de eka bogne, buinggoemenyne.}      \extran{I will speak in English, making jokes.}}    \example{
}}  \sense{\textbf{2.} 
    \definition{to blaspheme; insult, mock}    \note{Oral Traditions}    \example{
      \exsen{Bogo Adi bom inggoemeny eran.}      \extran{He is blaspheming God.}}    \example{
      \exsen{Sabi tamenyang a inggoemeny eka de däpanyaemneyo.}      \extran{The teachers of the law were saying insults.}}}
  \variants{\SpellingVariant \vartext{ingguimeny}}}

\entry{inggol}{\headword{inggol}  \pronnote{iŋɡol}
  \pos{Intransitive verb}  \sense{
    \definition{to move}    \example{
      \exsen{Mälla da angde ikop dägagän, llɨg bo dowae e guinggolän.}      \extran{When the woman saw it, she moved closer to the boy.}}}
  \variants{\Dialectal Variant \vartext{inggoe}}}

\entry{ingguimeny}{\headword{ingguimeny}
  \variantof{SpellingVariant of \vartext{inggoemeny}}}

\entry{ingneningnen}{\headword{ingneningnen}
  \pos{Adverb}  \sense{
    \definition{dancing around}    \example{
      \exsen{Bogo ingnenignen giddollag dan.}      \extran{He lives dancing about.}}    \example{
      \exsen{Ngäna ingnenignenang ibi allan gänyaolle.}      \extran{I am coming this way dancing.}}}}

\entry{ingoll}{\headword{ingoll}
  \pos{Noun}  \sense{\textbf{1.} 
    \definition{face}    \example{
      \exsen{Angde bongo bäne mälla bo ingoll papa eralle, bongo God bo ingoll papa eralle.}      \extran{When you hit your wife in the face, you are hitting God in the face.}}    \example{
      \exsen{Ubi obo ingoll de iddpo alle dämllaeyo.}      \extran{They tied his face with cloth (i.e. they blindfolded him).}}}  \sense{\textbf{2.} 
    \definition{front}    \example{
      \exsen{Ttongo lla da agbänmällnan a ttongo lla da nindugan obo ingoll me.}      \extran{One man was jumping and another man ran in front of him.}}    \example{
      \exsen{Ulle binang itrel peyang lla bo ingoll alle eka tameny a mudan.}      \extran{Don't converse in front of a very sick man.}}}}

\entry{=ingoll}{\headword{=ingoll}  \pronnote{iŋoɽ}
  \pos{Nominal enclitic}  \sense{
    \definition{similative clitic; like}    \example{
      \exsen{Bänz a pam nidol ingoll dagaeya.}      \extran{The mosquitoes were like needles.}}    \example{
      \exsen{Täme da käza ingoll ddäddäg be ap me giddollag dan.}      \extran{The goanna is crocodile-like animal but it lives in the grassland.}}}}

\entry{ingong}{\headword{ingong}  \pronnote{iŋoŋ}
  \pos{Noun}  \sense{\textbf{1.} 
    \definition{dance}    \example{
      \exsen{ingong täräp}      \extran{dance time}}}  \sense{\textbf{2.} 
    \definition{sing-sing}}  \sense{\textbf{3.} 
    \definition{to dance}    \example{
      \exsen{Obo moko da ingong e dan.}      \extran{He wants to dance.}}    \example{
      \exsen{Ngäna gänyaolle bengne.}      \extran{I will dance this way.}}    \example{
      \exsen{Ngäna bam ingongang nallan.}      \extran{I am making you dance.}}    \example{
      \exsen{Lla da ada dänyängnän menae me adawatta bogo ngattong abal gongttägän ttamän ma ngättma we.}      \extran{People were dancing like this on the sides because he got to the finish line first.}}}}

\entry{inkäm}{\headword{inkäm}  \pronnote{inkəm}
  \pos{Noun}  \sense{\textbf{1.} 
    \definition{voice}    \example{
      \exsen{Ngämi bam ge ddobae ddonddo nalla adawatta bäne ge mer abal inkäm dan.}      \extran{We are very proud of you because of your very nice voice.}}}  \sense{\textbf{2.} 
    \definition{neck}    \example{
      \exsen{Obo inkäm a tubutubu dan.}      \extran{His neck is short.}}}}

\entry{inkäm tän}{\headword{inkäm tän}
  \pos{Noun}  \sense{
}}

\entry{inkäm tär}{\headword{inkäm tär}  \note{Parts of a bird}  \pronnote{inkəm tər}
  \pos{Noun}  \sense{
    \definition{throat}    \note{Parts of a bird}    \example{
      \exsen{Pa bo inkäm tär a ddäddäg ma dan.}      \extran{A bird's throat is edible.}}    \example{
}    \example{
}}}

\entry{inkätt}{\headword{inkätt}  \note{Musical instruments}  \pronnote{inkəʈ͡ʂ}
  \pos{Noun}  \sense{\textbf{1.} 
    \definition{neck; throat}    \note{Musical instruments}    \example{
      \exsen{Bogo Zon ine kämbägag bom inkätt danddugän.}      \extran{He sliced through the throat of John the Baptist.}}}  \sense{\textbf{2.} 
    \definition{hollow of drum}    \note{Musical instruments}    \example{
}    \example{
}}}

\entry{inmol}{\headword{inmol}  \note{Trees}  \pronnote{inmol}
  \pos{Noun}  \sense{
    \definition{type of small tree that grows in the bush along creeks; used to weave grass skirts after being pounded}    \note{Trees}    \example{
}    \example{
}    \example{
}}}

\entry{inmoll}{\headword{inmoll}  \pronnote{inmoɽ}
  \pos{Transitive verb}  \sense{\textbf{1.} 
    \definition{to cover}    \example{
      \exsen{Ddapall a säremang gogon a deyainmollän.}      \extran{The sky got dark and overcast.}}}  \sense{\textbf{2.} 
    \definition{to step on; vanquish}    \example{
      \exsen{Ngäna bäne mäkang lla de mängallmeny bägneg a bäne ttäle koll igiigi me bowattäll ubira inmoll e.}      \extran{I will make your enemies powerless and put them beneath your feet, vanquished.}}}}

\entry{inngoeinngoe}{\headword{inngoeinngoe}
  \pos{Modifier}  \sense{
    \definition{shaky}    \example{
      \exsen{Ngämi tätäm era inngoeinngoe mo dae daupeya.}      \extran{Yesterday, we (excl.) crossed a shaky bridge.}}}}

\entry{inpiak}{\headword{inpiak}  \pronnote{inpiak}
  \pos{Noun}  \sense{
    \definition{whistling kite}    \scientificname{Haliastur sphenurus}    \example{
}    \example{
}    \example{
}}}

\entry{Inpiakma}{\headword{Inpiakma}  \note{Places around Limol}  \pronnote{inpiakma}
  \pos{Proper noun}  \sense{
    \definition{Inpiakma (toponym)}    \note{Places around Limol}    \example{
}    \example{
}    \example{
}}}

\entry{Inpir}{\headword{Inpir}
  \pos{Proper noun}  \sense{
    \definition{Makayam/Tirio language}}}

\entry{intot}{\headword{intot}  \pronnote{intot}
  \pos{Noun}  \sense{
    \definition{brow bone}    \example{
}    \example{
}    \example{
}}}

\entry{intot kom}{\headword{intot kom}  \note{Parts of the body}  \pronnote{intot kom}
  \pos{Noun}  \sense{
    \definition{eyebrow}    \note{Parts of the body}    \example{
}    \example{
}    \example{
}}}

\entry{inttemängg}{\headword{inttemängg}  \pronnote{inʈ͡ʂeməŋɡ}
  \pos{Intransitive verb}  \sense{\textbf{1.} 
    \definition{to part ways, leave}    \example{
      \exsen{Eka tamenyatt me, ubi guinttemänggeyo a oba ma we deyareyo.}      \extran{After discussing, they parted ways and went to their own homes.}}}  \sense{\textbf{2.} 
    \definition{to leave, see off, release, set free}    \example{
      \exsen{Kottllam bom dätrameyo bem ulle we dainttemonggeyo.}      \extran{They carried the turtle to the ocean and set him free.}}}}

\entry{inttemongg}{\headword{inttemongg}
  \pos{Transitive verb}  \sense{
    \definition{to let go}}}

\entry{inttoemängg}{\headword{inttoemängg}
  \pos{Transitive verb}  \sense{
    \definition{to see (someone) off}}}

\entry{inu}{\headword{inu}  \pronnote{inu}
  \pos{Intransitive verb}  \sense{\textbf{1.} 
    \definition{to sleep}    \example{
      \exsen{Ngäna inunin allan.}      \extran{I am sleeping.}}    \example{
      \exsen{Gänyme inuag!}      \extran{[You] sleep here!}}    \example{
      \exsen{Bogo gotarnän pollon me.}      \extran{He was sleeping in a bush.}}}  \sense{\textbf{2.} 
    \definition{sleep}    \example{
      \exsen{Bogo inu atta angällbänan.}      \extran{He awoke from sleep.}}}  \sense{\textbf{3.} 
    \definition{asleep, sleeping}    \example{
      \exsen{Bundae inu dan.}      \extran{Bundae is asleep.}}}  \sense{\textbf{4.} 
    \definition{night}    \example{
      \exsen{Dädme ngämi kumuddäga inu gogmam.}      \extran{We (excl.) stayed there three nights.}}}
  \variants{\SpellingVariant \vartext{yinu}}
  \variants{\Unspecified Variant \vartext{yunu}}}

\entry{inu kuddäll}{\headword{inu kuddäll}  \pronnote{inu kuɖ͡ʐəɽ}
  \pos{Modifier}  \sense{
    \definition{fast asleep, dead asleep}    \example{
      \exsen{Llamäg a inu kuddäll daeya.}      \extran{The old man was fast asleep.}}    \example{
}    \example{
}}}

\entry{inu tanter}{\headword{inu tanter}  \lexeme{tanter}  \pronnote{tanter}
  \pos{Transitive verb}  \sense{
    \definition{to guard in one's sleep, sleep with}    \example{
      \exsen{Bogo gotarnän ma da däbe yunu danteralle.}      \extran{She slept over, guarding that house.}}}}

\entry{inuinu}{\headword{inuinu}  \pronnote{inuinu}
  \pos{Noun}  \sense{
    \definition{white-faced robin}    \scientificname{Tregellasia leucops}}}

\entry{inuinu}{\headword{inuinu}  \pronnote{inuinu}
  \pos{Adverb}  \sense{
    \definition{sleepy}    \example{
      \exsen{Ngäna bam inuinu nallan.}      \extran{I'm making you sleepy.}}    \example{
}    \example{
}}}

\entry{inungoe}{\headword{inungoe}  \pronnote{inuŋoe}
  \pos{Transitive verb}  \sense{
    \definition{to shake (when dancing)}    \example{
      \exsen{Bongo nainungoenalleǃ}      \extran{You were shaking!}}}}

\entry{ioläm}{\headword{ioläm}  \note{Birds}  \pronnote{ioləm}
  \pos{Noun}  \sense{
    \definition{type of bird}    \note{Birds}    \example{
}    \example{
}    \example{
}}}

\entry{ip}{\headword{ip}  \pronnote{ip}
  \pos{Noun}  \sense{
    \definition{type of tree that grows in the grassland with bark that is chewed and sap used as an adhesive or poured on a spear to strengthen it}    \example{
}    \example{
}    \example{
}}}

\entry{Ipott}{\headword{Ipott}  \note{Places around Limol}  \pronnote{ipoʈ͡ʂ}
  \pos{Proper noun}  \sense{
    \definition{Ipott (on the road to Bisuaka near Limol ma kuddäll)}    \note{Places around Limol}    \example{
}    \example{
}    \example{
}}}

\entry{Iräm}{\headword{Iräm}  \note{Place names around Limol}  \pronnote{irəm}
  \pos{Proper noun}  \sense{
    \definition{creek and washing place (in Limol, AZ95)}    \note{Place names around Limol}    \example{
}    \example{
}    \example{
}}
  \variants{\SpellingVariant \vartext{yirɨm}}}

\entry{Iräm ine}{\headword{Iräm ine}  \pronnote{irəm ine}
  \pos{Proper noun}  \sense{
    \definition{Iräm creek}}}

\entry{Iräm tatuma}{\headword{Iräm tatuma}  \note{Places around Limol}  \pronnote{irəm tatuma}
  \pos{Proper noun}  \sense{
    \definition{Iräm washing place}    \note{Places around Limol}    \example{
}    \example{
}    \example{
}}}

\entry{irängän}{\headword{irängän}  \pronnote{irəŋən}
  \pos{Intransitive verb}  \sense{\textbf{1.} 
    \definition{to come out}    \example{
      \exsen{Lla da ine ik atta guirängänän.}      \extran{The man came out from the water.}}}  \sense{\textbf{2.} 
    \definition{to get out, lift out}    \example{
      \exsen{Angde ngäna net de tuk i dirängän, käza da net ik i gozenän.}      \extran{When I lifted the net up, a crocodile went inside it.}}    \example{
      \exsen{Mana ddone mullae gogon obo känyer, ede adawatta bogo Lama bälle wälle gogon, Emi bälle yirngän e.}      \extran{Mana wasn't able to herself, so she called for Lama to lift Emi out.}}}
  \variants{\Unspecified Variant \vartext{yerngän}}
  \variants{\Fast Speech Variant \vartext{irngän}}}

\entry{iräpang}{\headword{iräpang}  \pronnote{irəpaŋ}
  \pos{Modifier}  \sense{
    \definition{dirty}    \example{
      \exsen{Bogo ine me gongkäbägän, ede obo inggol iräpang agan.}      \extran{He dived into the water, so his face is dirty.}}}}

\entry{irngän}{\headword{irngän}
  \variantof{Fast Speech Variant of \vartext{irängän}}}

\entry{irwe}{\headword{irwe}  \note{Trees}  \pronnote{irwe}
  \pos{Noun}  \sense{
    \definition{type of cultivated tree with white flowers and juicy, red and white fruit with two seeds; used to treat cough}    \note{Trees}    \example{
}    \example{
}    \example{
}}}

\entry{irwema}{\headword{irwema}  \note{Friendship/social terms of endearment}  \pronnote{irwema}
  \pos{Noun}  \sense{
    \definition{two female friends who share a twin fruit from the irwe tree}    \note{Friendship/social terms of endearment}    \example{
}    \example{
}    \example{
}}}

\entry{Isago}{\headword{Isago}  \pronnote{isaɡo}
  \pos{Proper noun}  \sense{
    \definition{Isago (big island, perhaps in Fly River)}    \example{
}    \example{
}    \example{
}}}

\entry{Isaiah}{\headword{Isaiah}
  \variantof{SpellingVariant of \vartext{Azaya}}}

\entry{isipi}{\headword{isipi}
  \variantof{Dialectal Variant of \vartext{English loanword}}}

\entry{ist}{\headword{ist}
  \pos{Noun}  \sense{
    \definition{yeast}}}

\entry{it}{\headword{it}
  \variantof{Dialectal Variant of \vartext{English loanword}}}

\entry{ita}{\headword{ita}  \pronnote{ita}
  \pos{Noun}  \sense{
    \definition{type of sedge}}}

\entry{itaita}{\headword{itaita}  \note{Trees}  \pronnote{itaita}
  \pos{Noun}  \sense{
    \definition{type of big tree that grows in the bush with red fruit and a straight trunk used for timber}    \note{Trees}    \example{
}    \example{
}    \example{
}}}

\entry{itbonmäll}{\headword{itbonmäll}
  \pos{Noun}  \sense{
    \definition{dirt}    \example{
      \exsen{Ngämlle ge kaptte da ibonmällang agnegan, ede ngänäm tatuma natram.}      \extran{These clothes got dirty on me, so [you] show me to the washing place.}}}}

\entry{itrae}{\headword{itrae}
  \pos{Intransitive verb}  \sense{
    \definition{to hang}}}

\entry{itrell}{\headword{itrell}  \pronnote{itreɽ}
  \pos{Noun}  \sense{\textbf{1.} 
    \definition{disease, illness, sickness}    \example{
      \exsen{Kandärmang säre bablle pätt a ddone abal agan itrell atta.}      \extran{Sorry that the illness has made your body so unwell.}}}  \sense{\textbf{2.} 
    \definition{to get hurt}    \example{
      \exsen{Da ubi gullem de bänglläbaebneyo, ubi mɨnyi ddone buitrellnegnän.}      \extran{If they pick up the snakes, they won't get hurt.}}}}

\entry{itrellang}{\headword{itrellang}
  \pos{Modifier}  \sense{
    \definition{ill, sick}    \example{
      \exsen{Yesu yuwog abal itrellang lla de ai dägnegnän.}      \extran{Jesus healed many sick men.}}}}

\entry{ittal}{\headword{ittal}
  \variantof{Unspecified Variant of \vartext{ittall}}}

\entry{ittall}{\headword{ittall}  \note{Verbs}  \pronnote{iʈ͡ʂaɽ}
  \pos{Transitive/Intransitive coverb}  \sense{\textbf{1.} 
    \definition{to hang}    \note{Verbs}    \example{
      \exsen{Kaptte da ittananang da pätanan e.}      \extran{The clothes are hanging to dry.}}    \example{
      \exsen{Mälla da net de mɨnyi bonyän baittallän.}      \extran{The woman will take the net and hang it.}}    \example{
      \exsen{Molemoleg a llo ttam me itraenen anggan.}      \extran{Caterpillars hang from tree leaves.}}}  \sense{\textbf{2.} 
    \definition{to follow closely, stalk (+ peyang)}    \note{Verbs}    \example{
      \exsen{Ngäna Wagibo bo peyang buittall.}      \extran{I will stalk Wagibo.}}}
  \variants{\Unspecified Variant \vartext{ittal}}}

\entry{ittitt}{\headword{ittitt}
  \variantof{SpellingVariant of \vartext{ittɨtt}}}

\entry{ittɨtt}{\headword{ittɨtt}  \pronnote{iʈ͡ʂɪʈ͡ʂ}
  \pos{Transitive verb}  \sense{
    \definition{to catch (an aquatic animal)}    \example{
      \exsen{Angde toto gognän, Kila a Wala gall alle tudi ittaenen e deyareyo bem e.}      \extran{When evening came, Kila and Wala went to the sea to fish with the canoe.}}    \example{
      \exsen{Bogo komlla guzi de deyaittän.}      \extran{She caught two crawfish.}}    \example{
      \exsen{Ttongdae kollba da guittän.}      \extran{One fish was caught.}}}
  \variants{\SpellingVariant \vartext{ittitt}}}

\entry{ittma}{\headword{ittma}
  \pos{Noun}  \sense{
    \definition{husband's house}    \example{
      \exsen{ittma we zan}      \extran{ceremony in which a woman enters her husband's household}}}}

\entry{ittpo}{\headword{ittpo}
  \variantof{SpellingVariant of \vartext{iddpo}}}

\entry{Ivan}{\headword{Ivan}  \pronnote{ivan}
  \pos{Proper noun}  \sense{
    \definition{male personal name}}}

\entry{iwae}{\headword{iwae}  \note{Tools}  \pronnote{iwae}
  \pos{Noun}  \sense{
    \definition{type of weapon}    \note{Tools}    \example{
}    \example{
}    \example{
}}}

\entry{iya}{\headword{iya}  \pronnote{ija}
  \pos{Noun}  \sense{
    \definition{Australian masked owl}    \scientificname{Tyto novaehollandiae calabyi}    \example{
}    \example{
}    \example{
}}}

\entry{iyeiyem}{\headword{iyeiyem}  \note{Birds}  \pronnote{ijeijem}
  \pos{Noun}  \sense{
    \definition{common emerald dove}    \scientificname{Chalcophaps indica}    \note{Birds}    \example{
}    \example{
}    \example{
}}}

\entry{Izag}{\headword{Izag}
  \pos{Noun}  \sense{
    \definition{name of clan}}}

\entry{izig}{\headword{izig}  \lexeme{izig}  \pronnote{iziɡ}
  \pos{Kinship noun}  \sense{\textbf{1.} 
    \definition{co-sister-in-law (woman's husband's brother's wife)}}  \sense{\textbf{2.} 
    \definition{co-wife (another woman married to the same husband)}}}

\entry{izigag}{\headword{izigag}  \pronnote{iziɡaɡ}
  \pos{Noun}  \sense{
    \definition{polygynous marriage}}}

\entry{iziz}{\headword{iziz}  \note{Government}  \pronnote{iziz}
  \pos{Noun}  \sense{
    \definition{bribery}    \note{Government}    \example{
}    \example{
}    \example{
}}}

\entry{izozallo}{\headword{izozallo}  \pronnote{izozaɽo}
  \pos{Noun}  \sense{
    \definition{animal friend}}}

\chapter{II}


\entry{iid}{\headword{iid}
  \variantof{SpellingVariant of \vartext{yid}}}

\chapter{Ɨ}


\entry{~ɨt}{\headword{~ɨt}
  \variantof{Inflected Form of \vartext{~RED}}}

\chapter{J}


\entry{Jack}{\headword{Jack}
  \pos{Proper noun}  \sense{
    \definition{male personal name}}}

\entry{Jackae}{\headword{Jackae}  \pronnote{d͡ʒakae}
  \pos{Proper noun}  \sense{
    \definition{male personal name}}}

\entry{Jacklin}{\headword{Jacklin}
  \pos{Proper noun}  \sense{
    \definition{female personal name}}}

\entry{Jae}{\headword{Jae}  \pronnote{d͡ʒae}
  \pos{Proper noun}  \sense{
    \definition{male personal name}}}

\entry{James}{\headword{James}
  \variantof{SpellingVariant of \vartext{Zeims}}}

\entry{Jamila}{\headword{Jamila}
  \variantof{SpellingVariant of \vartext{Jamilah}}  \pronnote{d͡ʒamila}}

\entry{Jamilah}{\headword{Jamilah}  \pronnote{d͡ʒamila}
  \pos{Proper noun}  \sense{
    \definition{female personal name}}
  \variants{\SpellingVariant \vartext{Jamila}}}

\entry{Jane}{\headword{Jane}  \pronnote{d͡ʒen}
  \pos{Proper noun}  \sense{
    \definition{female personal name}}}

\entry{Jeff}{\headword{Jeff}  \pronnote{d͡ʒef}
  \pos{Proper noun}  \sense{
    \definition{male personal name}}}

\entry{Jeiks}{\headword{Jeiks}
  \pos{Proper noun}  \sense{
    \definition{male personal name}}}

\entry{Jeimi}{\headword{Jeimi}
  \pos{Proper noun}  \sense{
    \definition{male personal name}}}

\entry{Jemila}{\headword{Jemila}
  \pos{Proper noun}  \sense{
    \definition{female personal name}}
  \variants{\SpellingVariant \vartext{Zemila}}}

\entry{Jen}{\headword{Jen}
  \pos{Proper noun}  \sense{
    \definition{male personal name}}}

\entry{Jeped}{\headword{Jeped}  \pronnote{d͡ʒeped}
  \pos{Proper noun}  \sense{
    \definition{male personal name}}}

\entry{Jerry}{\headword{Jerry}  \pronnote{d͡ʒɛri}
  \pos{Proper noun}  \sense{
    \definition{male personal name}}
  \variants{\SpellingVariant \vartext{Zeri}}}

\entry{Jo}{\headword{Jo}
  \pos{Proper noun}  \sense{
    \definition{male personal name}}}

\entry{Joanna}{\headword{Joanna}  \pronnote{d͡ʒoana}
  \pos{Proper noun}  \sense{
    \definition{female personal name}}}

\entry{Joden}{\headword{Joden}
  \pos{Proper noun}  \sense{
    \definition{male personal name}}}

\entry{Joe}{\headword{Joe}
  \variantof{SpellingVariant of \vartext{Zo}}  \pronnote{d͡ʒo}}

\entry{Joebert}{\headword{Joebert}  \pronnote{d͡ʒobert}
  \pos{Proper noun}  \sense{
    \definition{male personal name}}}

\entry{Joe-noh}{\headword{Joe-noh}  \pronnote{d͡ʒo-no}
  \pos{Proper noun}  \sense{
    \definition{male personal name}}}

\entry{John}{\headword{John}  \pronnote{d͡ʒan}
  \pos{Proper noun}  \sense{
    \definition{male personal name}}
  \variants{\SpellingVariant \vartext{Zon}}}

\entry{Johnny}{\headword{Johnny}  \pronnote{d͡ʒani}
  \pos{Proper noun}  \sense{
    \definition{male personal name}}
  \variants{\SpellingVariant \vartext{Joni, Zoni}}}

\entry{Jonas}{\headword{Jonas}
  \variantof{SpellingVariant of \vartext{Zonas}}}

\entry{Jonathan}{\headword{Jonathan}  \pronnote{d͡ʒanaθɪn}
  \pos{Proper noun}  \sense{
    \definition{male personal name}}}

\entry{Joni}{\headword{Joni}
  \variantof{SpellingVariant of \vartext{Johnny}}}

\entry{Jordan}{\headword{Jordan}  \pronnote{d͡ʒordan}
  \pos{Proper noun}  \sense{
    \definition{male personal name}}}

\entry{Joseph}{\headword{Joseph}
  \variantof{SpellingVariant of \vartext{Zosep}}  \pronnote{d͡ʒosef}}

\entry{Joshua}{\headword{Joshua}  \pronnote{d͡ʒoʃua}
  \pos{Proper noun}  \sense{
    \definition{male personal name}}}

\entry{Jowanang}{\headword{Jowanang}
  \pos{Proper noun}  \sense{
    \definition{male personal name}}}

\entry{Joy-Lin}{\headword{Joy-Lin}  \pronnote{d͡ʒoi-lin}
  \pos{Proper noun}  \sense{
    \definition{female personal name}}}

\entry{Joys}{\headword{Joys}  \pronnote{d͡ʒois}
  \pos{Proper noun}  \sense{
    \definition{female personal name}}}

\entry{Jubli}{\headword{Jubli}
  \pos{Proper noun}  \sense{
    \definition{male personal name}}}

\entry{Judas}{\headword{Judas}
  \variantof{SpellingVariant of \vartext{Zudas}}}

\entry{Jugu}{\headword{Jugu}  \pronnote{d͡ʒuɡu}
  \pos{Proper noun}  \sense{
    \definition{male personal name}}}

\entry{Julia}{\headword{Julia}  \pronnote{d͡ʒulia}
  \pos{Proper noun}  \sense{
    \definition{female personal name}}}

\entry{Julie}{\headword{Julie}
  \variantof{SpellingVariant of \vartext{Zuli}}}

\entry{Julienne}{\headword{Julienne}  \pronnote{d͡ʒulien}
  \pos{Proper noun}  \sense{
    \definition{female personal name}}}

\entry{Junior}{\headword{Junior}  \pronnote{d͡ʒunjər}
  \pos{Proper noun}  \sense{
    \definition{male personal name}}}

\chapter{K}


\entry{ka}{\headword{ka}  \pronnote{ka}
  \pos{TAM particle}  \sense{\textbf{1.} 
    \definition{counterfactual particle (often preceded by ada)}    \example{
      \exsen{Ngämaene gall a ada ka bopänayän, adawatta däbe za da ngämim deyangkollmällnän.}      \extran{We thought our (excl.) boat was going to capsize, because that thing was following us.}}    \example{
      \exsen{Ddia da ttam agan adawatta ngäna bam ada ka bongo walle we aspunalle.}      \extran{The deer is alive because I thought that you fell into the water.}}    \example{
      \exsen{Mer llo popo alle bongo darbänen alle, ada ka män duwar da bun mi popo de nowattällan.}      \extran{You decorate with nice tree flowers like a young girl puts flowers in her hair.}}    \example{
      \exsen{Da klloklloe eka bäpanyeyo, ka Ende eka mɨnyi budabän.}      \extran{If they speak a mixed language, it seems the Ende language will be lost.}}}  \sense{\textbf{2.} 
    \definition{question particle}    \example{
      \exsen{E pazi me ka?}      \extran{In what year?}}}  \sense{\textbf{3.} 
    \definition{emphatic particle}    \example{
      \exsen{Däbe tätämatt winy a ka tomowangtomowang moko allan.}      \extran{That honey from yesterday, it's tasting a bit bitter.}}}}

\entry{ka}{\headword{ka}  \pronnote{ka}
  \pos{Interjection}  \sense{
    \definition{no}    \example{
      \exsen{― Skul e angosan. ― Ka, ma me daden.}      \extran{― He returned to school. ― No, he's at home.}}}}

\entry{kab}{\headword{kab}  \pronnote{kab}
  \pos{Noun}  \sense{
    \definition{string, rope, fiber}    \example{
      \exsen{kab llatet att tär}      \extran{string made from twisted fiber}}    \example{
      \exsen{Ddone aeya mullae gogalle oblle kab alle mällamälla we.}      \extran{No one was able to tie him up with a rope.}}}}

\entry{käba}{\headword{käba}  \pronnote{kəba}
  \pos{Noun}  \sense{
    \definition{solo hunt early in the morning}    \example{
      \exsen{Bogo käba walle ddäddäg gäzag deya.}      \extran{He killed animals on hunts.}}    \example{
      \exsen{Ngäna käba ma ibi allan.}      \extran{I'm going hunting.}}}}

\entry{käbab}{\headword{käbab}
  \pos{Transitive verb}  \sense{\textbf{1.} 
    \definition{to stop}    \example{
      \exsen{Ngämo moko da käbab e dan.}      \extran{I want to stop.}}    \example{
      \exsen{Okbab.}      \extran{[You] stop.}}}  \sense{\textbf{2.} 
    \definition{to stop}    \example{
      \exsen{Näkbabeyo!}      \extran{[You all] stop it!}}}}

\entry{käbädral}{\headword{käbädral}  \note{Trees}  \pronnote{kəbədral}
  \pos{Noun}  \sense{
    \definition{type of tree that grows in the bush with sturdy wood}    \note{Trees}    \example{
}    \example{
}    \example{
}}}

\entry{kabadu}{\headword{kabadu}  \pronnote{kabadu}
  \pos{Noun}  \sense{
    \definition{traditional dish consisting of coconut cream, sago, meat, and tulip greens}}}

\entry{kabag}{\headword{kabag}  \pronnote{kabaɡ}
  \pos{Noun}  \sense{
    \definition{type of thick grass that grows in swamps}}}

\entry{käbäll}{\headword{käbäll}  \note{Trees}  \pronnote{kəbəɽ}
  \pos{Noun}  \sense{
    \definition{type of tree that grows in the bush with soft wood used for canoe paddles}    \note{Trees}    \example{
}    \example{
}    \example{
}}}

\entry{Käbäll}{\headword{Käbäll}
  \variantof{Unspecified Variant of \vartext{Käball}}}

\entry{Käball}{\headword{Käball}
  \pos{Proper noun}  \sense{
    \definition{Kaball (toponym)}}
  \variants{\Unspecified Variant \vartext{Käbäll}}}

\entry{Käballag}{\headword{Käballag}  \note{Dialect names}  \pronnote{kəbaɽaɡ}
  \pos{Proper noun}  \sense{\textbf{1.} 
    \definition{Ende dialect spoken in Käball}    \note{Dialect names}    \example{
}    \example{
}    \example{
}}  \sense{\textbf{2.} 
    \definition{Ende clan with the dog totem}    \note{Dialect names}}}

\entry{käban}{\headword{käban}  \pronnote{kəban}
  \pos{Noun}  \sense{
    \definition{louse}    \example{
      \exsen{Ngämi nazernan gänyme käban namoeaemnalla.}      \extran{We were here, hunting lice.}}    \example{
}    \example{
}}}

\entry{käban käp}{\headword{käban käp}  \pronnote{kəban kəp}
  \pos{Noun}  \sense{
    \definition{nit, louse egg}    \example{
}    \example{
}    \example{
}}}

\entry{kabär}{\headword{kabär}  \note{Types of Taro}  \pronnote{kabər}
  \pos{Noun}  \sense{
    \definition{type of big taro}    \note{Types of Taro}    \example{
}    \example{
}    \example{
}}}

\entry{kabkab}{\headword{kabkab}  \pronnote{kabkab}
  \pos{Noun}  \sense{
    \definition{type of vine-like plant}    \example{
      \exsen{gärep kabkab}      \extran{grapevine}}    \example{
}    \example{
}}}

\entry{kädab}{\headword{kädab}  \note{synthetic form, analytic form unknown}  \pronnote{kədab}
  \pos{Transitive verb}  \sense{\textbf{1.} 
    \definition{to break, split}    \note{synthetic form, analytic form unknown}    \example{
      \exsen{Ngäna nge de dakdab.}      \extran{I split the coconut open.}}    \example{
      \exsen{Mälla da lla de bun dakdaebaebneyo.}      \extran{The women broke the men's skulls.}}}  \sense{\textbf{2.} 
    \definition{to share, split, portion}    \note{synthetic form, analytic form unknown}    \example{
      \exsen{Wayati de nakdab.}      \extran{[You] portion the watermelon.}}    \example{
      \exsen{Ttängäm dakdaeballo ada ttongdae bo gänyan, ttongdae bo gänyan, do tämamae lla da mullae täräp gogmallo.}      \extran{They split the garden like here's one for him, one for him, until everyone got their fair share.}}}}

\entry{kädakäde}{\headword{kädakäde}
  \pos{Noun}  \sense{
    \definition{pieces}    \example{
      \exsen{sana kädakäde}      \extran{pieces of sago}}}}

\entry{Kadawa}{\headword{Kadawa}  \pronnote{kadawa}
  \pos{Proper noun}  \sense{
    \definition{Kadawa (in Kiwai Rural LLG)}    \example{
}    \example{
}    \example{
}}}

\entry{kädbae}{\headword{kädbae}
  \pos{Transitive verb}  \sense{
    \definition{to test, try}    \example{
      \exsen{Ddia midd de kälakälae dakädbaeneyo.}      \extran{Little by little, they tried the deer meat.}}}}

\entry{käde}{\headword{käde}
  \pos{Adverb}  \sense{
}}

\entry{kädebällag mälla}{\headword{kädebällag mälla}  \note{Types of Taro}  \pronnote{kədebəɽaɡ məɽa}
  \pos{Noun}  \sense{
    \definition{type of big taro}    \note{Types of Taro}    \example{
}    \example{
}    \example{
}}}

\entry{kädgal}{\headword{kädgal}  \pronnote{kədɡal}
  \pos{Noun}  \sense{
    \definition{type of tree that grows in the bush}    \example{
}    \example{
}    \example{
}}}

\entry{kädkäd}{\headword{kädkäd}  \pronnote{kədkəd}
  \pos{Transitive verb}  \sense{
    \definition{to remove bark, debark}    \example{
      \exsen{Ngäna put de näkädan.}      \extran{I removed the bark.}}}}

\entry{kädkäd}{\headword{kädkäd}  \pronnote{kədkəd}
  \pos{Noun}  \sense{
    \definition{type of initiation that one must do before you get something from someone}    \example{
      \exsen{Bobzag a kädkäd de oba zime ngasnges allo.}      \extran{Bobzag clan already did their kädkäd.}}}}

\entry{kädkäd}{\headword{kädkäd}  \note{Adjectives}  \pronnote{kədkəd}
  \pos{Modifier}  \sense{
    \definition{cold}    \note{Adjectives}    \example{
      \exsen{Män a kädkädag walle we agäbänan.}      \extran{The girl jumped in the cold water.}}    \example{
      \exsen{Ngämo ttäle da kädkädag agallo.}      \extran{My legs got cold.}}    \example{
}}}

\entry{kae}{\headword{kae}  \note{pronounced as two syllables, not a diphthong. See O01-0234 in the EMKP collection.}
  \pos{Noun}  \sense{
    \definition{a type of plant traditionally chewed and used to welcome people}    \note{pronounced as two syllables, not a diphthong. See O01-0234 in the EMKP collection.}}}

\entry{kae ine}{\headword{kae ine}
  \pos{Noun}  \sense{
    \definition{wine}}}

\entry{kaeg}{\headword{kaeg}  \note{Friendship/social terms of endearment}  \pronnote{kaeɡ}
  \pos{Noun}  \sense{
    \definition{close male friend of the same age who went through initiation at the same time}    \note{Friendship/social terms of endearment}    \example{
}    \example{
}    \example{
}}}

\entry{kaekae}{\headword{kaekae}  \note{Trees}  \pronnote{kaekae}
  \pos{Noun}  \sense{
    \definition{type of small plant with blue and purple flowers and black fruit that cassowaries eat}    \note{Trees}    \example{
}    \example{
}    \example{
}}}

\entry{kaekep}{\headword{kaekep}  \pronnote{kaekep}
  \pos{Transitive verb}  \sense{\textbf{1.} 
    \definition{to chew}    \example{
      \exsen{Gae da nge ingollae dan, be obo täma da kaekep ma dan.}      \extran{The gae palm is like coconut, but its husk is for chewing.}}}  \sense{\textbf{2.} 
    \definition{to struggle to cut (e.g. with a dull knife)}}}

\entry{kaembre}{\headword{kaembre}  \note{Trees}  \pronnote{kaembre}
  \pos{Noun}  \sense{
    \definition{type of tree}    \note{Trees}    \example{
}    \example{
}    \example{
}}}

\entry{kaemne}{\headword{kaemne}  \note{Nature}  \pronnote{kaemne}
  \pos{Noun}  \sense{
    \definition{bee}    \note{Nature}    \example{
      \exsen{Kaemne de ttang alle mällam a mudan, adawatta da bam ttang me naddägän.}      \extran{Don't hold bees with your hand, because they might sting you in the hand.}}    \example{
}    \example{
}}}

\entry{kaen}{\headword{kaen}  \pronnote{kaen}
  \pos{Transitive coverb}  \sense{
    \definition{to wrap up}    \example{
      \exsen{Ngämi käza de net alle kaen eralla, mängallmeny agan zime käza da.}      \extran{We (excl.) wrap up the crocodile with the net; the crocodile is already weak.}}}}

\entry{kaepse}{\headword{kaepse}  \note{Trees}  \pronnote{kaepse}
  \pos{Noun}  \sense{
    \definition{type of tree that grows in the bush with white flowers and big, yellow fruit that cassowaries and deer eat}    \note{Trees}    \example{
}    \example{
}    \example{
}}}

\entry{käg}{\headword{käg}  \pronnote{kəɡ}
  \pos{Noun}  \sense{\textbf{1.} 
    \definition{type of palm with wood used for flooring}    \example{
      \exsen{Llɨg a walle menae me mälla da käg toko me daeya.}      \extran{A boy was beside the water; a woman was on top of the käg palm.}}}  \sense{\textbf{2.} 
    \definition{floor}    \example{
      \exsen{Llɨg a käg me nädämawan.}      \extran{The children sat on the floor.}}}}

\entry{käg}{\headword{käg}  \pronnote{kəɡ}
  \pos{Property noun}  \sense{
    \definition{vague container-like thing}    \example{
      \exsen{sana käg}      \extran{sago container}}    \example{
      \exsen{Däbe llɨg kälsre de däbe käg ik i däzanän.}      \extran{She put that small boy in the container.}}    \example{
      \exsen{Bogo käg me adämenan ekaklle me dɨba katkatre toko dowae me.}      \extran{He sat in the seat on the floor, near the tabletop.}}}}

\entry{käg bänbänang}{\headword{käg bänbänang}  \note{Traditional Arrows}  \pronnote{kəɡ bənbənaŋ}
  \pos{Noun}  \sense{
    \definition{type of spear}    \note{Traditional Arrows}    \example{
}    \example{
}    \example{
}}}

\entry{käg botta}{\headword{käg botta}  \note{Everything in the house}  \pronnote{kəɡ boʈ͡ʂa}
  \pos{Noun}  \sense{
    \definition{horizontal boards on which the floor goes}    \note{Everything in the house}    \example{
}    \example{
}    \example{
}}}

\entry{käg drol}{\headword{käg drol}
  \pos{Noun}  \sense{
    \definition{container}}}

\entry{käg manas}{\headword{käg manas}  \note{Tools}  \pronnote{kəɡ manas}
  \pos{Noun}  \sense{
    \definition{container for squeezing sago}    \note{Tools}    \example{
}    \example{
}    \example{
}}}

\entry{käg tater}{\headword{käg tater}  \note{Everything in the house}  \pronnote{kəɡ tater}
  \pos{Noun}  \sense{
    \definition{flooring, floor mat}    \note{Everything in the house}    \example{
      \exsen{Käg tater de dikomeyo.}      \extran{They brought a floor mat.}}    \example{
}    \example{
}}}

\entry{Kagär}{\headword{Kagär}  \pronnote{kaɡər}
  \pos{Proper noun}  \sense{
    \definition{female personal name}}}

\entry{Kago}{\headword{Kago}
  \pos{Proper noun}  \sense{
    \definition{male personal name}}}

\entry{kak}{\headword{kak}  \lexeme{kak}  \note{Family ties}  \pronnote{kak}
  \pos{Kinship noun}  \sense{\textbf{1.} 
    \definition{grandmother (one's parent's mother; reciprocal)}    \note{Family ties}    \example{
      \exsen{Kak säre zäme ause abal allan.}      \extran{Grandmother is already growing very old.}}    \example{
}}  \sense{\textbf{2.} 
    \definition{grandchild (woman's child's child; reciprocal)}    \note{Family ties}}  \sense{\textbf{3.} 
    \definition{mother-in-law (woman's husband's mother; reciprocal)}    \note{Family ties}}  \sense{\textbf{4.} 
    \definition{daughter-in-law (woman's son's wife; reciprocal)}    \note{Family ties}}}

\entry{käk}{\headword{käk}  \pronnote{kək}
  \pos{Noun}  \sense{
    \definition{bubble}    \example{
      \exsen{Obo enda käk agnan walle mäg me.}      \extran{Something is making bubbles in the creek.}}}}

\entry{kakab}{\headword{kakab}  \pronnote{kakab}
  \pos{Noun}  \sense{\textbf{1.} 
    \definition{leftovers, remainder, remnant}    \example{
      \exsen{Ddob kakab de dowattälleya ttongo ebdo we.}      \extran{They left some leftovers for another day.}}}  \sense{\textbf{2.} 
    \definition{leftover, remaining}    \example{
      \exsen{up wo kakab}      \extran{remaining ripe bananas}}}}

\entry{kakak}{\headword{kakak}
  \pos{Noun}  \sense{
    \definition{nonsingular form of kak}}}

\entry{kakakän}{\headword{kakakän}  \note{Swadesh}  \pronnote{kakakən}
  \pos{Noun}  \sense{\textbf{1.} 
    \definition{(outer) space}    \note{Swadesh}    \example{
      \exsen{Kok a wa piro da ddapall me kakakän giddollnen eran.}      \extran{The moon and stars are in space.}}}  \sense{\textbf{2.} 
    \definition{loose, floating, unbound}    \note{Swadesh}    \example{
      \exsen{nil kakakän}      \extran{loose nail}}    \example{
      \exsen{Kollba da kakakän ibnin allan.}      \extran{The fish is floating on the water.}}    \example{
      \exsen{Ttongdae kui kälsre dae obo känyer kakänkakän gognän.}      \extran{A small island was all alone.}}}  \sense{\textbf{3.} 
    \definition{to float}    \note{Swadesh}}
  \variants{\Unspecified Variant \vartext{kakänkakän}}}

\entry{kakal}{\headword{kakal}  \pronnote{kakal}
  \pos{Intransitive verb}  \sense{\textbf{1.} 
    \definition{to enter, board, go in}    \example{
      \exsen{Bogo gall e gokakalän.}      \extran{She got in the canoe.}}    \example{
      \exsen{Ngäna sɨmell lelang atta llo we okakalan.}      \extran{I was afraid of the pig, so I went in the tree.}}}  \sense{\textbf{2.} 
}}

\entry{kakällät}{\headword{kakällät}  \pronnote{kakəɽət}
  \pos{Transitive verb}  \sense{
    \definition{to split a leaf}}}

\entry{käkäm}{\headword{käkäm}  \pronnote{kəkəm}
  \pos{Noun}  \sense{
    \definition{young leaf}}}

\entry{kakän~}{\headword{kakän~}
  \variantof{freevarlab of \vartext{RED~}}}

\entry{käkan}{\headword{käkan}  \lexeme{käkan}  \pronnote{kəkan}
  \pos{Noun}  \sense{\textbf{1.} 
    \definition{tide}}  \sense{\textbf{2.} 
    \definition{vague amorphous and/or liquid thing}    \example{
      \exsen{bällma käkan, ddapall käkan}      \extran{saliva, cloud}}}}

\entry{kakänkakän}{\headword{kakänkakän}
  \variantof{Unspecified Variant of \vartext{kakakän}}  \pronnote{kakənkakən}}

\entry{käkäp}{\headword{käkäp}  \pronnote{kəkəp}
  \pos{Noun}  \sense{
    \definition{half}    \example{
      \exsen{Mäse ada llo me gomllamalle llo käkäp ada goddälläbnalle.}      \extran{I tried to hold onto the tree, but it broke in half.}}}}

\entry{käkäpyo}{\headword{käkäpyo}  \note{Trees}  \pronnote{kəkəpjo}
  \pos{Noun}  \sense{
    \definition{type of tree that grows in the grassland and savannah with flowers that wallaby and deer eat}    \note{Trees}    \example{
}    \example{
}    \example{
}}}

\entry{Kakaya}{\headword{Kakaya}
  \pos{Proper noun}  \sense{
    \definition{Kakaya (toponym)}}}

\entry{kakayam}{\headword{kakayam}  \pronnote{kakajam}
  \pos{Noun}  \sense{
    \definition{greater bird-of-paradise}    \scientificname{Paradisaea apoda}    \example{
}    \example{
}    \example{
}}}

\entry{Kakayam}{\headword{Kakayam}  \pronnote{kakajam}
  \pos{Proper noun}  \sense{
    \definition{female personal name}}}

\entry{kakep}{\headword{kakep}
  \pos{Noun}  \sense{
    \definition{pain}    \example{
      \exsen{Bam tutu kälngkäl me ttäle kakep a nallan?}      \extran{Are your legs hurting (lit. is the leg pain getting you) from climbing the mountain?}}}}

\entry{Kakeya}{\headword{Kakeya}  \note{Places around Limol}  \pronnote{kakeja}
  \pos{Proper noun}  \sense{
    \definition{Kakeya (bush, camping and sago place of Baba Zi from Upiara; on the road to Bisuaka)}    \note{Places around Limol}    \example{
}    \example{
}    \example{
}}}

\entry{käkle}{\headword{käkle}
  \variantof{Fast Speech Variant of \vartext{kälekäle}}}

\entry{käklläp}{\headword{käklläp}  \note{Gardens}  \pronnote{kəkɽəp}
  \pos{Transitive verb}  \sense{\textbf{1.} 
    \definition{to weed for the second time (easy work to remove the remaining plants)}    \note{Gardens}    \example{
      \exsen{Bogo ttängäm de bäklläpän.}      \extran{He will weed the garden.}}    \example{
      \exsen{Ngämo mokowang melem a klläpnan.}      \extran{My favorite work is easy weeding.}}    \example{
}}  \sense{\textbf{2.} 
    \definition{to sting, be painful}    \note{Gardens}}  \sense{\textbf{3.} 
    \definition{to be amused}    \note{Gardens}}}

\entry{käkllätt}{\headword{käkllätt}  \pronnote{kəkɽəʈ͡ʂ}
  \pos{Transitive verb}  \sense{\textbf{1.} 
    \definition{to weed roughly (leaving bits behind)}    \example{
      \exsen{Ngäna tätäm towall de däkllätt.}      \extran{Yesterday, I roughly weeded the grass.}}}  \sense{\textbf{2.} 
    \definition{to fight, argue}    \example{
      \exsen{Tätäm däkllättnegeya ngämo mäg alle.}      \extran{Yesterday, my mother and I had an argument.}}}}

\entry{kakne}{\headword{kakne}
  \pos{Modifier}  \sense{
    \definition{floating, loose}    \example{
      \exsen{kakne mättnan ma iddpo}      \extran{cloak (lit. 'garment for wearing loose')}}}}

\entry{käkoll}{\headword{käkoll}  \note{Clothing}  \pronnote{kəkoɽ}
  \pos{Noun}  \sense{
    \definition{baby mat}    \note{Clothing}    \example{
}    \example{
}    \example{
}}}

\entry{kakoll}{\headword{kakoll}  \pronnote{kakoɽ}
  \pos{Noun}  \sense{\textbf{1.} 
    \definition{dish}    \example{
      \exsen{Ttongo lla da ttongo za de nanyagaenan kakoll me.}      \extran{A man was stirring something in a dish.}}}  \sense{\textbf{2.} 
    \definition{cup}    \example{
      \exsen{ine nane ma kakoll}      \extran{cup for drinking water}}}}

\entry{kakoll kutt}{\headword{kakoll kutt}  \note{Everyday words and phrases}  \pronnote{kakoɽ kuʈ͡ʂ}
  \pos{Noun}  \sense{
    \definition{endocarp of coconut}    \note{Everyday words and phrases}    \example{
}    \example{
}    \example{
}}}

\entry{Kakos}{\headword{Kakos}  \pronnote{kakos}
  \pos{Proper noun}  \sense{
    \definition{female personal name}}}

\entry{käkpäl}{\headword{käkpäl}  \note{Trees}  \pronnote{kəkpəl}
  \pos{Noun}  \sense{
    \definition{type of tree that grows in the bush; burned to fertilize the ground}    \note{Trees}    \example{
}    \example{
}    \example{
}}}

\entry{käkpäl}{\headword{käkpäl}  \note{Trees}  \pronnote{kəkpəl}
  \pos{Noun}  \sense{
    \definition{sago cooked directly over the fire}    \note{Trees}}}

\entry{kakud}{\headword{kakud}  \lexeme{kakud}  \pronnote{kakud}
  \pos{Noun}  \sense{
    \definition{type of thin, curved, and long yam without thorns}}}

\entry{kala}{\headword{kala}
  \pos{Noun}  \sense{
    \definition{color, pigment, dye}    \example{
      \exsen{Ngäna ag me kala de yu nägnegan.}      \extran{I boiled the pigments in the morning.}}}}

\entry{käla~}{\headword{käla~}
  \variantof{Derived Variant of \vartext{RED~}}}

\entry{kälae}{\headword{kälae}  \pronnote{kəlae}
  \pos{Modifier}  \sense{\textbf{1.} 
    \definition{small, little}    \example{
      \exsen{ttoenttoen kälae}      \extran{short little story}}    \example{
      \exsen{Däbe källäm a ade yäbäd ulle atta dallän do kälae abal gowensegän ine da.}      \extran{That pond became very small from the strong sun.}}}  \sense{\textbf{2.} 
    \definition{little, few}    \example{
      \exsen{kälae eka kutt}      \extran{a few words}}    \example{
      \exsen{Ine da kälae dan.}      \extran{The amount of water is little.}}}}

\entry{kälaekälae}{\headword{kälaekälae}
  \variantof{SpellingVariant of \vartext{kälakälae}}  \pronnote{kəlaekəlae}}

\entry{kälaepot}{\headword{kälaepot}  \pronnote{kəlaepot}
  \pos{Noun}  \sense{
    \definition{tiptoes}    \example{
      \exsen{Kälaepot alle ibi allan.}      \extran{She's tiptoeing.}}}}

\entry{kälakälae}{\headword{kälakälae}
  \pos{Adverb}  \sense{\textbf{1.} 
    \definition{a little}    \example{
      \exsen{Ngäna kälakälae gontämon do yäbäd a angde goklawän dam ngäna ibi de gongkam.}      \extran{I waited a little until the sun set; then I started walking.}}    \example{
      \exsen{Ende eka da kälakälae duli mäzi budabän.}      \extran{The Ende language might get lost in the near future.}}    \example{
      \exsen{Eka kutt a ddob kälakälae sapasapang a dadegaeya.}      \extran{There were other slight differences in their words.}}}  \sense{\textbf{2.} 
    \definition{into pieces}    \example{
      \exsen{Gollab a ekaklle we enanae gollomän kälakälae.}      \extran{The turtle shell got smashed into pieces on the ground.}}}
  \variants{\Fast Speech Variant \vartext{klaklae}}
  \variants{\SpellingVariant \vartext{kälaekälae}}}

\entry{kälakäle}{\headword{kälakäle}  \pronnote{kəlakəle}
  \pos{Intransitive verb}  \sense{
    \definition{to set (of the sun)}    \example{
      \exsen{Yäbäd a goklanän.}      \extran{The sun was setting.}}}}

\entry{Kalamato}{\headword{Kalamato}  \pronnote{kalamato}
  \pos{Proper noun}  \sense{
    \definition{female personal name}}}

\entry{kälämpag}{\headword{kälämpag}  \pronnote{kələmpaɡ}
  \pos{Transitive verb}  \sense{\textbf{1.} 
    \definition{to pad (with cloth or leaves to support a heavy load)}}  \sense{\textbf{2.} 
    \definition{to patch (e.g. roof)}}}

\entry{kälängkäl}{\headword{kälängkäl}
  \variantof{Unspecified Variant of \vartext{kängkäl}}  \pronnote{kələŋkəl}}

\entry{kälas}{\headword{kälas}
  \variantof{SpellingVariant of \vartext{klas}}}

\entry{kälbae}{\headword{kälbae}  \pronnote{kəlbae}
  \pos{Transitive verb}  \sense{
    \definition{to singe (use brief heat to remove hair or down)}    \example{
      \exsen{Mälla da ttägäll baugän bottamänän, mätta de bakälbeaebneyo.}      \extran{The woman will make mumu, finish, and then singe the yams.}}}}

\entry{Kaldon}{\headword{Kaldon}  \pronnote{kaldon}
  \pos{Proper noun}  \sense{
    \definition{male personal name}}}

\entry{käle~}{\headword{käle~}
  \variantof{Plural reduplicant of \vartext{RED~}}}

\entry{kälekäle}{\headword{kälekäle}  \pronnote{kəlekəle}
  \pos{Modifier}  \sense{
    \definition{nonsingular form of kälae}    \example{
      \exsen{ngämi angde käkle dagaeya}      \extran{when we (excl.) were little}}    \example{
      \exsen{Ge nyäng a klekle dag.}      \extran{These bags are small.}}    \example{
      \exsen{Oba llɨg kälekäle tämamae ma me dowattälleyo.}      \extran{They left all their small children at home.}}}
  \variants{\Fast Speech Variant \vartext{käkle, kläkle, klekle}}}

\entry{kalemtoe}{\headword{kalemtoe}  \note{Trees}  \pronnote{kalemtoe}
  \pos{Noun}  \sense{
    \definition{type of tree that grows in the bush with small, long, thin edible nuts that must be broken with stone; nuts are edible when steamed and are also eaten by cassowaries and doves; flesh is used for cream; similar to toe tree}    \note{Trees}    \example{
}    \example{
}    \example{
}}}

\entry{kälepalle}{\headword{kälepalle}  \pronnote{kəlepaɽe}
  \pos{Manner adverb}  \sense{
    \definition{slowly}    \example{
      \exsen{Ngäna kälepalle danyroene, angde dowae mäse däga, ddia da  dindugän.}      \extran{I was creeping slowly, and when I almost had it, the deer fled.}}    \example{
}    \example{
}}
  \variants{\Fast Speech Variant \vartext{kälpalle}}}

\entry{kälkäl}{\headword{kälkäl}  \pronnote{kəlkəl}
  \pos{Transitive verb}  \sense{
    \definition{to lie, fabricate}    \example{
      \exsen{Ngäna obo pallall kälkäl eran.}      \extran{I'm making up stories about her.}}}}

\entry{kalkmo}{\headword{kalkmo}  \note{Parts of the body}  \pronnote{kalkmo}
  \pos{Noun}  \sense{
    \definition{joints}    \note{Parts of the body}    \example{
}    \example{
}    \example{
}}}

\entry{käll}{\headword{käll}  \lexeme{käll}  \pronnote{kəɽ}
  \pos{Noun}  \sense{
    \definition{spleen}    \example{
}}}

\entry{källa}{\headword{källa}  \pronnote{kəɽa}
  \pos{Noun}  \sense{\textbf{1.} 
    \definition{feces, poop, waste}    \example{
      \exsen{Däräng källa de ngätt att de baglloeabnalla.}      \extran{We will take dog poop out from the yards.}}    \example{
      \exsen{Däräng a yuwog ngätt me källa bebeyag dan.}      \extran{The dog takes craps in many yards.}}    \example{
}}  \sense{\textbf{2.} 
    \definition{intestines, bowels, guts, innards (of an animal)}    \example{
      \exsen{Sɨmell de enanae dukolldänegän idoidog alle, källa de tämamae däträpnegän.}      \extran{He shot the pig dead with the spear; its guts were all cut open.}}}  \sense{\textbf{3.} 
    \definition{vague amorphous and/or soft thing, lump (e.g. cloud, fish tail, lump of earwax)}}}

\entry{kalla}{\headword{kalla}
  \pos{Intransitive verb}  \sense{
    \definition{to lie down}    \example{
      \exsen{Ngäna obom dangkallängg.}      \extran{I made him lie down.}}}}

\entry{källa mit}{\headword{källa mit}  \pronnote{kəɽa mit}
  \pos{Noun}  \sense{
    \definition{bottom of fence}    \example{
}    \example{
}    \example{
}}}

\entry{kallag}{\headword{kallag}
  \variantof{Unspecified Variant of \vartext{kalläg}}}

\entry{kalläg}{\headword{kalläg}  \pronnote{kaɽəɡ}
  \pos{Noun}  \sense{
    \definition{type of edible, fatty fish with big scales; found in the swamp, similar to barramundi}}
  \variants{\Unspecified Variant \vartext{kallag}}}

\entry{Källäk}{\headword{Källäk}  \pronnote{kəɽək}
  \pos{Proper noun}  \sense{
    \definition{Kallak (toponym)}}}

\entry{kallakalla}{\headword{kallakalla}
  \pos{Noun}  \sense{
}}

\entry{källakällae}{\headword{källakällae}  \pronnote{kəɽakəɽae}
  \pos{Modifier}  \sense{
    \definition{hospitable}    \example{
}    \example{
}}}

\entry{källakällawang}{\headword{källakällawang}
  \variantof{Unspecified Variant of \vartext{källakällong}}}

\entry{källakälle}{\headword{källakälle}  \note{Things a person can do}  \pronnote{kəɽakəɽe}
  \pos{Transitive verb}  \sense{\textbf{1.} 
    \definition{to poison the river}    \note{Things a person can do}    \example{
      \exsen{walle kllaklle}      \extran{to poison the river}}    \example{
}    \example{
}}  \sense{\textbf{2.} 
    \definition{to scrape food off the fire, e.g. banana, taro, yam}    \note{Things a person can do}}}

\entry{källakällong}{\headword{källakällong}  \pronnote{kəɽakoɽoŋ}
  \pos{Adverb}  \sense{\textbf{1.} 
    \definition{while pooping; constantly needing to poop}}  \sense{\textbf{2.} 
    \definition{traveling urgently, arduously, or nonstop as if to a toilet}    \example{
      \exsen{Bogo källkällong dindu allanǃ}      \extran{He's running like he needs to poop!}}    \example{
      \exsen{Ubi tätäm ddia de kapu digeyo källakällong ma we.}      \extran{Yesterday, they carried the deer all the way home.}}}
  \variants{\Unspecified Variant \vartext{källakällawang}}}

\entry{källäll}{\headword{källäll}
  \variantof{Unspecified Variant of \vartext{kɨllɨll}}}

\entry{källäm}{\headword{källäm}  \note{Water}  \pronnote{kəɽəm}
  \pos{Noun}  \sense{
    \definition{pond; lagoon}    \note{Water}    \example{
      \exsen{Kottllam a wa kollba da ttongo källäm me dazernän.}      \extran{Turtles and fish were living in a pond.}}    \example{
}}
  \variants{\Unspecified Variant \vartext{kɨllɨm}}}

\entry{källama kup}{\headword{källama kup}  \note{Everything in the house}  \pronnote{kəɽama kup}
  \pos{Noun}  \sense{
    \definition{toilet (lit. hole for pooping)}    \note{Everything in the house}    \example{
}    \example{
}    \example{
}}}

\entry{kallamatt}{\headword{kallamatt}  \pronnote{kaɽamaʈ͡ʂ}
  \pos{Noun}  \sense{
    \definition{Oriental dollarbird}    \scientificname{Eurystomus orientalis waigiouensis}    \example{
}    \example{
}    \example{
}}}

\entry{källamitang}{\headword{källamitang}  \pronnote{kəɽamitaŋ}
  \pos{Noun}  \sense{
    \definition{the first stage of sago growth in which the root has developed following planting and the plant is ~1m tall}    \example{
      \exsen{Obo sana da källamitang gogon.}}}}

\entry{källän}{\headword{källän}  \note{Clothing}  \pronnote{kəɽən}
  \pos{Noun}  \sense{\textbf{1.} 
    \definition{belt}    \note{Clothing}    \example{
      \exsen{Zon bo gablle mättnan ma källän a ddäddäg ttoe dädär alle ngasngesatt daeya.}      \extran{John's belt for wearing around his waist was made from dry animal hide.}}    \example{
}    \example{
}}  \sense{\textbf{2.} 
    \definition{waist}    \note{Clothing}}}

\entry{källäng}{\headword{källäng}  \note{Animals/Plants}  \pronnote{kəɽəŋ}
  \pos{Noun}  \sense{
    \definition{type of animal}    \note{Animals/Plants}    \example{
}    \example{
}    \example{
}}}

\entry{kallänggag}{\headword{kallänggag}  \note{Stages of life}  \pronnote{kaɽəŋɡaɡ}
  \pos{Noun}  \sense{
    \definition{stage of life, when a baby is still lying down.}    \note{Stages of life}    \example{
}    \example{
}    \example{
}}}

\entry{Källängmäll}{\headword{Källängmäll}  \note{Places around Limol}  \pronnote{kəɽəŋməɽ}
  \pos{Proper noun}  \sense{
    \definition{sago place near Old Kibobma}    \note{Places around Limol}    \example{
}    \example{
}    \example{
}}}

\entry{kalläntäg}{\headword{kalläntäg}  \pronnote{kaɽəntəɡ}
  \pos{Transitive verb}  \sense{
    \definition{to split}    \example{
      \exsen{Bogo mätta ulleulle de deyangkälltägalle.}      \extran{She split the big yams.}}}}

\entry{kallat}{\headword{kallat}  \pronnote{kaɽat}
  \pos{}  \sense{
    \definition{burp}}}

\entry{källatolma}{\headword{källatolma}  \note{Numbers}  \pronnote{kəɽatolma}
  \pos{Noun}  \sense{\textbf{1.} 
    \definition{middle finger}    \note{Numbers}}  \sense{\textbf{2.} 
    \definition{three (lit. middle finger; body counting numeral)}    \note{Numbers}    \example{
}    \example{
}    \example{
}}}

\entry{källayoyo}{\headword{källayoyo}  \note{Trees}  \pronnote{kəɽajojo}
  \pos{Noun}  \sense{
    \definition{type of tree that grows in the bush with leaves used as toilet paper}    \note{Trees}    \example{
}    \example{
}    \example{
}}}

\entry{kallekalle}{\headword{kallekalle}  \lexeme{kallekalle}  \pronnote{kaɽekaɽe}
  \pos{Noun}  \sense{
    \definition{type of yam with a white interior, red skin, and hairs}}}

\entry{Källid}{\headword{Källid}
  \pos{Proper noun}  \sense{
    \definition{female personal name}}}

\entry{kallɨngg}{\headword{kallɨngg}
  \pos{Transitive verb}  \sense{
    \definition{to kill}    \example{
      \exsen{Sɨmell de mɨnyi bangkallɨgiyu.}      \extran{They will kill the pig.}}}}

\entry{källkae}{\headword{källkae}
  \pos{Noun}  \sense{\textbf{1.} 
    \definition{future}    \example{
      \exsen{Ende tän a ako ngänam eran ada llɨg klekle a komuniti da bo källkae dag.}      \extran{The Ende tribe also understands that children are the community's future.}}}  \sense{\textbf{2.} 
    \definition{later, in the future}    \example{
      \exsen{Aska källkae ngämo llɨg a mɨnyi pällämpälläm eka de bepänyneyo.}      \extran{I think my children will speak in English in the future.}}}}

\entry{kallkäll}{\headword{kallkäll}  \note{Adjectives}  \pronnote{kaɽkəɽ}
  \pos{Noun}  \sense{
    \definition{cold}    \note{Adjectives}    \example{
      \exsen{Kallkäll ngänäm dagän.}      \extran{I got cold (lit. the cold got me).}}    \example{
      \exsen{Gongkäbägän, abo kallkäll atta guiringänän.}      \extran{He dived, then got out because of the cold.}}}}

\entry{källkäll}{\headword{källkäll}
  \variantof{Dialectal Variant of \vartext{kɨllkɨll}}  \pronnote{kəɽkəɽ}}

\entry{kallkäll gllaengg}{\headword{kallkäll gllaengg}  \pronnote{kaɽkəɽ ɡɽaeŋɡ}
  \pos{Intransitive compound verb}  \sense{
    \definition{to warm up, get warm, warm oneself}    \example{
      \exsen{Bogo yu me kallkäll gonggllaemenynän.}      \extran{He was warming up by the fire.}}    \example{
      \exsen{Da ngäna ddone yu bipellknegalle bablle, bongo ddone gänya mer yu mi kallkäll bonggllaegalle.}      \extran{If I hadn't chopped wood for you, you couldn't get warm by this nice fire.}}}}

\entry{kallkäll sod}{\headword{kallkäll sod}  \note{Clothing}  \pronnote{kaɽkəɽ sod}
  \pos{Noun}  \sense{
    \definition{long-sleeve shirt; coat}    \note{Clothing}    \example{
}    \example{
}    \example{
}}}

\entry{kallkell}{\headword{kallkell}  \pronnote{kaɽkeɽ}
  \pos{Transitive verb}  \sense{
    \definition{to deleaf a plant to reveal the shoot, or to take the skin off}}}

\entry{källmakällme}{\headword{källmakällme}  \pronnote{kəɽmakəɽme}
  \pos{Intransitive verb}  \sense{
    \definition{to survive}    \example{
      \exsen{Da ttongo kantri me lla da bompallängkmenynegän obaoba komlla kullum e, obaoba däbe kantri da mɨnyi ddone bakällmewän.}      \extran{If the people in a country divide themselves into two groups, that country will not survive.}}}}

\entry{kallmo}{\headword{kallmo}  \note{Birds}  \pronnote{kaɽmo}
  \pos{Noun}  \sense{
    \definition{butcherbird (black-backed, hooded)}    \scientificname{Cracticus mentalis, C. cassicus}    \note{Birds}    \example{
}    \example{
}    \example{
}}}

\entry{Källnyam}{\headword{Källnyam}
  \pos{Proper noun}  \sense{
    \definition{male personal name}}}

\entry{källpalle}{\headword{källpalle}
  \variantof{SpellingVariant of \vartext{kälpalle}}  \pronnote{kəɽpaɽe}}

\entry{Källtae}{\headword{Källtae}
  \pos{Proper noun}  \sense{
    \definition{female personal name}}
  \variants{\SpellingVariant \vartext{Källtai}}}

\entry{Källtai}{\headword{Källtai}
  \variantof{SpellingVariant of \vartext{Källtae}}}

\entry{kalltam}{\headword{kalltam}  \pronnote{kaɽtam}
  \pos{Transitive verb}  \sense{
    \definition{to split}    \example{
      \exsen{Keta da sanador kalltaematt dan.}      \extran{A roof post is a split sago stalk.}}}}

\entry{källttakälltte}{\headword{källttakälltte}  \pronnote{kəɽʈ͡ʂakəɽʈ͡ʂe}
  \pos{Transitive verb}  \sense{
    \definition{to cast away, expel}    \example{
      \exsen{Yesu gagäll anyke de däkällttawän lla bo guwo watt de.}      \extran{Jesus cast the bad spirit out from the man's heart.}}}}

\entry{kalmoe}{\headword{kalmoe}  \pronnote{kalmoe}
  \pos{Modifier}  \sense{
    \definition{pliable, bendable}}}

\entry{kalmokalmoe}{\headword{kalmokalmoe}  \note{good characteristics of people}  \pronnote{kalmokalmoe}
  \pos{Modifier}  \sense{
    \definition{welcoming}    \note{good characteristics of people}    \example{
      \exsen{Ngämo yae lla kalmokalmoe deya.}      \extran{My mother was a welcoming person.}}    \example{
}    \example{
}}}

\entry{kälngkäl}{\headword{kälngkäl}
  \variantof{Unspecified Variant of \vartext{kängkäl}}}

\entry{Kälnyam}{\headword{Kälnyam}
  \pos{Proper noun}  \sense{
    \definition{male personal name}}}

\entry{kalokalo}{\headword{kalokalo}  \pronnote{kalokalo}
  \pos{Modifier}  \sense{
    \definition{pliable, flexible, bendable}}}

\entry{kälpalle}{\headword{kälpalle}
  \variantof{Fast Speech Variant of \vartext{kälepalle}}  \pronnote{kəlpaɽe}
  \variants{\SpellingVariant \vartext{källpalle}}}

\entry{kälsäre}{\headword{kälsäre}  \note{Adjectives}  \pronnote{kəlsəre}
  \pos{Modifier}  \sense{
    \definition{small, little}    \note{Adjectives}    \example{
      \exsen{Ge mätta da käläsre dan be nge da ulle dan.}      \extran{This yam is smaller than the coconut.}}    \example{
}    \example{
}}
  \variants{\Fast Speech Variant \vartext{kälsre}}}

\entry{kälsre}{\headword{kälsre}
  \variantof{Fast Speech Variant of \vartext{kälsäre}}}

\entry{kälsrong}{\headword{kälsrong}  \pronnote{kəlsroŋ}
  \pos{Modifier}  \sense{
    \definition{nice singing voice}}}

\entry{kaltakaltamang}{\headword{kaltakaltamang}  \note{Traditional Arrows}  \pronnote{kaltakaltamaŋ}
  \pos{Noun}  \sense{
    \definition{type of spear}    \note{Traditional Arrows}    \example{
}    \example{
}    \example{
}}}

\entry{käm}{\headword{käm}  \note{Parts of the body}  \pronnote{kəm}
  \pos{Noun}  \sense{\textbf{1.} 
    \definition{stomach}    \note{Parts of the body}    \example{
      \exsen{Ngämo käm a ulle agan.}      \extran{I'm full (lit. my stomach got big).}}    \example{
      \exsen{Sɨmell bo käm de dupotän.}      \extran{He sliced the pig's stomach open.}}    \example{
      \exsen{Kämang mälla da gänya itrell peyang a mänyi oba käm me llɨg abira basiaemneyo.}      \extran{Pregnant women with this disease will give it to their children in their womb.}}}  \sense{\textbf{2.} 
    \definition{womb}    \note{Parts of the body}    \example{
      \exsen{Kämang mälla da gänya itrell peyang a mänyi oba käm me llɨg abira basiaemneyo.}      \extran{Pregnant women with this disease will give it to their children in their womb.}}}}

\entry{käm}{\headword{käm}  \pronnote{kəm}
  \pos{Noun}  \sense{
    \definition{love, enjoyment}    \example{
      \exsen{eka kämang}      \extran{talkative}}    \example{
      \exsen{Ddobae togotogol tongoe kämang dagwaeya.}      \extran{They (du.) really loved to play hide-and-seek.}}}}

\entry{kam}{\headword{kam}  \note{Water}  \pronnote{kam}
  \pos{Transitive verb}  \sense{\textbf{1.} 
    \definition{to start, begin}    \note{Water}    \example{
      \exsen{kam ma ngättma}      \extran{place for starting (i.e. starting line)}}    \example{
      \exsen{Yogoll a gongkamän.}      \extran{The rain started.}}    \example{
      \exsen{Obo ma we ibi we kam allan.}      \extran{He is starting to go home.}}}  \sense{\textbf{2.} 
    \definition{to start}    \note{Water}    \example{
      \exsen{Edeb a ine kämbmeny de dängkamän.}      \extran{Edeb started diving in the water.}}}}

\entry{kam}{\headword{kam}  \pronnote{kam}
  \pos{Transitive verb}  \sense{
    \definition{to cut (e.g. meat, skin)}    \example{
      \exsen{Ngäna ngämo tupi di nakaman.}      \extran{I cut my index finger.}}}}

\entry{Käm}{\headword{Käm}
  \pos{Proper noun}  \sense{
    \definition{female personal name}}}

\entry{käm}{\headword{käm}  \pronnote{kəm}
  \pos{Locational}  \sense{
    \definition{underneath, under, beneath}    \example{
      \exsen{Mälla da ma käm me dan.}      \extran{The woman is beneath the house.}}}}

\entry{käma}{\headword{käma}
  \variantof{Fast Speech Variant of \vartext{ngasekäma}}  \pronnote{kəma}}

\entry{kämag}{\headword{kämag}  \note{Weather}  \pronnote{kəmaɡ}
  \pos{Modifier}  \sense{\textbf{1.} 
    \definition{west, western}    \note{Weather}    \example{
      \exsen{kämag ingong}      \extran{a dance from the west}}}  \sense{\textbf{2.} 
    \definition{round dance with singers and kundu drum in the middle}    \note{Weather}}  \sense{\textbf{3.} 
    \definition{west wind; windy storm from the west}    \note{Weather}    \example{
      \exsen{Tɨtɨm kämag peyang daeya.}      \extran{Yesterday, there was a windy storm.}}}  \sense{\textbf{4.} 
    \definition{season characterized by windy storms from the west (first season; corresponds to January)}    \note{Weather}    \example{
      \exsen{Kämag amne me, bängnameyo abo ada mätta da zäme redi allan.}      \extran{In the middle of kämag, they will understand that the yams are ready.}}    \example{
      \exsen{Kämag a mätta polle me mondre täräp dan.}}}}

\entry{kämag ma}{\headword{kämag ma}
  \variantof{SpellingVariant of \vartext{kämagma}}}

\entry{kämag wel}{\headword{kämag wel}  \pronnote{kəmaɡ wel}
  \pos{Noun}  \sense{
    \definition{Big winds in January, a big wind from the west that blows trees down}    \example{
      \exsen{Kämag wel a era za gagällang dan a ai dan lla de ade bäbäddän.}      \extran{This big wind is a bad thing and it can kill people.}}    \example{
}    \example{
}}}

\entry{kämagma}{\headword{kämagma}  \pronnote{kəmaɡma}
  \pos{Noun}  \sense{
    \definition{west}}
  \variants{\SpellingVariant \vartext{kämag ma}}}

\entry{kamäkamät}{\headword{kamäkamät}  \note{Trees}  \pronnote{kaməkamət}
  \pos{Noun}  \sense{
    \definition{type of big tree that grows in the bush with yellow flowers and big yellow fruit that are collected}    \note{Trees}    \example{
}    \example{
}    \example{
}}}

\entry{kämakäme}{\headword{kämakäme}  \pronnote{kəm}
  \pos{Transitive verb}  \sense{
    \definition{to heal}    \example{
      \exsen{Itrellang lla de omäg alle dakämeyo.}      \extran{They healed the sick man with magic.}}}}

\entry{=kämall}{\headword{=kämall}
  \pos{Nominal enclitic}  \sense{
    \definition{comitative case clitic; with (only used after oba, ama, and nouns followed by =aba)}    \example{
      \exsen{Deyarän däba ine nanen e oba kämall.}      \extran{He went there to drink water with them.}}    \example{
      \exsen{Ge lla da ngäna ama kämall skul att dan.}      \extran{These people are who I went to school with.}}    \example{
      \exsen{Ngämi gobllab Eramang gall tapma we ddob llayaba kämall.}      \extran{We (excl.) went to Eramang canoe dock with some other people.}}}
  \variants{\Dialectal Variant \vartext{=kämäll}}}

\entry{=kämäll}{\headword{=kämäll}
  \variantof{Dialectal Variant of \vartext{=kämall}}}

\entry{kamän}{\headword{kamän}  \pronnote{kamən}
  \pos{Transitive verb}  \sense{\textbf{1.} 
    \definition{to kill at once}    \example{
      \exsen{Ngäna ddäddäg de näkamänan.}      \extran{I killed the animal straight off.}}}  \sense{\textbf{2.} 
    \definition{to pick (e.g. mushrooms)}}}

\entry{käman}{\headword{käman}  \pronnote{kəman}
  \pos{Noun}  \sense{
    \definition{traditional type of cassava}}}

\entry{kämang}{\headword{kämang}  \note{Birth}  \pronnote{kəmaŋ}
  \pos{Modifier}  \sense{
    \definition{pregnant}    \note{Birth}    \example{
      \exsen{Da män a kämang allan, ede llɨg a zeg allan.}      \extran{If a woman gets preganant, a child will be born.}}    \example{
}    \example{
}}}

\entry{kämany}{\headword{kämany}  \note{Friendship/social terms of endearment}  \pronnote{kəmaɲ}
  \pos{Noun}  \sense{
    \definition{type of friendship}    \note{Friendship/social terms of endearment}    \example{
}    \example{
}    \example{
}}}

\entry{kämätt}{\headword{kämätt}  \note{Parts of the Body}  \pronnote{kəməʈ͡ʂ}
  \pos{Noun}  \sense{
    \definition{testicle}    \note{Parts of the Body}    \example{
}    \example{
}    \example{
}}}

\entry{kambäg}{\headword{kambäg}  \pronnote{kambəɡ}
  \pos{Intransitive verb}  \sense{
    \definition{to escape, flee}    \example{
      \exsen{Ngäna gänyeri angkebägan.}      \extran{I fled this way.}}}}

\entry{kämbäg}{\headword{kämbäg}  \note{Diving}  \pronnote{kəmbəɡ}
  \pos{Intransitive verb}  \sense{\textbf{1.} 
    \definition{to dive}    \note{Diving}    \example{
      \exsen{Ngämi ine kämbmeny de dängkameya.}      \extran{We (excl.) started diving.}}    \example{
      \exsen{Wiya angkäbäg!}      \extran{[You] come dive!}}    \example{
}}  \sense{\textbf{2.} 
    \definition{to baptize}    \note{Diving}    \example{
      \exsen{Zon ine kämbägag}      \extran{John the Baptist}}    \example{
      \exsen{Pasta da obom ine dängkäbägän dedme.}      \extran{The pastor baptized him there.}}}}

\entry{kame}{\headword{kame}  \pronnote{kame}
  \pos{Noun}  \sense{\textbf{1.} 
    \definition{ignorance, non-knowing; incomprehension, non-understanding}    \example{
      \exsen{Ende eka kame lla}      \extran{person who doesn't speak Ende}}    \example{
      \exsen{Obo kame dan amom bällädän.}      \extran{She doesn't know (lit. her ignorance is) whom she will marry.}}    \example{
      \exsen{Ngäma kame da ada käza da daden walle ik mi.}      \extran{We (excl.) don't realize (lit. it's our ignorance) that there's a crocodile in the water.}}    \example{
      \exsen{Ngämo kame dan ematta bongo mikutt alle.}      \extran{I don't see (lit. my incomprehension is) why you're mad.}}}  \sense{\textbf{2.} 
    \definition{unknown, unfamiliar}    \example{
      \exsen{kame lla}      \extran{stranger}}    \example{
      \exsen{Da kame ttängämang lla de bognegän, ede ngäna mɨnyi Ingglis eka walle bäntameny.}      \extran{If it's a man from an unfamiliar village, then I will converse using English.}}    \example{
}}}

\entry{kame}{\headword{kame}  \pronnote{kame}
  \pos{Temporal adverb}  \sense{
    \definition{again}    \example{
      \exsen{Kame ge ttoen bongo ddone mɨnyi nanges.}      \extran{You must not do this thing again.}}    \example{
      \exsen{Kumuddäga we ebdo me bogo mɨnyi kame ttam bogon kuddäll atta.}      \extran{In three days, he will rise from the dead and live again.}}}}

\entry{kame eka}{\headword{kame eka}  \pronnote{kame eka}
  \pos{Verb}  \sense{
    \definition{focus}}}

\entry{käme näkäp me ddäganen}{\headword{käme näkäp me ddäganen}  \note{School}  \pronnote{kəme nəkəp me ɖ͡ʐəɡanen}
  \pos{Transitive verb}  \sense{
    \definition{to memorize}    \note{School}    \example{
}    \example{
}    \example{
}}}

\entry{kame tameny ma skul}{\headword{kame tameny ma skul}  \note{School}  \pronnote{kame tameɲ ma skul}
  \pos{}  \sense{
    \definition{open distance school}    \note{School}    \example{
}    \example{
}    \example{
}}}

\entry{kamebi}{\headword{kamebi}
  \variantof{SpellingVariant of \vartext{kemibi}}}

\entry{kameka}{\headword{kameka}
  \pos{Verb}  \sense{
    \definition{to ignore}    \example{
      \exsen{Kameka da mudan.}      \extran{Listen up (lit. Don't ignore!).}}}}

\entry{kamekame}{\headword{kamekame}  \pronnote{kamekame}
  \pos{Adverb}  \sense{
    \definition{unknowingly, mistakenly, wrongly}    \example{
      \exsen{Ngämi kamekame ada ka up a walle ik mi dan.}      \extran{We (excl.) mistakenly thought the banana was inside the water.}}}}

\entry{kamekong}{\headword{kamekong}  \pronnote{kamekoŋ}
  \pos{Modifier}  \sense{
    \definition{busy}    \example{
      \exsen{Mosen da angde kamekong gogon kaptte gllämnan me, mänyan da bi komllaebe ada tongoenen me kamekong gogeyo.}      \extran{While the eldest was busy washing clothes, the two younger ones were busy playing.}}}}

\entry{kameny}{\headword{kameny}
  \pos{Noun}  \sense{
    \definition{absence}    \example{
      \exsen{Ngämo kameny me de ngämo masar a gagäll gogon.}      \extran{My grandfather passed away in my absence (i.e. before I was born).}}}}

\entry{kämgag}{\headword{kämgag}  \pronnote{kəmɡaɡ}
  \pos{Noun}  \sense{
    \definition{type of lizard}    \example{
}    \example{
}    \example{
}}}

\entry{kämgall}{\headword{kämgall}  \pronnote{kəmɡaɽ}
  \pos{Noun}  \sense{
    \definition{underneath}}}

\entry{kamibi}{\headword{kamibi}
  \variantof{SpellingVariant of \vartext{kemibi}}}

\entry{kämlla}{\headword{kämlla}  \lexeme{kämlla}  \pronnote{kəmɽa}
  \pos{Noun}  \sense{
    \definition{short-beaked echidna}    \scientificname{Tachyglossus aculeatus}    \example{
      \exsen{Kämlla da täkäll peyang dan.}      \extran{The echidna has spines.}}}}

\entry{kämo}{\headword{kämo}
  \variantof{Unspecified Variant of \vartext{komo}}  \note{Poison}  \pronnote{kəmo}}

\entry{kamo}{\headword{kamo}  \note{Friendship/social terms of endearment}  \pronnote{kamo}
  \pos{Noun}  \sense{
    \definition{reciprocal term for the young man and the older man that takes him through initiation}    \note{Friendship/social terms of endearment}    \example{
}    \example{
}    \example{
}}}

\entry{kämoe}{\headword{kämoe}
  \pos{Noun}  \sense{
}}

\entry{kämongg}{\headword{kämongg}
  \pos{Transitive verb}  \sense{
    \definition{to get from someone}}}

\entry{kampani}{\headword{kampani}
  \pos{Noun}  \sense{
    \definition{company}}}

\entry{kampentri}{\headword{kampentri}
  \pos{Noun}  \sense{
    \definition{carpentry}}}

\entry{kämser käpang}{\headword{kämser käpang}  \note{Traditional Arrows}  \pronnote{kəmser kəpaŋ}
  \pos{Noun}  \sense{
    \definition{type of arrow}    \note{Traditional Arrows}    \example{
}    \example{
}    \example{
}}}

\entry{kämsir}{\headword{kämsir}  \note{Trees}  \pronnote{kəmsir}
  \pos{Noun}  \sense{
    \definition{type of tree the grows in the bush with fruit that is black outside and green and red inside and is eaten by cassowaries; similar to sir tree fruit}    \note{Trees}    \example{
}    \example{
}    \example{
}}}

\entry{kämsir budar}{\headword{kämsir budar}  \note{Grubs}  \pronnote{kəmsir budar}
  \pos{Noun}  \sense{
    \definition{type of grub}    \note{Grubs}    \example{
}    \example{
}    \example{
}}}

\entry{kämtupi}{\headword{kämtupi}  \note{Animals/Plants}  \pronnote{kəmtupi}
  \pos{Noun}  \sense{
    \definition{type of ginger}    \note{Animals/Plants}    \example{
}    \example{
}    \example{
}}}

\entry{kamuka}{\headword{kamuka}  \note{Trees}  \pronnote{kamuka}
  \pos{Noun}  \sense{
    \definition{type of cultivated thorny citrus tree with small white flowers and softball-size fruit with thick green skin}    \note{Trees}    \example{
}    \example{
}    \example{
}}}

\entry{kän}{\headword{kän}  \pronnote{kən}
  \pos{Transitive verb}  \sense{\textbf{1.} 
    \definition{to withdraw, come out}    \example{
      \exsen{Ngäna skul atta gongkän greid paeb me.}      \extran{I withdrew from school in grade five.}}}  \sense{\textbf{2.} 
    \definition{to remove, take out, take off, undo}    \example{
      \exsen{Bäne sod de nängkän.}      \extran{Take off your shirt.}}    \example{
      \exsen{Ngäna net de apte tär nängkänan.}      \extran{I undid the string on one side of the net.}}    \example{
      \exsen{Gänyme nge kopnen ma patt de aeya dängkänän?}      \extran{Who took out the coconut dehusker (from the ground) here?}}}}

\entry{kän}{\headword{kän}  \lexeme{kän}  \pronnote{kən}
  \pos{Noun}  \sense{
    \definition{type of big, round yam with a white interior and thorns}}}

\entry{känaebag}{\headword{känaebag}
  \variantof{Dialectal Variant of \vartext{känazbag}}  \pronnote{kənaebaɡ}}

\entry{känakone}{\headword{känakone}
  \variantof{Unspecified Variant of \vartext{konakone}}}

\entry{känameny}{\headword{känameny}
  \pos{Transitive verb}  \sense{
}}

\entry{känär}{\headword{känär}  \note{Grubs}  \pronnote{kənər}
  \pos{Noun}  \sense{\textbf{1.} 
    \definition{type of edible grub found in the bush}    \note{Grubs}    \example{
}    \example{
}    \example{
}}  \sense{\textbf{2.} 
    \note{Grubs}}}

\entry{kanas}{\headword{kanas}  \note{Traditional Arrows}  \pronnote{kanas}
  \pos{Noun}  \sense{
    \definition{type of basic arrow}    \note{Traditional Arrows}    \example{
      \exsen{Masamasar täräp me ngattong kanas toboll de dokomneyo.}      \extran{In the old times, people first carried kanas spears.}}    \example{
}    \example{
}}}

\entry{kanas ma}{\headword{kanas ma}  \pronnote{kanas ma}
  \pos{Noun}  \sense{
    \definition{type of pointed roof not found in Limol}}}

\entry{känazbag}{\headword{känazbag}  \pronnote{kənazbaɡ}
  \pos{Temporal adverb}  \sense{
    \definition{tomorrow}    \example{
      \exsen{Abo känazbag ikop bogmam.}      \extran{We will see each other tomorrow.}}}
  \variants{\Dialectal Variant \vartext{känaebag}}}

\entry{kandärmang}{\headword{kandärmang}  \pronnote{kandərmaŋ}
  \pos{Modifier}  \sense{\textbf{1.} 
    \definition{sorry, apologetic, regretful}    \example{
      \exsen{Kandärmang nägagallo.}      \extran{They were feeling sorry for them.}}    \example{
      \exsen{Kandärmang allan ngäna bäne pate, ddone mulldae dan.}      \extran{I feel sorry towards you, but I'm unable.}}    \example{
      \exsen{Bogo kandärmang me dan.}      \extran{She's feeling regretful.}}}  \sense{\textbf{2.} 
    \definition{sorry}}  \sense{\textbf{3.} 
    \definition{pity, sympathy}    \example{
      \exsen{Kandärmang a era ulle ttoen dan.}      \extran{Sympathy is an important habit.}}}  \sense{\textbf{4.} 
    \definition{sorry, pitiful, unfortunate, sad}    \example{
      \exsen{Ge lla ddobae kandärmang dan, bogo ekameny dan, llan ttomoll dan.}      \extran{This man is very unfortunate; he is mute and deaf.}}    \example{
      \exsen{Kandärmang, llamda da bälle gazen ma da ddone mullae gogon.}      \extran{Sadly, the old man didn't have a way to get out.}}    \example{
      \exsen{Emi ngänaeka gogon kandärmang ingoll gogon.}      \extran{Emi cried pitifully.}}}}

\entry{kandärmang}{\headword{kandärmang}  \pronnote{kandərmaŋ}
  \pos{Noun}  \sense{
    \definition{present, gift}    \example{
      \exsen{Ngämo kandärmang a gänyan bablle.}      \extran{My present for you is here.}}    \example{
      \exsen{Eso ulle bänene kandärmang de.}      \extran{Thanks a lot for your gift.}}}}

\entry{kandärmang eka}{\headword{kandärmang eka}
  \pos{Noun}  \sense{
    \definition{apology}    \example{
      \exsen{Ngäna bänene kandärmang eka de malam eran.}      \extran{I accept your apology.}}}}

\entry{Kanengga}{\headword{Kanengga}
  \pos{Proper noun}  \sense{
    \definition{male personal name}}}

\entry{kang}{\headword{kang}  \pronnote{kaŋ}
  \pos{Noun}  \sense{
    \definition{sucker (additional unwanted shoot that grows by the base of a tree)}    \example{
      \exsen{Kak sana kang de ttaengän eran.}      \extran{Grandmother pulls a sago sucker.}}}}

\entry{Kange}{\headword{Kange}  \pronnote{kaŋe}
  \pos{Proper noun}  \sense{
    \definition{male personal name}}}

\entry{Kangge}{\headword{Kangge}
  \pos{Proper noun}  \sense{
    \definition{male personal name}}}

\entry{kangkäg}{\headword{kangkäg}  \pronnote{kaŋkəɡ}
  \pos{Transitive verb}  \sense{
    \definition{to carry, bear (a load)}}}

\entry{kängkäl}{\headword{kängkäl}  \pronnote{kəŋkəl}
  \pos{Transitive verb}  \sense{\textbf{1.} 
    \definition{to ascend, climb, go up, rise}    \example{
      \exsen{Yäbäd a dingkälän tuk e.}      \extran{The sun rose up into the sky.}}}  \sense{\textbf{2.} 
    \definition{to ascend, climb, go up}    \example{
      \exsen{Lla da llo de kälngkäl eran.}      \extran{The man is climbing the tree.}}    \example{
      \exsen{Bogo manggo de dängkälän.}      \extran{She climbed the mango tree.}}}
  \variants{\Unspecified Variant \vartext{kälngkäl, kälängkäl}}}

\entry{kängkäm}{\headword{kängkäm}  \pronnote{kəŋkəm}
  \pos{Transitive verb}  \sense{
    \definition{to squeeze, press}    \example{
      \exsen{Ge kup a gärep käp yid kämnan e dan.}      \extran{This hole is for pressing grape juice (i.e. a winepress).}}    \example{
      \exsen{Ubi sana kämnanang dag.}      \extran{They are squeezing sago.}}}}

\entry{Kaninga}{\headword{Kaninga}  \pronnote{kaniŋa}
  \pos{}  \sense{
    \definition{male personal name}}}

\entry{kanken}{\headword{kanken}  \note{Mushrooms}  \pronnote{kanken}
  \pos{Noun}  \sense{
    \definition{type of mushroom}    \note{Mushrooms}    \example{
}    \example{
}    \example{
}}}

\entry{kannas}{\headword{kannas}  \note{Hunting}  \pronnote{kannas}
  \pos{Noun}  \sense{
    \definition{type of bow made out of pitpit}    \note{Hunting}    \example{
}    \example{
}    \example{
}}}

\entry{kanoe}{\headword{kanoe}  \pronnote{kanoe}
  \pos{Noun}  \sense{
    \definition{type of tree with fruits that cassowary, pigs, and deer eat and bark that is put in the water to kill fish}}}

\entry{kanong}{\headword{kanong}
  \variantof{Dialectal Variant of \vartext{tanong}}  \pronnote{kanoŋ}}

\entry{kansel}{\headword{kansel}
  \pos{Noun}  \sense{
    \definition{counsel}    \example{
}}}

\entry{kantärpie}{\headword{kantärpie}  \note{Water}  \pronnote{kantərpie}
  \pos{Noun}  \sense{
    \definition{log stuck in the water}    \note{Water}    \example{
}    \example{
}    \example{
}}}

\entry{käntrokäntrom nane}{\headword{käntrokäntrom nane}  \pronnote{kəntrokəntrom nane}
  \pos{Ideophone}  \sense{
    \definition{gulp}    \example{
      \exsen{Ngäna ine de käntrokäntrom nänawan.}      \extran{I gulped the water.}}}}

\entry{känttatt}{\headword{känttatt}  \note{Everything in the house}  \pronnote{kənʈ͡ʂaʈ͡ʂ}
  \pos{Noun}  \sense{
    \definition{bedroom; room, chamber}    \note{Everything in the house}    \example{
      \exsen{Bogo mɨnyi tuk me ulle abal känttatt de yantepägän.}      \extran{He will show you (pl.) a very big room above.}}    \example{
}    \example{
}}}

\entry{känyäkänyär}{\headword{känyäkänyär}
  \variantof{Fast Speech Variant of \vartext{känyärkänyär}}}

\entry{kanyam}{\headword{kanyam}  \pronnote{kaɲam}
  \pos{Intransitive verb}  \sense{
    \definition{bend}}
  \variants{\Dialectal Variant \vartext{känyam}}}

\entry{känyam}{\headword{känyam}
  \variantof{Dialectal Variant of \vartext{kanyam}}}

\entry{känyär}{\headword{känyär}  \pronnote{kəɲər}
  \pos{Modifier}  \sense{\textbf{1.} 
    \definition{quiet}    \example{
      \exsen{Dagirne dirom a känyär dae gongttägän.}      \extran{I stayed there and a cassowary came up quietly.}}    \example{
      \exsen{Känyär giddollag be ddobae melemang deya.}      \extran{He was humble (lit. one who lives quietly) but hardworking.}}}  \sense{\textbf{2.} 
    \definition{secret, secretly}    \example{
      \exsen{Bongo känyär nonttog, Godd aebe umllang bogon.}      \extran{If you give it secretly, only God will know.}}    \example{
      \exsen{Obo pate oba känyär me gobällän a obom dangnoeyo.}      \extran{They went to him in secret and asked him.}}}  \sense{\textbf{3.} 
    \definition{alone, by oneself (follows a genitive noun)}    \example{
      \exsen{Ngämi ngäma känyär ulleulle gogmam.}      \extran{We (excl.) grew up by ourselves.}}    \example{
      \exsen{Bogo ma me daeya obo känyär.}      \extran{He was home alone.}}    \example{
      \exsen{Bina känyär wiyamom ge ngättma we.}      \extran{Come by yourselves to this place.}}}
  \variants{\SpellingVariant \vartext{känyer}}}

\entry{känyär eka}{\headword{känyär eka}  \pronnote{kəɲər eka}
  \pos{Noun}  \sense{
    \definition{secret}}}

\entry{känyär ttoen}{\headword{känyär ttoen}  \pronnote{kəɲər ʈ͡ʂoen}
  \pos{Noun}  \sense{
    \definition{secret}}}

\entry{känyärkänyär}{\headword{känyärkänyär}  \pronnote{kəɲərkəɲər}
  \pos{Adverb}  \sense{
    \definition{secretly}    \example{
      \exsen{Tämamae oba tatuma ibi ttoen de täbetäbe de dätäbeyo känyärkänyär.}      \extran{They all secretly planned a way to go to their washing place.}}}
  \variants{\Fast Speech Variant \vartext{känyäkänyär}}}

\entry{känyärkänyär}{\headword{känyärkänyär}  \pronnote{kəɲərkəɲər}
  \pos{Transitive coverb}  \sense{
    \definition{to soothe}    \example{
      \exsen{Ause da llɨg kälsre de mäse känyerkänyer dägnän, be ddone mullae gogon.}      \extran{The old woman was trying to soothe the boy, but she couldn't.}}}
  \variants{\Fast Speech Variant \vartext{känyäkänyär}}}

\entry{känyärtto}{\headword{känyärtto}  \pronnote{kəɲərʈ͡ʂo}
  \pos{Modifier}  \sense{
    \definition{silent, quiet}    \example{
      \exsen{Imanuel gall guwo me känyärtto gogän.}      \extran{Imanuel was silent in the heart of the canoe.}}    \example{
      \exsen{Yuwog abal lla da obom umllang dägaeyo känyärtto we.}      \extran{Many people told him to be quiet.}}}
  \variants{\Unspecified Variant \vartext{känyertto}}}

\entry{kanye}{\headword{kanye}
  \variantof{Infinitival reduplicant of \vartext{RED~}}}

\entry{kanyekanye}{\headword{kanyekanye}  \pronnote{kaɲekaɲe}
  \pos{Intransitive coverb}  \sense{\textbf{1.} 
    \definition{to move around, travel}    \example{
      \exsen{Lla da yuwog abal dagaeya a dazirnän ada sip a alla ubira pänggmenyang llameny kanyekanye bognegnän.}      \extran{There were so many people gathered, there like how sheep will move around without a shepherd.}}    \example{
      \exsen{Ngämo kanyekanye att a ngättma da gänyag ada, Balimo, Daru, Mospi.}      \extran{Here are the places I've travelled to: Balimo, Daru, and Port Moresby.}}}  \sense{\textbf{2.} 
    \definition{to go hunting}    \example{
      \exsen{Ngäna iddob kanyekanye we gotäba.}      \extran{I planned to go on a hunt at night.}}}}

\entry{känyer}{\headword{känyer}
  \variantof{SpellingVariant of \vartext{känyär}}  \pronnote{kəɲer}}

\entry{känyertto}{\headword{känyertto}
  \variantof{Unspecified Variant of \vartext{känyärtto}}}

\entry{Känykäny}{\headword{Känykäny}
  \variantof{Dialectal Variant of \vartext{Kinykiny}}}

\entry{känz}{\headword{känz}  \pronnote{kənz}
  \pos{Intransitive verb}  \sense{
    \definition{to go aside, go off course}    \example{
      \exsen{Bogo llo gäbagäba de ikop dägagän a dädme towall gongkäzän a ttäle gottkewän.}      \extran{He saw a shady tree and went off into the grass there and folded his legs.}}}}

\entry{Kaoga}{\headword{Kaoga}
  \pos{Proper noun}  \sense{
    \definition{male personal name}}}

\entry{kaonggall}{\headword{kaonggall}  \pronnote{kaoŋɡaɽ}
  \pos{Noun}  \sense{
    \definition{yellow-faced myna}    \scientificname{Mino dumontii}}}

\entry{käp}{\headword{käp}  \lexeme{käp}
  \pos{Noun}  \sense{\textbf{1.} 
    \definition{fruit}    \example{
      \exsen{Llɨg a llo käp de notan.}      \extran{The boy ate the fruit.}}}  \sense{\textbf{2.} 
    \definition{egg}    \example{
      \exsen{Ge pa da käp tumang zazeag dan.}      \extran{This bird lays many eggs.}}}  \sense{\textbf{3.} 
    \definition{vague relatively small and round thing}    \example{
      \exsen{wayati käp, dädär käp, ttägäll käp, batri käp, tudi käp}      \extran{watermelon, rock, money, battery, fish hook}}    \example{
      \exsen{Malla ngäna tumang kollba de naittnegan, tri käpdae.}      \extran{I didn't catch many fish, only three.}}}  \sense{\textbf{4.} 
}}

\entry{käp}{\headword{käp}
  \variantof{SpellingVariant of \vartext{kap}}}

\entry{kap}{\headword{kap}
  \pos{Noun}  \sense{
}
  \variants{\SpellingVariant \vartext{käp}}}

\entry{kapa}{\headword{kapa}
  \pos{Noun}  \sense{
}}

\entry{käpäkäpät}{\headword{käpäkäpät}  \note{Adjectives (Swadesh)}  \pronnote{kəpəkəpət}
  \pos{Adverb}  \sense{
    \definition{wet, soaked}    \note{Adjectives (Swadesh)}    \example{
      \exsen{Bogo walle we guspunän, käsre ma we käptakäptae gongosän.}      \extran{He fell in the water, so he came back home soaking wet.}}    \example{
}    \example{
}}
  \variants{\Derived Variant \vartext{käptakäptae}}}

\entry{Kapal}{\headword{Kapal}  \pronnote{kapal}
  \pos{Proper noun}  \sense{
    \definition{Kapal (Wipi- and Kawam-speaking village in Oriomo-Bituri Rural LLG; has airstrip, aid post, primary school; from Limol, one must pass through Bisuaka)}    \example{
}    \example{
}    \example{
}}}

\entry{kapalla}{\headword{kapalla}  \note{Water}  \pronnote{kapaɽa}
  \pos{Noun}  \sense{
    \definition{floating grass}    \note{Water}    \example{
      \exsen{Ngämlle kapalla ttam alle sana yuatt a ddobae moko dan.}      \extran{I like to eat sago wrapped in kapalla leaves.}}    \example{
}    \example{
}}}

\entry{kapän}{\headword{kapän}
  \pos{Noun}  \sense{
}}

\entry{käpang}{\headword{käpang}
  \pos{Modifier}  \sense{
    \definition{bearing fruit}    \example{
      \exsen{Tomato da ddone käpang gogon.}      \extran{The tomato plant had no fruit.}}}}

\entry{kapang}{\headword{kapang}  \note{Medicine}  \pronnote{kapaŋ}
  \pos{Noun}  \sense{
    \definition{Acacia}    \scientificname{Acacia}    \note{Medicine}    \example{
}    \example{
}    \example{
}}}

\entry{kapang bile}{\headword{kapang bile}  \note{Medicine}  \pronnote{kapaŋ bile}
  \pos{Noun}  \sense{
    \definition{type of medicine}    \note{Medicine}    \example{
}    \example{
}    \example{
}}}

\entry{kapang budar}{\headword{kapang budar}  \note{Grubs}  \pronnote{kapaŋ budar}
  \pos{Noun}  \sense{
    \definition{type of grub}    \note{Grubs}    \example{
}    \example{
}    \example{
}}}

\entry{Kapangang bun}{\headword{Kapangang bun}  \note{Places around Limol}  \pronnote{kapaŋaŋ bun}
  \pos{Proper noun}  \sense{
    \definition{Kapangang bun (sago place on the way to Egapo)}    \note{Places around Limol}    \example{
}    \example{
}    \example{
}}}

\entry{kapangmändär}{\headword{kapangmändär}  \note{Mushrooms}  \pronnote{kapaŋməndər}
  \pos{Noun}  \sense{
    \definition{type of mushroom}    \note{Mushrooms}    \example{
}    \example{
}    \example{
}}}

\entry{Kaparnaom}{\headword{Kaparnaom}  \pronnote{kaparnaom}
  \pos{Proper noun}  \sense{
    \definition{Capernaum}}}

\entry{käpät}{\headword{käpät}  \pronnote{kəpət}
  \pos{Noun}  \sense{
    \definition{moisture}    \example{
}    \example{
}    \example{
}}}

\entry{käpätang}{\headword{käpätang}
  \pos{Modifier}  \sense{
    \definition{wet}    \example{
      \exsen{Ngämo tater a käptang agan.}      \extran{My mat get wet.}}}
  \variants{\Fast Speech Variant \vartext{käptang}}}

\entry{käpe sen}{\headword{käpe sen}  \pronnote{kəpe sen}
  \pos{Verb}  \sense{
    \definition{come out}}}

\entry{kapera}{\headword{kapera}  \note{Friendship/social terms of endearment}  \pronnote{kapera}
  \pos{Noun}  \sense{
    \definition{male friend or partner who comes from out of town}    \note{Friendship/social terms of endearment}    \example{
}    \example{
}    \example{
}}}

\entry{kapi~}{\headword{kapi~}
  \variantof{Infinitival reduplicant of \vartext{RED~}}}

\entry{kapi}{\headword{kapi}
  \variantof{freevarlab of \vartext{English loanword}}}

\entry{kapkap}{\headword{kapkap}  \lexeme{kapkap}  \pronnote{kapkap}
  \pos{Noun}  \sense{
    \definition{mudskipper}}}

\entry{käpkumett}{\headword{käpkumett}  \note{Trees}  \pronnote{kəpkumeʈ͡ʂ}
  \pos{Noun}  \sense{
    \definition{type of tree}    \note{Trees}    \example{
}    \example{
}    \example{
}}}

\entry{käpom}{\headword{käpom}  \note{Trees}  \pronnote{kəpom}
  \pos{Noun}  \sense{
    \definition{type of big tree that grows in the bush with white flowers and edible white fruit; used as medicine for cough}    \note{Trees}    \example{
}    \example{
}    \example{
}}}

\entry{kapräl}{\headword{kapräl}  \note{Trees}  \pronnote{kaprəl}
  \pos{Noun}  \sense{
    \definition{type of tree}    \note{Trees}    \example{
}    \example{
}    \example{
}}}

\entry{käpre}{\headword{käpre}
  \pos{Noun}  \sense{
    \definition{type of big yam with a white interior, thorns, and few hairs}}}

\entry{käptakäptae}{\headword{käptakäptae}
  \variantof{Derived Variant of \vartext{käpäkäpät}}}

\entry{käptang}{\headword{käptang}
  \variantof{Fast Speech Variant of \vartext{käpätang}}}

\entry{kaptte}{\headword{kaptte}  \note{Clothing}  \pronnote{kapʈ͡ʂe}
  \pos{Noun}  \sense{\textbf{1.} 
    \definition{cloth}    \note{Clothing}    \example{
      \exsen{Mälla da kaptte gullun alle obo bun di dɨmllawän.}      \extran{The woman tied her head with an old cloth.}}}  \sense{\textbf{2.} 
    \definition{clothing, clothes; piece of clothing, garment}    \note{Clothing}    \example{
      \exsen{Kaptte de kotang de ngäsangngäsang de mättnan a mudan.}      \extran{Don't wear dirty clothes over and over.}}    \example{
      \exsen{Ngämo kaptte da era komlla dageyo.}      \extran{I have two pieces of clothing.}}}}

\entry{kaptte dodro ma}{\headword{kaptte dodro ma}  \note{Everything in the house}  \pronnote{kapʈ͡ʂe dodro ma}
  \pos{Noun}  \sense{
    \definition{washing board}    \note{Everything in the house}    \example{
}    \example{
}    \example{
}}}

\entry{kaptte gällämnan}{\headword{kaptte gällämnan}  \note{Things a person can do}  \pronnote{kapʈ͡ʂe ɡəɽəmnan}
  \pos{Noun}  \sense{
    \definition{laundry, washing clothes}    \note{Things a person can do}    \example{
      \exsen{Lama kamekong gogon kaptte gällämnan me.}      \extran{Lama was busy doing the laundry.}}    \example{
}    \example{
}}}

\entry{kaptte ittal ma}{\headword{kaptte ittal ma}  \note{Everything in the house}  \pronnote{kapʈ͡ʂe iʈ͡ʂal ma}
  \pos{Noun}  \sense{
    \definition{clothes line}    \note{Everything in the house}    \example{
}    \example{
}    \example{
}}}

\entry{kaptte tubutubu}{\headword{kaptte tubutubu}  \note{Clothing}  \pronnote{kapʈ͡ʂe tubutubu}
  \pos{Noun}  \sense{
    \definition{short trousers}    \note{Clothing}    \example{
}    \example{
}    \example{
}}}

\entry{kaptte tupi}{\headword{kaptte tupi}  \note{Clothing}  \pronnote{kapʈ͡ʂe tupi}
  \pos{Noun}  \sense{
    \definition{long trousers}    \note{Clothing}    \example{
}    \example{
}    \example{
}}}

\entry{kapu}{\headword{kapu}  \note{Verbs of Motion}  \pronnote{kapu}
  \pos{Transitive coverb}  \sense{
    \definition{to carry}    \note{Verbs of Motion}    \example{
      \exsen{Ngäna llɨg di kapu iran.}      \extran{I am carrying the baby.}}    \example{
}}
  \variants{\Baby Talk Variant \vartext{papu}}}

\entry{kär pipiem}{\headword{kär pipiem}  \pronnote{kər pipiem}
  \pos{Noun}  \sense{
    \definition{purple-tailed imperial pigeon}    \scientificname{Ducula rufigaster}}}

\entry{karado}{\headword{karado}  \pronnote{karado}
  \pos{Noun}  \sense{
    \definition{long spear for fishing}    \example{
}    \example{
}    \example{
}}}

\entry{kärakäralmeny}{\headword{kärakäralmeny}  \pronnote{kərakəralmeɲ}
  \pos{Modifier}  \sense{
    \definition{swirling}    \example{
      \exsen{Kärakäralmeny agnan wel a.}      \extran{The wind was swirling.}}}}

\entry{Karama}{\headword{Karama}  \pronnote{karama}
  \pos{Proper noun}  \sense{
    \definition{Karama (swamp and canoe place in Limol)}    \example{
      \exsen{Ngämi ttongo ebdo me Karama we gobllab inu we, Eramang gall tapma we ddob llayaba kämäll.}      \extran{One day, we went to Karama to sleep, to the canoe place with some people.}}    \example{
      \exsen{Karama kona me giddollag dan.}      \extran{He lives in Karama district.}}    \example{
      \exsen{Bobag daya dedam karama walle da adawatta yogoll ulle da dämanän.}      \extran{Karama river was flooded because there had been big rains.}}}}

\entry{karama}{\headword{karama}  \pronnote{karama}
  \pos{Noun}  \sense{
    \definition{swamp}    \example{
      \exsen{Karama da ade petapeta abal gogon.}      \extran{The swamp also became very shallow.}}}}

\entry{Karama Popo}{\headword{Karama Popo}
  \variantof{SpellingVariant of \vartext{Karamapopo}}}

\entry{Karamapopo}{\headword{Karamapopo}  \pronnote{karamapopo}
  \pos{Proper noun}  \sense{
    \definition{female personal name}}
  \variants{\SpellingVariant \vartext{Karama Popo}}}

\entry{Karao}{\headword{Karao}  \pronnote{karao}
  \pos{Proper noun}  \sense{
    \definition{male personal name}}
  \variants{\SpellingVariant \vartext{Karau}}}

\entry{Karau}{\headword{Karau}
  \variantof{SpellingVariant of \vartext{Karao}}  \pronnote{karau}}

\entry{Karea}{\headword{Karea}  \pronnote{karea}
  \pos{Proper noun}  \sense{
    \definition{male personal name}}}

\entry{Karen}{\headword{Karen}  \pronnote{karen}
  \pos{Proper noun}  \sense{
    \definition{female personal name}}}

\entry{Kares}{\headword{Kares}
  \pos{Proper noun}  \sense{
    \definition{female personal name}}}

\entry{kargeam}{\headword{kargeam}  \note{Types of Taro}  \pronnote{karɡeam}
  \pos{Noun}  \sense{
    \definition{type of big taro}    \note{Types of Taro}    \example{
}    \example{
}    \example{
}}}

\entry{Karis}{\headword{Karis}  \pronnote{karis}
  \pos{Proper noun}  \sense{
    \definition{female personal name}}}

\entry{karita}{\headword{karita}  \note{Names of bananas}  \pronnote{karita}
  \pos{Noun}  \sense{
    \definition{type of introduced banana}    \note{Names of bananas}    \example{
}    \example{
}    \example{
}}}

\entry{karpo}{\headword{karpo}
  \pos{Noun}  \sense{
    \definition{jar}}}

\entry{Kas}{\headword{Kas}
  \pos{Proper noun}  \sense{
    \definition{male personal name}}}

\entry{Kasakmai}{\headword{Kasakmai}  \note{Places around Limol}  \pronnote{kasakmai}
  \pos{Proper noun}  \sense{\textbf{1.} 
    \definition{Kasakmai (toponym)}    \note{Places around Limol}    \example{
}    \example{
}    \example{
}}  \sense{\textbf{2.} 
    \definition{name of a person}    \note{Places around Limol}}}

\entry{Kasimap}{\headword{Kasimap}  \note{Place names}  \pronnote{kasimap}
  \pos{Proper noun}  \sense{
    \definition{Kasimap (Abom-speaking village in Gogodala Rural LLG, near Zanor; has an elementary school)}    \note{Place names}    \example{
}    \example{
}    \example{
}}}

\entry{Kasir}{\headword{Kasir}
  \pos{Proper noun}  \sense{
    \definition{male personal name}}}

\entry{Kaso}{\headword{Kaso}  \pronnote{kaso}
  \pos{Proper noun}  \sense{
    \definition{male personal name}}}

\entry{käsre}{\headword{käsre}  \pronnote{kəsre}
  \pos{Adverb}  \sense{
    \definition{then}    \example{
      \exsen{Polle da popang gogän, käsre ge ddäddäg a gazgez de gongkaemnegän.}      \extran{The fence got a hole, and then this animal started to escape.}}}}

\entry{kästom}{\headword{kästom}
  \variantof{Unspecified Variant of \vartext{kastom}}}

\entry{kastom}{\headword{kastom}
  \pos{Noun}  \sense{
    \definition{custom}}
  \variants{\Unspecified Variant \vartext{kästom}}}

\entry{kätam}{\headword{kätam}  \pronnote{kətam}
  \pos{Intransitive verb}  \sense{\textbf{1.} 
    \definition{to splash (of an aquatic animal)}    \example{
      \exsen{Tärko da mängalae källa gokätaemän.}      \extran{The small fish quickly splashed its tail.}}}  \sense{\textbf{2.} 
    \definition{to flip, do a cartwheel}}}

\entry{Katama}{\headword{Katama}  \pronnote{katama}
  \pos{Proper noun}  \sense{
    \definition{male personal name}}}

\entry{kätang}{\headword{kätang}  \pronnote{kətaŋ}
  \pos{Modifier}  \sense{
    \definition{blind}}}

\entry{kätäräl}{\headword{kätäräl}  \note{Colors}  \pronnote{kətərəl}
  \pos{Noun}  \sense{
    \definition{color}    \note{Colors}    \example{
      \exsen{Obo ikop kätäräl bätbät att a opnaeyan pällampällam e.}      \extran{His eyes changed from black to white.}}}
  \variants{\Fast Speech Variant \vartext{käträl}}}

\entry{Kate}{\headword{Kate}  \pronnote{kate}
  \pos{Proper noun}  \sense{
    \definition{female personal name}}
  \variants{\SpellingVariant \vartext{Keit}}}

\entry{Katherine}{\headword{Katherine}  \pronnote{kaθrin}
  \pos{Proper noun}  \sense{
    \definition{female personal name}}}

\entry{katkatre}{\headword{katkatre}
  \variantof{Fast Speech Variant of \vartext{katrekatre}}  \pronnote{katkatre}}

\entry{käträkäträl}{\headword{käträkäträl}
  \pos{Modifier}  \sense{
    \definition{multicolored}    \example{
      \exsen{Apapi da ddob erag käträkäträl dag.}      \extran{There are some butterflies which are multicolored.}}}}

\entry{käträl}{\headword{käträl}
  \variantof{Fast Speech Variant of \vartext{kätäräl}}}

\entry{katre~}{\headword{katre~}
  \variantof{Derivational reduplicant of \vartext{RED~}}}

\entry{katre}{\headword{katre}  \pronnote{katre}
  \pos{Noun}  \sense{\textbf{1.} 
    \definition{board; flooring}    \example{
      \exsen{Llɨg a katre me dämanenang nazernan.}      \extran{The children were sitting on the floor.}}}  \sense{\textbf{2.} 
    \definition{raised, elevated}    \example{
      \exsen{Yesu ubim dokomän obo peyang ada Pita, Zeims, a Zon tuk abal tutu katre we.}      \extran{Jesus took Peter, James, and John with him up a very tall mountain.}}}}

\entry{katre ma}{\headword{katre ma}
  \pos{Noun}  \sense{
    \definition{raised house, stilt house}    \example{
      \exsen{Bogo ikop täbaeb e katre ma de dogowän adawatta lla da gämäll bognegnän.}      \extran{He built a raised house for keeping watch, because people would steal.}}}}

\entry{katrekatre}{\headword{katrekatre}  \note{Everything in the house}  \pronnote{katrekatre}
  \pos{Noun}  \sense{\textbf{1.} 
    \definition{table; desk}    \note{Everything in the house}    \example{
      \exsen{Mälla da kakoll de katkatre toko me naspunan.}      \extran{The woman threw the plate onto the table.}}}  \sense{\textbf{2.} 
    \definition{shelf}    \note{Everything in the house}    \example{
      \exsen{Abo ddäddäg de nokomeyo a katrekatre toko me nowattälleyo.}      \extran{You (pl.) must take the meat and put it on the shelf.}}}
  \variants{\Fast Speech Variant \vartext{katkatre}}}

\entry{katt}{\headword{katt}  \pronnote{kaʈ͡ʂ}
  \pos{Noun}  \sense{
    \definition{Meyer's friarbird}    \scientificname{Philemon meyeri}}}

\entry{kätt}{\headword{kätt}  \pronnote{kəʈ͡ʂ}
  \pos{Noun}  \sense{\textbf{1.} 
    \definition{bivalve; shell of a mollusc}}  \sense{\textbf{2.} 
    \definition{shell blade}}  \sense{\textbf{3.} 
    \definition{(slang) vagina, vulva}}  \sense{\textbf{4.} 
    \definition{vague hard and thin thing (e.g. a blade)}    \example{
      \exsen{tubukätt}      \extran{kneecap}}}}

\entry{katt}{\headword{katt}  \pronnote{kaʈ͡ʂ}
  \pos{Noun}  \sense{
    \definition{type of medium-sized rodent that lives in the bush}}}

\entry{kätt ine}{\headword{kätt ine}  \note{Colors}  \pronnote{kəʈ͡ʂ ine}
  \pos{Color term}  \sense{
    \definition{blue (lit. 'shell water')}    \note{Colors}    \example{
}    \example{
}    \example{
}}}

\entry{kätt tongoe}{\headword{kätt tongoe}  \note{Games}  \pronnote{kəʈ͡ʂ toŋoe}
  \pos{Noun}  \sense{
    \definition{type of game played with shells}    \note{Games}    \example{
}    \example{
}    \example{
}}}

\entry{kättapun}{\headword{kättapun}  \note{Water}  \pronnote{kəʈ͡ʂapun}
  \pos{Noun}  \sense{
    \definition{type of reed}    \note{Water}    \example{
}    \example{
}    \example{
}}}

\entry{kättekätte}{\headword{kättekätte}  \note{Birds}  \pronnote{kəʈ͡ʂekəʈ͡ʂe}
  \pos{Noun}  \sense{
    \definition{red-cheeked parrot}    \scientificname{Geoffroyus geoffroyi aruensis}    \note{Birds}    \example{
}    \example{
}    \example{
}}}

\entry{kättkätt}{\headword{kättkätt}  \note{Weaving}  \pronnote{kəʈ͡ʂkəʈ͡ʂ}
  \pos{Noun}  \sense{\textbf{1.} 
    \definition{weaving pattern with alternating cross and chevron}    \note{Weaving}}  \sense{\textbf{2.} 
    \definition{to start a new weaving pattern}    \note{Weaving}    \example{
      \exsen{Ibetty tätäm wipellgallagallab alle nyäng de däkättän.}      \extran{Yesterday, Ibetty switched from the wipellgallagallab design to a new one.}}}}

\entry{kättkätt}{\headword{kättkätt}  \pronnote{kəʈ͡ʂkəʈ͡ʂ}
  \pos{Transitive verb}  \sense{
    \definition{to fence, wall, build a fence or wall}    \example{
      \exsen{Ngäna polle de käg alle kättkätt eran.}      \extran{I'm making a fence with the palm.}}    \example{
      \exsen{Bogo dädär alle polle de däkättän.}      \extran{He built a wall out of stone.}}}}

\entry{kättlla}{\headword{kättlla}  \note{Trees}  \pronnote{kəʈ͡ʂɽa}
  \pos{Noun}  \sense{
    \definition{type of tree that grows in the bush with white flowers}    \note{Trees}    \example{
}    \example{
}    \example{
}}}

\entry{Kättpälläk bällämang}{\headword{Kättpälläk bällämang}  \note{Place names around Limol}  \pronnote{kəʈ͡ʂpəɽək bəɽəmaŋ}
  \pos{Proper noun}  \sense{
    \definition{Jerry Dareda's sacred place (near ttälebun, on the road to Kinkin, near Binyomoll)}    \note{Place names around Limol}    \example{
}    \example{
}    \example{
}}}

\entry{kau}{\headword{kau}  \note{Clothing}  \pronnote{kau}
  \pos{Noun}  \sense{
    \definition{wrestling clothes}    \note{Clothing}    \example{
}    \example{
}    \example{
}}}

\entry{Kauga}{\headword{Kauga}  \pronnote{kauɡa}
  \pos{Proper noun}  \sense{
    \definition{male personal name}}}

\entry{kausar gllall}{\headword{kausar gllall}
  \pos{Noun}  \sense{
}}

\entry{Kawa}{\headword{Kawa}
  \pos{Proper noun}  \sense{
    \definition{male personal name}}}

\entry{kawa}{\headword{kawa}  \pronnote{kawa}
  \pos{Noun}  \sense{\textbf{1.} 
    \definition{announcement, notice; plea}    \example{
}    \example{
}    \example{
}}  \sense{\textbf{2.} 
    \definition{to preach}    \example{
      \exsen{kawawang lla}      \extran{preacher}}    \example{
      \exsen{Komllakomlla melem agnan ada, ag me medicine melem, toto we God bäne eka kawa gogalle.}      \extran{They did both jobs: medicine in the morning, and preaching God's word in the evening.}}}}

\entry{Kawam}{\headword{Kawam}  \note{Language Names}  \pronnote{kawam}
  \pos{Proper noun}  \sense{
    \definition{Kawam language (Pahoturi River language spoken in Wim)}    \note{Language Names}    \example{
}    \example{
}    \example{
}}}

\entry{Kawiapo}{\headword{Kawiapo}  \pronnote{kawiapo}
  \pos{Proper noun}  \sense{
    \definition{Kaviapu (village in Gogodala Rural LLG, near Tapila)}    \example{
}    \example{
}    \example{
}}}

\entry{Kawito}{\headword{Kawito}
  \pos{Proper noun}  \sense{
    \definition{Kawito (station in Gogodala Rural LLG)}}}

\entry{Kaya}{\headword{Kaya}  \pronnote{kaja}
  \pos{Proper noun}  \sense{
    \definition{male personal name}}}

\entry{Kaysy}{\headword{Kaysy}  \pronnote{kesi}
  \pos{Proper noun}  \sense{
    \definition{female personal name}}}

\entry{käza}{\headword{käza}  \pronnote{kəza}
  \pos{Noun}  \sense{\textbf{1.} 
    \definition{Hall's New Guinea crocodile}    \scientificname{Crocodylus halli}    \example{
      \exsen{Ngäna käza de nägliban.}      \extran{I scared away the crocodile.}}}  \sense{\textbf{2.} 
    \definition{iguana}    \example{
}    \example{
}    \example{
}}}

\entry{käza allko}{\headword{käza allko}  \note{Animals/Plants}  \pronnote{kəza aɽko}
  \pos{Noun}  \sense{
    \definition{type of fly that antagonizes other flies}    \note{Animals/Plants}    \example{
}    \example{
}    \example{
}}}

\entry{käza bädma}{\headword{käza bädma}  \pronnote{kəza bədma}
  \pos{Noun}  \sense{
    \definition{type of plant that only the crocodile clan wears when going hunting for crocodiles; also a traditional medicine}}}

\entry{käza burala}{\headword{käza burala}  \pronnote{kəza burala}
  \pos{Noun}  \sense{
    \definition{water lily}}}

\entry{käza da nägagan}{\headword{käza da nägagan}  \note{Birth}  \pronnote{kəza da nəɡaɡan}
  \pos{Phrase}  \sense{
    \definition{to miscarry}    \note{Birth}    \example{
}    \example{
}    \example{
}}}

\entry{Käza kup ine ma}{\headword{Käza kup ine ma}  \note{Places around Limol}  \pronnote{kəza kup ine ma}
  \pos{Proper noun}  \sense{
    \definition{well and sago place of Kwakmae in Limol (behind aid post)}    \note{Places around Limol}    \example{
}    \example{
}    \example{
}}}

\entry{käza wirwir}{\headword{käza wirwir}  \pronnote{kəza wirwir}
  \pos{Noun}  \sense{
    \definition{frilled monarch}    \scientificname{Arses telescopthalmus}}}

\entry{käza wirwir}{\headword{käza wirwir}  \pronnote{kəza wirwir}
  \pos{Intransitive coverb}  \sense{
    \definition{to call a crocodile out from the water}    \example{
      \exsen{Niki sisri ag me käza wirwir agnan walle menae me.}      \extran{This morning, Niki was calling the crocodile on the side of the water.}}}}

\entry{käzabun}{\headword{käzabun}
  \pos{Noun}  \sense{
    \definition{large sago flower}}}

\entry{käzapig}{\headword{käzapig}  \note{Trees}  \pronnote{kəzapiɡ}
  \pos{Noun}  \sense{
    \definition{type of tree with big, sour, black fruit and white and blue flowers}    \note{Trees}    \example{
}    \example{
}    \example{
}}}

\entry{ke}{\headword{ke}
  \pos{Question particle}  \sense{\textbf{1.} 
    \definition{question particle}    \example{
      \exsen{Ewatta ke bibi ddobaeddobae lel amalla?}      \extran{Why are you all so afraid?}}    \example{
      \exsen{Aenen ke ngämo mäg a?}      \extran{Who is my mother?}}}  \sense{\textbf{2.} 
    \definition{counterfactual}}
  \variants{\Unspecified Variant \vartext{eke}}}

\entry{keam}{\headword{keam}  \note{Animal verbs}  \pronnote{keam}
  \pos{Ideophone}  \sense{
    \definition{sound made by deer}    \note{Animal verbs}    \example{
}    \example{
}    \example{
}}}

\entry{Keisi}{\headword{Keisi}
  \pos{Proper noun}  \sense{
    \definition{female personal name}}}

\entry{Keit}{\headword{Keit}
  \variantof{SpellingVariant of \vartext{Kate}}}

\entry{Keith}{\headword{Keith}  \pronnote{kiθ}
  \pos{Proper noun}  \sense{
    \definition{male personal name}}
  \variants{\SpellingVariant \vartext{Kit}}}

\entry{kek}{\headword{kek}  \pronnote{kek}
  \pos{Noun}  \sense{
    \definition{orange-footed scrubfowl}    \scientificname{Megapodius reinwardt}    \example{
}    \example{
}    \example{
}}}

\entry{Keke}{\headword{Keke}  \pronnote{keke}
  \pos{Proper noun}  \sense{
    \definition{female personal name}}}

\entry{Keks}{\headword{Keks}  \pronnote{keks}
  \pos{Proper noun}  \sense{
    \definition{personal name}}}

\entry{kemibi}{\headword{kemibi}  \pronnote{kemibi}
  \pos{Quantifier}  \sense{
    \definition{many}    \example{
      \exsen{Ngäma ma me nongg a kemibi abal dag.}      \extran{In my house, there are many geckos.}}}
  \variants{\SpellingVariant \vartext{kamibi, kamebi}}}

\entry{kemol}{\headword{kemol}
  \pos{Noun}  \sense{
    \definition{camel}}}

\entry{kemp}{\headword{kemp}
  \pos{Noun}  \sense{
    \definition{camp}}}

\entry{Kemu}{\headword{Kemu}
  \pos{Proper noun}  \sense{
    \definition{male personal name}}}

\entry{Kename}{\headword{Kename}  \pronnote{kename}
  \pos{Proper noun}  \sense{
    \definition{Kename (village in Gogodala Rural LLG; on an island in the Fly River)}    \example{
}    \example{
}    \example{
}}}

\entry{Keni}{\headword{Keni}
  \variantof{SpellingVariant of \vartext{Kenny}}}

\entry{Kenny}{\headword{Kenny}  \pronnote{keni}
  \pos{Proper noun}  \sense{
    \definition{male personal name}}
  \variants{\SpellingVariant \vartext{Keni}}}

\entry{kep}{\headword{kep}  \pronnote{kep}
  \pos{Noun}  \sense{
    \definition{hip}}}

\entry{kep kutt}{\headword{kep kutt}  \pronnote{kep kuʈ͡ʂ}
  \pos{Noun}  \sense{
    \definition{hip bone}    \example{
}    \example{
}    \example{
}}}

\entry{keräma pudd}{\headword{keräma pudd}  \pronnote{kerəma puɖ͡ʐ}
  \pos{Noun}  \sense{
    \definition{type of small yam with a white interior}}}

\entry{kerema}{\headword{kerema}  \pronnote{kerema}
  \pos{Noun}  \sense{
    \definition{type of taro}}}

\entry{Keren}{\headword{Keren}
  \pos{Proper noun}  \sense{
    \definition{female personal name}}}

\entry{Kergowa}{\headword{Kergowa}  \note{Place names}  \pronnote{kerɡowa}
  \pos{Proper noun}  \sense{
    \definition{Kergowa (in Gogodala Rural LLG; near Balimo)}    \note{Place names}    \example{
}    \example{
}    \example{
}}}

\entry{Keriso}{\headword{Keriso}  \pronnote{keriso}
  \pos{Proper noun}  \sense{
    \definition{Christ}}}

\entry{Kesa}{\headword{Kesa}
  \pos{Proper noun}  \sense{
    \definition{male personal name}}}

\entry{Kesama}{\headword{Kesama}  \pronnote{kesama}
  \pos{Proper noun}  \sense{
    \definition{male personal name}}}

\entry{keta}{\headword{keta}  \pronnote{keta}
  \pos{Noun}  \sense{
    \definition{roof post}    \example{
      \exsen{Täne pittnen e keta da dagaya ottamänan.}      \extran{The stalk weaving has finished for the roof.}}}}

\entry{Keti}{\headword{Keti}  \pronnote{keti}
  \pos{Proper noun}  \sense{
    \definition{Kurupel Täräp (Limol village), which was moved from Old Limol to Old Man Kurupel's camping place approximately four generations before 2015}}
  \variants{\SpellingVariant \vartext{KT}}}

\entry{ketmar}{\headword{ketmar}  \note{Trees}  \pronnote{ketmar}
  \pos{Noun}  \sense{\textbf{1.} 
    \definition{type of tree that grows in old gardens; after being skinned and soaked, it is weaved into skirts}    \note{Trees}    \example{
}    \example{
}    \example{
}}  \sense{\textbf{2.} 
    \definition{two yams hanging from a stick in the center of a yam counting pile}    \note{Trees}}}

\entry{ketol}{\headword{ketol}
  \pos{Noun}  \sense{
    \definition{kettle}}}

\entry{Ketrin}{\headword{Ketrin}
  \pos{Proper noun}  \sense{
    \definition{female personal name}}}

\entry{ketrop}{\headword{ketrop}  \note{Seasons}  \pronnote{ketrop}
  \pos{Noun}  \sense{\textbf{1.} 
    \definition{season when rain clouds form but are blown away by the wind}    \note{Seasons}    \example{
}    \example{
}    \example{
}}  \sense{\textbf{2.} 
    \definition{to stop the rain}    \note{Seasons}    \example{
      \exsen{Bogo mɨnyi yogoll de ketrop bägagän.}      \extran{She will stop the rain.}}}}

\entry{Kevelyn}{\headword{Kevelyn}  \pronnote{kɛvɛlɪn}
  \pos{Proper noun}  \sense{
    \definition{female personal name}}}

\entry{Kewameyato}{\headword{Kewameyato}  \pronnote{kewamejato}
  \pos{Proper noun}  \sense{
    \definition{female personal name}}}

\entry{keyadaola}{\headword{keyadaola}  \note{Birds}  \pronnote{kejadaola}
  \pos{Noun}  \sense{
    \definition{type of bird}    \note{Birds}    \example{
}    \example{
}    \example{
}}}

\entry{Kiata}{\headword{Kiata}  \pronnote{kiata}
  \pos{Proper noun}  \sense{
    \definition{male personal name}}}

\entry{Kibobma}{\headword{Kibobma}  \note{Place names}  \pronnote{kibobma}
  \pos{Proper noun}  \sense{
    \definition{Kibobma (previous settlement of Limol village; on the road to Kinkin, near the creeks)}    \note{Place names}    \example{
}    \example{
}    \example{
}}}

\entry{Kibul}{\headword{Kibul}
  \variantof{Unspecified Variant of \vartext{Kibuli}}}

\entry{Kibuli}{\headword{Kibuli}  \pronnote{kibuli}
  \pos{Proper noun}  \sense{
    \definition{Kibuli (Em-speaking village in Oriomo-Bituri Rural LLG; near Kurunti)}    \example{
}    \example{
}    \example{
}}
  \variants{\Unspecified Variant \vartext{Kibul}}}

\entry{Kidarga}{\headword{Kidarga}  \pronnote{kidarɡa}
  \pos{Proper noun}  \sense{
    \definition{male personal name}}}

\entry{kidwe}{\headword{kidwe}  \note{Poison}  \pronnote{kidwe}
  \pos{Noun}  \sense{
    \definition{millipede}    \note{Poison}    \example{
}    \example{
}    \example{
}}}

\entry{kikiem}{\headword{kikiem}  \note{Birds}  \pronnote{kikiem}
  \pos{Noun}  \sense{
    \definition{type of bird}    \note{Birds}    \example{
}    \example{
}    \example{
}}}

\entry{kiklem}{\headword{kiklem}  \pronnote{kiklem}
  \pos{Noun}  \sense{
    \definition{type of small edible rodent}    \example{
}    \example{
}    \example{
}}}

\entry{Kikori}{\headword{Kikori}  \pronnote{kikori}
  \pos{Proper noun}  \sense{\textbf{1.} 
    \definition{Kikori (town in Kikori District, located on the Kikori Delta)}    \example{
}    \example{
}}  \sense{\textbf{2.} 
    \definition{Kikori (river that flows into the Gulf of Papua)}}}

\entry{Kila}{\headword{Kila}  \pronnote{kila}
  \pos{Proper noun}  \sense{
    \definition{personal name}}}

\entry{kili}{\headword{kili}  \pronnote{kili}
  \pos{Noun}  \sense{\textbf{1.} 
    \definition{happiness, joy}    \example{
      \exsen{Kili da daden.}      \extran{There is joy.}}}  \sense{\textbf{2.} 
    \definition{happy}    \example{
      \exsen{Ngämi tämamae ddone ada kili gogmam.}      \extran{We (excl.) were all so happy.}}    \example{
      \exsen{Ngäna bäne pate ddone kili allan.}      \extran{I am not happy with you.}}}  \sense{\textbf{3.} 
    \definition{to praise}    \example{
      \exsen{Bogo kili nägagan.}      \extran{He praised him.}}}}

\entry{kiliang}{\headword{kiliang}  \note{good characteristics of people}  \pronnote{kiliaŋ}
  \pos{Modifier}  \sense{
    \definition{happy}    \note{good characteristics of people}    \example{
      \exsen{Obo ingoll a ddone kiliang agan.}      \extran{His face was not happy.}}    \example{
}    \example{
}}}

\entry{kilikili}{\headword{kilikili}  \pronnote{kilikili}
  \pos{Transitive coverb}  \sense{\textbf{1.} 
    \definition{to greet}    \example{
      \exsen{Bogo mɨnyi ballän kilikili bägawän a ttang bänttepägän.}      \extran{He will go greet him and shake his hand.}}}  \sense{\textbf{2.} 
    \definition{to rejoice, celebrate}    \example{
      \exsen{Däbe amom de kilikili nägnallo didri?}      \extran{Whom are they celebrating?}}    \example{
      \exsen{Lla da bobällnän kilikili bägayaebneyo.}      \extran{People will go and celebrate.}}}}

\entry{kilikiliang}{\headword{kilikiliang}  \pronnote{kilikiliaŋ}
  \pos{Adverb}  \sense{
    \definition{happily, joyfully}    \example{
      \exsen{Meri ma we kilikiliangae dallän.}      \extran{Mary went home happily.}}}}

\entry{kilimeny}{\headword{kilimeny}  \pronnote{kilimeɲ}
  \pos{Modifier}  \sense{\textbf{1.} 
    \definition{unhappy, sad, annoyed}    \example{
      \exsen{Män kälsre da kilimeny agan.}      \extran{The little girl was sad.}}    \example{
      \exsen{Bogo kilimeny gongäsän ma we.}      \extran{He went home unhappy.}}}  \sense{\textbf{2.} 
    \definition{to insult}    \example{
      \exsen{Bogo kilimeny nägagan.}      \extran{He insulted him.}}    \example{
}    \example{
}}}

\entry{kina}{\headword{kina}
  \pos{Noun}  \sense{
    \definition{kina (PGK, the currency of Papua New Guinea)}}}

\entry{kinekineang}{\headword{kinekineang}  \note{good characteristics of people}  \pronnote{kinekineaŋ}
  \pos{Modifier}  \sense{
    \definition{smart}    \note{good characteristics of people}    \example{
}    \example{
}}}

\entry{Kingsli}{\headword{Kingsli}
  \pos{Proper noun}  \sense{
    \definition{male personal name}}}

\entry{Kini}{\headword{Kini}  \pronnote{kini}
  \pos{Proper noun}  \sense{
    \definition{Kini (in Gogodala Rural LLG; near Balimo and Awaba)}    \example{
}    \example{
}    \example{
}}}

\entry{Kinkin}{\headword{Kinkin}  \pronnote{kinkin}
  \pos{Proper noun}  \sense{
    \definition{Kinkin (Ende- and Taeme-speaking village in Oriomo-Bituri Rural LLG, near Limol)}    \example{
}    \example{
}    \example{
}}
  \variants{\Unspecified Variant \vartext{Kinykiny}}}

\entry{kinpop}{\headword{kinpop}  \note{Trees}  \pronnote{kinpop}
  \pos{Noun}  \sense{
    \definition{type of tree that grows in the bush with white flowers and wood used for house sticks and rafters}    \note{Trees}    \example{
}    \example{
}    \example{
}}}

\entry{Kinykiny}{\headword{Kinykiny}
  \variantof{Unspecified Variant of \vartext{Kinkin}}
  \variants{\Dialectal Variant \vartext{Känykäny}}}

\entry{Kiongga}{\headword{Kiongga}
  \pos{Proper noun}  \sense{
    \definition{Kiongga (toponym)}}}

\entry{kip}{\headword{kip}  \pronnote{kip}
  \pos{Noun}  \sense{
    \definition{top of a plant}    \example{
      \exsen{llo kip}      \extran{treetop}}}}

\entry{kip papa}{\headword{kip papa}  \note{Gardens}  \pronnote{kip papa}
  \pos{Noun}  \sense{
    \definition{top of fence}    \note{Gardens}    \example{
}    \example{
}    \example{
}}}

\entry{Kiplin}{\headword{Kiplin}
  \pos{Proper noun}  \sense{
    \definition{male personal name}}}

\entry{Kipling}{\headword{Kipling}  \pronnote{kipliŋ}
  \pos{Proper noun}  \sense{
    \definition{male personal name}}}

\entry{kire}{\headword{kire}  \pronnote{kire}
  \pos{Modifier}  \sense{
    \definition{unripe; raw, fresh}}}

\entry{kirekire}{\headword{kirekire}  \pronnote{kirekire}
  \pos{Color term}  \sense{
    \definition{green}}}

\entry{kisin}{\headword{kisin}
  \pos{Noun}  \sense{
    \definition{kitchen}}}

\entry{Kit}{\headword{Kit}
  \variantof{SpellingVariant of \vartext{Keith}}}

\entry{kitapatt}{\headword{kitapatt}  \pronnote{kitapaʈ͡ʂ}
  \pos{Noun}  \sense{
    \definition{type of bandicoot-like animal}    \example{
}    \example{
}    \example{
}}}

\entry{kitar}{\headword{kitar}
  \pos{Noun}  \sense{
    \definition{floating grass}}}

\entry{kito}{\headword{kito}  \pronnote{kito}
  \pos{Noun}  \sense{
    \definition{type of black palm; in this immature stage, the shoots eaten as medicine and used to make baskets Used to build house and mat in the bush}    \example{
}    \example{
}    \example{
}}}

\entry{Kiwae}{\headword{Kiwae}
  \variantof{SpellingVariant of \vartext{Kiwai}}}

\entry{Kiwai}{\headword{Kiwai}  \note{Language Names}  \pronnote{kiwai}
  \pos{Proper noun}  \sense{
    \definition{Kiwai language (offical language of the region, native language of Daru; children's school songs are sometimes in this language)}    \note{Language Names}    \example{
}    \example{
}    \example{
}}
  \variants{\SpellingVariant \vartext{Kiwae}}}

\entry{kiwale wel}{\headword{kiwale wel}  \note{?}  \pronnote{kiwale wel}
  \pos{Noun}  \sense{
    \definition{type of wind}    \note{?}    \example{
}    \example{
}    \example{
}}}

\entry{kiyaddadda}{\headword{kiyaddadda}  \pronnote{kijaɖ͡ʐaɖ͡ʐa}
  \pos{Noun}  \sense{
    \definition{paradise kingfisher (common, buff-breasted, little)}    \scientificname{Tanysiptera galatea, T. sylvia, T. hydrocharis}    \example{
}    \example{
}    \example{
}}}

\entry{kɨk}{\headword{kɨk}  \pronnote{kɪk}
  \pos{Verb}  \sense{
    \definition{to crumple}}}

\entry{kɨllakɨlle}{\headword{kɨllakɨlle}  \pronnote{kɪɽakɪɽe}
  \pos{Transitive verb}  \sense{
    \definition{to scrape}}}

\entry{kɨllɨll}{\headword{kɨllɨll}  \note{Snakes}  \pronnote{kɪɽɪɽ}
  \pos{Noun}  \sense{
    \definition{type of snake}    \note{Snakes}    \example{
      \exsen{Ngäna kɨllɨll di näddägan.}      \extran{I ate the killill snake.}}    \example{
}    \example{
}}
  \variants{\Unspecified Variant \vartext{källäll}}}

\entry{kɨllɨm}{\headword{kɨllɨm}
  \variantof{Unspecified Variant of \vartext{källäm}}  \note{Water}  \pronnote{kɪɽɪm}}

\entry{kɨllkɨll}{\headword{kɨllkɨll}  \pronnote{kɪɽkɪɽ}
  \pos{Transitive verb}  \sense{
    \definition{to dig}    \example{
      \exsen{Auma däkällallo.}      \extran{They dig a grave.}}    \example{
      \exsen{Gottamänän dibaballe kup di däkällneg a pos de däganeg.}      \extran{I finished, and then I dug the holes and put the posts in.}}}
  \variants{\Dialectal Variant \vartext{källkäll}}}

\entry{kla~}{\headword{kla~}
  \variantof{Infinitival reduplicant of \vartext{RED~}}}

\entry{klak}{\headword{klak}  \note{Government}  \pronnote{klak}
  \pos{Noun}  \sense{
    \definition{court clerk}    \note{Government}    \example{
}    \example{
}    \example{
}}}

\entry{klak}{\headword{klak}  \pronnote{klak}
  \pos{Noun}  \sense{
    \definition{type of harpoon for fishing}}}

\entry{klaklae}{\headword{klaklae}
  \variantof{Fast Speech Variant of \vartext{kälakälae}}  \pronnote{klaklae}}

\entry{kläkle}{\headword{kläkle}
  \variantof{Fast Speech Variant of \vartext{kälekäle}}  \pronnote{kləkle}}

\entry{klas}{\headword{klas}  \pronnote{klas}
  \pos{Noun}  \sense{
    \definition{class}}
  \variants{\SpellingVariant \vartext{kälas}}}

\entry{klasrum}{\headword{klasrum}
  \pos{Noun}  \sense{
    \definition{classroom}}}

\entry{kle~}{\headword{kle~}
  \variantof{Plural reduplicant of \vartext{RED~}}}

\entry{klekle}{\headword{klekle}
  \variantof{Fast Speech Variant of \vartext{kälekäle}}}

\entry{klloklloe}{\headword{klloklloe}
  \pos{Transitive verb}  \sense{\textbf{1.} 
    \definition{to gather, join, come together}    \example{
      \exsen{Ubi goklloeyän Malläm ttängäm e.}      \extran{They gathered in Malam village.}}}  \sense{\textbf{2.} 
    \definition{to gather, bring together, mix}    \example{
      \exsen{Chairman a obo melem a ada dan ge, lla de baklloenegnän eka tameny e.}      \extran{This is the chairman's job: he will gather the people for discussions.}}    \example{
      \exsen{Ngämi eka de klloklloe panynen eralla.}      \extran{We are speaking mixing languages.}}}
  \variants{\SpellingVariant \vartext{kolloekolloe, kollokolloe}}}

\entry{kllomokllomoll}{\headword{kllomokllomoll}  \pronnote{kɽomokɽomoɽ}
  \pos{Adverb}  \sense{
    \definition{downwind, leeward}    \example{
      \exsen{Däräng dedme dinduag a ada gogon, wel kllomokllomoll dallän.}      \extran{Dog went running like this; he went downwind.}}}}

\entry{kllum}{\headword{kllum}
  \variantof{Fast Speech Variant of \vartext{kullum}}  \pronnote{kɽum}}

\entry{ko}{\headword{ko}
  \variantof{Fast Speech Variant of \vartext{ako}}}

\entry{ko}{\headword{ko}
  \pos{TAM particle}  \sense{
    \definition{until}}}

\entry{ko~}{\headword{ko~}
  \variantof{Infinitival reduplicant of \vartext{RED~}}}

\entry{kobädd}{\headword{kobädd}  \note{Trees}  \pronnote{kobəɖ͡ʐ}
  \pos{Noun}  \sense{
    \definition{type of tree}    \note{Trees}    \example{
}    \example{
}    \example{
}}}

\entry{Kobam}{\headword{Kobam}  \pronnote{kobam}
  \pos{Proper noun}  \sense{
    \definition{male personal name}}}

\entry{Kobddag}{\headword{Kobddag}
  \pos{Proper noun}  \sense{
    \definition{Kobddag (toponym)}}}

\entry{kobe}{\headword{kobe}  \lexeme{kobe}  \pronnote{kobe}
  \pos{Noun}  \sense{
    \definition{type of tree that grows along creeks (~ 3 m) with white and red or blue and purple flowers and edible red fruit with black or white stripes and 4-6 seeds inside}}}

\entry{Kobe}{\headword{Kobe}
  \pos{Proper noun}  \sense{
    \definition{female personal name}}}

\entry{kobeam}{\headword{kobeam}  \lexeme{kobeam}  \note{Marriage}  \pronnote{kobeam}
  \pos{Kinship noun}  \sense{
    \definition{co-brother-in-law (man's wife's sister's husband)}    \note{Marriage}}
  \variants{\SpellingVariant \vartext{kobeyam}}}

\entry{Kobemitang}{\headword{Kobemitang}  \note{Places around Limol}  \pronnote{kobemitaŋ}
  \pos{Proper noun}  \sense{
    \definition{Kobemitang (toponym)}    \note{Places around Limol}    \example{
}    \example{
}    \example{
}}}

\entry{kobeyam}{\headword{kobeyam}
  \variantof{SpellingVariant of \vartext{kobeam}}}

\entry{Koboddag}{\headword{Koboddag}  \note{Place names around Limol}  \pronnote{koboɖ͡ʐaɡ}
  \pos{Proper noun}  \sense{
    \definition{Koboddag (toponym)}    \note{Place names around Limol}    \example{
}    \example{
}    \example{
}}}

\entry{kodor}{\headword{kodor}  \pronnote{kodor}
  \pos{Noun}  \sense{
    \definition{piece, lump}    \example{
      \exsen{ekaklle kodor}      \extran{plot of land}}    \example{
      \exsen{Alla ubira sana kodor de basiaemeya?}      \extran{Can you give them a piece of sago?}}}}

\entry{kodowa}{\headword{kodowa}  \pronnote{kodowa}
  \pos{Noun}  \sense{
    \definition{dish consisting of sago cooked in leaves on the fire}}}

\entry{Koe}{\headword{Koe}  \pronnote{koe}
  \pos{Proper noun}  \sense{
    \definition{male personal name}}}

\entry{Koebänang}{\headword{Koebänang}  \note{Places around Limol}  \pronnote{koebənaŋ}
  \pos{Proper noun}  \sense{
    \definition{Koebänang (sago and hunting place; old settlement near Buddobuddog)}    \note{Places around Limol}    \example{
}    \example{
}    \example{
}}
  \variants{\Fast Speech Variant \vartext{Koebnang}}}

\entry{Koebnang}{\headword{Koebnang}
  \variantof{Fast Speech Variant of \vartext{Koebänang}}}

\entry{koeme}{\headword{koeme}  \note{Trees}  \pronnote{koeme}
  \pos{Noun}  \sense{
    \definition{type of tree that grows along creeks with edible, round red fruit}    \note{Trees}    \example{
}    \example{
}    \example{
}}}

\entry{koemekoeme}{\headword{koemekoeme}  \note{Trees}  \pronnote{koemekoeme}
  \pos{Noun}  \sense{
    \definition{type of tree that grows along creeks with inedible, small red fruit}    \note{Trees}    \example{
}    \example{
}    \example{
}}}

\entry{koen}{\headword{koen}  \pronnote{koen}
  \pos{Intransitive verb}  \sense{
    \definition{to turn back}    \example{
      \exsen{Ngäna llowam atta koen allan.}      \extran{I'm turning back due to fatigue.}}}}

\entry{koenbäll}{\headword{koenbäll}  \note{Trees}  \pronnote{koenbəɽ}
  \pos{Noun}  \sense{
    \definition{type of tree that grows in the bush (especially in old gardens) with hanging green fruit and liquid used to treat sores}    \note{Trees}    \example{
}    \example{
}    \example{
}}}

\entry{Koenbäll kutt}{\headword{Koenbäll kutt}  \note{Places around Limol}  \pronnote{koenbəɽ kuʈ͡ʂ}
  \pos{Proper noun}  \sense{
    \definition{Koenbäll kutt (sago and washing place of Paine and Warama Kurupel)}    \note{Places around Limol}    \example{
}    \example{
}    \example{
}}}

\entry{koenmäll}{\headword{koenmäll}  \pronnote{koenməɽ}
  \pos{Transitive verb}  \sense{
    \definition{to chase}    \example{
      \exsen{Bogo obom nängkoenmällan.}      \extran{He chased him.}}    \example{
      \exsen{Ttongo täräp me llɨg kälekäle da gotäbanegän pa koenmäll e.}      \extran{One time, little boys planned to chase birds.}}    \example{
      \exsen{Dadabi, mäsemäse llɨg a dongkoenmällaemneyo.}      \extran{The boys chased them naked.}}    \example{
      \exsen{Lla da poura de nängkoenmällnegan.}      \extran{The man chased four chickens.}}}
  \variants{\SpellingVariant \vartext{koinmäll}}}

\entry{koepang}{\headword{koepang}  \pronnote{koepaŋ}
  \pos{Modifier}  \sense{
    \definition{sour}    \example{
      \exsen{Ddob sana gudne da koepangkoepang mokowang dag.}      \extran{Some old pieces of sago taste sour.}}}}

\entry{koinmäll}{\headword{koinmäll}
  \variantof{SpellingVariant of \vartext{koenmäll}}  \pronnote{koinməɽ}}

\entry{kok}{\headword{kok}  \lexeme{kok}  \pronnote{kok}
  \pos{Noun}  \sense{\textbf{1.} 
    \definition{moon}    \example{
      \exsen{Kok ekaklle de nganaenen eran.}      \extran{The moon spins around Earth.}}}  \sense{\textbf{2.} 
    \definition{month}    \example{
      \exsen{Ada gongnomenyne, ada ka ttongdae kok me daeya sukul a.}      \extran{I used to think that school was one month long.}}}}

\entry{kok}{\headword{kok}  \lexeme{kok}  \pronnote{kok}
  \pos{Kinship noun}  \sense{\textbf{1.} 
    \definition{grandchild (one's child's child)}}  \sense{\textbf{2.} 
    \definition{daughter-in-law (one's son's wife)}}}

\entry{kok}{\headword{kok}  \lexeme{kok}  \pronnote{kok}
  \pos{Noun}  \sense{
    \definition{grasshopper}    \example{
      \exsen{Mälla da sisri ag me kok de nägäddaeballo kollba tokong e.}      \extran{This morning, the women caught grasshoppers for fishing bait.}}}}

\entry{kok a nganzig}{\headword{kok a nganzig}  \note{Birth}  \pronnote{kok a ŋanziɡ}
  \pos{Phrase}  \sense{
    \definition{to go a month without menstruating}    \note{Birth}}}

\entry{kok apte}{\headword{kok apte}  \note{Moon phases}  \pronnote{kok apte}
  \pos{Noun}  \sense{
    \definition{half moon}    \note{Moon phases}    \example{
}    \example{
}    \example{
}}}

\entry{kok indrang}{\headword{kok indrang}
  \pos{Noun}  \sense{
    \definition{moonlight}}}

\entry{kok kätt pälläk}{\headword{kok kätt pälläk}
  \pos{Noun}  \sense{
}}

\entry{kok kuddäll}{\headword{kok kuddäll}  \note{Moon phases}  \pronnote{kok kuɖ͡ʐəɽ}
  \pos{Noun}  \sense{
    \definition{new moon}    \note{Moon phases}    \example{
}    \example{
}    \example{
}}}

\entry{kok pälläk}{\headword{kok pälläk}  \note{Moon phases}  \pronnote{kok pəɽək}
  \pos{Noun}  \sense{
    \definition{first quarter moon}    \note{Moon phases}    \example{
}    \example{
}    \example{
}}}

\entry{kok patar}{\headword{kok patar}  \note{Clothing}  \pronnote{kok patar}
  \pos{Noun}  \sense{
    \definition{necklace}    \note{Clothing}    \example{
}    \example{
}    \example{
}}}

\entry{kok pllayang}{\headword{kok pllayang}  \note{Moon phases}  \pronnote{kok pɽajaŋ}
  \pos{Noun}  \sense{
    \definition{full moon}    \note{Moon phases}    \example{
}    \example{
}    \example{
}}}

\entry{kok sisor}{\headword{kok sisor}  \note{Moon phases}  \pronnote{kok sisor}
  \pos{Noun}  \sense{
    \definition{first light (the end of a new moon)}    \note{Moon phases}    \example{
}    \example{
}    \example{
}}}

\entry{kok tärangg}{\headword{kok tärangg}
  \pos{Noun}  \sense{
    \definition{fast}    \example{
      \exsen{Ubi mɨnyi kok tärangg me bazärnän ge wik me.}      \extran{They will be fasting this week.}}}}

\entry{kok to}{\headword{kok to}  \pronnote{kok to}
  \pos{Noun}  \sense{
    \definition{orgy}    \example{
}    \example{
}    \example{
}}}

\entry{kokäl}{\headword{kokäl}  \lexeme{kokäl}  \pronnote{kokəl}
  \pos{Noun}  \sense{
    \definition{small mudcrab}}}

\entry{kokall}{\headword{kokall}  \note{Types of palms}  \pronnote{kokaɽ}
  \pos{Noun}  \sense{
    \definition{type of palm with fruit that are smaller than coconuts}    \note{Types of palms}    \example{
}    \example{
}    \example{
}}}

\entry{kokallkokall}{\headword{kokallkokall}  \note{Trees}  \pronnote{kokaɽkokaɽ}
  \pos{Noun}  \sense{
    \definition{type of tree}    \note{Trees}    \example{
}    \example{
}    \example{
}}}

\entry{kokallma}{\headword{kokallma}  \note{Friendship/social terms of endearment}  \pronnote{kokaɽma}
  \pos{Noun}  \sense{
    \definition{two friends who share a twin fruit from the kokall tree (palm tree in the bush)}    \note{Friendship/social terms of endearment}    \example{
}    \example{
}    \example{
}}}

\entry{Koke}{\headword{Koke}
  \pos{Proper noun}  \sense{
    \definition{Koke (toponym)}}}

\entry{kokeya moleg}{\headword{kokeya moleg}
  \pos{Noun}  \sense{
    \definition{lost sacred stone with markings on it; represents a girl that belongs to the Dumollang clan}}}

\entry{kokllo}{\headword{kokllo}  \note{?}  \pronnote{kokɽo}
  \pos{Transitive verb}  \sense{
    \definition{to scratch, scrape}    \note{?}    \example{
      \exsen{Ngäna nge de kokllo eran.}      \extran{I'm scraping a coconut.}}    \example{
      \exsen{Lla ulle da bandrabandrog bun akllonan.}      \extran{The big man was scratching his head while dancing.}}    \example{
}}}

\entry{Kokma}{\headword{Kokma}
  \pos{Proper noun}  \sense{
    \definition{Kokma (toponym)}}}

\entry{kokmer}{\headword{kokmer}
  \pos{Noun}  \sense{
    \definition{when a hunter calls out his sister's name after shooting an animal (she will get the back of the animal if it's killed)}}}

\entry{kokne}{\headword{kokne}  \pronnote{kokne}
  \pos{Noun}  \sense{
    \definition{type of tree that grows in the grassland with blue flowers and edible blue fruit}    \example{
}    \example{
}    \example{
}}}

\entry{kokngal}{\headword{kokngal}  \note{Hunting}  \pronnote{kokŋal}
  \pos{Noun}  \sense{\textbf{1.} 
    \definition{hunting in the rain}    \note{Hunting}    \example{
      \exsen{Kokngal e ibi täräp a eran yogoll täräp me dan.}      \extran{The time to go kokngal hunting is when it's rainy.}}    \example{
}    \example{
}}  \sense{\textbf{2.} 
    \definition{season characterized by heavy rain, hunting, and fishing with lines and nets (third season; corresponds to mid-February)}    \note{Hunting}}}

\entry{koko}{\headword{koko}
  \pos{Noun}  \sense{
    \definition{shoot (of a plant)}    \example{
      \exsen{Mätta koko da dade de nganae eran.}      \extran{The yam shoot is wrapping around the yam stick.}}}}

\entry{koko}{\headword{koko}
  \pos{Transitive verb}  \sense{
    \definition{to cut (flesh or meat), butcher}    \example{
      \exsen{Bogo ddäddäg de dokonggän oblle.}      \extran{He cut the meat for her.}}    \example{
      \exsen{Ddia koko we ngäna llɨg di dokom.}      \extran{I took the boys to butcher the deer.}}}}

\entry{kokoang}{\headword{kokoang}
  \variantof{SpellingVariant of \vartext{kokowang}}}

\entry{kokok}{\headword{kokok}  \lexeme{kokok}  \pronnote{kokok}
  \pos{Kinship noun}  \sense{
    \definition{nonsingular form of kok}}}

\entry{kokol}{\headword{kokol}  \note{Names of bananas}  \pronnote{kokol}
  \pos{Noun}  \sense{
    \definition{type of introduced banana}    \note{Names of bananas}    \example{
}    \example{
}    \example{
}}}

\entry{kokop}{\headword{kokop}  \pronnote{kokop}
  \pos{Transitive verb}  \sense{
    \definition{to peel, skin, husk}    \example{
      \exsen{Ge lla da nae de kokop eran.}      \extran{This man is peeling a sweet potato.}}    \example{
      \exsen{Ttongo lla ulle da up o de näkopnan.}      \extran{A big man peeled a ripe banana.}}}}

\entry{kokopasi}{\headword{kokopasi}  \pronnote{kokopasi}
  \pos{Noun}  \sense{
    \definition{shrikethrush (Arafura, rufous)}    \scientificname{Colluricincla megarhyncha aruensis, C. rufogaster goodsoni}    \example{
}    \example{
}    \example{
}}}

\entry{kokowang}{\headword{kokowang}  \pronnote{kokowaŋ}
  \pos{Modifier}  \sense{\textbf{1.} 
    \definition{sprouted}    \example{
      \exsen{Nge kokowang a ibeny ma dan a ddob otnan ma dag.}      \extran{Some sprouted cocounts are for planting and others are for eating.}}}  \sense{\textbf{2.} 
    \definition{the first stage of coconut growth in which the shoot has just emerged (before planting)}}
  \variants{\SpellingVariant \vartext{kokoang}}}

\entry{kokpe}{\headword{kokpe}  \pronnote{kokpe}
  \pos{Noun}  \sense{
    \definition{type of tree with yellow flowers and leaves that are used to wrap food for cooking on the fire}}}

\entry{kokta}{\headword{kokta}  \pronnote{kokta}
  \pos{Noun}  \sense{
    \definition{moonlight}    \example{
      \exsen{Kokta da näbdaban.}      \extran{Moonlight is dawning.}}    \example{
      \exsen{Kokta da indrang abal bogon.}      \extran{Moon shines so bright.}}}}

\entry{koktakokta}{\headword{koktakokta}  \note{Trees}  \pronnote{koktakokta}
  \pos{Noun}  \sense{
    \definition{when dead trees or tree branches become dry and shine in the night}    \note{Trees}}}

\entry{kokto}{\headword{kokto}  \note{Diving}  \pronnote{kokto}
  \pos{Transitive verb}  \sense{
    \definition{to bail (water)}    \note{Diving}    \example{
      \exsen{Ngäna gall e godagal, ine dokto.}      \extran{I got in the canoe and bailed the water out.}}    \example{
}    \example{
}}}

\entry{kol}{\headword{kol}  \pronnote{kol}
  \pos{Noun}  \sense{
    \definition{sago pith}}}

\entry{kolem}{\headword{kolem}  \lexeme{kolem}  \pronnote{kolem}
  \pos{Noun}  \sense{
    \definition{type of palm with coconuts that come in red or green varieties}}}

\entry{koles}{\headword{koles}
  \variantof{SpellingVariant of \vartext{kolos}}}

\entry{koll}{\headword{koll}  \note{Hunting}  \pronnote{koɽ}
  \pos{Locational}  \sense{\textbf{1.} 
    \definition{inner part}    \note{Hunting}    \example{
      \exsen{ttang, ttang koll}      \extran{hand, palm}}}  \sense{\textbf{2.} 
    \definition{part of a bow}    \note{Hunting}    \example{
}    \example{
}    \example{
}}}

\entry{kollam}{\headword{kollam}  \pronnote{koɽam}
  \pos{Transitive verb}  \sense{
    \definition{to stand up}    \example{
      \exsen{Ngäna ma de dakollam.}      \extran{I stood the house up.}}}}

\entry{kollba}{\headword{kollba}  \note{Market}  \pronnote{koɽba}
  \pos{Noun}  \sense{
    \definition{fish}    \note{Market}    \example{
      \exsen{Bäne kollba mu da auli dan?}      \extran{How much for the fish?}}    \example{
      \exsen{Ngämo kollba mu da ttongo ttang lläpät dan.}      \extran{My fish is five kina.}}    \example{
      \exsen{Kollba da tämamae dadrowän a be kottllam aebe ttam dagirnän.}      \extran{The fish all died but the turtle lived on.}}    \example{
      \exsen{Bogo komlla kollba ulle de deyaittän.}      \extran{He caught two big fish.}}}}

\entry{kollba täkäll}{\headword{kollba täkäll}  \note{Poison}  \pronnote{koɽba təkəɽ}
  \pos{Noun}  \sense{
    \definition{fish fin}    \note{Poison}    \example{
}    \example{
}    \example{
}}}

\entry{kolldab}{\headword{kolldab}
  \pos{Verb}  \sense{
}}

\entry{kolldän}{\headword{kolldän}  \pronnote{koɽdən}
  \pos{Transitive verb}  \sense{
    \definition{to shoot; stab}    \example{
      \exsen{Sɨmell de enanae dukollädnegän pakos alle.}      \extran{He shot the pig straight through with the arrow.}}}}

\entry{kollko}{\headword{kollko}  \note{Medicine}  \pronnote{koɽko}
  \pos{Noun}  \sense{
    \definition{breadfruit}    \scientificname{Artocarpus altilis}    \note{Medicine}    \example{
}    \example{
}}}

\entry{kollkokollko}{\headword{kollkokollko}  \note{Trees}  \pronnote{koɽkokoɽko}
  \pos{Noun}  \sense{
    \definition{type of tree that grows in bush with edible red fruit with an edible, big round nut in the seed}    \note{Trees}    \example{
}    \example{
}    \example{
}}}

\entry{kollkollatt}{\headword{kollkollatt}
  \variantof{freevarlab of \vartext{kullkullatt}}}

\entry{kollmäll}{\headword{kollmäll}  \pronnote{koɽməɽ}
  \pos{Transitive verb}  \sense{
    \definition{to follow}    \example{
      \exsen{kollmällang e ebdo}      \extran{the following day}}    \example{
      \exsen{Llɨg a mänmän de dongkollmällaemeyo.}      \extran{The boys followed the girls.}}    \example{
      \exsen{Abo bongo mamamamall ngämim yangkollmällneg.}      \extran{You must follow us quickly.}}    \example{
      \exsen{Lama ngämo kollmällang dan.}      \extran{Lama is the second child after me (lit. Lama's following me).}}}
  \variants{\Unspecified Variant \vartext{kollmoll}}}

\entry{kollmällang}{\headword{kollmällang}  \pronnote{koɽməɽaŋ}
  \pos{Noun}  \sense{
    \definition{follower, disciple}    \example{
      \exsen{Yesu duwem gognän obo kolmällang aba peyang.}      \extran{Jesus ate with his disciples.}}}}

\entry{kollmoll}{\headword{kollmoll}
  \variantof{Unspecified Variant of \vartext{kollmäll}}}

\entry{kollmos ttam}{\headword{kollmos ttam}  \pronnote{koɽmos ʈ͡ʂam}
  \pos{Noun}  \sense{
    \definition{soul}    \example{
      \exsen{Ngämo kollmos ttam de nänddämän kuddäll ngättma we.}      \extran{It sunk my spirit to a place of death.}}}}

\entry{kolloekolloe}{\headword{kolloekolloe}
  \variantof{SpellingVariant of \vartext{klloklloe}}}

\entry{kollokolloe}{\headword{kollokolloe}
  \variantof{SpellingVariant of \vartext{klloklloe}}}

\entry{kollong}{\headword{kollong}  \note{Trees}  \pronnote{koɽoŋ}
  \pos{Noun}  \sense{
    \definition{type of tree that grows in the bush; used for grass skirts}    \note{Trees}    \example{
}    \example{
}    \example{
}}}

\entry{Kollwam}{\headword{Kollwam}
  \pos{Proper noun}  \sense{
    \definition{male personal name}}}

\entry{kollwany}{\headword{kollwany}  \pronnote{koɽwaɲ}
  \pos{Transitive verb}  \sense{
    \definition{to hang}    \example{
      \exsen{Pakätt de llo ngoeang me dukullwenyallo.}      \extran{They hung the robe on the tree branch.}}}}

\entry{Kolmet}{\headword{Kolmet}  \pronnote{kolmet}
  \pos{Proper noun}  \sense{
    \definition{female personal name}}}

\entry{Koloam}{\headword{Koloam}  \pronnote{koloam}
  \pos{Proper noun}  \sense{
    \definition{male personal name}}}

\entry{kolos}{\headword{kolos}
  \pos{Noun}  \sense{
    \definition{constable uniform}}
  \variants{\SpellingVariant \vartext{koles}}}

\entry{Kols}{\headword{Kols}  \pronnote{kols}
  \pos{Proper noun}  \sense{
    \definition{male personal name}}}

\entry{kolwa}{\headword{kolwa}  \note{Trees}  \pronnote{kolwa}
  \pos{Noun}  \sense{
    \definition{type of tree that grows in grassland with a trunk that has a diameter of ~9 cm, but very strong and used for spears and weapons}    \note{Trees}    \example{
}    \example{
}    \example{
}}}

\entry{kom}{\headword{kom}  \note{Parts of the body}  \pronnote{kom}
  \pos{Noun}  \sense{\textbf{1.} 
    \definition{hair; hair-like thing}    \note{Parts of the body}    \example{
      \exsen{bun kom, ttatt kom, mätta kom}      \extran{hair on head, beard, yam root hair}}    \example{
}}  \sense{\textbf{2.} 
    \definition{feather}    \note{Parts of the body}    \example{
      \exsen{Bogo ansi gogon a kom a pädrällag gognegän.}      \extran{He sneezed and the feathers flew.}}}}

\entry{komälle}{\headword{komälle}  \note{Games}  \pronnote{koməɽe}
  \pos{Noun}  \sense{
    \definition{type of game played with string}    \note{Games}    \example{
}    \example{
}    \example{
}}}

\entry{Kombosie}{\headword{Kombosie}  \pronnote{kombosie}
  \pos{Proper noun}  \sense{
    \definition{male personal name}}}

\entry{komene}{\headword{komene}  \note{Birth}  \pronnote{komene}
  \pos{Noun}  \sense{
    \definition{postpartum period during which a woman warms herself by the fire}    \note{Birth}    \example{
}    \example{
}    \example{
}}}

\entry{komengg}{\headword{komengg}
  \pos{Intransitive verb}  \sense{
    \definition{to go and come back}}}

\entry{komett}{\headword{komett}  \note{Government}  \pronnote{komeʈ͡ʂ}
  \pos{Noun}  \sense{
    \definition{committee}    \note{Government}    \example{
}    \example{
}    \example{
}}}

\entry{kominiti}{\headword{kominiti}
  \pos{Noun}  \sense{
    \definition{community}}
  \variants{\SpellingVariant \vartext{komuniti}}}

\entry{komiti}{\headword{komiti}
  \pos{Noun}  \sense{
    \definition{committee}}}

\entry{komlla}{\headword{komlla}  \note{Numbers}  \pronnote{komɽa}
  \pos{Numeral}  \sense{
    \definition{two (yam counting numeral; also general)}    \note{Numbers}    \example{
}    \example{
}    \example{
}}}

\entry{komlla komlla}{\headword{komlla komlla}  \note{Numbers}  \pronnote{komɽa komɽa}
  \pos{Numeral}  \sense{
    \definition{four (yam counting numeral; 2+2)}    \note{Numbers}    \example{
}    \example{
}    \example{
}}}

\entry{komlla komlla ttongo dduma}{\headword{komlla komlla ttongo dduma}  \note{Numbers}  \pronnote{komɽa komɽa ʈ͡ʂoŋo ɖ͡ʐuma}
  \pos{Numeral}  \sense{
    \definition{five (yam counting numeral; 2+2+1)}    \note{Numbers}    \example{
}    \example{
}    \example{
}}}

\entry{komlla putt}{\headword{komlla putt}  \note{Numbers}  \pronnote{komɽa puʈ͡ʂ}
  \pos{Numeral}  \sense{
    \definition{twelve (yam counting numeral; 2×6)}    \note{Numbers}    \example{
}    \example{
}    \example{
}}}

\entry{komlla ttuim mängallmeny}{\headword{komlla ttuim mängallmeny}  \note{Birth}  \pronnote{komɽa ʈ͡ʂuim məŋaɽmeɲ}
  \pos{}  \sense{
    \definition{two thighs weak}    \note{Birth}    \example{
}    \example{
}    \example{
}}}

\entry{komlla we}{\headword{komlla we}  \pronnote{komɽa we}
  \pos{Ordinal numeral}  \sense{
    \definition{second}    \example{
      \exsen{Komlla we nge da ngäna notan.}      \extran{This is the second coconut I ate today.}}    \example{
}    \example{
}}}

\entry{komlla yangete}{\headword{komlla yangete}  \note{Numbers}  \pronnote{komɽa jaŋete}
  \pos{Numeral}  \sense{
    \definition{doubling two, e.g. shoes tied together}    \note{Numbers}    \example{
}    \example{
}    \example{
}}}

\entry{komllaebe}{\headword{komllaebe}  \pronnote{komɽaebe}
  \pos{Numeral}  \sense{
    \definition{exactly two; both}    \example{
      \exsen{Mänyanda bi komllaebe ade tatu de gongkameyo.}      \extran{The younger two also started to wash.}}}
  \variants{\freevarlab \vartext{komllaebme}}}

\entry{komllaebekomllaebe}{\headword{komllaebekomllaebe}
  \pos{Numeral}  \sense{
    \definition{exactly four (2+2)}}}

\entry{komllaebme}{\headword{komllaebme}
  \variantof{freevarlab of \vartext{komllaebe}}}

\entry{komllakomlla}{\headword{komllakomlla}  \pronnote{komɽakomɽa}
  \pos{Quantifier}  \sense{
    \definition{both}    \example{
      \exsen{Bogo komllakomlla panypenyang deya ada Tame eka da wa Ende eka da.}      \extran{She spoke both Taeme and Ende languages.}}}}

\entry{komllazegatt}{\headword{komllazegatt}
  \pos{Noun}  \sense{
}}

\entry{komlle}{\headword{komlle}  \pronnote{komɽe}
  \pos{Noun}  \sense{
    \definition{type of game played with string}}}

\entry{komo}{\headword{komo}  \note{Meresen}  \pronnote{komo}
  \pos{Noun}  \sense{
    \definition{ginger}    \note{Meresen}}}

\entry{komo}{\headword{komo}  \note{Meresen}  \pronnote{komo}
  \pos{Noun}  \sense{\textbf{1.} 
    \definition{centipede}    \note{Meresen}    \example{
      \exsen{Maryanne bom komo da näddägan.}      \extran{Maryanne was bit by a centipede.}}    \example{
}}  \sense{\textbf{2.} 
    \definition{scorpion}    \note{Meresen}}
  \variants{\Unspecified Variant \vartext{kämo}}}

\entry{komo mogmog}{\headword{komo mogmog}  \note{Poison}  \pronnote{komo moɡmoɡ}
  \pos{Noun}  \sense{
    \definition{type of black scorpion}    \note{Poison}    \example{
}    \example{
}    \example{
}}}

\entry{komo takle}{\headword{komo takle}  \note{Poison}  \pronnote{komo takle}
  \pos{Noun}  \sense{
    \definition{type of scorpion}    \note{Poison}    \example{
}    \example{
}    \example{
}}}

\entry{komongg}{\headword{komongg}
  \pos{Ditransitive verb}  \sense{
    \definition{causative-applicative form of ony (to make carry; bring for, take to; take from)}    \example{
      \exsen{Grace ngänäm gull de dangkomonggän.}      \extran{Grace had me bring the net.}}    \example{
      \exsen{Bong obo pope sɨmell midd de dangkomonggän.}      \extran{Bong brought the pork for his uncle.}}    \example{
      \exsen{Ttall a obo za de dangkämonggän.}      \extran{The wallaby took her thing.}}}}

\entry{komony}{\headword{komony}
  \pos{Noun}  \sense{
}}

\entry{komotupi}{\headword{komotupi}  \pronnote{komotupi}
  \pos{Noun}  \sense{
    \definition{long ginger}    \example{
}    \example{
}    \example{
}}}

\entry{komotupi molle molle}{\headword{komotupi molle molle}  \note{Dancing}  \pronnote{komotupi moɽe moɽe}
  \pos{Noun}  \sense{
    \definition{type of dancing game that involves ginger}    \note{Dancing}    \example{
}    \example{
}    \example{
}}}

\entry{komuniti}{\headword{komuniti}
  \variantof{SpellingVariant of \vartext{kominiti}}}

\entry{kona}{\headword{kona}
  \pos{Noun}  \sense{\textbf{1.} 
    \definition{corner}    \example{
      \exsen{Obo ma da duli kona me dan.}      \extran{His house is over there, on the corner.}}}  \sense{\textbf{2.} 
    \definition{district, section, area (of a settlement)}    \example{
      \exsen{Ngäna sisri Llimoll me giddollnen allan Basido kona me.}      \extran{I am now living in Limol, in the Basido district.}}}}

\entry{konakone}{\headword{konakone}  \pronnote{konakone}
  \pos{Transitive verb}  \sense{\textbf{1.} 
    \definition{to cover}    \example{
      \exsen{Yäbäd de ddapall käkan da dakonewän.}      \extran{Clouds covered the sun.}}}  \sense{\textbf{2.} 
    \definition{cover, sheet, blanket}    \example{
      \exsen{Däbe män a konakone alle gokonewän.}      \extran{That girl covered up with a blanket.}}}
  \variants{\Unspecified Variant \vartext{känakone}}}

\entry{Kondobäll}{\headword{Kondobäll}
  \variantof{Unspecified Variant of \vartext{Kondobol}}}

\entry{Kondobol}{\headword{Kondobol}  \pronnote{kondobol}
  \pos{Proper noun}  \sense{
    \definition{Kondobol (Taeme-speaking village in Morehead Rural LLG; from Limol, one must pass through Kinkin)}    \example{
}    \example{
}    \example{
}}
  \variants{\Unspecified Variant \vartext{Kondobäll, Kondoboll}}}

\entry{Kondoboll}{\headword{Kondoboll}
  \variantof{Unspecified Variant of \vartext{Kondobol}}
  \variants{\SpellingVariant \vartext{Kondobäll}}}

\entry{Kondobu}{\headword{Kondobu}  \pronnote{kondobu}
  \pos{Proper noun}  \sense{
    \definition{Konedobu (in Gogodala Rural LLG)}    \example{
}    \example{
}    \example{
}}}

\entry{kongkottom}{\headword{kongkottom}  \pronnote{koŋkoʈ͡ʂom}
  \pos{Transitive verb}  \sense{
    \definition{to gulp, drink quickly}}}

\entry{konkon}{\headword{konkon}  \pronnote{konkon}
  \pos{Modifier}  \sense{\textbf{1.} 
    \definition{crazy, mad, insane, mentally ill}    \example{
      \exsen{Ge ause da konkonang daeya.}      \extran{This old woman was crazy.}}    \example{
      \exsen{Ngäna konkon allan.}      \extran{I'm getting mad.}}    \example{
      \exsen{Ngämo yae konkonang gogon.}      \extran{My mother became mentally ill.}}}  \sense{\textbf{2.} 
    \definition{stupid, ignorant, foolish}    \example{
      \exsen{Ngäna konkonang dan, ngämo kame dan bina eka da.}      \extran{I am ignorant; I don't know your (pl.) language.}}}  \sense{\textbf{3.} 
    \definition{intoxicated, intoxicating, consciousness-altering, drunk}    \example{
      \exsen{ine konkon me}      \extran{drunk}}    \example{
      \exsen{Ine da ngänäm gagäll konkonang dagän.}      \extran{The drink made me really drunk.}}}}

\entry{konskak}{\headword{konskak}  \note{Types of Taro}  \pronnote{konskak}
  \pos{Noun}  \sense{
    \definition{type of big taro}    \note{Types of Taro}    \example{
}    \example{
}    \example{
}}}

\entry{konykony}{\headword{konykony}  \note{Poison}  \pronnote{koɲkoɲ}
  \pos{Noun}  \sense{
    \definition{type of stinging insect that lives in the ground}    \note{Poison}    \example{
}    \example{
}    \example{
}}}

\entry{konymad}{\headword{konymad}  \pronnote{koɲmad}
  \pos{Noun}  \sense{
    \definition{man who steals a woman's things during her initiation ceremony}}}

\entry{konzar}{\headword{konzar}  \lexeme{konzar}  \pronnote{konzar}
  \pos{Noun}  \sense{
    \definition{type of small yam with a white interior, hairs, and small thorns}}}

\entry{kopae därängmeny}{\headword{kopae därängmeny}  \pronnote{kopae dərəŋmeɲ}
  \pos{Noun}  \sense{
    \definition{to go to another village to trade}    \example{
      \exsen{Llimollang a kopae därängmeny e Wim e abällan.}      \extran{Limol people went to Wim to trade.}}}}

\entry{kopek}{\headword{kopek}  \pronnote{kopek}
  \pos{Noun}  \sense{\textbf{1.} 
    \definition{valley}    \example{
      \exsen{Ngängälatt dädär a mänyi tutu atta bolldaeyän kopek e.}      \extran{A round rock will roll from the hill to the valley.}}}  \sense{\textbf{2.} 
    \definition{pit, hole}    \example{
      \exsen{källama kopek}      \extran{toilet (lit. 'hole for pooping')}}}}

\entry{kopllalle}{\headword{kopllalle}  \pronnote{kopɽaɽe}
  \pos{Noun}  \sense{
    \definition{oriole (brown, olive-backed, green)}    \scientificname{Oriolus szalayi, O. sagittatus, O. flavocinctus}    \example{
}    \example{
}    \example{
}}}

\entry{koplle}{\headword{koplle}  \note{Trees}  \pronnote{kopɽe}
  \pos{Noun}  \sense{
    \definition{type of big tree that grows in the bush with fruit that cassowaries eat}    \note{Trees}    \example{
}    \example{
}    \example{
}}}

\entry{koplle täkmäl täkmäl}{\headword{koplle täkmäl täkmäl}  \note{Games}  \pronnote{kopɽe təkməl təkməl}
  \pos{Noun}  \sense{
    \definition{type of game involving throwing koplle fruit}    \note{Games}    \example{
}    \example{
}    \example{
}}}

\entry{Koreya}{\headword{Koreya}  \pronnote{koreja}
  \pos{Proper noun}  \sense{
    \definition{Korea}}}

\entry{kormas}{\headword{kormas}  \note{Birds}  \pronnote{kormas}
  \pos{Noun}  \sense{
    \definition{type of bird}    \note{Birds}    \example{
}    \example{
}    \example{
}}}

\entry{korolläm}{\headword{korolläm}  \note{Palms}  \pronnote{koroɽəm}
  \pos{Noun}  \sense{
    \definition{type of vine that grows beside creeks. Leaves are used to wrap things.}    \note{Palms}    \example{
}    \example{
}    \example{
}}}

\entry{kos}{\headword{kos}
  \pos{Noun}  \sense{
    \definition{course}}}

\entry{kosrom}{\headword{kosrom}  \note{Mushrooms}  \pronnote{kosrom}
  \pos{Noun}  \sense{
    \definition{type of large mushroom that grows in the winter}    \note{Mushrooms}    \example{
}    \example{
}    \example{
}}}

\entry{kot}{\headword{kot}
  \variantof{Unspecified Variant of \vartext{kwatt}}}

\entry{kot}{\headword{kot}  \pronnote{kot}
  \pos{Noun}  \sense{
    \definition{dirt}}}

\entry{kotang}{\headword{kotang}  \note{Adjectives}  \pronnote{kotaŋ}
  \pos{Modifier}  \sense{
    \definition{dirty, unclean}    \note{Adjectives}    \example{
      \exsen{ttang kotang peyang duwem}      \extran{eating with dirty hands}}    \example{
      \exsen{Ngäna sisri kaptte kotang de gällämnan anggan.}      \extran{I'm cleaning the dirty clothes now.}}}}

\entry{kote}{\headword{kote}  \note{Parts of the Body}  \pronnote{kote}
  \pos{Noun}  \sense{
    \definition{nape}    \note{Parts of the Body}    \example{
}    \example{
}    \example{
}}}

\entry{kote kutt}{\headword{kote kutt}  \note{Parts of a reptile}  \pronnote{kote kuʈ͡ʂ}
  \pos{Noun}  \sense{
    \definition{cervical vertebrae}    \note{Parts of a reptile}    \example{
}    \example{
}    \example{
}}}

\entry{kote ttäkam}{\headword{kote ttäkam}
  \pos{Intransitive compound verb}  \sense{
    \definition{to hang one's head}    \example{
      \exsen{Lla da kote gottkamän.}      \extran{The man hung his head.}}}}

\entry{kotkot}{\headword{kotkot}  \pronnote{kotkot}
  \pos{Modifier}  \sense{
    \definition{dirty, unclean}    \example{
      \exsen{gagäll kotkot anyke}      \extran{bad, unclean spirit}}}}

\entry{kotmeny}{\headword{kotmeny}  \note{Water}  \pronnote{kotmeɲ}
  \pos{Modifier}  \sense{
    \definition{clean, pure}    \note{Water}    \example{
      \exsen{Ngämo umllang dan bongo aenen, bongo Adi bo Kotmeny dan!}      \extran{I know who you are; you are God's pure one!}}    \example{
}    \example{
}}}

\entry{kotol}{\headword{kotol}  \note{Birth}  \pronnote{kotol}
  \pos{Noun}  \sense{\textbf{1.} 
    \definition{traditional practice of sterilizing women after having 6–7 kids; the placenta is buried and a coconut is planted on top}    \note{Birth}}  \sense{\textbf{2.} 
    \definition{traditional practice of stopping smoking}    \note{Birth}}}

\entry{kotom}{\headword{kotom}  \pronnote{kotom}
  \pos{Noun}  \sense{
    \definition{head covering, crown}}}

\entry{kott}{\headword{kott}  \pronnote{koʈ͡ʂ}
  \pos{Noun}  \sense{
    \definition{testicles}}}

\entry{kottllam}{\headword{kottllam}  \pronnote{koʈ͡ʂɽam}
  \pos{Noun}  \sense{
    \definition{turtle, tortoise}    \example{
      \exsen{Kottllam a ttongo källäm me dazernän.}      \extran{The turtles were living in a pond.}}}}

\entry{kottma}{\headword{kottma}  \note{Parts of the body}  \pronnote{koʈ͡ʂma}
  \pos{Noun}  \sense{
    \definition{body part}    \note{Parts of the body}    \example{
}    \example{
}    \example{
}}}

\entry{kowa}{\headword{kowa}
  \pos{Noun}  \sense{
    \definition{upper back}}}

\entry{kowatt}{\headword{kowatt}
  \variantof{SpellingVariant of \vartext{kwatt}}}

\entry{kowatt ikopang}{\headword{kowatt ikopang}  \note{Government}  \pronnote{kowaʈ͡ʂ ikopaŋ}
  \pos{Noun}  \sense{
    \definition{law and order}    \note{Government}    \example{
}    \example{
}    \example{
}}}

\entry{Kral}{\headword{Kral}  \pronnote{kral}
  \pos{Proper noun}  \sense{
    \definition{female personal name}}}

\entry{krismas}{\headword{krismas}  \pronnote{krismas}
  \pos{Noun}  \sense{
    \definition{Christmas}}}

\entry{kristen}{\headword{kristen}
  \pos{Modifier}  \sense{
    \definition{Christian}}
  \variants{\SpellingVariant \vartext{kristin}}}

\entry{kristin}{\headword{kristin}
  \variantof{SpellingVariant of \vartext{kristen}}}

\entry{Kristina}{\headword{Kristina}
  \pos{Proper noun}  \sense{
    \definition{female personal name}}}

\entry{KT}{\headword{KT}
  \variantof{SpellingVariant of \vartext{Keti}}  \pronnote{kt}}

\entry{ku}{\headword{ku}  \pronnote{ku}
  \pos{Locational}  \sense{
    \definition{center, core, middle}    \example{
      \exsen{nyongo ku me}      \extran{in the middle of the road}}}
  \variants{\Unspecified Variant \vartext{kuu}}}

\entry{ku}{\headword{ku}  \pronnote{ku}
  \pos{Intransitive verb}  \sense{
    \definition{to wake up}    \example{
      \exsen{Ddia da dukunegalle inu atta.}      \extran{The deer woke up.}}}}

\entry{kuang}{\headword{kuang}  \pronnote{kuaŋ}
  \pos{Noun}  \sense{
    \definition{the ninth stage of coconut growth in which the endosperm of the fruit is fully formed and coconut water remains}    \example{
      \exsen{Da bongo nge de nängkäl, abo kuang dae nädd.}}}}

\entry{kub}{\headword{kub}  \pronnote{kub}
  \pos{Noun}  \sense{\textbf{1.} 
    \definition{flesh, meat}}  \sense{\textbf{2.} 
    \definition{shooting star}}  \sense{\textbf{3.} 
    \definition{burning red fire}}}

\entry{kubäll}{\headword{kubäll}
  \variantof{SpellingVariant of \vartext{kubull}}  \pronnote{kubəɽ}}

\entry{kube}{\headword{kube}  \pronnote{kube}
  \pos{Noun}  \sense{
    \definition{bucket}    \example{
      \exsen{Masar ine kube de ngämo nangesan.}      \extran{Grandfather made me a water bucket.}}}}

\entry{kubllu}{\headword{kubllu}  \lexeme{kubllu}  \pronnote{kubɽu}
  \pos{Noun}  \sense{
    \definition{type of tree that grows in the bush and savannah with bark used for string; similar to kapang}}}

\entry{kubull}{\headword{kubull}  \note{Animals/Plants}  \pronnote{kubuɽ}
  \pos{Noun}  \sense{
    \definition{bush wallaby, dusky pademelon}    \scientificname{Thylogale brunii}    \note{Animals/Plants}    \example{
}    \example{
}    \example{
}}
  \variants{\SpellingVariant \vartext{kubäll}}}

\entry{kud}{\headword{kud}  \note{Type of pandanus}  \pronnote{kud}
  \pos{Noun}  \sense{
    \definition{type of pandanus with fat triangular fruit}    \note{Type of pandanus}    \example{
}    \example{
}    \example{
}}}

\entry{kud dämadämar}{\headword{kud dämadämar}
  \pos{Noun}  \sense{
    \definition{dragonfly}}}

\entry{kudäd}{\headword{kudäd}  \note{Birth}  \pronnote{kudəd}
  \pos{Modifier}  \sense{
    \definition{barren}    \note{Birth}    \example{
}    \example{
}    \example{
}}}

\entry{kudädäri}{\headword{kudädäri}  \pronnote{kudədəri}
  \pos{Noun}  \sense{
    \definition{Zoe's imperial pigeon}    \scientificname{Ducula zoeae}    \example{
}    \example{
}    \example{
}}}

\entry{kuddäb}{\headword{kuddäb}  \pronnote{kuɖ͡ʐəb}
  \pos{Noun}  \sense{
    \definition{raft}}}

\entry{kuddäkuddäll}{\headword{kuddäkuddäll}  \pronnote{kuɖ͡ʐəkuɖ͡ʐəɽ}
  \pos{Modifier}  \sense{\textbf{1.} 
    \definition{soft}    \example{
      \exsen{Mätta da kuddällkuddäll bognegän.}      \extran{The yams will become soft.}}}  \sense{\textbf{2.} 
    \definition{easy}    \example{
      \exsen{Orpmang eka da kuddäkuddäll a era ngäna ai dan bondär.}      \extran{Wipi language is easy; I can understand it.}}    \example{
      \exsen{Kuddällkuddäll dan baob wätät e.}      \extran{It's easy to eat waterlily.}}}
  \variants{\Unspecified Variant \vartext{kuddällkuddäll, kuddukuddäll}}
  \variants{\SpellingVariant \vartext{kuddukuddull}}}

\entry{kuddäll}{\headword{kuddäll}  \pronnote{kuɖ͡ʐəɽ}
  \pos{Noun}  \sense{\textbf{1.} 
    \definition{death}    \example{
      \exsen{Gagäll dan lla da bem e buspunän kuddäll e.}      \extran{It is bad for a man to throw himself into the ocean to his death.}}}  \sense{\textbf{2.} 
    \definition{to die}    \example{
      \exsen{Gullem a kuddäll agan.}      \extran{The snake died.}}}  \sense{\textbf{3.} 
    \definition{dead}    \example{
      \exsen{sɨmell kuddäll}      \extran{dead pig}}    \example{
      \exsen{kuddäll pätt}      \extran{dead body}}}  \sense{\textbf{4.} 
    \definition{abandoned}    \example{
      \exsen{ma kuddäll, ttängäm kuddäll}      \extran{abandoned house, abandoned garden}}}
  \variants{\SpellingVariant \vartext{kuddɨll}}}

\entry{Kuddäll Opap}{\headword{Kuddäll Opap}
  \pos{Proper noun}  \sense{
    \definition{Passover}}}

\entry{kuddäll ttam}{\headword{kuddäll ttam}
  \pos{Adverb}  \sense{
    \definition{life-or-death, to death, as if one may die}    \example{
      \exsen{Kuddäll ttam ngäna dindug.}      \extran{I ran for my life.}}    \example{
      \exsen{Ubi kuddäll ttam gotongoenegnän.}      \extran{They were laughing to death.}}}}

\entry{kuddällkuddäll}{\headword{kuddällkuddäll}
  \variantof{Unspecified Variant of \vartext{kuddäkuddäll}}}

\entry{kuddɨll}{\headword{kuddɨll}
  \variantof{SpellingVariant of \vartext{kuddäll}}}

\entry{kuddukuddäll}{\headword{kuddukuddäll}
  \variantof{Unspecified Variant of \vartext{kuddäkuddäll}}  \pronnote{kuɖ͡ʐukuɖ͡ʐəɽ}}

\entry{kuddukuddull}{\headword{kuddukuddull}
  \variantof{SpellingVariant of \vartext{kuddäkuddäll}}}

\entry{kudu}{\headword{kudu}
  \pos{Noun}  \sense{
    \definition{corner}}}

\entry{Kudurwe}{\headword{Kudurwe}  \pronnote{kudurwe}
  \pos{Proper noun}  \sense{
    \definition{female personal name}}}

\entry{Kui}{\headword{Kui}
  \pos{Proper noun}  \sense{
    \definition{Kui (toponym)}}}

\entry{kui}{\headword{kui}  \pronnote{kui}
  \pos{Noun}  \sense{
    \definition{island}    \example{
      \exsen{Ubi deagerneyo ttongo kui me.}      \extran{They (two) lived on an island.}}}}

\entry{kuib}{\headword{kuib}  \note{Musical instruments}  \pronnote{kuib}
  \pos{Noun}  \sense{
    \definition{type of drum}    \note{Musical instruments}    \example{
}    \example{
}    \example{
}}}

\entry{kuibiag}{\headword{kuibiag}  \lexeme{kuibiag}  \pronnote{kuibiaɡ}
  \pos{Noun}  \sense{
    \definition{Papuan black snake}    \scientificname{Pseudechis papuanus}    \example{
      \exsen{Llɨg kälekäle, towall e ibnin a mudag, kuibiag a amom bäddägän.}      \extran{Children, do not walk in the grass; Papuan black snacks will bite anyone.}}}}

\entry{kuimang}{\headword{kuimang}
  \pos{Transitive coverb}  \sense{
    \definition{to knock on}    \example{
      \exsen{Obo ma ud de kuimang dägagän.}      \extran{She knocked on his door.}}}}

\entry{Kuiwang}{\headword{Kuiwang}  \note{Place names}  \pronnote{kuiwaŋ}
  \pos{Proper noun}  \sense{
    \definition{Kuiwang (Taeme-speaking village in Morehead Rural LLG; from Limol, one must pass through Malam)}    \note{Place names}    \example{
}    \example{
}    \example{
}}}

\entry{kuki}{\headword{kuki}  \pronnote{kuki}
  \pos{Modifier}  \sense{\textbf{1.} 
    \definition{false, deceptive}    \example{
      \exsen{kuki ttoen}      \extran{trick}}    \example{
      \exsen{kukiang keriso}      \extran{false messiah}}}  \sense{\textbf{2.} 
    \definition{to deceive, trick, lie to}    \example{
      \exsen{Bogo gongnamän ada mälla da obom era kuki dägagän.}      \extran{He realized that the woman deceived him.}}    \example{
      \exsen{Ddob kukiang lla da mɨnyi bälltaemnän.}      \extran{Other deceiving men will be coming.}}    \example{
      \exsen{Ai dan God ngänäm märal bagän da kuki allan.}      \extran{I swear to God I'm not lying (lit. God can curse me if I'm lying).}}}}

\entry{kuki eka}{\headword{kuki eka}
  \pos{Noun}  \sense{
    \definition{lie}    \example{
      \exsen{Ubi sapasapang kuki eka de däpanyaemeyo Yesu bo pallall.}      \extran{They told various lies about Jesus.}}}}

\entry{kukiny}{\headword{kukiny}  \lexeme{kukiny}  \pronnote{kukiɲ}
  \pos{Noun}  \sense{\textbf{1.} 
    \definition{type of long-leaf grass}}  \sense{\textbf{2.} 
    \definition{type of spear made from kukiny grass that is used to hunt birds}}}

\entry{kukoll}{\headword{kukoll}  \pronnote{kukoɽ}
  \pos{Modifier}  \sense{
    \definition{healthy, fertile, vibrant}    \example{
      \exsen{Kutt a kukollang ekaklle me guspullän a mermerae mängallang kukollang däpäddän.}      \extran{The seed fell to the fertile ground and sprouted vibrantly.}}}}

\entry{kukollkukoll}{\headword{kukollkukoll}  \note{Colors}  \pronnote{kukoɽkukoɽ}
  \pos{Color term}  \sense{
    \definition{green}    \note{Colors}    \example{
}    \example{
}    \example{
}}}

\entry{kukpi}{\headword{kukpi}  \note{Trees}  \pronnote{kukpi}
  \pos{Noun}  \sense{
    \definition{type of tree that grows in the bush with big fruit that children play with}    \note{Trees}    \example{
}    \example{
}    \example{
}}}

\entry{Kukpikukpi}{\headword{Kukpikukpi}  \pronnote{kukpikukpi}
  \pos{Proper noun}  \sense{
    \definition{Kukpikukpi (toponym)}}}

\entry{kukrub}{\headword{kukrub}  \pronnote{kukrub}
  \pos{Noun}  \sense{\textbf{1.} 
    \definition{immature coconut that is fully liquid inside}    \example{
      \exsen{Ubi tätäm nge kukrub doklloaebeyo a nge dädär peyang gongkänemenynegnän.}      \extran{Yesterday, they scratched water-filled coconuts and ate them with dry coconut.}}}  \sense{\textbf{2.} 
    \definition{the seventh stage of coconut growth in which the endosperm of the fruit is completely liquid}}}

\entry{Kuks}{\headword{Kuks}
  \pos{Proper noun}  \sense{
    \definition{male personal name}}}

\entry{kuku}{\headword{kuku}
  \pos{Noun}  \sense{
    \definition{type of grass}}}

\entry{Kukua}{\headword{Kukua}  \pronnote{kukua}
  \pos{Proper noun}  \sense{
    \definition{male personal name}}}

\entry{kukull}{\headword{kukull}
  \variantof{Fast Speech Variant of \vartext{kullkull}}  \pronnote{kukuɽ}}

\entry{Kukumi}{\headword{Kukumi}
  \pos{Proper noun}  \sense{
    \definition{male personal name}}}

\entry{kul}{\headword{kul}  \pronnote{kul}
  \pos{Transitive coverb}  \sense{
    \definition{to smash food}    \example{
      \exsen{Pamker a mätta de kul bägayaebneyo.}      \extran{They will be smashing the yams and pumpkins together.}}}}

\entry{kulläb}{\headword{kulläb}  \note{Parts of a fire}  \pronnote{kuɽəb}
  \pos{Noun}  \sense{
    \definition{large black termite mound}    \note{Parts of a fire}    \example{
}    \example{
}    \example{
}}}

\entry{Kulläntti}{\headword{Kulläntti}
  \variantof{Unspecified Variant of \vartext{Kurunti}}}

\entry{Kullintti}{\headword{Kullintti}
  \variantof{Unspecified Variant of \vartext{Kurunti}}}

\entry{kullkull}{\headword{kullkull}  \pronnote{kuɽkuɽ}
  \pos{Noun}  \sense{
    \definition{grassfire; burnt grass}    \example{
      \exsen{Kullkull a dallän llamda pate.}      \extran{The grassfire moved towards the old man.}}    \example{
      \exsen{Ngäna deyarne kullkull att dae.}      \extran{I was returning to the burnt grass.}}}
  \variants{\Fast Speech Variant \vartext{kukull}}}

\entry{kullkullatt}{\headword{kullkullatt}  \pronnote{kuɽkuɽaʈ͡ʂ}
  \pos{Modifier}  \sense{
    \definition{burnt (of an area)}    \example{
      \exsen{Ttongo däräng a kullkullat enddäna me pollpoll dängkamän.}      \extran{A dog started the bark in the burnt-up clearing.}}}
  \variants{\freevarlab \vartext{kollkollatt}}}

\entry{Kullme}{\headword{Kullme}  \note{Places around Limol}  \pronnote{kuɽme}
  \pos{Proper noun}  \sense{
    \definition{Kullme (garden place near Egapo; filled with abandoned rubber trees)}    \note{Places around Limol}    \example{
}    \example{
}    \example{
}}}

\entry{Kullopang}{\headword{Kullopang}  \note{Places around Limol}  \pronnote{kuɽopaŋ}
  \pos{Proper noun}  \sense{
    \definition{Kullopang (sago and garden place of Kaoga Dobola; on the shortcut road to Kinkin)}    \note{Places around Limol}    \example{
}    \example{
}    \example{
}}}

\entry{kullum}{\headword{kullum}  \pronnote{kuɽum}
  \pos{Noun}  \sense{
    \definition{group}    \example{
      \exsen{Da ttongo kantri me lla da bompallängkmenynegän obaoba komlla kullum e a bogäddeyo obaoba, däbe kantri da mɨnyi ddone bakällmewän.}      \extran{If in a country, the people divide themselves into two groups and fight each other, that country will not survive.}}}
  \variants{\Fast Speech Variant \vartext{kllum}}}

\entry{Kullwam}{\headword{Kullwam}
  \pos{Proper noun}  \sense{
    \definition{male personal name}}}

\entry{kum}{\headword{kum}  \pronnote{kum}
  \pos{Noun}  \sense{\textbf{1.} 
    \definition{buttocks, butt}    \example{
}}  \sense{\textbf{2.} 
    \definition{abdomen of an insect}    \example{
}}  \sense{\textbf{3.} 
    \definition{butt, base, bottom, lower or back part of something}    \example{
      \exsen{Mätta bo kum a angttägan matu ik dae.}      \extran{The base of the yam comes up through the ground.}}}}

\entry{kumattkumatt}{\headword{kumattkumatt}
  \pos{Adverb}  \sense{
    \definition{back}    \example{
      \exsen{Bogo kumattkumatt gongosän.}      \extran{He turned and went back.}}}}

\entry{kumddäga}{\headword{kumddäga}
  \variantof{Fast Speech Variant of \vartext{kumuddäga}}  \pronnote{kumɖ͡ʐəɡa}}

\entry{kumi}{\headword{kumi}  \pronnote{kumi}
  \pos{Noun}  \sense{
    \definition{central roof beam}}}

\entry{kumi kakälläp}{\headword{kumi kakälläp}  \pronnote{kumi kakəɽəp}
  \pos{Noun}  \sense{
    \definition{bark on top of roof}}}

\entry{kumi spun}{\headword{kumi spun}  \pronnote{kumi spun}
  \pos{Intransitive compound verb}  \sense{
    \definition{to put the center roof post (kumi) on.}}}

\entry{kumie}{\headword{kumie}
  \variantof{SpellingVariant of \vartext{kumiye}}  \pronnote{kumie}}

\entry{kumiye}{\headword{kumiye}
  \pos{Noun}  \sense{
    \definition{cough}    \example{
      \exsen{Llamda da kumiye agan.}      \extran{The old man coughed.}}}
  \variants{\SpellingVariant \vartext{kumye, kumie}}}

\entry{kumiye käp}{\headword{kumiye käp}  \pronnote{kumije kəp}
  \pos{Noun}  \sense{\textbf{1.} 
    \definition{phlegm}}  \sense{\textbf{2.} 
    \definition{part of a fish}}  \sense{\textbf{3.} 
    \definition{bark on roof ridge that prevents rain from coming in}    \example{
}    \example{
}    \example{
}}}

\entry{kumpit}{\headword{kumpit}  \note{Poison}  \pronnote{kumpit}
  \pos{Noun}  \sense{
    \definition{stinger (of an insect)}    \note{Poison}    \example{
}    \example{
}    \example{
}}}

\entry{kumuddäga}{\headword{kumuddäga}  \note{Numbers}  \pronnote{kumuɖ͡ʐəɡa}
  \pos{Numeral}  \sense{
    \definition{three (yam counting numeral; also general)}    \note{Numbers}    \example{
      \exsen{Ag mullaemullae kumuddäga täräp pauro eka watt me.}      \extran{Every morning, I wake up after the third rooster crow.}}    \example{
}    \example{
}}
  \variants{\Fast Speech Variant \vartext{kumddäga}}}

\entry{kumuddäga we}{\headword{kumuddäga we}  \note{Everyday words and phrases}  \pronnote{kumuɖ͡ʐəɡa we}
  \pos{Ordinal numeral}  \sense{
    \definition{third}    \note{Everyday words and phrases}    \example{
}    \example{
}    \example{
}}}

\entry{kumuddägaebe}{\headword{kumuddägaebe}
  \pos{Numeral}  \sense{
    \definition{exactly three}}}

\entry{kumuddägakumuddäga}{\headword{kumuddägakumuddäga}
  \pos{Numeral}  \sense{
    \definition{six (3+3)}}}

\entry{Kumull}{\headword{Kumull}
  \pos{Proper noun}  \sense{
    \definition{Kumull (toponym)}}}

\entry{kumye}{\headword{kumye}
  \variantof{SpellingVariant of \vartext{kumiye}}}

\entry{Kuna}{\headword{Kuna}
  \pos{Proper noun}  \sense{
    \definition{male personal name}}}

\entry{kunen}{\headword{kunen}
  \pos{Intransitive verb}  \sense{
    \definition{to flee, scatter, run away}    \example{
      \exsen{Nikol angde walle menae e nallan, ddone ada kollba nukuaemallo.}      \extran{When Nikol went to the side of the river, many fish scattered away.}}}}

\entry{kungge}{\headword{kungge}  \note{Trees}  \pronnote{kuŋɡe}
  \pos{Noun}  \sense{
    \definition{type of tree}    \note{Trees}    \example{
}    \example{
}    \example{
}}}

\entry{kungge}{\headword{kungge}  \note{Poison}  \pronnote{kuŋɡe}
  \pos{Noun}  \sense{
    \definition{spider}    \note{Poison}    \example{
      \exsen{Da bam kungge da naddägän, bongo mɨnyi ttattlle ampug.}      \extran{If a spider bites you, you will be very ill.}}    \example{
}    \example{
}}}

\entry{kunglle}{\headword{kunglle}
  \pos{Ideophone}  \sense{
    \definition{sound made by flying foxes}}}

\entry{kunob}{\headword{kunob}  \note{Trees}  \pronnote{kunob}
  \pos{Noun}  \sense{
    \definition{type of tree that grows around creeks with light wood used for house sticks}    \note{Trees}    \example{
}    \example{
}    \example{
}}}

\entry{kunu}{\headword{kunu}  \pronnote{kunu}
  \pos{Noun}  \sense{
    \definition{type of short tree that grows in the grassland with poisonous bark for stunning fish}    \example{
}    \example{
}    \example{
}}}

\entry{kunun}{\headword{kunun}  \note{Seasons}  \pronnote{kunun}
  \pos{Noun}  \sense{
    \definition{season when crops are ready to be harvested (sixth season; corresponds to April)}    \note{Seasons}    \example{
}    \example{
}    \example{
}}}

\entry{kunur}{\headword{kunur}  \lexeme{kunur}  \pronnote{kunur}
  \pos{Noun}  \sense{
    \definition{corn, maize}    \example{
      \exsen{Ttongo lla da kunur däm kutt de ekaklle me bibewän.}      \extran{A man will plant the corn kernels in the ground.}}}}

\entry{kunur käp}{\headword{kunur käp}
  \pos{Noun}  \sense{
    \definition{corn kernel}}}

\entry{kunuwälläb}{\headword{kunuwälläb}  \note{Types of Taro}  \pronnote{kunuwəɽəb}
  \pos{Noun}  \sense{
    \definition{type of big taro}    \note{Types of Taro}    \example{
}    \example{
}    \example{
}}}

\entry{Kunyemäll}{\headword{Kunyemäll}  \note{Places around Limol}  \pronnote{kuɲeməɽ}
  \pos{Proper noun}  \sense{
    \definition{Kunyemäll (on the road to Malam near Zarma; filled with black palms that were cut for the school)}    \note{Places around Limol}    \example{
}    \example{
}    \example{
}}}

\entry{kup}{\headword{kup}  \pronnote{kup}
  \pos{Noun}  \sense{\textbf{1.} 
    \definition{hole, pit}    \example{
      \exsen{Ngäna ngämo män kälsäre de umllang dägag, ada ibi sisri mɨnyi kup di bäklleya.}      \extran{I told my little girl that we are going to dig a hole now.}}}  \sense{\textbf{2.} 
    \definition{well}    \example{
      \exsen{Ine da bodo dan kup mi.}      \extran{The well is full of water.}}}  \sense{\textbf{3.} 
    \definition{valley}    \example{
      \exsen{Obo ma duli dan kup mi.}      \extran{His house is there in the valley.}}}
  \variants{\Unspecified Variant \vartext{kwäp}}}

\entry{kup~}{\headword{kup~}
  \variantof{Derived Variant of \vartext{RED~}}}

\entry{kupi käp}{\headword{kupi käp}  \pronnote{kupi kəp}
  \pos{Noun}  \sense{
    \definition{type of fruit (~4 cm in diameter) that children shoot at birds}}}

\entry{kupkup}{\headword{kupkup}
  \pos{Noun}  \sense{
    \definition{small hole, pothole}}}

\entry{kupoll}{\headword{kupoll}
  \variantof{Unspecified Variant of \vartext{kupull}}  \pronnote{kupoɽ}}

\entry{kupull}{\headword{kupull}  \lexeme{kupull}  \pronnote{kupuɽ}
  \pos{Noun}  \sense{
    \definition{earthworm, worm}}
  \variants{\Unspecified Variant \vartext{kupoll}}}

\entry{Kur}{\headword{Kur}  \note{Place names}  \pronnote{kur}
  \pos{Proper noun}  \sense{
    \definition{Kur (Wipi-speaking village in Oriomo-Bituri Rural LLG; on the road to Oriomo)}    \note{Place names}    \example{
}    \example{
}    \example{
}}}

\entry{kuram}{\headword{kuram}  \note{Food}  \pronnote{kuram}
  \pos{Noun}  \sense{
    \definition{type of long yam with a white interior, hairs, and thorns}    \note{Food}    \example{
}    \example{
}    \example{
}}}

\entry{kurikuri}{\headword{kurikuri}  \note{Games}  \pronnote{kurikuri}
  \pos{Noun}  \sense{
    \definition{game involving spinning tree fruit on one's hand}    \note{Games}    \example{
}    \example{
}    \example{
}}}

\entry{Kurinti}{\headword{Kurinti}
  \variantof{Unspecified Variant of \vartext{Kurunti}}  \note{Place names}  \pronnote{kurinti}}

\entry{Kurintti}{\headword{Kurintti}
  \variantof{Unspecified Variant of \vartext{Kurunti}}}

\entry{kurkur}{\headword{kurkur}  \note{Birds}  \pronnote{kurkur}
  \pos{Noun}  \sense{
    \definition{type of bird}    \note{Birds}    \example{
}    \example{
}    \example{
}}}

\entry{kurmirang}{\headword{kurmirang}  \note{Traditional Arrows}  \pronnote{kurmiraŋ}
  \pos{Noun}  \sense{
    \definition{type of spear}    \note{Traditional Arrows}    \example{
}    \example{
}    \example{
}}}

\entry{Kurunti}{\headword{Kurunti}
  \pos{Proper noun}  \sense{
    \definition{Kurunti (Em-speaking village in Oriomo-Bituri Rural LLG; from Limol, one must pass through Malam)}}
  \variants{\Unspecified Variant \vartext{Kurintti, Kullintti, Kulläntti, Kurinti}}}

\entry{Kurupel}{\headword{Kurupel}  \pronnote{kurupel}
  \pos{Proper noun}  \sense{
    \definition{male personal name}}}

\entry{kus}{\headword{kus}  \note{Birds}  \pronnote{kus}
  \pos{Noun}  \sense{
    \definition{type of bird}    \note{Birds}    \example{
}    \example{
}    \example{
}}}

\entry{kusi}{\headword{kusi}
  \pos{Postposition}  \sense{
}}

\entry{kut}{\headword{kut}  \pronnote{kut}
  \pos{Noun}  \sense{
    \definition{raincloud}    \example{
      \exsen{Kut a ddone bogon.}      \extran{They aren't rainclouds.}}}}

\entry{kutae}{\headword{kutae}  \lexeme{kutae}  \pronnote{kutae}
  \pos{Noun}  \sense{
    \definition{type of small yam}}}

\entry{kutkut}{\headword{kutkut}
  \pos{Noun}  \sense{
    \definition{name of clan}}}

\entry{Kutpi Käp}{\headword{Kutpi Käp}  \pronnote{kutpi kəp}
  \pos{Proper noun}  \sense{
    \definition{Kutpi Käp (toponym)}}}

\entry{kutt}{\headword{kutt}  \lexeme{kutt}  \pronnote{kuʈ͡ʂ}
  \pos{Noun}  \sense{\textbf{1.} 
    \definition{bone}    \example{
      \exsen{ddäg kutt}      \extran{backbone}}    \example{
      \exsen{Lepresi da kutt ngämnanang itrell dan.}      \extran{Leprosy is a disease that can affect the bones.}}}  \sense{\textbf{2.} 
    \definition{seed, core}    \example{
      \exsen{Yokon kutt di dibewän.}      \extran{Yokon plants the seeds.}}}  \sense{\textbf{3.} 
    \definition{vague hard thing}    \example{
      \exsen{ine kutt, pane kutt}      \extran{water bucket, pan}}}}

\entry{kutt gugu}{\headword{kutt gugu}  \note{Conflict/war}  \pronnote{kuʈ͡ʂ ɡuɡu}
  \pos{Noun}  \sense{
    \definition{type of peace restoration}    \note{Conflict/war}    \example{
}    \example{
}    \example{
}}}

\entry{kutt llo}{\headword{kutt llo}  \note{Trees}  \pronnote{kuʈ͡ʂ ɽo}
  \pos{Noun}  \sense{
    \definition{type of tree that grows in the bush with bark that is scraped and rubbed on sores}    \note{Trees}    \example{
}    \example{
}    \example{
}}}

\entry{kutt snameny}{\headword{kutt snameny}  \lexeme{snameny}  \pronnote{snameɲ}
  \pos{Transitive verb}  \sense{
    \definition{a funeral tradition, where you take the body to an isolated place and one person will stay with the body and a spirit will come to tell them how they died.}}}

\entry{kutt tataem}{\headword{kutt tataem}
  \pos{Intransitive compound verb}  \sense{
    \definition{to shiver}    \example{
      \exsen{Ikopse agaebne adawede ge ttoen a ddone kutt tataemang kallkällang abal täräp me bongesän.}      \extran{[You all] pray that this thing doesn't happen at a shiveringly cold time.}}}}

\entry{kuttakuttang}{\headword{kuttakuttang}  \note{Adjectives}  \pronnote{kuʈ͡ʂakuʈ͡ʂaŋ}
  \pos{Modifier}  \sense{
    \definition{thin, bony}    \note{Adjectives}    \example{
}    \example{
}    \example{
}}}

\entry{kuyu}{\headword{kuyu}  \note{Trees}  \pronnote{kuju}
  \pos{Noun}  \sense{
    \definition{type of tree that grows in the bush or along creeks with soft wood that rots easily}    \note{Trees}    \example{
}    \example{
}    \example{
}}}

\entry{Kuyu}{\headword{Kuyu}  \note{Places around Limol}  \pronnote{kuju}
  \pos{Proper noun}  \sense{
    \definition{garden place of Matthew Bodog and Kaoga Dobola in Limol}    \note{Places around Limol}    \example{
}    \example{
}    \example{
}}}

\entry{kuu}{\headword{kuu}
  \variantof{Unspecified Variant of \vartext{ku}}}

\entry{kwa}{\headword{kwa}  \lexeme{kwa}  \pronnote{kwa}
  \pos{Ideophone}  \sense{
    \definition{sound made by a dog in pain}}}

\entry{kwae}{\headword{kwae}  \pronnote{kwae}
  \pos{Noun}  \sense{
    \definition{type of yam with a white or purple interior}}}

\entry{kwaena kutt}{\headword{kwaena kutt}  \pronnote{kwaena kuʈ͡ʂ}
  \pos{Noun}  \sense{
    \definition{hip bone}    \example{
}    \example{
}    \example{
}}}

\entry{kwakall}{\headword{kwakall}  \note{Trees}  \pronnote{kwakaɽ}
  \pos{Noun}  \sense{
    \definition{type of tree}    \note{Trees}    \example{
}    \example{
}    \example{
}}}

\entry{kwakasru}{\headword{kwakasru}  \note{Snakes}  \pronnote{kwakasru}
  \pos{Noun}  \sense{
    \definition{type of snake}    \note{Snakes}    \example{
}    \example{
}    \example{
}}}

\entry{Kwakmae}{\headword{Kwakmae}  \pronnote{kwakmae}
  \pos{Proper noun}  \sense{
    \definition{female personal name}}}

\entry{Kwalde}{\headword{Kwalde}  \pronnote{kwalde}
  \pos{Proper noun}  \sense{
    \definition{male personal name}}}

\entry{Kwale}{\headword{Kwale}  \pronnote{kwale}
  \pos{Proper noun}  \sense{
    \definition{female personal name}}}

\entry{kwallang}{\headword{kwallang}  \note{Trees}  \pronnote{kwaɽaŋ}
  \pos{Noun}  \sense{
    \definition{type of bush used for posts}    \note{Trees}    \example{
}    \example{
}    \example{
}}}

\entry{Kwallangkäbäll}{\headword{Kwallangkäbäll}  \note{Places around Limol}  \pronnote{kwaɽaŋkəbəɽ}
  \pos{Proper noun}  \sense{
    \definition{community garden place in Limol}    \note{Places around Limol}    \example{
}    \example{
}    \example{
}}}

\entry{Kwangka}{\headword{Kwangka}
  \pos{}  \sense{
    \definition{Kwangka (toponym)}}}

\entry{kwangka}{\headword{kwangka}  \pronnote{kwaŋka}
  \pos{Noun}  \sense{
    \definition{Torresian crow}    \scientificname{Corvus orru}    \example{
}    \example{
}    \example{
}}}

\entry{kwangka lläkäm}{\headword{kwangka lläkäm}  \note{Mushrooms}  \pronnote{kwaŋka ɽəkəm}
  \pos{Noun}  \sense{
    \definition{type of inedible mushroom}    \note{Mushrooms}    \example{
}    \example{
}    \example{
}}}

\entry{Kwangkangatt}{\headword{Kwangkangatt}
  \pos{Proper noun}  \sense{
    \definition{sacred place of Dobola (on the road to Malam)}}}

\entry{kwantta}{\headword{kwantta}  \note{Trees}  \pronnote{kwanʈ͡ʂa}
  \pos{Noun}  \sense{
    \definition{type of tree that grows in the grassland (~9 m) with green and yellow leaves used as bow pigment and hard fruit used to play a hockey-like game; also used for posts; liquid is extracted from the bark and given to dogs when they show signs of sickness}    \note{Trees}    \example{
}    \example{
}    \example{
}}}

\entry{kwäp}{\headword{kwäp}
  \variantof{Unspecified Variant of \vartext{kup}}}

\entry{Kwara}{\headword{Kwara}  \pronnote{kwara}
  \pos{Proper noun}  \sense{
    \definition{female personal name}}}

\entry{kwarakwara}{\headword{kwarakwara}  \pronnote{kwarakwara}
  \pos{Noun}  \sense{
    \definition{eastern hooded pitta}    \scientificname{Pitta novaeguineae}    \example{
}    \example{
}    \example{
}}}

\entry{kwaratang}{\headword{kwaratang}  \note{Adjectives}  \pronnote{kwarataŋ}
  \pos{Modifier}  \sense{
    \definition{thin, skinny, slim (animate)}    \note{Adjectives}    \example{
      \exsen{Ge däräng a kwaratang abal dan.}      \extran{This dog is very thin.}}    \example{
}    \example{
}}}

\entry{kwas}{\headword{kwas}  \lexeme{kwas}  \pronnote{kwas}
  \pos{Noun}  \sense{
    \definition{type of taro kongkong}}}

\entry{kwata}{\headword{kwata}  \note{Trees}  \pronnote{kwata}
  \pos{Noun}  \sense{
    \definition{type of tree that grows in the bush; used for house sticks}    \note{Trees}    \example{
}    \example{
}    \example{
}}}

\entry{kwätäs}{\headword{kwätäs}  \note{Trees}  \pronnote{kwətəs}
  \pos{Noun}  \sense{
    \definition{type of tree}    \note{Trees}    \example{
}    \example{
}    \example{
}}}

\entry{kwatt}{\headword{kwatt}
  \pos{Noun}  \sense{
    \definition{court}    \example{
      \exsen{Angde ubi bibim kwatt e yatraemneyo, mudag näkäp ttomoenen a.}      \extran{When they take you all to court, do not worry.}}    \example{
      \exsen{Yesu bom dämllaeyo a diba imnealle Paelet pate dätrameyo kwatt e.}      \extran{They tied up Jesus and after that, they took him to Pilate in court.}}}
  \variants{\Unspecified Variant \vartext{kot}}
  \variants{\SpellingVariant \vartext{kowatt}}}

\entry{Kwe}{\headword{Kwe}  \pronnote{kwe}
  \pos{Proper noun}  \sense{
    \definition{male personal name}}}

\entry{kwib}{\headword{kwib}  \note{Source: Jack Dipa}
  \pos{Noun}  \sense{
    \definition{charcoal made from a particular tree called upiye, used for painting during dance and initiation ceremony}    \note{Source: Jack Dipa}}}

\entry{Kwinton}{\headword{Kwinton}
  \variantof{SpellingVariant of \vartext{Quinton}}}

\chapter{L}


\entry{labalaba}{\headword{labalaba}
  \pos{Noun}  \sense{
    \definition{lap-lap}    \example{
      \exsen{Yäbäd angde ttänttämang abal gogon, däbe lla da obo labalaba de dängkänän ttänttäm atta.}      \extran{When the sun got very hot, that man took off his lap-lap because of the heat.}}}}

\entry{labelabet}{\headword{labelabet}  \note{Games}  \pronnote{labelabet}
  \pos{Noun}  \sense{
    \definition{type of game involving planting a stick in the middle of a circle}    \note{Games}    \example{
}    \example{
}    \example{
}}}

\entry{laem}{\headword{laem}  \pronnote{laem}
  \pos{Transitive verb}  \sense{\textbf{1.} 
    \definition{to roll, wrap}    \example{
      \exsen{laemnenatt otät}      \extran{rolled food, wrap}}    \example{
      \exsen{Bogo ttäkäll de dälemnegän.}      \extran{She rolled up balls of yam and coconut.}}}  \sense{\textbf{2.} 
    \definition{to argue}    \example{
}}}

\entry{laem~}{\headword{laem~}
  \variantof{freevarlab of \vartext{RED~}}}

\entry{laen}{\headword{laen}
  \pos{Noun}  \sense{
    \definition{line}}}

\entry{läläm}{\headword{läläm}  \pronnote{lələm}
  \pos{Noun}  \sense{
    \definition{muddy spot}}}

\entry{Lama}{\headword{Lama}  \pronnote{lama}
  \pos{Proper noun}  \sense{
    \definition{female personal name}}}

\entry{Lamlam}{\headword{Lamlam}  \pronnote{lamlam}
  \pos{Proper noun}  \sense{
    \definition{name of a female ancestor (sister of Moli)}}}

\entry{land}{\headword{land}
  \variantof{Dialectal Variant of \vartext{English loanword}}}

\entry{language}{\headword{language}
  \variantof{Dialectal Variant of \vartext{English loanword}}}

\entry{läpon}{\headword{läpon}  \pronnote{ləpon}
  \pos{Intransitive verb}  \sense{\textbf{1.} 
    \definition{to be amazed, be in awe}    \example{
      \exsen{Ubira tämamae lla yaba läponang gogon.}      \extran{All of them were amazed.}}    \example{
      \exsen{Ubi daläponeyo.}      \extran{They were amazed.}}}  \sense{\textbf{2.} 
    \definition{to doubt}}}

\entry{last}{\headword{last}
  \variantof{Dialectal Variant of \vartext{English loanword}}}

\entry{Lauren}{\headword{Lauren}  \pronnote{lauren}
  \pos{Proper noun}  \sense{
    \definition{female personal name}}}

\entry{Lei}{\headword{Lei}
  \pos{Proper noun}  \sense{
    \definition{Lei (toponym)}}}

\entry{Lek Märi}{\headword{Lek Märi}
  \pos{Proper noun}  \sense{
    \definition{Lake Murray (in Lake Murray Rural LLG)}}}

\entry{lel}{\headword{lel}  \note{Verbs}  \pronnote{lel}
  \pos{Noun}  \sense{\textbf{1.} 
    \definition{fear}    \note{Verbs}    \example{
      \exsen{Ubi dindugeyo lel me.}      \extran{They ran in fear.}}}  \sense{\textbf{2.} 
    \definition{to fear, be afraid of}    \note{Verbs}    \example{
      \exsen{Lla da sɨmell de lel dägagän.}      \extran{The man feared the pig.}}    \example{
      \exsen{Ge kungge da bäne lel ma za dan.}      \extran{This spider is what you fear.}}    \example{
}}  \sense{\textbf{3.} 
    \definition{shame}    \note{Verbs}}  \sense{\textbf{4.} 
    \definition{to be ashamed (of)}    \note{Verbs}    \example{
      \exsen{Bogo lel allan ngämo mit me.}      \extran{She is ashamed because of me.}}    \example{
      \exsen{Baba da mɨnyi obom lel bägagän.}      \extran{Father will be ashamed of him.}}}}

\entry{lelang}{\headword{lelang}  \pronnote{lelaŋ}
  \pos{Modifier}  \sense{\textbf{1.} 
    \definition{afraid, scared}    \example{
      \exsen{Bogo deya lelang llo toko e gogäbänän.}      \extran{He was so scared he jumped into the treetop.}}    \example{
}}  \sense{\textbf{2.} 
    \definition{shameful}    \example{
      \exsen{Däbe mɨnyi ddobae lelang abal ttoen da bogalle obo kuki tab e.}      \extran{If he broke the promise, that would be a very shameful thing.}}}}

\entry{lelmeny}{\headword{lelmeny}
  \pos{Modifier}  \sense{
    \definition{brave, fearless, bold}}}

\entry{lepade}{\headword{lepade}  \note{Trees}  \pronnote{lepade}
  \pos{Noun}  \sense{
    \definition{type of cultivated tree with purple and white flowers and edible black fruit that kids like to eat}    \note{Trees}    \example{
}    \example{
}    \example{
}}}

\entry{lepresi}{\headword{lepresi}
  \pos{Noun}  \sense{
    \definition{leprosy}}}

\entry{lesna}{\headword{lesna}  \note{Trees}  \pronnote{lesna}
  \pos{Noun}  \sense{
    \definition{type of tree}    \note{Trees}    \example{
}    \example{
}    \example{
}}}

\entry{Letai}{\headword{Letai}  \pronnote{letai}
  \pos{Proper noun}  \sense{
    \definition{female personal name}}}

\entry{Lewada}{\headword{Lewada}  \pronnote{lewada}
  \pos{Proper noun}  \sense{
    \definition{Lewada (Makayam-speaking village in Gogodala Rural LLG, on the Fly River; GPS: 8.327787, 142.785487)}    \example{
}    \example{
}    \example{
}}}

\entry{lid}{\headword{lid}
  \pos{Noun}  \sense{
    \definition{lid}}}

\entry{lida}{\headword{lida}
  \pos{Noun}  \sense{
    \definition{leader}}}

\entry{Lidiya}{\headword{Lidiya}
  \pos{Proper noun}  \sense{
    \definition{female personal name}}}

\entry{liglig}{\headword{liglig}  \note{Verbs of Motion}  \pronnote{liɡliɡ}
  \pos{Intransitive verb}  \sense{
    \definition{to have sex, copulate}    \note{Verbs of Motion}    \example{
      \exsen{Bogo liglig eran.}      \extran{He is having sex.}}    \example{
      \exsen{Ubi do me alignallo.}      \extran{They two had sex there.}}    \example{
      \exsen{Ubi iddob me ddone bulignegnän.}      \extran{They will not be having sex at night.}}}}

\entry{liglig suang}{\headword{liglig suang}  \pronnote{liɡliɡ suaŋ}
  \pos{Modifier}  \sense{
    \definition{having a lot of sex}}}

\entry{liglig ttoen}{\headword{liglig ttoen}  \pronnote{liɡliɡ ʈ͡ʂoen}
  \pos{Noun}  \sense{
    \definition{sexual relations}    \example{
      \exsen{Bäne gullbe o mälla walle mae liglig ttoen de nangesnegneyo.}      \extran{Only have sexual relations with your husband or wife.}}    \example{
}    \example{
}}}

\entry{Lili}{\headword{Lili}
  \pos{Proper noun}  \sense{
    \definition{female personal name}}
  \variants{\SpellingVariant \vartext{Lily}}}

\entry{Lilian}{\headword{Lilian}  \pronnote{lilian}
  \pos{Proper noun}  \sense{
    \definition{female personal name}}}

\entry{Lily}{\headword{Lily}
  \variantof{SpellingVariant of \vartext{Lili}}  \pronnote{lili}}

\entry{Limoll}{\headword{Limoll}  \pronnote{limoɽ}
  \pos{Proper noun}  \sense{
    \definition{Limol (Ende-speaking village in Morehead Rural LLG; GPS: -8.641783, 142.682533)}}
  \variants{\Dialectal Variant \vartext{Llimoll}}}

\entry{Limollang}{\headword{Limollang}  \note{Dialect names}  \pronnote{limoɽaŋ}
  \pos{Noun}  \sense{\textbf{1.} 
    \definition{Limol villager, person from Limol}    \note{Dialect names}}  \sense{\textbf{2.} 
    \definition{Ende dialect spoken in Limol}    \note{Dialect names}    \example{
}    \example{
}    \example{
}}}

\entry{Linda}{\headword{Linda}  \pronnote{linda}
  \pos{Proper noun}  \sense{
    \definition{female personal name}}}

\entry{Linette}{\headword{Linette}  \pronnote{lineʈ͡ʂe}
  \pos{Proper noun}  \sense{
    \definition{female personal name}}
  \variants{\SpellingVariant \vartext{Lynet}}}

\entry{linge}{\headword{linge}  \note{Palms}  \pronnote{liŋe}
  \pos{Noun}  \sense{
    \definition{type of palm tree with blue flowers}    \note{Palms}    \example{
}    \example{
}    \example{
}}}

\entry{linguistic}{\headword{linguistic}
  \variantof{freevarlab of \vartext{English loanword}}}

\entry{Liseng}{\headword{Liseng}
  \pos{Proper noun}  \sense{
}}

\entry{lizom}{\headword{lizom}  \note{Mushrooms}  \pronnote{lizom}
  \pos{Noun}  \sense{
    \definition{type of mushroom}    \note{Mushrooms}    \example{
}    \example{
}    \example{
}}}

\entry{lɨklɨk}{\headword{lɨklɨk}
  \pos{Intransitive verb}  \sense{
    \definition{to melt}    \example{
      \exsen{Ngämo ine kutt a daeya yu da lɨklɨk eran.}      \extran{Fire is melting my water container.}}}}

\entry{ll}{\headword{ll}  \pronnote{ɽ}
  \pos{Auxiliary verb}  \sense{
    \definition{aux}}}

\entry{lla}{\headword{lla}  \lexeme{lla}  \pronnote{ɽa}
  \pos{Noun}  \sense{\textbf{1.} 
    \definition{man, male}    \anthronote{hand sign for man is to pull on chin as if on beard}    \example{
      \exsen{Lla pättäpättäk a Meri bom dangnoeyän llädäd e.}      \extran{The short man asked Mary to marry him.}}}  \sense{\textbf{2.} 
    \definition{husband}}  \sense{\textbf{3.} 
    \definition{person, human being}    \example{
      \exsen{Obo pallall lla da ddone ada kilikili gognegnän a tuk i gogäbnän.}      \extran{The people next to him were so happy and were jumping in the air.}}}  \sense{\textbf{4.} 
}}

\entry{lla bombllo}{\headword{lla bombllo}  \pronnote{ɽa bombɽo}
  \pos{Noun}  \sense{
    \definition{generation}    \example{
      \exsen{Ewatta ke gänya täräp me lla bombllo da wasnen anggan ttowamang ttoen tonton ikop e ddapall att?}      \extran{Why does the current generation beg to directly see miraculous things from heaven?}}}}

\entry{lla diben}{\headword{lla diben}  \note{Snakes}  \pronnote{ɽa diben}
  \pos{Noun}  \sense{
    \definition{type of snake}    \note{Snakes}    \example{
}    \example{
}    \example{
}}}

\entry{lla gugu}{\headword{lla gugu}  \note{Conflict/war}  \pronnote{ɽa ɡuɡu}
  \pos{Noun}  \sense{
    \definition{restoring peace after a murder by trading a young girl in the victim's place}    \note{Conflict/war}    \example{
}    \example{
}    \example{
}}}

\entry{lla gul}{\headword{lla gul}  \pronnote{ɽa ɡul}
  \pos{Noun}  \sense{
    \definition{crowd}    \example{
      \exsen{Ddobag a bi dingismällalle täräb ma we, lla gulag a erame dag.}      \extran{The others return to the funeral home, where the crowd is.}}}}

\entry{lla ikoikopang}{\headword{lla ikoikopang}  \note{Dancing}  \pronnote{ɽa ikoikopaŋ}
  \pos{Noun}  \sense{
    \definition{audience}    \note{Dancing}    \example{
}    \example{
}    \example{
}}}

\entry{lla llɨg}{\headword{lla llɨg}  \pronnote{ɽa ɽɪɡ}
  \pos{Noun}  \sense{
    \definition{boy}    \example{
      \exsen{Zon lla llɨg dan.}      \extran{John is a boy.}}}}

\entry{lla päpätt sämpallmeny alla}{\headword{lla päpätt sämpallmeny alla}  \note{Birth}  \pronnote{ɽa pəpəʈ͡ʂ səmpaɽmeɲ aɽa}
  \pos{}  \sense{
    \definition{false labor}    \note{Birth}    \example{
}    \example{
}    \example{
}}}

\entry{lla up}{\headword{lla up}  \note{Names of bananas}  \pronnote{ɽa up}
  \pos{Noun}  \sense{
    \definition{type of introduced banana}    \note{Names of bananas}    \example{
}    \example{
}    \example{
}}}

\entry{lla winy}{\headword{lla winy}  \note{Bees}  \pronnote{ɽa wiɲ}
  \pos{Noun}  \sense{
    \definition{type of biting bee that lives in trees}    \note{Bees}    \example{
}    \example{
}    \example{
}}}

\entry{lläb}{\headword{lläb}  \pronnote{ɽəb}
  \pos{Intransitive verb}  \sense{
    \definition{to go under}    \example{
      \exsen{Kungge da llo ttam dädär ik i golläbän.}      \extran{The spider went under the dry tree leaves.}}}}

\entry{lläbe}{\headword{lläbe}  \note{Parts of the Body}  \pronnote{ɽəbe}
  \pos{Noun}  \sense{
    \definition{nail, claw (of a finger or toe)}    \note{Parts of the Body}    \example{
}    \example{
}    \example{
}}}

\entry{lläblläb}{\headword{lläblläb}
  \pos{Transitive verb}  \sense{
    \definition{to hide}}}

\entry{llabun}{\headword{llabun}  \pronnote{ɽabun}
  \pos{Noun}  \sense{
    \definition{relative, kinsman, clansman}    \example{
      \exsen{Kemu bo llabun a gänyme a Malläm me dadeg.}      \extran{Kemu has relatives here and in Malam.}}}}

\entry{llädäd}{\headword{llädäd}  \pronnote{ɽədəd}
  \pos{Transitive verb}  \sense{\textbf{1.} 
    \definition{to grab, get, catch}    \example{
      \exsen{Ngäna walle we gäbän allan ddia llädäd e.}      \extran{I am jumping in the water to grab the deer.}}    \example{
      \exsen{Ubi guspullän, be ubi obaoba gollädnegän.}      \extran{They fell, but they caught themselves.}}    \example{
      \exsen{Iba pätt a ttongdae dan, ako kame ddone bällädeya.}      \extran{We (incl.) only have one body; we don't get another.}}}  \sense{\textbf{2.} 
    \definition{to buy}    \example{
      \exsen{Da ttongo lla da bäne baba bälle ten taosen mani bonttogän, ende bogo bällädän?}      \extran{If someone gave your dad ten thousand kina, what would he buy?}}    \example{
      \exsen{Tu andred alle dällädeyo ttoe de.}      \extran{He bought the skin for two hundred.}}}  \sense{\textbf{3.} 
    \definition{to marry}    \example{
      \exsen{E mälla de ke bongo dälläd?}      \extran{Which woman did you marry?}}    \example{
      \exsen{Da bongo lla tupi di nälläd, bongo mɨnyi llɨg kamebiag ag.}      \extran{If you marry the tall man, you will have many children.}}}  \sense{\textbf{4.} 
    \definition{to come, arrive (figuratively)}    \example{
      \exsen{Angde obo täräp a ngäs e dällädalle, bogo mɨnyi do alle olle gogalle oba pate.}      \extran{When his time to return came, he would call out to them.}}}}

\entry{llädae}{\headword{llädae}
  \pos{Transitive verb}  \sense{
    \definition{to roll}    \example{
      \exsen{Llɨg a bal de llädaenen eran.}      \extran{The boy is rolling the ball.}}    \example{
      \exsen{Ubi dädär ulle de dalldaeyo auma ud de däpakeyo.}      \extran{They rolled a large rock and blocked the grave entrance.}}}}

\entry{lläg}{\headword{lläg}
  \variantof{SpellingVariant of \vartext{llɨg}}  \pronnote{ɽəɡ}}

\entry{lläk~}{\headword{lläk~}
  \variantof{freevarlab of \vartext{RED~}}}

\entry{llakällakätt}{\headword{llakällakätt}  \note{Trees}  \pronnote{ɽakəɽakəʈ͡ʂ}
  \pos{Noun}  \sense{
    \definition{type of tree that grows in the bush; used for house sticks}    \note{Trees}    \example{
}    \example{
}    \example{
}}}

\entry{lläkäm}{\headword{lläkäm}  \note{Medicine}  \pronnote{ɽəkəm}
  \pos{Noun}  \sense{
    \definition{mushroom}    \note{Medicine}    \example{
      \exsen{Lläkäm a mamall zozoang dan.}      \extran{A mushroom goes rotten quickly.}}    \example{
}    \example{
}}
  \variants{\SpellingVariant \vartext{llɨkɨm}}}

\entry{llakek}{\headword{llakek}  \pronnote{ɽakek}
  \pos{Intransitive verb}  \sense{
    \definition{to collapse}    \example{
      \exsen{Särem a ollakan.}      \extran{Darkness collapses.}}}}

\entry{lläklläk}{\headword{lläklläk}
  \pos{Transitive verb}  \sense{
    \definition{to spread fire}    \example{
      \exsen{Sali yu di mermerae lläklläk eran.}      \extran{Sali is spreading the fire nicely.}}}}

\entry{llakllek}{\headword{llakllek}  \pronnote{ɽakɽek}
  \pos{Transitive verb}  \sense{
    \definition{to destroy}    \example{
      \exsen{Tätäm ttäle bun parga de bob a era dällekän.}      \extran{Yesterday the flood destroyed the Ttäle Bun bridge.}}}}

\entry{llällam}{\headword{llällam}  \pronnote{llä'llam}
  \pos{Ideophone}  \sense{
    \definition{rustling (of plants); roaring (of water)}    \example{
      \exsen{Gullem a sisri ag me towall llällallällamang yaran.}      \extran{This morning the snake came rustling the grass.}}    \example{
      \exsen{Däbaballe ine llällam de dandärän.}      \extran{From there, he heard the water roaring.}}}}

\entry{llällamang}{\headword{llällamang}  \pronnote{ɽəɽamaŋ}
  \pos{Modifier}  \sense{
    \definition{noisy}}}

\entry{llälläp}{\headword{llälläp}  \lexeme{llälläp}  \pronnote{ɽəɽəp}
  \pos{Noun}  \sense{
    \definition{type of small snake that catches frogs}}}

\entry{llallem}{\headword{llallem}  \pronnote{ɽaɽem}
  \pos{Transitive verb}  \sense{
    \definition{to finish weaving}}}

\entry{llama}{\headword{llama}
  \pos{Modifier}  \sense{\textbf{1.} 
    \definition{unsatisfied}    \example{
      \exsen{Kwalde säre llama näkäp me gongosmällnän.}      \extran{Kwalde sadly returned unsatisfied.}}}  \sense{\textbf{2.} 
    \definition{hesitant, reluctant}    \example{
      \exsen{Däräng a llama näkäp me gongezänmällnän mama watta.}      \extran{The dog was coming out of the grass pile hesitantly.}}}}

\entry{llama}{\headword{llama}  \pronnote{ɽama}
  \pos{Modifier}  \sense{
    \definition{other people's, of others}    \example{
      \exsen{Ende eka da ngäma eka dan, Ingglis a llama eka dan.}      \extran{Ende is our (excl.) language; English is the language of other people.}}    \example{
      \exsen{Llama ttoen da ge, be ngämi zäme ngämira dägageya ge ttoen de.}      \extran{This is another people's custom, but we are already doing this custom.}}}}

\entry{llamäg}{\headword{llamäg}  \pronnote{ɽaməɡ}
  \pos{Noun}  \sense{
    \definition{old man}    \example{
      \exsen{llamäg a auseause}      \extran{old men and women}}}}

\entry{llamda}{\headword{llamda}  \note{Family ties}  \pronnote{ɽamda}
  \pos{Noun}  \sense{
    \definition{old man}    \note{Family ties}    \example{
      \exsen{Llamda da kumiye agan.}      \extran{The old man coughed.}}}}

\entry{llame}{\headword{llame}  \pronnote{ɽame}
  \pos{Manner adverb}  \sense{
    \definition{together}    \example{
      \exsen{Ngämi llame deyareya mamoema.}      \extran{We (excl.) went hunting together.}}    \example{
      \exsen{Ngämi llameae gobällne.}      \extran{We (excl.) were going together.}}    \example{
      \exsen{Oba za sinen a llamealle mae.}      \extran{They shared their things together.}}}}

\entry{llan}{\headword{llan}  \pronnote{ɽan}
  \pos{Noun}  \sense{
    \definition{ear}    \example{
      \exsen{Obo llan a ulle abal dan.}      \extran{His ear is very big.}}}}

\entry{llan gonggo}{\headword{llan gonggo}  \pronnote{ɽan ɡoŋɡo}
  \pos{Transitive compound verb}  \sense{
    \definition{to turn one's ear, listen intently}    \example{
      \exsen{Ubi tätäm kaemne de llan donggoeyo.}      \extran{Yesterday, they listened intently to the bee.}}}}

\entry{llan ik}{\headword{llan ik}  \note{Parts of the Body}  \pronnote{ɽan ik}
  \pos{Noun}  \sense{
    \definition{eardrum}    \note{Parts of the Body}    \example{
}    \example{
}    \example{
}}}

\entry{llan källa}{\headword{llan källa}
  \pos{Noun}  \sense{
    \definition{earwax}}}

\entry{llan kollkoll}{\headword{llan kollkoll}  \note{Parts of a fish}  \pronnote{ɽan koɽkoɽ}
  \pos{Noun}  \sense{
    \definition{gill}    \note{Parts of a fish}    \example{
}    \example{
}    \example{
}}}

\entry{llan kom}{\headword{llan kom}  \note{Parts of the Body}  \pronnote{ɽan kom}
  \pos{Noun}  \sense{
    \definition{ear hair}    \note{Parts of the Body}    \example{
}    \example{
}    \example{
}}}

\entry{llan mit}{\headword{llan mit}  \pronnote{ɽan mit}
  \pos{Noun}  \sense{
    \definition{operculum (gill cover)}    \example{
}    \example{
}    \example{
}}}

\entry{llanded}{\headword{llanded}  \pronnote{ɽanded}
  \pos{Transitive verb}  \sense{\textbf{1.} 
    \definition{to clear}    \example{
      \exsen{Oblle nyongo de dallendneyo.}      \extran{They were clearing the road for her.}}}  \sense{\textbf{2.} 
    \definition{to clarify, make clear, explain}    \example{
      \exsen{Bogo eka midd de dallendän.}      \extran{He clarified the meaning.}}}}

\entry{llandräg}{\headword{llandräg}  \pronnote{ɽandrəɡ}
  \pos{Intransitive coverb}  \sense{
    \definition{to listen}    \example{
      \exsen{Bongo mɨnyi llandräg ag?}      \extran{Are you going to listen?}}    \example{
      \exsen{Abo llandräg ag Malläm att eka we.}      \extran{Expect some news from Malam.}}    \example{
      \exsen{Daeya do llandär gognän.}      \extran{She was there listening.}}}}

\entry{lläng}{\headword{lläng}  \note{Adjectives}  \pronnote{ɽəŋ}
  \pos{Modifier}  \sense{
    \definition{sharp}    \note{Adjectives}    \example{
      \exsen{Ttam a ddobae lläng dag.}      \extran{The leaves are very sharp.}}    \example{
}    \example{
}}
  \variants{\Unspecified Variant \vartext{llɨng}}}

\entry{llängmeny}{\headword{llängmeny}  \note{Adjectives}  \pronnote{ɽəŋmeɲ}
  \pos{Modifier}  \sense{
    \definition{dull, blunt}    \note{Adjectives}    \example{
      \exsen{Gudne giri da llängmeny dan.}      \extran{The old knife is dull.}}    \example{
}    \example{
}}}

\entry{lläntäg}{\headword{lläntäg}  \pronnote{ɽəntəɡ}
  \pos{Transitive verb}  \sense{
    \definition{to tell}    \example{
      \exsen{Golläntmenyän tämamae ngalen de.}      \extran{He told him everything.}}}}

\entry{lläpän}{\headword{lläpän}  \pronnote{ɽəpən}
  \pos{Transitive verb}  \sense{
    \definition{to dig, harvest, unearth (a tuber or corm)}    \example{
      \exsen{Mälla da dagaeya ttängäm e abällan mätta lläplläp e.}      \extran{The women went to the garden to dig the yams.}}    \example{
      \exsen{Bogo biye kemibi dälläpnegalle.}      \extran{He used to harvest many taro roots.}}}}

\entry{lläpät}{\headword{lläpät}  \note{Market}  \pronnote{ɽəpət}
  \pos{Noun}  \sense{
    \definition{digit}    \note{Market}    \example{
      \exsen{ttang lläpät, nying lläpät}      \extran{finger, toe}}}
  \variants{\Dialectal Variant \vartext{llɨpɨt}}}

\entry{llapu}{\headword{llapu}  \note{Trees}  \pronnote{ɽapu}
  \pos{Noun}  \sense{
    \definition{type of big tree that grows in the bush with red fruit and white flowers}    \note{Trees}    \example{
}    \example{
}    \example{
}}}

\entry{llapuyurwe}{\headword{llapuyurwe}  \note{Trees}  \pronnote{ɽapujurwe}
  \pos{Noun}  \sense{
    \definition{type of tree with seedless, edible red fruit}    \note{Trees}    \example{
}    \example{
}    \example{
}}}

\entry{llät}{\headword{llät}
  \variantof{Fast Speech Variant of \vartext{llɨtɨt}}  \pronnote{ɽət}}

\entry{llätät}{\headword{llätät}
  \variantof{Dialectal Variant of \vartext{llɨtɨt}}  \pronnote{ɽətət}}

\entry{llätät}{\headword{llätät}  \pronnote{ɽətət}
  \pos{Transitive verb}  \sense{
    \definition{to get rid of oven stones}    \example{
      \exsen{Llɨg kälakäle, ttägäll käp de nälläteyoǃ}      \extran{Small children, get rid of these mumu stones!}}}}

\entry{llatat}{\headword{llatat}  \pronnote{ɽatat}
  \pos{Transitive verb}  \sense{
    \definition{to twist, wrap}    \example{
      \exsen{Gull a kab llatnenatt alle iatt kollba gäddnan ma dan.}      \extran{A net is a thing for catching fish, woven from twisted string.}}}
  \variants{\Unspecified Variant \vartext{llatet}}}

\entry{llatat}{\headword{llatat}  \pronnote{ɽatat}
  \pos{Transitive verb}  \sense{
    \definition{to trace, track, follow the blood (of an animal to kill it)}    \example{
      \exsen{Gongosän ma we, däräng a wa kumuddäga lla da dokomän, ddia llatat e.}      \extran{He returned home, and took his dog and three men to track the deer.}}}
  \variants{\SpellingVariant \vartext{llätät}}}

\entry{llatata}{\headword{llatata}
  \pos{Noun}  \sense{
    \definition{Food (such as sago, ripe bananas, and coconut cream, or yams and coconut cream) wrapped in a woven cococnut leaf (with a banana leaf within it) and cooked in a mumu}}}

\entry{llatet}{\headword{llatet}
  \variantof{Unspecified Variant of \vartext{llatat}}}

\entry{llätmäll}{\headword{llätmäll}
  \pos{Transitive verb}  \sense{
    \definition{to choose, select}    \example{
      \exsen{Bogo tuwelb lla de dällätmällnegän.}      \extran{He chose twelve men.}}}}

\entry{llätt}{\headword{llätt}  \pronnote{ɽəʈ͡ʂ}
  \pos{Noun}  \sense{\textbf{1.} 
    \definition{end}    \example{
      \exsen{ekaklle llätt e}      \extran{to the ends of the earth}}    \example{
      \exsen{pazi llätt me}      \extran{at the end of the year}}    \example{
      \exsen{Ge ttoen a llätt e gongttägän.}      \extran{This story is coming to an end.}}}  \sense{\textbf{2.} 
    \definition{to stop, end, finish}    \example{
      \exsen{Oba abo dindu da dädme llätt gogon.}      \extran{Their race stopped there.}}    \example{
      \exsen{Oba nag ttoen a llätt gogon.}      \extran{Their friendship ended.}}    \example{
      \exsen{Llɨtt agmom!}      \extran{(You all) stop!}}}
  \variants{\SpellingVariant \vartext{llɨtt}}}

\entry{llättmeny}{\headword{llättmeny}  \pronnote{ɽəʈ͡ʂmeɲ}
  \pos{Adverb}  \sense{
    \definition{nonstop}    \example{
      \exsen{Ngäna dinduangmeae gogne llättmenyae.}      \extran{I was running nonstop.}}    \example{
}    \example{
}}}

\entry{llimba}{\headword{llimba}  \pronnote{ɽimba}
  \pos{Noun}  \sense{
    \definition{fingernail}    \example{
      \exsen{Ngämo llimba da agällbänan.}      \extran{My fingernail broke.}}}}

\entry{Llimoll}{\headword{Llimoll}
  \variantof{Dialectal Variant of \vartext{Limoll}}}

\entry{llɨg}{\headword{llɨg}  \lexeme{llɨg}  \note{Nature}  \pronnote{ɽɪɡ}
  \pos{Noun}  \sense{\textbf{1.} 
    \definition{boy}    \note{Nature}    \example{
      \exsen{llɨg mänmän}      \extran{boys and girls}}    \example{
}    \example{
}}  \sense{\textbf{2.} 
    \definition{son}    \note{Nature}    \example{
      \exsen{Godd bo Llɨg indrang da.}      \extran{The Son of God is light.}}}  \sense{\textbf{3.} 
    \definition{child}    \note{Nature}    \example{
      \exsen{Bäne llɨg bo bin a ainin?}      \extran{What is your child's name?}}}
  \variants{\SpellingVariant \vartext{lläg}}}

\entry{llɨg dum}{\headword{llɨg dum}  \note{Birth}  \pronnote{ɽɪɡ dum}
  \pos{Noun}  \sense{\textbf{1.} 
    \definition{womb, uterus}    \note{Birth}}  \sense{\textbf{2.} 
    \definition{placenta}    \note{Birth}    \example{
}    \example{
}    \example{
}}}

\entry{llɨg ikopang}{\headword{llɨg ikopang}  \pronnote{ɽɪɡ ikopaŋ}
  \pos{Noun}  \sense{
    \definition{boyfriend}    \example{
      \exsen{Malla ngämo llɨg ikopang da,}}}}

\entry{llɨg mapät}{\headword{llɨg mapät}  \pronnote{ɽɪɡ mapət}
  \pos{Noun}  \sense{
    \definition{infant}    \example{
      \exsen{Ddäma da llɨg mapät komnen ma dan.}      \extran{A basket is for carrying infants.}}}}

\entry{llɨgllɨg}{\headword{llɨgllɨg}  \pronnote{ɽɪɡɽɪɡ}
  \pos{Noun}  \sense{
    \definition{nonsingular form of llɨg}    \example{
      \exsen{Mäkat llɨgllɨgag a dan bäne gaguma me.}      \extran{There is a rat with babies in your yamhouse.}}}}

\entry{llɨgmeny}{\headword{llɨgmeny}  \note{Birth}  \pronnote{ɽɪɡmeɲ}
  \pos{Modifier}  \sense{
    \definition{childless}    \note{Birth}    \example{
      \exsen{Ngäna llɨgmeny dan.}      \extran{I don't have children. / I'm childless.}}    \example{
}    \example{
}}}

\entry{llɨgmeny agallo}{\headword{llɨgmeny agallo}  \note{Birth}  \pronnote{ɽɪɡmeɲ aɡaɽo}
  \pos{}  \sense{
    \definition{to plan for not having children}    \note{Birth}    \example{
}    \example{
}    \example{
}}}

\entry{llɨgmeny e wätäba}{\headword{llɨgmeny e wätäba}  \note{Birth}  \pronnote{ɽɪɡmeɲ e wətəba}
  \pos{}  \sense{
    \definition{to plan for not having children}    \note{Birth}    \example{
}    \example{
}    \example{
}}}

\entry{llɨkɨm}{\headword{llɨkɨm}
  \variantof{SpellingVariant of \vartext{lläkäm}}}

\entry{llɨng}{\headword{llɨng}
  \variantof{Unspecified Variant of \vartext{lläng}}}

\entry{llɨpɨt}{\headword{llɨpɨt}
  \variantof{Dialectal Variant of \vartext{lläpät}}  \pronnote{ɽɪpɪt}}

\entry{llɨtɨn}{\headword{llɨtɨn}
  \pos{Transitive verb}  \sense{
    \definition{to select}}}

\entry{llɨtɨt}{\headword{llɨtɨt}  \note{?}  \pronnote{ɽɪtɪt}
  \pos{Transitive verb}  \sense{\textbf{1.} 
    \definition{to tell, report, say}    \note{?}    \example{
      \exsen{Bogo eka de llɨtɨtang meae nägnan.}      \extran{He told the same story repeatedly.}}    \example{
      \exsen{Män a llɨg bälle tongoeang ttoen di dɨllɨtän.}      \extran{The girl told the boy a funny story.}}    \example{
      \exsen{Ubi ddone amlle ende golläntmenynegän adawatta ddobaeddobae lel gognegnän.}      \extran{They didn't say anything to anyone because they were very afraid.}}}  \sense{\textbf{2.} 
    \definition{to sing}    \note{?}    \example{
      \exsen{Wagiba bandra de nɨllɨtnan sisri ag me.}      \extran{Wagiba was singing this morning.}}    \example{
}}
  \variants{\Dialectal Variant \vartext{llätät}}
  \variants{\Fast Speech Variant \vartext{llät}}}

\entry{llɨtɨt}{\headword{llɨtɨt}  \pronnote{ɽɪtɪt}
  \pos{Transitive verb}  \sense{
    \definition{to butcher, cut}    \example{
      \exsen{Ddia ddäg de nällɨt.}      \extran{[You] cut the deer back.}}}
  \variants{\Dialectal Variant \vartext{llätät}}}

\entry{llɨtt}{\headword{llɨtt}
  \variantof{SpellingVariant of \vartext{llätt}}}

\entry{llɨtt}{\headword{llɨtt}  \pronnote{ɽɪʈ͡ʂ}
  \pos{Intransitive verb}  \sense{
    \definition{to get worse, worsen}    \example{
      \exsen{Lla da angde itrel a dägagalle, itrel atta bollɨttnän bollɨttnän angde ako kuddäll a damllamalle, kuddäll gogalle.}      \extran{When a man gets sick, and he gets worse and worse, then death catches him and he dies.}}}}

\entry{llo}{\headword{llo}  \lexeme{llo}  \pronnote{ɽo}
  \pos{Noun}  \sense{\textbf{1.} 
    \definition{tree}    \example{
      \exsen{Ngäna llo de plengg eran.}      \extran{I am chopping down the tree.}}    \example{
      \exsen{Tätäm Matthew llo de dängkälän.}      \extran{Yesterday, Matthew climbed the tree.}}}  \sense{\textbf{2.} 
    \definition{wood}    \example{
      \exsen{llo pallkoll tubu}      \extran{short piece of wood}}    \example{
      \exsen{Melemang dag ubi, llo ma abällaban.}      \extran{They are working; they went for wood.}}}  \sense{\textbf{3.} 
    \definition{stick}    \example{
      \exsen{Lla da llo alle ekaklle me dadräbän.}      \extran{The man wrote on the ground with a stick.}}}}

\entry{llo batt}{\headword{llo batt}
  \pos{Noun}  \sense{
    \definition{wooden stick}    \example{
      \exsen{Ngäna ttongo llo batt e dängällbän.}      \extran{I got one stick.}}}}

\entry{llo ddage}{\headword{llo ddage}  \lexeme{llo ddage}  \pronnote{ɽo ɖ͡ʐaɡe}
  \pos{Noun}  \sense{
    \definition{tree branch}    \example{
      \exsen{Ngäna llo ddage me ingongang dan.}      \extran{I am dancing on the tree branch.}}}}

\entry{llo ngoeang}{\headword{llo ngoeang}  \lexeme{llo ngoeang}  \pronnote{ɽo ŋojaŋ}
  \pos{Noun}  \sense{
    \definition{forked tree branch}}}

\entry{llo patipati}{\headword{llo patipati}  \lexeme{llo patipati}  \pronnote{ɽo patipati}
  \pos{Noun}  \sense{
    \definition{tree goanna}}}

\entry{llo patt}{\headword{llo patt}  \lexeme{llo patt}  \pronnote{ɽo paʈ͡ʂ}
  \pos{Noun}  \sense{
    \definition{fallen tree}    \example{
      \exsen{Ngäna llo patt toko me godmen.}      \extran{I sat on the fallen tree.}}}}

\entry{llo pätt}{\headword{llo pätt}
  \pos{Noun}  \sense{
    \definition{tree trunk}}}

\entry{llo ttam}{\headword{llo ttam}  \lexeme{llo ttam}  \pronnote{ɽo ʈ͡ʂam}
  \pos{Noun}  \sense{
    \definition{tree leaf}    \example{
      \exsen{Molemoleg a llo ttam otnanang dag.}      \extran{Caterpillars eat tree leaves.}}}}

\entry{llo ttoe}{\headword{llo ttoe}
  \pos{Noun}  \sense{
    \definition{tree bark}}}

\entry{llo tubu}{\headword{llo tubu}
  \pos{Noun}  \sense{
    \definition{tree stump}}}

\entry{llokollokott}{\headword{llokollokott}  \pronnote{ɽokoɽokoʈ͡ʂ}
  \pos{Adverb}  \sense{
    \definition{firmly, tightly, strongly}    \example{
      \exsen{Ngäna sɨmell de llokollokott dinam.}      \extran{I pressed the pig down tightly.}}}}

\entry{llokollokott}{\headword{llokollokott}  \pronnote{ɽokoɽokoʈ͡ʂ}
  \pos{Modifier}  \sense{\textbf{1.} 
    \definition{hard}    \example{
      \exsen{mamlla llokollokott}      \extran{hard rope}}}  \sense{\textbf{2.} 
    \definition{difficult}    \example{
      \exsen{bo ttam llokollokott}      \extran{her difficult life}}}  \sense{\textbf{3.} 
    \definition{stubborn}    \example{
      \exsen{Oba tikop a llokollokott dagaeya.}      \extran{Their hearts were stubborn.}}}}

\entry{llokott}{\headword{llokott}  \pronnote{ɽokoʈ͡ʂ}
  \pos{Modifier}  \sense{\textbf{1.} 
    \definition{hard, firm}    \example{
      \exsen{Ekaklle da llokttang dan.}      \extran{The ground is hard.}}    \example{
      \exsen{Däbe lla da obo labalaba de llokttangae damllamän.}      \extran{That man held his lap-lap firmly.}}}  \sense{\textbf{2.} 
    \definition{difficult, hard}    \example{
      \exsen{Oba Balimo me ttam giddoll a ddobae llokott daeya.}      \extran{Their life in Balimo was very difficult.}}}  \sense{\textbf{3.} 
    \definition{strong}    \example{
      \exsen{Imomdae lla da llokottang a ainin, ibi wi do ma.}      \extran{He who is truly strong is going to the house.}}}}

\entry{llollo}{\headword{llollo}  \pronnote{ɽoɽo}
  \pos{Verb}  \sense{
    \definition{to pay respects}    \example{
      \exsen{Llabunda bi llollo e gobällaballe.}      \extran{The relatives would go to pay their respects.}}}}

\entry{llollom}{\headword{llollom}  \pronnote{ɽoɽom}
  \pos{Intransitive verb}  \sense{\textbf{1.} 
    \definition{to break, be damaged}    \example{
      \exsen{Ekaklle we enanae gollomän kälakälae.}      \extran{It hit the ground, smashing into pieces.}}}  \sense{\textbf{2.} 
    \definition{to break}    \example{
      \exsen{Bogo karpo inkätt de dullomän.}      \extran{She broke the jar seal.}}}}

\entry{llowam}{\headword{llowam}  \pronnote{ɽowam}
  \pos{Noun}  \sense{\textbf{1.} 
    \definition{fatigue, tiredness}    \example{
      \exsen{Mark bom llowam da dägagän.}      \extran{Mark was tired (lit. fatigue got Mark).}}}  \sense{\textbf{2.} 
    \definition{disdain, hate}    \example{
      \exsen{Ngänäm llowam da nallan sana wätät a.}      \extran{I don't want to eat sago (lit. disdain to eat sago gets me).}}}}

\entry{llowamang}{\headword{llowamang}
  \pos{Modifier}  \sense{\textbf{1.} 
    \definition{tired}    \example{
      \exsen{Ngämo pätt a llowamang gogon.}      \extran{My body got tired.}}}  \sense{\textbf{2.} 
    \definition{unpreferable, unpleasant, tiresome (to someone in the dative)}    \example{
      \exsen{Ddob llayabira sosoga sana kängkäm a ddobae llowamang dan.}      \extran{Squeezing sosoga sago is really unpleasant to some people.}}    \example{
      \exsen{Ngämlle llowamang dan gänya llɨg bo peyang eka tameny a, be de maigag de aya näbäddan.}      \extran{I don't want to talk with this boy (lit. talking with this boy is unpreferable to me), but with the one who killed the bandicoot.}}    \example{
}}  \sense{\textbf{3.} 
    \definition{upset, annoyed}    \example{
      \exsen{Meri dallän Yokon pate ako llowamang dägagän obom.}      \extran{Mary went to Yokon and then got upset at her.}}}}

\entry{llowawi}{\headword{llowawi}  \note{Trees}  \pronnote{ɽowawi}
  \pos{Noun}  \sense{
    \definition{type of tree}    \note{Trees}    \example{
}    \example{
}    \example{
}}}

\entry{llullu}{\headword{llullu}  \pronnote{ɽuɽu}
  \pos{Verb}  \sense{
    \definition{to cut a tree so that one big tree falls onto many smaller trees and breaks them}}}

\entry{llune}{\headword{llune}  \pronnote{ɽune}
  \pos{Modifier}  \sense{
    \definition{wild}    \example{
      \exsen{Llune ddäddäg a dadegaeya.}      \extran{There were wild animals.}}}}

\entry{llupi}{\headword{llupi}
  \pos{Noun}  \sense{
    \definition{branch}    \example{
      \exsen{Mälla da llupi de mɨnyi bäträpnegnän giri alle.}      \extran{The woman will cut the branches with a knife.}}}}

\entry{llupi ttäganen ttägnen}{\headword{llupi ttäganen ttägnen}  \note{Games}  \pronnote{ɽupi ʈ͡ʂəɡanen ʈ͡ʂəɡnen}
  \pos{Noun}  \sense{
    \definition{type of game involving hiding a tree branch in the water}    \note{Games}    \example{
}    \example{
}    \example{
}}}

\entry{lluwam}{\headword{lluwam}  \pronnote{ɽuwam}
  \pos{Transitive verb}  \sense{
    \definition{to incapacitate, make unable}    \example{
      \exsen{Ngänäm wän a nying me nalluwaman Malläm e ibi we.}      \extran{This boil on my foot has made me unable to walk to Malam.}}}}

\entry{lo}{\headword{lo}
  \pos{Noun}  \sense{
    \definition{law}}}

\entry{Lois}{\headword{Lois}
  \pos{Proper noun}  \sense{
    \definition{female personal name}}}

\entry{lok}{\headword{lok}
  \pos{Transitive coverb}  \sense{
    \definition{lock}}}

\entry{loli}{\headword{loli}
  \pos{Noun}  \sense{
    \definition{candy}}}

\entry{Lomae}{\headword{Lomae}  \pronnote{lomae}
  \pos{Proper noun}  \sense{
    \definition{female personal name}}}

\entry{lomenyang}{\headword{lomenyang}  \pronnote{lomeɲaŋ}
  \pos{Noun}  \sense{
    \definition{someone who complains, complainer}}}

\entry{longgo}{\headword{longgo}  \pronnote{loŋɡo}
  \pos{Noun}  \sense{
    \definition{noise}    \example{
      \exsen{Llɨg obaene longgo de dändär, gongällbän a.}      \extran{I heard the kids making noise and woke up.}}}}

\entry{Loni}{\headword{Loni}  \pronnote{loni}
  \pos{Proper noun}  \sense{
    \definition{female personal name}}}

\entry{lonsis}{\headword{lonsis}  \note{Trees}  \pronnote{lonsis}
  \pos{Noun}  \sense{
    \definition{lemon plant}    \note{Trees}    \example{
}    \example{
}    \example{
}}}

\entry{Lovelyn}{\headword{Lovelyn}  \pronnote{ləvlɪn}
  \pos{Proper noun}  \sense{
    \definition{female personal name}}}

\entry{Ludwina}{\headword{Ludwina}  \pronnote{ludwina}
  \pos{Proper noun}  \sense{
    \definition{female personal name}}}

\entry{lugoe}{\headword{lugoe}
  \pos{Transitive verb}  \sense{
    \definition{to drag}}}

\entry{Luke}{\headword{Luke}
  \pos{Proper noun}  \sense{
    \definition{male personal name}}}

\entry{lulu}{\headword{lulu}  \pronnote{lulu}
  \pos{Transitive verb}  \sense{
    \definition{to gossip about, tell rumors about}    \example{
      \exsen{Luluang lla da ai dan gagäll ud de bätramän.}      \extran{One who gossips can open the door to bad things.}}    \example{
      \exsen{Ngäna dandärmällnegne mälla de ami bam daluneyo.}      \extran{I heard the women who were gossiping about you.}}    \example{
}}}

\entry{Lulu}{\headword{Lulu}  \pronnote{lulu}
  \pos{Proper noun}  \sense{
    \definition{female personal name}}}

\entry{lus}{\headword{lus}  \pronnote{lus}
  \pos{Intransitive coverb}  \sense{
    \definition{lose}}}

\entry{Lydia}{\headword{Lydia}  \pronnote{lidia}
  \pos{Proper noun}  \sense{
    \definition{female personal name}}}

\entry{Lynet}{\headword{Lynet}
  \variantof{SpellingVariant of \vartext{Linette}}}

\entry{Lyneth}{\headword{Lyneth}  \pronnote{linɛθ}
  \pos{Proper noun}  \sense{
    \definition{female personal name}}}

\chapter{M}


\entry{ma}{\headword{ma}  \pronnote{ma}
  \pos{Noun}  \sense{\textbf{1.} 
    \definition{house, home}    \example{
      \exsen{Ubi dadeg ma me?}      \extran{Are they home?}}    \example{
      \exsen{Abo bongo ttongo ma sisor de nogo.}      \extran{You must build a new house.}}}  \sense{\textbf{2.} 
    \definition{place, location}    \example{
      \exsen{kwatt ma, ine ma, Ende ma}      \extran{courthouse, well, Ende-speaking place}}}  \sense{\textbf{3.} 
    \definition{community}    \example{
      \exsen{Bin lla da mer de ngasnen erallo ma makäp me.}      \extran{Law enforcers do good in the community.}}}
  \variants{\SpellingVariant \vartext{maa}}}

\entry{=ma}{\headword{=ma}
  \pos{Clitic}  \sense{\textbf{1.} 
    \definition{characteristic case clitic (indicates purpose, source, or reason)}    \example{
      \exsen{ddäddäg, ddäddäg ma}      \extran{eat, edible}}    \example{
      \exsen{ikopse, ikopse ma ma}      \extran{prayer, house for prayer}}    \example{
      \exsen{Widere da gall gllae ma ttoen dan.}      \extran{An oar is a thing to paddle a canoe.}}}  \sense{\textbf{2.} 
    \definition{actor nominalizer}    \example{
      \exsen{kaptte ittall, kaptte ittall ma}      \extran{hanging clothes, clothesline}}}}

\entry{ma~}{\headword{ma~}
  \variantof{Infinitival reduplicant of \vartext{RED~}}}

\entry{ma kisin}{\headword{ma kisin}  \pronnote{ma kisin}
  \pos{Noun}  \sense{
    \definition{traditional house built directly on the ground}}}

\entry{ma ttängäm}{\headword{ma ttängäm}  \lexeme{ma ttängäm}  \pronnote{ma ʈ͡ʂəŋəm}
  \pos{Noun}  \sense{
    \definition{type of tree}}}

\entry{mab}{\headword{mab}  \note{Type of pandanus}  \pronnote{mab}
  \pos{Noun}  \sense{
    \definition{pandanus}    \scientificname{Pandanus}    \note{Type of pandanus}    \example{
      \exsen{Bogo tätäm mab de däddänän.}      \extran{Yesterday, he picked pandanus.}}    \example{
}    \example{
}}}

\entry{Mabudawan}{\headword{Mabudawan}  \note{Place names}  \pronnote{mabudawan}
  \pos{Proper noun}  \sense{
    \definition{Mabudawan/Mabaduan (Agob-speaking village in Kiwai Rural LLG near Saibai Island)}    \note{Place names}    \example{
}    \example{
}    \example{
}}}

\entry{mabun}{\headword{mabun}  \note{Oral Traditions}  \pronnote{mabun}
  \pos{Noun}  \sense{\textbf{1.} 
    \definition{clan totem (the sacred symbol of a clan group; usually a plant, animal, or body part)}    \note{Oral Traditions}    \example{
      \exsen{Obo mabun a dirom.}      \extran{Her totem is cassowary.}}}  \sense{\textbf{2.} 
    \definition{sacred}    \note{Oral Traditions}    \example{
      \exsen{mabun ttoen}      \extran{sacred story}}    \example{
}    \example{
}}}

\entry{mäd kubull}{\headword{mäd kubull}  \note{Animal names}  \pronnote{məd kubuɽ}
  \pos{Noun}  \sense{
    \definition{black bush wallaby}    \note{Animal names}    \example{
}    \example{
}    \example{
}}}

\entry{madä}{\headword{madä}
  \variantof{Dialectal Variant of \vartext{mäda}}  \pronnote{madə}}

\entry{mäda}{\headword{mäda}  \pronnote{məda}
  \pos{Kinship noun}  \sense{\textbf{1.} 
    \definition{father}    \example{
      \exsen{Ge obo mäda wainen.}      \extran{This is his father.}}}  \sense{\textbf{2.} 
    \definition{owner (of an animal)}    \example{
      \exsen{Däräng a oba mäda Kwalde bäne ibiatt de dongkollmälleyo.}      \extran{The dogs followed their owner Kwalde's footprints.}}}
  \variants{\Unspecified Variant \vartext{medä}}
  \variants{\Dialectal Variant \vartext{meda, madä}}
  \variants{\SpellingVariant \vartext{mäde}}}

\entry{mädameny}{\headword{mädameny}  \pronnote{mədameɲ}
  \pos{Modifier}  \sense{
    \definition{fatherless}}}

\entry{mädaolle}{\headword{mädaolle}
  \variantof{SpellingVariant of \vartext{mädaulle}}}

\entry{mädaulle}{\headword{mädaulle}  \lexeme{mädaulle}  \pronnote{mədauɽe}
  \pos{Kinship noun}  \sense{\textbf{1.} 
    \definition{aunt (one's father's elder sister)}}  \sense{\textbf{2.} 
    \definition{uncle (one's father's elder brother)}    \example{
      \exsen{Mädaulle daeya pu wi nallan.}      \extran{Uncle went to the swamp garden.}}}  \sense{\textbf{3.} 
    \definition{brother-in-law (woman's husband's elder brother)}}  \sense{\textbf{4.} 
    \definition{sister-in-law (woman's husband's elder sister or woman's younger brother's wife; reciprocal)}}
  \variants{\SpellingVariant \vartext{mädaolle}}}

\entry{mäde}{\headword{mäde}
  \variantof{SpellingVariant of \vartext{mäda}}  \pronnote{məde}}

\entry{mädi}{\headword{mädi}
  \variantof{Dialectal Variant of \vartext{mɨnyi}}}

\entry{madik}{\headword{madik}  \note{Tools}  \pronnote{madik}
  \pos{Noun}  \sense{
    \definition{type of tool}    \note{Tools}    \example{
}    \example{
}    \example{
}}}

\entry{Madima}{\headword{Madima}  \pronnote{madima}
  \pos{Proper noun}  \sense{
    \definition{female personal name}}}

\entry{Madlin}{\headword{Madlin}
  \pos{Proper noun}  \sense{
    \definition{female personal name}}}

\entry{madmed}{\headword{madmed}  \note{Trees}  \pronnote{madmed}
  \pos{Noun}  \sense{
    \definition{type of big tree that grows in the bush}    \note{Trees}    \example{
}    \example{
}    \example{
}}}

\entry{Mado}{\headword{Mado}  \pronnote{mado}
  \pos{Proper noun}  \sense{
    \definition{male personal name}}}

\entry{madu ma}{\headword{madu ma}
  \variantof{SpellingVariant of \vartext{maduma}}}

\entry{maduma}{\headword{maduma}  \pronnote{maduma}
  \pos{Noun}  \sense{
    \definition{village}    \example{
      \exsen{Obo dindugän Llimoll maduma we.}      \extran{He ran to Limol village.}}}
  \variants{\SpellingVariant \vartext{madu ma, maaduma}}}

\entry{madumame maket}{\headword{madumame maket}  \note{Market}  \pronnote{madumame maket}
  \pos{Noun}  \sense{
    \definition{black market}    \note{Market}    \example{
}    \example{
}    \example{
}}}

\entry{Madura}{\headword{Madura}  \pronnote{madura}
  \pos{Proper noun}  \sense{
    \definition{male personal name}}}

\entry{madura}{\headword{madura}  \note{Names of bananas}  \pronnote{madura}
  \pos{Noun}  \sense{
    \definition{type of introduced banana}    \note{Names of bananas}    \example{
}    \example{
}    \example{
}}}

\entry{=mae}{\headword{=mae}  \pronnote{mae}
  \pos{Nominal enclitic}  \sense{
    \definition{perlative case clitic; through}    \example{
      \exsen{Däbe dirom ulle de dazu kuddäll e mae.}      \extran{I shot that big cassowary to death.}}}
  \variants{\Dialectal Variant \vartext{=meae}}}

\entry{=mae}{\headword{=mae}  \pronnote{mae}
  \pos{Nominal enclitic}  \sense{
    \definition{restrictive clitic; only}    \example{
      \exsen{Ngäna Ende eka walle mae eka allan.}      \extran{I only speak in Ende.}}}
  \variants{\Dialectal Variant \vartext{=meae}}}

\entry{mae}{\headword{mae}
  \variantof{Dialectal Variant of \vartext{=meae}}}

\entry{maebo}{\headword{maebo}  \pronnote{maebo}
  \pos{Noun}  \sense{
    \definition{type of spear made from sago}    \example{
}    \example{
}    \example{
}}}

\entry{maemae}{\headword{maemae}  \lexeme{maemae}  \pronnote{maemae}
  \pos{Noun}  \sense{
    \definition{type of big tree that grows in places where yam gardens are planted; used for making canoes and paddles}}}

\entry{maer}{\headword{maer}  \pronnote{maer}
  \pos{Noun}  \sense{
    \definition{myrrh}}}

\entry{mag}{\headword{mag}
  \pos{}}

\entry{mäg}{\headword{mäg}  \lexeme{mäg}  \pronnote{məɡ}
  \pos{Kinship noun}  \sense{\textbf{1.} 
    \definition{mother}    \example{
      \exsen{Ngämo mäg bo bin a Sana.}      \extran{My mother's name is Sana.}}    \example{
}    \example{
}}  \sense{\textbf{2.} 
    \definition{source}    \example{
      \exsen{Ge buk a ibim taempmeny anggan mägda abal eka midd de.}      \extran{This book shows us the original meaning.}}}}

\entry{mäg ulle}{\headword{mäg ulle}  \pronnote{məɡ uɽe}
  \pos{Kinship noun}  \sense{
    \definition{aunt (one's father's elder brother's wife)}    \example{
}    \example{
}    \example{
}}}

\entry{mäga}{\headword{mäga}  \note{Sago palms}  \pronnote{məɡa}
  \pos{Noun}  \sense{
    \definition{type of sago}    \note{Sago palms}    \example{
}    \example{
}    \example{
}}}

\entry{Mägaem}{\headword{Mägaem}
  \variantof{Unspecified Variant of \vartext{Mäkayam}}}

\entry{mägäll}{\headword{mägäll}  \note{Trees}  \pronnote{məɡəɽ}
  \pos{Noun}  \sense{\textbf{1.} 
    \definition{type of tree with bark used to tie spears, white flowers, edible leaves, and finger-sized, edible red fruit with edible seeds}    \note{Trees}    \example{
}    \example{
}    \example{
}}  \sense{\textbf{2.} 
    \definition{strength}    \note{Trees}}}

\entry{mägällma}{\headword{mägällma}  \note{Friendship/social terms of endearment}  \pronnote{məɡəɽma}
  \pos{Noun}  \sense{
    \definition{friends who share a twin fruit from the mägäll tree}    \note{Friendship/social terms of endearment}    \example{
}    \example{
}    \example{
}}}

\entry{Mägayam}{\headword{Mägayam}
  \variantof{Unspecified Variant of \vartext{Mäkayam}}}

\entry{mägda}{\headword{mägda}  \pronnote{məɡda}
  \pos{Noun}  \sense{\textbf{1.} 
    \definition{coconut body (as opposed to the shoot)}}  \sense{\textbf{2.} 
}}

\entry{mägda koll}{\headword{mägda koll}  \pronnote{məɡda koɽ}
  \pos{Noun}  \sense{
    \definition{plant term}    \example{
      \exsen{Gaguma me mägda koll otät dag.}      \extran{The very big yams can be preserved in the yam house.}}}}

\entry{Maggie}{\headword{Maggie}
  \variantof{SpellingVariant of \vartext{Megi}}  \pronnote{maɡɡie}}

\entry{mägmeny}{\headword{mägmeny}  \pronnote{məɡmeɲ}
  \pos{Modifier}  \sense{
    \definition{motherless}}}

\entry{mai}{\headword{mai}  \note{Sago palms}  \pronnote{ma'i}
  \pos{Noun}  \sense{
    \definition{type of sago}    \note{Sago palms}    \example{
}    \example{
}    \example{
}}}

\entry{maigag}{\headword{maigag}  \note{Names of bananas}  \pronnote{maiɡaɡ}
  \pos{Noun}  \sense{
    \definition{type of introduced banana}    \note{Names of bananas}    \example{
}}}

\entry{maigag}{\headword{maigag}  \note{Names of bananas}  \pronnote{maiɡaɡ}
  \pos{Noun}  \sense{
    \definition{northern brown bandicoot}    \scientificname{Isoodon macrourus}    \note{Names of bananas}}}

\entry{maiwa}{\headword{maiwa}
  \pos{Noun}  \sense{
    \definition{type of pandanus with a long, smooth fruit cooked in mumu}    \example{
}}}

\entry{maiya}{\headword{maiya}  \pronnote{maija}
  \pos{Noun}  \sense{
    \definition{bulbil (fruit of yam that is planted)}}}

\entry{mäk}{\headword{mäk}  \pronnote{mək}
  \pos{Noun}  \sense{
    \definition{war}    \example{
      \exsen{Ngäma yae bi dazernän, ge mäk a angde gongesän.}      \extran{Our (excl.) mother lived there when this war happened.}}    \example{
}    \example{
}}}

\entry{mak}{\headword{mak}  \pronnote{mak}
  \pos{Noun}  \sense{\textbf{1.} 
    \definition{mark}}  \sense{\textbf{2.} 
    \definition{to choose}    \example{
      \exsen{Ubi komllaebmae llo de mak dägaeyo.}      \extran{They each chose a tree.}}}}

\entry{Mak}{\headword{Mak}
  \pos{Proper noun}  \sense{
    \definition{male personal name}}
  \variants{\SpellingVariant \vartext{Mark}}}

\entry{mäk lla}{\headword{mäk lla}  \note{Conflict/war}  \pronnote{mək ɽa}
  \pos{Noun}  \sense{
    \definition{soldier}    \note{Conflict/war}    \example{
      \exsen{Paelet Yesu bom Rom mäk lla yaba pate däntäpeyamän.}      \extran{Pilate handed Jesus over to the Roman soldiers.}}}}

\entry{mäk lla ami doko}{\headword{mäk lla ami doko}  \note{Conflict/war}  \pronnote{mək ɽa ami doko}
  \pos{Noun}  \sense{
    \definition{veteran}    \note{Conflict/war}    \example{
}    \example{
}    \example{
}}}

\entry{mäk lla piar doko}{\headword{mäk lla piar doko}  \note{Conflict/war}  \pronnote{mək ɽa piar doko}
  \pos{Noun}  \sense{
    \definition{veteran}    \note{Conflict/war}    \example{
}    \example{
}    \example{
}}}

\entry{mäka}{\headword{mäka}  \pronnote{məka}
  \pos{Noun}  \sense{
    \definition{plantar wart (on the soles of feet; caused by a virus)}}}

\entry{Makaka}{\headword{Makaka}
  \pos{Proper noun}  \sense{
    \definition{female personal name}}}

\entry{makäm}{\headword{makäm}  \note{Swadesh}  \pronnote{makəm}
  \pos{Noun}  \sense{
    \definition{space under a house}    \note{Swadesh}    \example{
      \exsen{Ngämi makäm me gägaebllaeb eran.}      \extran{We (excl.) are standing underneath the house.}}    \example{
}    \example{
}}}

\entry{makäm katrekatre}{\headword{makäm katrekatre}  \pronnote{makəm katrekatre}
  \pos{Noun}  \sense{
    \definition{area to sit or rest underneath the house}}}

\entry{mäkämäkäp}{\headword{mäkämäkäp}  \pronnote{məkəməkəp}
  \pos{Noun}  \sense{
    \definition{dreadlocks}}}

\entry{mäkamäke}{\headword{mäkamäke}  \pronnote{məkaməke}
  \pos{Transitive verb}  \sense{
    \definition{to use}    \example{
      \exsen{Ibi ge up de mullamulla we mäkanen eralla.}      \extran{We use this banana for medicine.}}}}

\entry{mäkan}{\headword{mäkan}
  \pos{Noun}  \sense{
    \definition{desire (of someone)}    \example{
      \exsen{Ngämo mäkan a pällämpälläm eka umllang e dan.}      \extran{I want (lit. my desire is) to learn English.}}}}

\entry{mäkäp}{\headword{mäkäp}  \pronnote{məkəp}
  \pos{Noun}  \sense{\textbf{1.} 
    \definition{knot}}  \sense{\textbf{2.} 
    \definition{to knot}}}

\entry{makäp}{\headword{makäp}  \pronnote{makəp}
  \pos{Locational}  \sense{\textbf{1.} 
    \definition{inside, in, within, among}    \example{
      \exsen{Kumuddäga ebdo makäp me, melem bättemän.}      \extran{I will finish the work in three days.}}    \example{
      \exsen{Mɨnyi bina tuwelb makäp att ttongo da sänge bagän ngänäm.}      \extran{One among you twelve will betray me.}}    \example{
      \exsen{Angde dälltamän, ngäna wayati de dänye ngämo nagnag aba makäp me.}      \extran{When they arrived, I shared the watermelon with my friends.}}}  \sense{\textbf{2.} 
    \definition{duration}    \example{
      \exsen{Diba mamoe makäp me, ngämi misde ddäddäg abaene towall llälläm emae lländräg gogaebne.}      \extran{During the hunt, we were just listening for the rustling of the animals in the grass.}}    \example{
      \exsen{Bogo dallnän yuwog abal dektta yaba pate oblle ngämingg e, be bogo ako orwa gognän oba ngämingg makäp me.}      \extran{She would go to many doctors so they could help her, but during their treatment, she would suffer again.}}}}

\entry{mäkat}{\headword{mäkat}  \note{Animals/Plants}  \pronnote{məkat}
  \pos{Noun}  \sense{
    \definition{rat}    \note{Animals/Plants}    \example{
      \exsen{Mäkat llɨgllɨgag a dan bäne gaguma me.}      \extran{There is a rat with babies in your yamhouse.}}    \example{
}    \example{
}}}

\entry{Mäkayam}{\headword{Mäkayam}  \note{Language Names}  \pronnote{məkajam}
  \pos{Proper noun}  \sense{
    \definition{Makayam/Tirio language}    \note{Language Names}    \example{
}    \example{
}    \example{
}}
  \variants{\Unspecified Variant \vartext{Mägaem, Mägayam}}}

\entry{maket}{\headword{maket}
  \pos{Noun}  \sense{
    \definition{market}    \example{
      \exsen{Ngäna maket me melemang dan.}      \extran{I work in the market.}}}
  \variants{\freevarlab \vartext{makett}}
  \variants{\SpellingVariant \vartext{market}}}

\entry{makett}{\headword{makett}
  \variantof{freevarlab of \vartext{maket}}}

\entry{Maki}{\headword{Maki}  \pronnote{maki}
  \pos{Proper noun}  \sense{
    \definition{female personal name}}}

\entry{malam}{\headword{malam}  \pronnote{malam}
  \pos{Transitive verb}  \sense{\textbf{1.} 
    \definition{to obey, follow}    \example{
      \exsen{Lla da obäne eka de dämalaemneyo.}      \extran{People used to obey his words.}}}  \sense{\textbf{2.} 
    \definition{to accept}    \example{
      \exsen{Ngäna bänene kandärmang eka de malam eran.}      \extran{I accept your apology.}}}}

\entry{mälamäle}{\headword{mälamäle}  \pronnote{məlaməle}
  \pos{Transitive verb}  \sense{\textbf{1.} 
    \definition{to patch}}  \sense{\textbf{2.} 
    \definition{to dress (a wound)}    \example{
      \exsen{Känaebag, ngäna mɨnyi ute de bamle.}      \extran{Tomorrow, I will dress the wound.}}}}

\entry{mälanzeg}{\headword{mälanzeg}
  \pos{Transitive verb}  \sense{
    \definition{to tie}}}

\entry{mall~}{\headword{mall~}
  \variantof{Inflected Form of \vartext{RED~}}}

\entry{mälla}{\headword{mälla}  \lexeme{mälla}  \pronnote{məɽa}
  \pos{Noun}  \sense{\textbf{1.} 
    \definition{woman, female}    \anthronote{the hand sign for woman is to hold up the breasts}}  \sense{\textbf{2.} 
    \definition{wife}    \example{
      \exsen{Bäne baba auli mälla peyang dan?}      \extran{How many wives does your father have?}}}}

\entry{mälla~}{\headword{mälla~}
  \variantof{Plural reduplicant of \vartext{RED~}}}

\entry{malla}{\headword{malla}  \pronnote{maɽa}
  \pos{Negation particle}  \sense{
    \definition{not}    \example{
      \exsen{Ngäna tatuma ibi allan be malla mɨnyi tatuma me de.}      \extran{I am going to wash, but not at the washing place.}}    \example{
      \exsen{Bogo malla gänyme zegatt dan.}      \extran{He was not born here.}}}}

\entry{mälla ause}{\headword{mälla ause}
  \variantof{freevarlab of \vartext{mällause}}}

\entry{mälla gämäll}{\headword{mälla gämäll}  \note{Conflict/war}  \pronnote{məɽa ɡəməɽ}
  \pos{}  \sense{
    \definition{adultery (of a man)}    \note{Conflict/war}    \example{
}    \example{
}    \example{
}}}

\entry{mälla wareka}{\headword{mälla wareka}  \note{Conflict/war}  \pronnote{məɽa wareka}
  \pos{}  \sense{
    \definition{jealous of wife}    \note{Conflict/war}    \example{
}    \example{
}    \example{
}}}

\entry{mälla yae}{\headword{mälla yae}  \lexeme{mälla yae}  \pronnote{məɽa jae}
  \pos{Kinship noun}  \sense{
    \definition{aunt (one's mother's elder sister)}}}

\entry{mallakläk}{\headword{mallakläk}  \pronnote{maɽaklək}
  \pos{}  \sense{
    \definition{collapse}}}

\entry{mällakutang}{\headword{mällakutang}  \note{Trees}  \pronnote{məɽakutaŋ}
  \pos{Noun}  \sense{
    \definition{type of tree that grows in the bush with white flowers, black bark, and wood used for house sticks}    \note{Trees}    \example{
}    \example{
}    \example{
}}}

\entry{mälläll}{\headword{mälläll}
  \variantof{Unspecified Variant of \vartext{märäl}}}

\entry{mällam}{\headword{mällam}  \pronnote{məɽam}
  \pos{Noun}  \sense{
    \definition{type of plant}}}

\entry{mällam}{\headword{mällam}  \pronnote{məɽam}
  \pos{Transitive verb}  \sense{
    \definition{to hold; get, grab, catch}    \example{
      \exsen{Rop de namllam ttang me.}      \extran{[You] hold the rope in your hand.}}    \example{
      \exsen{Ttongo ngäna skul me tresära bin di mällaem eran.}      \extran{I hold the position of treasurer at a school.}}    \example{
      \exsen{Bogo llo käp de naspunan nag da pate be nag da ddone namllaman.}      \extran{She threw the fruit to her friend but they didn't catch it.}}    \example{
      \exsen{Kallkäll a damllamän obom.}      \extran{He got cold (lit. the cold got him).}}}}

\entry{mallam}{\headword{mallam}
  \variantof{Unspecified Variant of \vartext{Malläm}}  \pronnote{maɽam}}

\entry{Malläm}{\headword{Malläm}  \pronnote{maɽəm}
  \pos{Proper noun}  \sense{
    \definition{Malam (Ende-speaking village in Morehead Rural LLG; a two-hour walk 9.3km) from Limol; GPS: -8.712716, 142.656097)}}
  \variants{\Unspecified Variant \vartext{mallam}}}

\entry{mällamälla}{\headword{mällamälla}  \note{Verbs}  \pronnote{məɽaməɽa}
  \pos{Transitive verb}  \sense{
    \definition{to tie}    \note{Verbs}    \example{
      \exsen{Yu wi daspuniyu, mällamällawatt de, ada oba yu a bättämän.}      \extran{They threw him on the fire and tied him so that the fire would burn him.}}    \example{
      \exsen{Ttongo lla da ddone mullae dan mängallang lla bo ma we bozenän obo za gämäll e, be dam de da bogo ngattong däbe lla de bämllawän.}      \extran{A person cannot enter the house of a powerful man to steal his things unless they tie that man up first.}}}
  \variants{\Unspecified Variant \vartext{mällamälle}}}

\entry{mällamälle}{\headword{mällamälle}
  \variantof{Unspecified Variant of \vartext{mällamälla}}}

\entry{mällamollang}{\headword{mällamollang}  \note{Numbers}  \pronnote{məɽamoɽaŋ}
  \pos{}  \sense{
    \definition{more than two couples}    \note{Numbers}    \example{
}    \example{
}    \example{
}}}

\entry{mälläng}{\headword{mälläng}  \note{Parts of a bird}  \pronnote{məɽəŋ}
  \pos{Noun}  \sense{
    \definition{nose}    \note{Parts of a bird}    \example{
      \exsen{Ngämo mälläng a nongonongorang dan.}      \extran{My nose is itchy.}}    \example{
}    \example{
}}
  \variants{\SpellingVariant \vartext{mällɨng, mɨllɨng}}}

\entry{mälläng ik}{\headword{mälläng ik}  \pronnote{məɽəŋ ik}
  \pos{Noun}  \sense{
    \definition{nostril}    \example{
      \exsen{Obo mälläng ik a ulleulle abal dagaeya.}      \extran{His nostrils were very big.}}}}

\entry{mälläng kobllokobllom}{\headword{mälläng kobllokobllom}  \note{Parts of the Body}  \pronnote{məɽəŋ kobɽokobɽom}
  \pos{Noun}  \sense{
    \definition{part of the body}    \note{Parts of the Body}    \example{
}    \example{
}    \example{
}}}

\entry{mälläng kom}{\headword{mälläng kom}  \note{Parts of an insect}  \pronnote{məɽəŋ kom}
  \pos{Noun}  \sense{\textbf{1.} 
    \definition{nose hair}    \note{Parts of an insect}}  \sense{\textbf{2.} 
    \definition{antenna}    \note{Parts of an insect}    \example{
}    \example{
}    \example{
}}}

\entry{mälläng obod}{\headword{mälläng obod}  \pronnote{məɽəŋ obod}
  \pos{Noun}  \sense{
    \definition{runny nose}    \example{
}    \example{
}    \example{
}}}

\entry{mälläng seam}{\headword{mälläng seam}  \pronnote{məɽəŋ seam}
  \pos{Intransitive compound verb}  \sense{
    \definition{to wrinkle up your nose and breathe out, reaction to a bad smell.}}}

\entry{mälläng sobasoba}{\headword{mälläng sobasoba}  \note{Parts of a reptile}  \pronnote{məɽəŋ sobasoba}
  \pos{Noun}  \sense{
    \definition{snout}    \note{Parts of a reptile}    \example{
}    \example{
}    \example{
}}}

\entry{mälläng wällong aga}{\headword{mälläng wällong aga}  \pronnote{məɽəŋ wəɽoŋ aɡa}
  \pos{Noun}  \sense{
    \definition{"nose swelled up", very serious, quiet, ignoring people}}}

\entry{mällänggäbe}{\headword{mällänggäbe}  \note{Water}  \pronnote{məɽəŋɡəbe}
  \pos{Noun}  \sense{
    \definition{type of plant that grows along the riverside}    \note{Water}    \example{
}    \example{
}    \example{
}}}

\entry{mällät}{\headword{mällät}  \lexeme{mällät}  \pronnote{məɽət}
  \pos{Noun}  \sense{
    \definition{type of tree that grows in the grassland with fruit used to ignite fires and yellow flowers with edible nectar}}}

\entry{mällät käp}{\headword{mällät käp}  \pronnote{məɽət kəp}
  \pos{Noun}  \sense{
    \definition{second stage of sago growth}    \example{
      \exsen{Sana da mällät käpang dan.}      \extran{The sago is in its second stage of growth.}}}}

\entry{mällät käp}{\headword{mällät käp}  \note{confirm species. it was previously called "white dot python" and this is a species with a characteristic white dot pattern only around the mouth}  \pronnote{məɽət kəp}
  \pos{Noun}  \sense{
    \definition{D'Albertis python}    \scientificname{Leiopython albertisii}    \note{confirm species. it was previously called "white dot python" and this is a species with a characteristic white dot pattern only around the mouth}    \example{
      \exsen{Ag me ibi näddägallo mällät käp gullem de.}      \extran{This morning, we (incl.) ate a mällät käp snake.}}}}

\entry{mällätbudar}{\headword{mällätbudar}  \note{Grubs}  \pronnote{məɽətbudar}
  \pos{Noun}  \sense{
    \definition{type of grub}    \note{Grubs}    \example{
}    \example{
}    \example{
}}}

\entry{mällätgugu}{\headword{mällätgugu}  \note{Types of Taro}  \pronnote{məɽətɡuɡu}
  \pos{Noun}  \sense{
    \definition{type of taro}    \note{Types of Taro}    \example{
}    \example{
}    \example{
}}}

\entry{mällause}{\headword{mällause}  \note{Family ties}  \pronnote{məɽause}
  \pos{Noun}  \sense{
    \definition{old woman}    \note{Family ties}    \example{
      \exsen{Mällause bo auma da dan de.}      \extran{The old woman's grave is there.}}    \example{
}    \example{
}}
  \variants{\freevarlab \vartext{mälla ause}}}

\entry{mällawang}{\headword{mällawang}  \pronnote{məɽawaŋ}
  \pos{Modifier}  \sense{
    \definition{married (of a male)}}
  \variants{\Fast Speech Variant \vartext{mällong}}}

\entry{mallet}{\headword{mallet}
  \variantof{Unspecified Variant of \vartext{maret}}  \pronnote{maɽet}}

\entry{mällɨng}{\headword{mällɨng}
  \variantof{SpellingVariant of \vartext{mälläng}}}

\entry{mällɨng pän}{\headword{mällɨng pän}  \lexeme{pän}  \pronnote{pən}
  \pos{Intransitive verb}  \sense{\textbf{1.} 
    \definition{to ripen, e.g. of a papaya. Literally: to decorate a nose.}}  \sense{\textbf{2.} 
    \definition{to decorate nose}}}

\entry{mällka~}{\headword{mällka~}
  \variantof{Derived Variant of \vartext{RED~}}}

\entry{mällkae}{\headword{mällkae}
  \pos{Intransitive verb}  \sense{
    \definition{to bend}    \example{
      \exsen{Kote dämällkaenän obo pate.}      \extran{He was bowing his head (lit. bent nape) to her.}}}}

\entry{mällkakallamatt}{\headword{mällkakallamatt}  \note{Snakes}  \pronnote{məɽkakaɽamaʈ͡ʂ}
  \pos{Noun}  \sense{
    \definition{type of venomous snake}    \note{Snakes}    \example{
}    \example{
}    \example{
}}}

\entry{mällkam}{\headword{mällkam}
  \pos{Intransitive verb}  \sense{
    \definition{to bend}    \example{
      \exsen{Kote gomällkamän ngänaeka de gongkamän.}      \extran{He hung his head (lit. bent nape) and started to cry.}}}}

\entry{mällkamällkam}{\headword{mällkamällkam}  \pronnote{məɽkaməɽkam}
  \pos{Adverb}  \sense{
    \definition{bending over}}}

\entry{mallmell}{\headword{mallmell}
  \variantof{Unspecified Variant of \vartext{malmal}}}

\entry{mällong}{\headword{mällong}
  \variantof{Fast Speech Variant of \vartext{mällawang}}  \pronnote{məɽoŋ}}

\entry{mällpa}{\headword{mällpa}  \note{Family ties}  \pronnote{məɽpa}
  \pos{Kinship noun}  \sense{
    \definition{aunt (one's mother's sister)}    \note{Family ties}    \example{
}    \example{
}    \example{
}}}

\entry{mällpe}{\headword{mällpe}
  \pos{Noun}  \sense{
}}

\entry{malmal}{\headword{malmal}  \note{Gardens}  \pronnote{malmal}
  \pos{Transitive verb}  \sense{\textbf{1.} 
    \definition{to mark land (for planting or settlement)}    \note{Gardens}    \example{
      \exsen{Bobäll ttängäm malmal e.}      \extran{Let's go mark a garden spot.}}    \example{
      \exsen{Ge ma de dämaleyo.}      \extran{They chose to settle this area.}}}  \sense{\textbf{2.} 
    \definition{to accuse, blame}    \note{Gardens}    \example{
      \exsen{Ubi yuwog abal ttoen e bam malmal nallo.}      \extran{They are accusing you of so many things.}}}
  \variants{\Unspecified Variant \vartext{mallmell}}}

\entry{mälmäl}{\headword{mälmäl}  \note{Verbs}  \pronnote{məlməl}
  \pos{Transitive verb}  \sense{\textbf{1.} 
    \definition{to squeeze}    \note{Verbs}    \example{
}    \example{
}    \example{
}}  \sense{\textbf{2.} 
    \note{Verbs}}}

\entry{malmel}{\headword{malmel}  \pronnote{malmel}
  \pos{Verb}  \sense{
    \definition{to squeeze (e.g. a hand)}}}

\entry{malonäbe}{\headword{malonäbe}  \lexeme{malonäbe}  \pronnote{malonəbe}
  \pos{Noun}  \sense{
    \definition{type of reed that is too soft for weaving}}}

\entry{mam}{\headword{mam}  \note{Parts of the Body}  \pronnote{mam}
  \pos{Noun}  \sense{\textbf{1.} 
    \definition{blood}    \note{Parts of the Body}    \example{
      \exsen{Ngäna ddia mam de dalltäne.}      \extran{I was following the blood.}}    \example{
}    \example{
}}  \sense{\textbf{2.} 
    \definition{to bleed}    \note{Parts of the Body}    \example{
      \exsen{Da ngäna volleyball botngoe känaebag, tupi da mɨnyi mam bogon.}      \extran{If I play volleyball tomorrow, my pointer finger will bleed.}}}  \sense{\textbf{3.} 
    \definition{red; pink}    \note{Parts of the Body}    \example{
      \exsen{Emi ngänaeka gogon a obo ikop a mamang gognegän.}      \extran{Emi cried and her eyes were red.}}}
  \variants{\Dialectal Variant \vartext{mem}}}

\entry{mam titi popo}{\headword{mam titi popo}  \pronnote{mam titi popo}
  \pos{Color term}  \sense{
    \definition{dark red}}}

\entry{mama}{\headword{mama}  \pronnote{mama}
  \pos{Noun}  \sense{
    \definition{food (baby talk word)}}}

\entry{mama}{\headword{mama}  \pronnote{mama}
  \pos{Noun}  \sense{\textbf{1.} 
    \definition{grass pile}    \example{
      \exsen{Däräng dallän ada kukiny mama de ikop dägagän.}      \extran{The dog went and saw a pile of kukiny grass.}}}  \sense{\textbf{2.} 
}}

\entry{mama}{\headword{mama}  \pronnote{mama}
  \pos{Kinship noun}  \sense{
    \definition{mother}}}

\entry{mamaemamae}{\headword{mamaemamae}  \pronnote{mamaemamae}
  \pos{Modifier}  \sense{
    \definition{bloody}    \example{
      \exsen{mamaemamae kanas}      \extran{bloody spear}}}
  \variants{\Fast Speech Variant \vartext{mamamamae}}}

\entry{mamal}{\headword{mamal}
  \variantof{Unspecified Variant of \vartext{mamall}}}

\entry{mamall}{\headword{mamall}  \pronnote{mamaɽ}
  \pos{Adverb}  \sense{
    \definition{quickly}    \example{
      \exsen{Lläkäm a ttongo mamall zozoang dan.}      \extran{A mushroom goes rotten quickly.}}}
  \variants{\Unspecified Variant \vartext{mamal}}}

\entry{mamalldae}{\headword{mamalldae}
  \pos{Adverb}  \sense{
    \definition{immediately, right away}    \example{
      \exsen{Itrell me mamalldae meresin e gogmallo.}      \extran{In sickness, they go for medicine immediately.}}    \example{
      \exsen{Mamalldae wiyamom!}      \extran{(You all) come right away!}}}
  \variants{\Unspecified Variant \vartext{mamälldae}}}

\entry{mamälldae}{\headword{mamälldae}
  \variantof{Unspecified Variant of \vartext{mamalldae}}}

\entry{mamam}{\headword{mamam}  \note{Colors}  \pronnote{mamam}
  \pos{Color term}  \sense{
    \definition{red}    \note{Colors}    \example{
      \exsen{Ngämo mamam mätta de aya notan?}      \extran{Who ate my red yam?}}    \example{
}    \example{
}}}

\entry{mamamam}{\headword{mamamam}  \pronnote{mamamam}
  \pos{Color term}  \sense{
    \definition{red}    \example{
      \exsen{Gabma yaba pätt a ai dan mamall yäbäd att a mamamamang bogon.}      \extran{White people's bodies can quickly become red from the sun.}}}}

\entry{mamamamae}{\headword{mamamamae}
  \variantof{Fast Speech Variant of \vartext{mamaemamae}}  \pronnote{mamamamae}}

\entry{mamamamall}{\headword{mamamamall}  \pronnote{mamamamaɽ}
  \pos{Adverb}  \sense{
    \definition{quickly}    \example{
      \exsen{Abo bongo mamamamall ngämim yangkollmällneg.}      \extran{You must follow us quickly.}}}}

\entry{mamamemett}{\headword{mamamemett}  \note{Musical instruments}  \pronnote{mamamemeʈ͡ʂ}
  \pos{Modifier}  \sense{
    \definition{fast beat}    \note{Musical instruments}    \example{
}    \example{
}    \example{
}}}

\entry{mambag}{\headword{mambag}  \pronnote{mambaɡ}
  \pos{Noun}  \sense{
    \definition{type of goanna}    \example{
      \exsen{Sali bi mambag de ttägäll ik mi näbäddallo.}      \extran{Saly and them killed the mambag in the anthill.}}}}

\entry{mämbär}{\headword{mämbär}  \note{Food}  \pronnote{məmbər}
  \pos{Noun}  \sense{
    \definition{type of big, round yam with a white interior, hairs, and no thorns}    \note{Food}    \example{
}    \example{
}    \example{
}}}

\entry{Mame}{\headword{Mame}
  \pos{Proper noun}  \sense{
    \definition{female personal name}}}

\entry{mame}{\headword{mame}  \pronnote{mame}
  \pos{Intransitive verb}  \sense{
    \definition{to fall, rain}    \example{
      \exsen{Yogoll ulle da dämanän.}      \extran{It was raining a lot.}}}}

\entry{mameat}{\headword{mameat}  \pronnote{mameat}
  \pos{Noun}  \sense{
    \definition{papaya, pawpaw}    \scientificname{Carica papaya}    \example{
}    \example{
}    \example{
}}}

\entry{Mamen}{\headword{Mamen}  \note{Place names}  \pronnote{mamen}
  \pos{Proper noun}  \sense{
    \definition{Mamen (toponym)}    \note{Place names}    \example{
}    \example{
}    \example{
}}}

\entry{mami}{\headword{mami}  \pronnote{mami}
  \pos{Kinship noun}  \sense{
    \definition{mommy, mummy}}}

\entry{mamkiel}{\headword{mamkiel}  \note{Names of bananas}  \pronnote{mamkiel}
  \pos{Noun}  \sense{
    \definition{type of native banana}    \note{Names of bananas}    \example{
}    \example{
}    \example{
}}}

\entry{mamlla}{\headword{mamlla}  \lexeme{mamlla}  \pronnote{mamɽa}
  \pos{Noun}  \sense{
    \definition{rope, string}    \example{
      \exsen{Obo llan de mamlla walle dämällanegän.}      \extran{He tied his ears with string.}}}}

\entry{mamoe}{\headword{mamoe}  \pronnote{mamoe}
  \pos{Transitive verb}  \sense{\textbf{1.} 
    \definition{to hunt, go hunting}    \example{
      \exsen{Obo moko da mamoe e dan.}      \extran{He wants to go hunting.}}    \example{
      \exsen{Ubi därängmeny damoenegnän.}      \extran{They were hunting without dogs.}}    \example{
      \exsen{Namoeaemnalla.}      \extran{We were hunting them.}}}  \sense{\textbf{2.} 
    \definition{hunt}    \example{
      \exsen{Mamoe a angde gongkamalle, ddob lla da mɨnyi po we ngattong e gobällalle.}      \extran{When the hunt starts, some men will go first to block the way.}}}}

\entry{mamon}{\headword{mamon}  \pronnote{mamon}
  \pos{Transitive verb}  \sense{\textbf{1.} 
    \definition{to string (e.g. a bow, sago beater)}    \example{
      \exsen{Tätäm ngäna bägäl de damon.}      \extran{Yesterday, I strung the bow.}}    \example{
      \exsen{Bägäl a zime amonan?}      \extran{Is the bow strung?}}    \example{
      \exsen{Ttäle da amonan.}      \extran{My leg is strung (i.e. it's tight and I can't flex it).}}}  \sense{\textbf{2.} 
    \definition{to fashion, shape, make}    \example{
      \exsen{Mälla da mɨnyi abor de bamonän.}      \extran{The woman will fashion a sago beater.}}}}

\entry{mamos}{\headword{mamos}  \note{Birds}  \pronnote{mamos}
  \pos{Noun}  \sense{
    \definition{comb-crested jacana}    \scientificname{Irediparra gallinacea}    \note{Birds}    \example{
}    \example{
}    \example{
}}}

\entry{mamos}{\headword{mamos}  \note{Birds}  \pronnote{mamos}
  \pos{Noun}  \sense{
    \definition{village constable}    \note{Birds}    \example{
      \exsen{Mamos a ttängäm ikopang lla deya.}      \extran{The constable was a man who looked after the village.}}}}

\entry{mamos bo bun}{\headword{mamos bo bun}  \note{Government}  \pronnote{mamos bo bun}
  \pos{Noun}  \sense{
    \definition{term in government}    \note{Government}    \example{
}    \example{
}    \example{
}}}

\entry{mämrem}{\headword{mämrem}  \pronnote{məmrem}
  \pos{Intransitive coverb}  \sense{
    \definition{to growl}    \example{
      \exsen{Bogo mämrem peyang eka gogon.}      \extran{He spoke with a growl.}}}}

\entry{män}{\headword{män}  \lexeme{män}  \pronnote{mən}
  \pos{Noun}  \sense{
    \definition{type of big tree cultivated in the bush with white flowers and blue and purple fruit that cassowary eat; tobacco is planted in the soil near this tree after it is burned}}}

\entry{män~}{\headword{män~}
  \variantof{Plural reduplicant of \vartext{RED~}}}

\entry{män}{\headword{män}
  \pos{Intransitive verb}  \sense{
    \definition{to settle}}}

\entry{män}{\headword{män}
  \pos{Noun}  \sense{\textbf{1.} 
    \definition{girl}    \example{
      \exsen{Ngäna ge män de dangem.}      \extran{I adopted this girl.}}}  \sense{\textbf{2.} 
    \definition{sister}    \example{
}}  \sense{\textbf{3.} 
    \definition{She's my daughter.}    \example{
      \exsen{Bogo ngämo män kälsre dan.}}}}

\entry{män duwar}{\headword{män duwar}
  \pos{Noun}  \sense{
    \definition{young girl}    \example{
      \exsen{Män duwar da bun mi popo de nowattällan.}      \extran{The young girl put flowers in her hair.}}}
  \variants{\SpellingVariant \vartext{mänduar}}}

\entry{Mana}{\headword{Mana}  \pronnote{mana}
  \pos{Proper noun}  \sense{
    \definition{female personal name}}}

\entry{Manaleato}{\headword{Manaleato}  \pronnote{manaleato}
  \pos{Proper noun}  \sense{
    \definition{female personal name}}
  \variants{\SpellingVariant \vartext{Manaliato}}}

\entry{Manaliato}{\headword{Manaliato}
  \variantof{SpellingVariant of \vartext{Manaleato}}}

\entry{mänang}{\headword{mänang}  \lexeme{mänang}  \note{Family ties}  \pronnote{mənaŋ}
  \pos{Kinship noun}  \sense{\textbf{1.} 
    \definition{father-in-law (man's wife's father; reciprocal)}    \note{Family ties}}  \sense{\textbf{2.} 
    \definition{mother-in-law (man's wife's mother; reciprocal)}    \note{Family ties}}  \sense{\textbf{3.} 
    \definition{son-in-law (one's daughter's husband; reciprocal)}    \note{Family ties}}  \sense{\textbf{4.} 
    \definition{brother-in-law (one's younger sister's husband)}    \note{Family ties}}}

\entry{Manang}{\headword{Manang}
  \pos{Proper noun}  \sense{
    \definition{male personal name}}}

\entry{mänänggmit}{\headword{mänänggmit}  \note{Parts of the Body}  \pronnote{mənəŋɡmit}
  \pos{Noun}  \sense{
    \definition{body part}    \note{Parts of the Body}    \example{
}    \example{
}    \example{
}}}

\entry{mänaulle}{\headword{mänaulle}  \pronnote{mənauɽe}
  \pos{Temporal adverb}  \sense{
    \definition{every time}}}

\entry{mänd~}{\headword{mänd~}
  \variantof{Inflected Form of \vartext{RED~}}}

\entry{mända}{\headword{mända}  \note{Numbers}  \pronnote{mənda}
  \pos{Noun}  \sense{\textbf{1.} 
    \definition{thumb; big toe}    \note{Numbers}}  \sense{\textbf{2.} 
    \definition{five (lit. thumb; body counting numeral)}    \note{Numbers}    \example{
      \exsen{mända ollong}      \extran{five times}}    \example{
}    \example{
}}}

\entry{mändd~}{\headword{mändd~}
  \variantof{Infinitival reduplicant of \vartext{RED~}}}

\entry{mandde}{\headword{mandde}  \note{Part of the week}  \pronnote{manɖ͡ʐe}
  \pos{Noun}  \sense{
    \definition{Monday}    \note{Part of the week}    \example{
}    \example{
}    \example{
}}}

\entry{mänddmändd}{\headword{mänddmändd}  \note{Water}  \pronnote{mənɖ͡ʐmənɖ͡ʐ}
  \pos{Transitive verb}  \sense{\textbf{1.} 
    \definition{to drown, struggle in water}    \note{Water}    \example{
      \exsen{Bogo mänddmändd e gongkamän walle we.}      \extran{She started to drown in the water.}}    \example{
      \exsen{Bogo aoli gomänddän walle me.}      \extran{She almost drowned in the water.}}    \example{
}    \example{
}}  \sense{\textbf{2.} 
    \definition{to drown, force underwater}    \note{Water}    \example{
      \exsen{Bogo ddia de ine dämänddän.}      \extran{He drowned the deer in the water.}}}  \sense{\textbf{3.} 
    \note{Water}}}

\entry{mande gara}{\headword{mande gara}  \note{Water}  \pronnote{mande ɡara}
  \pos{Noun}  \sense{
    \definition{water leech}    \note{Water}    \example{
}    \example{
}    \example{
}}}

\entry{mändmänd}{\headword{mändmänd}  \pronnote{məndmənd}
  \pos{Transitive verb}  \sense{\textbf{1.} 
    \definition{to raise, rear}    \example{
      \exsen{Ubi obom damändeyo.}      \extran{They raised him.}}    \example{
      \exsen{Indonesia me gänya ddäddäg de damändaebneyo.}      \extran{In Indonesia, they tamed this animal.}}    \example{
}}  \sense{\textbf{2.} 
    \definition{to feed}    \example{
      \exsen{Ngäna obom kumuddäga nge alle bamänd.}      \extran{I will feed him three coconuts.}}}  \sense{\textbf{3.} 
}}

\entry{mandri}{\headword{mandri}  \note{Trees}  \pronnote{mandri}
  \pos{Noun}  \sense{
    \definition{cultivated lemon tree}    \note{Trees}    \example{
}    \example{
}    \example{
}}}

\entry{mänduar}{\headword{mänduar}
  \variantof{SpellingVariant of \vartext{män duwar}}}

\entry{Maneya}{\headword{Maneya}  \pronnote{maneja}
  \pos{Proper noun}  \sense{
    \definition{name of a person}}}

\entry{mang}{\headword{mang}  \pronnote{maŋ}
  \pos{Kinship noun}  \sense{
    \definition{brother (of a woman)}    \example{
      \exsen{Niki era Nikol bo mang dan.}      \extran{Nicky is Nicole's brother.}}}}

\entry{mang~}{\headword{mang~}
  \variantof{Plural reduplicant of \vartext{RED~}}}

\entry{Mang}{\headword{Mang}  \pronnote{maŋ}
  \pos{Proper noun}  \sense{
    \definition{Mang (toponym)}}}

\entry{mängal}{\headword{mängal}  \pronnote{məŋal}
  \pos{Modifier}  \sense{
    \definition{quick}    \example{
      \exsen{Ge towall a mängalae penongg allan.}      \extran{This grass is burning quickly.}}}}

\entry{mängäll}{\headword{mängäll}
  \variantof{Unspecified Variant of \vartext{mängall}}}

\entry{mängall}{\headword{mängall}  \pronnote{məŋaɽ}
  \pos{Noun}  \sense{\textbf{1.} 
    \definition{strength, power}    \example{
      \exsen{Bäne mängall a ulle dan, be ngämo kälsre dan.}      \extran{You are stronger than me (lit. your strength is big, but mine is small).}}}  \sense{\textbf{2.} 
    \definition{to struggle}    \example{
      \exsen{Däbe mälla da ada mängall gognän tudi nyongkoe e.}      \extran{That woman was struggling to pull the fishing line.}}}  \sense{\textbf{3.} 
}
  \variants{\Unspecified Variant \vartext{mängäll}}}

\entry{mängall eka}{\headword{mängall eka}  \pronnote{məŋaɽ eka}
  \pos{Noun}  \sense{
    \definition{encouragement, cheer}    \example{
      \exsen{Obo pallall lla da mäse obom mängall eka dägneyo.}      \extran{His supporters tried to cheer him on.}}}}

\entry{mängallang}{\headword{mängallang}  \note{good characteristics of people}  \pronnote{məŋaɽaŋ}
  \pos{Modifier}  \sense{
    \definition{strong}    \note{good characteristics of people}    \example{
      \exsen{ddobae mängallang lla}      \extran{very strong person}}    \example{
}    \example{
}}}

\entry{mängallmeny}{\headword{mängallmeny}  \pronnote{məŋaɽmeɲ}
  \pos{Modifier}  \sense{
    \definition{weak}    \example{
      \exsen{Sisri obo pätt a mängallmeny allan.}      \extran{His body is weak now.}}    \example{
      \exsen{Wel ulle da mängallmeny gogon.}      \extran{The storm grew weak.}}    \example{
}}}

\entry{mängalmängal}{\headword{mängalmängal}  \pronnote{məŋalməŋal}
  \pos{Modifier}  \sense{
    \definition{quick}    \example{
      \exsen{Ngämi wätät de mängamängalae yu dägaebeya.}      \extran{We (excl.) quickly cooked the food.}}    \example{
}    \example{
}}
  \variants{\Fast Speech Variant \vartext{mängamängal}}}

\entry{mängamängal}{\headword{mängamängal}
  \variantof{Fast Speech Variant of \vartext{mängalmängal}}}

\entry{Mangel}{\headword{Mangel}
  \pos{Proper noun}  \sense{
    \definition{Mangel (toponym)}}}

\entry{Manggeya}{\headword{Manggeya}  \pronnote{maŋɡeja}
  \pos{Proper noun}  \sense{
    \definition{female personal name}}}

\entry{manggo}{\headword{manggo}  \note{Medicine}  \pronnote{maŋɡo}
  \pos{Noun}  \sense{
    \definition{mango tree}    \note{Medicine}    \example{
      \exsen{manggo däg}      \extran{bunch of mango}}    \example{
}    \example{
}}}

\entry{mangkimangki}{\headword{mangkimangki}  \note{Games}  \pronnote{maŋkimaŋki}
  \pos{Noun}  \sense{
    \definition{type of game involving chasing}    \note{Games}    \example{
}    \example{
}    \example{
}}}

\entry{Mangkol}{\headword{Mangkol}  \pronnote{maŋkol}
  \pos{Proper noun}  \sense{
    \definition{female personal name}}}

\entry{manglle}{\headword{manglle}  \pronnote{maŋɽe}
  \pos{Transitive coverb}  \sense{
    \definition{to lure}    \example{
      \exsen{Ngäna ddia llɨg eka däga adawede oba manglle bäga.}      \extran{I made baby deer sounds to lure her.}}}}

\entry{manglle kämang}{\headword{manglle kämang}  \note{Numbers}  \pronnote{maŋɽe kəmaŋ}
  \pos{Modifier}  \sense{
    \definition{always together}    \note{Numbers}    \example{
}    \example{
}    \example{
}}}

\entry{mangllemanglle}{\headword{mangllemanglle}  \note{Numbers}  \pronnote{maŋɽemaŋɽe}
  \pos{Modifier}  \sense{
    \definition{always together}    \note{Numbers}    \example{
}    \example{
}    \example{
}}}

\entry{mangllong}{\headword{mangllong}  \pronnote{maŋɽoŋ}
  \pos{Modifier}  \sense{
    \definition{friendly, appealing, sociable}    \example{
      \exsen{Ngäna lla bo peyang mangllong dan.}      \extran{I am friendly with people.}}}}

\entry{mangmang}{\headword{mangmang}  \pronnote{maŋmaŋ}
  \pos{Kinship noun}  \sense{
    \definition{nonsingular form of mang}}}

\entry{mani}{\headword{mani}  \pronnote{mani}
  \pos{Noun}  \sense{\textbf{1.} 
    \definition{money}    \example{
      \exsen{Bongo ge käza de mani muang nägag.}      \extran{You should sell this crocodile for money.}}}  \sense{\textbf{2.} 
    \definition{kina (PGK, the currency of Papua New Guinea)}}}

\entry{mani käp}{\headword{mani käp}
  \pos{Noun}  \sense{
    \definition{money}}}

\entry{manika}{\headword{manika}  \lexeme{manika}  \pronnote{manika}
  \pos{Noun}  \sense{
    \definition{cassava}    \scientificname{Manihot esculenta}}}

\entry{mankäp}{\headword{mankäp}  \note{Parts of the Body}  \pronnote{mankəp}
  \pos{Noun}  \sense{
    \definition{calf (back of leg)}    \note{Parts of the Body}    \example{
}    \example{
}    \example{
}}}

\entry{mänkot}{\headword{mänkot}
  \pos{Noun}  \sense{\textbf{1.} 
    \definition{uninitiated person}}  \sense{\textbf{2.} 
    \definition{non-believer}}}

\entry{mänmän}{\headword{mänmän}  \lexeme{mänmän}  \pronnote{mənmən}
  \pos{Noun}  \sense{
    \definition{nonsingular form of män}    \example{
      \exsen{Mänmän a sana de zazaba we dändäraebeyo.}      \extran{Girls stuffed sago into sago bags.}}    \example{
      \exsen{Ngämo mänmän a aba bin a ada, Ina, Lulu, Jamila.}      \extran{My sisters' names are Ina, Lulu, and Jamila.}}}}

\entry{mänmän duwaduwar}{\headword{mänmän duwaduwar}
  \pos{Noun}  \sense{
    \definition{nonsingular form of män duwar}}}

\entry{mänmänpitepite}{\headword{mänmänpitepite}  \note{Animals/Plants}  \pronnote{mənmənpitepite}
  \pos{Noun}  \sense{
    \definition{type of plant that grows by the river; eaten by wallabies}    \note{Animals/Plants}    \example{
}    \example{
}    \example{
}}}

\entry{mantär}{\headword{mantär}  \note{Tools}  \pronnote{mantər}
  \pos{Noun}  \sense{
    \definition{type of flat woven rope made from tree bark.}    \note{Tools}    \example{
}    \example{
}    \example{
}}}

\entry{many}{\headword{many}
  \variantof{Dialectal Variant of \vartext{English loanword}}}

\entry{mänya~}{\headword{mänya~}
  \variantof{freevarlab of \vartext{RED~}}}

\entry{mänyamänyan}{\headword{mänyamänyan}
  \variantof{Fast Speech Variant of \vartext{mänyanmänyan}}}

\entry{mänyän}{\headword{mänyän}  \pronnote{məɲən}
  \pos{Noun}  \sense{
    \definition{type of fish}}}

\entry{mänyan}{\headword{mänyan}  \note{Family ties}  \pronnote{məɲan}
  \pos{Kinship noun}  \sense{\textbf{1.} 
    \definition{younger sibling of the same-sex (man's younger brother or woman's younger sister)}    \note{Family ties}    \example{
}    \example{
}    \example{
}}  \sense{\textbf{2.} 
    \definition{co-sibling-in-law (man's wife's younger sister's husband or woman's husband's younger brother's wife)}    \note{Family ties}}}

\entry{mänyanmänyan}{\headword{mänyanmänyan}  \pronnote{məɲanməɲan}
  \pos{Kinship noun}  \sense{
    \definition{younger siblings}}
  \variants{\Fast Speech Variant \vartext{mänyamänyan}}}

\entry{mänyi}{\headword{mänyi}
  \variantof{SpellingVariant of \vartext{mɨnyi}}  \pronnote{məɲi}}

\entry{mänymäny}{\headword{mänymäny}
  \pos{Verb}  \sense{
    \definition{to vomit}}}

\entry{märal}{\headword{märal}  \pronnote{məral}
  \pos{Intransitive coverb}  \sense{\textbf{1.} 
    \definition{to curse, swear}}  \sense{\textbf{2.} 
    \definition{to curse}    \example{
      \exsen{Bongo ikopse alle lla de ddone märal näga adawede obo kuddäll bogon.}      \extran{You cannot curse a person in prayer so that they die.}}    \example{
}    \example{
}}}

\entry{märäl}{\headword{märäl}
  \pos{Noun}  \sense{
    \definition{size}    \example{
      \exsen{allko märäl ziz}      \extran{a fly-sized insect}}    \example{
      \exsen{Lla zazer ma mälläll be obo pätt a ulle abal daeya.}      \extran{The size was so that a person could enter, but his body was really big.}}}
  \variants{\Unspecified Variant \vartext{mälläll, märäll}}}

\entry{märäl}{\headword{märäl}
  \pos{Noun}  \sense{
    \definition{age-mate, someone of the same age}    \example{
      \exsen{Warani ngämo märäl dan.}      \extran{Warani is my age-mate.}}    \example{
      \exsen{Ngämo märäl a aenen ke?}      \extran{Who's the same age as me (lit. my age-mate)?}}}
  \variants{\Unspecified Variant \vartext{mälläll, märäll}}}

\entry{märäll}{\headword{märäll}
  \variantof{Unspecified Variant of \vartext{märäl}}  \pronnote{mərəɽ}}

\entry{marat}{\headword{marat}  \note{Weaving}  \pronnote{marat}
  \pos{Noun}  \sense{
    \definition{weaving pattern with a circular bottom and rectangular sides}    \note{Weaving}    \example{
}    \example{
}    \example{
}}}

\entry{mare}{\headword{mare}  \note{Type of pandanus}  \pronnote{mare}
  \pos{Noun}  \sense{
    \definition{type of pandanus}    \note{Type of pandanus}    \example{
}    \example{
}    \example{
}}}

\entry{Mareas}{\headword{Mareas}
  \pos{Proper noun}  \sense{
    \definition{personal name}}}

\entry{Marega}{\headword{Marega}  \pronnote{mareɡa}
  \pos{Proper noun}  \sense{
    \definition{male personal name}}}

\entry{maret}{\headword{maret}
  \pos{Noun}  \sense{
    \definition{marriage}}
  \variants{\Unspecified Variant \vartext{mallet}}}

\entry{Mareyas}{\headword{Mareyas}
  \pos{Proper noun}  \sense{
    \definition{Mareyas (toponym)}}}

\entry{Maria}{\headword{Maria}  \pronnote{maria}
  \pos{Proper noun}  \sense{
    \definition{female personal name}}}

\entry{Marian}{\headword{Marian}
  \pos{Proper noun}  \sense{
    \definition{female personal name}}}

\entry{Marias}{\headword{Marias}
  \pos{Proper noun}  \sense{
    \definition{male personal name}}}

\entry{maribärät}{\headword{maribärät}  \note{Yams}  \pronnote{maribərət}
  \pos{Noun}  \sense{
    \definition{type of long yam with a white interior and hairs}    \note{Yams}    \example{
}    \example{
}    \example{
}}}

\entry{Marion}{\headword{Marion}  \pronnote{marion}
  \pos{Proper noun}  \sense{
    \definition{female personal name}}}

\entry{Mark}{\headword{Mark}
  \variantof{SpellingVariant of \vartext{Mak}}}

\entry{markae}{\headword{markae}
  \pos{Noun}  \sense{\textbf{1.} 
    \definition{white person}}  \sense{\textbf{2.} 
    \definition{white, Western; foreign}}}

\entry{markae bägäl}{\headword{markae bägäl}
  \pos{Noun}  \sense{
    \definition{gun}}}

\entry{markae eka}{\headword{markae eka}
  \pos{Proper noun}  \sense{
    \definition{English}}}

\entry{markaebärät}{\headword{markaebärät}  \note{Yams}  \pronnote{markaebərət}
  \pos{Noun}  \sense{
    \definition{type of medium-sized yam with a white interior, thorns, and no hairs}    \note{Yams}    \example{
}    \example{
}    \example{
}}}

\entry{market}{\headword{market}
  \variantof{SpellingVariant of \vartext{maket}}  \note{Market}  \pronnote{market}}

\entry{Martha}{\headword{Martha}  \pronnote{marθa}
  \pos{Proper noun}  \sense{
    \definition{female personal name}}}

\entry{Martin}{\headword{Martin}
  \variantof{SpellingVariant of \vartext{Maten}}  \pronnote{martin}}

\entry{Mary}{\headword{Mary}
  \variantof{SpellingVariant of \vartext{Meri}}}

\entry{Maryanne}{\headword{Maryanne}
  \variantof{SpellingVariant of \vartext{Merian}}}

\entry{mas}{\headword{mas}  \pronnote{mas}
  \pos{Noun}  \sense{
    \definition{small bone found in cassowaries}}}

\entry{Mas}{\headword{Mas}
  \pos{Proper noun}  \sense{
    \definition{male personal name}}}

\entry{masa~}{\headword{masa~}
  \variantof{freevarlab of \vartext{RED~}}}

\entry{masaka}{\headword{masaka}  \lexeme{masaka}  \pronnote{masaka}
  \pos{Noun}  \sense{
    \definition{single stem plant with inedible fruit}}}

\entry{Masam}{\headword{Masam}
  \pos{Proper noun}  \sense{
    \definition{Bitur language}}}

\entry{=masäm}{\headword{=masäm}
  \variantof{Dialectal Variant of \vartext{=mattäm}}}

\entry{masamasar}{\headword{masamasar}  \pronnote{masamasar}
  \pos{Kinship noun}  \sense{\textbf{1.} 
    \definition{grandparents}    \example{
      \exsen{Ngämo masamasar a ngämo pate gognegnän Ende eka walle de.}      \extran{My grandparents spoke to me in Ende.}}}  \sense{\textbf{2.} 
    \definition{forefathers, ancestors}    \example{
      \exsen{God bo eka da angde Ende tän e gozenän, lla da Ende tän bo masamasar ngasnen de däpnaeaemeyo.}      \extran{When God's word entered the Ende tribe, people turned away from the ways of the Ende tribe's ancestors.}}}}

\entry{masar}{\headword{masar}  \lexeme{masar}  \note{Family ties}  \pronnote{masar}
  \pos{Kinship noun}  \sense{\textbf{1.} 
    \definition{grandfather (one's parent's father; reciprocal); ancestor}    \note{Family ties}    \example{
      \exsen{Masar, bablle ngoe a ddone dag.}      \extran{Grandfather, you have no teeth.}}    \example{
}}  \sense{\textbf{2.} 
    \definition{grandchild (man's child's child; reciprocal)}    \note{Family ties}}  \sense{\textbf{3.} 
    \definition{father-in-law (woman's husband's father; reciprocal)}    \note{Family ties}}  \sense{\textbf{4.} 
    \definition{daughter-in-law (man's son's wife; reciprocal)}    \note{Family ties}}  \sense{\textbf{5.} 
    \definition{uncle-in-law (woman's husband's mother's younger brother; reciprocal)}    \note{Family ties}}  \sense{\textbf{6.} 
    \definition{niece-in-law (man's elder sister's son's wife; reciprocal)}    \note{Family ties}}  \sense{\textbf{7.} 
    \definition{ancestral, old (of one's grandparents' time)}    \note{Family ties}    \example{
      \exsen{masar täräp me}      \extran{in the olden days}}    \example{
      \exsen{masar eka}      \extran{ancestral language}}}}

\entry{mäsdae}{\headword{mäsdae}
  \variantof{Dialectal Variant of \vartext{misdae}}  \pronnote{məsdae}}

\entry{mäsde}{\headword{mäsde}
  \variantof{Dialectal Variant of \vartext{misdae}}}

\entry{mäse}{\headword{mäse}  \pronnote{məse}
  \pos{TAM particle}  \sense{\textbf{1.} 
    \definition{imminent particle (indicates that something about to take place)}    \example{
      \exsen{Ngäna ako mäse tutu de dängkälne, gänyaollemae ikop dige sɨmell da angttägan.}      \extran{As I was about to climb the mountain, I saw a pig coming towards me.}}}  \sense{\textbf{2.} 
    \definition{conative particle (indicates an attempt; to try to do something)}    \example{
      \exsen{Ngäna ngängälläb e mäse däga.}      \extran{I tried to lift it.}}}  \sense{\textbf{3.} 
    \definition{not yet full, unsatisfied}    \example{
      \exsen{Ngämlle kame ako otät de nasi, ngämo käm a mäse agan.}      \extran{Give me more food; my stomach isn't full yet.}}}}

\entry{=masem}{\headword{=masem}
  \variantof{Unspecified Variant of \vartext{=mattäm}}}

\entry{mäsemäse}{\headword{mäsemäse}  \pronnote{məseməse}
  \pos{Adverb}  \sense{\textbf{1.} 
    \definition{naked}    \example{
      \exsen{Dadabi, mäsemäse llɨg a dongkoenmällaemneyo.}      \extran{The boys chased them naked.}}}  \sense{\textbf{2.} 
    \definition{empty-handed}    \example{
      \exsen{Mäsemäse ma gongttägeya.}      \extran{We arrived home empty-handed.}}}}

\entry{=masen}{\headword{=masen}
  \variantof{Unspecified Variant of \vartext{=mattäm}}}

\entry{Masingara}{\headword{Masingara}  \pronnote{masiŋara}
  \pos{Proper noun}  \sense{
    \definition{Masingara (Bine-speaking village in Oriomo-Bituri Rural LLG)}    \example{
}    \example{
}    \example{
}}}

\entry{Masta}{\headword{Masta}
  \pos{Proper noun}  \sense{
    \definition{female personal name}}}

\entry{mäta}{\headword{mäta}  \pronnote{məta}
  \pos{Noun}  \sense{
    \definition{type of tree with a red stem, young green leaves, and bark used as rope when old; big trees are used to light fire; liquid is extracted and drunk to treat cough}}}

\entry{Mata}{\headword{Mata}  \pronnote{mata}
  \pos{Proper noun}  \sense{
    \definition{Mata (in Morehead Rural LLG)}    \example{
}    \example{
}    \example{
}}}

\entry{mäta mamlla}{\headword{mäta mamlla}  \note{Gardens}  \pronnote{məta mamɽa}
  \pos{Noun}  \sense{
    \definition{rope made from mäta bark (used to secure fencing)}    \note{Gardens}    \example{
}    \example{
}    \example{
}}}

\entry{matamata}{\headword{matamata}  \note{Trees}  \pronnote{matamata}
  \pos{Noun}  \sense{
    \definition{type of tree that grows along swamps with white and blue flowers and edible fruit (black outside, red inside) that ripen during January and February}    \note{Trees}    \example{
}    \example{
}    \example{
}}}

\entry{mätämätär}{\headword{mätämätär}  \note{Trees}  \pronnote{mətəmətər}
  \pos{Noun}  \sense{
    \definition{type of tree}    \note{Trees}    \example{
}    \example{
}    \example{
}}}

\entry{Mätär}{\headword{Mätär}  \note{Places around Limol}  \pronnote{mətər}
  \pos{Proper noun}  \sense{
    \definition{Mätär (toponym)}    \note{Places around Limol}    \example{
}    \example{
}    \example{
}}}

\entry{Matär}{\headword{Matär}
  \pos{Proper noun}  \sense{
    \definition{male personal name}}}

\entry{mätar}{\headword{mätar}
  \variantof{Unspecified Variant of \vartext{mätaru}}}

\entry{mätar onyang}{\headword{mätar onyang}  \note{Government}  \pronnote{mətar oɲaŋ}
  \pos{Noun}  \sense{
    \definition{peace officer}    \note{Government}    \example{
}    \example{
}    \example{
}}}

\entry{mätar onyang}{\headword{mätar onyang}  \note{Government}  \pronnote{mətar oɲaŋ}
  \pos{Noun}  \sense{
    \definition{type of tree}    \note{Government}}}

\entry{mätaru}{\headword{mätaru}  \note{Conflict/war}  \pronnote{mətaru}
  \pos{Modifier}  \sense{
    \definition{calm, peaceful, quiet}    \note{Conflict/war}    \example{
      \exsen{Tämamae ttängäm a mätaru abal gogon.}      \extran{The whole place became very calm.}}}
  \variants{\Unspecified Variant \vartext{mätar}}}

\entry{Mataru}{\headword{Mataru}
  \pos{Proper noun}  \sense{
    \definition{personal name}}}

\entry{mätemäte}{\headword{mätemäte}  \note{Names of bananas}  \pronnote{məteməte}
  \pos{Noun}  \sense{
    \definition{type of introduced banana}    \note{Names of bananas}    \example{
}    \example{
}    \example{
}}}

\entry{Maten}{\headword{Maten}
  \pos{Proper noun}  \sense{
    \definition{male personal name}}
  \variants{\SpellingVariant \vartext{Martin}}}

\entry{Mathias}{\headword{Mathias}
  \variantof{SpellingVariant of \vartext{Matias}}  \pronnote{maθias}}

\entry{Mathilda}{\headword{Mathilda}
  \pos{Proper noun}  \sense{
    \definition{female personal name}}
  \variants{\SpellingVariant \vartext{Matilda}}}

\entry{Matias}{\headword{Matias}
  \pos{Proper noun}  \sense{
    \definition{male personal name}}
  \variants{\SpellingVariant \vartext{Mathias, Maties}}}

\entry{Maties}{\headword{Maties}
  \variantof{SpellingVariant of \vartext{Matias}}}

\entry{Matilda}{\headword{Matilda}
  \variantof{SpellingVariant of \vartext{Mathilda}}}

\entry{mätka}{\headword{mätka}  \lexeme{mätka}  \pronnote{mətka}
  \pos{Noun}  \sense{
    \definition{type of tall grass with a single stem and bunches of up to 12 small, edible fruit that grow near the base}}}

\entry{mätkakallamatt}{\headword{mätkakallamatt}  \note{Snakes}  \pronnote{mətkakaɽamaʈ͡ʂ}
  \pos{Noun}  \sense{
    \definition{type of snake}    \note{Snakes}    \example{
}    \example{
}    \example{
}}}

\entry{mätkin}{\headword{mätkin}  \note{Numbers}  \pronnote{mətkin}
  \pos{Noun}  \sense{\textbf{1.} 
    \definition{ring finger}    \note{Numbers}    \example{
}}  \sense{\textbf{2.} 
    \definition{two (lit. ring finger; body counting numeral)}    \note{Numbers}    \example{
}    \example{
}    \example{
}}}

\entry{mätta}{\headword{mätta}  \lexeme{mätta}  \pronnote{mä'tta}
  \pos{Noun}  \sense{
    \definition{lesser yam}    \scientificname{Dioscorea esculenta}    \example{
      \exsen{Mätta da nyäng de bodo dägagän.}      \extran{Yams filled the basket.}}}}

\entry{matta}{\headword{matta}  \pronnote{maʈ͡ʂa}
  \pos{Noun}  \sense{\textbf{1.} 
    \definition{shoulder}}  \sense{\textbf{2.} 
    \definition{to shoulder, carry on one's shoulder}    \example{
      \exsen{Angde llo täraem a gottamänän, ngäna dibaballe matta däganegne ma we.}      \extran{When the tree cutting was finished, I then shouldered them to the house.}}}  \sense{\textbf{3.} 
    \definition{eight (lit. shoulder; body counting numeral)}}}

\entry{mätta pirpir}{\headword{mätta pirpir}  \pronnote{məʈ͡ʂa pirpir}
  \pos{Noun}  \sense{
    \definition{rainbow bee-eater}    \scientificname{Merops ornatus}    \example{
}    \example{
}    \example{
}}
  \variants{\Unspecified Variant \vartext{mätta pirpirik}}}

\entry{mätta pirpirik}{\headword{mätta pirpirik}
  \variantof{Unspecified Variant of \vartext{mätta pirpir}}  \pronnote{məʈ͡ʂa pirpirik}}

\entry{mättae}{\headword{mättae}
  \pos{Transitive verb}  \sense{
    \definition{to threaten, raise fists}    \example{
      \exsen{Ngämo mälla bom mättaenen nängkaman.}      \extran{I started to threaten my wife.}}    \example{
      \exsen{Ngäna obom damttaene.}      \extran{I was threatening him.}}}}

\entry{=mattäm}{\headword{=mattäm}
  \pos{Nominal enclitic}  \sense{
    \definition{ablative case clitic; from}}
  \variants{\Unspecified Variant \vartext{=masen, =masem}}
  \variants{\Dialectal Variant \vartext{=masäm}}}

\entry{mattgal}{\headword{mattgal}  \pronnote{maʈ͡ʂɡal}
  \pos{Transitive verb}  \sense{
    \definition{to put in fire}    \example{
      \exsen{Yu de daudewän a otät de yu wi dämattgalän.}      \extran{He lit a fire and put the food in the fire.}}}}

\entry{Matthew}{\headword{Matthew}
  \variantof{SpellingVariant of \vartext{Metyu}}}

\entry{mättmätt}{\headword{mättmätt}  \pronnote{məʈ͡ʂməʈ͡ʂ}
  \pos{Intransitive verb}  \sense{\textbf{1.} 
    \definition{to wear, dress oneself, put on}    \example{
      \exsen{Bogo banggu mättmätt allan.}      \extran{He is wearing a headdress.}}    \example{
      \exsen{Mägda sisor kaptte amättan.}      \extran{Mother put on new clothes.}}}  \sense{\textbf{2.} 
    \definition{to dress}    \example{
      \exsen{Ngäna obom mättmätt eran.}      \extran{I am dressing her.}}}}

\entry{mattmett}{\headword{mattmett}
  \pos{Noun}  \sense{
    \definition{plain, flat area}}}

\entry{mattmett}{\headword{mattmett}
  \pos{Transitive verb}  \sense{
    \definition{to put in oven}    \example{
      \exsen{Sɨmell gullbe de ddob sana peyang dazgiaebeya, ada dämetteyo ttägäll ik i.}      \extran{We wrapped the boar with some sago and put it in the mumu.}}    \example{
}}}

\entry{matu}{\headword{matu}  \pronnote{matu}
  \pos{Noun}  \sense{\textbf{1.} 
    \definition{ground}    \example{
      \exsen{Baet abo gongetamän matu we.}      \extran{Cuscus came down to the ground.}}}  \sense{\textbf{2.} 
    \definition{lower part}    \example{
      \exsen{Matu me kollko ddage me dedme gotogolän.}      \extran{He hid there on the lowest branch.}}}}

\entry{Mavis}{\headword{Mavis}  \pronnote{mavis}
  \pos{Proper noun}  \sense{
    \definition{female personal name}}}

\entry{mawa}{\headword{mawa}
  \pos{Noun}  \sense{
    \definition{magic}    \example{
      \exsen{mawa lla}      \extran{fairy, elf}}}}

\entry{Max}{\headword{Max}
  \variantof{SpellingVariant of \vartext{Meks}}  \pronnote{maks}}

\entry{maza}{\headword{maza}
  \pos{Noun}  \sense{
    \definition{reef}}}

\entry{mäzi}{\headword{mäzi}
  \variantof{Dialectal Variant of \vartext{mɨnyi}}}

\entry{maa}{\headword{maa}
  \variantof{SpellingVariant of \vartext{ma}}}

\entry{maaduma}{\headword{maaduma}
  \variantof{SpellingVariant of \vartext{maduma}}}

\entry{me}{\headword{me}  \pronnote{me}
  \pos{Noun}  \sense{
    \definition{rainbow}}}

\entry{me}{\headword{me}
  \variantof{freevarlab of \vartext{English loanword}}}

\entry{me}{\headword{me}
  \pos{Noun}  \sense{
}}

\entry{=me}{\headword{=me}  \pronnote{me}
  \pos{Nominal enclitic}  \sense{\textbf{1.} 
    \definition{locative case clitic; in, at, on}    \example{
      \exsen{Ngäna ngämo pällämpälläm palle de ngämo mittapa me nowattällan.}      \extran{I put my white clay on my cheeks.}}    \example{
      \exsen{bem ulle me}      \extran{in the big sea}}    \example{
      \exsen{däbe täräp me}      \extran{at that time}}}  \sense{\textbf{2.} 
    \definition{during, while}    \example{
      \exsen{Bogo iang me bandra allan.}      \extran{He was singing while weaving.}}}
  \variants{\SpellingVariant \vartext{=mee}}}

\entry{=meae}{\headword{=meae}
  \variantof{Dialectal Variant of \vartext{=mae}}  \pronnote{=meae}
  \variants{\Dialectal Variant \vartext{mae}}}

\entry{Mecklyn}{\headword{Mecklyn}  \pronnote{mɛklɪn}
  \pos{Proper noun}  \sense{
    \definition{female personal name}}}

\entry{medä}{\headword{medä}
  \variantof{Unspecified Variant of \vartext{mäda}}  \pronnote{medə}}

\entry{meda}{\headword{meda}
  \variantof{Dialectal Variant of \vartext{mäda}}}

\entry{Medang}{\headword{Medang}
  \pos{Proper noun}  \sense{
    \definition{Medang (toponym)}}}

\entry{medol}{\headword{medol}
  \pos{Noun}  \sense{
    \definition{medal}}}

\entry{Megam}{\headword{Megam}  \pronnote{meɡam}
  \pos{Proper noun}  \sense{
    \definition{male personal name}}}

\entry{Megi}{\headword{Megi}  \pronnote{meɡi}
  \pos{Proper noun}  \sense{
    \definition{female personal name}}
  \variants{\SpellingVariant \vartext{Maggie}}}

\entry{mekae}{\headword{mekae}  \note{Trees}  \pronnote{mekae}
  \pos{Noun}  \sense{
    \definition{type of tree with white flowers and edible nuts that children eat}    \note{Trees}    \example{
}    \example{
}    \example{
}}}

\entry{mekewa}{\headword{mekewa}  \note{Names of bananas}  \pronnote{mekewa}
  \pos{Noun}  \sense{
    \definition{type of introduced banana}    \note{Names of bananas}    \example{
}    \example{
}    \example{
}}}

\entry{Meklin}{\headword{Meklin}
  \pos{Proper noun}  \sense{
    \definition{female personal name}}}

\entry{Meks}{\headword{Meks}
  \pos{Proper noun}  \sense{
    \definition{male personal name}}
  \variants{\SpellingVariant \vartext{Max}}}

\entry{melem}{\headword{melem}  \pronnote{melem}
  \pos{Noun}  \sense{
    \definition{work}    \example{
      \exsen{Oba melem dättemäneyo a gongoseyo ma we.}      \extran{They finished their work and returned home.}}    \example{
      \exsen{Obo peyang mɨnyi melem bog.}      \extran{I will work with her.}}}}

\entry{melem duag}{\headword{melem duag}  \note{Stages of life}  \pronnote{melem duaɡ}
  \pos{Noun}  \sense{
    \definition{stage of life when you're able to do all the work}    \note{Stages of life}    \example{
}    \example{
}    \example{
}}}

\entry{melemang}{\headword{melemang}  \note{good characteristics of people}  \pronnote{melemaŋ}
  \pos{Modifier}  \sense{\textbf{1.} 
    \definition{hardworking}    \note{good characteristics of people}    \example{
      \exsen{Bogo ddobae melemang daeya.}      \extran{He was very hardworking.}}    \example{
}    \example{
}}  \sense{\textbf{2.} 
    \definition{servant, laborer}    \note{good characteristics of people}    \example{
      \exsen{Ttängäm mondrog lla da melemang de dällädeyo a gagagäll duwebeyo.}      \extran{The garden workers grabbed the servant and beat him badly.}}}}

\entry{Meliye}{\headword{Meliye}  \pronnote{melije}
  \pos{Proper noun}  \sense{
    \definition{Meliye (toponym)}    \example{
}    \example{
}    \example{
}}}

\entry{Melvin}{\headword{Melvin}  \pronnote{melvin}
  \pos{Proper noun}  \sense{
    \definition{male personal name}}}

\entry{mem}{\headword{mem}
  \variantof{Dialectal Variant of \vartext{mam}}  \pronnote{mem}}

\entry{memba}{\headword{memba}
  \pos{Noun}  \sense{
    \definition{member}}}

\entry{memram}{\headword{memram}  \pronnote{memram}
  \pos{Noun}  \sense{
    \definition{sweat}    \example{
      \exsen{Memram a dägagän.}      \extran{He got sweaty (lit. sweat got him).}}}}

\entry{mena}{\headword{mena}
  \pos{Transitive verb}  \sense{
    \definition{to scorch}    \example{
      \exsen{Yäbäd ttänttäm a sisor koko de dämenanegän.}      \extran{The heat of the sun scorched the young shoots.}}}}

\entry{menae}{\headword{menae}  \pronnote{menae}
  \pos{Locational}  \sense{
    \definition{side, edge; vicinity}    \example{
      \exsen{walle menae me}      \extran{by the water}}    \example{
      \exsen{nyongo menae me}      \extran{on the roadside}}}}

\entry{mend}{\headword{mend}  \note{Birds}  \pronnote{mend}
  \pos{Noun}  \sense{
    \definition{type of bird}    \note{Birds}    \example{
}    \example{
}    \example{
}}}

\entry{mengkmeny}{\headword{mengkmeny}
  \pos{Modifier}  \sense{
    \definition{aching}}}

\entry{menizment}{\headword{menizment}
  \pos{Noun}  \sense{
    \definition{management}}}

\entry{menttäg}{\headword{menttäg}  \pronnote{menʈ͡ʂəɡ}
  \pos{Intransitive verb}  \sense{
    \definition{to shoulder, carry one one's shoulders or head}    \example{
      \exsen{Nyäng goimänttmenyeyo deyarneyo abo ma we.}      \extran{They put the baskets on their shoulders and headed home.}}}}

\entry{menttlog}{\headword{menttlog}
  \pos{Transitive verb}  \sense{
    \definition{to agree}}}

\entry{=meny}{\headword{=meny}  \pronnote{meɲ}
  \pos{Nominal enclitic}  \sense{
    \definition{privative clitic; without}    \example{
      \exsen{mängall, mängallmeny; mu, miny}      \extran{strength, weak; payment, in debt}}}}

\entry{mer}{\headword{mer}  \pronnote{mer}
  \pos{Noun}  \sense{
    \definition{type of spear}}}

\entry{mer}{\headword{mer}  \pronnote{mer}
  \pos{Modifier}  \sense{\textbf{1.} 
    \definition{good}    \example{
      \exsen{Ngäna ttongo mer kollba de näddägan.}      \extran{I ate one of the good fish.}}}  \sense{\textbf{2.} 
    \definition{holy}    \example{
      \exsen{Mer Anyke}      \extran{Holy Spirit}}}  \sense{\textbf{3.} 
    \definition{to bless}    \example{
      \exsen{Yesu llɨg kälekäle de mer dägnegän.}      \extran{Jesus blessed the children.}}}}

\entry{mer ag}{\headword{mer ag}  \lexeme{mer ag}  \pronnote{mer aɡ}
  \pos{Phrase}  \sense{
    \definition{good morning}}}

\entry{mer awi}{\headword{mer awi}  \note{?}  \pronnote{mer awi}
  \pos{Phrase}  \sense{
    \definition{good evening}    \note{?}    \example{
}    \example{
}    \example{
}}}

\entry{mer iddob}{\headword{mer iddob}  \note{?}  \pronnote{mer iɖ͡ʐob}
  \pos{Phrase}  \sense{
    \definition{good night}    \note{?}    \example{
}    \example{
}    \example{
}}}

\entry{mer toto}{\headword{mer toto}  \note{?}  \pronnote{mer toto}
  \pos{Phrase}  \sense{
    \definition{good afternoon}    \note{?}    \example{
}    \example{
}    \example{
}}}

\entry{mer ttoen gagällag lla}{\headword{mer ttoen gagällag lla}  \note{Bad characteristics of people}  \pronnote{mer ʈ͡ʂoen ɡaɡəɽaɡ ɽa}
  \pos{Noun}  \sense{
    \definition{messy person}    \note{Bad characteristics of people}    \example{
}    \example{
}    \example{
}}}

\entry{mer ttoenang lla}{\headword{mer ttoenang lla}  \note{good characteristics of people}  \pronnote{mer ʈ͡ʂoenaŋ ɽa}
  \pos{Noun}  \sense{
    \definition{good character}    \note{good characteristics of people}    \example{
}    \example{
}    \example{
}}}

\entry{Meragag}{\headword{Meragag}
  \pos{Proper noun}  \sense{
    \definition{Meragag (toponym)}}}

\entry{Meramerall}{\headword{Meramerall}
  \pos{Proper noun}  \sense{
    \definition{Meramerall (toponym)}}}

\entry{meresen}{\headword{meresen}
  \variantof{SpellingVariant of \vartext{meresin}}}

\entry{meresin}{\headword{meresin}  \pronnote{meresin}
  \pos{Noun}  \sense{
    \definition{medicine}    \example{
}}
  \variants{\SpellingVariant \vartext{meresen}}}

\entry{Meri}{\headword{Meri}
  \pos{Proper noun}  \sense{
    \definition{female personal name}}
  \variants{\SpellingVariant \vartext{Mary}}}

\entry{Merian}{\headword{Merian}
  \pos{Proper noun}  \sense{
    \definition{female personal name}}
  \variants{\SpellingVariant \vartext{Merien, Maryanne}}}

\entry{Merien}{\headword{Merien}
  \variantof{SpellingVariant of \vartext{Merian}}}

\entry{mermer}{\headword{mermer}  \pronnote{mermer}
  \pos{Verb}  \sense{
    \definition{to shrink, shrivel}    \example{
      \exsen{Ge mätta da mermer allan.}      \extran{This yam is shrinking.}}}}

\entry{mermer}{\headword{mermer}  \pronnote{mermer}
  \pos{Adverb}  \sense{
    \definition{properly, nicely, well}    \example{
      \exsen{Ngäna ngämo toboll de mermer met däganeg.}      \extran{I sharpened my spears well.}}}}

\entry{Meroka}{\headword{Meroka}
  \pos{Proper noun}  \sense{
    \definition{female personal name}}}

\entry{Merol}{\headword{Merol}  \pronnote{merol}
  \pos{Proper noun}  \sense{
    \definition{female personal name}}}

\entry{Mesa}{\headword{Mesa}  \pronnote{mesa}
  \pos{Proper noun}  \sense{
    \definition{male personal name}}}

\entry{met}{\headword{met}  \pronnote{met}
  \pos{Transitive coverb}  \sense{
    \definition{to sharpen}    \example{
      \exsen{Ttongo giri de met dägageyo.}      \extran{They sharpened a knife.}}}}

\entry{metar}{\headword{metar}  \lexeme{metar}  \pronnote{metar}
  \pos{Noun}  \sense{
    \definition{type of long yam with a white interior, white skin, and hairs}}}

\entry{Metiyu}{\headword{Metiyu}
  \variantof{SpellingVariant of \vartext{Metyu}}}

\entry{metmäll}{\headword{metmäll}
  \pos{Transitive verb}  \sense{
    \definition{to beat, flog, hit}}}

\entry{Metyu}{\headword{Metyu}
  \pos{Proper noun}  \sense{
    \definition{male personal name}}
  \variants{\SpellingVariant \vartext{Metiyu, Matthew}}}

\entry{Mewato}{\headword{Mewato}
  \pos{Proper noun}  \sense{
    \definition{female personal name}}}

\entry{meyang}{\headword{meyang}  \lexeme{meyang}  \pronnote{mejaŋ}
  \pos{Kinship noun}  \sense{\textbf{1.} 
    \definition{uncle (one's father's younger brother)}    \example{
      \exsen{Meyang daeya tatu ma nallan.}      \extran{Uncle went to wash.}}}  \sense{\textbf{2.} 
    \definition{brother-in-law (woman's husband's younger brother)}}}

\entry{=mee}{\headword{=mee}
  \variantof{SpellingVariant of \vartext{=me}}}

\entry{Michael}{\headword{Michael}  \pronnote{maikəl}
  \pos{Proper noun}  \sense{
    \definition{male personal name}}}

\entry{Michaelyn}{\headword{Michaelyn}  \pronnote{maikəlɪn}
  \pos{Proper noun}  \sense{
    \definition{female personal name}}}

\entry{Michelle}{\headword{Michelle}  \pronnote{miʃel}
  \pos{Proper noun}  \sense{
    \definition{female personal name}}}

\entry{midd}{\headword{midd}  \pronnote{miɖ͡ʐ}
  \pos{Noun}  \sense{\textbf{1.} 
    \definition{meat}    \example{
      \exsen{sɨmell midd}      \extran{pork}}    \example{
}    \example{
}}  \sense{\textbf{2.} 
    \definition{meaning; core, essence}    \example{
      \exsen{eka bo midd}      \extran{a word's meaning}}}}

\entry{midi}{\headword{midi}
  \pos{Noun}  \sense{
}}

\entry{Migul}{\headword{Migul}
  \pos{Proper noun}  \sense{
    \definition{male personal name}}}

\entry{mik}{\headword{mik}  \lexeme{mik}  \pronnote{mik}
  \pos{Noun}  \sense{
    \definition{widow, widower}    \example{
      \exsen{Ttongo llɨgmeny mik da oblle damändän.}      \extran{A childless widow raised him.}}}}

\entry{mikmik}{\headword{mikmik}
  \pos{Noun}  \sense{
    \definition{nonsingular form of mik}}}

\entry{mikmik}{\headword{mikmik}
  \pos{Verb}  \sense{
    \definition{to suck}}}

\entry{miks}{\headword{miks}
  \pos{Noun}  \sense{\textbf{1.} 
    \definition{mix, mixture}    \example{
      \exsen{Ngäma eka da miks ingoll gogon.}      \extran{Our (excl.) language has become like a mix.}}}  \sense{\textbf{2.} 
    \definition{to mix}}}

\entry{mikutt}{\headword{mikutt}  \note{Things a person can do}  \pronnote{mikuʈ͡ʂ}
  \pos{Modifier}  \sense{
    \definition{angry, mad}    \note{Things a person can do}    \example{
      \exsen{Obo mikutt a ttäntämang e gopnaeän.}      \extran{His anger made him hot.}}    \example{
      \exsen{Da ttongo lla da bäne za de gämäll bogon, bongo mɨnyi mikutt ag.}      \extran{If someone steals your thing, you will be angry.}}    \example{
      \exsen{Yae mikutt allan.}      \extran{Mother is mad.}}    \example{
      \exsen{Bogo mikutt gogon.}      \extran{He got angry.}}}}

\entry{mikuttang}{\headword{mikuttang}  \note{Verbs}  \pronnote{mikuʈ͡ʂaŋ}
  \pos{Modifier}  \sense{\textbf{1.} 
    \definition{angry, short-tempered}    \note{Verbs}    \example{
      \exsen{Bogo era ddobae mikuttang lla dan.}      \extran{He's a very angry person.}}}  \sense{\textbf{2.} 
    \definition{to hate}    \note{Verbs}    \example{
      \exsen{Mikuttang eran.}      \extran{He hates him.}}    \example{
}    \example{
}    \example{
}}}

\entry{mikuttangae}{\headword{mikuttangae}  \pronnote{mikuʈ͡ʂaŋae}
  \pos{Adverb}  \sense{
    \definition{angrily}    \example{
      \exsen{Nga bogo mikuttangae dagernän.}      \extran{He stayed angry for a long time.}}}}

\entry{mikuttmeny}{\headword{mikuttmeny}  \pronnote{mikuʈ͡ʂmeɲ}
  \pos{Modifier}  \sense{
    \definition{calm}    \example{
      \exsen{Bogo era mikuttmeny lla dan.}      \extran{He's a calm fellow.}}}}

\entry{mimi}{\headword{mimi}  \pronnote{mimi}
  \pos{Noun}  \sense{
    \definition{pig (hunting word)}}}

\entry{mimidämäll}{\headword{mimidämäll}  \note{Snakes}  \pronnote{mimidəməɽ}
  \pos{Noun}  \sense{
    \definition{type of snake}    \note{Snakes}    \example{
}    \example{
}    \example{
}}}

\entry{minggore manggo}{\headword{minggore manggo}  \note{Trees}  \pronnote{miŋɡore maŋɡo}
  \pos{Noun}  \sense{
    \definition{type of tree}    \note{Trees}    \example{
}    \example{
}    \example{
}}}

\entry{Mingkällbun}{\headword{Mingkällbun}  \note{Places around Limol}  \pronnote{miŋkəɽbun}
  \pos{Proper noun}  \sense{
    \definition{Mingkällbun (toponym)}    \note{Places around Limol}    \example{
}    \example{
}    \example{
}}}

\entry{ministri}{\headword{ministri}
  \pos{Noun}  \sense{
    \definition{ministry}    \example{
}}}

\entry{Minkäm}{\headword{Minkäm}
  \pos{Proper noun}  \sense{
    \definition{Minkam (toponym)}}
  \variants{\Unspecified Variant \vartext{Minkom}}}

\entry{Minkom}{\headword{Minkom}
  \variantof{Unspecified Variant of \vartext{Minkäm}}}

\entry{Minkomminkomang}{\headword{Minkomminkomang}
  \pos{Proper noun}  \sense{
    \definition{Minkomminkomang (toponym)}}}

\entry{Minong}{\headword{Minong}  \pronnote{minoŋ}
  \pos{Proper noun}  \sense{
    \definition{male personal name}}}

\entry{mintor}{\headword{mintor}  \note{Trees}  \pronnote{mintor}
  \pos{Noun}  \sense{
    \definition{type of tree with yellow flowers that bloom in June and July and a root is used as a kind of hockey stick}    \note{Trees}    \example{
}    \example{
}    \example{
}}}

\entry{miny}{\headword{miny}  \pronnote{miɲ}
  \pos{Noun}  \sense{
    \definition{ant egg/larvae (small)}}
  \variants{\Unspecified Variant \vartext{minyminy}}}

\entry{minyminy}{\headword{minyminy}
  \variantof{Unspecified Variant of \vartext{miny}}}

\entry{mipdab}{\headword{mipdab}  \pronnote{mipdab}
  \pos{Transitive coverb}  \sense{
    \definition{to accommodate; offer food to a visitor}    \example{
      \exsen{Ngänäm Malläm me mipdab dageyo.}      \extran{They accomodated me in Malam.}}}}

\entry{Miriang}{\headword{Miriang}
  \pos{Proper noun}  \sense{
    \definition{male personal name}}}

\entry{miriwa}{\headword{miriwa}  \note{Type of pandanus}  \pronnote{miriwa}
  \pos{Noun}  \sense{
    \definition{type of pandanus}    \note{Type of pandanus}    \example{
}    \example{
}    \example{
}}}

\entry{mirmir}{\headword{mirmir}  \pronnote{mirmir}
  \pos{Transitive verb}  \sense{
    \definition{to lick}}}

\entry{miroli}{\headword{miroli}  \pronnote{miroli}
  \pos{Noun}  \sense{
    \definition{black-capped lory}    \scientificname{Lorius lory erythrothorax}    \example{
}    \example{
}    \example{
}}}

\entry{misdae}{\headword{misdae}  \pronnote{misdae}
  \pos{Adverb}  \sense{
    \definition{just, simply}    \example{
      \exsen{Ngäna misdae ada dagerne.}      \extran{I was just staying like this.}}    \example{
      \exsen{Ttomoe a mudan, misdae känyer ag.}      \extran{Don't you worry, just be quiet.}}}
  \variants{\Fast Speech Variant \vartext{misde}}
  \variants{\SpellingVariant \vartext{mistae}}
  \variants{\Dialectal Variant \vartext{mäsde, mäsdae}}}

\entry{misde}{\headword{misde}
  \variantof{Fast Speech Variant of \vartext{misdae}}}

\entry{mise}{\headword{mise}  \pronnote{mise}
  \pos{Noun}  \sense{
    \definition{common cicadabird}    \scientificname{Edolisoma tenuirostre}    \example{
}    \example{
}    \example{
}}}

\entry{misituryam}{\headword{misituryam}  \lexeme{misituryam}  \pronnote{misiturjam}
  \pos{Noun}  \sense{
    \definition{type of long yam with a white or purple interior}}}

\entry{mislok}{\headword{mislok}  \note{Names of bananas}  \pronnote{mislok}
  \pos{Noun}  \sense{
    \definition{type of introduced banana}    \note{Names of bananas}    \example{
}    \example{
}    \example{
}}}

\entry{mismis}{\headword{mismis}  \note{Trees}  \pronnote{mismis}
  \pos{Noun}  \sense{
    \definition{type of tree used for firewood}    \note{Trees}    \example{
}    \example{
}    \example{
}}}

\entry{Misseilene}{\headword{Misseilene}  \pronnote{misseilene}
  \pos{Proper noun}  \sense{
    \definition{female personal name}}}

\entry{mista}{\headword{mista}
  \pos{Noun}  \sense{
    \definition{mister}}}

\entry{mistae}{\headword{mistae}
  \variantof{SpellingVariant of \vartext{misdae}}}

\entry{mit~}{\headword{mit~}
  \variantof{freevarlab of \vartext{RED~}}}

\entry{mit}{\headword{mit}  \pronnote{mit}
  \pos{Noun}  \sense{\textbf{1.} 
    \definition{base (of a plant)}    \example{
      \exsen{Ubi llo mit mi dagaeya.}      \extran{They were at the base of the tree.}}}  \sense{\textbf{2.} 
    \definition{reason, sake}    \example{
      \exsen{Mit a gänyan.}      \extran{Here is the reason.}}    \example{
      \exsen{Ttongo lla da säremang ma me daeya kuddäll e lla gäz mit me.}      \extran{A man was in prison because he beat someone to death.}}    \example{
      \exsen{Iba mit me däpleon.}      \extran{He died for our (incl.) sake.}}}  \sense{\textbf{3.} 
    \definition{origin, source}    \example{
      \exsen{Ge Ende eka da bo mit eka dan.}      \extran{This is the origin story of the Ende language.}}}  \sense{\textbf{4.} 
}}

\entry{mitamitang}{\headword{mitamitang}  \pronnote{mitamitaŋ}
  \pos{Transitive coverb}  \sense{
    \definition{to blame}    \example{
      \exsen{Ngänäm mitamitang dagän.}      \extran{They were blaming me.}}}}

\entry{mitin}{\headword{mitin}
  \variantof{SpellingVariant of \vartext{miting}}}

\entry{miting}{\headword{miting}
  \pos{Noun}  \sense{
    \definition{meeting}}
  \variants{\Unspecified Variant \vartext{mittin}}
  \variants{\SpellingVariant \vartext{mitin}}}

\entry{mitmit}{\headword{mitmit}  \note{Verbs}  \pronnote{mitmit}
  \pos{Noun}  \sense{
    \definition{blunt axe}    \note{Verbs}}}

\entry{mitmit}{\headword{mitmit}  \note{Verbs}  \pronnote{mitmit}
  \pos{Transitive coverb}  \sense{
    \definition{to miss, long for}    \note{Verbs}    \example{
      \exsen{Lla da obom ddobae ddobae mitmit nägagallo.}      \extran{They were thinking of him when he left.}}    \example{
      \exsen{Da medäda erowe gopällttänalle llɨg da mɨnyi mitmit bäganän ngänaeka peyang.}      \extran{‎If father set off somewhere, the children will miss him in tears.}}    \example{
}}}

\entry{mitoem}{\headword{mitoem}  \note{Mushrooms}  \pronnote{mitoem}
  \pos{Noun}  \sense{
    \definition{type of mushroom}    \note{Mushrooms}    \example{
}    \example{
}    \example{
}}}

\entry{mittäpa}{\headword{mittäpa}
  \variantof{Unspecified Variant of \vartext{mittapa}}}

\entry{mittapa}{\headword{mittapa}  \note{Parts of the Body}  \pronnote{miʈ͡ʂapa}
  \pos{Noun}  \sense{
    \definition{extremities of face (i.e. temples, cheek, chin)}    \note{Parts of the Body}    \example{
}    \example{
}    \example{
}}
  \variants{\Unspecified Variant \vartext{mittäpa}}}

\entry{mittin}{\headword{mittin}
  \variantof{Unspecified Variant of \vartext{miting}}  \pronnote{miʈ͡ʂin}}

\entry{miwiwi}{\headword{miwiwi}  \pronnote{miwiwi}
  \pos{Noun}  \sense{\textbf{1.} 
    \definition{dabbling duck (pacific black duck, grey teal)}    \scientificname{Anas superciliosa superciliosa, A. gracilis}    \example{
}    \example{
}    \example{
}}  \sense{\textbf{2.} 
    \definition{spotted whistling duck}    \scientificname{Dendrocygna guttata}}}

\entry{mizi}{\headword{mizi}  \pronnote{mizi}
  \pos{Manner adverb}  \sense{
    \definition{usually}    \example{
      \exsen{Mizi mällaeyaba dan ge melem a.}      \extran{Usually, this is women's work.}}}}

\entry{mɨka}{\headword{mɨka}  \note{Names of bananas}  \pronnote{mɪka}
  \pos{Noun}  \sense{
    \definition{type of introduced banana}    \note{Names of bananas}    \example{
}    \example{
}    \example{
}}}

\entry{mɨllɨng}{\headword{mɨllɨng}
  \variantof{SpellingVariant of \vartext{mälläng}}}

\entry{mɨnyi}{\headword{mɨnyi}  \pronnote{mɪɲi}
  \pos{TAM particle}  \sense{
    \definition{future tense particle}    \example{
      \exsen{Mɨnyi mer abal bogon.}      \extran{It will be very good.}}}
  \variants{\SpellingVariant \vartext{mänyi}}
  \variants{\Dialectal Variant \vartext{mäzi, mädi}}}

\entry{mo}{\headword{mo}  \note{Everything in the house}  \pronnote{mo}
  \pos{Noun}  \sense{\textbf{1.} 
    \definition{step, stair(s)}    \note{Everything in the house}    \example{
      \exsen{Medäda mo toko alle ikop eran llɨgda bom.}      \extran{The father is looking at his son from the top of the stairs.}}    \example{
      \exsen{Bogo gompedägän mo me.}      \extran{She tripped on the stairs.}}    \example{
}}  \sense{\textbf{2.} 
    \definition{bridge}    \note{Everything in the house}    \example{
      \exsen{Ubi tätäm ttongo mo sisor de daulliyu.}      \extran{Yesterday, they crossed over a new bridge.}}}}

\entry{mo katrekatre}{\headword{mo katrekatre}  \note{Everything in the house}  \pronnote{mo katrekatre}
  \pos{Noun}  \sense{
    \definition{staircase}    \note{Everything in the house}    \example{
}    \example{
}    \example{
}}}

\entry{mobera}{\headword{mobera}  \note{Water}  \pronnote{mobera}
  \pos{Noun}  \sense{
    \definition{outrigger}    \note{Water}    \example{
}    \example{
}    \example{
}}}

\entry{modaeb}{\headword{modaeb}  \pronnote{modaeb}
  \pos{Transitive verb}  \sense{
    \definition{to feel}}}

\entry{modgae}{\headword{modgae}  \pronnote{modɡae}
  \pos{Verb}  \sense{
    \definition{to chase away, potentially an Agob word?}}}

\entry{Moed}{\headword{Moed}
  \pos{Proper noun}  \sense{
    \definition{Moed (toponym)}}}

\entry{Moem}{\headword{Moem}  \pronnote{moem}
  \pos{Proper noun}  \sense{
    \definition{Moem (toponym)}}}

\entry{moep}{\headword{moep}  \pronnote{moep}
  \pos{Noun}  \sense{
    \definition{charcoal dust}}}

\entry{moepang}{\headword{moepang}  \note{Stages of life}  \pronnote{moepaŋ}
  \pos{Modifier}  \sense{\textbf{1.} 
    \definition{matured}    \note{Stages of life}    \example{
      \exsen{Bogo moepang agan.}      \extran{He is all grown up.}}    \example{
}    \example{
}}  \sense{\textbf{2.} 
    \definition{charcoal-stained}    \note{Stages of life}}}

\entry{moepo}{\headword{moepo}  \note{Poison}  \pronnote{moepo}
  \pos{Noun}  \sense{
    \definition{type of tree that grows in the bush with white flowers, inedible red fruit, and poisonous seeds and bark; an indicator that the soil is fertile and good for making a garden}    \note{Poison}    \example{
}    \example{
}    \example{
}}}

\entry{moepotatae}{\headword{moepotatae}  \note{Trees}  \pronnote{moepotatae}
  \pos{Noun}  \sense{
    \definition{type of tree}    \note{Trees}    \example{
}    \example{
}    \example{
}}}

\entry{mok}{\headword{mok}  \pronnote{mok}
  \pos{Noun}  \sense{
    \definition{friarbird (noisy, little, helmeted)}    \scientificname{Philemon corniculatus, P. citreogularis, P. buceroides}    \example{
}    \example{
}    \example{
}}}

\entry{mokmok}{\headword{mokmok}  \note{Verbs}  \pronnote{mokmok}
  \pos{Verb}  \sense{
    \definition{to wipe; dry by wiping}    \note{Verbs}    \example{
}    \example{
}}}

\entry{moko}{\headword{moko}  \pronnote{moko}
  \pos{Noun}  \sense{\textbf{1.} 
    \definition{desire, want; love}    \example{
      \exsen{Ngäna market e ibi allan, bäne moko da ewe dan?}      \extran{I am going to the market; do you want anything (lit. does your desire exist for anything)?}}    \example{
      \exsen{Endan bäne moko da ngäna bablle banges}      \extran{What do you want me to do for you?}}    \example{
      \exsen{Ngämi moko gogeya.}      \extran{We (excl.) two fell in love.}}    \example{
      \exsen{Ngäna up i moko allan.}      \extran{I like bananas.}}}  \sense{\textbf{2.} 
    \definition{taste, flavor}    \example{
      \exsen{Solt bo moko da ddone budabän.}      \extran{The salt won't lose its taste.}}    \example{
      \exsen{Sana gudne da koepang moko allan.}      \extran{The old sago tastes sour.}}    \example{
      \exsen{Ngatengate da biratt a ddobe moko dan.}      \extran{Roasted possum is very tasty.}}}}

\entry{mokoang}{\headword{mokoang}  \pronnote{mokoaŋ}
  \pos{Modifier}  \sense{\textbf{1.} 
    \definition{desirable, enjoyable; preferred, favorite}    \example{
      \exsen{Ngäna mokoang eran ngämo melem de.}      \extran{I love my job.}}    \example{
      \exsen{Ngäna bam mokoang nallan.}      \extran{I love you.}}    \example{
      \exsen{Ngäna ngämo mokowang bandra de llɨtnen eran.}      \extran{I am singing my favorite song.}}}  \sense{\textbf{2.} 
    \definition{tasty, delicious, flavorful; sweet}    \example{
      \exsen{Ge wätät a mokoang dan.}      \extran{This food is tasty.}}}
  \variants{\SpellingVariant \vartext{mokowang}}
  \variants{\Fast Speech Variant \vartext{mokwang}}}

\entry{mokoll}{\headword{mokoll}  \note{Trees}  \pronnote{mokoɽ}
  \pos{Noun}  \sense{
    \definition{type of tree that grows in the bush with thick bark, white flowers, and small green fruit}    \note{Trees}    \example{
}    \example{
}    \example{
}}}

\entry{mokon}{\headword{mokon}  \pronnote{mokon}
  \pos{Transitive verb}  \sense{
    \definition{to anoint}    \example{
      \exsen{Bogo itrellang lla de owel alle domokonmällaemneyo.}      \extran{He was anointing the sick men with oil.}}}}

\entry{mokowang}{\headword{mokowang}
  \variantof{SpellingVariant of \vartext{mokoang}}  \pronnote{mokowaŋ}}

\entry{moksir}{\headword{moksir}  \note{Trees}  \pronnote{moksir}
  \pos{Noun}  \sense{
    \definition{type of tree}    \note{Trees}    \example{
}    \example{
}    \example{
}}}

\entry{mokwang}{\headword{mokwang}
  \variantof{Fast Speech Variant of \vartext{mokoang}}}

\entry{molemoleg}{\headword{molemoleg}  \note{Grubs}  \pronnote{molemoleɡ}
  \pos{Noun}  \sense{
    \definition{type of grub}    \note{Grubs}    \example{
}    \example{
}    \example{
}}}

\entry{Moli}{\headword{Moli}  \pronnote{moli}
  \pos{Proper noun}  \sense{
    \definition{name of a male ancestor (brother of Lamlam)}}}

\entry{moll}{\headword{moll}  \note{Medicine}  \pronnote{moɽ}
  \pos{Noun}  \sense{
    \definition{type of tree that grows around Malam with white flowers}    \note{Medicine}    \example{
}    \example{
}    \example{
}}}

\entry{molle}{\headword{molle}  \note{Smell/scents}  \pronnote{moɽe}
  \pos{Noun}  \sense{
    \definition{scent, odor, smell}    \note{Smell/scents}    \example{
      \exsen{kuddäll molle}      \extran{rotten smell}}    \example{
      \exsen{Llo popo molle da mer dan.}      \extran{The scent of the flowers is good.}}    \example{
      \exsen{Ngäna otät molle de dändär eran.}      \extran{I am smelling food.}}}}

\entry{molle dändär}{\headword{molle dändär}  \pronnote{moɽe dəndər}
  \pos{Transitive compound verb}  \sense{
    \definition{to smell, perceive a smell}    \example{
      \exsen{Ngäna sɨmell de gänyaollemae molle deyandärneg.}      \extran{I smelled pigs coming this way.}}}}

\entry{molle dungg}{\headword{molle dungg}  \note{Verbs}  \pronnote{moɽe duŋɡ}
  \pos{Transitive compound verb}  \sense{
    \definition{to smell, take a whiff}    \note{Verbs}    \example{
}    \example{
}}}

\entry{molleang}{\headword{molleang}  \pronnote{moɽeaŋ}
  \pos{Modifier}  \sense{
    \definition{scented}    \example{
      \exsen{mer molleang owel}      \extran{fragrant oil}}}
  \variants{\Fast Speech Variant \vartext{mollong}}}

\entry{mollemeny}{\headword{mollemeny}  \pronnote{moɽemeɲ}
  \pos{Modifier}  \sense{
    \definition{odorless}}}

\entry{mollemolle}{\headword{mollemolle}
  \pos{Transitive coverb}  \sense{
    \definition{to sniff, smell}    \example{
      \exsen{Däräng a llɨg di mollemolle nägagan.}      \extran{The dog sniffed the boy.}}}}

\entry{mollok}{\headword{mollok}  \note{Names of bananas}  \pronnote{moɽok}
  \pos{Noun}  \sense{
    \definition{type of introduced banana}    \note{Names of bananas}    \example{
}    \example{
}    \example{
}}}

\entry{mollong}{\headword{mollong}
  \variantof{Fast Speech Variant of \vartext{molleang}}  \pronnote{moɽoŋ}}

\entry{momana ma}{\headword{momana ma}
  \pos{Noun}  \sense{
    \definition{hideout made of leaves for hunting cassowaries}    \example{
      \exsen{Ddob llaeyabaene momana ma alle dirom gädnanatt a kemibi dag.}      \extran{Many people kill cassowaries by using a hideout.}}}}

\entry{Mome}{\headword{Mome}  \pronnote{mome}
  \pos{Proper noun}  \sense{
    \definition{female personal name}}}

\entry{momea gäl}{\headword{momea gäl}  \note{Trees}  \pronnote{momea ɡəl}
  \pos{Noun}  \sense{
    \definition{type of tree}    \note{Trees}    \example{
}    \example{
}    \example{
}}}

\entry{Momeya}{\headword{Momeya}
  \pos{Proper noun}  \sense{
    \definition{Momeya (toponym)}}}

\entry{momolltätän}{\headword{momolltätän}  \note{Names of bananas}  \pronnote{momoɽtətən}
  \pos{Noun}  \sense{
    \definition{type of introduced banana}    \note{Names of bananas}    \example{
}    \example{
}    \example{
}}}

\entry{mompara}{\headword{mompara}  \note{Trees}  \pronnote{mompara}
  \pos{Noun}  \sense{
    \definition{type of tree that grows near swamps and creeks with white flowers and yellow and brown fruit that ripen in April and May and are eaten by deer}    \note{Trees}    \example{
}    \example{
}    \example{
}}}

\entry{mompel}{\headword{mompel}  \note{Trees}  \pronnote{mompel}
  \pos{Noun}  \sense{
    \definition{aibika}    \scientificname{Abelmoschus manihot}    \note{Trees}    \anthronote{shrub whose leaves are eaten with sago}    \example{
}    \example{
}    \example{
}}}

\entry{Mompelang}{\headword{Mompelang}  \note{Places around Limol}  \pronnote{mompelaŋ}
  \pos{Proper noun}  \sense{
    \definition{Mompelang (toponym)}    \note{Places around Limol}    \example{
}    \example{
}    \example{
}}}

\entry{mondo}{\headword{mondo}  \note{Trees}  \pronnote{mondo}
  \pos{Noun}  \sense{
    \definition{type of tree that grows in the bush with green or purple fruits that animals eat}    \note{Trees}    \example{
}    \example{
}    \example{
}}}

\entry{mondre}{\headword{mondre}  \note{Things a person can do}  \pronnote{mondre}
  \pos{Noun}  \sense{
    \definition{gardening, garden work}    \note{Things a person can do}    \example{
      \exsen{Bogo mondre bognän.}      \extran{She will be gardening.}}    \example{
}    \example{
}}}

\entry{mondreag}{\headword{mondreag}  \pronnote{mondreaɡ}
  \pos{Modifier}  \sense{
    \definition{gardener, having a green thumb}    \example{
      \exsen{Ddobae mondrog lla deya.}      \extran{He was a great gardener.}}}
  \variants{\Fast Speech Variant \vartext{mondrog}}}

\entry{mondremondre}{\headword{mondremondre}  \pronnote{mondremondre}
  \pos{Transitive/Intransitive coverb}  \sense{
    \definition{to move}    \example{
      \exsen{Bäne ttang de ade mɨnyi mondremondre bägagän.}      \extran{Your hand will also be moving (lit. movement will also get your hand).}}    \example{
      \exsen{Däbe lla da mondremondre gogon.}      \extran{That man moved.}}}}

\entry{mondrog}{\headword{mondrog}
  \variantof{Fast Speech Variant of \vartext{mondreag}}  \pronnote{mondroɡ}}

\entry{mopmop}{\headword{mopmop}  \note{Trees}  \pronnote{mopmop}
  \pos{Noun}  \sense{
    \definition{type of tree that grows in the grassland (~3 m) with white flowers and red fruit that are eaten to treat cough}    \note{Trees}    \example{
}    \example{
}    \example{
}}}

\entry{mormor}{\headword{mormor}  \lexeme{mormor}  \pronnote{mormor}
  \pos{Noun}  \sense{
    \definition{type of small herb with white flowers and ovate leaves}}}

\entry{mos}{\headword{mos}  \lexeme{mos}  \pronnote{mos}
  \pos{Noun}  \sense{
    \definition{type of goanna}}}

\entry{Mosbi}{\headword{Mosbi}
  \variantof{Fast Speech Variant of \vartext{Pot Mosbi}}
  \variants{\SpellingVariant \vartext{Mozbi}}
  \variants{\Unspecified Variant \vartext{Mospi}}}

\entry{mosemosen}{\headword{mosemosen}
  \variantof{Fast Speech Variant of \vartext{mosenmosen}}  \pronnote{mosemosen}}

\entry{mosen}{\headword{mosen}  \note{Family ties}  \pronnote{mosen}
  \pos{Kinship noun}  \sense{\textbf{1.} 
    \definition{older sibling of the same-sex (man's older brother or woman's older sister)}    \note{Family ties}    \example{
      \exsen{Paul ngämo mosen dan.}      \extran{Paul is my older brother (said by a male speaker).}}}  \sense{\textbf{2.} 
    \definition{eldest, firstborn}    \note{Family ties}    \example{
      \exsen{Ngäna mosen llɨg dan.}      \extran{I'm the firstborn.}}}}

\entry{mosen~}{\headword{mosen~}
  \variantof{Unspecified Variant of \vartext{RED~}}}

\entry{mosenmosen}{\headword{mosenmosen}
  \pos{Kinship noun}  \sense{\textbf{1.} 
    \definition{older siblings}}  \sense{\textbf{2.} 
    \definition{elders, seniors}    \example{
      \exsen{Ngäma kame dan, be mosenmosen lla da mɨnyi ngämim umllang beyageyo.}      \extran{We (excl.) don't know, but our elders will tell us.}}}
  \variants{\Fast Speech Variant \vartext{mosemosen}}}

\entry{Moses}{\headword{Moses}  \pronnote{moses}
  \pos{Proper noun}  \sense{
    \definition{male personal name}}}

\entry{Mospi}{\headword{Mospi}
  \variantof{Unspecified Variant of \vartext{Mosbi}}  \pronnote{mospi}}

\entry{motom}{\headword{motom}  \note{Types of Taro}  \pronnote{motom}
  \pos{Noun}  \sense{
    \definition{type of small taro}    \note{Types of Taro}    \example{
}    \example{
}    \example{
}}}

\entry{Motu}{\headword{Motu}  \note{Language Names}  \pronnote{motu}
  \pos{Proper noun}  \sense{
    \definition{Hiri Motu language}    \note{Language Names}    \example{
}    \example{
}    \example{
}}
  \variants{\SpellingVariant \vartext{Mutu}}
  \variants{\Dialectal Variant \vartext{Muti}}}

\entry{Motu loanword}{\headword{Motu loanword}
  \pos{}  \sense{
}}

\entry{Moyabag}{\headword{Moyabag}
  \pos{Proper noun}  \sense{
    \definition{male personal name}}}

\entry{Moyäm}{\headword{Moyäm}  \pronnote{mojəm}
  \pos{Noun}  \sense{
    \definition{Moyäm (toponym)}}}

\entry{moza}{\headword{moza}  \lexeme{moza}  \pronnote{moza}
  \pos{Noun}  \sense{
    \definition{type of venomous snake}}}

\entry{mozaeya}{\headword{mozaeya}
  \variantof{SpellingVariant of \vartext{mozaya}}}

\entry{mozaya}{\headword{mozaya}
  \pos{Noun}  \sense{
    \definition{type of large, edible fish found in the swamp}}
  \variants{\SpellingVariant \vartext{mozaeya}}}

\entry{Mozbi}{\headword{Mozbi}
  \variantof{SpellingVariant of \vartext{Mosbi}}}

\entry{mu}{\headword{mu}  \pronnote{mu}
  \pos{Noun}  \sense{\textbf{1.} 
    \definition{payment, price, value}    \example{
      \exsen{skul mu}      \extran{school fees}}    \example{
      \exsen{Bablle mɨnyi mu banttog.}      \extran{I will pay you (lit. give you payment).}}    \example{
      \exsen{300 kina mu me dägageyo däbe käza ttoe de.}      \extran{That crocodile skin sold for 300 kina.}}}  \sense{\textbf{2.} 
    \definition{response, reply, answer; repayment, revenge}    \example{
      \exsen{Obo moko da mu wi dan.}      \extran{He wants to respond.}}    \example{
      \exsen{Polis a ngämo gullbe de namllamallo, ngämo mu di nangesallo.}      \extran{The police got my husband; they carried out my revenge.}}}  \sense{\textbf{3.} 
    \definition{valuable}    \example{
      \exsen{Obo ttoe a ddobae mu dan.}      \extran{Its skin is very valuable.}}}
  \variants{\SpellingVariant \vartext{muu}}}

\entry{mubine}{\headword{mubine}
  \pos{Noun}  \sense{
    \definition{dish consisting of food (such as banana, yam, or sweet potato) with coconut cream. up mubine, mätta mubine, nai mubine}}}

\entry{muda}{\headword{muda}
  \pos{TAM particle}  \sense{
    \definition{prohibitive}}}

\entry{mudag}{\headword{mudag}  \pronnote{mudaɡ}
  \pos{Copular verb}  \sense{
    \definition{present plural form of mudan}    \example{
      \exsen{Abo källkae ibi da mudag.}      \extran{[You all] never go there again.}}    \example{
      \exsen{Bibi eka da mudag.}      \extran{You all, don't talk.}}}}

\entry{mudageyo}{\headword{mudageyo}
  \pos{Copular verb}  \sense{
    \definition{present dual form of mudan}    \example{
      \exsen{Bibi komllaebe gäddnan a mudageyo.}      \extran{You two, don't fight.}}}}

\entry{mudan}{\headword{mudan}  \pronnote{mudan}
  \pos{Copular verb}  \sense{
    \definition{prohibitive copula (present singular form)}    \example{
      \exsen{Ngänäm kollmäl a mudan.}      \extran{[You] don't follow me.}}    \example{
      \exsen{Mudan lel a.}      \extran{[You] don't be afraid.}}}}

\entry{mugbusu}{\headword{mugbusu}  \note{Types of Taro}  \pronnote{muɡbusu}
  \pos{Noun}  \sense{
    \definition{type of taro}    \note{Types of Taro}    \example{
}    \example{
}    \example{
}}}

\entry{Mugi}{\headword{Mugi}
  \pos{Proper noun}  \sense{
    \definition{male personal name}}}

\entry{Muidebag}{\headword{Muidebag}  \note{Places around Limol}  \pronnote{muidebaɡ}
  \pos{Proper noun}  \sense{
    \definition{Muidebag (toponym)}    \note{Places around Limol}    \example{
}    \example{
}    \example{
}}}

\entry{Mul}{\headword{Mul}  \note{Place names around Limol}  \pronnote{mul}
  \pos{Proper noun}  \sense{
    \definition{Mul (toponym)}    \note{Place names around Limol}    \example{
}    \example{
}    \example{
}}}

\entry{mulkul}{\headword{mulkul}  \note{Parts of the body}  \pronnote{mulkul}
  \pos{Noun}  \sense{
    \definition{brain}    \note{Parts of the body}    \example{
}    \example{
}    \example{
}}}

\entry{Mull}{\headword{Mull}
  \pos{Proper noun}  \sense{
    \definition{Mull (toponym)}}}

\entry{mullae}{\headword{mullae}  \pronnote{muɽae}
  \pos{Modifier}  \sense{\textbf{1.} 
    \definition{enough}    \example{
      \exsen{Iba ge kollba da mullae dag.}      \extran{We (incl.) have enough fish.}}}  \sense{\textbf{2.} 
    \definition{able, can, be allowed}    \example{
      \exsen{Llamda da bälle gazenma da ddone mullae gogon.}      \extran{The old man was unable to escape.}}    \example{
      \exsen{Bäne ma me ngäna mullae dan bodmen?}      \extran{Can I stay at your house?}}}
  \variants{\Unspecified Variant \vartext{mulldae}}}

\entry{mullae~}{\headword{mullae~}
  \variantof{Dialectal Variant of \vartext{RED~}}}

\entry{mullaemullae}{\headword{mullaemullae}  \pronnote{muɽaemuɽae}
  \pos{Modifier}  \sense{
    \definition{nonsingular form of mullae}    \example{
      \exsen{Tämamae ttoen a dedme mullamullae dag.}      \extran{There's enough of everything there.}}}
  \variants{\Fast Speech Variant \vartext{mullamullae}}}

\entry{mullaemullae}{\headword{mullaemullae}
  \pos{Quantifier}  \sense{
    \definition{every}    \example{
      \exsen{Ngäna pazi mullamullae mänmän de ngällabnanangmeae anggan.}      \extran{I pick up the girls every year.}}}
  \variants{\Fast Speech Variant \vartext{mullamullae}}}

\entry{mullamulla}{\headword{mullamulla}  \pronnote{muɽamuɽa}
  \pos{Noun}  \sense{
    \definition{medicine}}
  \variants{\SpellingVariant \vartext{muramura}}}

\entry{mullamullae}{\headword{mullamullae}
  \variantof{Fast Speech Variant of \vartext{mullaemullae}}}

\entry{mulldae}{\headword{mulldae}
  \variantof{Unspecified Variant of \vartext{mullae}}  \pronnote{muɽdae}}

\entry{mulmul}{\headword{mulmul}
  \pos{Noun}  \sense{
    \definition{rite of passage}}}

\entry{muminy}{\headword{muminy}  \note{Market}  \pronnote{mumiɲ}
  \pos{Modifier}  \sense{
    \definition{in debt}    \note{Market}    \example{
      \exsen{Ngämo ttongo ttang lläpät a muminy dan.}      \extran{I owe five. / I am five in debt.}}    \example{
}    \example{
}}}

\entry{Munu}{\headword{Munu}  \pronnote{munu}
  \pos{Proper noun}  \sense{
    \definition{male personal name}}}

\entry{mupni}{\headword{mupni}  \note{Trees}  \pronnote{mupni}
  \pos{Noun}  \sense{
    \definition{type of cultivated mango tree with fruit with yellow skin and a white interior that is juiced}    \note{Trees}    \example{
}    \example{
}    \example{
}}}

\entry{Mur}{\headword{Mur}  \pronnote{mur}
  \pos{Proper noun}  \sense{
    \definition{Mur (toponym)}    \example{
}    \example{
}    \example{
}}}

\entry{mur}{\headword{mur}  \pronnote{mur}
  \pos{Transitive coverb}  \sense{
    \definition{to clean, tidy}}}

\entry{muramura}{\headword{muramura}
  \variantof{SpellingVariant of \vartext{mullamulla}}}

\entry{muro}{\headword{muro}
  \pos{Noun}  \sense{
    \definition{magic}}}

\entry{Musato}{\headword{Musato}  \pronnote{musato}
  \pos{Proper noun}  \sense{
    \definition{female personal name}}}

\entry{must}{\headword{must}
  \variantof{freevarlab of \vartext{English loanword}}}

\entry{mutae}{\headword{mutae}  \lexeme{mutae}  \pronnote{mutae}
  \pos{Noun}  \sense{
    \definition{type of yam with a yellow interior, no thorns, and a vine that grows clockwise}}}

\entry{Muti}{\headword{Muti}
  \variantof{Dialectal Variant of \vartext{Motu}}}

\entry{mutt}{\headword{mutt}  \lexeme{mutt}  \pronnote{muʈ͡ʂ}
  \pos{Noun}  \sense{
    \definition{river source}    \example{
      \exsen{Ngäna walle mutt i ibi allan.}      \extran{I am going upstream.}}}}

\entry{muttbul}{\headword{muttbul}
  \pos{Modifier}  \sense{
    \definition{quiet}}}

\entry{muttmutt}{\headword{muttmutt}  \note{Poison}  \pronnote{muʈ͡ʂmuʈ͡ʂ}
  \pos{Noun}  \sense{
    \definition{type of grub that lives in the ground and eats crops}    \note{Poison}    \example{
}    \example{
}    \example{
}}}

\entry{Mutu}{\headword{Mutu}
  \variantof{SpellingVariant of \vartext{Motu}}}

\entry{Muyabag}{\headword{Muyabag}  \pronnote{mujabaɡ}
  \pos{Proper noun}  \sense{
    \definition{male personal name}}}

\entry{muu}{\headword{muu}
  \variantof{SpellingVariant of \vartext{mu}}}

\chapter{N}


\entry{n-}{\headword{n-}
  \pos{Verb}  \sense{
}}

\entry{n-}{\headword{n-}  \pronnote{n}
  \pos{Verb}  \sense{\textbf{1.} 
    \definition{marks a singular or third person plural patient argument of a transitive verb in the recent past tense}    \example{
      \exsen{Otät a kuddäll e nabäddan.}      \extran{Food killed me.}}}  \sense{\textbf{2.} 
    \definition{marks a singular or third person plural patient argument of a transitive verb in the future tense (/irrelais mood), but only when a third person acts on a second person, or a second person acts on any person.}}  \sense{\textbf{3.} 
}  \sense{\textbf{4.} 
}  \sense{\textbf{5.} 
}  \sense{\textbf{6.} 
}  \sense{\textbf{7.} 
}  \sense{\textbf{8.} 
}  \sense{\textbf{9.} 
}}

\entry{=n}{\headword{=n}
  \pos{Verbal enclitic}  \sense{
    \definition{cop.prs.sgS}}
  \variants{\Dialectal Variant \vartext{=wän, =nän}}
  \variants{\freevarlab \vartext{=0}}}

\entry{na~}{\headword{na~}
  \variantof{Infinitival reduplicant of \vartext{RED~}}}

\entry{nadum}{\headword{nadum}
  \pos{Noun}  \sense{
    \definition{namesake}    \example{
      \exsen{Ngämo nadum a daden.}      \extran{I have a namesake.}}}}

\entry{nae}{\headword{nae}  \lexeme{nae}  \pronnote{nae}
  \pos{Noun}  \sense{
    \definition{sweet potato}}}

\entry{naeka}{\headword{naeka}
  \variantof{Fast Speech Variant of \vartext{ngänaeka}}  \pronnote{naeka}}

\entry{Naemäll}{\headword{Naemäll}
  \pos{Proper noun}  \sense{
    \definition{female personal name}}}

\entry{naen}{\headword{naen}  \pronnote{naen}
  \pos{Numeral}  \sense{
    \definition{nine (English numeral; also general)}}
  \variants{\SpellingVariant \vartext{nain}}}

\entry{naenti}{\headword{naenti}
  \pos{Numeral}  \sense{
    \definition{ninety}}
  \variants{\SpellingVariant \vartext{nainti, ninety}}}

\entry{naentin}{\headword{naentin}
  \pos{Numeral}  \sense{
    \definition{nineteen (English numeral)}}
  \variants{\SpellingVariant \vartext{naintin}}}

\entry{nag}{\headword{nag}  \pronnote{naɡ}
  \pos{Kinship noun}  \sense{
    \definition{friend}    \example{
}    \example{
}    \example{
}}}

\entry{nag~}{\headword{nag~}
  \variantof{Plural reduplicant of \vartext{RED~}}}

\entry{nag ttoen}{\headword{nag ttoen}  \pronnote{naɡ ʈ͡ʂoen}
  \pos{Noun}  \sense{
    \definition{friendship}    \example{
      \exsen{Oba nag ttoen a dädme llätt gogon.}      \extran{Their friendship ended there.}}}}

\entry{Nagab}{\headword{Nagab}  \pronnote{naɡab}
  \pos{Proper noun}  \sense{
    \definition{male personal name}}}

\entry{Nägäm}{\headword{Nägäm}  \pronnote{nəɡəm}
  \pos{Proper noun}  \sense{
    \definition{male personal name}}}

\entry{Nagat}{\headword{Nagat}
  \pos{Proper noun}  \sense{
    \definition{male personal name}}}

\entry{Nageg}{\headword{Nageg}  \pronnote{naɡeɡ}
  \pos{Proper noun}  \sense{
    \definition{male personal name}}}

\entry{nagnag}{\headword{nagnag}  \pronnote{naɡnaɡ}
  \pos{Noun}  \sense{
    \definition{nonsingular form of nag}}}

\entry{naigae}{\headword{naigae}  \pronnote{naiɡae}
  \pos{Modifier}  \sense{
    \definition{south}    \example{
      \exsen{Ngäna naigae pallall e ibi allan.}      \extran{I am heading south.}}}}

\entry{nain}{\headword{nain}
  \variantof{SpellingVariant of \vartext{naen}}}

\entry{nainti}{\headword{nainti}
  \variantof{SpellingVariant of \vartext{naenti}}}

\entry{naintin}{\headword{naintin}
  \variantof{SpellingVariant of \vartext{naentin}}}

\entry{näk}{\headword{näk}  \pronnote{nək}
  \pos{Intransitive coverb}  \sense{
    \definition{to hiccup, or to have a pressure feeling on your throat or chest from eating too many bananas}}}

\entry{Nakaku}{\headword{Nakaku}
  \pos{Proper noun}  \sense{
    \definition{Nakaku (toponym)}}}

\entry{näkäp}{\headword{näkäp}  \pronnote{nəkəp}
  \pos{Noun}  \sense{
    \definition{mind, mindset, consciousness; thoughts}    \example{
      \exsen{Ngämo näkäp e gongttägän ada ngaska zime iddob amänong agan.}      \extran{It came to my mind that maybe it's already past midnight.}}    \example{
      \exsen{Oblle mer näkäp a ddone ngänttägmäll allan.}      \extran{He is not having good thoughts.}}}}

\entry{näkäp de pänae}{\headword{näkäp de pänae}  \note{Everyday words and phrases}  \pronnote{nəkəp de pənae}
  \pos{Phrase}  \sense{
    \definition{to change one's mind}    \note{Everyday words and phrases}    \example{
      \exsen{Ngäna mäse karama we ibi we daeya, be näkäp de näpänaeyan wälläng e.}      \extran{I was supposed to go to the canoe place, but I changed my mind to go to the bush.}}    \example{
}    \example{
}}}

\entry{näkäp ttomoe}{\headword{näkäp ttomoe}  \pronnote{nəkəp ʈomoe}
  \pos{Intransitive compound verb}  \sense{
    \definition{to worry}    \example{
      \exsen{Bogo näkäp ttomoenen me giddollnen allan.}      \extran{He is living worrying.}}}}

\entry{näkäpang}{\headword{näkäpang}  \note{School}  \pronnote{nəkəpaŋ}
  \pos{Modifier}  \sense{
    \definition{smart}    \note{School}    \example{
}    \example{
}    \example{
}}}

\entry{näkäpmeny}{\headword{näkäpmeny}
  \pos{Modifier}  \sense{
    \definition{foolish}}}

\entry{näkäpngon}{\headword{näkäpngon}
  \pos{Noun}  \sense{
    \definition{thought, knowledge}    \example{
      \exsen{Obo näkäpngon tämamae ai da erageya audaebnegan.}      \extran{He lost all his thoughts that were good.}}    \example{
      \exsen{Obo näkäpngon alle eka panypeny eran.}      \extran{She is speaking with wisdom.}}}}

\entry{Naklae}{\headword{Naklae}  \pronnote{naklae}
  \pos{Proper noun}  \sense{
    \definition{male personal name}}
  \variants{\Unspecified Variant \vartext{Nakllae}}}

\entry{Nakllae}{\headword{Nakllae}
  \variantof{Unspecified Variant of \vartext{Naklae}}}

\entry{Nakuri}{\headword{Nakuri}  \pronnote{nakuri}
  \pos{Proper noun}  \sense{
    \definition{male personal name}}}

\entry{nälan}{\headword{nälan}
  \pos{Manner adverb}  \sense{
    \definition{accidentally}    \example{
      \exsen{Bogo ankom de näbäddan nälan ttang alle.}      \extran{He accidentally killed the ant with his hand.}}}}

\entry{=nalla}{\headword{=nalla}  \pronnote{naɽa}
  \pos{Auxiliary verb}  \sense{
    \definition{aux.prs.1|2nsgA>1|2sgP}}}

\entry{=nallan}{\headword{=nallan}  \pronnote{naɽan}
  \pos{Auxiliary verb}  \sense{\textbf{1.} 
    \definition{aux.prs.1|3sgA>1|2sgP}}  \sense{\textbf{2.} 
}  \sense{\textbf{3.} 
}}

\entry{=nalle}{\headword{=nalle}
  \pos{Auxiliary verb}  \sense{
    \definition{aux.prs.2sgA>1sgP}}}

\entry{nallib}{\headword{nallib}  \note{Traditional Arrows}  \pronnote{naɽib}
  \pos{Noun}  \sense{
    \definition{type of spear}    \note{Traditional Arrows}    \example{
}    \example{
}    \example{
}}}

\entry{=nallo}{\headword{=nallo}  \pronnote{naɽo}
  \pos{Auxiliary verb}  \sense{\textbf{1.} 
    \definition{aux.prs.3nsgA>1|2sgP}}  \sense{\textbf{2.} 
}}

\entry{Nalon}{\headword{Nalon}
  \pos{Proper noun}  \sense{
    \definition{male personal name}}}

\entry{Nama}{\headword{Nama}  \pronnote{nama}
  \pos{Proper noun}  \sense{
    \definition{male personal name}}}

\entry{Namaya}{\headword{Namaya}  \pronnote{namaja}
  \pos{Proper noun}  \sense{
    \definition{female personal name}}}

\entry{namba}{\headword{namba}
  \variantof{Dialectal Variant of \vartext{English loanword}}}

\entry{=nän}{\headword{=nän}
  \variantof{Dialectal Variant of \vartext{=n}}}

\entry{nänäm}{\headword{nänäm}  \note{Weaving}  \pronnote{nənəm}
  \pos{Noun}  \sense{
    \definition{diagonal checkerboard weaving pattern}    \note{Weaving}    \example{
}    \example{
}    \example{
}}}

\entry{Nancy}{\headword{Nancy}
  \variantof{SpellingVariant of \vartext{Nensi}}}

\entry{nane}{\headword{nane}  \lexeme{nane}  \pronnote{nane}
  \pos{Kinship noun}  \sense{
    \definition{aunt (one's parent's younger sister)}    \example{
      \exsen{Nensi bi nane alle känaebag mɨnyi tudi ma beyareyo}      \extran{Nancy will go fishing tomorrow with her aunt.}}}}

\entry{nane}{\headword{nane}  \note{Things a person can do}  \pronnote{nane}
  \pos{Transitive verb}  \sense{
    \definition{to drink}    \note{Things a person can do}    \example{
      \exsen{Bogo nane eran.}      \extran{He is drinking.}}    \example{
      \exsen{Ngäna konkonang ine de bäna.}      \extran{I will drink the beer.}}    \example{
      \exsen{Ngämi dänaeyaemnalla.}      \extran{We (excl.) were drinking.}}}}

\entry{Nänga}{\headword{Nänga}
  \pos{Proper noun}  \sense{
    \definition{female personal name}}}

\entry{nängga}{\headword{nängga}  \pronnote{nəŋɡa}
  \pos{Noun}  \sense{
    \definition{type of pandanus with bunches of fruit with strong shells; they fall one at a time}    \example{
}    \example{
}    \example{
}}}

\entry{Nanggon}{\headword{Nanggon}
  \pos{Proper noun}  \sense{
    \definition{male personal name}}}

\entry{nängttäg}{\headword{nängttäg}
  \pos{Intransitive verb}  \sense{
    \definition{to come}}}

\entry{Naomi}{\headword{Naomi}  \pronnote{naomi}
  \pos{Proper noun}  \sense{
    \definition{female personal name}}}

\entry{Narma}{\headword{Narma}  \pronnote{narma}
  \pos{Proper noun}  \sense{
    \definition{male personal name}}}

\entry{Nasma}{\headword{Nasma}  \pronnote{nasma}
  \pos{Proper noun}  \sense{
    \definition{male personal name}}}

\entry{nätnät}{\headword{nätnät}
  \pos{Noun}  \sense{
    \definition{insects}}}

\entry{Nazaret}{\headword{Nazaret}  \pronnote{nazaret}
  \pos{Proper noun}  \sense{
    \definition{Nazareth}}}

\entry{Nedlyn}{\headword{Nedlyn}  \pronnote{nedlin}
  \pos{Proper noun}  \sense{
    \definition{female personal name}}}

\entry{Nensi}{\headword{Nensi}  \pronnote{nensi}
  \pos{Proper noun}  \sense{
    \definition{female personal name}}
  \variants{\SpellingVariant \vartext{Nancy}}}

\entry{nepi}{\headword{nepi}  \note{Clothing}  \pronnote{nepi}
  \pos{Noun}  \sense{
    \definition{napkins}    \note{Clothing}    \example{
}    \example{
}    \example{
}}}

\entry{nes}{\headword{nes}
  \pos{Noun}  \sense{
    \definition{nurse}}}

\entry{net}{\headword{net}
  \pos{Noun}  \sense{
    \definition{net}    \example{
      \exsen{Ngäna ttongo net de dirängän ine watt.}      \extran{I lifted a net out from the water.}}}
  \variants{\Unspecified Variant \vartext{nett}}}

\entry{nett}{\headword{nett}
  \variantof{freevarlab of \vartext{English loanword}}}

\entry{New}{\headword{New}
  \variantof{Dialectal Variant of \vartext{English loanword}}}

\entry{ngä~}{\headword{ngä~}
  \variantof{Infinitival reduplicant of \vartext{RED~}}}

\entry{nga}{\headword{nga}  \pronnote{ŋa}
  \pos{TAM particle}  \sense{
    \definition{marks immediate or near future}    \example{
      \exsen{Nga yu ttätta de däbe gänyaolle iweny.}      \extran{Now [you] bring that piece of burned wood over here.}}    \example{
      \exsen{Nga bibi ddone ngänaem amalla?}      \extran{Are you all not understanding?}}}}

\entry{nga~}{\headword{nga~}
  \variantof{Infinitival reduplicant of \vartext{RED~}}}

\entry{ngäd~}{\headword{ngäd~}
  \variantof{freevarlab of \vartext{RED~}}}

\entry{ngädnan~}{\headword{ngädnan~}
  \variantof{freevarlab of \vartext{RED~}}}

\entry{ngädngäd}{\headword{ngädngäd}  \pronnote{ŋədŋəd}
  \pos{Intransitive verb}  \sense{\textbf{1.} 
    \definition{to roll up, curl}    \example{
      \exsen{Gullem a gongädän.}      \extran{The snake coiled up.}}}  \sense{\textbf{2.} 
    \definition{to roll up, curl}    \example{
      \exsen{Ngämi tater nängädaeballa.}      \extran{We (excl.) rolled up the mats.}}}  \sense{\textbf{3.} 
    \definition{to fold}    \example{
      \exsen{Bogo ttang nängädnegan.}      \extran{She folded her arms.}}}}

\entry{ngaeka}{\headword{ngaeka}
  \variantof{Fast Speech Variant of \vartext{ngänaeka}}}

\entry{ngak}{\headword{ngak}
  \variantof{Unspecified Variant of \vartext{ngok}}}

\entry{ngal}{\headword{ngal}  \pronnote{ŋal}
  \pos{Intransitive coverb}  \sense{
    \definition{to waddle, shuffle}}}

\entry{ngalbongalboe}{\headword{ngalbongalboe}  \pronnote{ŋalboŋalboe}
  \pos{Adverb}  \sense{
    \definition{on the side}    \example{
      \exsen{Bogo ngalbongalboe bognegnän.}      \extran{He will be on the side.}}}}

\entry{ngalen}{\headword{ngalen}  \pronnote{ŋalen}
  \pos{Noun}  \sense{\textbf{1.} 
    \definition{way, habit, manner, custom}    \example{
      \exsen{Gabma yaba ngalen a gongttägän ngäma pate.}      \extran{White people's customs came to us (excl.).}}}  \sense{\textbf{2.} 
    \definition{thing}    \example{
      \exsen{Golläntmenyän tämamae ngalen de.}      \extran{He told him everything.}}}
  \variants{\SpellingVariant \vartext{ngälen}}}

\entry{ngälen}{\headword{ngälen}
  \variantof{SpellingVariant of \vartext{ngalen}}  \pronnote{ŋəlen}}

\entry{ngäll~}{\headword{ngäll~}
  \variantof{freevarlab of \vartext{RED~}}}

\entry{ngälla~}{\headword{ngälla~}
  \variantof{freevarlab of \vartext{RED~}}}

\entry{ngälladab}{\headword{ngälladab}
  \pos{Verb}  \sense{
}}

\entry{ngällae}{\headword{ngällae}  \pronnote{ŋəɽae}
  \pos{Intransitive verb}  \sense{
    \definition{to look back}    \example{
      \exsen{Ngäna ngllae allan.}      \extran{I am looking back.}}    \example{
      \exsen{Bogo mäse dinduag me ako gongllaeyän däräng bälle, be ddone obom ikop dägagän.}      \extran{While running, she tried to look back for the dog, but she couldn't see him.}}}}

\entry{ngällban}{\headword{ngällban}
  \variantof{Unspecified Variant of \vartext{ngällbän}}  \pronnote{ŋəɽban}}

\entry{ngällbän}{\headword{ngällbän}  \pronnote{ŋəɽbən}
  \pos{Intransitive verb}  \sense{\textbf{1.} 
    \definition{to rise, arise, come up}    \example{
      \exsen{Ine ulle da gongälllbänmällnän.}      \extran{The big wave was rising.}}}  \sense{\textbf{2.} 
    \definition{to awake, wake up, get up}    \example{
      \exsen{Ddobag a gudae ngällbaenen eran yunu atta.}      \extran{Some wake up from sleep early.}}}  \sense{\textbf{3.} 
    \definition{to lift}    \example{
      \exsen{Kottllam de kapu wi dängällbäneyo.}      \extran{They lifted the turtle up to carry.}}}  \sense{\textbf{4.} 
    \definition{to get, take}    \example{
      \exsen{teks mani ngällbänang lla}      \extran{tax collector}}    \example{
      \exsen{Dindugän, do täl pallkoll dingällbänän.}      \extran{He ran and got a piece of bamboo.}}}  \sense{\textbf{5.} 
    \definition{to wake up, to wake}    \example{
      \exsen{Ngäna bam danglläbän.}      \extran{I woke you up.}}}
  \variants{\Unspecified Variant \vartext{ngällban}}}

\entry{ngallmeny}{\headword{ngallmeny}  \pronnote{ŋaɽmeɲ}
  \pos{Transitive verb}  \sense{
    \definition{to advise}    \example{
      \exsen{Llɨg ngallmeny a era kälsre meae dan mer a.}      \extran{The best time to advise children is in childhood.}}}}

\entry{ngällngäll}{\headword{ngällngäll}  \pronnote{ŋəɽŋəɽ}
  \pos{Transitive verb}  \sense{
    \definition{to produce, yield, bear (fruit)}    \example{
      \exsen{Wit däm a ulleulle gognegän a käp de yuwog gongällanegän.}      \extran{The wheat grew big and produced lots of grain.}}    \example{
      \exsen{Nge da ngällngällang dan.}      \extran{The coconut tree is bearing fruit.}}}}

\entry{ngallngall}{\headword{ngallngall}  \pronnote{ŋaɽŋaɽ}
  \pos{Noun}  \sense{
    \definition{catbird (spotted, ochre-breasted)}    \scientificname{Aliruoedus maculosus, A. stonii}    \example{
}    \example{
}    \example{
}}}

\entry{ngälngäl}{\headword{ngälngäl}  \note{Adjectives}  \pronnote{ŋəlŋəl}
  \pos{Modifier}  \sense{
    \definition{round}    \note{Adjectives}    \example{
      \exsen{Ngälngälatt dädär a mänyi tutu atta bolldaeyan kopek e.}      \extran{A round rock will roll from the hill to the valley.}}    \example{
}    \example{
}}}

\entry{ngam}{\headword{ngam}  \pronnote{ŋam}
  \pos{Noun}  \sense{\textbf{1.} 
    \definition{breast}    \example{
      \exsen{llɨg ngam sinenang mälla}      \extran{breastfeeding woman}}}  \sense{\textbf{2.} 
    \definition{nine (lit. breast; body counting numeral)}    \example{
}    \example{
}    \example{
}}}

\entry{ngam indäb}{\headword{ngam indäb}  \note{Parts of the Body}  \pronnote{ŋam indəb}
  \pos{Noun}  \sense{
    \definition{nipple, teat}    \note{Parts of the Body}    \example{
}    \example{
}    \example{
}}}

\entry{ngam koll}{\headword{ngam koll}  \pronnote{ŋam koɽ}
  \pos{Noun}  \sense{
    \definition{breast}}}

\entry{ngäma}{\headword{ngäma}  \pronnote{ŋəma}
  \pos{Personal pronoun}  \sense{
    \definition{possessive form of ngämi}    \example{
      \exsen{Bäne ngämingg a ulle abal dan ngäma pate.}      \extran{Your help towards us is great.}}}}

\entry{ngämae}{\headword{ngämae}
  \pos{Intransitive verb}  \sense{
    \definition{to go around, take the long way}    \example{
      \exsen{Kurupel gungmaeyän ttongo nyongo dae dallän do auma.}      \extran{Kurupel went using another path to the grave.}}}}

\entry{ngämaene}{\headword{ngämaene}
  \pos{Personal pronoun}  \sense{
    \definition{ablative-possessive form of ngämi}    \example{
      \exsen{Ngäna bablle mɨnyi ngämaene ine kämbmenyatt eka de bɨllɨtne.}      \extran{I will be telling you our diving story.}}}}

\entry{ngämälle}{\headword{ngämälle}
  \variantof{SpellingVariant of \vartext{ngämlle}}}

\entry{ngämäne}{\headword{ngämäne}
  \variantof{SpellingVariant of \vartext{ngämene}}}

\entry{ngämangäma}{\headword{ngämangäma}
  \pos{Personal pronoun}  \sense{
}}

\entry{ngämar}{\headword{ngämar}
  \pos{Transitive verb}  \sense{
    \definition{to haul}    \example{
      \exsen{Mälla da wayati deyangmereyo.}      \extran{The women hauled over the watermelons.}}}}

\entry{ngämen}{\headword{ngämen}  \pronnote{ŋəmen}
  \pos{Transitive verb}  \sense{
    \definition{to reach, catch up to}    \example{
      \exsen{Ngäna bam mɨnyi bangmen.}      \extran{I will catch up to you.}}}}

\entry{ngämene}{\headword{ngämene}
  \pos{Personal pronoun}  \sense{
    \definition{ablative-possessive form of ngäna}    \example{
      \exsen{Bongo ngämene eka de malaem anggan.}      \extran{You obey my words.}}}
  \variants{\SpellingVariant \vartext{ngämäne}}
  \variants{\Fast Speech Variant \vartext{ngämne}}}

\entry{ngämengäme}{\headword{ngämengäme}  \pronnote{ŋəmeŋəme}
  \pos{Noun}  \sense{
    \definition{type of vine with sweet, edible fruit that are yellow when ripe}}}

\entry{ngämi}{\headword{ngämi}  \pronnote{ŋəmi}
  \pos{Personal pronoun}  \sense{
    \definition{we (first person nonsingular exclusive pronoun, nominative form)}    \example{
      \exsen{Ngämi mɨnyi bam batrameya.}      \extran{We will carry you.}}}}

\entry{ngämim}{\headword{ngämim}  \pronnote{ŋəmim}
  \pos{Personal pronoun}  \sense{
    \definition{accusative form of ngämi}    \example{
      \exsen{Abo bongo ngämim yangkollmällneg.}      \extran{You must follow us.}}}}

\entry{ngämingg}{\headword{ngämingg}  \pronnote{ŋəmiŋɡ}
  \pos{Intransitive verb}  \sense{\textbf{1.} 
    \definition{to help}    \example{
      \exsen{Lla da ttongo lla de nangminggan towall ttokoe e.}      \extran{A man helps another man to cut grass.}}    \example{
      \exsen{Ngäna bam mɨnyi bangmingg.}      \extran{I will help you.}}    \example{
      \exsen{Ibi mɨnyi iba mälla de bangmiminyaemnalla.}      \extran{We (incl.) will be supporting our women.}}    \example{
}}  \sense{\textbf{2.} 
    \definition{to answer}    \example{
      \exsen{Bogo do gungminggän ada, "Aoǃ"}      \extran{There he answered, "Yesǃ"}}}}

\entry{ngäminggag}{\headword{ngäminggag}  \note{good characteristics of people}  \pronnote{ŋəmiŋɡaɡ}
  \pos{Modifier}  \sense{
    \definition{helpful, supportive}    \note{good characteristics of people}    \example{
      \exsen{Bogo ngäminggag yae dan.}      \extran{She is a supportive mother.}}    \example{
}    \example{
}}}

\entry{ngämira}{\headword{ngämira}
  \pos{Personal pronoun}  \sense{
    \definition{dative form of ngämi}    \example{
      \exsen{Da ngämi e we bam baweseya, ngäma moko da nangesneg ngämira.}      \extran{Whatever we may ask of you, we want you to do it for us.}}}}

\entry{ngämlla}{\headword{ngämlla}
  \variantof{Dialectal Variant of \vartext{ngämlle}}}

\entry{ngämlle}{\headword{ngämlle}  \pronnote{ŋəmɽe}
  \pos{Personal pronoun}  \sense{
    \definition{dative form of ngäna}    \example{
      \exsen{Bäne däräng a daden, be ngämlle ddone dan.}      \extran{You have a dog, but I don't (lit. a dog exists to you, but not to me).}}}
  \variants{\SpellingVariant \vartext{ngämälle}}
  \variants{\Dialectal Variant \vartext{ngämlla}}}

\entry{ngämne}{\headword{ngämne}
  \variantof{Fast Speech Variant of \vartext{ngämene}}}

\entry{ngämo}{\headword{ngämo}  \pronnote{ŋəmo}
  \pos{Personal pronoun}  \sense{
    \definition{possessive form of ngäna}    \example{
      \exsen{Ngämo zegatt ttängäm a Malläm.}      \extran{My birthplace is Malam.}}}}

\entry{ngämral}{\headword{ngämral}  \pronnote{ŋəmral}
  \pos{Noun}  \sense{
    \definition{type of vine that grows in the bush and irritates skin}}}

\entry{ngamtep}{\headword{ngamtep}  \note{Mushrooms}  \pronnote{ŋamtep}
  \pos{Noun}  \sense{
    \definition{type of mushroom}    \note{Mushrooms}    \example{
}    \example{
}    \example{
}}}

\entry{ngän}{\headword{ngän}  \pronnote{ŋən}
  \pos{Noun}  \sense{
    \definition{boundary}    \example{
      \exsen{Ende tän bo ekaklle ngän}      \extran{the boundary of Ende lands}}}}

\entry{ngäna}{\headword{ngäna}
  \variantof{Unspecified Variant of \vartext{ngänaeka}}}

\entry{ngäna}{\headword{ngäna}  \pronnote{ŋəna}
  \pos{Personal pronoun}  \sense{
    \definition{I (first person singular pronoun, nominative form)}    \example{
      \exsen{Ngäna gotar.}      \extran{I slept.}}}}

\entry{ngäna~}{\headword{ngäna~}
  \variantof{freevarlab of \vartext{RED~}}}

\entry{nganae}{\headword{nganae}  \pronnote{ŋanae}
  \pos{Transitive verb}  \sense{\textbf{1.} 
    \definition{to coil, go around}    \example{
      \exsen{Mätta da dade de nganae eran.}      \extran{The yam is coiling around the yam stick.}}}  \sense{\textbf{2.} 
    \definition{to spin, rotate}    \example{
      \exsen{Ekaklle da nganaenen allan.}      \extran{The earth is spinning.}}}}

\entry{ngänaeka}{\headword{ngänaeka}  \note{Verbs of Motion}  \pronnote{ŋənaeka}
  \pos{Transitive coverb}  \sense{
    \definition{to cry (about, for)}    \note{Verbs of Motion}    \example{
      \exsen{Bongo ewatta ngänaeka alle?}      \extran{Why are you crying?}}    \example{
      \exsen{Bom ngänaeka dägagallo.}      \extran{They cry for him.}}}
  \variants{\Unspecified Variant \vartext{ngäna}}
  \variants{\Fast Speech Variant \vartext{naeka, ngaeka}}}

\entry{ngänaeka pollongg}{\headword{ngänaeka pollongg}  \pronnote{ŋənaeka poɽoŋɡ}
  \pos{Intransitive compound verb}  \sense{
    \definition{to burst into tears}    \example{
      \exsen{Bogo ngänaeka gompllogon.}      \extran{He burst into tears.}}}}

\entry{ngänaeka täräm}{\headword{ngänaeka täräm}
  \pos{Noun}  \sense{
    \definition{tear}}}

\entry{ngänaekangänaeka}{\headword{ngänaekangänaeka}  \pronnote{ŋənaekaŋənaeka}
  \pos{Intransitive coverb}  \sense{
    \definition{nonsingular form of ngänaeka}    \example{
      \exsen{Tämamae llɨg de papa dägnegnän, ngänaekangänaeka gognän.}      \extran{He hit all the boys and they cried.}}}}

\entry{ngänaekongänaekong}{\headword{ngänaekongänaekong}
  \pos{Adverb}  \sense{
    \definition{crying, in tears}}
  \variants{\Unspecified Variant \vartext{ngänangänong}}}

\entry{nganaenganae}{\headword{nganaenganae}  \pronnote{ŋanaeŋanae}
  \pos{Modifier}  \sense{
    \definition{wrapping around}}}

\entry{ngänäm}{\headword{ngänäm}  \pronnote{ŋənəm}
  \pos{Personal pronoun}  \sense{
    \definition{accusative form of ngäna}    \example{
      \exsen{Ngänäm kollmäll a mudan.}      \extran{Don't follow me.}}}}

\entry{ngänam}{\headword{ngänam}  \pronnote{ŋənam}
  \pos{Transitive verb}  \sense{\textbf{1.} 
    \definition{to understand, recognize}    \example{
      \exsen{Ubi mɨnyi eka midd de bängnameyo.}      \extran{They will understand the meaning.}}    \example{
      \exsen{Bibi ddone ke ngänaem amalla?}      \extran{Don't you all understand?}}    \example{
      \exsen{Bogo gongnamän ada, nagda obom era kuki dägagän.}      \extran{He recognized that his friend had deceived him.}}}  \sense{\textbf{2.} 
    \definition{understanding}}}

\entry{=ngänäm}{\headword{=ngänäm}  \pronnote{ŋənəm}
  \pos{Nominal enclitic}  \sense{
    \definition{similative case clitic; like}    \example{
      \exsen{Ttongo lla de dirom ngänäm zuwoe a mudan.}      \extran{Don't shoot another person like they are a cassowary.}}}}

\entry{ngänameny}{\headword{ngänameny}
  \pos{Modifier}  \sense{
    \definition{unrecognizable, obscure}    \example{
      \exsen{Ngaskäma Ende eka da mɨnyi ngänameny bogon källkae.}      \extran{Maybe the Ende language will become unrecognizable later.}}}
  \variants{\SpellingVariant \vartext{ngänammeny}}}

\entry{ngänammeny}{\headword{ngänammeny}
  \variantof{SpellingVariant of \vartext{ngänameny}}}

\entry{ngänangänong}{\headword{ngänangänong}
  \variantof{Unspecified Variant of \vartext{ngänaekongänaekong}}}

\entry{ngänaurur}{\headword{ngänaurur}  \pronnote{ŋənaurur}
  \pos{Intransitive coverb}  \sense{
    \definition{to mourn}}}

\entry{ngänawa}{\headword{ngänawa}  \pronnote{ŋənawa}
  \pos{Personal pronoun}  \sense{
    \definition{emphatic form of ngäna}    \example{
      \exsen{Ngänawa lla de umllang anggan.}      \extran{It is I that tells people.}}}}

\entry{ngänawaeben}{\headword{ngänawaeben}
  \pos{Copular verb}  \sense{
    \definition{restrictive copular form of ngäna (present form)}    \example{
      \exsen{Ngäna ngämo känyer ngämo pemliangae ngänawaeben ngämo ma me.}      \extran{It is only me with my family alone in my house.}}}}

\entry{ngänawaenen}{\headword{ngänawaenen}  \pronnote{ŋənawaenen}
  \pos{Copular verb}  \sense{
    \definition{copular form of ngäna (present form)}    \example{
      \exsen{Gänyaolle gall zizag a ngänawaenen.}      \extran{The owner of this canoe is me.}}}}

\entry{ngängäll käm}{\headword{ngängäll käm}  \note{Parts of the Body}  \pronnote{ŋəŋəɽ kəm}
  \pos{Noun}  \sense{
    \definition{body part}    \note{Parts of the Body}    \example{
}    \example{
}    \example{
}}}

\entry{ngange}{\headword{ngange}  \pronnote{ŋaŋe}
  \pos{Transitive verb}  \sense{\textbf{1.} 
    \definition{to communicate, deliver a message}    \example{
      \exsen{ngangema ttoen}      \extran{method of communication}}    \example{
      \exsen{Lla da mɨnyi utt alle eka de bängaeyo.}      \extran{Men will deliver the message with the conch shell.}}}  \sense{\textbf{2.} 
}}

\entry{ngangem}{\headword{ngangem}
  \pos{Transitive verb}  \sense{
    \definition{to adopt}    \example{
      \exsen{Ngäna ge män de obaene de dangem.}      \extran{I adopted this girl from them.}}}}

\entry{ngangeyang}{\headword{ngangeyang}  \pronnote{ŋaŋejaŋ}
  \pos{Noun}  \sense{
    \definition{messenger}    \example{
      \exsen{Ngangeyang abo umllang dägaeballo aeya kuddäll agan.}      \extran{Then the messengers will say who died.}}}}

\entry{ngänglam}{\headword{ngänglam}  \note{Animal verbs}  \pronnote{ŋəŋlam}
  \pos{Intransitive verb}  \sense{
    \definition{to buzz}    \note{Animal verbs}    \example{
      \exsen{Ngänglamang dan.}      \extran{It's buzzing.}}    \example{
}    \example{
}}}

\entry{ngangleb}{\headword{ngangleb}  \pronnote{ŋaŋleb}
  \pos{Transitive verb}  \sense{
    \definition{to look for, search for}    \example{
      \exsen{Lla da mälla de dängalbänän.}      \extran{The man searched for the woman.}}}}

\entry{ngänngän}{\headword{ngänngän}  \pronnote{ŋənŋən}
  \pos{Intransitive verb}  \sense{
    \definition{to swear}}}

\entry{ngänongg}{\headword{ngänongg}
  \variantof{Unspecified Variant of \vartext{ngonongg}}  \note{Swadesh}  \pronnote{ŋənoŋɡ}}

\entry{ngänttäg}{\headword{ngänttäg}  \pronnote{ŋənʈ͡ʂəɡ}
  \pos{Intransitive verb}  \sense{\textbf{1.} 
    \definition{to arrive, return}    \example{
      \exsen{Ge lla ngänttäg allan ttongo ttängäm atta kilikiliang.}      \extran{This man is returning from another place happily.}}    \example{
      \exsen{Ngämo näkäp e gongttägän.}      \extran{It came to my mind.}}    \example{
      \exsen{Mɨnyi ddob sisor eka da ako bälltaemnän.}      \extran{Then other new languages will arrive.}}}  \sense{\textbf{2.} 
    \definition{to bring, carry}    \example{
      \exsen{Llɨg kälsre de ine da do ddob llayaba pate dängttägän.}      \extran{The water carried the boy over to some people over there.}}}}

\entry{ngänyanyemmeny}{\headword{ngänyanyemmeny}  \pronnote{ŋəɲaɲemmeɲ}
  \pos{Modifier}  \sense{
    \definition{mean, unkind}    \example{
}    \example{
}}}

\entry{ngänyngäny}{\headword{ngänyngäny}  \pronnote{ŋəɲŋəɲ}
  \pos{Transitive verb}  \sense{
    \definition{to swallow}    \example{
      \exsen{Amtet de nängänymällneg!}      \extran{Take (lit. swallow) a breath!}}}}

\entry{nganzig}{\headword{nganzig}  \pronnote{ŋanziɡ}
  \pos{Transitive verb}  \sense{
    \definition{to pass, overtake}    \example{
      \exsen{Ddia da angde inu kuddäll gogon, kottllam a deyanzigän.}      \extran{While Deer was dead sleep, Tortoise passed right by him.}}}}

\entry{Ngao}{\headword{Ngao}  \note{Place names}  \pronnote{ŋao}
  \pos{Proper noun}  \sense{
    \definition{Ngao (toponym)}    \note{Place names}    \example{
}    \example{
}    \example{
}}}

\entry{ngarängg}{\headword{ngarängg}  \pronnote{'ngarängg}
  \pos{Transitive verb}  \sense{
    \definition{to encounter, meet, run into}    \example{
      \exsen{Da ngäna lla de bangerängg nyongo me, mɨnyi Ende eka walle eka bäntameny.}      \extran{If I encounter someone on the street, I will speak in Ende to them.}}    \example{
      \exsen{Gänya ngata me ddädddäg de ddone dangerängg.}      \extran{I didn't encounter any animals at that spot.}}}}

\entry{ngas~}{\headword{ngas~}
  \variantof{Infinitival reduplicant of \vartext{RED~}}}

\entry{ngäs}{\headword{ngäs}  \pronnote{ŋəs}
  \pos{Intransitive verb}  \sense{\textbf{1.} 
    \definition{to return, come back}    \example{
      \exsen{Käsre llɨg a ngänangänong ma we dängäsmällän.}      \extran{Then the boys went home crying.}}    \example{
      \exsen{Bol a angosan.}      \extran{The ball came back.}}}  \sense{\textbf{2.} 
    \definition{to return, bring back}    \example{
      \exsen{Joshua däbe ttägäll käp de deyangosän ma we.}      \extran{Joshua brought that money back home.}}}}

\entry{ngäs~}{\headword{ngäs~}
  \variantof{Inflected Form of \vartext{RED~}}}

\entry{ngasa}{\headword{ngasa}
  \variantof{Fast Speech Variant of \vartext{ngasekäma}}}

\entry{ngäsangngäsang}{\headword{ngäsangngäsang}
  \pos{Adverb}  \sense{
    \definition{repeatedly}    \example{
      \exsen{Kaptte de kotang de ngäsangngäsang mättnan a mudan.}      \extran{Don't wear dirty clothes over and over.}}}}

\entry{ngase}{\headword{ngase}
  \pos{TAM particle}  \sense{
    \definition{hortative/optative particle}    \example{
      \exsen{Ngase ibi bobäll bem de apte menae e baupeya.}      \extran{Let us go over to the other side of the sea.}}    \example{
      \exsen{Ngase Ende eka de mɨnyi mokowang bägawän.}      \extran{Let them love to speak Ende.}}}}

\entry{ngäse~}{\headword{ngäse~}
  \variantof{Inflected Form of \vartext{RED~}}}

\entry{ngaseka}{\headword{ngaseka}
  \variantof{Fast Speech Variant of \vartext{ngasekäma}}}

\entry{ngasekäma}{\headword{ngasekäma}
  \pos{TAM particle}  \sense{\textbf{1.} 
    \definition{potential marker; may, could, might}    \example{
      \exsen{Erame dan ngasekäma?}      \extran{Where could they be?}}    \example{
      \exsen{Da ttongo lla da bäne baba bälle 1000 mani bonttogän, ende bäne baba ngasekäma bällädän?}      \extran{If someone gave your father 1000 kina, what might he buy?}}    \example{
}    \example{
}}  \sense{\textbf{2.} 
    \definition{maybe}    \example{
      \exsen{Mɨnyi ngaska llame duwem bognegnän oba peyang.}      \extran{Maybe they will eat together with them.}}}
  \variants{\Fast Speech Variant \vartext{käma, ngaseka, ngaskäma, aska, asa, ngasa, ngaskma, ngaska}}}

\entry{ngäsengäse}{\headword{ngäsengäse}  \pronnote{ŋəseŋəse}
  \pos{Adverb}  \sense{
    \definition{back, the other way}    \example{
      \exsen{Ubi komlla gopnaeyo a ngäsengäse dindugiyu.}      \extran{Those two turned around and ran back.}}}}

\entry{ngaska}{\headword{ngaska}
  \variantof{Fast Speech Variant of \vartext{ngasekäma}}}

\entry{ngaskäma}{\headword{ngaskäma}
  \variantof{Fast Speech Variant of \vartext{ngasekäma}}  \pronnote{ŋaskəma}}

\entry{ngaskma}{\headword{ngaskma}
  \variantof{Fast Speech Variant of \vartext{ngasekäma}}}

\entry{ngasnges}{\headword{ngasnges}
  \pos{Transitive verb}  \sense{\textbf{1.} 
    \definition{to do}    \example{
      \exsen{Ende bongo ngasnges eralle ttongo lla bo pate, bongo ngasnges eralle God bo pate.}      \extran{What you do to another person, you do to God.}}    \example{
      \exsen{Bongo bäne kumuddäga melem de nangesneg.}      \extran{Do your three chores.}}}  \sense{\textbf{2.} 
    \definition{to make}    \example{
      \exsen{Ibi ibra mɨnyi ttongo sisor bikwem de bangeseya.}      \extran{We (excl.) will make a new fireplace for us.}}    \example{
      \exsen{Masar ine kube de ngämo nangesan.}      \extran{Grandfather made my water bucket.}}}  \sense{\textbf{3.} 
    \definition{to happen}    \example{
      \exsen{Ge ttoen a gongesän ngämo patme.}      \extran{This thing happened to me.}}}}

\entry{ngata}{\headword{ngata}  \pronnote{ŋata}
  \pos{Noun}  \sense{
    \definition{spot}    \example{
      \exsen{Ge ttongo ngata me dädme gongkäbmenyne ako.}      \extran{We dived again at this other spot.}}}}

\entry{ngata~}{\headword{ngata~}
  \variantof{Derived Variant of \vartext{RED~}}}

\entry{ngata me}{\headword{ngata me}
  \pos{Adverb}  \sense{
    \definition{around, about, approximately}    \example{
      \exsen{Ngaska sikstis, däba ngata me ngäna gozeg.}      \extran{Maybe sometime around the sixties I was born.}}    \example{
      \exsen{Sɨmell a tu taosen namba ngata me dagaeya.}      \extran{There were about two thousand pigs.}}}}

\entry{ngatae}{\headword{ngatae}
  \pos{Transitive verb}  \sense{
    \definition{to visit}    \example{
      \exsen{Ngämi mɨnyi ngataenen e buddällmam.}      \extran{We (excl.) will come to visit.}}}}

\entry{ngätae}{\headword{ngätae}
  \pos{Transitive verb}  \sense{
    \definition{to check}    \example{
      \exsen{Lla da mɨnyi ngätaenen e ballnän.}      \extran{The man will be going to check.}}}}

\entry{ngatangata}{\headword{ngatangata}
  \pos{Noun}  \sense{
    \definition{small space}    \example{
      \exsen{Ngäna lla yaba ngatangata dae ibi allan.}      \extran{I am going through a tight space crowded with people.}}}}

\entry{ngatengate}{\headword{ngatengate}  \note{Animals/Plants}  \pronnote{ŋateŋate}
  \pos{Noun}  \sense{
    \definition{gliding possum, sugar glider}    \scientificname{Petaurus breviceps flavidus}    \note{Animals/Plants}    \example{
      \exsen{Ngatengate da biratt a ddobae moko dan.}      \extran{Roasted possum is very tasty.}}    \example{
}    \example{
}}}

\entry{ngätt}{\headword{ngätt}  \pronnote{ŋəʈ͡ʂ}
  \pos{Noun}  \sense{
    \definition{yard}    \example{
      \exsen{Däräng a yuwog ngätt me källa bebeyag dan.}      \extran{The dog takes craps in many yards.}}}}

\entry{ngatt}{\headword{ngatt}  \note{Trees}  \pronnote{ŋaʈ͡ʂ}
  \pos{Noun}  \sense{
    \definition{type of tree}    \note{Trees}    \example{
}    \example{
}    \example{
}}}

\entry{ngättäma}{\headword{ngättäma}  \pronnote{ŋəʈ͡ʂəma}
  \pos{Noun}  \sense{
    \definition{place, spot}    \example{
      \exsen{tudi ma ngättma}      \extran{place for fishing}}    \example{
      \exsen{Däba ngättäma we gotarän.}      \extran{In that place, he slept.}}}
  \variants{\Fast Speech Variant \vartext{ngättma}}}

\entry{ngättangätta}{\headword{ngättangätta}  \pronnote{ŋəʈ͡ʂaŋəʈ͡ʂa}
  \pos{Transitive verb}  \sense{\textbf{1.} 
    \definition{to block, obscure}    \example{
      \exsen{Bongo ddone ikop dagne, llo da ngänäm dangttawalle.}      \extran{You couldn't see me; the tree was obscuring me.}}}  \sense{\textbf{2.} 
    \definition{to occupy, take over}    \example{
      \exsen{Särem a dangttanän ttängäm de.}      \extran{Darkness took over the village.}}}}

\entry{ngättangätte}{\headword{ngättangätte}  \note{School}  \pronnote{ŋəʈ͡ʂaŋəʈ͡ʂe}
  \pos{Transitive verb}  \sense{\textbf{1.} 
    \definition{to count}    \note{School}    \example{
      \exsen{Ingglis eka walle dangättnealle.}      \extran{I was counting in English.}}    \example{
}    \example{
}}  \sense{\textbf{2.} 
    \definition{to read}    \note{School}}}

\entry{ngättma}{\headword{ngättma}
  \variantof{Fast Speech Variant of \vartext{ngättäma}}  \pronnote{ŋəʈ͡ʂma}}

\entry{ngattong}{\headword{ngattong}  \pronnote{ŋaʈ͡ʂoŋ}
  \pos{Ordinal numeral}  \sense{\textbf{1.} 
    \definition{first}    \example{
      \exsen{ngattong e ebdo me}      \extran{on the first day}}    \example{
      \exsen{Ge obo ngattong täräp dan.}      \extran{This is his first time.}}}  \sense{\textbf{2.} 
    \definition{first, at first, initially, previously, before}    \example{
      \exsen{Ngäna ngattong botogol.}      \extran{I will hide first.}}    \example{
      \exsen{Ngattong skulmeny dagaeya.}      \extran{Before, there was no school.}}}  \sense{\textbf{3.} 
    \definition{beginning; past}    \example{
      \exsen{ngattong att otät yu ngalen}      \extran{cooking methods from the past}}    \example{
      \exsen{Ngattong att gagäll ttoen de bälnan amallo.}      \extran{They're remembering bad things from before.}}}  \sense{\textbf{4.} 
    \definition{front}    \example{
      \exsen{Gall e ngämi godagaleya, ngäna imne we, ede Imanuel ngattong e godmenän.}      \extran{We (excl.) boarded the canoe with me in the back; Imanuel sat in front.}}    \example{
      \exsen{Bogo tubutubu gogon obo ngattong alle.}      \extran{He kneeled before him.}}}}

\entry{ngattong mullmull}{\headword{ngattong mullmull}  \pronnote{ŋaʈ͡ʂoŋ muɽmuɽ}
  \pos{Noun}  \sense{
    \definition{the very beginning}    \example{
      \exsen{ngattong mullmull e}      \extran{at the very beginning}}    \example{
      \exsen{Ngäna mɨnyi bongkam ngattong mullmull abal ttängäm de alla ingollang ngasnges erallo.}      \extran{I will start from the very beginning and show how we make a garden.}}}}

\entry{ngattongattong}{\headword{ngattongattong}
  \pos{Modifier}  \sense{
    \definition{very beginning, very first}    \example{
      \exsen{ngattongattong täräp me}      \extran{in the very beginning}}}}

\entry{ngattongattong}{\headword{ngattongattong}
  \pos{Adverb}  \sense{
    \definition{ahead, in front}    \example{
      \exsen{Ddob imneimne dag a ddob ngattongattong dag.}      \extran{Some are behind and some are ahead.}}    \example{
      \exsen{Obo ngattong täräp daeya ngattongattong ibnin e.}      \extran{It was his first time going in front.}}}}

\entry{nge}{\headword{nge}  \lexeme{nge}  \pronnote{ŋe}
  \pos{Noun}  \sense{\textbf{1.} 
    \definition{coconut palm}    \scientificname{Cocos nucifera}    \example{
      \exsen{Ngäna mulldae dan ge nge de bängkäll.}      \extran{I can climb this coconut palm.}}}  \sense{\textbf{2.} 
    \definition{coconut}    \example{
      \exsen{nge däg}      \extran{coconut bunch}}    \example{
      \exsen{Ngämi nge de diwenyeya.}      \extran{We (excl.) brought over the coconut.}}}}

\entry{nge dara dara pite}{\headword{nge dara dara pite}  \pronnote{ŋe dara dara pite}
  \pos{Noun}  \sense{
    \definition{coconut grass skirt}    \example{
}    \example{
}    \example{
}}}

\entry{nge dɨdɨr}{\headword{nge dɨdɨr}  \pronnote{ŋe dɪdɪr}
  \pos{Noun}  \sense{\textbf{1.} 
    \definition{fully dry coconut}    \example{
      \exsen{Nge dɨdɨr bo ngasnen a kemibi dag, ada otät sespen a, mubine da, kul ngasnen a ako ddob kemibi ttoen a.}      \extran{Dry coconut has many uses: you can boil it, cook it in a cake, making balls, and few other things.}}}  \sense{\textbf{2.} 
    \definition{the eleventh and final stage of coconut growth in which the fruit is completely dry and brown}}}

\entry{nge ibeny}{\headword{nge ibeny}  \note{Conflict/war}  \pronnote{ŋe ibeɲ}
  \pos{Noun}  \sense{
    \definition{planting a coconut as a gesture of peace}    \note{Conflict/war}    \example{
}    \example{
}    \example{
}}}

\entry{nge id}{\headword{nge id}  \pronnote{ŋe id}
  \pos{Noun}  \sense{
    \definition{coconut cream}    \example{
}    \example{
}    \example{
}}}

\entry{nge ine}{\headword{nge ine}  \lexeme{nge ine}  \pronnote{ŋe ine}
  \pos{Noun}  \sense{
    \definition{coconut water}}}

\entry{nge pollgo}{\headword{nge pollgo}  \lexeme{nge pollgo}  \pronnote{ŋe poɽɡo}
  \pos{Noun}  \sense{
    \definition{coconut frog, green}}}

\entry{nge popo}{\headword{nge popo}  \pronnote{ŋe popo}
  \pos{Noun}  \sense{\textbf{1.} 
    \definition{coconut flower}    \example{
      \exsen{Niki tämamae apte ttang llɨpɨt pa de dägäddnegän nge popo me.}      \extran{Niki killed all five birds in the coconut flowers.}}}  \sense{\textbf{2.} 
    \definition{the fourth stage of coconut growth in which flowers appear}}}

\entry{nge tärmir}{\headword{nge tärmir}  \note{Snakes}  \pronnote{ŋe tərmir}
  \pos{Noun}  \sense{
    \definition{type of snake}    \note{Snakes}    \example{
}    \example{
}    \example{
}}}

\entry{nge tikop}{\headword{nge tikop}  \pronnote{ŋe tikop}
  \pos{Noun}  \sense{\textbf{1.} 
    \definition{tiny coconut, baby coconut}    \example{
      \exsen{Ddob lla da gae nge tikop de kaepnen amallo.}      \extran{Some people chew the baby coconuts of the gae palm.}}}  \sense{\textbf{2.} 
    \definition{the sixth stage of coconut growth in which the fruits have a diameter of ~5 cm}}}

\entry{nge ttam}{\headword{nge ttam}  \lexeme{nge ttam}  \pronnote{ŋe ʈ͡ʂam}
  \pos{Noun}  \sense{
    \definition{coconut frond}}}

\entry{Ngeba}{\headword{Ngeba}  \pronnote{ŋeba}
  \pos{Proper noun}  \sense{
    \definition{Ngeba (toponym)}}}

\entry{ngem}{\headword{ngem}  \note{wrong species name? native to North America}  \pronnote{ŋem}
  \pos{Noun}  \sense{
    \definition{virginia oppossum}    \note{wrong species name? native to North America}    \example{
}    \example{
}    \example{
}}}

\entry{Ngerbab}{\headword{Ngerbab}  \pronnote{ŋerbab}
  \pos{Proper noun}  \sense{
    \definition{male personal name}}}

\entry{ngetae}{\headword{ngetae}  \pronnote{ŋetae}
  \pos{Transitive verb}  \sense{\textbf{1.} 
    \definition{to touch}    \example{
      \exsen{Ttongo lla da mälla da walle nägabällallo ddɨg alle angeteallo.}      \extran{A man was standing with his wife with their backs touching.}}}  \sense{\textbf{2.} 
    \definition{to get engaged}    \example{
      \exsen{Lla da medäda wa mägda de wanseg yaran a obo mälla peyang ngetae allan.}      \extran{The man leaves his father and mother and gets engaged to his woman.}}}  \sense{\textbf{3.} 
    \definition{to unite, join}    \example{
      \exsen{Lla da ddone mɨnyi sapasapang beyagän Adi bäne enanae ngetaeatt de.}      \extran{Man shall not separate what God has joined together for good.}}}  \sense{\textbf{4.} 
    \definition{to grip, hold onto}    \example{
      \exsen{Ngäma ttang de deyangeteyo.}      \extran{They gripped our (excl.) hands.}}}}

\entry{ngetam}{\headword{ngetam}  \pronnote{ŋetam}
  \pos{Intransitive verb}  \sense{
    \definition{to come down, descend}    \example{
      \exsen{Angetam.}      \extran{[You] come down.}}    \example{
      \exsen{Baet gongetamän matu we.}      \extran{The cuscus came down to the ground.}}}}

\entry{ngleg}{\headword{ngleg}  \pronnote{ŋleɡ}
  \pos{Transitive verb}  \sense{
    \definition{to clear the ground before building}    \example{
      \exsen{Banglegeya.}      \extran{Let's clear the ground.}}}}

\entry{ngllongg}{\headword{ngllongg}
  \pos{Transitive verb}  \sense{
    \definition{to forget}    \example{
      \exsen{Ubi gongllomenynegän otät kongkom e.}      \extran{They forgot to bring food.}}}}

\entry{ngoe}{\headword{ngoe}
  \variantof{SpellingVariant of \vartext{ngoi}}}

\entry{ngoeang}{\headword{ngoeang}  \pronnote{ŋoeaŋ}
  \pos{Noun}  \sense{
    \definition{traditional Y-shaped house post (used prior to nails)}}}

\entry{ngoetrangoetra}{\headword{ngoetrangoetra}  \lexeme{ngoetrangoetra}  \pronnote{ŋoetraŋoetra}
  \pos{Noun}  \sense{
    \definition{type of ant that doesn't bite}}}

\entry{ngoi}{\headword{ngoi}  \pronnote{ŋoi}
  \pos{Noun}  \sense{
    \definition{tooth}    \example{
}    \example{
}}
  \variants{\SpellingVariant \vartext{ngoe}}}

\entry{ngoi}{\headword{ngoi}  \pronnote{ŋoi}
  \pos{Transitive coverb}  \sense{
    \definition{to crowd, surround}    \example{
      \exsen{Ngoi dägaebeyo tämamae mänmän de.}      \extran{They surrounded all the girls.}}}}

\entry{ngok}{\headword{ngok}  \pronnote{ŋok}
  \pos{Ideophone}  \sense{
    \definition{oink, snort}}
  \variants{\Unspecified Variant \vartext{ngak}}}

\entry{ngokngok}{\headword{ngokngok}  \pronnote{ŋokŋok}
  \pos{Noun}  \sense{
    \definition{boobook (Australian, barking [barking owl])}    \scientificname{Ninox boobook pusilla, N. connivens}    \example{
}    \example{
}    \example{
}}}

\entry{ngolloe}{\headword{ngolloe}
  \pos{Transitive verb}  \sense{
    \definition{to look up at}    \example{
      \exsen{Gugiag giddollnen alla ngonomeny allan tuk i ikop allan kok de ngolloenen eran.}      \extran{He's standing staying, thinking, looking up, looking at the moon.}}}}

\entry{ngollot}{\headword{ngollot}
  \pos{Transitive verb}  \sense{
    \definition{to make burst, make explode}    \example{
      \exsen{Ddol a nyäng de bungolltän.}      \extran{The gas made the bag explode.}}}}

\entry{ngolo}{\headword{ngolo}  \lexeme{ngolo}  \pronnote{ŋolo}
  \pos{Noun}  \sense{
    \definition{type of big yam with a purple interior, thorns, and hairs}}}

\entry{ngon}{\headword{ngon}  \pronnote{ŋon}
  \pos{Intransitive coverb}  \sense{
    \definition{to walk in a stylish fashion, to strut}}}

\entry{ngonenngonen duwem}{\headword{ngonenngonen duwem}  \pronnote{ŋonenŋonen duwem}
  \pos{Transitive/Intransitive coverb}  \sense{
    \definition{to scarf down, eat very quickly}}}

\entry{ngongo}{\headword{ngongo}  \pronnote{ŋoŋo}
  \pos{Transitive verb}  \sense{
    \definition{to smooth, sand (a surface)}    \example{
      \exsen{Turik de nongowan.}      \extran{I smoothed the axe.}}}}

\entry{ngongop}{\headword{ngongop}  \note{Verbs}  \pronnote{ŋoŋop}
  \pos{Transitive verb}  \sense{
    \definition{to hug, embrace}    \note{Verbs}    \example{
      \exsen{Mäg da llɨg kälsre de nongopan.}      \extran{The mother hugged her little boy.}}    \example{
}}}

\entry{ngonoe}{\headword{ngonoe}  \pronnote{ŋonoe}
  \pos{Transitive verb}  \sense{
    \definition{to ask}    \example{
      \exsen{Llamda da abo ubim ngonoenen de deyangkamän.}      \extran{The old man started to question those two.}}    \example{
      \exsen{Ubi obaoba gongnoemenynegnän ada, "Ge endan?"}      \extran{They were asking themselves, "What is this?"}}}}

\entry{ngonoe eka}{\headword{ngonoe eka}
  \pos{Noun}  \sense{
    \definition{question}    \example{
      \exsen{Daballe ttongo ngonoe eka de nangnoeyan.}      \extran{After that, she asked another question.}}}}

\entry{ngonongg}{\headword{ngonongg}  \pronnote{ŋonoŋɡ}
  \pos{Intransitive verb}  \sense{\textbf{1.} 
    \definition{to think}    \example{
      \exsen{Bogo ada gongnomenynän ada ka däräng a pollnenang me erallo.}      \extran{He was thinking that the dogs were barking.}}    \example{
      \exsen{Emi ulle abal gongnomenyän.}      \extran{Emi was thinking hard.}}}  \sense{\textbf{2.} 
    \definition{to think about, think of, recall, remember}    \example{
      \exsen{Ngäna bam angnonggan.}      \extran{I thought of you.}}    \example{
      \exsen{Eka de tämamae ngäna era ddone ngonomeny allan.}      \extran{I can't recall all the words.}}}
  \variants{\Unspecified Variant \vartext{ngänongg}}}

\entry{ngowangowe}{\headword{ngowangowe}  \note{Adjectives}  \pronnote{ŋowaŋowe}
  \pos{Modifier}  \sense{
    \definition{narrow}    \note{Adjectives}    \example{
      \exsen{Ge walle da ngowangowe dan gall ibi wi.}      \extran{The river is too narrow for the boat to pass.}}    \example{
}    \example{
}}}

\entry{Nicky}{\headword{Nicky}
  \variantof{SpellingVariant of \vartext{Niki}}  \pronnote{niki}}

\entry{Nicole}{\headword{Nicole}
  \variantof{SpellingVariant of \vartext{Nikol}}}

\entry{nidol}{\headword{nidol}
  \pos{Noun}  \sense{
    \definition{needle}}}

\entry{nikbin}{\headword{nikbin}
  \pos{Noun}  \sense{
    \definition{nickname}}}

\entry{Niki}{\headword{Niki}  \pronnote{niki}
  \pos{Proper noun}  \sense{
    \definition{male personal name}}
  \variants{\SpellingVariant \vartext{Nicky}}}

\entry{Nikol}{\headword{Nikol}  \pronnote{nikol}
  \pos{Proper noun}  \sense{
    \definition{female personal name}}
  \variants{\SpellingVariant \vartext{Nicole}}}

\entry{nil}{\headword{nil}  \pronnote{nil}
  \pos{Noun}  \sense{
    \definition{nail (made of metal)}}}

\entry{ninety}{\headword{ninety}
  \variantof{SpellingVariant of \vartext{naenti}}}

\entry{nineyem}{\headword{nineyem}  \pronnote{ninejem}
  \pos{Noun}  \sense{
    \definition{whisper}    \example{
      \exsen{Angde llamda da togolag dagirnän, llayabaene eka nineyem de dandärän.}      \extran{While the old man was hiding, he heard whispers from people.}}}}

\entry{Niniab}{\headword{Niniab}  \pronnote{niniab}
  \pos{Proper noun}  \sense{
    \definition{male personal name}}}

\entry{Nixon}{\headword{Nixon}
  \pos{Proper noun}  \sense{
    \definition{male personal name}}}

\entry{no}{\headword{no}
  \variantof{Dialectal Variant of \vartext{English loanword}}}

\entry{Noar}{\headword{Noar}  \pronnote{noar}
  \pos{Proper noun}  \sense{
    \definition{female personal name}}
  \variants{\SpellingVariant \vartext{Nowar}}}

\entry{Nogat}{\headword{Nogat}  \pronnote{noɡat}
  \pos{Proper noun}  \sense{
    \definition{male personal name}}}

\entry{Nolin}{\headword{Nolin}  \pronnote{nolin}
  \pos{Proper noun}  \sense{
    \definition{female personal name}}}

\entry{nongg}{\headword{nongg}  \lexeme{nongg}  \pronnote{noŋɡ}
  \pos{Noun}  \sense{
    \definition{type of gecko}    \example{
      \exsen{Ngäma ma me nongg a kemibi abal dag.}      \extran{In my house there are many geckos.}}}}

\entry{nongonongor}{\headword{nongonongor}  \pronnote{noŋonoŋor}
  \pos{Modifier}  \sense{
    \definition{itchy}    \example{
      \exsen{Mälläng a nongonongorang gogon, käsre ansi gogon.}      \extran{His nose got itchy; then he sneezed.}}}}

\entry{nono}{\headword{nono}  \pronnote{nono}
  \pos{Transitive coverb}  \sense{
    \definition{to crumble, make soft}}}

\entry{nop mu}{\headword{nop mu}  \note{Marriage}  \pronnote{nop mu}
  \pos{Noun}  \sense{
    \definition{inter-tribe payment}    \note{Marriage}    \example{
}    \example{
}    \example{
}}}

\entry{Nope}{\headword{Nope}
  \pos{Proper noun}  \sense{
    \definition{male personal name}}}

\entry{nopmäg}{\headword{nopmäg}
  \pos{Kinship noun}  \sense{
    \definition{one's mother's exchange sister}}}

\entry{noponda}{\headword{noponda}  \note{Family ties}  \pronnote{noponda}
  \pos{Kinship noun}  \sense{
    \definition{one's father's exchange brother (still owes him)}    \note{Family ties}    \example{
}    \example{
}    \example{
}}}

\entry{nora}{\headword{nora}  \note{Trees}  \pronnote{nora}
  \pos{Noun}  \sense{
    \definition{type of big introduced tree with red flowers that attract birds}    \note{Trees}    \example{
}    \example{
}    \example{
}}}

\entry{Norma}{\headword{Norma}  \pronnote{norma}
  \pos{Proper noun}  \sense{
    \definition{female personal name}}}

\entry{not}{\headword{not}
  \variantof{freevarlab of \vartext{English loanword}}}

\entry{nothing}{\headword{nothing}
  \variantof{Dialectal Variant of \vartext{English loanword}}}

\entry{noto}{\headword{noto}
  \variantof{Dialectal Variant of \vartext{Taeme loanword}}}

\entry{now}{\headword{now}
  \variantof{Dialectal Variant of \vartext{English loanword}}}

\entry{Nowar}{\headword{Nowar}
  \variantof{SpellingVariant of \vartext{Noar}}  \pronnote{nowar}}

\entry{Nuam}{\headword{Nuam}
  \pos{Proper noun}  \sense{
    \definition{female personal name}}}

\entry{Nugini}{\headword{Nugini}
  \pos{Proper noun}  \sense{
    \definition{New Guinea}}}

\entry{Nuopin}{\headword{Nuopin}  \pronnote{nuopin}
  \pos{Proper noun}  \sense{
    \definition{female personal name}}}

\entry{nurae}{\headword{nurae}  \note{Birds}  \pronnote{nurae}
  \pos{Noun}  \sense{
    \definition{type of large bird}    \note{Birds}    \example{
}    \example{
}    \example{
}}}

\entry{nyägae}{\headword{nyägae}  \pronnote{ɲəɡae}
  \pos{Intransitive verb}  \sense{\textbf{1.} 
    \definition{to flail}    \example{
      \exsen{Guzi da nyägaenen allan.}      \extran{The crayfish was flailing around.}}}  \sense{\textbf{2.} 
    \definition{to stir}    \example{
      \exsen{Mälla da wätät de nanyägaenen spun alle.}      \extran{The woman was stirring the food with a spoon.}}    \example{
      \exsen{Danygaeaebneyo sana peyang.}      \extran{They were stirring them together with sago.}}}}

\entry{nyäkiyam}{\headword{nyäkiyam}
  \pos{Noun}  \sense{
}}

\entry{nyaknen}{\headword{nyaknen}
  \pos{Intransitive verb}  \sense{
}}

\entry{nyäkukub}{\headword{nyäkukub}  \lexeme{nyäkukub}  \pronnote{ɲəkukub}
  \pos{Noun}  \sense{
    \definition{sago washing basket}}}

\entry{nyämae~}{\headword{nyämae~}
  \variantof{freevarlab of \vartext{RED~}}}

\entry{nyämaenyämae}{\headword{nyämaenyämae}  \pronnote{ɲəmaeɲəmae}
  \pos{Transitive verb}  \sense{
    \definition{to surround}    \example{
      \exsen{Lla da obom danymaeyo.}      \extran{People surrounded him.}}}}

\entry{nyamäll}{\headword{nyamäll}  \pronnote{ɲaməɽ}
  \pos{Intransitive verb}  \sense{
    \definition{to be late}    \example{
      \exsen{Ubi dänyamällän.}      \extran{They were late.}}}}

\entry{nyamällatt}{\headword{nyamällatt}  \lexeme{nyamällatt}  \note{Family ties}  \pronnote{ɲaməɽaʈ͡ʂ}
  \pos{Kinship noun}  \sense{
    \definition{woman's exchange sibling's child}    \note{Family ties}}}

\entry{nyamnen}{\headword{nyamnen}
  \pos{Transitive verb}  \sense{
}}

\entry{nyan}{\headword{nyan}  \pronnote{ɲan}
  \pos{Noun}  \sense{
    \definition{type of tree with small leaves and bark used for mumu and topping a roof}}}

\entry{nyäng}{\headword{nyäng}  \pronnote{ɲəŋ}
  \pos{Intransitive verb}  \sense{
    \definition{to dance}    \example{
      \exsen{Dänyängnän obo peyang.}      \extran{They were dancing with him.}}}}

\entry{nyäng}{\headword{nyäng}  \pronnote{ɲəŋ}
  \pos{Noun}  \sense{
    \definition{basket; bag}    \example{
      \exsen{Nyäng a mätta alle bodo gogon.}      \extran{The bag became filled with yams.}}}}

\entry{nyäng tär}{\headword{nyäng tär}
  \pos{Noun}  \sense{
    \definition{woven strap or handle of a nyäng}}}

\entry{nyängkallbidd}{\headword{nyängkallbidd}  \pronnote{ɲəŋkaɽbiɖ͡ʐ}
  \pos{Noun}  \sense{
    \definition{type of bag}    \example{
      \exsen{Mätta komnen e nyängkallbidd a mer dan.}      \extran{Nyängkallbidd bags are good for carrying yams.}}}}

\entry{nyängkallddäb}{\headword{nyängkallddäb}  \note{Tools}  \pronnote{ɲəŋkaɽɖ͡ʐəb}
  \pos{Noun}  \sense{
    \definition{type of basket}    \note{Tools}    \example{
}    \example{
}    \example{
}}}

\entry{nyängoe}{\headword{nyängoe}  \pronnote{ɲəŋoe}
  \pos{Transitive verb}  \sense{
    \definition{to walk slowly with someone whose leg is hurt}    \example{
      \exsen{Danyängoeneyo gänye ma dälltam.}      \extran{We walked slowly and came up to the house.}}}}

\entry{nyanun}{\headword{nyanun}  \pronnote{ɲanun}
  \pos{Noun}  \sense{
    \definition{action}}}

\entry{nyäny}{\headword{nyäny}  \pronnote{ɲəɲ}
  \pos{Transitive verb}  \sense{\textbf{1.} 
    \definition{to paint}    \example{
      \exsen{Ubi pale alle nänyallo.}      \extran{They painted themselves with white clay.}}}  \sense{\textbf{2.} 
    \definition{to anoint}    \example{
      \exsen{Ubi mer mollong owel de dällädaebeyo obo pätt nyäny e.}      \extran{They bought fragrant oils for anointing his body.}}}}

\entry{nyänye}{\headword{nyänye}
  \pos{Transitive verb}  \sense{
    \definition{to share, split}    \example{
      \exsen{Ngämi dänyeya däbe ttägäll käp de.}      \extran{We (excl.) shared that money.}}    \example{
      \exsen{Otät de nyänan amallo.}      \extran{They're sharing the food.}}}}

\entry{nyanyem}{\headword{nyanyem}  \note{Things a person can do}  \pronnote{ɲaɲem}
  \pos{Transitive verb}  \sense{
    \definition{to respect}    \note{Things a person can do}    \example{
      \exsen{Bogo mälla de obo ddone mermerae nyanyem eran.}      \extran{He isn't properly respecting his wife.}}    \example{
      \exsen{Ngämo llɨg de mɨnyi banyemeyo.}      \extran{They will respect my son.}}}}

\entry{nyanyu}{\headword{nyanyu}  \note{Dancing}  \pronnote{ɲaɲu}
  \pos{Noun}  \sense{\textbf{1.} 
    \definition{action}    \note{Dancing}    \example{
      \exsen{Bogo gagäll nyanyu di ngasnges eran mällada pate.}      \extran{He's doing bad things to his wife.}}}  \sense{\textbf{2.} 
    \definition{to act, dramatize}    \note{Dancing}    \example{
      \exsen{Ge lla ekameny dan, be mɨnyi eka de ttang alle mae nanyunegneyo.}      \extran{This man is mute, but just his hands will act out the words.}}}}

\entry{nyärab}{\headword{nyärab}  \pronnote{ɲərab}
  \pos{Intransitive verb}  \sense{
    \definition{to go up, ascend}    \example{
      \exsen{Tutu gonyärabeya.}      \extran{We ascended the hill.}}}}

\entry{nyaraman}{\headword{nyaraman}  \note{Weather}  \pronnote{ɲaraman}
  \pos{Noun}  \sense{
    \definition{fine day}    \note{Weather}    \example{
}    \example{
}    \example{
}}}

\entry{nyärmeny}{\headword{nyärmeny}  \note{Conflict/war}  \pronnote{ɲərmeɲ}
  \pos{Noun}  \sense{
    \definition{little argument between children}    \note{Conflict/war}    \example{
}    \example{
}    \example{
}}}

\entry{nyäroe}{\headword{nyäroe}  \pronnote{ɲəroe}
  \pos{Intransitive verb}  \sense{
    \definition{to creep}    \example{
      \exsen{Angde ngäna ikop däga ada ddia dan, ngäna nyäroe de dängkam.}      \extran{When I saw the deer, I started to creep up.}}}}

\entry{nyärpae}{\headword{nyärpae}  \pronnote{ɲərpae}
  \pos{Intransitive verb}  \sense{
    \definition{to refuse (someone in the dative)}    \example{
      \exsen{Däbe män a oblle nyärpae ma da ddone mulldae gogon.}      \extran{That girl couldn't refuse him.}}}}

\entry{nyeny}{\headword{nyeny}  \lexeme{nyeny}  \pronnote{ɲeɲ}
  \pos{Noun}  \sense{
    \definition{type of large tree that grows along the river and in swamps with white flowers, brown fruit, and white bark used for making mumu and topping a roof}}}

\entry{nying}{\headword{nying}  \pronnote{ɲiŋ}
  \pos{Noun}  \sense{
    \definition{foot}    \example{
      \exsen{Baba bom iddob sana täkäll da nazunan nying me.}      \extran{Last night, a sago thorn got Dad in the foot.}}}}

\entry{nying kollop}{\headword{nying kollop}  \note{Clothing}  \pronnote{ɲiŋ koɽop}
  \pos{Noun}  \sense{
    \definition{sandal}    \note{Clothing}    \example{
}    \example{
}    \example{
}}}

\entry{nying lläbe}{\headword{nying lläbe}
  \pos{Noun}  \sense{
    \definition{toenail}}}

\entry{nying lläpät}{\headword{nying lläpät}
  \pos{Noun}  \sense{
    \definition{toe}}}

\entry{nying tongoe}{\headword{nying tongoe}  \pronnote{ɲiŋ toŋoe}
  \pos{Noun}  \sense{
    \definition{type of game involving kicking}}}

\entry{nying ttoe}{\headword{nying ttoe}
  \pos{Noun}  \sense{
    \definition{shoe}}}

\entry{nyinggulgul}{\headword{nyinggulgul}  \note{Birds}  \pronnote{ɲiŋɡulɡul}
  \pos{Noun}  \sense{
    \definition{type of bird}    \note{Birds}    \example{
}    \example{
}    \example{
}}}

\entry{nyingnying}{\headword{nyingnying}  \note{Gardens}  \pronnote{ɲiŋɲiŋ}
  \pos{Noun}  \sense{
    \definition{foundation of fence}    \note{Gardens}    \example{
}    \example{
}    \example{
}}}

\entry{nylib}{\headword{nylib}
  \pos{Noun}  \sense{
    \definition{type of arrow}}}

\entry{nyonga}{\headword{nyonga}  \pronnote{ɲoŋa}
  \pos{Noun}  \sense{
    \definition{Triton cockatoo}    \scientificname{Cacatua galerita triton}}}

\entry{nyongkoe}{\headword{nyongkoe}  \note{Verbs}  \pronnote{ɲoŋkoe}
  \pos{Transitive verb}  \sense{
    \definition{to pull}    \note{Verbs}    \example{
      \exsen{Angde tudi gognän, obo tudi di kollba ulle da danykoeyän.}      \extran{While fishing, a big fish pulled on her line.}}    \example{
}    \example{
}}}

\entry{nyongo}{\headword{nyongo}  \pronnote{ɲoŋo}
  \pos{Noun}  \sense{\textbf{1.} 
    \definition{road, path, way}    \example{
      \exsen{Gänya nyongo dae nalle.}      \extran{[You] take this road.}}}  \sense{\textbf{2.} 
    \definition{way, method}    \example{
      \exsen{Erodias nyongo de deyagnegnän Zon bälle gäz e.}      \extran{Herodias was searching for a way to kill John.}}}}

\entry{nyongo taempägag}{\headword{nyongo taempägag}
  \pos{Noun}  \sense{
    \definition{navigator (of a canoe)}    \anthronote{There are a variety of hand signs used by the navigator to communicate with the operator (gall onyang). These include changing the way you're facing to mean "veer right or left", to point left, right, or straight to indicate a turn or the lack of one and to speed up. A large wave of the hand indicates a big turn or a series of turns. A raised closed hand indicates to slow down or to stop.}}}

\entry{nyukukum}{\headword{nyukukum}  \note{Tools}  \pronnote{ɲukukum}
  \pos{Noun}  \sense{
    \definition{type of bag, used for squeezing sago, a thing woven from tree bark or reeds to wash sago or to fill with sago.}    \note{Tools}    \example{
      \exsen{Bäne aoli tot nyukukum dag?}      \extran{How many nyukukum bags do you have?}}    \example{
}    \example{
}}}

\chapter{O}


\entry{o}{\headword{o}  \pronnote{o}
  \pos{Modifier}  \sense{
    \definition{ripe}    \example{
      \exsen{Mätta ttam a o agnegan.}      \extran{The yam leaves will be ripe.}}}
  \variants{\SpellingVariant \vartext{wo}}}

\entry{o}{\headword{o}
  \pos{Interjection}  \sense{
    \definition{oh}}}

\entry{o-}{\headword{o-}  \pronnote{o}
  \pos{Intransitive verb}  \sense{
    \definition{root.ext.intr}}}

\entry{=o}{\headword{=o}  \pronnote{o}
  \pos{Discourse clitic}  \sense{
    \definition{vocative marker}}}

\entry{o}{\headword{o}  \pronnote{o}
  \pos{Coordinating connective}  \sense{
    \definition{or}    \example{
      \exsen{ik mi o upe me}      \extran{inside or outside}}}}

\entry{o klak}{\headword{o klak}
  \pos{Adverb}  \sense{
    \definition{o'clock}    \example{
      \exsen{Tri o klak gogän.}      \extran{It's three o'clock.}}}
  \variants{\SpellingVariant \vartext{o klok}}}

\entry{o klok}{\headword{o klok}
  \variantof{SpellingVariant of \vartext{o klak}}}

\entry{oba}{\headword{oba}
  \pos{Subordinating connective}  \sense{
    \definition{so that}    \example{
      \exsen{Bongo ine ma nalle oba bongo ine neyatan.}      \extran{You, go to the well so you can fetch water.}}}}

\entry{obä}{\headword{obä}
  \pos{Personal pronoun}  \sense{
    \definition{instrumental-comitative form of bogo}    \example{
      \exsen{Kukumi oballe beyareyo do Daru wi.}      \extran{He will go with Kukumi to Daru.}}}}

\entry{oba}{\headword{oba}  \pronnote{oba}
  \pos{Personal pronoun}  \sense{
    \definition{possessive form of ubi}    \example{
      \exsen{Oba mokwang tongoe a togotogol daeya.}      \extran{Their favorite game was hide-and-seek.}}}}

\entry{obaene}{\headword{obaene}
  \pos{Personal pronoun}  \sense{
    \definition{ablative-possessive form of ubi}    \example{
      \exsen{Llɨg obaene longgo de dandär.}      \extran{I heard noise from the children.}}}}

\entry{Obama}{\headword{Obama}
  \pos{Proper noun}  \sense{
    \definition{Obama (toponym)}}}

\entry{obäne}{\headword{obäne}
  \variantof{freevarlab of \vartext{obene}}}

\entry{obaoba}{\headword{obaoba}  \pronnote{obaoba}
  \pos{Personal pronoun}  \sense{\textbf{1.} 
    \definition{each other (reciprocal pronoun)}    \example{
      \exsen{Ubi obaoba tatu anggan.}      \extran{They are washing each other.}}    \example{
      \exsen{Llɨg a wa män a ubi obaoba umllang dagwaeya.}      \extran{The boy and the girl knew each other.}}}  \sense{\textbf{2.} 
    \definition{themselves (reflexive form of ubi)}    \example{
      \exsen{Ubi obaoba tatu anggan.}      \extran{They are washing themselves.}}    \example{
      \exsen{Llɨg a obaoba gognegän.}      \extran{The boys were by themselves.}}}}

\entry{obazaga}{\headword{obazaga}
  \pos{Personal pronoun}  \sense{
    \definition{themselves (reflexive form of ubi)}}}

\entry{obene}{\headword{obene}  \pronnote{obene}
  \pos{Personal pronoun}  \sense{
    \definition{ablative-possessive form of bogo}    \example{
      \exsen{Ngämo baba da ddobae mamoeang daeya. Ngäna obene ddäddäg de däddägnegne.}      \extran{My father was a very good hunter. I ate meat from his kills.}}}
  \variants{\freevarlab \vartext{obäne}}
  \variants{\Fast Speech Variant \vartext{obne}}}

\entry{obergaban}{\headword{obergaban}  \note{Weather}  \pronnote{oberɡaban}
  \pos{Noun}  \sense{
    \definition{clear day}    \note{Weather}    \example{
}    \example{
}    \example{
}}}

\entry{Obewa}{\headword{Obewa}  \pronnote{obewa}
  \pos{Proper noun}  \sense{
    \definition{male personal name}}}

\entry{obi}{\headword{obi}
  \variantof{Dialectal Variant of \vartext{ubi}}  \pronnote{obi}}

\entry{oblle}{\headword{oblle}  \pronnote{obɽe}
  \pos{Personal pronoun}  \sense{
    \definition{dative form of bogo}    \example{
      \exsen{Ngäna oblle toboll de nanttogan.}      \extran{I gave her the spear.}}}}

\entry{obne}{\headword{obne}
  \variantof{Fast Speech Variant of \vartext{obene}}}

\entry{obo}{\headword{obo}  \pronnote{obo}
  \pos{Personal pronoun}  \sense{
    \definition{possessive form of bogo}    \example{
      \exsen{Gänyan obo sana da.}      \extran{Here is his sago.}}}}

\entry{obo zaga}{\headword{obo zaga}
  \variantof{SpellingVariant of \vartext{obozaga}}}

\entry{oboll}{\headword{oboll}  \note{Trees}  \pronnote{oboɽ}
  \pos{Noun}  \sense{
    \definition{type of small tree that grows in the grassland (~2 m) with white flowers and red fruit with a large brown seed; fruit is eaten by birds; wood is used for sturdy posts}    \note{Trees}    \example{
}    \example{
}    \example{
}}}

\entry{obom}{\headword{obom}  \pronnote{obom}
  \pos{Personal pronoun}  \sense{
    \definition{accusative form of bogo}    \example{
      \exsen{Ngäna obom dandär.}      \extran{I heard him.}}}}

\entry{obosasa}{\headword{obosasa}  \pronnote{obosasa}
  \pos{Noun}  \sense{\textbf{1.} 
    \definition{Australian rufous fantail}    \scientificname{Rhipidura rufifrons}    \example{
}    \example{
}    \example{
}}  \sense{\textbf{2.} 
    \definition{blue jewel-babbler}    \scientificname{Ptilorrhoa caelurescens}}}

\entry{obozaga}{\headword{obozaga}
  \pos{Personal pronoun}  \sense{
    \definition{oneself, himself, herself (reflexive form of bogo)}    \example{
      \exsen{Yesu obozaga bo kuddäll de gomballägän.}      \extran{Jesus predicted his own death.}}}
  \variants{\Fast Speech Variant \vartext{obzaga}}
  \variants{\SpellingVariant \vartext{obo zaga}}}

\entry{oboobo}{\headword{oboobo}  \pronnote{oboobo}
  \pos{Personal pronoun}  \sense{
    \definition{oneself, himself, herself (reflexive form of bogo)}    \example{
      \exsen{Mägda oboobo tatu agan.}      \extran{Mother washed herself.}}}}

\entry{obzaga}{\headword{obzaga}
  \variantof{Fast Speech Variant of \vartext{obozaga}}}

\entry{od}{\headword{od}  \pronnote{od}
  \pos{Noun}  \sense{
    \definition{type of fatty freshwater fish}}}

\entry{odaode}{\headword{odaode}
  \variantof{Dialectal Variant of \vartext{udaude}}  \pronnote{odaode}}

\entry{odoolo}{\headword{odoolo}  \note{Birds}  \pronnote{odoolo}
  \pos{Noun}  \sense{
    \definition{type of bird}    \note{Birds}    \example{
}    \example{
}    \example{
}}}

\entry{of}{\headword{of}
  \variantof{Dialectal Variant of \vartext{English loanword}}}

\entry{officer}{\headword{officer}
  \variantof{Dialectal Variant of \vartext{English loanword}}}

\entry{Ogbaperma}{\headword{Ogbaperma}  \note{Places around Limol}  \pronnote{oɡbaperma}
  \pos{Proper noun}  \sense{
    \definition{Ogbaperma (camping place)}    \note{Places around Limol}    \example{
}    \example{
}    \example{
}}}

\entry{ogo}{\headword{ogo}
  \pos{Verb}  \sense{
    \definition{hoard}}}

\entry{Ogoa}{\headword{Ogoa}  \pronnote{oɡoa}
  \pos{Proper noun}  \sense{
    \definition{male personal name}}}

\entry{ogog}{\headword{ogog}  \note{Birds}  \pronnote{oɡoɡ}
  \pos{Noun}  \sense{
    \definition{grey-crowned babbler}    \scientificname{Pomatostomus temporalis temporalis}    \note{Birds}    \example{
}    \example{
}    \example{
}}}

\entry{oi}{\headword{oi}  \pronnote{oi}
  \pos{Interjection}  \sense{
    \definition{oh}}}

\entry{ok}{\headword{ok}
  \variantof{SpellingVariant of \vartext{oke}}  \pronnote{ok}}

\entry{oke}{\headword{oke}  \pronnote{oke}
  \pos{Interjection}  \sense{
    \definition{okay}}
  \variants{\SpellingVariant \vartext{ok}}}

\entry{oknen ma}{\headword{oknen ma}  \note{School}  \pronnote{oknen ma}
  \pos{Noun}  \sense{
    \definition{duster, eraser}    \note{School}    \example{
}    \example{
}    \example{
}}}

\entry{okok}{\headword{okok}  \note{Verbs}  \pronnote{okok}
  \pos{Transitive verb}  \sense{
    \definition{to rub}    \note{Verbs}    \example{
}    \example{
}    \example{
}}}

\entry{okokol}{\headword{okokol}  \pronnote{okokol}
  \pos{Transitive coverb}  \sense{
    \definition{to welcome}    \example{
      \exsen{Bogo mɨnyi llɨg kälsre de okokol bägagän.}      \extran{She will welcome the little boy.}}}
  \variants{\SpellingVariant \vartext{okookol}}}

\entry{okookol}{\headword{okookol}
  \variantof{SpellingVariant of \vartext{okokol}}  \note{Things a person can do}  \pronnote{okookol}}

\entry{ol}{\headword{ol}
  \pos{Noun}  \sense{
    \definition{hall}}}

\entry{ol}{\headword{ol}
  \variantof{Unspecified Variant of \vartext{wäl}}}

\entry{Olalea}{\headword{Olalea}  \pronnote{olalea}
  \pos{Proper noun}  \sense{
    \definition{female personal name}}}

\entry{old}{\headword{old}
  \variantof{Dialectal Variant of \vartext{English loanword}}}

\entry{Old Maoto}{\headword{Old Maoto}  \pronnote{old maoto}
  \pos{Proper noun}  \sense{
    \definition{Old Maoto (toponym)}    \example{
}    \example{
}    \example{
}}}

\entry{olib}{\headword{olib}
  \pos{Noun}  \sense{
    \definition{olive}}}

\entry{oll}{\headword{oll}  \note{Food}  \pronnote{oɽ}
  \pos{Noun}  \sense{
    \definition{sugarcane}    \scientificname{Saccharum}    \note{Food}    \example{
}    \example{
}    \example{
}}
  \variants{\Unspecified Variant \vartext{wäll}}}

\entry{olle}{\headword{olle}  \pronnote{oɽe}
  \pos{Intransitive coverb}  \sense{\textbf{1.} 
    \definition{to shout, yell, call out}    \example{
      \exsen{Olle de gongkamän.}      \extran{She started to shout.}}}  \sense{\textbf{2.} 
    \definition{to call (over), summon}    \example{
      \exsen{Bongo ngänäm olle nagalle gänyaolle.}      \extran{You called me over here.}}}
  \variants{\Unspecified Variant \vartext{wolle, wälle}}}

\entry{=olle}{\headword{=olle}
  \pos{Nominal enclitic}  \sense{
    \definition{allative case clitic; to, into, towards}}}

\entry{ollondd}{\headword{ollondd}  \pronnote{oɽonɖ͡ʐ}
  \pos{Noun}  \sense{
    \definition{root}}
  \variants{\Dialectal Variant \vartext{wälländd}}}

\entry{ollong}{\headword{ollong}  \pronnote{oɽoŋ}
  \pos{Noun}  \sense{
    \definition{time, occasion}    \example{
      \exsen{ttongdae ollong, komlla ollong, kumuddäga ollong, tupi ollong, mända ollong}      \extran{once, twice, thrice, four times, five times}}    \example{
      \exsen{Tätäm ngäna komlla ollong räba de daspun.}      \extran{Yesterday, I threw the eraser twice.}}}}

\entry{olmapänyik}{\headword{olmapänyik}  \note{Parts of the Body}  \pronnote{olmapəɲik}
  \pos{Noun}  \sense{
    \definition{armpit}    \note{Parts of the Body}    \example{
      \exsen{olmapänyik kom}      \extran{armpit hair}}    \example{
}    \example{
}}}

\entry{olmopäga}{\headword{olmopäga}  \note{Sago palms}  \pronnote{olmopəɡa}
  \pos{Noun}  \sense{
    \definition{type of sago}    \note{Sago palms}    \example{
}    \example{
}    \example{
}}}

\entry{omad}{\headword{omad}  \note{Friendship/social terms of endearment}  \pronnote{omad}
  \pos{Noun}  \sense{
    \definition{type of friend}    \note{Friendship/social terms of endearment}    \example{
}    \example{
}    \example{
}}}

\entry{omäg}{\headword{omäg}  \pronnote{oməɡ}
  \pos{Noun}  \sense{
    \definition{magic}    \example{
      \exsen{omäg ttoen}      \extran{sorcery}}    \example{
      \exsen{Digezänän omäg pawa walle.}      \extran{He took it out with magic power.}}    \example{
}}
  \variants{\SpellingVariant \vartext{omog}}}

\entry{omägag}{\headword{omägag}  \note{Things a person can do}  \pronnote{oməɡaɡ}
  \pos{Noun}  \sense{
    \definition{magician, sorcerer, fortune-teller}    \note{Things a person can do}    \example{
}    \example{
}}}

\entry{omawe}{\headword{omawe}  \note{Trees}  \pronnote{omawe}
  \pos{Noun}  \sense{
    \definition{type of tree}    \note{Trees}    \example{
}    \example{
}    \example{
}}}

\entry{omgälgäl}{\headword{omgälgäl}  \note{Poison}  \pronnote{omɡəlɡəl}
  \pos{Noun}  \sense{
    \definition{red ants}    \note{Poison}    \example{
}    \example{
}    \example{
}}}

\entry{omog}{\headword{omog}
  \variantof{SpellingVariant of \vartext{omäg}}}

\entry{omom}{\headword{omom}  \pronnote{omom}
  \pos{Transitive verb}  \sense{
    \definition{to sweep}    \example{
      \exsen{Ngäna ag me ma de noman.}      \extran{I swept the house in the morning.}}}}

\entry{one}{\headword{one}
  \variantof{SpellingVariant of \vartext{wan}}}

\entry{Ongg}{\headword{Ongg}  \pronnote{oŋɡ}
  \pos{Proper noun}  \sense{
    \definition{male personal name}}}

\entry{onggall}{\headword{onggall}  \pronnote{oŋɡaɽ}
  \pos{Noun}  \sense{
    \definition{frilled lizard}    \scientificname{Chlamydosaurus kingii}    \example{
      \exsen{Iba malla onggall ddägnan a umllang dag.}      \extran{We (incl.) do not eat frilled lizard.}}}}

\entry{Ono}{\headword{Ono}  \note{Places around Limol}  \pronnote{ono}
  \pos{Proper noun}  \sense{
    \definition{Ono (toponym)}    \note{Places around Limol}    \example{
}    \example{
}    \example{
}}}

\entry{ony}{\headword{ony}  \pronnote{oɲ}
  \pos{Transitive verb}  \sense{
    \definition{to carry; get; bring}    \example{
      \exsen{sabi onyang lla}      \extran{person who carries out the law}}    \example{
      \exsen{Era ebdo me mɨnyi ibi bobollab ttängäm e mätta kongkom e?}      \extran{On which day will we (incl.) go to the garden to get the yams?}}    \example{
      \exsen{Ngämi nge de diwenyeya.}      \extran{We (excl.) brought over the coconut.}}    \example{
      \exsen{Bibra otät de aya bikomän?}      \extran{Who is bringing over your (pl.) food?}}}
  \variants{\SpellingVariant \vartext{wäny}}}

\entry{opa}{\headword{opa}  \note{Trees}  \pronnote{opa}
  \pos{Noun}  \sense{
    \definition{type of tree}    \note{Trees}    \example{
}    \example{
}    \example{
}}}

\entry{opa ttangtte}{\headword{opa ttangtte}  \note{Trees}  \pronnote{opa ʈ͡ʂaŋʈ͡ʂe}
  \pos{Noun}  \sense{
    \definition{type of tree}    \note{Trees}    \example{
}    \example{
}    \example{
}}}

\entry{opap}{\headword{opap}  \pronnote{opap}
  \pos{Transitive verb}  \sense{
    \definition{to cross over, pass over, move across}    \example{
      \exsen{Ibi bopäll bem de apte menae e daupeya.}      \extran{Let's cross the sea to the other side.}}}}

\entry{ope}{\headword{ope}
  \variantof{freevarlab of \vartext{upe}}}

\entry{open}{\headword{open}
  \variantof{freevarlab of \vartext{English loanword}}}

\entry{Opo}{\headword{Opo}  \note{Places around Limol}  \pronnote{opo}
  \pos{Proper noun}  \sense{
    \definition{Opo (toponym)}    \note{Places around Limol}    \example{
}    \example{
}    \example{
}}}

\entry{opodo}{\headword{opodo}  \note{Weaving}  \pronnote{opodo}
  \pos{Noun}  \sense{
    \definition{weaving pattern with checks}    \note{Weaving}    \example{
}    \example{
}    \example{
}}}

\entry{opop}{\headword{opop}
  \pos{Verb}  \sense{\textbf{1.} 
}  \sense{\textbf{2.} 
}}

\entry{or}{\headword{or}
  \variantof{Dialectal Variant of \vartext{English loanword}}}

\entry{orbam}{\headword{orbam}  \note{Water}  \pronnote{orbam}
  \pos{Noun}  \sense{
    \definition{marsh}    \note{Water}    \example{
}    \example{
}    \example{
}}
  \variants{\SpellingVariant \vartext{worbam}}}

\entry{Oriomo}{\headword{Oriomo}
  \pos{Proper noun}  \sense{
    \definition{Oriomo (Wipi-speaking village in Oriomo-Bituri Rural LLG)}}}

\entry{Orpmang}{\headword{Orpmang}
  \pos{Proper noun}  \sense{
    \definition{Wipi language}}}

\entry{orwa}{\headword{orwa}  \pronnote{orwa}
  \pos{Noun}  \sense{
    \definition{suffering}    \example{
      \exsen{Bäne orwa atta agezän.}      \extran{Come out from your suffering.}}    \example{
      \exsen{Bogo itrell atta orwa gognän.}      \extran{He was suffering from illness.}}}}

\entry{osne}{\headword{osne}  \note{Types of Taro}  \pronnote{osne}
  \pos{Noun}  \sense{
    \definition{type of small taro}    \note{Types of Taro}    \example{
}    \example{
}    \example{
}}}

\entry{ospel}{\headword{ospel}  \pronnote{ospel}
  \pos{Noun}  \sense{
    \definition{hospital; aid post}}}

\entry{ospitol}{\headword{ospitol}
  \variantof{Dialectal Variant of \vartext{English loanword}}}

\entry{otal}{\headword{otal}  \pronnote{otal}
  \pos{Intransitive verb}  \sense{
    \definition{to perch}    \example{
      \exsen{Pa da llo me gotalalle.}      \extran{The bird perched in the tree.}}}}

\entry{otät}{\headword{otät}  \pronnote{otət}
  \pos{Transitive verb}  \sense{\textbf{1.} 
    \definition{to eat}    \example{
      \exsen{Diba toto we, bogo sana dorko de doton.}      \extran{That evening, she ate dry sago.}}    \example{
      \exsen{Ubi llo käp otnanang dag.}      \extran{They eat tree fruit.}}}  \sense{\textbf{2.} 
    \definition{food}    \example{
      \exsen{Bogo mer wätät de yu dägagän.}      \extran{She cooked nice food.}}}  \sense{\textbf{3.} 
    \definition{hunger}    \example{
      \exsen{Otät atta ngämo kutt a mängallmeny agmallo.}      \extran{My bones became weak from hunger.}}}
  \variants{\Unspecified Variant \vartext{wätät}}}

\entry{otät ngänaeka ttoen}{\headword{otät ngänaeka ttoen}  \pronnote{otət ŋənaeka ʈ͡ʂoen}
  \pos{Noun}  \sense{
    \definition{starvation}}}

\entry{otät yu ma}{\headword{otät yu ma}  \note{Everything in the house}  \pronnote{otət ju ma}
  \pos{Noun}  \sense{
    \definition{kitchen}    \note{Everything in the house}    \example{
}    \example{
}    \example{
}}}

\entry{otnan ma}{\headword{otnan ma}  \pronnote{otnan ma}
  \pos{Modifier}  \sense{
    \definition{edible}}}

\entry{ottam}{\headword{ottam}
  \pos{Verb}  \sense{
    \definition{find}}}

\entry{Ouli}{\headword{Ouli}  \pronnote{ouli}
  \pos{Proper noun}  \sense{
    \definition{female personal name}}}

\entry{owel}{\headword{owel}
  \pos{Noun}  \sense{
    \definition{oil}}}

\chapter{OO}


\entry{O-o-o}{\headword{O-o-o}
  \pos{Interjection}  \sense{
}}

\chapter{P}


\entry{pa~}{\headword{pa~}
  \variantof{Inflected Form of \vartext{RED~}}}

\entry{pa}{\headword{pa}  \lexeme{pa}  \pronnote{pa}
  \pos{Noun}  \sense{
    \definition{bird}    \example{
      \exsen{Bogo pa näddägan ngämäne gäzatt de.}      \extran{He ate the bird that I killed.}}}}

\entry{pa kom}{\headword{pa kom}
  \pos{Noun}  \sense{
    \definition{feather}}}

\entry{päd}{\headword{päd}  \pronnote{pəd}
  \pos{Noun}  \sense{
    \definition{scar}    \example{
}    \example{
}    \example{
}}}

\entry{päddab}{\headword{päddab}  \note{Could this be two verbs - päddab and päddae?}  \pronnote{pəɖ͡ʐab}
  \pos{Intransitive verb}  \sense{\textbf{1.} 
    \definition{to grow}    \note{Could this be two verbs - päddab and päddae?}    \example{
      \exsen{Tomato da däpddabän.}      \extran{The tomato plant grew.}}}  \sense{\textbf{2.} 
    \definition{to be cooked, be done}    \note{Could this be two verbs - päddab and päddae?}    \example{
      \exsen{Ngäma wätät a zime opddaenegan.}      \extran{Our (excl.) food has finished cooking.}}}
  \variants{\Unspecified Variant \vartext{pɨddab}}}

\entry{päddpädd}{\headword{päddpädd}
  \pos{Intransitive verb}  \sense{\textbf{1.} 
    \definition{to grow}    \example{
      \exsen{ine toko me päddnenang towall}      \extran{grass growing on the water's surface}}    \example{
      \exsen{Däm kutt a ekaklle me guspullän a däpäddän.}      \extran{The seed fell to the ground and grew.}}}  \sense{\textbf{2.} 
    \definition{to explode}    \example{
      \exsen{Yu me pattlle da päddpädd allan.}      \extran{The bamboo explodes in the fire.}}}
  \variants{\SpellingVariant \vartext{pɨddpɨdd}}}

\entry{padiem}{\headword{padiem}  \note{Names of Bananas}  \pronnote{padiem}
  \pos{Noun}  \sense{
    \definition{type of introduced banana}    \note{Names of Bananas}    \example{
}    \example{
}    \example{
}}}

\entry{pädoe}{\headword{pädoe}  \pronnote{pədoe}
  \pos{Transitive verb}  \sense{\textbf{1.} 
    \definition{to blow}    \example{
      \exsen{Wel a gopädoeän.}      \extran{The wind blew.}}}  \sense{\textbf{2.} 
    \definition{to blow}    \example{
      \exsen{Mälla da wa lla da utt de dapädoenegeyo.}      \extran{The man and women blew conch shells.}}}}

\entry{pädpäd}{\headword{pädpäd}  \note{Verbs}  \pronnote{pədpəd}
  \pos{Noun}  \sense{
    \definition{tattoo}    \note{Verbs}    \example{
}    \example{
}    \example{
}}}

\entry{pädrall}{\headword{pädrall}  \pronnote{pədraɽ}
  \pos{Intransitive verb}  \sense{\textbf{1.} 
    \definition{to spread (out), scatter, disperse}    \example{
      \exsen{Lla da oba mällamällong gupädrellän otät ma.}      \extran{The men and their wives spread out to get food.}}    \example{
      \exsen{Oba näkäpngon a tämamae apädrellan.}      \extran{Their thoughts are all scattered.}}}  \sense{\textbf{2.} 
    \definition{to spread (out), scatter, sow}    \example{
      \exsen{däm kutt pädrallag lla}      \extran{man sowing seeds}}    \example{
      \exsen{Bogo up di dotnegnän a ttoe dapädrällnegnän.}      \extran{He was eating bananas and scattering the peels.}}}
  \variants{\SpellingVariant \vartext{pudrell}}}

\entry{paeb}{\headword{paeb}  \pronnote{paeb}
  \pos{Numeral}  \sense{
    \definition{five (English numeral)}}
  \variants{\SpellingVariant \vartext{faeb, five, paib}}}

\entry{Paeke}{\headword{Paeke}
  \pos{Proper noun}  \sense{
    \definition{female personal name}}}

\entry{Paelet}{\headword{Paelet}  \pronnote{paelet}
  \pos{Proper noun}  \sense{
    \definition{Pilate}}}

\entry{paen}{\headword{paen}
  \pos{Noun}  \sense{
    \definition{fine}}}

\entry{paengg}{\headword{paengg}  \pronnote{paeŋɡ}
  \pos{Transitive verb}  \sense{
    \definition{to guess}    \example{
      \exsen{Ngäna paengg eran ada Jugu mɨnyi sisri bongttägän Wim atta.}      \extran{I'm guessing that Jugu will arrive from Wim now.}}}}

\entry{pag}{\headword{pag}  \pronnote{paɡ}
  \pos{Noun}  \sense{
    \definition{salt}}}

\entry{pägamän}{\headword{pägamän}
  \pos{Noun}  \sense{
    \definition{type of yam}}}

\entry{paib}{\headword{paib}
  \variantof{SpellingVariant of \vartext{paeb}}}

\entry{Paine}{\headword{Paine}  \pronnote{paine}
  \pos{Proper noun}  \sense{
    \definition{male personal name}}}

\entry{paitt}{\headword{paitt}
  \pos{Noun}  \sense{
    \definition{fight}}}

\entry{päk}{\headword{päk}  \pronnote{pək}
  \pos{Locational}  \sense{
    \definition{edge}    \example{
      \exsen{gall päk}      \extran{edge of the canoe}}}}

\entry{päkam}{\headword{päkam}  \pronnote{pəkam}
  \pos{Verb}  \sense{
    \definition{to bend down, keel over (due to sickness)}}}

\entry{pakätt}{\headword{pakätt}  \note{Clothing}  \pronnote{pakəʈ͡ʂ}
  \pos{Noun}  \sense{
    \definition{widow's robe}    \note{Clothing}    \example{
}    \example{
}    \example{
}}}

\entry{paklle}{\headword{paklle}  \note{Snakes}  \pronnote{pakɽe}
  \pos{Noun}  \sense{
    \definition{type of snake}    \note{Snakes}    \example{
}    \example{
}    \example{
}}}

\entry{Pakllepakllemäll}{\headword{Pakllepakllemäll}  \note{Places around Limol}  \pronnote{pakɽepakɽeməɽ}
  \pos{Proper noun}  \sense{
    \definition{Pakllepakllemäll (camping place)}    \note{Places around Limol}    \example{
}    \example{
}    \example{
}}}

\entry{pakos}{\headword{pakos}  \pronnote{pakos}
  \pos{Noun}  \sense{
    \definition{type of spear}    \example{
      \exsen{Madura pakos de daspunän.}      \extran{Madura shot the spear.}}}}

\entry{palament}{\headword{palament}
  \pos{Noun}  \sense{
}}

\entry{pälan}{\headword{pälan}
  \pos{Noun}  \sense{
    \definition{plan}    \example{
      \exsen{Mälla da pälan gognegän.}      \extran{The women made a plan.}}}}

\entry{pale}{\headword{pale}  \note{Clothing}  \pronnote{pale}
  \pos{Noun}  \sense{
    \definition{type of white clay used for painting babies}    \note{Clothing}    \example{
}    \example{
}    \example{
}}}

\entry{pälengg}{\headword{pälengg}  \pronnote{pəleŋɡ}
  \pos{Intransitive verb}  \sense{
    \definition{to drop down, fall}    \example{
      \exsen{Angde gogllaeaebne, yogoll ulle da dipliwän.}      \extran{While we were paddling, a heavy rain started coming down.}}}}

\entry{pälkom}{\headword{pälkom}  \note{Parts of the Body}  \pronnote{pəlkom}
  \pos{Noun}  \sense{
    \definition{body hair}    \note{Parts of the Body}    \example{
      \exsen{Ngämo pälkom da guirewän.}      \extran{My hairs stood.}}    \example{
}    \example{
}}
  \variants{\Unspecified Variant \vartext{pällkom}}}

\entry{pall}{\headword{pall}
  \pos{Noun}  \sense{
    \definition{varieties of palms with coconuts with a red exocarp}}}

\entry{pall kottllam}{\headword{pall kottllam}
  \pos{Noun}  \sense{
    \definition{red-bellied short-necked turtle}    \scientificname{Emydura subglobosa}}}

\entry{pall kubull}{\headword{pall kubull}  \note{Animal names}  \pronnote{paɽ kubuɽ}
  \pos{Noun}  \sense{
    \definition{red-legged pademelon}    \scientificname{Thylogale stigmatica orimo}    \note{Animal names}    \example{
}    \example{
}    \example{
}}}

\entry{pall tawe}{\headword{pall tawe}  \pronnote{paɽ tawe}
  \pos{Noun}  \sense{
    \definition{type of large palm that grows in the grassland with white flowers and coconuts with a red exocarp; small sticks are good for houses, and bark is used for walling}}}

\entry{pälläk}{\headword{pälläk}  \note{Names of Bananas}  \pronnote{pəɽək}
  \pos{Noun}  \sense{
    \definition{type of introduced banana}    \note{Names of Bananas}    \example{
}    \example{
}    \example{
}}}

\entry{pallall}{\headword{pallall}  \pronnote{paɽaɽ}
  \pos{Noun}  \sense{\textbf{1.} 
    \definition{direction; area}    \example{
      \exsen{Ngäna naigae pallall e ibi allan.}      \extran{I'm going south.}}    \example{
      \exsen{Adawalle de ddia da gupudrellän iba pallall e.}      \extran{From there, the deer spread in our (incl.) direction.}}}  \sense{\textbf{2.} 
    \definition{side}    \example{
      \exsen{sawe pallall}      \extran{left side}}    \example{
      \exsen{Ubi guspullnän ekaklle we obo ngattong pallall alle.}      \extran{They were falling to the ground in front of him.}}}  \sense{\textbf{3.} 
    \definition{about}    \example{
      \exsen{Män a llɨg pate däräng pallall e eka gogon.}      \extran{The girl talked to the boy about the dog.}}}}

\entry{pälläm}{\headword{pälläm}  \pronnote{pəɽəm}
  \pos{Modifier}  \sense{
    \definition{visible}    \example{
      \exsen{Ddäg kutt dae pälläm gognän.}      \extran{Only the backbone was visible.}}}}

\entry{pallam}{\headword{pallam}  \pronnote{paɽam}
  \pos{Transitive verb}  \sense{
    \definition{to cut open}    \example{
      \exsen{Obom dapellaemnän.}      \extran{He was cutting him open (i.e. for surgery).}}}}

\entry{pällämpälläm}{\headword{pällämpälläm}  \note{Colors}  \pronnote{pəɽəmpəɽəm}
  \pos{Color term}  \sense{\textbf{1.} 
    \definition{white, bright}    \note{Colors}    \example{
      \exsen{Obo iddpo da pällämpälläm a gombenmällnegnän.}      \extran{His clothes were white and shining.}}    \example{
      \exsen{Zime abo ballme da gopnän pällämpälläm gognän ttängäm a.}      \extran{Dawn had already broken and the village was becoming bright.}}    \example{
}}  \sense{\textbf{2.} 
    \definition{white, Western}    \note{Colors}}}

\entry{pällämpälläm eka}{\headword{pällämpälläm eka}  \note{School}  \pronnote{pəɽəmpəɽəm eka}
  \pos{Proper noun}  \sense{
    \definition{English}    \note{School}    \example{
}    \example{
}    \example{
}}}

\entry{pällämpälläm yurwe}{\headword{pällämpälläm yurwe}  \note{Trees}  \pronnote{pəɽəmpəɽəm jurwe}
  \pos{Noun}  \sense{
    \definition{type of tree with white flowers and edible white fruit.}    \note{Trees}    \example{
}    \example{
}    \example{
}}}

\entry{pallängkmeny}{\headword{pallängkmeny}  \pronnote{paɽəŋkmeɲ}
  \pos{Transitive verb}  \sense{\textbf{1.} 
    \definition{to divide}    \example{
      \exsen{Da ttongo kantri me lla da bompallängkmenynegän obaoba komlla kullum e a bogäddeyo obaoba, däbe kantri da mɨnyi ddone bakällmewän.}      \extran{If the people in a country divide themselves into two groups and fight each other, that country will not survive.}}}  \sense{\textbf{2.} 
    \definition{to judge}    \example{
      \exsen{God bo enanae ttoen pallängkmeny}      \extran{God's Final Judgment}}}}

\entry{pällganen}{\headword{pällganen}
  \pos{Transitive verb}  \sense{
    \definition{to hang}    \example{
      \exsen{Ttongo mälla da iddpo de nällpäganegnan.}      \extran{A woman was hanging the clothes.}}}}

\entry{pällkam}{\headword{pällkam}  \pronnote{pəɽkam}
  \pos{Transitive verb}  \sense{
    \definition{to split, break}    \example{
      \exsen{Ngäna kakoll de napällkaman.}      \extran{I broke the plate.}}    \example{
      \exsen{Mälla da täl de dapällkaemaemneyo.}      \extran{The women were splitting the bamboo.}}    \example{
      \exsen{Ngäna ngämo ikop glas de kamekame napällkaman.}      \extran{I accidentally broke my glasses.}}}}

\entry{pällkapällkae}{\headword{pällkapällkae}  \pronnote{pəɽkapəɽkae}
  \pos{Noun}  \sense{
    \definition{pieces, shards}}
  \variants{\Unspecified Variant \vartext{pallkepallke}}}

\entry{pallkeakeya}{\headword{pallkeakeya}  \note{Mushrooms}  \pronnote{paɽkeakeja}
  \pos{Noun}  \sense{
    \definition{type of red mushroom that grows in the grassland}    \note{Mushrooms}    \example{
}    \example{
}    \example{
}}}

\entry{pallkepallke}{\headword{pallkepallke}
  \variantof{Unspecified Variant of \vartext{pällkapällkae}}  \pronnote{paɽkepaɽke}}

\entry{pallkoll}{\headword{pallkoll}  \pronnote{paɽkoɽ}
  \pos{Noun}  \sense{
    \definition{piece}    \example{
      \exsen{yu pallkoll}      \extran{piece of firewood}}}
  \variants{\SpellingVariant \vartext{pällkoll}}}

\entry{pällkoll}{\headword{pällkoll}
  \variantof{SpellingVariant of \vartext{pallkoll}}  \pronnote{pəɽkoɽ}}

\entry{pällkom}{\headword{pällkom}
  \variantof{Unspecified Variant of \vartext{pälkom}}}

\entry{Pällmang}{\headword{Pällmang}  \note{Places around Limol}  \pronnote{pəɽmaŋ}
  \pos{Proper noun}  \sense{
    \definition{Pällmang (toponym)}    \note{Places around Limol}    \example{
}    \example{
}    \example{
}}}

\entry{pällnampällnam}{\headword{pällnampällnam}  \pronnote{pəɽnampəɽnam}
  \pos{Adverb}  \sense{
    \definition{squatting, crouching}    \example{
      \exsen{Bogo käg me pällnampällnam adämenan.}      \extran{She squatted down on the floor.}}}}

\entry{pällonpällon}{\headword{pällonpällon}
  \variantof{Unspecified Variant of \vartext{pollonpollon}}  \pronnote{pəɽonpəɽon}}

\entry{pallta}{\headword{pallta}
  \pos{Modifier}  \sense{
}}

\entry{palltapallta}{\headword{palltapallta}
  \pos{Modifier}  \sense{
    \definition{flat}    \example{
      \exsen{llo palltapallta torpoll}      \extran{flat piece of wood}}}}

\entry{pällttän}{\headword{pällttän}  \note{Updated 2017}  \pronnote{pəɽʈ͡ʂən}
  \pos{Intransitive verb}  \sense{
    \definition{to set off, start walking}    \note{Updated 2017}    \example{
      \exsen{Däpällätt abo do Taolang.}      \extran{I started walking to Taolang.}}}}

\entry{pällulle}{\headword{pällulle}  \pronnote{pəɽuɽe}
  \pos{Noun}  \sense{
    \definition{lung}    \example{
}    \example{
}    \example{
}}
  \variants{\Unspecified Variant \vartext{pllulle, pllulli}}}

\entry{Palsa}{\headword{Palsa}
  \pos{Proper noun}  \sense{
    \definition{Palsa (toponym)}}}

\entry{pam}{\headword{pam}  \pronnote{pam}
  \pos{Noun}  \sense{
    \definition{injection, shot, vaccine}    \example{
      \exsen{Ospel e mängalae gobäll, llɨg di dokomeya pam alle zuwoenen e.}      \extran{We quickly went to the hospital, bringing the children to get shots.}}}}

\entry{pama}{\headword{pama}
  \pos{Noun}  \sense{
    \definition{farmer}}
  \variants{\SpellingVariant \vartext{fama}}}

\entry{pämbäll}{\headword{pämbäll}
  \pos{Noun}  \sense{
}}

\entry{pambu}{\headword{pambu}  \note{Gardens}  \pronnote{pambu}
  \pos{Noun}  \sense{
    \definition{hoe (tool)}    \note{Gardens}    \example{
}    \example{
}    \example{
}}}

\entry{pameny}{\headword{pameny}
  \pos{Intransitive verb}  \sense{
    \definition{to scream}    \example{
      \exsen{Lla da ompamenyan.}      \extran{The man screamed.}}}}

\entry{pamker}{\headword{pamker}  \lexeme{pamker}  \pronnote{pamker}
  \pos{Noun}  \sense{
    \definition{pumpkin}}}

\entry{pamli}{\headword{pamli}
  \variantof{SpellingVariant of \vartext{pemli}}}

\entry{pampem}{\headword{pampem}
  \pos{Transitive verb}  \sense{
    \definition{to fish}    \example{
      \exsen{Ag me källäm pamnen e gobäll.}      \extran{In the morning, we went to fish at the pond.}}}}

\entry{pämpllätt}{\headword{pämpllätt}  \pronnote{pəmpɽəʈ͡ʂ}
  \pos{Intransitive verb}  \sense{
    \definition{to become healthy; grow}}}

\entry{pana}{\headword{pana}  \note{Friendship/social terms of endearment}  \pronnote{pana}
  \pos{Noun}  \sense{
    \definition{type of relationship}    \note{Friendship/social terms of endearment}    \example{
}    \example{
}    \example{
}}}

\entry{pänae}{\headword{pänae}  \pronnote{pənae}
  \pos{Transitive verb}  \sense{\textbf{1.} 
    \definition{to turn back, turn around}    \example{
      \exsen{Ubi komlla gopnaeyo a ngäsengäse dindugiyu.}      \extran{The two of them turned around and ran the other way.}}}  \sense{\textbf{2.} 
    \definition{to capsize, turn over, flip over}    \example{
      \exsen{Gopnaeyän.}      \extran{It capsized.}}    \example{
      \exsen{Ngämaene gall a ada ka bopänayän, adawatta däbe za da ngämim ge deyangkollmällnän.}      \extran{Our (excl.) boat nearly capsized because that thing was following us.}}}  \sense{\textbf{3.} 
    \definition{to turn, become, transform}    \example{
      \exsen{Obo mikutt a ttänttämang e gopnaeyän.}      \extran{His anger made him hot.}}    \example{
      \exsen{Lla daeya Moli, be angde walle we guspunän, bunkuttang e gopnaen.}      \extran{Moli was a man, but when he jumped into the water, he turned into a catfish.}}}  \sense{\textbf{4.} 
    \definition{to translate}    \example{
      \exsen{Ngäna ttoenttoen de näpänaeyan Ende we.}      \extran{I translated this story into Ende.}}}  \sense{\textbf{5.} 
    \definition{zenith (position of the sun at high noon)}}}

\entry{pänaemeny}{\headword{pänaemeny}  \pronnote{pənaemeɲ}
  \pos{Transitive verb}  \sense{
    \definition{to check}    \example{
      \exsen{Mälla da mɨnyi bobällnän mätta mɨnyi bampnaemenyaemneyo ada mätta da zäme käp agan abo.}      \extran{The women will go and check the yams, whether they are ready yet.}}}}

\entry{pänaengg}{\headword{pänaengg}
  \pos{Intransitive verb}  \sense{
    \definition{to move opposite}}}

\entry{Panakawa}{\headword{Panakawa}
  \pos{Proper noun}  \sense{
    \definition{Panakawa (toponym)}}}

\entry{pänayang}{\headword{pänayang}  \note{Stages of life}  \pronnote{pənajaŋ}
  \pos{Noun}  \sense{
    \definition{stage of life when a child begins to turn over}    \note{Stages of life}    \example{
}    \example{
}    \example{
}}}

\entry{pänbäll}{\headword{pänbäll}  \note{Fishing}  \pronnote{pənbəɽ}
  \pos{Noun}  \sense{
    \definition{poisonous vine or root (used in fishing to stun fish)}    \note{Fishing}    \example{
}    \example{
}    \example{
}}}

\entry{panda}{\headword{panda}  \note{Trees}  \pronnote{panda}
  \pos{Noun}  \sense{
    \definition{type of tree}    \note{Trees}    \example{
}    \example{
}    \example{
}}}

\entry{pänddäg}{\headword{pänddäg}
  \pos{Transitive verb}  \sense{
    \definition{causative-applicative form of päddpädd}    \example{
      \exsen{Lla da täl de dämpäddmenyaemallo yu alle.}      \extran{People used to blow up bamboo with fire.}}}}

\entry{pänddäg}{\headword{pänddäg}  \pronnote{pənɖ͡ʐəɡ}
  \pos{Intransitive verb}  \sense{\textbf{1.} 
    \definition{to start a grassfire}    \example{
      \exsen{Bogo gompäddägän.}      \extran{He started the grassfire.}}}  \sense{\textbf{2.} 
    \definition{to hatch}    \example{
      \exsen{Pauro käp a gompäddägän.}      \extran{The chicken egg started hatching.}}}  \sense{\textbf{3.} 
    \definition{to break water}}  \sense{\textbf{4.} 
    \definition{to mature (of a plant or a person)}    \example{
}}}

\entry{pane}{\headword{pane}  \pronnote{pane}
  \pos{Noun}  \sense{
    \definition{pot}    \example{
      \exsen{Ddob mälla da pane moepang de ikrol alle gällämnan allo.}      \extran{Some women wash dirty pots using ash.}}}}

\entry{pängän}{\headword{pängän}
  \pos{Intransitive verb}  \sense{
    \definition{to disappear}    \example{
      \exsen{Bogo apngänän.}      \extran{He disappeared.}}}}

\entry{pänggmeny}{\headword{pänggmeny}  \pronnote{pəŋɡmeɲ}
  \pos{Transitive verb}  \sense{
    \definition{to protect, look after, take care of}    \example{
      \exsen{Ede adawatta ibi mɨnyi iba pätt de mermerae bampägmenyaemnalla.}      \extran{This is why we must take care of our bodies.}}}}

\entry{panggopanggo}{\headword{panggopanggo}  \lexeme{panggopanggo}  \pronnote{paŋɡopaŋɡo}
  \pos{Noun}  \sense{
    \definition{type of round yam with a white interior and hairs}}}

\entry{pängmeny}{\headword{pängmeny}  \pronnote{pəŋmeɲ}
  \pos{Verb}  \sense{
    \definition{to scope out}    \example{
}}}

\entry{pani}{\headword{pani}
  \pos{Modifier}  \sense{
    \definition{funny}}}

\entry{panis}{\headword{panis}
  \pos{Transitive coverb}  \sense{
    \definition{to punish}}}

\entry{pänmällang mälla}{\headword{pänmällang mälla}  \note{Types of Taro}  \pronnote{pənməɽaŋ məɽa}
  \pos{Noun}  \sense{
    \definition{type of big taro}    \note{Types of Taro}    \example{
}    \example{
}    \example{
}}}

\entry{pänongg}{\headword{pänongg}
  \pos{Transitive verb}  \sense{\textbf{1.} 
    \definition{to awake, wake up, get up}    \example{
      \exsen{Ngäna gomponongg.}      \extran{I woke up.}}}  \sense{\textbf{2.} 
    \definition{to wake up, wake}    \example{
      \exsen{Ngäna obom pänongg eran.}      \extran{I am waking him up.}}}}

\entry{pänpän}{\headword{pänpän}
  \pos{Noun}  \sense{\textbf{1.} 
    \definition{dust}    \example{
      \exsen{Bibi mɨnyi bina nying pänpän de näpttayaemneyo.}      \extran{You (all) will be shaking the dust from your feet.}}    \example{
}}  \sense{\textbf{2.} 
    \definition{calcium hydroxide, lime}    \anthronote{consumed with areca nut and betel leaf}}
  \variants{\Unspecified Variant \vartext{pɨnpɨn}}}

\entry{panya}{\headword{panya}  \note{Food}  \pronnote{paɲa}
  \pos{Noun}  \sense{
    \definition{pineapple}    \note{Food}    \anthronote{ripe in December and January}    \example{
}    \example{
}    \example{
}}}

\entry{pänyae}{\headword{pänyae}  \note{Animal verbs}  \pronnote{pəɲae}
  \pos{Intransitive verb}  \sense{
    \definition{to hop}    \note{Animal verbs}    \example{
      \exsen{Bogo pänyaemällang dan.}      \extran{He hops slowly.}}    \example{
}    \example{
}}}

\entry{pänyanz}{\headword{pänyanz}  \pronnote{pəɲanz}
  \pos{Verb}  \sense{
    \definition{to grow}    \example{
      \exsen{Saly ddone ada lla tupi da pänyanz eran.}      \extran{Saly is growing to be very tall.}}}}

\entry{pänyik}{\headword{pänyik}  \pronnote{pəɲik}
  \pos{Intransitive coverb}  \sense{
    \definition{to whisper}}}

\entry{panypeny}{\headword{panypeny}  \note{School}  \pronnote{paɲpeɲ}
  \pos{Transitive verb}  \sense{
    \definition{to speak, talk, say}    \note{School}    \example{
      \exsen{Ngäna Ende eka de panypeny eran.}      \extran{I speak Ende.}}    \example{
      \exsen{Ngäna Ingglis eka, Motu eka de ddone bäpanyneg.}      \extran{I can't speak English or Hiri Motu.}}    \example{
      \exsen{Bogo ge eka de däpanyän.}      \extran{She said this.}}}}

\entry{paopao}{\headword{paopao}
  \pos{Transitive verb}  \sense{
    \definition{to cut around, prune (a plant)}    \example{
      \exsen{Däbe kapang ulle de näpao!}      \extran{[You] prune the big acacia tree!}}}}

\entry{pap}{\headword{pap}  \pronnote{pap}
  \pos{Noun}  \sense{
    \definition{type of round mutae yam}}}

\entry{päp}{\headword{päp}  \pronnote{pəp}
  \pos{Transitive verb}  \sense{\textbf{1.} 
    \definition{to draw or mark}}  \sense{\textbf{2.} 
    \definition{to speed up}}}

\entry{päpa}{\headword{päpa}
  \pos{Noun}  \sense{
    \definition{line}}}

\entry{papa}{\headword{papa}  \pronnote{papa}
  \pos{Transitive coverb}  \sense{
    \definition{to hit, beat}    \example{
      \exsen{Papa nägagan.}      \extran{He hit him.}}}}

\entry{papälläk}{\headword{papälläk}  \pronnote{papəɽək}
  \pos{Transitive verb}  \sense{
    \definition{to split, chop}    \example{
      \exsen{Bogo yu di däpellknegän.}      \extran{She split the wood.}}}
  \variants{\Unspecified Variant \vartext{papllek}}}

\entry{pape}{\headword{pape}  \pronnote{pape}
  \pos{Transitive verb}  \sense{\textbf{1.} 
    \definition{to cut (grass)}    \example{
      \exsen{Kitar pu de däpeyo.}      \extran{They cut the kitar grass.}}    \example{
}    \example{
}}  \sense{\textbf{2.} 
    \definition{to smash, crush (something soft, but not a flower)}    \example{
      \exsen{Ngäna mameat de näpewan.}      \extran{I smashed the papaya.}}    \example{
      \exsen{Ubi mätta de bäpeaemneyo.}      \extran{They will be smashing yams.}}}}

\entry{papek}{\headword{papek}  \pronnote{papek}
  \pos{Noun}  \sense{\textbf{1.} 
    \definition{dam, blockage; wall}}  \sense{\textbf{2.} 
    \definition{to block; close}    \example{
      \exsen{Ud a papekatt dan.}      \extran{The door has been closed.}}    \example{
      \exsen{Bogo de papek allan ttang alle.}      \extran{He blocked him with his hand.}}}  \sense{\textbf{3.} 
    \definition{to wall, make a wall}    \example{
      \exsen{llo ttoe kädnanatt ma paknen e}      \extran{a thing that is taken off trees to make walls}}    \example{
      \exsen{Oba ma da pall put alle papekatt dan.}      \extran{Their house is walled with pall bark.}}}}

\entry{papiye}{\headword{papiye}  \pronnote{papije}
  \pos{Noun}  \sense{
    \definition{animal tracks}    \example{
      \exsen{Ddone ada ddäddäg papiye dag sisri wälläng me.}      \extran{There are many animal tracks in the bush now.}}}}

\entry{paplläg}{\headword{paplläg}  \note{Verbs of Motion}  \pronnote{papɽəɡ}
  \pos{Intransitive verb}  \sense{
    \definition{to fly}    \note{Verbs of Motion}    \example{
      \exsen{Ngämi pällgaeb amalla.}      \extran{We are flying.}}    \example{
      \exsen{Pa da dapllägän matu we.}      \extran{The bird flew down.}}}}

\entry{papllek}{\headword{papllek}
  \variantof{Unspecified Variant of \vartext{papälläk}}}

\entry{papoe}{\headword{papoe}  \pronnote{papoe}
  \pos{Transitive verb}  \sense{
    \definition{to pierce, make a hole}    \example{
      \exsen{Bogo obo llan de napoeyan.}      \extran{She pierced her ear.}}    \example{
      \exsen{Ngäna nyäng tär ngasnges e nyäng de napoenan waeya käsre alle.}      \extran{I used the small wire to make holes in the bag to put the new bag string.}}}}

\entry{Papon}{\headword{Papon}  \pronnote{papon}
  \pos{Proper noun}  \sense{
    \definition{male personal name}}}

\entry{papu}{\headword{papu}
  \variantof{Baby Talk Variant of \vartext{kapu}}  \pronnote{papu}}

\entry{Papua}{\headword{Papua}
  \variantof{Dialectal Variant of \vartext{English loanword}}}

\entry{Papua Niugini}{\headword{Papua Niugini}
  \pos{Proper noun}  \sense{
    \definition{Papua New Guinea}}
  \variants{\SpellingVariant \vartext{Papwa Niyu Gini, Papwa Nu Gini}}}

\entry{Papwa Niyu Gini}{\headword{Papwa Niyu Gini}
  \variantof{SpellingVariant of \vartext{Papua Niugini}}}

\entry{Papwa Nu Gini}{\headword{Papwa Nu Gini}
  \variantof{SpellingVariant of \vartext{Papua Niugini}}}

\entry{para}{\headword{para}  \pronnote{para}
  \pos{Noun}  \sense{
    \definition{event where harvest and hunting bounty are compared and gifted for bragging rights}}}

\entry{Parade}{\headword{Parade}
  \pos{Noun}  \sense{
    \definition{Parade (toponym)}}}

\entry{päräl}{\headword{päräl}  \pronnote{pərəl}
  \pos{Noun}  \sense{
    \definition{radjah shelduck}    \scientificname{Radjah radjah}}}

\entry{päräll}{\headword{päräll}  \note{Adjectives (Swadesh)}  \pronnote{pərəɽ}
  \pos{Modifier}  \sense{
    \definition{dry}    \note{Adjectives (Swadesh)}    \example{
}}}

\entry{Parama}{\headword{Parama}  \pronnote{parama}
  \pos{Proper noun}  \sense{
    \definition{Parama (in Kiwai Rural LLG)}    \example{
}    \example{
}    \example{
}}}

\entry{päre}{\headword{päre}  \pronnote{pəre}
  \pos{Adverb}  \sense{
    \definition{sadly, nonsingular form of säre}    \example{
      \exsen{Llokottang dag päre ngäma ttam a.}      \extran{Sadly, our (excl.) lives are difficult.}}}
  \variants{\Fast Speech Variant \vartext{pre}}}

\entry{parga}{\headword{parga}  \pronnote{parɡa}
  \pos{Noun}  \sense{
    \definition{bridge}    \example{
      \exsen{Bob a parga de dällekän.}      \extran{The flood destroyed the bridge.}}}}

\entry{paro kottllam}{\headword{paro kottllam}  \pronnote{paro koʈ͡ʂɽam}
  \pos{Noun}  \sense{
    \definition{type of turtle with red scales on neck}}}

\entry{parpar}{\headword{parpar}  \pronnote{parpar}
  \pos{Noun}  \sense{
    \definition{type of sago bundle wrapped in a long cylinder in sago leaves}}}

\entry{pärta}{\headword{pärta}  \note{Numbers}  \pronnote{pərta}
  \pos{Numeral}  \sense{
    \definition{thirty-six (yam counting numeral; 6^2)}    \note{Numbers}    \example{
}    \example{
}    \example{
}}
  \variants{\Unspecified Variant \vartext{purta}}}

\entry{päs}{\headword{päs}  \pronnote{pəs}
  \pos{Noun}  \sense{
    \definition{type of sugarcane-like plant with fruit on top}}}

\entry{pasis}{\headword{pasis}
  \pos{Noun}  \sense{
    \definition{deep}}}

\entry{Paskam}{\headword{Paskam}
  \pos{Proper noun}  \sense{
    \definition{male personal name}}}

\entry{paspas}{\headword{paspas}  \pronnote{paspas}
  \pos{Noun}  \sense{
    \definition{type of game where players try to keep a ball off the ground}}}

\entry{pasta}{\headword{pasta}  \pronnote{pasta}
  \pos{Noun}  \sense{
    \definition{pastor}    \example{
      \exsen{Ngämi sisri pasta melem de ngasnen eralla.}      \extran{Now we are doing pastor work.}}    \example{
}    \example{
}}}

\entry{Pasuwe}{\headword{Pasuwe}
  \pos{Proper noun}  \sense{
    \definition{company name}}}

\entry{pat}{\headword{pat}  \note{Types of Taro}  \pronnote{pat}
  \pos{Noun}  \sense{
    \definition{type of taro that is eaten from the suckers}    \note{Types of Taro}    \example{
}    \example{
}    \example{
}}}

\entry{pat}{\headword{pat}
  \pos{Noun}  \sense{
}}

\entry{Pata}{\headword{Pata}
  \pos{Proper noun}  \sense{
    \definition{male personal name}}}

\entry{=patalle}{\headword{=patalle}
  \variantof{Dialectal Variant of \vartext{=patatt}}}

\entry{pätangeny}{\headword{pätangeny}
  \pos{Verb}  \sense{
}}

\entry{pätapäta}{\headword{pätapäta}
  \pos{Modifier}  \sense{
    \definition{dry}}}

\entry{pätär}{\headword{pätär}  \note{Parts of the Body}  \pronnote{pətər}
  \pos{Noun}  \sense{
    \definition{white hair}    \note{Parts of the Body}    \example{
      \exsen{Bogo zäme pätär peyang allan.}      \extran{He already has white hair.}}    \example{
}    \example{
}}}

\entry{patara}{\headword{patara}  \pronnote{patara}
  \pos{Noun}  \sense{
    \definition{wall}    \example{
      \exsen{Komlla ebdo me abo ada käg a wa patara melem a dibem deyanges.}      \extran{Then, I made the floors and walls in two days.}}}}

\entry{pätaräll}{\headword{pätaräll}
  \pos{Verb}  \sense{
}}

\entry{patarapatara}{\headword{patarapatara}  \note{Games}  \pronnote{patarapatara}
  \pos{Noun}  \sense{
    \definition{type of hand game where players make an L shape with their fingers}    \note{Games}    \example{
}    \example{
}    \example{
}}}

\entry{=patatt}{\headword{=patatt}
  \pos{Nominal enclitic}  \sense{
    \definition{animate ablative case clitic; from (a person marked with the possessive)}    \example{
      \exsen{Gagäll anyke da bäne män bo patatt a zäme agezänan.}      \extran{The bad spirit has already left your daughter.}}}
  \variants{\Dialectal Variant \vartext{=patalle}}}

\entry{=pate}{\headword{=pate}
  \pos{Nominal enclitic}  \sense{
    \definition{animate allative case clitic; to, towards, at (a person optionally marked with the possessive)}    \example{
      \exsen{Meri dallän Yokon pate.}      \extran{Mary went to Yokon.}}    \example{
      \exsen{Llamda da ddone ada mikutt gogon oba pate.}      \extran{The old man got very angry at them.}}}}

\entry{patepate}{\headword{patepate}  \note{Dancing}  \pronnote{patepate}
  \pos{Noun}  \sense{
    \definition{bamboo sticks used for percussion}    \note{Dancing}    \example{
}    \example{
}    \example{
}}}

\entry{Patha}{\headword{Patha}  \pronnote{paθa}
  \pos{Proper noun}  \sense{
    \definition{male personal name}}}

\entry{patiti}{\headword{patiti}  \note{Birds}  \pronnote{patiti}
  \pos{Noun}  \sense{
    \definition{type of bird}    \note{Birds}    \example{
}    \example{
}    \example{
}}}

\entry{patkoll}{\headword{patkoll}  \pronnote{patkoɽ}
  \pos{Noun}  \sense{
    \definition{bundle}    \example{
      \exsen{Ngäna yu de papllekdae bäga, bongo abo pattkolldae ignig a ada ma we ikom.}      \extran{I'll just chop the wood; then you just make the bundles and carry them home.}}    \example{
      \exsen{Ngäna kollba de patkoll däga nyäng peyang.}      \extran{I bundled the fish with the bag.}}}
  \variants{\Unspecified Variant \vartext{pattkoll}}}

\entry{patkopatkoll}{\headword{patkopatkoll}
  \pos{Adverb}  \sense{
    \definition{carrying bundles, with bundles}}}

\entry{=patme}{\headword{=patme}  \pronnote{patme}
  \pos{Nominal enclitic}  \sense{
    \definition{animate locative case clitic; at someone's place; for (a person marked with the possessive)}    \example{
      \exsen{Ngämo za da dag do Pauma bo patme.}      \extran{My things are there at Pauma's place.}}    \example{
      \exsen{Ngäna oba patme ada ulle gog.}      \extran{I was raised by them (lit. I grew up at theirs).}}    \example{
      \exsen{Bogo däpleon iba patme.}      \extran{He died for us.}}}}

\entry{pätpät}{\headword{pätpät}  \note{Water}  \pronnote{pətpət}
  \pos{Transitive verb}  \sense{
    \definition{to dry}    \note{Water}    \example{
}    \example{
}    \example{
}}}

\entry{patrol}{\headword{patrol}
  \pos{Noun}  \sense{
    \definition{patrol}}}

\entry{pätt}{\headword{pätt}  \pronnote{pəʈ͡ʂ}
  \pos{Noun}  \sense{\textbf{1.} 
    \definition{body (of a person or animal)}    \example{
      \exsen{Obo pätt a ulle daeya.}      \extran{Her body was large.}}    \example{
      \exsen{Ngämo pätt a llowamang gogon.}      \extran{My body was tired.}}    \example{
      \exsen{Ngäna käza bo pätt de net alle kaen dägne.}      \extran{I wrapped the crocodile's body with a net.}}}  \sense{\textbf{2.} 
    \definition{trunk of a plant in the ground; a single plant}    \example{
      \exsen{Nge pätt a tupi dan.}      \extran{The coconut palm trunk is long.}}    \example{
      \exsen{Ubi up wo pätt dowae e gogeyo.}      \extran{The two of them got close to the ripe banana tree.}}    \example{
      \exsen{Käp ku da wit pätt atta teti gongesnegän, ddob wit pätt atta siksti.}      \extran{Thirty seeds came from the wheat plant; from another wheat plant, sixty.}}    \example{
      \exsen{Bogo kunur pätt de däträpnegalle.}      \extran{He reaped the corn plants.}}}}

\entry{patt}{\headword{patt}  \lexeme{patt}  \pronnote{paʈ͡ʂ}
  \pos{Noun}  \sense{\textbf{1.} 
    \definition{tree trunk (fallen)}    \example{
      \exsen{Ngäna llo patt toko me godmen.}      \extran{I sat on top of a tree trunk.}}}  \sense{\textbf{2.} 
    \definition{drum shell}}  \sense{\textbf{3.} 
    \definition{coconut husking stick}    \example{
      \exsen{Gänyme nge kopnen ma patt de aeya dängkänän?}      \extran{Who pulled out the coconut husking sticks?}}}}

\entry{patt}{\headword{patt}  \lexeme{patt}  \pronnote{paʈ͡ʂ}
  \pos{Intransitive coverb}  \sense{
    \definition{to wake up late (past 8 am)}}}

\entry{Pätta}{\headword{Pätta}  \note{Places around Limol}  \pronnote{pəʈ͡ʂa}
  \pos{Proper noun}  \sense{
    \definition{Pätta (toponym)}    \note{Places around Limol}    \example{
}    \example{
}    \example{
}}}

\entry{pättä~}{\headword{pättä~}
  \variantof{freevarlab of \vartext{RED~}}}

\entry{pättäk}{\headword{pättäk}  \note{Adjectives}  \pronnote{pəʈ͡ʂək}
  \pos{Modifier}  \sense{
    \definition{short}    \note{Adjectives}    \example{
      \exsen{Ngäna era mälla tupi dan, bongo era pättäk dan.}      \extran{I'm a tall woman, but you're short.}}    \example{
}}}

\entry{pättäkpättäk}{\headword{pättäkpättäk}  \pronnote{pəʈ͡ʂəkpəʈ͡ʂək}
  \pos{Modifier}  \sense{
    \definition{nonsingular form of pättäk}    \example{
      \exsen{Ge ma da tämamae pättäkpättäk dag.}      \extran{All these houses are short.}}}
  \variants{\Fast Speech Variant \vartext{pättäpättäk}}}

\entry{pättäkpättäk}{\headword{pättäkpättäk}  \pronnote{pəʈ͡ʂəkpəʈ͡ʂək}
  \pos{Modifier}  \sense{
    \definition{a bit short}}
  \variants{\Fast Speech Variant \vartext{pättäpättäk}}}

\entry{pättäl}{\headword{pättäl}  \note{Trees}  \pronnote{pəʈ͡ʂəl}
  \pos{Noun}  \sense{
    \definition{type of tree that grows near swamps with yellow flowers and bark used for rope to tie firewood}    \note{Trees}    \example{
}    \example{
}    \example{
}}}

\entry{pättäpättäk}{\headword{pättäpättäk}
  \variantof{Fast Speech Variant of \vartext{pättäkpättäk}}}

\entry{pättapätte}{\headword{pättapätte}  \pronnote{pəʈ͡ʂapəʈ͡ʂe}
  \pos{Transitive verb}  \sense{\textbf{1.} 
    \definition{to shake off}    \example{
      \exsen{Bibi mɨnyi bina nying pänpän de näpttayaemneyo.}      \extran{You (all) will be shaking the dust from your feet.}}}  \sense{\textbf{2.} 
    \definition{to hit grass with a stick to scare away animals}}}

\entry{pättkäp}{\headword{pättkäp}  \note{Hunting}  \pronnote{pəʈ͡ʂkəp}
  \pos{Noun}  \sense{\textbf{1.} 
    \definition{node (ring or line on bamboo or sugarcane that separates internodes)}    \note{Hunting}}  \sense{\textbf{2.} 
    \definition{part of a bow}    \note{Hunting}    \example{
}    \example{
}    \example{
}}}

\entry{pattkoll}{\headword{pattkoll}
  \variantof{Unspecified Variant of \vartext{patkoll}}}

\entry{pattlle}{\headword{pattlle}  \lexeme{pattlle}  \pronnote{paʈ͡ʂɽe}
  \pos{Noun}  \sense{
    \definition{type of small bamboo that grows in the bush and along creeks; used to cook sago and make flutes; may be burned for fertile land for planting}}}

\entry{pättol}{\headword{pättol}
  \pos{Transitive verb}  \sense{
    \definition{to start singing}    \example{
      \exsen{Bogo bandra de dipttolän.}      \extran{He starting singing a song.}}}}

\entry{Paul}{\headword{Paul}
  \variantof{SpellingVariant of \vartext{Pol}}  \pronnote{paul}}

\entry{Pauma}{\headword{Pauma}  \pronnote{pauma}
  \pos{Proper noun}  \sense{
    \definition{female personal name}}}

\entry{pauro}{\headword{pauro}  \note{Birds}  \pronnote{pauro}
  \pos{Noun}  \sense{
    \definition{chicken}    \scientificname{Gallus gallus domesticus}    \note{Birds}    \example{
}    \example{
}    \example{
}}}

\entry{paus}{\headword{paus}
  \pos{}}

\entry{pawa}{\headword{pawa}  \pronnote{pawa}
  \pos{Noun}  \sense{
    \definition{power}}}

\entry{Pawaturi}{\headword{Pawaturi}
  \pos{Proper noun}  \sense{
    \definition{Pahoturi River}}}

\entry{paya}{\headword{paya}
  \pos{Transitive coverb}  \sense{
    \definition{to shoot, fire}}}

\entry{päzäg}{\headword{päzäg}  \lexeme{päzäg}  \pronnote{pəzəɡ}
  \pos{Kinship noun}  \sense{\textbf{1.} 
    \definition{brother-in-law (man's wife's brother or man's sister's husband; reciprocal)}}  \sense{\textbf{2.} 
    \definition{sister-in-law (man's wife's sister)}}
  \variants{\SpellingVariant \vartext{päzɨg}}}

\entry{päzäpäzäg}{\headword{päzäpäzäg}
  \pos{Kinship noun}  \sense{
    \definition{in-laws}}}

\entry{pazi}{\headword{pazi}  \note{Seasons}  \pronnote{pazi}
  \pos{Noun}  \sense{
    \definition{year}    \note{Seasons}    \example{
      \exsen{da pazi me, de pazi me}      \extran{last year, next year}}    \example{
      \exsen{kumuddäga pazi imnealle}      \extran{three years later}}    \example{
      \exsen{Bäne pazi da auli dan?}      \extran{How old are you (lit. how many are your years)?}}}}

\entry{päzɨg}{\headword{päzɨg}
  \variantof{SpellingVariant of \vartext{päzäg}}}

\entry{peba}{\headword{peba}  \note{School}  \pronnote{peba}
  \pos{Noun}  \sense{
    \definition{paper}    \note{School}    \example{
}    \example{
}    \example{
}}}

\entry{pedae}{\headword{pedae}  \pronnote{pedae}
  \pos{Transitive verb}  \sense{
    \definition{to sweep}    \example{
      \exsen{Ma de dapedealle.}      \extran{He sweeped the house.}}}}

\entry{Pedaya}{\headword{Pedaya}  \pronnote{pedaja}
  \pos{Proper noun}  \sense{
    \definition{Pedaya (toponym)}    \example{
}    \example{
}    \example{
}}}

\entry{Pedro}{\headword{Pedro}  \pronnote{pedro}
  \pos{Proper noun}  \sense{
    \definition{male personal name}}}

\entry{pekpek}{\headword{pekpek}
  \pos{Noun}  \sense{
    \definition{place}}}

\entry{pemli}{\headword{pemli}
  \pos{Noun}  \sense{
    \definition{family}}
  \variants{\SpellingVariant \vartext{femli, pamli, family}}}

\entry{pen}{\headword{pen}  \pronnote{pen}
  \pos{Noun}  \sense{
    \definition{pen}}}

\entry{penangg}{\headword{penangg}
  \variantof{Unspecified Variant of \vartext{penongg}}  \note{Gardens}  \pronnote{penaŋɡ}}

\entry{pendäg}{\headword{pendäg}  \pronnote{pendəɡ}
  \pos{Intransitive verb}  \sense{\textbf{1.} 
    \definition{to trip}    \example{
      \exsen{Gopendägalle ekaklle we.}      \extran{He was tripping and falling to the ground.}}}  \sense{\textbf{2.} 
    \definition{to push to the ground, jostle, trip}    \example{
      \exsen{Lla gul ulle daeya, Yesu adame obo kollmällang de umllang dägnegän oblle gall kälsäre särämbae e, oba lla da ddone obom beyawändägallo a bimpidägmällnallo.}      \extran{The crowd was big, so Jesus told his disciples to prepare a small boat for him so that people would not come crowd him and be pushing him.}}}}

\entry{penganyäm}{\headword{penganyäm}
  \pos{Noun}  \sense{
    \definition{type of yam}}}

\entry{pengg}{\headword{pengg}
  \pos{Transitive verb}  \sense{
    \definition{to make the killing shot, deliver the final blow}    \example{
      \exsen{Ddäg aeya nampegan obo ddäg dan.}      \extran{The back [of an animal] belongs to him, the one who delivered the final blow.}}}}

\entry{pengkameny}{\headword{pengkameny}
  \pos{Kinship noun}  \sense{
    \definition{sister}}}

\entry{penmällpenmäll}{\headword{penmällpenmäll}  \note{Colors}  \pronnote{penməɽpenməɽ}
  \pos{Modifier}  \sense{
    \definition{spotted}    \note{Colors}    \example{
}    \example{
}    \example{
}}}

\entry{penongg}{\headword{penongg}  \pronnote{penoŋɡ}
  \pos{Intransitive verb}  \sense{\textbf{1.} 
    \definition{to burn, set on fire, ignite}    \example{
      \exsen{Ge towall a mängalae penongg allan.}      \extran{This grass is burning fast.}}    \example{
      \exsen{Yäbäd a aempononggan.}      \extran{The sun blazed.}}}  \sense{\textbf{2.} 
    \definition{to burn, set on fire}    \example{
      \exsen{Ngäna towall penameny gongkam.}      \extran{I started to ignite the grass.}}    \example{
      \exsen{Yu daempononggeyo ma de.}      \extran{They set fire to the house.}}    \example{
      \exsen{Ause da to dänglläbänän, kullkull de daimponomenyän tämamae.}      \extran{The old woman got the light and burned the whole garden.}}}
  \variants{\Unspecified Variant \vartext{penangg}}}

\entry{pentae}{\headword{pentae}
  \pos{Intransitive verb}  \sense{\textbf{1.} 
    \definition{to spread, be transmitted}    \example{
      \exsen{Ge itrell a mɨnyi bompetaeyän.}      \extran{This disease will spread.}}    \example{
      \exsen{Ge itrell a mängalae lla yaba pate pentaeang dan.}      \extran{This disease spreads to others quickly.}}}  \sense{\textbf{2.} 
    \definition{to transfer, transmit, spread, pass on, pass down}    \example{
      \exsen{Ende lla da oba omog ttoen tameny de oba llɨg aba pate dampetaeneyo.}      \extran{Ende people were passing their sorcery on to their children.}}    \example{
}}}

\entry{Pentae}{\headword{Pentae}
  \pos{Proper noun}  \sense{
    \definition{female personal name}}
  \variants{\SpellingVariant \vartext{Pentai}}}

\entry{Pentai}{\headword{Pentai}
  \variantof{SpellingVariant of \vartext{Pentae}}  \pronnote{pentai}}

\entry{pentapentae}{\headword{pentapentae}  \pronnote{pentapentae}
  \pos{Adverb}  \sense{
    \definition{directly, in person, personally}    \example{
      \exsen{Ubi upe me dägabällabän adawatta ma ik a gonttämawän lla walle, be eka de pentapentae dandmoeyo ma ik e obo pate.}      \extran{They stood outside because the inside of the house was full of people, but they sent word directly to him inside.}}}}

\entry{pentngeny}{\headword{pentngeny}
  \pos{Intransitive verb}  \sense{
    \definition{to trip}    \example{
      \exsen{Bogo säremsärem gompetngenyän yu pallkoll me.}      \extran{In the darkness, he tripped over a piece of wood.}}}}

\entry{Penz}{\headword{Penz}  \pronnote{penz}
  \pos{Noun}  \sense{
    \definition{name of a river}}}

\entry{pepätt}{\headword{pepätt}  \pronnote{pepəʈ͡ʂ}
  \pos{Noun}  \sense{\textbf{1.} 
    \definition{black-faced monarch}    \scientificname{Monarcha melanopsis}}  \sense{\textbf{2.} 
    \definition{spot-winged monarch}    \scientificname{Symposiachrus guttula}}}

\entry{pepeb}{\headword{pepeb}  \note{Oral Traditions}  \pronnote{pepeb}
  \pos{Noun}  \sense{
    \definition{folktale, legend}    \note{Oral Traditions}    \anthronote{story about unseen people/spirits, stories about them, told only by the old people. These stories are being passed down. They have not been written down. Old people keep them in their minds, older women especially, and tell them to children.}    \example{
      \exsen{Ngämlle ttongo pepeb de nällät.}      \extran{[You] tell me one folktale.}}    \example{
}    \example{
}}}

\entry{pepeb eka}{\headword{pepeb eka}  \pronnote{pepeb eka}
  \pos{Noun}  \sense{
    \definition{legend}}}

\entry{pepeb wup}{\headword{pepeb wup}  \note{Names of bananas}  \pronnote{pepeb wup}
  \pos{Noun}  \sense{
    \definition{type of introduced banana}    \note{Names of bananas}    \example{
}    \example{
}    \example{
}}}

\entry{pera}{\headword{pera}  \pronnote{pera}
  \pos{Noun}  \sense{
    \definition{bird perch}}}

\entry{peraengg}{\headword{peraengg}
  \variantof{SpellingVariant of \vartext{peraingg}}}

\entry{peraenggag}{\headword{peraenggag}  \note{Source: Uli and Diwa}
  \pos{Noun}  \sense{
    \definition{type of dance with two groups}    \note{Source: Uli and Diwa}}}

\entry{peraingg}{\headword{peraingg}  \pronnote{peraiŋɡ}
  \pos{Intransitive verb}  \sense{\textbf{1.} 
    \definition{to cross, intersect}    \example{
      \exsen{Tawe da goimpregeyo.}      \extran{The tawe palms were crossed.}}    \example{
      \exsen{Ngäna Karama atta yaran, ngämi Winson alle aimpregalla.}      \extran{I came back from Karama; Winson and I had intersecting paths (but we didn't see each other).}}}  \sense{\textbf{2.} 
    \definition{to protrude, stick out}    \example{
      \exsen{Sɨmell ada daeya, ngoi peraingg att.}      \extran{The pig, its tusks were sticking out like this.}}}
  \variants{\SpellingVariant \vartext{peraengg}}}

\entry{perälla}{\headword{perälla}  \note{Palms}  \pronnote{perəɽa}
  \pos{Noun}  \sense{
    \definition{type of vine}    \note{Palms}    \example{
}    \example{
}    \example{
}}}

\entry{perapera}{\headword{perapera}  \pronnote{perapera}
  \pos{Adverb}  \sense{
    \definition{perching}}}

\entry{pes}{\headword{pes}
  \pos{Ordinal numeral}  \sense{
    \definition{first}    \example{
      \exsen{pes täräp}      \extran{the first time}}}}

\entry{petapeta}{\headword{petapeta}  \note{Adjectives}  \pronnote{petapeta}
  \pos{Modifier}  \sense{\textbf{1.} 
    \definition{thin (inanimate)}    \note{Adjectives}    \example{
      \exsen{Ge käg a petapeta dan}      \extran{The floor is thin.}}    \example{
}    \example{
}}  \sense{\textbf{2.} 
    \definition{shallow}    \note{Adjectives}    \example{
      \exsen{Källäm a petapeta abal gogon.}      \extran{The ponds became very shallow.}}}}

\entry{Petepo}{\headword{Petepo}  \pronnote{petepo}
  \pos{Proper noun}  \sense{
    \definition{female personal name}}}

\entry{Peter}{\headword{Peter}
  \variantof{SpellingVariant of \vartext{Pita}}}

\entry{Petom}{\headword{Petom}  \pronnote{petom}
  \pos{Proper noun}  \sense{
    \definition{Petom (toponym)}    \example{
}    \example{
}    \example{
}}}

\entry{petron}{\headword{petron}  \pronnote{petron}
  \pos{Noun}  \sense{
    \definition{blood vessel; vein; artery}}}

\entry{pett}{\headword{pett}
  \pos{Noun}  \sense{
}}

\entry{pewäl}{\headword{pewäl}  \note{Palms}  \pronnote{pewəl}
  \pos{Noun}  \sense{
    \definition{type of palm}    \note{Palms}    \example{
}    \example{
}    \example{
}}}

\entry{pewälewäle}{\headword{pewälewäle}  \pronnote{pewəlewəle}
  \pos{Noun}  \sense{
    \definition{green pygmy goose}    \scientificname{Nettapus pulchellus}    \example{
}    \example{
}    \example{
}}}

\entry{Pewe}{\headword{Pewe}  \pronnote{pewe}
  \pos{Proper noun}  \sense{
    \definition{male personal name}}}

\entry{peyam}{\headword{peyam}  \pronnote{pejam}
  \pos{Intransitive verb}  \sense{\textbf{1.} 
    \definition{to come out from water, surface}    \example{
      \exsen{Däbe käza da walle amne me gopeyamän.}      \extran{That crocodile surfaced in the middle of the pond.}}}  \sense{\textbf{2.} 
    \definition{to pop up, reappear, resurface}    \example{
      \exsen{Bongo gänyame opeyamalleǃ}      \extran{You're here again!}}}}

\entry{=peyang}{\headword{=peyang}  \note{Adjectives (Swadesh)}  \pronnote{pejaŋ}
  \pos{Nominal enclitic}  \sense{
    \definition{comitative clitic; with (a noun marked with the possessive)}    \note{Adjectives (Swadesh)}    \example{
      \exsen{Ngäna Grace bo peyang duwem agan.}      \extran{I ate with Grace.}}}}

\entry{peyom}{\headword{peyom}  \pronnote{pejom}
  \pos{Modifier}  \sense{
    \definition{mating}}}

\entry{Piasorosoro}{\headword{Piasorosoro}  \pronnote{piasorosoro}
  \pos{Proper noun}  \sense{
    \definition{male personal name}}}

\entry{pid}{\headword{pid}  \lexeme{pid}  \pronnote{pid}
  \pos{Noun}  \sense{
    \definition{horsefly}}}

\entry{Pidgin}{\headword{Pidgin}
  \variantof{SpellingVariant of \vartext{Pizin}}  \pronnote{pidɡin}}

\entry{pidor}{\headword{pidor}  \note{Birds}  \pronnote{pidor}
  \pos{Noun}  \sense{
    \definition{white-bellied sea-eagle}    \scientificname{Icthyophaga leucogaster}    \note{Birds}    \example{
}    \example{
}    \example{
}}}

\entry{Pidortama}{\headword{Pidortama}  \note{Places around Limol}  \pronnote{pidortama}
  \pos{Proper noun}  \sense{
    \definition{Pidortama (toponym)}    \note{Places around Limol}    \example{
}    \example{
}    \example{
}}}

\entry{pidroll}{\headword{pidroll}  \note{Grubs}  \pronnote{pidroɽ}
  \pos{Noun}  \sense{
    \definition{black palm weevil}    \scientificname{Rhynchophorus bilineatus}    \note{Grubs}    \example{
}    \example{
}    \example{
}}}

\entry{pig}{\headword{pig}  \pronnote{piɡ}
  \pos{Noun}  \sense{
    \definition{type of cultivated tree with edible fruit}    \example{
}    \example{
}    \example{
}}}

\entry{piksa}{\headword{piksa}
  \pos{Noun}  \sense{
    \definition{picture}}}

\entry{pilat}{\headword{pilat}
  \variantof{Unspecified Variant of \vartext{pilatt}}}

\entry{pilatt}{\headword{pilatt}
  \pos{Noun}  \sense{
    \definition{plate, dish}    \example{
      \exsen{Ngäna mɨnyi pilatt de bänggllämneg.}      \extran{I will wash the dishes.}}}
  \variants{\Unspecified Variant \vartext{pilat}}}

\entry{pimbyom}{\headword{pimbyom}  \pronnote{pimbjom}
  \pos{Noun}  \sense{
    \definition{small pieces of bark}}}

\entry{pin}{\headword{pin}  \note{Medicine}  \pronnote{pin}
  \pos{Noun}  \sense{
    \definition{type of tree with white or red flowers and composite fruit that birds eat}    \note{Medicine}    \grammarnote{boiled and given to children with urine problems. type of bottlebrush, but species name unknown.}    \example{
}    \example{
}    \example{
}}}

\entry{Pinang}{\headword{Pinang}
  \pos{Proper noun}  \sense{
    \definition{Pinang (toponym)}}}

\entry{ping}{\headword{ping}  \note{Clothing}  \pronnote{piŋ}
  \pos{Noun}  \sense{
    \definition{baby pin}    \note{Clothing}    \example{
}    \example{
}    \example{
}}}

\entry{Pingam}{\headword{Pingam}  \pronnote{piŋam}
  \pos{Proper noun}  \sense{
    \definition{female personal name}}}

\entry{pinggudd}{\headword{pinggudd}  \pronnote{piŋɡuɖ͡ʐ}
  \pos{Noun}  \sense{
    \definition{skirt}    \example{
      \exsen{Bogo eraeya pinggudd ulle peyang daeya.}      \extran{She was wearing a long skirt.}}}}

\entry{pinsäg}{\headword{pinsäg}
  \pos{Intransitive verb}  \sense{
    \definition{to tear}    \example{
      \exsen{Iddpo da komlla gumpisägeyo tuk alle do matu.}      \extran{The cloth tore in two from top to bottom.}}}}

\entry{pintopinto}{\headword{pintopinto}  \pronnote{pintopinto}
  \pos{Transitive coverb}  \sense{
    \definition{to blow leaves}}}

\entry{pintta}{\headword{pintta}  \pronnote{pinʈ͡ʂa}
  \pos{Noun}  \sense{
    \definition{palm cockatoo}    \scientificname{Probosciger aterrimus}    \example{
}    \example{
}    \example{
}}}

\entry{piny}{\headword{piny}  \pronnote{piɲ}
  \pos{Noun}  \sense{
    \definition{kingfisher (azure, little)}    \scientificname{Ceyx azureus, C. pusillus}    \example{
}    \example{
}    \example{
}}}

\entry{pinya dorollog}{\headword{pinya dorollog}  \pronnote{piɲa doroɽoɡ}
  \pos{Noun}  \sense{
    \definition{sacred kingfisher}    \scientificname{Todiramphus sanctus sanctus}    \example{
}    \example{
}    \example{
}}}

\entry{Pinzin}{\headword{Pinzin}
  \variantof{Dialectal Variant of \vartext{Pizin}}}

\entry{pinzopinzo}{\headword{pinzopinzo}
  \pos{Noun}  \sense{
    \definition{curly hair}}}

\entry{pinzopinzo}{\headword{pinzopinzo}  \note{Insects}  \pronnote{pinzopinzo}
  \pos{Noun}  \sense{
    \definition{type of small insect that lives in the ground}    \note{Insects}    \example{
}    \example{
}    \example{
}}}

\entry{pip}{\headword{pip}  \lexeme{pip}  \pronnote{pip}
  \pos{Noun}  \sense{
    \definition{red bee}}}

\entry{pipi}{\headword{pipi}  \pronnote{pipi}
  \pos{Transitive verb}  \sense{
    \definition{to shoot, spear}    \example{
      \exsen{Ngäna ddob mozaya kollba de däpineg.}      \extran{I speared some mozaya fish.}}}}

\entry{Pipi}{\headword{Pipi}
  \pos{Proper noun}  \sense{
    \definition{female personal name}}}

\entry{Pipiato}{\headword{Pipiato}
  \pos{Proper noun}  \sense{
    \definition{female personal name}}}

\entry{pipllo}{\headword{pipllo}  \lexeme{pipllo}  \pronnote{pipɽo}
  \pos{Noun}  \sense{
    \definition{lizard}}}

\entry{pipllug}{\headword{pipllug}  \pronnote{pipɽuɡ}
  \pos{Intransitive verb}  \sense{
    \definition{to fly}    \example{
      \exsen{Bongo mɨnyi ikop nägneg ddob apapi di bipllugaebnän.}      \extran{You will see some butterflies flying.}}}}

\entry{pipti}{\headword{pipti}
  \pos{Numeral}  \sense{
    \definition{fifty}}
  \variants{\SpellingVariant \vartext{fifti}}}

\entry{piptin}{\headword{piptin}
  \pos{Numeral}  \sense{
    \definition{fifteen (English numeral)}}
  \variants{\SpellingVariant \vartext{fiftin}}}

\entry{pirik}{\headword{pirik}  \note{Tools}  \pronnote{pirik}
  \pos{Noun}  \sense{
    \definition{baton, stick}    \note{Tools}    \example{
      \exsen{Pirik a lla metmäll ma dan.}      \extran{A baton is for hitting people.}}    \example{
}    \example{
}}}

\entry{pirngän}{\headword{pirngän}
  \pos{Transitive verb}  \sense{
    \definition{to draw a weapon, take out}    \example{
      \exsen{Ngäna pakos de dapirngän.}      \extran{I took out a spear.}}}}

\entry{piro}{\headword{piro}  \lexeme{piro}  \pronnote{piro}
  \pos{Noun}  \sense{\textbf{1.} 
    \definition{star}}  \sense{\textbf{2.} 
    \definition{star weaving pattern}}}

\entry{piro kanas}{\headword{piro kanas}  \pronnote{piro kanas}
  \pos{Noun}  \sense{
    \definition{type of game where the first person to see a star in the sky wins}}}

\entry{piro ttängattänge}{\headword{piro ttängattänge}  \pronnote{piro ʈ͡ʂəŋaʈ͡ʂəŋe}
  \pos{Intransitive compound verb}  \sense{
    \definition{to become unconscious}    \example{
      \exsen{Bogo piro nangttängenan.}      \extran{He fell down unconcsious.}}}}

\entry{pisam}{\headword{pisam}  \pronnote{pisam}
  \pos{Intransitive verb}  \sense{\textbf{1.} 
    \definition{to tear (of something thin)}    \example{
      \exsen{Ngämo kaptte da apisaman.}      \extran{My garment tore.}}}  \sense{\textbf{2.} 
    \definition{to tear (something thin)}    \example{
      \exsen{Ngäna nyäg de napisaman.}      \extran{I broke the basket.}}    \example{
      \exsen{Ngäna peba de pisam eran.}      \extran{I am tearing the paper.}}    \example{
}}}

\entry{pisepise}{\headword{pisepise}  \pronnote{pisepise}
  \pos{Noun}  \sense{
    \definition{small pieces}    \example{
      \exsen{Llaeyabira ako midd pisepise dikomeya ma we.}      \extran{We brought small pieces of meat for everyone.}}}}

\entry{Pisi}{\headword{Pisi}  \pronnote{pisi}
  \pos{Proper noun}  \sense{
    \definition{Pisi (in Gogodala Rural LLG)}    \example{
}    \example{
}    \example{
}}}

\entry{Piskae}{\headword{Piskae}  \note{Places around Limol}  \pronnote{piskae}
  \pos{Proper noun}  \sense{
    \definition{Piskae (toponym)}    \note{Places around Limol}    \example{
}    \example{
}    \example{
}}}

\entry{pisor}{\headword{pisor}
  \pos{Noun}  \sense{
}}

\entry{pit}{\headword{pit}  \pronnote{pit}
  \pos{Noun}  \sense{
    \definition{small insect that lives in the swamp and eats taro and sweet potato}}}

\entry{Pita}{\headword{Pita}  \pronnote{pita}
  \pos{Proper noun}  \sense{
    \definition{male personal name}}
  \variants{\SpellingVariant \vartext{Peter}}}

\entry{pitatep}{\headword{pitatep}
  \pos{Verb}  \sense{\textbf{1.} 
    \definition{to lift an animal or person and throw them to the ground, hit them hard against the ground}}  \sense{\textbf{2.} 
    \definition{to tackle someone, as in rugby}}}

\entry{pite}{\headword{pite}  \note{Source: Jack Dipa}
  \pos{Noun}  \sense{
    \definition{large winged ankom}    \note{Source: Jack Dipa}}}

\entry{pite}{\headword{pite}  \note{Clothing}  \pronnote{pite}
  \pos{Noun}  \sense{\textbf{1.} 
    \definition{grass skirt}    \note{Clothing}    \example{
      \exsen{Ingnenang aba pite da mermerae abal beräberäl agnegnan iddob me ingnen täräp me.}      \extran{The dancers' grass skirts swung nicely during the night dance.}}}  \sense{\textbf{2.} 
    \definition{tailfeather}    \note{Clothing}    \example{
      \exsen{Kakayam pa bo pite da mer tutpi dag.}      \extran{The tailfeathers of a bird of paradise are nice and long.}}}}

\entry{pitepite}{\headword{pitepite}  \note{Clothing}  \pronnote{pitepite}
  \pos{Adverb}  \sense{
    \definition{under the skirt}    \note{Clothing}    \example{
}    \example{
}    \example{
}}}

\entry{Pitepo}{\headword{Pitepo}
  \pos{Proper noun}  \sense{
    \definition{female personal name}}}

\entry{pitkae}{\headword{pitkae}  \pronnote{pitkae}
  \pos{Transitive verb}  \sense{
    \definition{to untie}    \example{
      \exsen{Ngäna mɨnyi bomällkam a obo nying kollop tär de bäpitkaeneg.}      \extran{I will bend down and untie his sandal laces.}}}}

\entry{pitpit}{\headword{pitpit}  \pronnote{pitpit}
  \pos{Noun}  \sense{
    \definition{sugarcane}    \scientificname{Saccharum edule}}}

\entry{pitraempäg}{\headword{pitraempäg}  \pronnote{pitraempəɡ}
  \pos{Transitive verb}  \sense{
    \definition{to make a failed shot, shoot an arrow that falls}}}

\entry{pitratra}{\headword{pitratra}  \pronnote{pitratra}
  \pos{Noun}  \sense{
    \definition{masked lapwing}    \scientificname{Vanellus miles}    \example{
}    \example{
}    \example{
}}}

\entry{pitt}{\headword{pitt}  \pronnote{piʈ͡ʂ}
  \pos{Noun}  \sense{\textbf{1.} 
    \definition{arrowhead hafting string}}  \sense{\textbf{2.} 
    \definition{band}    \example{
      \exsen{ttang pitt, ttäle pitt}      \extran{bracelet, leg band}}}}

\entry{pittäpen}{\headword{pittäpen}  \pronnote{piʈ͡ʂəpen}
  \pos{Transitive verb}  \sense{
    \definition{to drop or bump someone by holding their two legs and hands. To hold an animal (such as a small wallaby) by two legs and swing and whack them against the ground or a tree to kill them}}
  \variants{\SpellingVariant \vartext{pittäpin}}}

\entry{pittäpin}{\headword{pittäpin}
  \variantof{SpellingVariant of \vartext{pittäpen}}}

\entry{pittbudar}{\headword{pittbudar}  \note{Grubs}  \pronnote{piʈ͡ʂbudar}
  \pos{Noun}  \sense{
    \definition{type of edible grub found in the bush and grassland}    \note{Grubs}    \example{
}    \example{
}    \example{
}}}

\entry{pittpitt}{\headword{pittpitt}  \note{Swadesh}  \pronnote{piʈ͡ʂpiʈ͡ʂ}
  \pos{Transitive verb}  \sense{\textbf{1.} 
    \definition{to weave}    \note{Swadesh}    \example{
      \exsen{Sana dor de adingoll däpittaemeya.}      \extran{We weaved the sago stalks like this.}}}  \sense{\textbf{2.} 
    \definition{to sew, stitch}    \note{Swadesh}    \example{
      \exsen{Amo umllang dan pittpitt a?}      \extran{Who here knows how to sew?}}    \example{
      \exsen{Pittnen de dängkaemnegän.}      \extran{She started to stitch (the wounds).}}    \example{
}}}

\entry{piya}{\headword{piya}  \lexeme{piya}  \pronnote{pija}
  \pos{Noun}  \sense{
    \definition{type of flowering plant}}}

\entry{piya gany}{\headword{piya gany}  \note{Conflict/war}  \pronnote{pjia ɡaɲ}
  \pos{Noun}  \sense{
    \definition{planting a piya plant as a gesture of peace}    \note{Conflict/war}    \example{
}    \example{
}    \example{
}}}

\entry{piya ibeny}{\headword{piya ibeny}  \pronnote{pija ibeɲ}
  \pos{Noun}  \sense{
    \definition{planting a piya plant as a gesture of peace}}}

\entry{piye}{\headword{piye}  \note{Parts of a bird}  \pronnote{pije}
  \pos{Noun}  \sense{\textbf{1.} 
    \definition{pus}    \note{Parts of a bird}}  \sense{\textbf{2.} 
    \definition{sperm}    \note{Parts of a bird}    \example{
}    \example{
}}}

\entry{piyepiye}{\headword{piyepiye}
  \pos{Noun}  \sense{
    \definition{blister}}}

\entry{piyupiyu}{\headword{piyupiyu}  \lexeme{piyupiyu}  \pronnote{pijupiju}
  \pos{Noun}  \sense{
    \definition{large vine that grows from tree to tree in the bush}}}

\entry{piyupiyu budar}{\headword{piyupiyu budar}  \note{Grubs}  \pronnote{pijupiju budar}
  \pos{Noun}  \sense{
    \definition{type of grub}    \note{Grubs}    \example{
}    \example{
}    \example{
}}}

\entry{Pizi}{\headword{Pizi}
  \pos{Proper noun}  \sense{
    \definition{Pizi (toponym)}}}

\entry{pizin}{\headword{pizin}
  \variantof{SpellingVariant of \vartext{Pizin}}}

\entry{Pizin}{\headword{Pizin}
  \pos{Proper noun}  \sense{
    \definition{Tok Pisin}}
  \variants{\Dialectal Variant \vartext{Pinzin}}
  \variants{\SpellingVariant \vartext{Pidgin, pizin}}}

\entry{pɨddab}{\headword{pɨddab}
  \variantof{Unspecified Variant of \vartext{päddab}}  \pronnote{pɪɖ͡ʐab}}

\entry{pɨddpɨdd}{\headword{pɨddpɨdd}
  \variantof{SpellingVariant of \vartext{päddpädd}}}

\entry{pɨnpɨn}{\headword{pɨnpɨn}
  \variantof{Unspecified Variant of \vartext{pänpän}}  \pronnote{pɪnpɪn}}

\entry{pɨnyapɨnye}{\headword{pɨnyapɨnye}  \pronnote{pɪɲapɪɲe}
  \pos{Noun}  \sense{
    \definition{area with burnt grass}    \example{
      \exsen{Ngäna pɨnyapɨnye koenmäll e ibi allan.}      \extran{I'm going to go hunting in the burnt grass area.}}}}

\entry{pɨnypɨny}{\headword{pɨnypɨny}  \pronnote{pɪɲpɪɲ}
  \pos{Transitive verb}  \sense{\textbf{1.} 
    \definition{to plant a lot}}  \sense{\textbf{2.} 
    \definition{many stars appearing}    \example{
      \exsen{Piro da pɨnypɨny eran.}      \extran{Many stars are coming out.}}}}

\entry{plaig}{\headword{plaig}  \note{School}  \pronnote{plaiɡ}
  \pos{Noun}  \sense{
    \definition{flag}    \note{School}    \example{
}    \example{
}    \example{
}}}

\entry{plein}{\headword{plein}
  \variantof{SpellingVariant of \vartext{plen}}}

\entry{plen}{\headword{plen}
  \pos{Noun}  \sense{
    \definition{airplane}}
  \variants{\SpellingVariant \vartext{plein}}}

\entry{plengg}{\headword{plengg}  \pronnote{pleŋɡ}
  \pos{Intransitive verb}  \sense{\textbf{1.} 
    \definition{to die}    \example{
      \exsen{Bogo däplewän.}      \extran{He died.}}}  \sense{\textbf{2.} 
    \definition{to cut down, cut off}    \example{
      \exsen{Ngäna llo de plengg eran.}      \extran{I am cutting down the tree.}}}}

\entry{plenngätt}{\headword{plenngätt}
  \pos{Noun}  \sense{
    \definition{airstrip}}}

\entry{plenplen}{\headword{plenplen}  \pronnote{plenplen}
  \pos{Noun}  \sense{
    \definition{pinwheel}}}

\entry{pllaengän}{\headword{pllaengän}
  \pos{Verb}  \sense{
}}

\entry{pllulle}{\headword{pllulle}
  \variantof{Unspecified Variant of \vartext{pällulle}}}

\entry{pllulleaga}{\headword{pllulleaga}
  \pos{Noun}  \sense{
    \definition{the seventh and final stage of sago growth during which the pith will not yield any starch}}}

\entry{pllulli}{\headword{pllulli}
  \variantof{Unspecified Variant of \vartext{pällulle}}}

\entry{pllutt}{\headword{pllutt}  \pronnote{pɽuʈ͡ʂ}
  \pos{Noun}  \sense{
    \definition{type of vine with fruit that are yellow when ripe; children eat them}}}

\entry{po}{\headword{po}  \pronnote{po}
  \pos{Noun}  \sense{
    \definition{mound of dirt around a plant}    \example{
      \exsen{Nae po de gagäll a mudan.}      \extran{Don't ruin the sweet potato mound.}}}}

\entry{po}{\headword{po}  \pronnote{po}
  \pos{Transitive verb}  \sense{\textbf{1.} 
    \definition{to pierce}    \example{
      \exsen{Mälläng e dapoeaemeyo.}      \extran{They pierced noses.}}}  \sense{\textbf{2.} 
    \definition{to split}}}

\entry{po}{\headword{po}  \pronnote{po}
  \pos{Numeral}  \sense{
    \definition{four (English numeral; also general)}}
  \variants{\SpellingVariant \vartext{four, fo, poo}}}

\entry{po~}{\headword{po~}
  \variantof{Infinitival reduplicant of \vartext{RED~}}}

\entry{po}{\headword{po}  \pronnote{po}
  \pos{Intransitive coverb}  \sense{
    \definition{to block an animal from escaping}    \example{
      \exsen{Bongo ngattong e po ag.}      \extran{You, block in front.}}    \example{
}}}

\entry{poapoa}{\headword{poapoa}  \pronnote{poapoa}
  \pos{Modifier}  \sense{\textbf{1.} 
    \definition{light (in weight)}}  \sense{\textbf{2.} 
    \definition{light, bright}    \example{
      \exsen{Bädab pwapwa agnan.}      \extran{Dawn was breaking.}}}  \sense{\textbf{3.} 
    \definition{easy}    \example{
      \exsen{Poapoa bogon.}      \extran{It will be easier.}}}
  \variants{\SpellingVariant \vartext{pwapwa}}}

\entry{pob}{\headword{pob}  \pronnote{pob}
  \pos{Noun}  \sense{
    \definition{savanna}}}

\entry{pobllem}{\headword{pobllem}  \note{Adjectives}  \pronnote{pobɽem}
  \pos{Modifier}  \sense{
    \definition{smooth}    \note{Adjectives}    \example{
      \exsen{pobllem ibima}      \extran{smooth path}}    \example{
}    \example{
}}}

\entry{Podar}{\headword{Podar}
  \variantof{Unspecified Variant of \vartext{Podare}}}

\entry{Podare}{\headword{Podare}  \pronnote{podare}
  \pos{Proper noun}  \sense{
    \definition{Podare (Wipi-speaking village in Oriomo-Bituri Rural LLG)}    \example{
}    \example{
}    \example{
}}
  \variants{\Unspecified Variant \vartext{Podar}}}

\entry{podd}{\headword{podd}  \pronnote{poɖ͡ʐ}
  \pos{Noun}  \sense{
    \definition{bald head, baldness}    \example{
      \exsen{Zäme podd a gongesän.}      \extran{He is already bald.}}}}

\entry{poddem}{\headword{poddem}  \pronnote{poɖ͡ʐem}
  \pos{Noun}  \sense{
    \definition{young or small mammal (e.g. deer, wallaby)}}}

\entry{poddpodd}{\headword{poddpodd}  \pronnote{poɖ͡ʐpoɖ͡ʐ}
  \pos{Noun}  \sense{
    \definition{plain, field, clearing}}}

\entry{poder}{\headword{poder}  \pronnote{poder}
  \pos{Noun}  \sense{
    \definition{scapula, shoulder blade}    \example{
}    \example{
}    \example{
}}}

\entry{Pogo}{\headword{Pogo}
  \pos{Proper noun}  \sense{
    \definition{Pogo (toponym)}}}

\entry{pogpog}{\headword{pogpog}
  \pos{Noun}  \sense{
}}

\entry{Pol}{\headword{Pol}  \pronnote{pol}
  \pos{Proper noun}  \sense{
    \definition{male personal name}}
  \variants{\SpellingVariant \vartext{Paul}}}

\entry{polen}{\headword{polen}  \note{School}  \pronnote{polen}
  \pos{Noun}  \sense{
    \definition{assembly}    \note{School}    \example{
}    \example{
}    \example{
}}}

\entry{Polin}{\headword{Polin}  \pronnote{polin}
  \pos{Proper noun}  \sense{
    \definition{female personal name}}}

\entry{polis}{\headword{polis}
  \pos{Noun}  \sense{
    \definition{police}    \example{
      \exsen{Polis a namllamallo ge lla de.}      \extran{Police seized this man.}}}}

\entry{Poll}{\headword{Poll}
  \pos{Proper noun}  \sense{
    \definition{male personal name}}}

\entry{poll~}{\headword{poll~}
  \variantof{Infinitival reduplicant of \vartext{RED~}}}

\entry{polle}{\headword{polle}  \lexeme{polle}  \pronnote{poɽe}
  \pos{Noun}  \sense{
    \definition{fence}    \example{
      \exsen{Ubi ttängäm kuddäll me polle de kame däkätteyo.}      \extran{They fenced the old garden.}}}}

\entry{polle~}{\headword{polle~}
  \variantof{Derivational reduplicant of \vartext{RED~}}}

\entry{polle bo}{\headword{polle bo}  \note{Gardens}  \pronnote{poɽe bo}
  \pos{Noun}  \sense{
    \definition{base of fence}    \note{Gardens}    \example{
}    \example{
}    \example{
}}}

\entry{pollepolle}{\headword{pollepolle}
  \pos{Noun}  \sense{
    \definition{small fence}}}

\entry{pollgo}{\headword{pollgo}  \lexeme{pollgo}  \pronnote{poɽɡo}
  \pos{Noun}  \sense{
    \definition{frog}    \example{
      \exsen{Pollgo da täl ik mi giddollag daeya.}      \extran{The frog lived in the bamboo.}}}}

\entry{pollgo suwe}{\headword{pollgo suwe}  \pronnote{poɽɡo suwe}
  \pos{Noun}  \sense{
    \definition{dirt left on skin after coming out of the water}    \example{
      \exsen{Bäne pätt a pollgo suwe peyang dan.}      \extran{Your body is covered in dirt from the water.}}}}

\entry{pollgo ttatta kuttang}{\headword{pollgo ttatta kuttang}  \note{Weaving}  \pronnote{poɽɡo ʈ͡ʂaʈ͡ʂa kuʈ͡ʂaŋ}
  \pos{Noun}  \sense{
    \definition{weaving pattern with concentric diamonds}    \note{Weaving}    \example{
}    \example{
}    \example{
}}}

\entry{pollo}{\headword{pollo}  \pronnote{poɽo}
  \pos{Modifier}  \sense{
    \definition{young}    \example{
      \exsen{pollo täräp}      \extran{youth}}    \example{
      \exsen{llɨg pollo}      \extran{young boy}}}}

\entry{pollon~}{\headword{pollon~}
  \variantof{Derived Variant of \vartext{RED~}}}

\entry{pollon}{\headword{pollon}  \pronnote{poɽon}
  \pos{Noun}  \sense{
    \definition{type of small bush}    \example{
      \exsen{Bogo gotarnän pollon me.}      \extran{He was sleeping in a small bush.}}}}

\entry{pollon bällabollott}{\headword{pollon bällabollott}  \note{Types of Taro}  \pronnote{poɽon bəɽaboɽoʈ͡ʂ}
  \pos{Noun}  \sense{
    \definition{type of big taro}    \note{Types of Taro}    \example{
}    \example{
}    \example{
}}}

\entry{pollonpollon}{\headword{pollonpollon}
  \pos{Noun}  \sense{
    \definition{nonsingular form of pollon}}
  \variants{\Unspecified Variant \vartext{pällonpällon}}}

\entry{pollpoll}{\headword{pollpoll}  \note{Animal verbs}  \pronnote{poɽpoɽ}
  \pos{Transitive verb}  \sense{
    \definition{to bark (at)}    \note{Animal verbs}    \example{
      \exsen{Ge däräng a sämell ulle de dupollneyo.}      \extran{These dogs were barking at the big pig.}}    \example{
      \exsen{Kwalde bogo ada gongnomenynän ada ka däräng a pollnenang me erallo däbe sämell de.}      \extran{Kwalde thought that the dogs were still barking at the pig.}}    \example{
}}}

\entry{poloengg}{\headword{poloengg}  \pronnote{poloeŋɡ}
  \pos{Transitive verb}  \sense{
    \definition{to insult}}}

\entry{poloengg eka}{\headword{poloengg eka}
  \pos{Noun}  \sense{
}}

\entry{pom}{\headword{pom}  \pronnote{pom}
  \pos{Noun}  \sense{
    \definition{type of long yam with a white interior, thorns, and hairs}    \example{
}    \example{
}    \example{
}}}

\entry{poma}{\headword{poma}  \note{Type of pandanus}  \pronnote{poma}
  \pos{Noun}  \sense{
    \definition{type of pandanus with leaves that women weave into mats}    \note{Type of pandanus}    \example{
}    \example{
}    \example{
}}}

\entry{pomana}{\headword{pomana}  \note{Fishing}  \pronnote{pomana}
  \pos{Noun}  \sense{
    \definition{fishing style in which men climb trees and shoot with ddangol spears}    \note{Fishing}    \example{
}    \example{
}    \example{
}}}

\entry{pomer}{\headword{pomer}  \pronnote{pomer}
  \pos{Noun}  \sense{
    \definition{whistle}    \example{
      \exsen{Ge lla da pomer alle bandra nällɨtan.}      \extran{This man whistled a song.}}}}

\entry{pomila}{\headword{pomila}  \note{Trees}  \pronnote{pomila}
  \pos{Noun}  \sense{
    \definition{type of citrus tree with big, edible fruit}    \note{Trees}    \example{
}    \example{
}    \example{
}}}

\entry{pondo}{\headword{pondo}  \note{Trees}  \pronnote{pondo}
  \pos{Noun}  \sense{
    \definition{type of tree that grows in the swamp and old gardens with yellow flowers and pod-shaped fruit with inedible seeds}    \note{Trees}    \example{
}    \example{
}    \example{
}}}

\entry{Pondollowang}{\headword{Pondollowang}  \note{Places around Limol}  \pronnote{pondoɽowaŋ}
  \pos{Proper noun}  \sense{
    \definition{Pondollowang (camping place)}    \note{Places around Limol}    \example{
}    \example{
}    \example{
}}}

\entry{ponen ma}{\headword{ponen ma}  \note{School}  \pronnote{ponen ma}
  \pos{Noun}  \sense{
    \definition{pencil sharpener}    \note{School}    \example{
}    \example{
}    \example{
}}}

\entry{ponganem}{\headword{ponganem}  \lexeme{ponganem}  \pronnote{poŋanem}
  \pos{Noun}  \sense{
    \definition{type of small yam with a white interior, hairs, and thorns}}}

\entry{Pongarke}{\headword{Pongarke}  \pronnote{poŋarke}
  \pos{Proper noun}  \sense{
    \definition{Pongariki (Nambo-speaking village in Morehead Rural LLG)}    \example{
}    \example{
}    \example{
}}}

\entry{pongo}{\headword{pongo}  \note{Snakes}  \pronnote{poŋo}
  \pos{Noun}  \sense{
    \definition{scent}    \note{Snakes}    \example{
}    \example{
}    \example{
}}}

\entry{pongoll}{\headword{pongoll}  \note{Food}  \pronnote{poŋoɽ}
  \pos{Noun}  \sense{
    \definition{type of yam with a white interior and no hairs or thorns}    \note{Food}    \example{
}    \example{
}}}

\entry{pongpongäll}{\headword{pongpongäll}  \note{Parts of the body}  \pronnote{poŋpoŋəɽ}
  \pos{Noun}  \sense{
    \definition{body part}    \note{Parts of the body}    \example{
}    \example{
}    \example{
}}}

\entry{ponong}{\headword{ponong}  \note{Trees}  \pronnote{ponoŋ}
  \pos{Noun}  \sense{
    \definition{type of tree that grows in grassland with white flowers and small green fruit and wood used for posts}    \note{Trees}    \example{
}    \example{
}    \example{
}}}

\entry{Ponongllowang}{\headword{Ponongllowang}  \note{Place names around Limol}  \pronnote{ponoŋɽowaŋ}
  \pos{Proper noun}  \sense{
    \definition{Ponongllowang (toponym)}    \note{Place names around Limol}    \example{
}    \example{
}    \example{
}}}

\entry{ponor}{\headword{ponor}  \pronnote{ponor}
  \pos{Intransitive verb}  \sense{
    \definition{to start running}    \example{
      \exsen{Angde ubi up wo pätt dowae e gogeyo, llamda da oba pate eraeranyang dipnurän.}      \extran{When they got close to the ripe banana tree, the old man began to run after them, yelling.}}}}

\entry{pop}{\headword{pop}  \pronnote{pop}
  \pos{Noun}  \sense{
    \definition{hole}    \example{
      \exsen{Tupi ttang llɨpɨt de pop me nowanseg.}      \extran{Put your long finger in the hole.}}}}

\entry{popang}{\headword{popang}  \pronnote{popaŋ}
  \pos{Modifier}  \sense{\textbf{1.} 
    \definition{holey, having a hole}    \example{
      \exsen{Ine kutt da bo lid a popang dan.}      \extran{The lid of the water container has a hole.}}}  \sense{\textbf{2.} 
    \definition{rubbish, darn}}}

\entry{pope}{\headword{pope}  \lexeme{pope}  \pronnote{pope}
  \pos{Kinship noun}  \sense{
    \definition{uncle (one's mother's brother)}    \example{
      \exsen{Lomae obo pope mu de zäme dangesän.}      \extran{Lomae already paid her uncle her birth payment.}}}}

\entry{pope mu}{\headword{pope mu}  \note{Conflict/war}  \pronnote{pope mu}
  \pos{Noun}  \sense{
    \definition{uncle payment}    \note{Conflict/war}    \example{
}    \example{
}    \example{
}}}

\entry{popel}{\headword{popel}
  \variantof{Unspecified Variant of \vartext{popell}}}

\entry{popell}{\headword{popell}  \note{Names of bananas}  \pronnote{popeɽ}
  \pos{Noun}  \sense{\textbf{1.} 
    \definition{type of introduced banana}    \note{Names of bananas}    \example{
}    \example{
}    \example{
}}  \sense{\textbf{2.} 
    \definition{dish consisting of ants and various types of bananas (incl. popell)}    \note{Names of bananas}}
  \variants{\Unspecified Variant \vartext{popel}}}

\entry{poper}{\headword{poper}  \note{Verbs}  \pronnote{poper}
  \pos{Modifier}  \sense{
    \definition{afraid, startled, scared}    \note{Verbs}    \example{
      \exsen{Pa poper gogon a ma ik atta gogezänän.}      \extran{The bird got scared and got out from inside the house.}}    \example{
      \exsen{Lelang eka dändär att me tikop poper agne.}      \extran{The heart will be racing after hearing a scary story.}}    \example{
}}}

\entry{poper ingkoeg}{\headword{poper ingkoeg}
  \pos{Transitive compound verb}  \sense{
    \definition{to scare, startle, frighten}    \example{
      \exsen{Poper naengkoegan.}      \extran{He scared her.}}    \example{
      \exsen{Mänmän de poper inkoegag bongttägän.}      \extran{He will arrive suddenly, startling the girls.}}}}

\entry{poper kwingg}{\headword{poper kwingg}  \pronnote{poper kwiŋɡ}
  \pos{Transitive compound verb}  \sense{
    \definition{to surprise}    \example{
      \exsen{Ngäna obom poper nangkwigan.}      \extran{I surprised her.}}}}

\entry{popllem}{\headword{popllem}
  \pos{Intransitive verb}  \sense{
    \definition{to flap}    \example{
      \exsen{Kaptte da wel me angde popllem anggan, dedam era mamaldae pätnan anggan.}      \extran{When the clothes flap in the wind, they dry quickly.}}}}

\entry{popllo}{\headword{popllo}
  \pos{Verb}  \sense{
}}

\entry{popo}{\headword{popo}  \lexeme{popo}  \pronnote{popo}
  \pos{Noun}  \sense{
    \definition{flower}}}

\entry{popo}{\headword{popo}  \lexeme{popo}  \pronnote{popo}
  \pos{Transitive verb}  \sense{\textbf{1.} 
    \definition{to sharpen}    \example{
      \exsen{Dompa de däpoaemneyo.}      \extran{They were sharpening the arrows.}}    \example{
      \exsen{Ge pensil a popoatt dan.}      \extran{This pencil is sharpened.}}}  \sense{\textbf{2.} 
    \definition{to carve}    \example{
      \exsen{Gall a llo popoatt dan.}      \extran{The canoe is carved wood.}}    \example{
      \exsen{Bägäl de popo eran.}      \extran{He shapes the bow.}}}}

\entry{porak}{\headword{porak}
  \pos{Noun}  \sense{
    \definition{type of spear}}}

\entry{porma}{\headword{porma}  \pronnote{porma}
  \pos{Noun}  \sense{
    \definition{traditional dish consisting of meat on top of sago cooked in sago or banana leaves}    \example{
}}}

\entry{pos}{\headword{pos}  \pronnote{pos}
  \pos{Noun}  \sense{
    \definition{post}}}

\entry{poses}{\headword{poses}  \pronnote{poses}
  \pos{Noun}  \sense{\textbf{1.} 
    \definition{martin (fairy, tree)}    \scientificname{Petrochelidon ariel, P. nigricans}    \example{
}    \example{
}    \example{
}}  \sense{\textbf{2.} 
    \definition{Pacific swift}    \scientificname{Apus pacificus}}}

\entry{pot}{\headword{pot}  \pronnote{pot}
  \pos{Noun}  \sense{\textbf{1.} 
    \definition{tip, end, point; base (of yam)}}  \sense{\textbf{2.} 
    \definition{talent, gift}}  \sense{\textbf{3.} 
    \definition{vagina of a mammal}}}

\entry{Pot Mosbi}{\headword{Pot Mosbi}
  \pos{Proper noun}  \sense{
    \definition{Port Moresby (the capital city of Papua New Guinea)}}
  \variants{\Fast Speech Variant \vartext{Mosbi}}}

\entry{poti}{\headword{poti}
  \pos{Numeral}  \sense{
    \definition{forty}}
  \variants{\SpellingVariant \vartext{forty, foti, forti}}}

\entry{potin}{\headword{potin}
  \pos{Numeral}  \sense{
    \definition{fourteen (English numeral)}}}

\entry{potkam}{\headword{potkam}  \note{Trees}  \pronnote{potkam}
  \pos{Noun}  \sense{
    \definition{type of cultivated tree that also grows in the savanna; treats cough and aches}    \note{Trees}    \example{
}    \example{
}    \example{
}}}

\entry{potne kätt}{\headword{potne kätt}
  \pos{Noun}  \sense{
}}

\entry{potopoto}{\headword{potopoto}  \note{Medicine}  \pronnote{potopoto}
  \pos{Noun}  \sense{
    \definition{type of tree that grows near swamps and creeks with yellow and white flowers and fruit that falls on the ground; animals eat it}    \note{Medicine}    \example{
}    \example{
}    \example{
}}}

\entry{potpot}{\headword{potpot}  \pronnote{potpot}
  \pos{Transitive verb}  \sense{
    \definition{to slice open}    \example{
      \exsen{Bogo sɨmell käm dupotän.}      \extran{He sliced open the pig stomach.}}}}

\entry{pott}{\headword{pott}  \note{Everything in the house}  \pronnote{poʈ͡ʂ}
  \pos{Noun}  \sense{
    \definition{toilet?}    \note{Everything in the house}    \example{
}    \example{
}    \example{
}}}

\entry{Pottängäm}{\headword{Pottängäm}  \note{Places around Limol}  \pronnote{poʈ͡ʂəŋəm}
  \pos{Proper noun}  \sense{
    \definition{Pottängäm (camp and garden place)}    \note{Places around Limol}    \example{
}    \example{
}    \example{
}}}

\entry{poo}{\headword{poo}
  \variantof{SpellingVariant of \vartext{po}}  \pronnote{poo}}

\entry{praedde}{\headword{praedde}  \note{Part of the week}  \pronnote{praeɖ͡ʐe}
  \pos{Noun}  \sense{
    \definition{Friday}    \note{Part of the week}    \example{
}    \example{
}    \example{
}}
  \variants{\SpellingVariant \vartext{Praede, priede, praede}}}

\entry{Praede}{\headword{Praede}
  \variantof{SpellingVariant of \vartext{praedde}}}

\entry{praede}{\headword{praede}
  \variantof{SpellingVariant of \vartext{praedde}}  \pronnote{praede}}

\entry{praemari skul}{\headword{praemari skul}  \note{School}  \pronnote{praemari skul}
  \pos{Noun}  \sense{
    \definition{primary school}    \note{School}    \example{
}    \example{
}    \example{
}}
  \variants{\SpellingVariant \vartext{praemeri skul}}}

\entry{praemeri skul}{\headword{praemeri skul}
  \variantof{SpellingVariant of \vartext{praemari skul}}}

\entry{pre}{\headword{pre}
  \variantof{Fast Speech Variant of \vartext{päre}}}

\entry{prep}{\headword{prep}
  \pos{Noun}  \sense{
    \definition{preparatory school}}}

\entry{pri}{\headword{pri}
  \pos{Modifier}  \sense{
    \definition{free}}}

\entry{priede}{\headword{priede}
  \variantof{SpellingVariant of \vartext{praedde}}  \pronnote{priede}}

\entry{Priscilla}{\headword{Priscilla}  \pronnote{prisila}
  \pos{Proper noun}  \sense{
    \definition{female personal name}}
  \variants{\SpellingVariant \vartext{Prisila}}}

\entry{Prisila}{\headword{Prisila}
  \variantof{SpellingVariant of \vartext{Priscilla}}}

\entry{Priski}{\headword{Priski}  \pronnote{priski}
  \pos{Proper noun}  \sense{
    \definition{female personal name}}}

\entry{probens}{\headword{probens}
  \pos{Noun}  \sense{
    \definition{province}}}

\entry{pu}{\headword{pu}  \pronnote{pu}
  \pos{Noun}  \sense{\textbf{1.} 
    \definition{floating grass or island in swamp}    \example{
      \exsen{Kapalla pu me mälla da mondre anggan.}      \extran{The woman gardens on the floating kapalla grass.}}}  \sense{\textbf{2.} 
    \definition{swamp garden}    \example{
      \exsen{Oba pu mi biye de ddone itrel a anggan.}      \extran{The disease didn't strike the taro in their swamp garden.}}}}

\entry{pud}{\headword{pud}
  \pos{Noun}  \sense{
    \definition{end (of a long object)}    \example{
      \exsen{pop peyang pud}      \extran{the end with a hole}}}}

\entry{pudd}{\headword{pudd}  \pronnote{puɖ͡ʐ}
  \pos{Noun}  \sense{
    \definition{place with flattened grass where wallabies sleep}}}

\entry{pudd}{\headword{pudd}  \pronnote{puɖ͡ʐ}
  \pos{Noun}  \sense{
    \definition{taro shoot}}}

\entry{puder}{\headword{puder}  \lexeme{puder}  \pronnote{puder}
  \pos{Noun}  \sense{
    \definition{type of long grass}}}

\entry{pudrell}{\headword{pudrell}
  \variantof{SpellingVariant of \vartext{pädrall}}  \pronnote{pudreɽ}}

\entry{pug}{\headword{pug}  \pronnote{puɡ}
  \pos{Noun}  \sense{
    \definition{bundle of bark, two or three inside.}}}

\entry{pugupugu}{\headword{pugupugu}  \pronnote{puɡupuɡu}
  \pos{Noun}  \sense{
    \definition{collared imperial pigeon}    \scientificname{Ducula mullerii}    \example{
}    \example{
}    \example{
}}}

\entry{puikmetutu}{\headword{puikmetutu}
  \pos{Noun}  \sense{
    \definition{type of gecko}}}

\entry{Puinde}{\headword{Puinde}  \pronnote{puinde}
  \pos{Proper noun}  \sense{
    \definition{male personal name}}}

\entry{pukong}{\headword{pukong}  \note{Adjectives}  \pronnote{pukoŋ}
  \pos{Modifier}  \sense{
    \definition{thick}    \note{Adjectives}    \example{
      \exsen{Ge peba da pukong dan.}      \extran{This paper is thick.}}    \example{
}    \example{
}}}

\entry{puku}{\headword{puku}  \note{Traditional Arrows}  \pronnote{puku}
  \pos{Noun}  \sense{
    \definition{type of arrow}    \note{Traditional Arrows}    \example{
}    \example{
}    \example{
}}}

\entry{puku kubull}{\headword{puku kubull}  \note{Animal names}  \pronnote{puku kubuɽ}
  \pos{Noun}  \sense{
    \definition{type of bush wallaby with white tail}    \note{Animal names}    \example{
}    \example{
}    \example{
}}}

\entry{pullkom}{\headword{pullkom}  \note{Parts of a bird}  \pronnote{puɽkom}
  \pos{Noun}  \sense{
    \definition{tailfeather}    \note{Parts of a bird}    \example{
}    \example{
}    \example{
}}}

\entry{pumi}{\headword{pumi}  \pronnote{pumi}
  \pos{Modifier}  \sense{
    \definition{exhausting, tiring, strenuous}    \example{
      \exsen{Pumi melem de nangesallo.}      \extran{They did exhausting work.}}    \example{
}    \example{
}}}

\entry{pumkak}{\headword{pumkak}  \pronnote{'pumkak}
  \pos{Noun}  \sense{
    \definition{type of sago}}}

\entry{pungg}{\headword{pungg}  \pronnote{puŋɡ}
  \pos{Intransitive verb}  \sense{\textbf{1.} 
    \definition{to get sick, be in pain}    \example{
      \exsen{Da bam kungge da naddägän, bongo mɨnyi ttattlle ampug.}      \extran{If a spider bites you, you will get very sick.}}}  \sense{\textbf{2.} 
    \definition{to take the first chop with the axe}}}

\entry{pupulläkäm}{\headword{pupulläkäm}  \note{Mushrooms}  \pronnote{pupuɽəkəm}
  \pos{Noun}  \sense{
    \definition{type of mushroom}    \note{Mushrooms}    \example{
}    \example{
}    \example{
}}}

\entry{purta}{\headword{purta}
  \variantof{Unspecified Variant of \vartext{pärta}}  \pronnote{purta}}

\entry{pus}{\headword{pus}  \lexeme{pus}  \pronnote{pus}
  \pos{Noun}  \sense{
    \definition{cat}}}

\entry{put}{\headword{put}  \pronnote{put}
  \pos{Noun}  \sense{
    \definition{type of brown bark}    \example{
      \exsen{Oba ma da pall put alle papekatt dan.}      \extran{Their house is walled with pall bark walling.}}}}

\entry{putt}{\headword{putt}  \note{Numbers}  \pronnote{puʈ͡ʂ}
  \pos{Noun}  \sense{\textbf{1.} 
    \definition{hoof}    \note{Numbers}    \example{
      \exsen{Ddia da putt gottkewän.}      \extran{The deer folded its hooves.}}}  \sense{\textbf{2.} 
    \definition{six (yam counting numeral; 6^1)}    \note{Numbers}    \example{
}    \example{
}    \example{
}}}

\entry{putt kämamkämam}{\headword{putt kämamkämam}  \pronnote{puʈ͡ʂ kəmamkəmam}
  \pos{Noun}  \sense{
    \definition{type of sago bundle wrapped with sago leaves and tied at one end}}}

\entry{puyem}{\headword{puyem}  \note{Hunting}  \pronnote{pujem}
  \pos{Noun}  \sense{
    \definition{hunting event in which one person hits the ground with a short stick}    \note{Hunting}    \example{
}    \example{
}    \example{
}}}

\entry{puzi}{\headword{puzi}  \pronnote{puzi}
  \pos{Noun}  \sense{
    \definition{type of bird with crown}    \example{
}    \example{
}    \example{
}}}

\entry{pwapwa}{\headword{pwapwa}
  \variantof{SpellingVariant of \vartext{poapoa}}  \pronnote{pwapwa}}

\chapter{Q}


\entry{Queenie}{\headword{Queenie}  \pronnote{kwueenie}
  \pos{Proper noun}  \sense{
    \definition{female personal name}}}

\entry{Quin}{\headword{Quin}
  \pos{Proper noun}  \sense{
    \definition{male personal name}}}

\entry{Quinten}{\headword{Quinten}
  \pos{Proper noun}  \sense{
    \definition{male personal name}}}

\entry{Quinteth}{\headword{Quinteth}  \pronnote{kwinteθ}
  \pos{Proper noun}  \sense{
    \definition{female personal name}}}

\entry{Quinton}{\headword{Quinton}  \pronnote{kwuinton}
  \pos{Proper noun}  \sense{
    \definition{male personal name}}
  \variants{\SpellingVariant \vartext{Kwinton}}}

\chapter{R}


\entry{räba}{\headword{räba}  \pronnote{rəba}
  \pos{Noun}  \sense{
    \definition{rubber; eraser}}
  \variants{\SpellingVariant \vartext{raba}}}

\entry{raba}{\headword{raba}
  \variantof{SpellingVariant of \vartext{räba}}}

\entry{Räba blok}{\headword{Räba blok}  \pronnote{rəba blok}
  \pos{Proper noun}  \sense{
    \definition{Rubber block (toponym)}}}

\entry{Rabaol}{\headword{Rabaol}
  \variantof{SpellingVariant of \vartext{Rabaul}}}

\entry{Rabaul}{\headword{Rabaul}
  \pos{Proper noun}  \sense{
    \definition{Rabaul (toponym)}}
  \variants{\SpellingVariant \vartext{Rabaol}}}

\entry{rabis}{\headword{rabis}
  \pos{Noun}  \sense{
    \definition{rubbish, trash}}}

\entry{rädräd}{\headword{rädräd}  \pronnote{rədrəd}
  \pos{Noun}  \sense{
    \definition{type of grey fish that lives in the swamp}}}

\entry{ran}{\headword{ran}
  \variantof{freevarlab of \vartext{English loanword}}}

\entry{Ranky}{\headword{Ranky}  \pronnote{raŋki}
  \pos{Proper noun}  \sense{
    \definition{male personal name}}}

\entry{rarae}{\headword{rarae}  \pronnote{rarae}
  \pos{Noun}  \sense{\textbf{1.} 
    \definition{hunting technique in which hunters line up to cover ground in the absence of dogs}}  \sense{\textbf{2.} 
    \definition{fishing technique in which ladies line up with nets}}}

\entry{Raroge}{\headword{Raroge}  \note{Place names}  \pronnote{raroɡe}
  \pos{Proper noun}  \sense{
    \definition{Raroge (toponym)}    \note{Place names}    \example{
}    \example{
}    \example{
}}}

\entry{Rasol}{\headword{Rasol}  \pronnote{rasol}
  \pos{Proper noun}  \sense{
    \definition{male personal name}}}

\entry{Rasti}{\headword{Rasti}  \pronnote{rasti}
  \pos{Proper noun}  \sense{\textbf{1.} 
    \definition{name of a dog in Limol, Rusty}}  \sense{\textbf{2.} 
    \definition{dog name}}}

\entry{rata}{\headword{rata}  \lexeme{rata}  \pronnote{rata}
  \pos{Noun}  \sense{
    \definition{type of flowering plant with red bracts and nectar that attracts insects}}}

\entry{Raynold}{\headword{Raynold}  \pronnote{renald}
  \pos{Proper noun}  \sense{
    \definition{male personal name}}}

\entry{RED~}{\headword{RED~}
  \pos{Adjective}  \sense{
    \definition{Reduplication of an infinitive (including its pluractional suffix) can be used to create an adjective describing someone's personality, i.e. what they characteristically do}}
  \variants{\Derivational reduplicant \vartext{kanye}}
  \variants{\freevarlab \vartext{laem~}}}

\entry{RED~}{\headword{RED~}
  \pos{Nominal}  \sense{
    \definition{attaches to a nominal to derive a new nominal with a smaller referent}}
  \variants{\Derived Variant \vartext{pollon~, auma~, ngata~, kup~, ma~}}
  \variants{\freevarlab \vartext{pättä~}}
  \variants{\Derivational reduplicant \vartext{katre~, polle~, ttoen~}}}

\entry{RED~}{\headword{RED~}
  \pos{}  \sense{
}
  \variants{\Dialectal Variant \vartext{togo~}}}

\entry{~RED}{\headword{~RED}
  \pos{Verb}  \sense{
}
  \variants{\Inflected Form \vartext{~ɨt}}}

\entry{RED~}{\headword{RED~}
  \pos{}  \sense{
    \definition{Certain adjectives and kinship terms reduplicate when the referent is nonsingular}}
  \variants{\Plural reduplicant \vartext{ulle~, kle~, män~, to~, mälla~, mang~, käle~, nag~, ko~}}
  \variants{\freevarlabs \vartext{mänya~, tubu~, masa~}}
  \variants{\Unspecified Variant \vartext{mosen~}}
  \variants{\Dialectal Variant \vartext{mullae~}}}

\entry{RED~}{\headword{RED~}
  \pos{Adverb}  \sense{
}
  \variants{\freevarlabs \vartext{ngälla~, era~, gäbä~, kakän~, ngädnan~, daram~, ddobae~, yuwet~, ttänganen~, ttonen~, ngäna~}}
  \variants{\Derived Variant \vartext{dune~, dowa~, mällka~, käla~, du~}}
  \variants{\Derivational reduplicant \vartext{imne~}}
  \variants{\Inflected Form \vartext{ngäse~}}}

\entry{RED~}{\headword{RED~}
  \pos{Verb}  \sense{
}
  \variants{\Inflected Form \vartext{pa~, mall~, sebor~, bäny~, ngäs~, mänd~, gu~, ttima~, dda~}}
  \variants{\Infinitival reduplicant \vartext{bälla~, za~, ngas~, ttä~, ngä~, kanye, yag~, ddä~, ta~, kapi~, po~, da~, eka~, ddo~, go~, gädd~, nga~, na~, kla~, ma~, poll~, mändd~, ko~}}
  \variants{\freevarlabs \vartext{ngäll~, mit~, lläk~, be~, ngäd~, ib~, badd~, nyämae~}}
  \variants{\Dialectal Variant \vartext{ibi, sä~}}
  \variants{\Derived Variant \vartext{ma~}}}

\entry{redi}{\headword{redi}
  \pos{Transitive coverb}  \sense{
    \definition{to prepare, make ready}}}

\entry{Redley}{\headword{Redley}  \pronnote{rɛdli}
  \pos{Proper noun}  \sense{
    \definition{male personal name}}}

\entry{Regina}{\headword{Regina}  \pronnote{reɡina}
  \pos{Proper noun}  \sense{
    \definition{female personal name}}}

\entry{rekoda}{\headword{rekoda}
  \pos{Noun}  \sense{
    \definition{recorder}}}

\entry{Reks}{\headword{Reks}
  \pos{Proper noun}  \sense{
    \definition{male personal name}}
  \variants{\SpellingVariant \vartext{Rex}}}

\entry{Rena}{\headword{Rena}
  \pos{Proper noun}  \sense{
    \definition{female personal name}}}

\entry{renz}{\headword{renz}  \pronnote{renz}
  \pos{Intransitive coverb}  \sense{
    \definition{to walk somewhere without carrying anything}}}

\entry{resourceful}{\headword{resourceful}
  \variantof{Dialectal Variant of \vartext{English loanword}}}

\entry{retam}{\headword{retam}  \note{Food}  \pronnote{retam}
  \pos{Noun}  \sense{
    \definition{type of big yam}    \note{Food}    \example{
}    \example{
}    \example{
}}}

\entry{Rex}{\headword{Rex}
  \variantof{SpellingVariant of \vartext{Reks}}  \pronnote{reks}}

\entry{Reend}{\headword{Reend}  \pronnote{reend}
  \pos{Proper noun}  \sense{
    \definition{male personal name}}}

\entry{Rhoda}{\headword{Rhoda}  \pronnote{rhoda}
  \pos{Proper noun}  \sense{
    \definition{female personal name}}
  \variants{\Unspecified Variant \vartext{Roda, Rhodda}}}

\entry{Rhodda}{\headword{Rhodda}
  \variantof{Unspecified Variant of \vartext{Rhoda}}}

\entry{Richard}{\headword{Richard}  \pronnote{ritʃard}
  \pos{Proper noun}  \sense{
    \definition{male personal name}}}

\entry{Rind}{\headword{Rind}  \pronnote{rind}
  \pos{Proper noun}  \sense{
    \definition{male personal name}}}

\entry{rir}{\headword{rir}  \pronnote{rir}
  \pos{Noun}  \sense{
    \definition{type of medium-sized bamboo}}}

\entry{rizen}{\headword{rizen}
  \variantof{Dialectal Variant of \vartext{English loanword}}}

\entry{rle}{\headword{rle}
  \pos{Verb}  \sense{
    \definition{to reproduce}}}

\entry{Roaele}{\headword{Roaele}  \pronnote{roaele}
  \pos{Proper noun}  \sense{
    \definition{male personal name}}}

\entry{Roak}{\headword{Roak}  \pronnote{roak}
  \pos{Proper noun}  \sense{
    \definition{male personal name}}}

\entry{Robae}{\headword{Robae}
  \pos{Proper noun}  \sense{
    \definition{female personal name}}
  \variants{\SpellingVariant \vartext{Robai}}}

\entry{Robai}{\headword{Robai}
  \variantof{SpellingVariant of \vartext{Robae}}  \pronnote{robai}}

\entry{Roda}{\headword{Roda}
  \variantof{Unspecified Variant of \vartext{Rhoda}}  \pronnote{roda}}

\entry{rokaroka}{\headword{rokaroka}  \note{Oral Traditions}  \pronnote{rokaroka}
  \pos{Noun}  \sense{\textbf{1.} 
    \definition{riddle}    \note{Oral Traditions}    \example{
}    \example{
}    \example{
}}  \sense{\textbf{2.} 
    \definition{fable}    \note{Oral Traditions}}}

\entry{rol}{\headword{rol}  \note{Poison}  \pronnote{rol}
  \pos{Noun}  \sense{
    \definition{type of hairy caterpillar}    \note{Poison}    \example{
}    \example{
}    \example{
}}}

\entry{rolkutt}{\headword{rolkutt}  \lexeme{rolkutt}  \pronnote{rolkuʈ͡ʂ}
  \pos{Noun}  \sense{
    \definition{crawling grass}}}

\entry{Rom}{\headword{Rom}  \pronnote{rom}
  \pos{Proper noun}  \sense{
    \definition{Rome}}}

\entry{rop}{\headword{rop}
  \pos{Noun}  \sense{
    \definition{rope}}}

\entry{rorwae}{\headword{rorwae}  \pronnote{rorwae}
  \pos{Verb}  \sense{
    \definition{to stagger (e.g. when drunk)}    \example{
      \exsen{Rorwae a mudan, llokttangae nägabäll.}      \extran{Don't stagger; stand straight.}}}}

\entry{rorwaerorwae}{\headword{rorwaerorwae}  \pronnote{rorwaerorwae}
  \pos{Adverb}  \sense{
    \definition{off balance, falteringly}    \example{
      \exsen{Maneya rorwaerorwae ibiag daeya.}      \extran{Maneya walked falteringly.}}}}

\entry{Ros}{\headword{Ros}
  \variantof{SpellingVariant of \vartext{Rose}}}

\entry{Rose}{\headword{Rose}  \pronnote{rose}
  \pos{Proper noun}  \sense{
    \definition{female personal name}}
  \variants{\SpellingVariant \vartext{Ros}}}

\entry{Rosela}{\headword{Rosela}  \pronnote{rosela}
  \pos{Proper noun}  \sense{
    \definition{female personal name}}}

\entry{Rowak}{\headword{Rowak}  \pronnote{rowak}
  \pos{Proper noun}  \sense{
    \definition{male personal name}}}

\entry{Rual}{\headword{Rual}  \pronnote{rual}
  \pos{Proper noun}  \sense{
    \definition{Rual (toponym)}    \example{
}    \example{
}    \example{
}}}

\entry{rubi}{\headword{rubi}  \note{Food}  \pronnote{rubi}
  \pos{Noun}  \sense{\textbf{1.} 
    \definition{type of long yam with a white interior, white skin, and no thorns}    \note{Food}    \example{
}    \example{
}    \example{
}}  \sense{\textbf{2.} 
    \definition{type of arrow}    \note{Food}}}

\entry{rullgoe}{\headword{rullgoe}
  \pos{Transitive verb}  \sense{
    \definition{to drag}    \example{
      \exsen{Tim ngänäm dangminggän ngämi däbe käza de tutu wi darullgoeya.}      \extran{Tim helped me drag that crocodile onto land.}}}}

\entry{rupi}{\headword{rupi}
  \pos{Noun}  \sense{
    \definition{type of leaf}}}

\entry{ruriruri}{\headword{ruriruri}  \pronnote{ruriruri}
  \pos{Noun}  \sense{
    \definition{earthquake}}}

\chapter{S}


\entry{1940s}{\headword{1940s}
  \variantof{freevarlab of \vartext{English loanword}}}

\entry{sa}{\headword{sa}
  \pos{Intransitive verb}  \sense{
    \definition{declare}}}

\entry{sä~}{\headword{sä~}
  \variantof{Dialectal Variant of \vartext{RED~}}}

\entry{sabana}{\headword{sabana}
  \pos{Noun}  \sense{
    \definition{savanna}}}

\entry{sabi}{\headword{sabi}  \pronnote{sabi}
  \pos{Noun}  \sense{\textbf{1.} 
    \definition{law, rule}    \example{
      \exsen{Ngämi bam da batrameya, abo bongo ge sabi de nongkollmäll.}      \extran{We will take you, but you must follow these rules.}}}  \sense{\textbf{2.} 
    \definition{taboo}    \example{
      \exsen{Ngämo mänang a ngänäm de aläm bagneyo ada, ngämo bin ttam a mudan, sabi.}      \extran{My in-laws will be warning me that it's taboo to call me by my name.}}}}

\entry{sabol}{\headword{sabol}  \note{Tools}  \pronnote{sabol}
  \pos{Noun}  \sense{
    \definition{shovel, spade}    \note{Tools}    \example{
}    \example{
}    \example{
}}}

\entry{Sadua}{\headword{Sadua}  \pronnote{sadua}
  \pos{Proper noun}  \sense{
    \definition{male personal name}}
  \variants{\Unspecified Variant \vartext{Saduwa}}}

\entry{Saduwa}{\headword{Saduwa}
  \variantof{Unspecified Variant of \vartext{Sadua}}  \pronnote{saduwa}}

\entry{sae}{\headword{sae}  \pronnote{sae}
  \pos{Transitive verb}  \sense{\textbf{1.} 
    \definition{to close, cover}    \example{
      \exsen{Nga bongo ikop nasǃ}      \extran{You, close your eyes!}}    \example{
      \exsen{Däräng a ikop dasän.}      \extran{The dog closed his eyes.}}}  \sense{\textbf{2.} 
    \definition{to extinguish, put out}    \example{
      \exsen{Yu da saemeny dan.}      \extran{The fire isn't put out.}}}}

\entry{saem}{\headword{saem}  \note{Stages of life}  \pronnote{saem}
  \pos{Verb}  \sense{
    \definition{to babble}    \note{Stages of life}    \example{
      \exsen{Bogo saem allan.}      \extran{He's babbling.}}    \example{
}    \example{
}}}

\entry{Saemon}{\headword{Saemon}
  \pos{Proper noun}  \sense{
    \definition{male personal name}}
  \variants{\SpellingVariant \vartext{Simon}}}

\entry{saen}{\headword{saen}
  \pos{Noun}  \sense{
    \definition{sign}    \example{
      \exsen{Mälla da saen dangesän.}      \extran{The woman made a sign.}}}}

\entry{saeten}{\headword{saeten}
  \pos{Proper noun}  \sense{
    \definition{Satan}}}

\entry{saeyang}{\headword{saeyang}  \pronnote{saejaŋ}
  \pos{Modifier}  \sense{
    \definition{closed}    \example{
      \exsen{Däräng täräpang dagirnän ikop saeyang.}      \extran{The dog stayed there a long time, eyes closed.}}}}

\entry{sägädag}{\headword{sägädag}
  \pos{Modifier}  \sense{
}}

\entry{sägädsägäd}{\headword{sägädsägäd}
  \variantof{SpellingVariant of \vartext{sägäsägäd}}}

\entry{sägäsägäd}{\headword{sägäsägäd}  \note{Colors}  \pronnote{səɡəsəɡəd}
  \pos{Color term}  \sense{
    \definition{yellow}    \note{Colors}    \example{
}    \example{
}}
  \variants{\SpellingVariant \vartext{sägädsägäd}}}

\entry{sägäsägäd manika}{\headword{sägäsägäd manika}  \pronnote{səɡəsəɡəd manika}
  \pos{Noun}  \sense{
    \definition{yellow cassava}}}

\entry{sagol}{\headword{sagol}  \note{Dancing}  \pronnote{saɡol}
  \pos{Noun}  \sense{
    \definition{men's dance}    \note{Dancing}    \example{
}    \example{
}    \example{
}}}

\entry{sagol nyanyu}{\headword{sagol nyanyu}  \note{Dancing}  \pronnote{saɡol ɲaɲu}
  \pos{Intransitive compound verb}  \sense{
    \definition{to learn the steps}    \note{Dancing}    \example{
}    \example{
}    \example{
}}}

\entry{Sägrep}{\headword{Sägrep}
  \pos{Proper noun}  \sense{
    \definition{male personal name}}}

\entry{sagwan}{\headword{sagwan}  \note{Tools}  \pronnote{saɡwan}
  \pos{Noun}  \sense{
    \definition{woodworking tool}    \note{Tools}    \example{
}    \example{
}    \example{
}}}

\entry{Saisiato}{\headword{Saisiato}  \pronnote{saisiato}
  \pos{Proper noun}  \sense{
    \definition{female personal name}}}

\entry{säkar}{\headword{säkar}  \note{Types of palms}  \pronnote{səkar}
  \pos{Noun}  \sense{
    \definition{type of big palm tree that grows near big rivers with white flowers and hanging red fruit that cassowaries eat}    \note{Types of palms}    \example{
}    \example{
}    \example{
}}}

\entry{sakar}{\headword{sakar}  \note{Type of pandanus}  \pronnote{sakar}
  \pos{Noun}  \sense{
    \definition{type of edible pandanus}    \note{Type of pandanus}    \example{
}    \example{
}    \example{
}}}

\entry{Saken}{\headword{Saken}
  \pos{Proper noun}  \sense{
    \definition{personal name}}}

\entry{säkmällsäkmäll}{\headword{säkmällsäkmäll}  \pronnote{səkməɽsəkməɽ}
  \pos{Adverb}  \sense{
    \definition{limping}    \example{
      \exsen{Bogo säkmällsäkmäll ibi allan.}      \extran{He is walking with a limp.}}}}

\entry{Sakoyaratt}{\headword{Sakoyaratt}  \note{Places around Limol}  \pronnote{sakojaraʈ͡ʂ}
  \pos{Proper noun}  \sense{
    \definition{Sakoyaratt (toponym)}    \note{Places around Limol}    \example{
}    \example{
}    \example{
}}}

\entry{säkpa mit}{\headword{säkpa mit}  \note{Conflict/war}  \pronnote{səkpa mit}
  \pos{Noun}  \sense{
    \definition{theft}    \note{Conflict/war}    \example{
}    \example{
}    \example{
}}}

\entry{saks}{\headword{saks}  \note{Clothing}  \pronnote{saks}
  \pos{Noun}  \sense{
    \definition{socks}    \note{Clothing}    \example{
}    \example{
}    \example{
}}}

\entry{Sali}{\headword{Sali}  \pronnote{sali}
  \pos{Proper noun}  \sense{
    \definition{male personal name}}
  \variants{\SpellingVariant \vartext{Saly}}}

\entry{Salome}{\headword{Salome}
  \pos{Proper noun}  \sense{
    \definition{female personal name}}}

\entry{Saly}{\headword{Saly}
  \variantof{SpellingVariant of \vartext{Sali}}  \pronnote{sali}}

\entry{Sam}{\headword{Sam}  \pronnote{sam}
  \pos{Proper noun}  \sense{
    \definition{male personal name}}}

\entry{Samae}{\headword{Samae}
  \pos{Proper noun}  \sense{
    \definition{male personal name}}
  \variants{\SpellingVariant \vartext{Shamae}}}

\entry{Samari}{\headword{Samari}  \pronnote{samari}
  \pos{Proper noun}  \sense{
    \definition{Samari (toponym)}    \example{
}    \example{
}    \example{
}}}

\entry{sämasäme}{\headword{sämasäme}
  \pos{Transitive verb}  \sense{
}}

\entry{Samat}{\headword{Samat}  \pronnote{samat}
  \pos{Proper noun}  \sense{
    \definition{female personal name}}}

\entry{sambuag}{\headword{sambuag}  \lexeme{sambuag}  \pronnote{sambuaɡ}
  \pos{Noun}  \sense{
    \definition{lobster; prawn}}}

\entry{sämell}{\headword{sämell}  \lexeme{sämell}
  \variantof{Unspecified Variant of \vartext{sɨmell}}  \pronnote{səmeɽ}}

\entry{samoa}{\headword{samoa}  \note{Names of bananas}  \pronnote{samoa}
  \pos{Noun}  \sense{
    \definition{type of introduced banana}    \note{Names of bananas}    \example{
}    \example{
}    \example{
}}}

\entry{sämongg}{\headword{sämongg}  \note{Verbs}  \pronnote{səmoŋɡ}
  \pos{Intransitive verb}  \sense{
    \definition{to feel}    \note{Verbs}    \example{
      \exsen{Ngäna mermerang ansomägan.}      \extran{I feel fine.}}    \example{
      \exsen{Obo pätt me gonsomonggän ada obo itrell de dektta llɨtt dägagän.}      \extran{He felt in his body that the doctor had stopped his illness.}}    \example{
}    \example{
}}}

\entry{Samson}{\headword{Samson}  \pronnote{samson}
  \pos{Proper noun}  \sense{
    \definition{male personal name}}}

\entry{Samuel}{\headword{Samuel}  \pronnote{samuel}
  \pos{Proper noun}  \sense{
    \definition{male personal name}}}

\entry{sana}{\headword{sana}  \lexeme{sana}  \pronnote{sana}
  \pos{Noun}  \sense{
    \definition{sago}    \example{
      \exsen{Sana tɨt iran.}      \extran{She is pounding the sago.}}    \example{
      \exsen{Sana ttäkoe eran.}      \extran{She is chopping the sago tree.}}}}

\entry{sana budar}{\headword{sana budar}  \note{Grubs}  \pronnote{sana budar}
  \pos{Noun}  \sense{
    \definition{sago grub (larva of black palm weevil)}    \note{Grubs}    \example{
}    \example{
}    \example{
}}}

\entry{sana marmar}{\headword{sana marmar}  \note{Grubs}  \pronnote{sana marmar}
  \pos{Noun}  \sense{
    \definition{sago grub (larva of black palm weevil)}    \note{Grubs}    \example{
}    \example{
}    \example{
}}}

\entry{sana tätäkang}{\headword{sana tätäkang}  \note{Traditional Arrows}  \pronnote{sana tətəkaŋ}
  \pos{Noun}  \sense{
    \definition{type of spear}    \note{Traditional Arrows}    \example{
}    \example{
}    \example{
}}}

\entry{sana wuttang}{\headword{sana wuttang}  \pronnote{sana wuʈ͡ʂaŋ}
  \pos{Noun}  \sense{
    \definition{the second stage of sago growth during which the plant is ~3 m tall}    \example{
}}}

\entry{sana wutwut}{\headword{sana wutwut}
  \pos{Noun}  \sense{
    \definition{sago dish cooked with pieces of banana leaf between the layers}}}

\entry{sanalläkäm}{\headword{sanalläkäm}  \note{Mushrooms}  \pronnote{sanaɽəkəm}
  \pos{Noun}  \sense{
    \definition{type of mushroom}    \note{Mushrooms}    \example{
}    \example{
}    \example{
}}}

\entry{sanasana}{\headword{sanasana}  \note{Sago palms}  \pronnote{sanasana}
  \pos{Noun}  \sense{
    \definition{type of edible sago}    \note{Sago palms}    \example{
}    \example{
}    \example{
}}}

\entry{sänasäne}{\headword{sänasäne}
  \pos{Transitive verb}  \sense{
    \definition{to take out}    \example{
      \exsen{Ngämi ddäddäg de mɨnyi bäsnaemeyo säspen att.}      \extran{We (excl.) will take out the meat from the pot.}}}}

\entry{sanateya}{\headword{sanateya}
  \pos{Noun}  \sense{
    \definition{sago pith}}}

\entry{sanawang}{\headword{sanawang}  \pronnote{sanawaŋ}
  \pos{Noun}  \sense{\textbf{1.} 
    \definition{area where sago grows}    \example{
      \exsen{Sanawang de yu da dättämän Kibobma me.}      \extran{The fire burned the sago growing area in Kibobma.}}}  \sense{\textbf{2.} 
    \definition{parable}}}

\entry{sänd}{\headword{sänd}
  \pos{Noun}  \sense{
    \definition{type of big yam with a white interior and few hairs}}}

\entry{sandana}{\headword{sandana}  \pronnote{sandana}
  \pos{Noun}  \sense{
    \definition{type of black palm; in this middle stage, the pith is chewed like sugarcane during the dry season for its moisture}}}

\entry{sände}{\headword{sände}
  \variantof{SpellingVariant of \vartext{sande}}}

\entry{Sände}{\headword{Sände}
  \variantof{SpellingVariant of \vartext{sande}}
  \variants{\SpellingVariant \vartext{sände}}}

\entry{sande}{\headword{sande}  \note{Part of the week}  \pronnote{sande}
  \pos{Noun}  \sense{\textbf{1.} 
    \definition{Sunday}    \note{Part of the week}    \example{
}    \example{
}    \example{
}}  \sense{\textbf{2.} 
    \definition{week}    \note{Part of the week}    \example{
      \exsen{Lla da ttongo sande makäp me mɨnyi täre we gonsärbemamalle.}      \extran{People will be preparing for the feast for a week.}}}
  \variants{\SpellingVariant \vartext{sände, Sände}}}

\entry{Sandra}{\headword{Sandra}  \pronnote{sandra}
  \pos{Proper noun}  \sense{
    \definition{female personal name}}}

\entry{säne}{\headword{säne}
  \pos{Noun}  \sense{
    \definition{type of yam}}}

\entry{Sanford}{\headword{Sanford}  \pronnote{sanford}
  \pos{Proper noun}  \sense{
    \definition{male personal name}}}

\entry{sanga}{\headword{sanga}  \pronnote{saŋa}
  \pos{Noun}  \sense{
    \definition{black-necked stork}    \scientificname{Ephippiorhynchus asiaticus australis}    \example{
}    \example{
}    \example{
}}}

\entry{sangapawi}{\headword{sangapawi}  \note{Food}  \pronnote{saŋapawi}
  \pos{Noun}  \sense{
    \definition{type of big, round yam with a white interior, white or red skin, hairs, and no thorns}    \note{Food}    \example{
}    \example{
}    \example{
}}}

\entry{sänge}{\headword{sänge}
  \pos{Transitive coverb}  \sense{
    \definition{to betray}    \example{
      \exsen{Bina makäp me ttongo aeya mɨnyi ngänäm sänge bagän, ttongo aeya duwem allan ngämo peyang.}      \extran{One among you will betray me: the one who eats with me.}}}}

\entry{saomasaoma}{\headword{saomasaoma}  \note{Clothing}  \pronnote{saomasaoma}
  \pos{Noun}  \sense{\textbf{1.} 
    \definition{armpit}    \note{Clothing}    \example{
}    \example{
}    \example{
}}  \sense{\textbf{2.} 
    \definition{dancing band}    \note{Clothing}    \example{
      \exsen{De lla da gagäll ikopang agan saomasaoma mättmättatt a.}      \extran{That man looks good (lit. bad) in his armbands.}}}}

\entry{säpalek}{\headword{säpalek}  \note{Everything in the house}  \pronnote{səpalek}
  \pos{Noun}  \sense{
    \definition{type of bag}    \note{Everything in the house}    \example{
      \exsen{Bibiae mätta de kemibi era säpaleng alle de ikoman.}      \extran{Bibiae brought many yams with the spalek bag.}}    \example{
}    \example{
}}
  \variants{\Fast Speech Variant \vartext{spalek}}
  \variants{\Dialectal Variant \vartext{säpaleng}}}

\entry{säpaleng}{\headword{säpaleng}
  \variantof{Dialectal Variant of \vartext{säpalek}}  \pronnote{səpaleŋ}}

\entry{sapang}{\headword{sapang}  \pronnote{sapaŋ}
  \pos{Modifier}  \sense{
    \definition{separate, apart, different; own, personal}    \example{
      \exsen{sapang ngättma we}      \extran{to a different place}}    \example{
      \exsen{Ddob oba masamasar a eragaeya, ddob oba sapang eka da dadegaeya.}      \extran{Some of their ancestors, some of them had separate languages.}}    \example{
      \exsen{Bogo kollba de sapang dowansegän.}      \extran{He put the fish aside.}}    \example{
      \exsen{Ubi sapang dag a ngämi ade sapang dag.}      \extran{They are on their own and we (excl.) are also on our own.}}    \example{
      \exsen{Puinde ade obo sapang nyängang dägagän.}      \extran{Puinde also had his personal bag.}}    \example{
}}}

\entry{sapangsapang}{\headword{sapangsapang}
  \variantof{Unspecified Variant of \vartext{sapasapang}}}

\entry{sapasapang}{\headword{sapasapang}  \pronnote{sapasapaŋ}
  \pos{Modifier}  \sense{
    \definition{various, many different}    \example{
      \exsen{sapasapang ttängäm att lla}      \extran{people from many different places}}    \example{
      \exsen{Ddob melem a sapasapang dadegän.}      \extran{There are a variety of other jobs.}}    \example{
      \exsen{Iba eka da gullem sapasapang pallall e dan.}      \extran{Our (incl.) story is about various kinds of snakes.}}}
  \variants{\Fast Speech Variant \vartext{sasapang}}
  \variants{\Unspecified Variant \vartext{sapangsapang}}}

\entry{sapebllabllot}{\headword{sapebllabllot}  \note{Types of Taro}  \pronnote{sapebɽabɽot}
  \pos{Noun}  \sense{
    \definition{type of taro}    \note{Types of Taro}    \example{
}    \example{
}    \example{
}}}

\entry{saper ine}{\headword{saper ine}  \note{Water}  \pronnote{saper ine}
  \pos{Noun}  \sense{
    \definition{clean water}    \note{Water}    \example{
}    \example{
}    \example{
}}}

\entry{sapiri}{\headword{sapiri}  \lexeme{sapiri}  \pronnote{sapiri}
  \pos{Noun}  \sense{
    \definition{type of flowering plant}}}

\entry{Sapusa}{\headword{Sapusa}  \pronnote{sapusa}
  \pos{Proper noun}  \sense{
    \definition{female personal name}}}

\entry{Sara}{\headword{Sara}  \pronnote{sara}
  \pos{Proper noun}  \sense{
    \definition{female personal name}}}

\entry{sära}{\headword{sära}  \note{Hunting}  \pronnote{səra}
  \pos{Noun}  \sense{\textbf{1.} 
    \definition{tail}    \note{Hunting}    \example{
      \exsen{Käza da kollba bo sära de däddägän.}      \extran{The crocodile ate the fish tail.}}    \example{
}    \example{
}}  \sense{\textbf{2.} 
    \definition{end of a bow}    \note{Hunting}}  \sense{\textbf{3.} 
    \definition{uncut grass skirt}    \note{Hunting}}}

\entry{sära mit}{\headword{sära mit}  \note{Parts of a reptile}  \pronnote{səra mit}
  \pos{Noun}  \sense{
    \definition{sacrum (bone)}    \note{Parts of a reptile}    \example{
}    \example{
}    \example{
}}}

\entry{sära pakall}{\headword{sära pakall}  \note{Parts of a reptile}  \pronnote{səra pakaɽ}
  \pos{Noun}  \sense{
    \definition{caudal fin}    \note{Parts of a reptile}    \example{
}    \example{
}    \example{
}}}

\entry{sära pättkäp}{\headword{sära pättkäp}  \note{Parts of an insect}  \pronnote{səra pəʈ͡ʂkəp}
  \pos{Noun}  \sense{
    \definition{mesothorax}    \note{Parts of an insect}    \example{
}    \example{
}    \example{
}}}

\entry{sära pipi}{\headword{sära pipi}
  \variantof{SpellingVariant of \vartext{särapipi}}}

\entry{sära pot}{\headword{sära pot}  \note{Parts of a reptile}  \pronnote{səra pot}
  \pos{Noun}  \sense{\textbf{1.} 
    \definition{tip of tail}    \note{Parts of a reptile}    \example{
}    \example{
}    \example{
}}  \sense{\textbf{2.} 
    \definition{lastborn, youngest}    \note{Parts of a reptile}}}

\entry{Sarah}{\headword{Sarah}
  \variantof{SpellingVariant of \vartext{Sera}}  \pronnote{sarah}}

\entry{särämbae}{\headword{särämbae}  \note{Everyday words and phrases}  \pronnote{sərəmbae}
  \pos{Transitive verb}  \sense{\textbf{1.} 
    \definition{to prepare, be prepared, be ready, get ready}    \note{Everyday words and phrases}    \example{
      \exsen{Tämamae ansärbemom!}      \extran{Everyone, get ready!}}    \example{
      \exsen{Ngäna ada ngonongg allan ada malla bongo särämbaeatt dan.}      \extran{I think that you are not prepared.}}}  \sense{\textbf{2.} 
    \definition{to prepare, arrange, get ready}    \note{Everyday words and phrases}    \example{
      \exsen{Ubi oba gall me dagwaeya, gull de dansärämbenegneyo.}      \extran{The two of them were in their canoe, preparing their nets.}}    \example{
      \exsen{Ngäna obo pätt de nansärämbeyan au we.}      \extran{I prepared his body for burial.}}}  \sense{\textbf{3.} 
    \definition{to fix, solve, resolve}    \note{Everyday words and phrases}    \example{
      \exsen{Ngämlle mokowang melem dan lla bo buddo särämbae.}      \extran{I enjoy the task of solving people's problems.}}}}

\entry{särapipi}{\headword{särapipi}
  \pos{Noun}  \sense{
    \definition{lastborn, youngest}}
  \variants{\SpellingVariant \vartext{sära pipi}}}

\entry{Sarbi}{\headword{Sarbi}  \pronnote{sarbi}
  \pos{Proper noun}  \sense{
    \definition{female personal name}}}

\entry{sare}{\headword{sare}  \note{Types of Taro}  \pronnote{sare}
  \pos{Noun}  \sense{
    \definition{type of big taro}    \note{Types of Taro}    \example{
}    \example{
}    \example{
}}}

\entry{säre}{\headword{säre}  \pronnote{səre}
  \pos{Adverb}  \sense{
    \definition{sadly}    \example{
      \exsen{Ause da säre kullkull de daempononggän tämamae.}      \extran{The old woman sadly burned all the grass.}}}}

\entry{särem}{\headword{särem}  \pronnote{sərem}
  \pos{Noun}  \sense{
    \definition{darkness, dark}    \example{
      \exsen{ag särem me}      \extran{at dawn (lit. in the dark of morning)}}    \example{
      \exsen{Särem a dangttanän ttängäm de.}      \extran{Darkness came over the village.}}}
  \variants{\SpellingVariant \vartext{sɨrem, sirem}}}

\entry{säremang}{\headword{säremang}  \pronnote{səremaŋ}
  \pos{Modifier}  \sense{
    \definition{dark, dim}    \example{
      \exsen{Iddob me ttängäm a säremang abal gogän.}      \extran{At night, the village became very dark.}}    \example{
      \exsen{Llamäg o mällause täräp me iba ikop a mɨnyi säremang bognegän.}      \extran{In old age, our (incl.) eyes will go dim.}}}}

\entry{säremang ma}{\headword{säremang ma}
  \pos{Noun}  \sense{
    \definition{prison, jail}}}

\entry{säremsärem}{\headword{säremsärem}
  \pos{Adverb}  \sense{
    \definition{in darkness, in the dark}    \example{
      \exsen{Ngämi säremsärem giddollnenang dagaeya.}      \extran{We (excl.) were living in darkness.}}}}

\entry{säresäremang}{\headword{säresäremang}  \pronnote{səresəremaŋ}
  \pos{Modifier}  \sense{
    \definition{cloudy, dim}}}

\entry{sari}{\headword{sari}
  \pos{Interjection}  \sense{
    \definition{sorry}}
  \variants{\SpellingVariant \vartext{sori, sorry}}}

\entry{sariem}{\headword{sariem}  \pronnote{sariem}
  \pos{Transitive coverb}  \sense{
    \definition{to neglect}}}

\entry{säroe}{\headword{säroe}
  \variantof{SpellingVariant of \vartext{soroe}}  \pronnote{səroe}}

\entry{säs}{\headword{säs}  \pronnote{səs}
  \pos{Noun}  \sense{
    \definition{type of sago wrapped in young leaves}}}

\entry{Sasa}{\headword{Sasa}
  \pos{Proper noun}  \sense{
    \definition{female personal name}}}

\entry{sasapang}{\headword{sasapang}
  \variantof{Fast Speech Variant of \vartext{sapasapang}}}

\entry{säsäri}{\headword{säsäri}
  \variantof{SpellingVariant of \vartext{sisri}}}

\entry{säsäs}{\headword{säsäs}
  \pos{Intransitive verb}  \sense{
    \definition{to rub}    \example{
      \exsen{Gosäsneyo obaoba.}      \extran{They were rubbing each other.}}}}

\entry{Sasi}{\headword{Sasi}  \pronnote{sasi}
  \pos{Proper noun}  \sense{
    \definition{female personal name}}}

\entry{sasor}{\headword{sasor}  \note{Traditional fishing ways and nets}  \pronnote{sasor}
  \pos{Noun}  \sense{
    \definition{trap for fish. Kawam word.}    \note{Traditional fishing ways and nets}    \example{
}    \example{
}    \example{
}}}

\entry{säspen}{\headword{säspen}  \pronnote{səspen}
  \pos{Noun}  \sense{\textbf{1.} 
    \definition{pot, saucepan}    \example{
      \exsen{Ngämi ddäddäg de mɨnyi bäsnaemeyo säspen att.}      \extran{We (excl.) will take out the meat from the pot.}}}  \sense{\textbf{2.} 
    \definition{to boil}    \example{
      \exsen{Ede ge ddia de däbänya wa dikomya ma we, ngäma sespen dägayaebeya wa däddägaebeya.}      \extran{So we (excl.) cut that deer, and brought it home, boiled it, and ate it.}}}
  \variants{\SpellingVariant \vartext{sespen}}}

\entry{säspull}{\headword{säspull}
  \variantof{Dialectal Variant of \vartext{sispull}}  \pronnote{səspuɽ}}

\entry{säsramsäsram}{\headword{säsramsäsram}  \pronnote{səsramsəsram}
  \pos{Adverb}  \sense{
    \definition{shuffling}}}

\entry{säsri}{\headword{säsri}
  \variantof{SpellingVariant of \vartext{sisri}}}

\entry{sat}{\headword{sat}
  \pos{Noun}  \sense{
}}

\entry{satade}{\headword{satade}  \pronnote{satade}
  \pos{Noun}  \sense{
    \definition{Saturday}}}

\entry{saus}{\headword{saus}  \lexeme{saus}  \pronnote{saus}
  \pos{Noun}  \sense{
    \definition{type of yam with a white or purple interior, thorns, and no hairs}}}

\entry{Sawa}{\headword{Sawa}  \pronnote{sawa}
  \pos{Proper noun}  \sense{
    \definition{male personal name}}}

\entry{Sawapo}{\headword{Sawapo}  \pronnote{sawapo}
  \pos{Proper noun}  \sense{
    \definition{male personal name}}}

\entry{sawasap}{\headword{sawasap}  \note{Trees}  \pronnote{sawasap}
  \pos{Noun}  \sense{
    \definition{type of cultivated tree with edible yellow or green fruit}    \note{Trees}    \example{
}    \example{
}    \example{
}}}

\entry{sawe}{\headword{sawe}  \note{Adjectives (Swadesh)}  \pronnote{sawe}
  \pos{Modifier}  \sense{
    \definition{left}    \note{Adjectives (Swadesh)}    \example{
      \exsen{Ngämlle sawe alle giri de nanttog.}      \extran{Give me the knife with your left hand.}}    \example{
}    \example{
}}}

\entry{Saweta}{\headword{Saweta}  \pronnote{saweta}
  \pos{Proper noun}  \sense{
    \definition{Saweta (toponym)}    \example{
}    \example{
}    \example{
}}}

\entry{sawis}{\headword{sawis}  \pronnote{sawis}
  \pos{Noun}  \sense{
    \definition{type of small yam with a long shape and purple interior}}}

\entry{sawiya}{\headword{sawiya}  \pronnote{sawija}
  \pos{Noun}  \sense{
    \definition{little egret}    \scientificname{Egretta garzetta}}}

\entry{SDA}{\headword{SDA}
  \pos{Noun}  \sense{
    \definition{SDA (Seventh-day Adventist Church)}}}

\entry{se}{\headword{se}  \pronnote{se}
  \pos{Noun}  \sense{
    \definition{yes}    \example{
      \exsen{Se imomdae dan.}      \extran{Yes, this story is true.}}}}

\entry{se}{\headword{se}
  \variantof{freevarlab of \vartext{ase}}  \pronnote{se}}

\entry{Sebe}{\headword{Sebe}  \pronnote{sebe}
  \pos{Proper noun}  \sense{
    \definition{Sebe (Bine-speaking village in Oriomo-Bituri Rural LLG)}    \example{
}    \example{
}    \example{
}}}

\entry{seben}{\headword{seben}  \note{Numbers}  \pronnote{seben}
  \pos{Numeral}  \sense{
    \definition{seven (English numeral; also general)}    \note{Numbers}    \example{
}    \example{
}    \example{
}}
  \variants{\SpellingVariant \vartext{seven}}}

\entry{sebenti}{\headword{sebenti}
  \pos{Numeral}  \sense{
    \definition{seventy}}
  \variants{\SpellingVariant \vartext{seventi}}}

\entry{sebentin}{\headword{sebentin}
  \pos{Numeral}  \sense{
    \definition{seventeen (English numeral)}}}

\entry{sebor~}{\headword{sebor~}
  \variantof{Inflected Form of \vartext{RED~}}}

\entry{sebor}{\headword{sebor}  \note{Polite words and conversation}  \pronnote{sebor}
  \pos{Transitive coverb}  \sense{
    \definition{to greet}    \note{Polite words and conversation}    \example{
      \exsen{Angde bogo ma we bongttägän tämamae lla de sebor bägnegän.}      \extran{When she comes to the village, she greets every person.}}    \example{
}    \example{
}}}

\entry{seborsebor}{\headword{seborsebor}
  \pos{Transitive coverb}  \sense{
    \definition{to welcome}    \example{
      \exsen{Mälla da llo de seborsebor bägayaebneyo.}      \extran{The women will welcome the men.}}}
  \variants{\Fast Speech Variant \vartext{sebosebor}}}

\entry{sebosebor}{\headword{sebosebor}
  \variantof{Fast Speech Variant of \vartext{seborsebor}}}

\entry{second}{\headword{second}
  \variantof{Dialectal Variant of \vartext{English loanword}}}

\entry{secretary}{\headword{secretary}
  \variantof{Dialectal Variant of \vartext{English loanword}}}

\entry{sek}{\headword{sek}
  \pos{Transitive coverb}  \sense{
    \definition{to check}    \example{
      \exsen{Ngämo nyäng sek nägaeyo.}      \extran{[You all] check my bag.}}}}

\entry{sel}{\headword{sel}
  \pos{Transitive coverb}  \sense{
    \definition{to sell}    \example{
      \exsen{Bogo obo otät de sel anggan.}      \extran{He's selling his food.}}}}

\entry{sel}{\headword{sel}
  \pos{Noun}  \sense{
    \definition{cell}    \example{
      \exsen{Polis a namllamallo lla de ngämo gänye nätramallo, sel ma ik i nazänallo.}      \extran{The police took my husband and put him in the cell.}}}}

\entry{sele}{\headword{sele}  \note{Medicine}  \pronnote{sele}
  \pos{Noun}  \sense{
    \definition{chili}    \note{Medicine}    \example{
}    \example{
}    \example{
}}
  \variants{\Unspecified Variant \vartext{sili}}}

\entry{sem}{\headword{sem}  \note{Trees}  \pronnote{sem}
  \pos{Noun}  \sense{
    \definition{type of tree that grows in the bush; used for making rope}    \note{Trees}    \example{
}    \example{
}    \example{
}}}

\entry{sem}{\headword{sem}
  \pos{Modifier}  \sense{
    \definition{same}    \example{
      \exsen{sem ekawang lla}      \extran{person who speaks the same language}}}}

\entry{senedrin}{\headword{senedrin}  \pronnote{senedrin}
  \pos{Noun}  \sense{
    \definition{chief priest}}}

\entry{sens}{\headword{sens}
  \pos{Transitive/Intransitive coverb}  \sense{\textbf{1.} 
    \definition{to change}    \example{
      \exsen{Bin a sens gogon.}      \extran{The name changed.}}    \example{
      \exsen{Eka kutt di sens erallo.}      \extran{They are changing the word.}}}  \sense{\textbf{2.} 
    \definition{to exchange (in marriage)}    \example{
      \exsen{Obom sens dägagän.}      \extran{He exchanged her for marriage.}}}}

\entry{Senti}{\headword{Senti}
  \pos{Proper noun}  \sense{
    \definition{male personal name}}}

\entry{Sera}{\headword{Sera}
  \pos{Proper noun}  \sense{
    \definition{female personal name}}
  \variants{\SpellingVariant \vartext{Sarah}}}

\entry{ses}{\headword{ses}
  \variantof{SpellingVariant of \vartext{sos}}}

\entry{seseyam}{\headword{seseyam}
  \pos{Ideophone}  \sense{
    \definition{swish, sound of legs moving}}}

\entry{sespen}{\headword{sespen}
  \variantof{SpellingVariant of \vartext{säspen}}  \pronnote{sespen}}

\entry{seven}{\headword{seven}
  \variantof{SpellingVariant of \vartext{seben}}}

\entry{seventi}{\headword{seventi}
  \variantof{SpellingVariant of \vartext{sebenti}}}

\entry{shake}{\headword{shake}
  \variantof{Dialectal Variant of \vartext{English loanword}}}

\entry{Shamae}{\headword{Shamae}
  \variantof{SpellingVariant of \vartext{Samae}}}

\entry{Sharon}{\headword{Sharon}  \pronnote{sharon}
  \pos{Proper noun}  \sense{
    \definition{female personal name}}}

\entry{Shim}{\headword{Shim}
  \pos{Proper noun}  \sense{
    \definition{male personal name}}}

\entry{Sibideri}{\headword{Sibideri}  \pronnote{sibideri}
  \pos{Proper noun}  \sense{
    \definition{Sibidiri (Idi-speaking village in Morehead Rural LLG)}    \example{
}    \example{
}    \example{
}}
  \variants{\Unspecified Variant \vartext{Sibidir}}}

\entry{Sibidir}{\headword{Sibidir}
  \variantof{Unspecified Variant of \vartext{Sibideri}}}

\entry{Sibiya}{\headword{Sibiya}
  \pos{Proper noun}  \sense{
    \definition{male personal name}}}

\entry{Sibne}{\headword{Sibne}  \note{Places around Limol}  \pronnote{sibne}
  \pos{Proper noun}  \sense{
    \definition{Sibne (toponym)}    \note{Places around Limol}    \example{
}    \example{
}    \example{
}}}

\entry{sidompa}{\headword{sidompa}  \note{Traditional Arrows}  \pronnote{sidompa}
  \pos{Noun}  \sense{
    \definition{type of spear}    \note{Traditional Arrows}    \example{
}    \example{
}    \example{
}}}

\entry{Siga}{\headword{Siga}
  \pos{Proper noun}  \sense{
    \definition{Siga (toponym)}}}

\entry{Sigabaduru}{\headword{Sigabaduru}  \pronnote{siɡabaduru}
  \pos{Proper noun}  \sense{
    \definition{Sigabaduru (in Kiwai Rural LLG)}    \example{
}    \example{
}    \example{
}}}

\entry{sigip}{\headword{sigip}  \note{Types of palms}  \pronnote{siɡip}
  \pos{Noun}  \sense{\textbf{1.} 
    \definition{type of palm with fruit hanging from long pedicels; people chew the fruit like betelnut}    \note{Types of palms}    \example{
}    \example{
}    \example{
}}  \sense{\textbf{2.} 
    \definition{type of big taro}    \note{Types of palms}}}

\entry{siklakla}{\headword{siklakla}  \pronnote{siklakla}
  \pos{Noun}  \sense{\textbf{1.} 
    \definition{golden-headed cisticola}    \scientificname{Cisticola exilis}}  \sense{\textbf{2.} 
    \definition{Australian reed warbler}    \scientificname{Acrocephalus australis}}  \sense{\textbf{3.} 
    \definition{singing bush lark}    \scientificname{Mirafra javanica}}}

\entry{siks}{\headword{siks}  \note{Numbers}  \pronnote{siks}
  \pos{Numeral}  \sense{
    \definition{six (English numeral; also general)}    \note{Numbers}    \example{
}    \example{
}    \example{
}}
  \variants{\SpellingVariant \vartext{six}}}

\entry{siksti}{\headword{siksti}
  \pos{Numeral}  \sense{
    \definition{sixty}}
  \variants{\SpellingVariant \vartext{sixty}}}

\entry{sikstin}{\headword{sikstin}
  \pos{Numeral}  \sense{
    \definition{sixteen (English numeral)}}}

\entry{Siku}{\headword{Siku}
  \pos{Proper noun}  \sense{
    \definition{male personal name}}}

\entry{sili}{\headword{sili}
  \variantof{Unspecified Variant of \vartext{sele}}  \note{Poison}  \pronnote{sili}}

\entry{Simon}{\headword{Simon}
  \variantof{SpellingVariant of \vartext{Saemon}}}

\entry{sinensinen}{\headword{sinensinen}  \note{good characteristics of people}  \pronnote{sinensinen}
  \pos{Modifier}  \sense{
    \definition{generous, giving}    \note{good characteristics of people}    \example{
}    \example{
}}}

\entry{singoll}{\headword{singoll}  \pronnote{siŋoɽ}
  \pos{Transitive coverb}  \sense{
    \definition{to give, provide, share}    \example{
      \exsen{Ngämlle mɨnyi ngämo nag wätät yuatt singoll bagän.}      \extran{My friend will share cooked food with me.}}}}

\entry{singosingol}{\headword{singosingol}  \pronnote{siŋosiŋol}
  \pos{Adverb}  \sense{
    \definition{upwind, windward}    \example{
      \exsen{Bogo wel singosingol dallän.}      \extran{He went windward.}}}
  \variants{\Dialectal Variant \vartext{singosingoll}}}

\entry{singosingoll}{\headword{singosingoll}
  \variantof{Dialectal Variant of \vartext{singosingol}}}

\entry{Sini}{\headword{Sini}
  \pos{Proper noun}  \sense{
    \definition{female personal name}}}

\entry{Sintia}{\headword{Sintia}  \pronnote{sintia}
  \pos{Proper noun}  \sense{
    \definition{female personal name}}
  \variants{\SpellingVariant \vartext{Cynthia}}}

\entry{sip}{\headword{sip}  \pronnote{sip}
  \pos{Noun}  \sense{
    \definition{chief}}}

\entry{sip}{\headword{sip}  \pronnote{sip}
  \pos{Noun}  \sense{
    \definition{sheep}    \example{
      \exsen{Ibi sisiri Yesu bo sip dag.}      \extran{We are now Jesus' sheep.}}}}

\entry{sipel}{\headword{sipel}  \pronnote{sipel}
  \pos{Intransitive coverb}  \sense{
    \definition{to rest}    \example{
      \exsen{Angde pos ganen a gottamänän ngäna sipel gog.}      \extran{When the post planting finished, I took a rest.}}}}

\entry{sipik}{\headword{sipik}  \pronnote{sipik}
  \pos{Noun}  \sense{
    \definition{type of large yam with a white and purple interior}}}

\entry{siporo}{\headword{siporo}  \note{Medicine}  \pronnote{siporo}
  \pos{Noun}  \sense{
    \definition{type of cultivated tree with thorns and sour yellow fruit}    \note{Medicine}    \example{
}    \example{
}    \example{
}}}

\entry{sir}{\headword{sir}  \note{Trees}  \pronnote{sir}
  \pos{Noun}  \sense{
    \definition{type of tree that grows in the bush with white flowers and edible black fruit with one seed inside; cassowary and pigs eat the fruit}    \note{Trees}    \example{
}    \example{
}    \example{
}}}

\entry{sirem}{\headword{sirem}
  \variantof{SpellingVariant of \vartext{särem}}  \pronnote{sirem}}

\entry{Sirmitang}{\headword{Sirmitang}  \note{Places around Limol}  \pronnote{sirmitaŋ}
  \pos{Proper noun}  \sense{
    \definition{Sirmitang (toponym)}    \note{Places around Limol}    \example{
}    \example{
}    \example{
}}}

\entry{sis}{\headword{sis}  \note{Seasons}  \pronnote{sis}
  \pos{Noun}  \sense{
    \definition{season when new gardens are cleared and fenced (fourteenth season; corresponds to early November)}    \note{Seasons}    \example{
}    \example{
}    \example{
}}}

\entry{sis}{\headword{sis}  \note{Seasons}  \pronnote{sis}
  \pos{Noun}  \sense{
    \definition{type of flying ant that come out from their flooded anthills in November}    \note{Seasons}}}

\entry{sisel}{\headword{sisel}  \note{Tools}  \pronnote{sisel}
  \pos{Noun}  \sense{
    \definition{chisel}    \note{Tools}    \example{
}}}

\entry{sisi}{\headword{sisi}  \note{Type of pandanus}  \pronnote{sisi}
  \pos{Noun}  \sense{\textbf{1.} 
    \definition{type of pandanus with a trunk used for wood}    \note{Type of pandanus}    \example{
}    \example{
}    \example{
}}  \sense{\textbf{2.} 
    \definition{the third stage of coconut growth in which the seedling has formed (after planting)}    \note{Type of pandanus}}}

\entry{sisiri}{\headword{sisiri}
  \variantof{SpellingVariant of \vartext{sisri}}  \pronnote{sisiri}}

\entry{sisor}{\headword{sisor}  \note{Adjectives}  \pronnote{sisor}
  \pos{Modifier}  \sense{\textbf{1.} 
    \definition{new}    \note{Adjectives}    \example{
      \exsen{Känazbag, ngäna mɨnyi bibra sisor ttoenttoen de bɨllɨt.}      \extran{Tomorrow, I will tell you a new story.}}    \example{
}    \example{
}}  \sense{\textbf{2.} 
    \definition{young}    \note{Adjectives}    \example{
      \exsen{sisor lla}      \extran{young people}}}}

\entry{sisor däp}{\headword{sisor däp}  \pronnote{sisor dəp}
  \pos{Noun}  \sense{
    \definition{plant term}    \example{
      \exsen{Nge koko da sisor däp dan.}}    \example{
      \exsen{Pope up däm sisor däp de nibenyan.}}}}

\entry{sisor llɨg}{\headword{sisor llɨg}  \note{Stages of life}  \pronnote{sisor ɽɪɡ}
  \pos{Noun}  \sense{
    \definition{newborn, infant}    \note{Stages of life}    \example{
}    \example{
}    \example{
}}}

\entry{sisor pazi}{\headword{sisor pazi}  \note{Seasons}  \pronnote{sisor pazi}
  \pos{Noun}  \sense{
    \definition{season of New Years (sixteenth season; corresponds to December)}    \note{Seasons}    \example{
}    \example{
}    \example{
}}}

\entry{sispull}{\headword{sispull}  \note{Dead things}  \pronnote{sispuɽ}
  \pos{Noun}  \sense{
    \definition{maggot}    \note{Dead things}    \example{
}    \example{
}    \example{
}}
  \variants{\Dialectal Variant \vartext{säspull}}}

\entry{sisri}{\headword{sisri}  \pronnote{sisri}
  \pos{Temporal adverb}  \sense{\textbf{1.} 
    \definition{now}    \example{
      \exsen{Mo da sisri ttäkam allan.}      \extran{The bridge is breaking now.}}}  \sense{\textbf{2.} 
    \definition{this; current}    \example{
}}
  \variants{\SpellingVariant \vartext{säsri, säsäri, sisiri}}}

\entry{sisri ag}{\headword{sisri ag}
  \pos{Noun}  \sense{
    \definition{this morning}}}

\entry{sisri ebdo}{\headword{sisri ebdo}
  \pos{Noun}  \sense{
    \definition{today}    \example{
}}}

\entry{sisri iddob}{\headword{sisri iddob}
  \pos{Noun}  \sense{
}}

\entry{sisri toto}{\headword{sisri toto}
  \pos{Noun}  \sense{
    \definition{this evening}}}

\entry{Sisuar}{\headword{Sisuar}  \pronnote{sisuar}
  \pos{Proper noun}  \sense{
    \definition{female personal name}}}

\entry{siti}{\headword{siti}
  \pos{Noun}  \sense{
    \definition{city}}}

\entry{six}{\headword{six}
  \variantof{SpellingVariant of \vartext{siks}}}

\entry{sixty}{\headword{sixty}
  \variantof{SpellingVariant of \vartext{siksti}}}

\entry{sɨmell}{\headword{sɨmell}  \pronnote{sɪmeɽ}
  \pos{Noun}  \sense{\textbf{1.} 
    \definition{pig}    \example{
      \exsen{sɨmell mägda}      \extran{sow}}}  \sense{\textbf{2.} 
    \definition{truck}}
  \variants{\Unspecified Variant \vartext{sämell}}}

\entry{sɨmell källamokott}{\headword{sɨmell källamokott}  \note{Trees}  \pronnote{sɨmeɽ kəɽamokoʈ͡ʂ}
  \pos{Noun}  \sense{
    \definition{type of spiky tree that grows in the bush with white flowers, yellow fruit, and hardwood used for firewood}    \note{Trees}    \example{
}    \example{
}    \example{
}}}

\entry{sɨmellkom}{\headword{sɨmellkom}  \note{Sago palms}  \pronnote{sɨmeɽkom}
  \pos{Noun}  \sense{
    \definition{type of sago}    \note{Sago palms}    \example{
}    \example{
}    \example{
}}}

\entry{sɨmellmak}{\headword{sɨmellmak}  \note{Weaving}  \pronnote{sɨmeɽmak}
  \pos{Noun}  \sense{
    \definition{weaving pattern with arrows pointing right}    \note{Weaving}    \example{
}    \example{
}    \example{
}}}

\entry{sɨngesɨnge}{\headword{sɨngesɨnge}  \pronnote{sɪŋesɪŋe}
  \pos{Noun}  \sense{\textbf{1.} 
    \definition{admiration}}  \sense{\textbf{2.} 
    \definition{to admire}}}

\entry{sɨrem}{\headword{sɨrem}
  \variantof{SpellingVariant of \vartext{särem}}  \pronnote{sɪrem}}

\entry{sɨs}{\headword{sɨs}  \pronnote{sɪs}
  \pos{Transitive verb}  \sense{
    \definition{to extinguish, turn off}    \example{
      \exsen{Ngäna tos de dɨs.}      \extran{I turned off the light.}}}}

\entry{sɨs}{\headword{sɨs}  \pronnote{sɪs}
  \pos{Interjection}  \sense{
    \definition{command given to a dog to chase an animal}}}

\entry{Skola}{\headword{Skola}  \pronnote{skola}
  \pos{Proper noun}  \sense{
    \definition{female personal name}}}

\entry{skul}{\headword{skul}  \note{Social problems}  \pronnote{skul}
  \pos{Noun}  \sense{
    \definition{school; education}    \note{Social problems}    \example{
}    \example{
}    \example{
}}
  \variants{\SpellingVariant \vartext{sukul}}}

\entry{skulang}{\headword{skulang}  \note{School}  \pronnote{skulaŋ}
  \pos{Noun}  \sense{
    \definition{student}    \note{School}    \example{
}    \example{
}    \example{
}}}

\entry{skulatt}{\headword{skulatt}
  \pos{Modifier}  \sense{
    \definition{educated}}}

\entry{slaslak}{\headword{slaslak}  \pronnote{slaslak}
  \pos{Noun}  \sense{
    \definition{red-winged parrot}    \scientificname{Aprosmictus erythropterus coccineopterus}    \example{
}    \example{
}    \example{
}}}

\entry{sllollongg}{\headword{sllollongg}  \pronnote{sɽoɽoŋɡ}
  \pos{Intransitive verb}  \sense{
    \definition{to sit close together}    \example{
      \exsen{Ubi komllaebmae ansllollonggallo duliballe.}      \extran{The two of them sat close together on that side.}}}}

\entry{so}{\headword{so}  \note{Types of palms}  \pronnote{so}
  \pos{Noun}  \sense{
    \definition{type of black palm; in this mature stage (~3 m), the wood is used for house flooring and containers for squeezing sago, and the pith is eaten}    \note{Types of palms}    \example{
}    \example{
}    \example{
}}}

\entry{so}{\headword{so}
  \variantof{Dialectal Variant of \vartext{English loanword}}}

\entry{Soba}{\headword{Soba}  \pronnote{soba}
  \pos{Proper noun}  \sense{
    \definition{male personal name}}}

\entry{Sobam}{\headword{Sobam}  \pronnote{sobam}
  \pos{Proper noun}  \sense{
    \definition{male personal name}}}

\entry{Sobeya}{\headword{Sobeya}
  \pos{Proper noun}  \sense{
    \definition{Sobeya (toponym)}}}

\entry{soccer}{\headword{soccer}
  \variantof{SpellingVariant of \vartext{soka}}}

\entry{sod}{\headword{sod}  \note{Clothing}  \pronnote{sod}
  \pos{Noun}  \sense{
    \definition{shirt}    \note{Clothing}    \example{
}    \example{
}    \example{
}}}

\entry{Sogale}{\headword{Sogale}  \pronnote{soɡale}
  \pos{Proper noun}  \sense{
    \definition{Sogale (Bine-speaking village in Oriomo-Bituri Rural LLG)}    \example{
}    \example{
}    \example{
}}}

\entry{soka}{\headword{soka}
  \pos{Noun}  \sense{
    \definition{soccer}}
  \variants{\SpellingVariant \vartext{soccer}}}

\entry{Soka}{\headword{Soka}  \pronnote{soka}
  \pos{Proper noun}  \sense{
    \definition{male personal name}}}

\entry{Sokola}{\headword{Sokola}
  \pos{Proper noun}  \sense{
    \definition{female personal name}}}

\entry{sokpa}{\headword{sokpa}  \pronnote{sokpa}
  \pos{Noun}  \sense{\textbf{1.} 
    \definition{tobacco}    \example{
}}  \sense{\textbf{2.} 
    \definition{cigarette}    \example{
      \exsen{Iba sokpa dag gänyme gall toko me nanen e}      \extran{Our (incl.) cigarettes are here on the canoe for smoking.}}}}

\entry{sokpa kllokllop}{\headword{sokpa kllokllop}
  \pos{Noun}  \sense{
    \definition{notebook-sized mat woven of tobacco}}}

\entry{sokpa llaweatt}{\headword{sokpa llaweatt}
  \pos{Noun}  \sense{
    \definition{tobacco woven like a rope}}}

\entry{solt}{\headword{solt}
  \pos{Noun}  \sense{
    \definition{salt}}}

\entry{Soma}{\headword{Soma}  \pronnote{soma}
  \pos{Proper noun}  \sense{
    \definition{male personal name}}}

\entry{somungg}{\headword{somungg}
  \pos{Transitive verb}  \sense{
    \definition{to realize}}}

\entry{Songno}{\headword{Songno}
  \pos{Proper noun}  \sense{
    \definition{female personal name}}}

\entry{sop}{\headword{sop}  \lexeme{sop}  \pronnote{sop}
  \pos{Noun}  \sense{
    \definition{soap}}}

\entry{sori}{\headword{sori}
  \variantof{SpellingVariant of \vartext{sari}}}

\entry{soroe}{\headword{soroe}  \note{Verbs}  \pronnote{soroe}
  \pos{Transitive verb}  \sense{\textbf{1.} 
    \definition{to try, attempt}    \note{Verbs}    \example{
      \exsen{Ongg kame ako dasäroeyän.}      \extran{Ongg tried it again.}}    \example{
      \exsen{Ubi mɨnyi oba mängall alle bosroenegnän ngämene eka ngänaem e.}      \extran{They will be trying hard to understand my words.}}    \example{
}}  \sense{\textbf{2.} 
    \definition{to challenge, try, test; tempt}    \note{Verbs}    \example{
      \exsen{Ubi ngonoe eka walle obom dasroeyo.}      \extran{They challenged him with a question.}}    \example{
      \exsen{Bogo dedme poti ebdo dagirnän a saeten dasroenän.}      \extran{He stayed there for forty days and Satan was tempting him.}}}  \sense{\textbf{3.} 
    \definition{test, exam}    \note{Verbs}    \example{
      \exsen{Amo bun da soroe me mullae gogalle, bogo mɨnyi dallalle tuk skul e.}      \extran{He who does well on the test will go to secondary school.}}}
  \variants{\SpellingVariant \vartext{säroe}}}

\entry{sorry}{\headword{sorry}
  \variantof{SpellingVariant of \vartext{sari}}}

\entry{sorsor}{\headword{sorsor}  \pronnote{sorsor}
  \pos{Verb}  \sense{
    \definition{to affect at a distance. For example, if poison is poured up stream, the fish down stream will also be affected.}}}

\entry{sos}{\headword{sos}  \pronnote{sos}
  \pos{Noun}  \sense{
    \definition{church}    \example{
      \exsen{Tämamae wiyamom sos e.}      \extran{Everyone, come to church.}}}
  \variants{\SpellingVariant \vartext{ses, ttärtt}}}

\entry{sosa}{\headword{sosa}
  \variantof{Dialectal Variant of \vartext{English loanword}}}

\entry{sosoga}{\headword{sosoga}  \note{Sago palms}  \pronnote{sosoɡa}
  \pos{Noun}  \sense{
    \definition{type of sago}    \note{Sago palms}    \example{
}    \example{
}    \example{
}}}

\entry{Sowa}{\headword{Sowa}
  \pos{Proper noun}  \sense{
    \definition{male personal name}}}

\entry{Sowati}{\headword{Sowati}  \pronnote{sowati}
  \pos{Proper noun}  \sense{
    \definition{male personal name}}}

\entry{spalek}{\headword{spalek}
  \variantof{Fast Speech Variant of \vartext{säpalek}}  \pronnote{spalek}}

\entry{spallma}{\headword{spallma}  \note{Friendship/social terms of endearment}  \pronnote{spaɽma}
  \pos{Noun}  \sense{
    \definition{two friends who split a twin coconut frond}    \note{Friendship/social terms of endearment}    \example{
}    \example{
}    \example{
}}}

\entry{spun}{\headword{spun}  \pronnote{spun}
  \pos{Intransitive verb}  \sense{\textbf{1.} 
    \definition{to fall; set}    \example{
      \exsen{Sɨmell a kuddäll a guspunän.}      \extran{The pig fell down dead.}}    \example{
      \exsen{Kakoll a obo ttang atta aspunan.}      \extran{The plate fell from his hand.}}    \example{
      \exsen{Yäbäd a zäme guspunmällnän.}      \extran{The sun was already setting.}}}  \sense{\textbf{2.} 
    \definition{to jump}    \example{
      \exsen{Däbe llɨg walle we guspunän.}      \extran{That boy jumped into the water.}}}  \sense{\textbf{3.} 
    \definition{to throw; shoot}    \example{
      \exsen{Ngämi tudi daspullaemnalla.}      \extran{We (excl.) started throwing fishing lines.}}    \example{
      \exsen{Ako kame kakoll de katkatre toko we naspullan mängall e enanae opallknegan.}      \extran{Again, she threw the dishes on top of the table forcefully and they finally broke.}}    \example{
      \exsen{Ibi mɨnyi damärärmae toboll baspunigalla.}      \extran{We (incl.) will both shoot arrows at the same time.}}}  \sense{\textbf{4.} 
    \definition{to remove an outer layer}    \example{
      \exsen{Ttoe de naspunan.}      \extran{He removed the skin.}}}
  \variants{\Dialectal Variant \vartext{supun}}}

\entry{spun}{\headword{spun}  \pronnote{spun}
  \pos{Noun}  \sense{
    \definition{spoon}}}

\entry{stael}{\headword{stael}
  \variantof{freevarlab of \vartext{English loanword}}}

\entry{Stanis}{\headword{Stanis}  \pronnote{stanis}
  \pos{Proper noun}  \sense{
    \definition{male personal name}}}

\entry{Stashalyn}{\headword{Stashalyn}  \pronnote{stæʃlɪn}
  \pos{Proper noun}  \sense{
    \definition{female personal name}}}

\entry{Steven}{\headword{Steven}
  \variantof{SpellingVariant of \vartext{Stibin}}  \pronnote{steven}}

\entry{Stibin}{\headword{Stibin}
  \pos{Proper noun}  \sense{
    \definition{male personal name}}
  \variants{\SpellingVariant \vartext{Steven, Stiven}}}

\entry{still}{\headword{still}
  \variantof{Dialectal Variant of \vartext{English loanword}}}

\entry{Stiven}{\headword{Stiven}
  \variantof{SpellingVariant of \vartext{Stibin}}}

\entry{stoa}{\headword{stoa}  \note{Market}  \pronnote{stoa}
  \pos{Noun}  \sense{
    \definition{store}    \note{Market}    \example{
      \exsen{Gull tupi di stoa me llädnan amallo.}      \extran{They bought long nets in the store.}}    \example{
}    \example{
}}
  \variants{\SpellingVariant \vartext{stowa}}}

\entry{stori}{\headword{stori}
  \pos{Noun}  \sense{
    \definition{story}}
  \variants{\SpellingVariant \vartext{story}}}

\entry{story}{\headword{story}
  \variantof{SpellingVariant of \vartext{stori}}}

\entry{stowa}{\headword{stowa}
  \variantof{SpellingVariant of \vartext{stoa}}}

\entry{student}{\headword{student}
  \variantof{Dialectal Variant of \vartext{English loanword}}}

\entry{su}{\headword{su}  \pronnote{su}
  \pos{Noun}  \sense{\textbf{1.} 
    \definition{prey}}  \sense{\textbf{2.} 
    \definition{secret}    \example{
      \exsen{Su eka midd a zäme bina pate indrang gognegän.}      \extran{The secret meaning has already been revealed to you all.}}}}

\entry{Suame}{\headword{Suame}  \note{Place names}  \pronnote{suame}
  \pos{Proper noun}  \sense{
    \definition{Suame (toponym)}    \note{Place names}    \example{
}    \example{
}    \example{
}}}

\entry{suang}{\headword{suang}  \note{Bad characteristics of people}  \pronnote{suaŋ}
  \pos{Noun}  \sense{
    \definition{selfish}    \note{Bad characteristics of people}    \example{
}    \example{
}}}

\entry{suga}{\headword{suga}  \pronnote{suɡa}
  \pos{Noun}  \sense{
    \definition{sugar}}}

\entry{suga galbe}{\headword{suga galbe}  \pronnote{suɡa ɡalbe}
  \pos{Noun}  \sense{
    \definition{type of large yam with a white interior}}}

\entry{Sui}{\headword{Sui}
  \variantof{SpellingVariant of \vartext{Suwi}}}

\entry{Suki}{\headword{Suki}  \pronnote{suki}
  \pos{Proper noun}  \sense{
    \definition{Suki (in Morehead Rural LLG)}    \example{
}    \example{
}    \example{
}}}

\entry{sukul}{\headword{sukul}
  \variantof{SpellingVariant of \vartext{skul}}  \pronnote{sukul}}

\entry{Suliki}{\headword{Suliki}  \pronnote{suliki}
  \pos{Proper noun}  \sense{
    \definition{male personal name}}}

\entry{sulut}{\headword{sulut}  \note{Types of Taro}  \pronnote{sulut}
  \pos{Noun}  \sense{
    \definition{type of taro}    \note{Types of Taro}    \example{
}    \example{
}    \example{
}}}

\entry{supun}{\headword{supun}
  \variantof{Dialectal Variant of \vartext{spun}}  \pronnote{supun}}

\entry{sur}{\headword{sur}  \lexeme{sur}  \pronnote{sur}
  \pos{Noun}  \sense{
    \definition{pushing tool}}}

\entry{surum}{\headword{surum}
  \pos{Noun}  \sense{
    \definition{sand}}}

\entry{surusuru}{\headword{surusuru}  \note{Trees}  \pronnote{surusuru}
  \pos{Noun}  \sense{
    \definition{type of tree that grows in the bush with white flowers and strong wood used for house posts and firewood; it was traditionally lit and used as a light at night because it burns slowly}    \scientificname{Cordia subcordata}    \note{Trees}    \example{
}    \example{
}    \example{
}}}

\entry{Susan}{\headword{Susan}  \pronnote{susan}
  \pos{Proper noun}  \sense{
    \definition{female personal name}}}

\entry{susu}{\headword{susu}  \pronnote{susu}
  \pos{Noun}  \sense{
    \definition{breastmilk}}
  \variants{\Baby Talk Variant \vartext{tutu}}}

\entry{suwe}{\headword{suwe}  \pronnote{suwe}
  \pos{Noun}  \sense{\textbf{1.} 
    \definition{urine, pee}    \example{
}    \example{
      \exsen{Yokon wälläng e dallän suwe ma.}      \extran{Yokon went to the bush to pee.}}    \example{
}}  \sense{\textbf{2.} 
    \definition{(euphemistic) testicles}}}

\entry{suwe gäd}{\headword{suwe gäd}  \note{Parts of the body}  \pronnote{suwe ɡəd}
  \pos{Noun}  \sense{
    \definition{body part}    \note{Parts of the body}    \example{
}    \example{
}    \example{
}}}

\entry{suwe tängg}{\headword{suwe tängg}  \lexeme{tängg}  \pronnote{təŋɡ}
  \pos{Transitive verb}  \sense{
    \definition{to urinate on, pee on}    \example{
      \exsen{Ekaklle de suwe näntägan.}      \extran{He peed on the ground.}}}}

\entry{suwe tɨngg}{\headword{suwe tɨngg}
  \pos{Transitive compound verb}  \sense{
    \definition{to urinate on}}}

\entry{suwe ttätt}{\headword{suwe ttätt}  \lexeme{ttätt}  \pronnote{ʈ͡ʂəʈ͡ʂ}
  \pos{Intransitive verb}  \sense{\textbf{1.} 
    \definition{to urinate}}  \sense{\textbf{2.} 
    \definition{surprise}}}

\entry{Suwede}{\headword{Suwede}  \pronnote{suwede}
  \pos{Proper noun}  \sense{
    \definition{male personal name}}}

\entry{Suwi}{\headword{Suwi}  \pronnote{suwi}
  \pos{Proper noun}  \sense{
    \definition{Sui (in Kiwai Rural LLG)}    \example{
}    \example{
}    \example{
}}
  \variants{\SpellingVariant \vartext{Sui}}}

\entry{Sylvien}{\headword{Sylvien}  \pronnote{sɪlvijɛn}
  \pos{Proper noun}  \sense{
    \definition{female personal name}}}

\chapter{T}


\entry{ta~}{\headword{ta~}
  \variantof{Infinitival reduplicant of \vartext{RED~}}}

\entry{täb}{\headword{täb}  \pronnote{təb}
  \pos{Noun}  \sense{
    \definition{type of tree}}}

\entry{tab}{\headword{tab}  \pronnote{tab}
  \pos{Noun}  \sense{\textbf{1.} 
    \definition{promise, oath; engagement}    \example{
      \exsen{Tab de obo pate dangesän.}      \extran{She made him a promise.}}}  \sense{\textbf{2.} 
    \definition{to promise}    \example{
      \exsen{Ngäna bibim tab anggan.}      \extran{I promise you all.}}}}

\entry{täbab}{\headword{täbab}
  \pos{Transitive verb}  \sense{
    \definition{to watch over}}}

\entry{täbädd}{\headword{täbädd}  \pronnote{təbəɖ͡ʐ}
  \pos{Noun}  \sense{
    \definition{guest, visitor, stranger}    \example{
      \exsen{Eran ngämo täbädd känttatt a erame?}      \extran{Where is my guest room?}}}}

\entry{täbäll}{\headword{täbäll}
  \variantof{SpellingVariant of \vartext{tobäll}}  \pronnote{təbəɽ}}

\entry{täbäll pud}{\headword{täbäll pud}  \note{Bees}  \pronnote{təbəɽ pud}
  \pos{Noun}  \sense{\textbf{1.} 
    \definition{type of biting bee found in trees}    \note{Bees}    \example{
}    \example{
}    \example{
}}  \sense{\textbf{2.} 
    \definition{type of tree}    \note{Bees}}}

\entry{täbatäbe}{\headword{täbatäbe}  \pronnote{təbatəbe}
  \pos{Intransitive verb}  \sense{\textbf{1.} 
    \definition{to plan}    \example{
      \exsen{Ngäna gotäba ngämo nge ibeny e.}      \extran{I planned to plant my coconut.}}}  \sense{\textbf{2.} 
    \definition{plan}    \example{
      \exsen{Ngämi diba toto me ngäma täbatäbe de dangeseya iddob me ddia ma koenmäll e.}      \extran{We (excl.) made the plan that evening to chase the deer at night.}}}}

\entry{täbe}{\headword{täbe}  \note{Trees}  \pronnote{təbe}
  \pos{Noun}  \sense{
    \definition{type of big tree that grows in the bush with white flowers and drupes that fall with an edible seed inside}    \note{Trees}    \example{
}    \example{
}    \example{
}}}

\entry{tabe}{\headword{tabe}  \pronnote{tabe}
  \pos{Noun}  \sense{\textbf{1.} 
    \definition{soft pad or support that placed between one's back and a spalek basket}}  \sense{\textbf{2.} 
    \definition{to pad your back from the spalek}}}

\entry{täbe budar}{\headword{täbe budar}  \note{Grubs}  \pronnote{təbe budar}
  \pos{Noun}  \sense{
    \definition{type of grub}    \note{Grubs}    \example{
}    \example{
}    \example{
}}}

\entry{täbe lläkäm}{\headword{täbe lläkäm}  \note{Mushrooms}  \pronnote{təbe ɽəkəm}
  \pos{Noun}  \sense{
    \definition{type of mushroom}    \note{Mushrooms}    \example{
}    \example{
}    \example{
}}}

\entry{täbie}{\headword{täbie}  \note{Birds}  \pronnote{təbie}
  \pos{Noun}  \sense{\textbf{1.} 
    \definition{black sunbird}    \scientificname{Leptocoma aspasia}    \note{Birds}    \example{
}    \example{
}    \example{
}}  \sense{\textbf{2.} 
    \definition{garden sunbird}    \scientificname{Cinnyris jugularis}    \note{Birds}}}

\entry{Tabita}{\headword{Tabita}
  \pos{Proper noun}  \sense{
    \definition{female personal name}}}

\entry{täbom}{\headword{täbom}  \note{Food}  \pronnote{təbom}
  \pos{Noun}  \sense{
    \definition{type of small yam with a white interior, hairs, and thorns}    \note{Food}    \example{
}    \example{
}    \example{
}}}

\entry{Tabubil}{\headword{Tabubil}
  \pos{Proper noun}  \sense{
    \definition{Tabubil (toponym)}}}

\entry{tada}{\headword{tada}  \note{Traditional fishing ways and nets}  \pronnote{tada}
  \pos{Noun}  \sense{
    \definition{fish trap}    \note{Traditional fishing ways and nets}    \example{
      \exsen{Ngäna kollba de tada walle iringän eran.}      \extran{I am catching the fish with the trap.}}    \example{
}    \example{
}}}

\entry{tae}{\headword{tae}  \lexeme{tae}  \pronnote{tae}
  \pos{Noun}  \sense{
    \definition{type of tall tree that grows by rivers with white flowers, green fruit, and sturdy wood that can be used as a bridge}}}

\entry{tae budar}{\headword{tae budar}  \note{Hunting}  \pronnote{tae budar}
  \pos{Noun}  \sense{
    \definition{type of edible grub}    \note{Hunting}    \example{
}    \example{
}    \example{
}}}

\entry{tae lläkäm}{\headword{tae lläkäm}  \pronnote{tae ɽəkəm}
  \pos{Noun}  \sense{
    \definition{type of mushroom}    \example{
}    \example{
}    \example{
}}}

\entry{taem}{\headword{taem}
  \pos{Noun}  \sense{
    \definition{time}    \example{
      \exsen{ddob taem me}      \extran{sometimes}}}
  \variants{\SpellingVariant \vartext{taim}}}

\entry{taemataema}{\headword{taemataema}  \note{Trees}  \pronnote{taemataema}
  \pos{Noun}  \sense{
    \definition{type of tree that grows near the swamp with long yellow flowers and leaves that are used to treat fungal skin infections}    \note{Trees}    \example{
}    \example{
}    \example{
}}}

\entry{Taeme}{\headword{Taeme}
  \variantof{SpellingVariant of \vartext{Tame}}  \pronnote{taeme}}

\entry{Taeme loanword}{\headword{Taeme loanword}
  \pos{}  \sense{
}
  \variants{\Dialectal Variant \vartext{noto}}}

\entry{taempäg}{\headword{taempäg}  \pronnote{taempəɡ}
  \pos{Transitive verb}  \sense{
    \definition{to show, indicate, reveal}    \example{
      \exsen{Bogo mɨnyi tuk me ulle abal känttatt de yantepägän.}      \extran{He will show you all a very large bedroom upstairs.}}    \example{
      \exsen{Adi gontepägän obo täbatäbe de.}      \extran{God revealed his plan.}}}}

\entry{taewa}{\headword{taewa}  \note{Birds}  \pronnote{taewa}
  \pos{Noun}  \sense{
    \definition{type of bird}    \note{Birds}    \example{
}    \example{
}    \example{
}}}

\entry{Tag}{\headword{Tag}  \pronnote{taɡ}
  \pos{Proper noun}  \sense{
    \definition{personal name}}}

\entry{tägab}{\headword{tägab}
  \pos{Transitive verb}  \sense{
    \definition{to turn upside down}    \example{
      \exsen{Ubi obom tutu we dirngäneyo a ingoll dätägabeyo ine gollab e obo bod att.}      \extran{They pulled her onto land and turned her upside down so the water would come out of her mouth.}}}}

\entry{Tai}{\headword{Tai}
  \variantof{SpellingVariant of \vartext{Tayi}}}

\entry{taim}{\headword{taim}
  \variantof{SpellingVariant of \vartext{taem}}  \pronnote{taim}}

\entry{tainäm}{\headword{tainäm}
  \pos{Noun}  \sense{
    \definition{mosquito net}}}

\entry{täk}{\headword{täk}  \pronnote{tək}
  \pos{Noun}  \sense{
    \definition{clitoris}}}

\entry{täkäll}{\headword{täkäll}  \pronnote{təkəɽ}
  \pos{Noun}  \sense{\textbf{1.} 
    \definition{fin}    \example{
      \exsen{kollba täkäll}      \extran{fish fin}}}  \sense{\textbf{2.} 
    \definition{horn}    \example{
      \exsen{ddia täkäll}      \extran{deer horn}}}  \sense{\textbf{3.} 
    \definition{thorn}    \example{
      \exsen{siporo täkäll}      \extran{lemon thorns}}}}

\entry{täkälltäkäll kutt}{\headword{täkälltäkäll kutt}  \note{Parts of a fish}  \pronnote{təkəɽtəkəɽ kuʈ͡ʂ}
  \pos{Noun}  \sense{
    \definition{spine}    \note{Parts of a fish}    \example{
}    \example{
}    \example{
}}}

\entry{täkälluit}{\headword{täkälluit}  \note{Parts of a fish}  \pronnote{təkəɽuit}
  \pos{Noun}  \sense{
    \definition{radial cartilage}    \note{Parts of a fish}    \example{
}    \example{
}    \example{
}}}

\entry{Takeya}{\headword{Takeya}
  \pos{Proper noun}  \sense{
    \definition{female personal name}}}

\entry{täkla}{\headword{täkla}  \note{Trees}  \pronnote{təkla}
  \pos{Noun}  \sense{
    \definition{tree type}    \note{Trees}    \example{
}    \example{
}    \example{
}}}

\entry{täko}{\headword{täko}  \pronnote{təko}
  \pos{Noun}  \sense{
    \definition{any body part}}}

\entry{täl}{\headword{täl}  \note{Trees}  \pronnote{təl}
  \pos{Noun}  \sense{
    \definition{type of large bamboo that grows anywhere; used for bows, bow strings, axe handles, spears, and clothing pegs}    \note{Trees}    \example{
}    \example{
}    \example{
}}}

\entry{täl wadär}{\headword{täl wadär}  \note{Hunting}  \pronnote{təl wadər}
  \pos{Noun}  \sense{
    \definition{bamboo string}    \note{Hunting}    \example{
}    \example{
}    \example{
}}}

\entry{talapia}{\headword{talapia}
  \pos{Noun}  \sense{
    \definition{tilapia}}}

\entry{täli}{\headword{täli}  \pronnote{təli}
  \pos{Transitive verb}  \sense{
    \definition{to repeat}    \example{
      \exsen{Bogo dallän ikopse we, dantlinigän däbem eka dae.}      \extran{He went to pray, repeating just those words.}}}}

\entry{tall}{\headword{tall}  \pronnote{taɽ}
  \pos{Noun}  \sense{
    \definition{type of tree with wood used for posts and easily-peeling bark used for walling; also lit as a torch}}}

\entry{Tallabunang}{\headword{Tallabunang}  \pronnote{taɽabunaŋ}
  \pos{Proper noun}  \sense{
    \definition{Tallabunang (toponym)}}}

\entry{talme}{\headword{talme}  \note{Birth}  \pronnote{talme}
  \pos{Transitive/Intransitive coverb}  \sense{
    \definition{to give birth (to)}    \note{Birth}    \example{
      \exsen{Ttongdae ebdo meae bode do Malläm me talme gogän.}      \extran{She gave birth in Malläm on the same day.}}    \example{
      \exsen{Obom talme dägagän.}      \extran{She gave birth to her.}}    \example{
}}}

\entry{tälpae}{\headword{tälpae}
  \pos{Verb}  \sense{
    \definition{to give oneself completely, even to death}}}

\entry{tälpe}{\headword{tälpe}  \pronnote{təlpe}
  \pos{Intransitive verb}  \sense{
    \definition{to volunteer}    \example{
      \exsen{Bogo atälpeyan.}      \extran{He volunteered.}}}}

\entry{tältäl}{\headword{tältäl}  \pronnote{təltəl}
  \pos{Noun}  \sense{
    \definition{type of grass}}}

\entry{Täm}{\headword{Täm}
  \pos{Proper noun}  \sense{
    \definition{female personal name}}}

\entry{täma}{\headword{täma}  \pronnote{təma}
  \pos{Noun}  \sense{
    \definition{husk, exocarp/mesocarp (of coconut)}}}

\entry{tämallang mälla}{\headword{tämallang mälla}  \note{Types of Taro}  \pronnote{təmaɽaŋ məɽa}
  \pos{Noun}  \sense{
    \definition{type of big taro}    \note{Types of Taro}    \example{
}    \example{
}    \example{
}}}

\entry{tämamae}{\headword{tämamae}  \pronnote{təmamae}
  \pos{Quantifier}  \sense{\textbf{1.} 
    \definition{all, every}    \example{
      \exsen{tämamae ngalen, tämamae lla}      \extran{everything, everyone}}    \example{
      \exsen{Kollba da tämamae dadrowän.}      \extran{All the fish died.}}}  \sense{\textbf{2.} 
    \definition{whole, entire}    \example{
      \exsen{tämamae bem gallgall}      \extran{the entire coastline}}}}

\entry{tämani}{\headword{tämani}  \lexeme{tämani}  \pronnote{təmani}
  \pos{Noun}  \sense{
    \definition{type of large yam with a white or white and red interior}}}

\entry{tamatama}{\headword{tamatama}  \lexeme{tamatama}  \pronnote{tamatama}
  \pos{Noun}  \sense{
    \definition{bean}}}

\entry{tämbag}{\headword{tämbag}  \pronnote{təmbaɡ}
  \pos{Transitive verb}  \sense{
    \definition{to add}}}

\entry{tämbameny}{\headword{tämbameny}  \pronnote{təmbameɲ}
  \pos{Transitive verb}  \sense{
    \definition{to instruct}    \example{
      \exsen{Ause da bom däntäbemenyän, "Abo bongo goeg de naemponomeny."}      \extran{The old woman instructed him, "You must burn the garden."}}}}

\entry{tame}{\headword{tame}  \note{Water}  \pronnote{tame}
  \pos{Noun}  \sense{
    \definition{wave (of water)}    \note{Water}    \example{
}    \example{
}    \example{
}}}

\entry{täme}{\headword{täme}  \note{Nature}  \pronnote{təme}
  \pos{Noun}  \sense{
    \definition{monitor lizard, goanna}    \scientificname{Varanus}    \note{Nature}    \example{
}    \example{
}    \example{
}}}

\entry{Tame}{\headword{Tame}  \note{Language Names}  \pronnote{tame}
  \pos{Proper noun}  \sense{
    \definition{Taeme language (Pahoturi River language spoken in Kinkin alongside Ende)}    \note{Language Names}    \example{
}    \example{
}    \example{
}}
  \variants{\Dialectal Variant \vartext{Teme}}
  \variants{\SpellingVariant \vartext{Taeme}}}

\entry{täme käp}{\headword{täme käp}  \note{Food}  \pronnote{təme kəp}
  \pos{Noun}  \sense{
    \definition{type of round yam with a white interior, red skin, and hairs}    \note{Food}    \example{
}    \example{
}    \example{
}}}

\entry{täme sära}{\headword{täme sära}  \note{Gardens}  \pronnote{təme səra}
  \pos{Noun}  \sense{
    \definition{bush rope}    \note{Gardens}    \example{
}    \example{
}    \example{
}}}

\entry{tameny}{\headword{tameny}  \note{School}  \pronnote{tameɲ}
  \pos{Transitive verb}  \sense{\textbf{1.} 
    \definition{to teach}    \note{School}    \example{
      \exsen{Ngäna bam bantemeny walle gllangglla de.}      \extran{I'm going to teach you to swim.}}    \example{
      \exsen{Yesu ubira yuwog abal ttoen de sanawang eka walle däntamenynegnän.}      \extran{Jesus was teaching them many things using parables.}}}  \sense{\textbf{2.} 
    \definition{to dicuss, converse}    \note{School}    \example{
}}  \sense{\textbf{3.} 
    \definition{to tell; converse with, speak to}    \note{School}}}

\entry{tamenyang}{\headword{tamenyang}  \note{School}  \pronnote{tameɲaŋ}
  \pos{Noun}  \sense{
    \definition{teacher}    \note{School}    \example{
}    \example{
}    \example{
}}}

\entry{tamllägtamlläg}{\headword{tamllägtamlläg}  \note{Poison}  \pronnote{tamɽəɡtamɽəɡ}
  \pos{Noun}  \sense{
    \definition{type of small caterpillar}    \note{Poison}    \example{
}    \example{
}    \example{
}}}

\entry{tamllängg}{\headword{tamllängg}  \pronnote{tamɽəŋɡ}
  \pos{Transitive verb}  \sense{
    \definition{to press onto something soft, maybe it will break (like the computer)}}}

\entry{tamongle}{\headword{tamongle}  \pronnote{tamoŋle}
  \pos{Noun}  \sense{
    \definition{type of mutae yam shaped like a foot}}}

\entry{tämpeyam}{\headword{tämpeyam}
  \pos{Transitive verb}  \sense{
    \definition{to abandon, give up}    \example{
      \exsen{Bongo ewatta ke ngänäm nantäpeyamalle?}      \extran{Why have you abandoned me?}}    \example{
      \exsen{Paelet Yesu bom Rom mäk lla yaba pate däntäpeyamän ddänggaddängge we.}      \extran{Pilate gave Jesus up to the Roman soldiers to be crucified.}}}}

\entry{tän}{\headword{tän}  \pronnote{tən}
  \pos{Noun}  \sense{\textbf{1.} 
    \definition{tribe, nation, people}    \example{
      \exsen{Ende tän}      \extran{the Ende people}}}  \sense{\textbf{2.} 
    \definition{clan}    \example{
      \exsen{Ngämi lla tän dag.}      \extran{We (excl.) are of the same clan.}}}  \sense{\textbf{3.} 
    \definition{stem}}}

\entry{tan}{\headword{tan}  \note{Tools}  \pronnote{tan}
  \pos{Noun}  \sense{\textbf{1.} 
    \definition{type of short plant with white flowers, brown fruit, and branches used as a broom}    \note{Tools}}  \sense{\textbf{2.} 
    \definition{broom, can come from different types of palm e.g. käg tan, mäta tan}    \note{Tools}    \example{
}    \example{
}    \example{
}}}

\entry{tänan}{\headword{tänan}
  \pos{Intransitive verb}  \sense{
}}

\entry{täne}{\headword{täne}
  \variantof{Unspecified Variant of \vartext{tine}}}

\entry{tängg}{\headword{tängg}
  \pos{Transitive verb}  \sense{
    \definition{to urinate on, pee on}}}

\entry{tänggag}{\headword{tänggag}  \pronnote{təŋɡaɡ}
  \pos{Transitive verb}  \sense{\textbf{1.} 
    \definition{to make a dog more sensitive to smells by rubbing lemongrass on their nose}}  \sense{\textbf{2.} 
    \definition{to steal}    \example{
}}}

\entry{Tanisha}{\headword{Tanisha}
  \pos{Proper noun}  \sense{
    \definition{female personal name}}}

\entry{tanong}{\headword{tanong}  \pronnote{tanoŋ}
  \pos{Adverb}  \sense{
    \definition{a little}    \example{
      \exsen{utale kanong}      \extran{a bit far}}}
  \variants{\Dialectal Variant \vartext{kanong}}}

\entry{tanteny}{\headword{tanteny}  \pronnote{tanteɲ}
  \pos{Noun}  \sense{
    \definition{type of tree}    \example{
}    \example{
}    \example{
}}
  \variants{\Unspecified Variant \vartext{tanyteny}}}

\entry{tanträm}{\headword{tanträm}  \pronnote{tantrəm}
  \pos{Transitive verb}  \sense{
    \definition{to smash}}}

\entry{täny}{\headword{täny}  \pronnote{təɲ}
  \pos{Noun}  \sense{\textbf{1.} 
    \definition{lesser fig-parrot}    \scientificname{Cyclopsitta}    \example{
}    \example{
}    \example{
}}  \sense{\textbf{2.} 
    \definition{red-flanked lorikeet}    \scientificname{Hypocharmosyna placentis placentis}}}

\entry{tanyäb}{\headword{tanyäb}
  \variantof{Dialectal Variant of \vartext{tanyib}}}

\entry{tanyib}{\headword{tanyib}
  \pos{Noun}  \sense{
    \definition{radioulna (fused radius and ulna, or arm bone) of a flying fox}    \example{
      \exsen{Ngäma mäda bi giri alle topoll bo tanyib de dupoaemallo a ngäma mɨllɨng e dapoaemallo.}      \extran{Our (excl.) fathers used to sharpen the phalanges of a flying fox with a knife and pierce our noses.}}}
  \variants{\Dialectal Variant \vartext{tanyäb}}}

\entry{tanyteny}{\headword{tanyteny}
  \variantof{Unspecified Variant of \vartext{tanteny}}}

\entry{Tao}{\headword{Tao}
  \pos{Proper noun}  \sense{
    \definition{male personal name}}}

\entry{Taolang}{\headword{Taolang}  \note{Places around Limol}  \pronnote{taolaŋ}
  \pos{Proper noun}  \sense{
    \definition{Taolang (camping place)}    \note{Places around Limol}    \example{
}    \example{
}    \example{
}}}

\entry{taosen}{\headword{taosen}  \pronnote{taosen}
  \pos{Numeral}  \sense{
    \definition{thousand}    \example{
      \exsen{ttongo taosen tukituki}      \extran{more than one thousand}}}
  \variants{\SpellingVariant \vartext{thousand}}}

\entry{täp}{\headword{täp}  \pronnote{təp}
  \pos{Noun}  \sense{
    \definition{sago shoot}    \example{
      \exsen{Sana täp a tupi dan.}      \extran{The sago shoot is long.}}}}

\entry{tap}{\headword{tap}
  \pos{Noun}  \sense{
    \definition{tunnel}}}

\entry{täp}{\headword{täp}
  \pos{Verb}  \sense{
}}

\entry{tap}{\headword{tap}  \lexeme{tap}  \pronnote{tap}
  \pos{Noun}  \sense{
    \definition{type of big yam with a white interior, thorns, and hairs}}}

\entry{tap}{\headword{tap}
  \pos{Transitive verb}  \sense{
    \definition{to harvest}    \example{
      \exsen{Tapnen de gongkaemallo.}      \extran{They started the harvest.}}    \example{
      \exsen{Dätepän.}      \extran{He harvested it.}}}}

\entry{tap}{\headword{tap}  \lexeme{tap}  \pronnote{tap}
  \pos{Transitive coverb}  \sense{
    \definition{to dock, land}    \example{
      \exsen{Gall de tap dägageya wa ma we deyareya inu we.}      \extran{We docked the boat and returned home to sleep.}}}}

\entry{täp gazenatt}{\headword{täp gazenatt}  \pronnote{təp ɡazenaʈ͡ʂ}
  \pos{Noun}  \sense{
    \definition{the fourth stage of sago growth in which the shoot emerges, indicating the pith is ready to be harvested}    \example{
      \exsen{Masar tätäm täp gazenatt sana de nagda bälle dantepägän.}      \extran{Yesterday, grandfather showed a flowering sago to his friend.}}}}

\entry{täpäll}{\headword{täpäll}  \pronnote{təpəɽ}
  \pos{Noun}  \sense{
    \definition{type of pandanus used in a traditional hairstyle and woven together to make a big mat}}}

\entry{täpe}{\headword{täpe}
  \pos{Noun}  \sense{
    \definition{mud}}}

\entry{täpetäpe}{\headword{täpetäpe}
  \pos{Modifier}  \sense{
    \definition{muddy}    \example{
      \exsen{Llɨg a täpetäpe gognegnän yogoll me.}      \extran{The kids were getting muddy in the rain.}}}}

\entry{Tapila}{\headword{Tapila}  \pronnote{tapila}
  \pos{Proper noun}  \sense{
    \definition{Tapila (Makayam-speaking village in Gogodala Rural LLG; GPS: -8.414202, 143.016867)}    \example{
}    \example{
}    \example{
}}}

\entry{tapma}{\headword{tapma}
  \pos{Noun}  \sense{
    \definition{dock, wharf}    \example{
      \exsen{Ubi mängalae gall tapma we dagllaeyo.}      \extran{We (incl.) paddled quickly to the canoe dock.}}}}

\entry{Tapma}{\headword{Tapma}  \pronnote{tapma}
  \pos{Proper noun}  \sense{
    \definition{Tapma (toponym)}    \example{
}}}

\entry{tär}{\headword{tär}  \pronnote{tər}
  \pos{Noun}  \sense{\textbf{1.} 
    \definition{string, line}    \example{
      \exsen{Ngäna obo nying kollop tär de mɨnyi bäpitkaeneg.}      \extran{I will untie his sandal laces.}}}  \sense{\textbf{2.} 
    \definition{strap or handle of a bag/basket (nyäng)}}}

\entry{tära}{\headword{tära}  \pronnote{təra}
  \pos{Transitive coverb}  \sense{
    \definition{to finish}    \example{
      \exsen{Ematta kollbameny abal angttägalle ma we? Bäne tämamae tära nägnegalle.}      \extran{Why did you return home without any fish? You finished them all.}}}}

\entry{täräb}{\headword{täräb}  \pronnote{tərəb}
  \pos{Noun}  \sense{
    \definition{a dead person's belongings (put outside at a funeral until the feast ends)}}}

\entry{täräb ma}{\headword{täräb ma}  \note{Dead things}  \pronnote{tərəb ma}
  \pos{Noun}  \sense{
    \definition{funeral home}    \note{Dead things}    \example{
}    \example{
}    \example{
}}}

\entry{tärabol}{\headword{tärabol}
  \pos{Noun}  \sense{
    \definition{trouble}}}

\entry{täräk}{\headword{täräk}  \pronnote{tərək}
  \pos{Transitive coverb}  \sense{
    \definition{something is wrong, not able to move because something is inside, for example a spear in your foot (täräk nägagän)}}}

\entry{tärak}{\headword{tärak}  \pronnote{tərak}
  \pos{Intransitive verb}  \sense{\textbf{1.} 
    \definition{to go inside}    \example{
      \exsen{Bogo dindugän do kukiny mama de ikop dägagän dibaeya däbe ik e gotärakän do gotägolän.}      \extran{He ran to a pile of tall grass that he saw and went inside and hid there.}}}  \sense{\textbf{2.} 
    \definition{to put in}    \example{
      \exsen{Ngäna pakos de pop nyongo dae dätrak.}      \extran{I put the pakos spear through the hole.}}}
  \variants{\Fast Speech Variant \vartext{trak}}}

\entry{täräk}{\headword{täräk}
  \pos{Noun}  \sense{
    \definition{bundle}    \example{
      \exsen{Bina patme auli sana täräk yuwatt dag?}      \extran{How many bundles of cooked sago do you all have?}}}
  \variants{\Fast Speech Variant \vartext{träk}}}

\entry{tarakoll}{\headword{tarakoll}  \pronnote{tarakoɽ}
  \pos{Noun}  \sense{
    \definition{protective outer wall built by ancestors}}}

\entry{täral}{\headword{täral}  \note{Trees}  \pronnote{təral}
  \pos{Noun}  \sense{
    \definition{type of tree that grows in the grassland with white flowers, brown fruit, and wood used for house posts}    \note{Trees}    \example{
}    \example{
}    \example{
}}}

\entry{täral pällämpälläm}{\headword{täral pällämpälläm}  \note{Trees}  \pronnote{təral pəɽəmpəɽəm}
  \pos{Noun}  \sense{
    \definition{type of tree}    \note{Trees}    \example{
}    \example{
}    \example{
}}}

\entry{täräll}{\headword{täräll}
  \pos{Transitive coverb}  \sense{
    \definition{to stick out}    \example{
      \exsen{Kurupel llamda de dägmar täräll dägnegän.}      \extran{Kurupel stuck out his tongue at the old man.}}}}

\entry{täram}{\headword{täram}
  \variantof{Unspecified Variant of \vartext{tɨram}}}

\entry{täräm}{\headword{täräm}  \pronnote{tərəm}
  \pos{Noun}  \sense{
    \definition{tear}}
  \variants{\Fast Speech Variant \vartext{träm}}}

\entry{täram}{\headword{täram}  \pronnote{təram}
  \pos{Transitive verb}  \sense{
    \definition{to lead, take, carry, collect}    \example{
      \exsen{Kwakmae obom dätramän walle we.}      \extran{Kwakmae led him to the water.}}    \example{
      \exsen{Ine da obom dätramän do Dowan tutu.}      \extran{The water carried him to Dowan mountain.}}    \example{
      \exsen{Ubi bibim mɨnyi kwatt e yatraemneyo.}      \extran{They will take you all to court.}}    \example{
      \exsen{Polis a obom täram erallo daolle.}      \extran{The police are taking him away.}}}
  \variants{\Fast Speech Variant \vartext{tram}}}

\entry{tarambobo}{\headword{tarambobo}
  \pos{Noun}  \sense{
    \definition{hooded butcherbird}    \scientificname{Cracticus cassicus}}}

\entry{tärämpmeny}{\headword{tärämpmeny}  \pronnote{tərəmpmeɲ}
  \pos{Noun}  \sense{
    \definition{sleeping with a dead person's belongings}}}

\entry{täränga}{\headword{täränga}  \pronnote{tərəŋa}
  \pos{Modifier}  \sense{
    \definition{low, scant, shallow}    \example{
      \exsen{Ine kutt me ine da täränga dan.}      \extran{The water in the bucket is low.}}}}

\entry{tärängesa}{\headword{tärängesa}
  \variantof{Unspecified Variant of \vartext{tärangesa}}}

\entry{tärangesa}{\headword{tärangesa}  \note{Numbers}  \pronnote{təraŋesa}
  \pos{Noun}  \sense{\textbf{1.} 
    \definition{pinky, little finger}    \note{Numbers}    \example{
      \exsen{Ge tärängesa da sɨmell bod ik i gozenän.}      \extran{This pinky went in the pig's mouth.}}}  \sense{\textbf{2.} 
    \definition{one (lit. pinky; body counting numeral)}    \note{Numbers}    \example{
}    \example{
}    \example{
}}
  \variants{\Unspecified Variant \vartext{tärängesa, tɨrangesa}}}

\entry{tärangg}{\headword{tärangg}
  \pos{Transitive verb}  \sense{
    \definition{to stop, hold back}    \example{
      \exsen{Grace kollba kälsäre de daspunän walle we, be mälla da tämamae angde ikop dägageyo, ubi Grace bom dantärageyo ada, "Kollba de supun a mudan!"}      \extran{Grace threw the small fish back into the water, but when all the women saw her, they held her back, saying, "Don't throw back the fish!"}}}
  \variants{\Fast Speech Variant \vartext{trangg}}}

\entry{täräp}{\headword{täräp}  \pronnote{tərəp}
  \pos{Noun}  \sense{
    \definition{hunting tally (collection of jaw bones often displayed in the front of people's homes)}}}

\entry{täräp}{\headword{täräp}  \pronnote{tərəp}
  \pos{Noun}  \sense{
    \definition{time}    \example{
      \exsen{Ma we ibi täräp.}      \extran{It's time to go home.}}}}

\entry{täräp}{\headword{täräp}  \pronnote{tərəp}
  \pos{Transitive coverb}  \sense{
    \definition{to portion, share, split, divide}    \example{
      \exsen{Sɨmell de täräp dägagän.}      \extran{She portioned out the pig.}}}}

\entry{täräpän}{\headword{täräpän}
  \pos{Transitive verb}  \sense{
    \definition{to cut}}}

\entry{Tärapang}{\headword{Tärapang}
  \pos{Proper noun}  \sense{
    \definition{female personal name}}}

\entry{tarasoso}{\headword{tarasoso}  \note{Birds}  \pronnote{tarasoso}
  \pos{Noun}  \sense{
    \definition{type of bird}    \note{Birds}    \example{
}    \example{
}    \example{
}}}

\entry{täratäre}{\headword{täratäre}  \pronnote{təratəre}
  \pos{Transitive verb}  \sense{
    \definition{to dig out, hollow out}}}

\entry{tärbatt}{\headword{tärbatt}
  \pos{Noun}  \sense{
}}

\entry{täre}{\headword{täre}  \note{Everyday words and phrases}  \pronnote{təre}
  \pos{Noun}  \sense{\textbf{1.} 
    \definition{feast}    \note{Everyday words and phrases}    \example{
      \exsen{Lla da ttongo sande makäp me mɨnyi täre we gonsärbemamalle.}      \extran{People took one week to prepare for the feast.}}    \example{
}    \example{
}}  \sense{\textbf{2.} 
    \definition{holy, sacred}    \note{Everyday words and phrases}}}

\entry{täre buk}{\headword{täre buk}
  \pos{Noun}  \sense{
    \definition{Scripture}}}

\entry{täre ma}{\headword{täre ma}
  \pos{Noun}  \sense{
    \definition{temple}}}

\entry{tärke käp}{\headword{tärke käp}  \pronnote{tərke kəp}
  \pos{Noun}  \sense{
    \definition{necklace}}}

\entry{tarketarke}{\headword{tarketarke}  \pronnote{tarketarke}
  \pos{Modifier}  \sense{
    \definition{brittle}    \example{
      \exsen{Llo ttam a tarketarke dag.}      \extran{The leaves are brittle.}}}
  \variants{\Fast Speech Variant \vartext{tatarke}}}

\entry{tarko}{\headword{tarko}
  \pos{Noun}  \sense{
    \definition{female friends who have shared a joined fruit or vegetable}    \example{
      \exsen{Kathy bo tarko da Mome aenen.}      \extran{Kathy's tarko is Mome.}}}}

\entry{tärko}{\headword{tärko}  \pronnote{tərko}
  \pos{Noun}  \sense{
    \definition{type of small fish}}}

\entry{tarme}{\headword{tarme}  \pronnote{tarme}
  \pos{Noun}  \sense{
    \definition{kookaburra (blue-winged; spangled)}    \scientificname{Dacelo leachii intermedia; Dacelo tyro archboldi}    \example{
}    \example{
}    \example{
}}}

\entry{tarme ballmenyang}{\headword{tarme ballmenyang}  \note{Seasons}  \pronnote{tarme baɽmeɲaŋ}
  \pos{Noun}  \sense{
    \definition{season when crops are bearing fruit (fifth season; corresponds to March)}    \note{Seasons}    \example{
}    \example{
}    \example{
}}}

\entry{tarme koeme}{\headword{tarme koeme}  \note{Trees}  \pronnote{tarme koeme}
  \pos{Noun}  \sense{
    \definition{type of tree}    \note{Trees}    \example{
}    \example{
}    \example{
}}}

\entry{tarmekälla}{\headword{tarmekälla}  \note{Types of Taro}  \pronnote{tarmekəɽa}
  \pos{Noun}  \sense{
    \definition{type of taro}    \note{Types of Taro}    \example{
}    \example{
}    \example{
}}}

\entry{tärnetärne}{\headword{tärnetärne}  \note{Animal verbs}  \pronnote{tərnetərne}
  \pos{Verb}  \sense{
    \definition{to peck}    \note{Animal verbs}    \example{
      \exsen{tärnetärneang dan}      \extran{he is pecking}}    \example{
}    \example{
}}}

\entry{taromba}{\headword{taromba}  \note{Numbers}  \pronnote{taromba}
  \pos{Numeral}  \sense{
    \definition{216 (yam counting numeral; 6^3)}    \note{Numbers}    \example{
}    \example{
}    \example{
}}}

\entry{tärpa}{\headword{tärpa}  \pronnote{tərpa}
  \pos{Transitive coverb}  \sense{\textbf{1.} 
    \definition{to overcook, burn}}  \sense{\textbf{2.} 
    \definition{yam skin}}  \sense{\textbf{3.} 
    \definition{leftovers}}}

\entry{tärpae}{\headword{tärpae}  \note{Food}  \pronnote{tərpae}
  \pos{Noun}  \sense{
    \definition{type of yam}    \note{Food}    \example{
}    \example{
}    \example{
}}}

\entry{tärpam}{\headword{tärpam}
  \pos{Transitive verb}  \sense{\textbf{1.} 
    \definition{to put across}    \example{
      \exsen{Ubi yuwog abal llo de datärpaemaemeyo walle me.}      \extran{They put many sticks across the water.}}}  \sense{\textbf{2.} 
    \definition{to crucify}    \example{
      \exsen{Ubi Yesu bom datärpameyo.}      \extran{They crucified Jesus.}}}}

\entry{tärpamatt}{\headword{tärpamatt}
  \pos{Noun}  \sense{
    \definition{cross}    \example{
      \exsen{llo tärpamatt toko me}      \extran{on top of the wooden cross}}}}

\entry{tärpan}{\headword{tärpan}
  \pos{Verb}  \sense{
}}

\entry{tärpi}{\headword{tärpi}  \pronnote{tərpi}
  \pos{Modifier}  \sense{
    \definition{slippery, smooth}    \example{
      \exsen{Nyongo da tärpi dan.}      \extran{The road is slippery.}}}}

\entry{tärpitärpi}{\headword{tärpitärpi}  \note{Adjectives (Swadesh)}  \pronnote{tərpitərpi}
  \pos{Modifier}  \sense{
    \definition{nonsingular form of tärpi}    \note{Adjectives (Swadesh)}    \example{
      \exsen{Ngäna kumuddäga ttägäll käp tärpitärpi de nabllenegan.}      \extran{I found three smooth stones.}}    \example{
}    \example{
}}}

\entry{tärpoll}{\headword{tärpoll}  \pronnote{tərpoɽ}
  \pos{Noun}  \sense{
    \definition{piece}    \example{
      \exsen{llo palltapallta torpoll me}      \extran{on a flat piece of wood}}}
  \variants{\SpellingVariant \vartext{torpoll}}}

\entry{tät}{\headword{tät}  \pronnote{tət}
  \pos{Noun}  \sense{\textbf{1.} 
    \definition{stretcher}    \example{
      \exsen{Ngänäm tät alle datrameyo.}      \extran{They carried me with a stretcher.}}}  \sense{\textbf{2.} 
    \definition{ladder}}}

\entry{tät}{\headword{tät}  \pronnote{tət}
  \pos{Noun}  \sense{
    \definition{type of insect}}}

\entry{tata}{\headword{tata}  \pronnote{tata}
  \pos{Noun}  \sense{
    \definition{junction, intersection}    \example{
      \exsen{Bogo nyongo tata me dägabällän.}      \extran{She stood at the road junction.}}}}

\entry{tata}{\headword{tata}
  \variantof{Baby Talk Variant of \vartext{ddäddäg}}  \pronnote{tata}}

\entry{tätäb}{\headword{tätäb}
  \pos{Noun}  \sense{
}}

\entry{tätäk}{\headword{tätäk}
  \pos{Noun}  \sense{
}}

\entry{tätäm}{\headword{tätäm}  \pronnote{tətəm}
  \pos{Temporal adverb}  \sense{
    \definition{yesterday}}}

\entry{tätämall}{\headword{tätämall}
  \pos{Noun}  \sense{
    \definition{type of small insect that glow green}}}

\entry{tätän}{\headword{tätän}  \note{Fishing}  \pronnote{tətən}
  \pos{Noun}  \sense{\textbf{1.} 
    \definition{rib, side, flank}    \note{Fishing}    \example{
      \exsen{sɨmell tätän}      \extran{pork rib}}}  \sense{\textbf{2.} 
    \definition{animal trap}    \note{Fishing}    \example{
      \exsen{Käza da aittan Martin bo tätän me.}      \extran{The crocodile was caught in Martin's trap.}}    \example{
}    \example{
}}}

\entry{Tätän}{\headword{Tätän}  \pronnote{tətən}
  \pos{Proper noun}  \sense{
    \definition{unisex personal name}}}

\entry{tätän kutt}{\headword{tätän kutt}  \note{Parts of the Body}  \pronnote{tətən kuʈ͡ʂ}
  \pos{Noun}  \sense{
    \definition{rib (bone)}    \note{Parts of the Body}    \example{
}    \example{
}    \example{
}}}

\entry{tätang}{\headword{tätang}  \pronnote{tətaŋ}
  \pos{Modifier}  \sense{
    \definition{rotten}}}

\entry{tatanggli}{\headword{tatanggli}  \pronnote{tataŋɡli}
  \pos{Noun}  \sense{
    \definition{willie wagtail}    \scientificname{Rhipidura leucophrys melaleuca}    \example{
}    \example{
}    \example{
}}}

\entry{tätäp}{\headword{tätäp}  \pronnote{tətəp}
  \pos{Noun}  \sense{
    \definition{pain}    \example{
      \exsen{Bogo tätäp de umllang gogän.}      \extran{He felt pain.}}    \example{
}    \example{
}}}

\entry{tatäraem}{\headword{tatäraem}  \pronnote{tatəraem}
  \pos{Noun}  \sense{
    \definition{noise}    \example{
      \exsen{Tatraem de gondärän.}      \extran{She heard the noise.}}    \example{
      \exsen{De dräm de aeya tatäraem eran?}      \extran{Who's making noise in the drum?}}}
  \variants{\Fast Speech Variant \vartext{tatraem}}}

\entry{tätäräk}{\headword{tätäräk}
  \variantof{SpellingVariant of \vartext{täträk}}}

\entry{tätärangae}{\headword{tätärangae}  \pronnote{tətəraŋae}
  \pos{Adverb}  \sense{
    \definition{following around}}}

\entry{tätäräp}{\headword{tätäräp}  \pronnote{tətərəp}
  \pos{Transitive verb}  \sense{\textbf{1.} 
    \definition{to cut}    \example{
      \exsen{Ibi goeg tatäräp e.}      \extran{Let's go to the garden to cut down trees.}}    \example{
      \exsen{Ngäna ngämo llan de yatärpänan.}      \extran{I cut my ears.}}}  \sense{\textbf{2.} 
    \definition{to cut across, take a shortcut}    \example{
      \exsen{Minong gänyaollemae ditäräpmällnän ngämo pate.}      \extran{Minong took a shortcut this way towards me.}}}  \sense{\textbf{3.} 
    \definition{man-made clearing}}
  \variants{\Fast Speech Variant \vartext{täträp}}
  \variants{\Unspecified Variant \vartext{tatäräp}}}

\entry{tatäräp}{\headword{tatäräp}
  \variantof{Unspecified Variant of \vartext{tätäräp}}  \pronnote{tatərəp}}

\entry{tätäräp}{\headword{tätäräp}
  \pos{Noun}  \sense{
    \definition{heat}    \example{
      \exsen{Ngäna yäbäd tätäräp atta gäba we inggol allan.}      \extran{I am moving to the shade because of the sun's heat.}}    \example{
      \exsen{Yäbäd a ebdo kälakälae tätäräpang dan, be iddob ag mer dan zanggae e.}      \extran{The sun is a bit hot during the day, but the night or morning are good for roaming.}}}
  \variants{\Fast Speech Variant \vartext{täträp}}}

\entry{tatarke}{\headword{tatarke}
  \variantof{Fast Speech Variant of \vartext{tarketarke}}}

\entry{tätärpeyam}{\headword{tätärpeyam}  \pronnote{tətərpejam}
  \pos{Noun}  \sense{
    \definition{type of grass (~1 m) used as a toy}}}

\entry{tater}{\headword{tater}  \note{Things a person can do}  \pronnote{tater}
  \pos{Noun}  \sense{\textbf{1.} 
    \definition{mat}    \note{Things a person can do}    \example{
      \exsen{Gaem ai tater gollaeb de inen anggan.}      \extran{Gaem is weaving many mats.}}    \example{
}}  \sense{\textbf{2.} 
    \definition{to spread out}    \note{Things a person can do}    \example{
      \exsen{Ubi oba iddpo de däternegeyo dongki bo ddäg toko me.}      \extran{They spread out their cloaks over the donkey's back.}}}}

\entry{tätkea}{\headword{tätkea}  \note{Sago palms}  \pronnote{tətkea}
  \pos{Noun}  \sense{
    \definition{type of sago}    \note{Sago palms}    \example{
}    \example{
}    \example{
}}}

\entry{tätmar}{\headword{tätmar}
  \pos{Noun}  \sense{
}}

\entry{tatraem}{\headword{tatraem}
  \variantof{Fast Speech Variant of \vartext{tatäraem}}}

\entry{täträk}{\headword{täträk}  \pronnote{tətrək}
  \pos{Transitive verb}  \sense{\textbf{1.} 
    \definition{to go underneath}    \example{
      \exsen{Käsre dibaballe ada gogon goträkän.}      \extran{Then from there, he went underneath.}}}  \sense{\textbf{2.} 
    \definition{to push in, push through}    \example{
      \exsen{Ngämo giri de käza bo ume ttäp e däträkne.}      \extran{I was pushing in my knife into the crocodile's mouth.}}}
  \variants{\SpellingVariant \vartext{tɨtɨrɨk, tätäräk}}}

\entry{täträp}{\headword{täträp}
  \variantof{Fast Speech Variant of \vartext{tätäräp}}  \note{Weather}  \pronnote{tətrəp}}

\entry{täträp}{\headword{täträp}  \note{Seasons}  \pronnote{tətrəp}
  \pos{Noun}  \sense{
    \definition{season of the end of harvesting and the start of hunting (eighth season; corresponds to late May)}    \note{Seasons}    \example{
}    \example{
}    \example{
}}}

\entry{tatratatraema}{\headword{tatratatraema}  \note{tatratatraem or tatratatraem=ma?}  \pronnote{tatratatraema}
  \pos{Noun}  \sense{
    \definition{music}    \note{tatratatraem or tatratatraem=ma?}    \example{
}    \example{
}    \example{
}}}

\entry{tatruk}{\headword{tatruk}  \note{Poison}  \pronnote{tatruk}
  \pos{Noun}  \sense{
    \definition{type of poisonous ant}    \note{Poison}    \example{
}    \example{
}    \example{
}}}

\entry{tatu}{\headword{tatu}  \lexeme{tatu}  \pronnote{tatu}
  \pos{Transitive/Intransitive coverb}  \sense{
    \definition{to bathe, wash oneself; wash (an animate object)}    \example{
      \exsen{Ubi obaoba tatu anggan.}      \extran{They are washing each other.}}    \example{
      \exsen{Bogo tatu agan.}      \extran{She bathed.}}}}

\entry{tatuma}{\headword{tatuma}  \lexeme{tatuma}  \pronnote{tatuma}
  \pos{Noun}  \sense{
    \definition{washing place, outdoor bathing area}}}

\entry{tatutatu}{\headword{tatutatu}
  \pos{Adverb}  \sense{
    \definition{cleanly}    \example{
      \exsen{tatutatu giddollnen}      \extran{living cleanly}}}}

\entry{tawa}{\headword{tawa}  \note{Water}  \pronnote{tawa}
  \pos{Noun}  \sense{
    \definition{swamp}    \note{Water}    \example{
}    \example{
}    \example{
}}}

\entry{tawa aeb}{\headword{tawa aeb}  \pronnote{tawa aeb}
  \pos{Noun}  \sense{
    \definition{Australasian swamphen}    \scientificname{Porphyrio melanotus}}}

\entry{Tawabo}{\headword{Tawabo}
  \pos{Proper noun}  \sense{
    \definition{Tawabo (toponym)}}}

\entry{tawar}{\headword{tawar}
  \pos{Noun}  \sense{
    \definition{totem symbol (thing that represents a clan)}}}

\entry{tawe}{\headword{tawe}  \note{Types of palms}  \pronnote{tawe}
  \pos{Noun}  \sense{
    \definition{type of slim, tall palm with coconuts that come in red or green varieties, bunches of yellow fruit that birds eat, and fronds used for camp flooring}    \note{Types of palms}    \example{
}    \example{
}    \example{
}}}

\entry{tawe ttäp}{\headword{tawe ttäp}  \note{Hunting}  \pronnote{tawe ʈ͡ʂəp}
  \pos{Noun}  \sense{
    \definition{type of spear}    \note{Hunting}    \example{
      \exsen{Ngäna tätäm kukullang me ttongo ttall de tawe ttäp alle dazu.}      \extran{Yesterday, I shot one wallaby in the grassfire with the tawe ttäp arrow.}}    \example{
}    \example{
}}}

\entry{tawekutt}{\headword{tawekutt}  \pronnote{tawekuʈ͡ʂ}
  \pos{Noun}  \sense{\textbf{1.} 
    \definition{partially dry coconut}    \example{
      \exsen{Sana peyang nyänggae ma nge da imomdae gänyan tawekutt a.}      \extran{The right coconut to mix with sago is this tawekutt.}}}  \sense{\textbf{2.} 
    \definition{the tenth stage of coconut growth during which the fruit begins to dry and brown}    \example{
}}}

\entry{Tawemitang}{\headword{Tawemitang}  \note{Places around Limol}  \pronnote{tawemitaŋ}
  \pos{Proper noun}  \sense{
    \definition{Tawemitang (toponym)}    \note{Places around Limol}    \example{
}    \example{
}    \example{
}}}

\entry{Tayi}{\headword{Tayi}  \pronnote{taji}
  \pos{Proper noun}  \sense{
    \definition{Tai (in Gogodala Rural LLG)}    \example{
}    \example{
}    \example{
}}
  \variants{\SpellingVariant \vartext{Tai}}}

\entry{teacher}{\headword{teacher}
  \variantof{Dialectal Variant of \vartext{English loanword}}}

\entry{Teachers}{\headword{Teachers}
  \variantof{freevarlab of \vartext{English loanword}}}

\entry{teb}{\headword{teb}  \note{Parts of a reptile}  \pronnote{teb}
  \pos{Noun}  \sense{
    \definition{mandible, jawbone}    \note{Parts of a reptile}    \example{
}    \example{
}    \example{
}}}

\entry{Tebar}{\headword{Tebar}
  \pos{Proper noun}  \sense{
    \definition{Tebar (toponym)}}}

\entry{Teks}{\headword{Teks}
  \pos{Proper noun}  \sense{
    \definition{male personal name}}}

\entry{teks}{\headword{teks}
  \pos{Noun}  \sense{
    \definition{tax}    \example{
      \exsen{teks mani ngällbänanang lla}      \extran{tax collector}}}}

\entry{Teme}{\headword{Teme}
  \variantof{Dialectal Variant of \vartext{Tame}}}

\entry{ten}{\headword{ten}  \note{Numbers}  \pronnote{ten}
  \pos{Numeral}  \sense{
    \definition{ten (English numeral, also general)}    \note{Numbers}    \example{
}    \example{
}    \example{
}}}

\entry{tentatente}{\headword{tentatente}  \note{Trees}  \pronnote{tentatente}
  \pos{Noun}  \sense{
    \definition{category of parasitic plants}    \note{Trees}    \example{
}    \example{
}    \example{
}}}

\entry{tep}{\headword{tep}  \pronnote{tep}
  \pos{Noun}  \sense{
    \definition{tree sap (used in making drums and arrows)}    \example{
}    \example{
}    \example{
}}}

\entry{tep}{\headword{tep}  \pronnote{tep}
  \pos{Transitive coverb}  \sense{
    \definition{to lie}}}

\entry{Tergo}{\headword{Tergo}  \pronnote{terɡo}
  \pos{Proper noun}  \sense{
    \definition{male personal name}}}

\entry{tergony}{\headword{tergony}
  \pos{Transitive verb}  \sense{
    \definition{to spread, unfold}    \example{
      \exsen{Lla da tater de mɨnyi bitergonyeyo.}      \extran{The people will spread mats.}}}}

\entry{Terrance}{\headword{Terrance}  \pronnote{terans}
  \pos{Proper noun}  \sense{
    \definition{male personal name}}}

\entry{tesde}{\headword{tesde}  \note{Part of the week}  \pronnote{tesde}
  \pos{Noun}  \sense{
    \definition{Thursday}    \note{Part of the week}    \example{
}    \example{
}    \example{
}}
  \variants{\SpellingVariant \vartext{teste}}}

\entry{teste}{\headword{teste}
  \variantof{SpellingVariant of \vartext{tesde}}}

\entry{teti}{\headword{teti}
  \pos{Numeral}  \sense{
    \definition{thirty}}
  \variants{\SpellingVariant \vartext{thirty}}}

\entry{tetin}{\headword{tetin}
  \pos{Numeral}  \sense{
    \definition{thirteen (English numeral)}}}

\entry{Tewa}{\headword{Tewa}  \pronnote{tewa}
  \pos{Proper noun}  \sense{
    \definition{male personal name}}}

\entry{Tewara}{\headword{Tewara}  \pronnote{tewara}
  \pos{Proper noun}  \sense{
    \definition{Tewara (Bitur-speaking village in Oriomo-Bitur Rural LLG)}    \example{
}    \example{
}    \example{
}}}

\entry{teweya}{\headword{teweya}  \pronnote{teweja}
  \pos{Noun}  \sense{
    \definition{brush cuckoo}    \scientificname{Cacomantis variolosus}}}

\entry{Teyapopo}{\headword{Teyapopo}  \note{Place names}  \pronnote{tejapopo}
  \pos{Proper noun}  \sense{
    \definition{Teyapopo (toponym)}    \note{Place names}    \example{
}    \example{
}    \example{
}}}

\entry{teyateyar}{\headword{teyateyar}  \pronnote{tejatejar}
  \pos{Modifier}  \sense{
    \definition{shallow}}}

\entry{the}{\headword{the}
  \variantof{freevarlab of \vartext{English loanword}}}

\entry{there}{\headword{there}
  \variantof{Dialectal Variant of \vartext{English loanword}}}

\entry{thirty}{\headword{thirty}
  \variantof{SpellingVariant of \vartext{teti}}}

\entry{this}{\headword{this}
  \variantof{Dialectal Variant of \vartext{English loanword}}}

\entry{Thomas}{\headword{Thomas}
  \variantof{SpellingVariant of \vartext{Tomas}}  \pronnote{tomas}}

\entry{thousand}{\headword{thousand}
  \variantof{SpellingVariant of \vartext{taosen}}}

\entry{three}{\headword{three}
  \variantof{SpellingVariant of \vartext{tri}}}

\entry{tibekllop}{\headword{tibekllop}  \pronnote{tibekɽop}
  \pos{Noun}  \sense{
    \definition{type of medicine}    \example{
}    \example{
}    \example{
}}}

\entry{tibi}{\headword{tibi}
  \pos{Noun}  \sense{
    \definition{tuberculosis}}}

\entry{tibra}{\headword{tibra}  \note{Names of Bananas}  \pronnote{tibra}
  \pos{Noun}  \sense{
    \definition{type of native banana}    \note{Names of Bananas}    \example{
}    \example{
}    \example{
}}}

\entry{tigi}{\headword{tigi}  \note{Snakes}  \pronnote{tiɡi}
  \pos{Noun}  \sense{
    \definition{type of snake}    \note{Snakes}    \example{
}    \example{
}    \example{
}}}

\entry{tikle}{\headword{tikle}  \pronnote{tikle}
  \pos{Noun}  \sense{
    \definition{type of small louse}    \example{
}    \example{
}    \example{
}}}

\entry{tiklem}{\headword{tiklem}  \note{Animal names}  \pronnote{tiklem}
  \pos{Noun}  \sense{
    \definition{type of animal}    \note{Animal names}    \example{
}    \example{
}    \example{
}}}

\entry{tikop}{\headword{tikop}  \pronnote{tikop}
  \pos{Noun}  \sense{\textbf{1.} 
    \definition{heart}    \example{
}}  \sense{\textbf{2.} 
    \definition{beloved, sweetheart, dear}    \example{
}}}

\entry{tikopmeny}{\headword{tikopmeny}  \pronnote{tikopmeɲ}
  \pos{Modifier}  \sense{
    \definition{startled, scared}    \example{
      \exsen{Lla da ttongo lla de aiyei abal tikopmeny nägagan.}      \extran{The man really scared the other man.}}    \example{
      \exsen{Madura tikopmeny gogän.}      \extran{Madura was startled.}}    \example{
}}}

\entry{tikuku}{\headword{tikuku}  \pronnote{tikuku}
  \pos{Noun}  \sense{\textbf{1.} 
    \definition{black-backed bittern}    \scientificname{Ixobrychus dubius}    \example{
}    \example{
}    \example{
}}  \sense{\textbf{2.} 
    \definition{pheasant coucal}    \scientificname{Centropus phasianinus}}}

\entry{Tim}{\headword{Tim}
  \pos{Proper noun}  \sense{
    \definition{female personal name}}}

\entry{timän}{\headword{timän}  \pronnote{timən}
  \pos{Transitive verb}  \sense{\textbf{1.} 
    \definition{to release, fire, shoot (an arrow)}    \example{
      \exsen{Ibi mɨnyi toboll de bitimänmällnegalla.}      \extran{We (incl.) will be firing arrows.}}}  \sense{\textbf{2.} 
    \definition{end}    \example{
      \exsen{Ekaklle ulle bo timän a dowae abal agan zime.}      \extran{The end of the world is near.}}}}

\entry{Tina}{\headword{Tina}  \pronnote{tina}
  \pos{Proper noun}  \sense{
    \definition{female personal name}}}

\entry{tine}{\headword{tine}  \pronnote{tine}
  \pos{Noun}  \sense{\textbf{1.} 
    \definition{sago leaf}}  \sense{\textbf{2.} 
    \definition{type of simple roof consisting of a single post covered in sago leaves}}
  \variants{\Unspecified Variant \vartext{täne, tɨne}}}

\entry{tintɨrmoll}{\headword{tintɨrmoll}
  \variantof{SpellingVariant of \vartext{tintromol}}}

\entry{tintromol}{\headword{tintromol}  \note{Poison}  \pronnote{tintromol}
  \pos{Noun}  \sense{
    \definition{type of black or red ants}    \note{Poison}    \example{
}    \example{
}    \example{
}}
  \variants{\SpellingVariant \vartext{tintɨrmoll}}}

\entry{tintromoltintromol}{\headword{tintromoltintromol}  \note{Games}  \pronnote{tintromoltintromol}
  \pos{Noun}  \sense{
    \definition{type of rhyming game}    \note{Games}    \example{
}    \example{
}    \example{
}}}

\entry{Tirere}{\headword{Tirere}  \note{Place names}  \pronnote{tirere}
  \pos{Proper noun}  \sense{
    \definition{Tirere/Tire'ere (Waboda-speaking village in Kiwai Rural LLG)}    \note{Place names}    \example{
}    \example{
}    \example{
}}}

\entry{tirit}{\headword{tirit}
  \pos{Noun}  \sense{
}}

\entry{Titi}{\headword{Titi}  \note{Places around Limol}  \pronnote{titi}
  \pos{Proper noun}  \sense{
    \definition{Titi (toponym)}    \note{Places around Limol}    \example{
}    \example{
}    \example{
}}}

\entry{titi}{\headword{titi}  \pronnote{titi}
  \pos{Noun}  \sense{
    \definition{bat}    \example{
      \exsen{Titi da ddobae sawasap otnan kämang dag.}      \extran{Bats eat lots of sawasap fruit.}}}}

\entry{titi}{\headword{titi}  \pronnote{titi}
  \pos{Noun}  \sense{\textbf{1.} 
    \definition{brown honeyeater}    \scientificname{Lichmera indistincta}    \example{
}    \example{
}    \example{
}}  \sense{\textbf{2.} 
    \definition{brown-backed honeyeater}    \scientificname{Ramsayornis modestus}}  \sense{\textbf{3.} 
    \definition{yellow-bellied longbill}    \scientificname{Toxorhamphus novaeguineae}}  \sense{\textbf{4.} 
    \definition{rufous-banded honeyeater}    \scientificname{Conopophila albogularis}}  \sense{\textbf{5.} 
    \definition{pygmy longbill}    \scientificname{Oedistoma pygmaeum}}}

\entry{Titus}{\headword{Titus}  \pronnote{titus}
  \pos{Proper noun}  \sense{
    \definition{male personal name}}}

\entry{Tizag}{\headword{Tizag}  \note{Diving}  \pronnote{tizaɡ}
  \pos{Proper noun}  \sense{
    \definition{Ende dialect}    \note{Diving}    \example{
}    \example{
}    \example{
}}
  \variants{\SpellingVariant \vartext{tizag}}}

\entry{tizag}{\headword{tizag}
  \variantof{SpellingVariant of \vartext{Tizag}}}

\entry{tɨke}{\headword{tɨke}
  \pos{Noun}  \sense{
}}

\entry{tɨn}{\headword{tɨn}  \pronnote{tɪn}
  \pos{Noun}  \sense{
    \definition{steam}}}

\entry{tɨne}{\headword{tɨne}
  \variantof{Unspecified Variant of \vartext{tine}}  \pronnote{tɪne}}

\entry{tɨram}{\headword{tɨram}  \pronnote{tɪram}
  \pos{Transitive verb}  \sense{
    \definition{to open}    \example{
      \exsen{Ge ud a täramang daeya, be sisri papekatt dan.}      \extran{This door was open, but now it's closed.}}}
  \variants{\Unspecified Variant \vartext{täram}}}

\entry{tɨrangesa}{\headword{tɨrangesa}
  \variantof{Unspecified Variant of \vartext{tärangesa}}}

\entry{tɨt}{\headword{tɨt}  \pronnote{tɪt}
  \pos{Noun}  \sense{
    \definition{type of tree found in the bush}}}

\entry{tɨt}{\headword{tɨt}  \pronnote{tɪt}
  \pos{Transitive verb}  \sense{\textbf{1.} 
    \definition{to beat sago, pound sago}    \example{
      \exsen{Ubi sana pätt att de tɨt erallo a kängkäm erallo.}      \extran{They pound the pith of the sago palm and squeeze it.}}    \example{
      \exsen{Bogo mɨnyi sana bɨtɨn.}      \extran{He will pound sago.}}}  \sense{\textbf{2.} 
    \definition{to hollow out, dig out}    \example{
      \exsen{Ngäna gall de tɨt iran.}      \extran{I am digging out a canoe.}}    \example{
      \exsen{Bäne turik de abo iweny gall tɨt i.}      \extran{Bring your axe to dig the canoe.}}}}

\entry{tɨtɨp}{\headword{tɨtɨp}
  \pos{Transitive verb}  \sense{
    \definition{to braid}    \example{
      \exsen{Bongo bäne bun de mermerae nätɨpnegalle.}      \extran{You braided your hair nicely.}}}}

\entry{tɨtɨrɨk}{\headword{tɨtɨrɨk}
  \variantof{SpellingVariant of \vartext{täträk}}  \pronnote{tɪtɪrɪk}}

\entry{tlläpmälltlläpmäll}{\headword{tlläpmälltlläpmäll}  \pronnote{tɽəpməɽtɽəpməɽ}
  \pos{Adverb}  \sense{
    \definition{nibbling}}}

\entry{to}{\headword{to}  \pronnote{to}
  \pos{Noun}  \sense{
    \definition{light}    \example{
      \exsen{Ge to ulle ngäna ikop däga ereya, ngäna lelang atta dindug.}      \extran{I saw this big light and ran out of fright.}}}}

\entry{to~}{\headword{to~}
  \variantof{Plural reduplicant of \vartext{RED~}}}

\entry{toatoall llɨg}{\headword{toatoall llɨg}  \pronnote{toatoaɽ ɽɪɡ}
  \pos{}  \sense{
    \definition{abandoned child}}}

\entry{toball}{\headword{toball}
  \variantof{SpellingVariant of \vartext{tobäll}}  \pronnote{tobaɽ}}

\entry{tobäll}{\headword{tobäll}  \note{Tools}  \pronnote{tobəɽ}
  \pos{Noun}  \sense{
    \definition{long spear shot like an arrow}    \note{Tools}    \example{
}    \example{
}    \example{
}}
  \variants{\SpellingVariant \vartext{toball, täbäll}}
  \variants{\Dialectal Variant \vartext{toboll}}}

\entry{tobäll abal}{\headword{tobäll abal}  \note{Traditional Spears}  \pronnote{tobəɽ abal}
  \pos{Noun}  \sense{
    \definition{type of spear}    \note{Traditional Spears}    \example{
}    \example{
}    \example{
}}}

\entry{tobäll käp}{\headword{tobäll käp}  \note{Tools}  \pronnote{tobəɽ kəp}
  \pos{Noun}  \sense{\textbf{1.} 
    \definition{arrowhead}    \note{Tools}}  \sense{\textbf{2.} 
    \definition{bullet}    \note{Tools}    \example{
}    \example{
}    \example{
}}}

\entry{toboll}{\headword{toboll}
  \variantof{Dialectal Variant of \vartext{tobäll}}  \pronnote{toboɽ}}

\entry{toe}{\headword{toe}  \note{Trees}  \pronnote{toe}
  \pos{Noun}  \sense{
    \definition{type of tree that grows in the bush with white flowers, black fruit, small edible nuts that doves and cassowaries eat, and pith that is steamed or extracted}    \note{Trees}    \example{
}    \example{
}    \example{
}}}

\entry{toelet}{\headword{toelet}
  \pos{Noun}  \sense{
    \definition{toilet}}
  \variants{\SpellingVariant \vartext{toilet}}}

\entry{toengg}{\headword{toengg}  \note{Water}  \pronnote{toeŋɡ}
  \pos{Noun}  \sense{
    \definition{confluence, junction (of a river)}    \note{Water}    \example{
}    \example{
}    \example{
}}}

\entry{Togllaema}{\headword{Togllaema}  \pronnote{toɡɽaema}
  \pos{Proper noun}  \sense{
    \definition{Togllaema (toponym)}}}

\entry{togo~}{\headword{togo~}
  \variantof{Dialectal Variant of \vartext{RED~}}}

\entry{togol}{\headword{togol}  \pronnote{toɡol}
  \pos{Transitive verb}  \sense{\textbf{1.} 
    \definition{to hide; go away to have sex secretly}    \example{
      \exsen{Ngäna ngattong botägol.}      \extran{I will hide first.}}    \example{
      \exsen{Bongo zime atägolalle?}      \extran{Are you already hiding?}}}  \sense{\textbf{2.} 
    \definition{to hide}    \example{
      \exsen{Ngämim deyatogolmällneyo ngäma mäg.}      \extran{Our (excl.) mothers were hiding us.}}}}

\entry{togotogol}{\headword{togotogol}  \pronnote{toɡotoɡol}
  \pos{Noun}  \sense{
    \definition{hide-and-seek}    \example{
      \exsen{Togotogol tongoe de ubi dangesneyo oba mokwang tongoe.}      \extran{They played hide-and-seek, their favorite game.}}    \example{
      \exsen{Ibi togotogol botongoenallaǃ}      \extran{Let's play hide and seek!}}}}

\entry{Togowa}{\headword{Togowa}  \pronnote{toɡowa}
  \pos{Proper noun}  \sense{
    \definition{Togowa (toponym)}    \example{
}    \example{
}    \example{
}}}

\entry{toilet}{\headword{toilet}
  \variantof{SpellingVariant of \vartext{toelet}}}

\entry{tok}{\headword{tok}  \pronnote{tok}
  \pos{Transitive coverb}  \sense{
    \definition{to wrap}    \example{
      \exsen{Bogo sɨmell midd di tok digewän llo ttam alle.}      \extran{He wraps the pork with leaves.}}}}

\entry{Tok pisin}{\headword{Tok pisin}  \note{Language Names}  \pronnote{tok pisin}
  \pos{Proper noun}  \sense{
    \definition{Tok Pisin}    \note{Language Names}    \example{
}    \example{
}    \example{
}}
  \variants{\Unspecified Variant \vartext{Tok pizin}}}

\entry{Tok pizin}{\headword{Tok pizin}
  \variantof{Unspecified Variant of \vartext{Tok pisin}}}

\entry{toko}{\headword{toko}  \pronnote{toko}
  \pos{Locational}  \sense{
    \definition{top}    \example{
      \exsen{Bogo llo toko me gotarän.}      \extran{She slept in the treetop.}}}}

\entry{tokom}{\headword{tokom}
  \pos{Noun}  \sense{
}}

\entry{tokong}{\headword{tokong}  \pronnote{tokoŋ}
  \pos{Noun}  \sense{
    \definition{bait}    \example{
      \exsen{Mälla da kok de nägäddaeballo kollba tokong e.}      \extran{The women caught grasshoppers for bait.}}}}

\entry{tokop}{\headword{tokop}  \note{Trees}  \pronnote{tokop}
  \pos{Noun}  \sense{\textbf{1.} 
    \definition{type of tree that grows in the bush; used as house sticks and medicine}    \note{Trees}    \example{
}    \example{
}    \example{
}}  \sense{\textbf{2.} 
    \definition{lump on skin}    \note{Trees}}}

\entry{toma}{\headword{toma}  \note{Parts of an insect}  \pronnote{toma}
  \pos{Noun}  \sense{
    \definition{wing}    \note{Parts of an insect}    \example{
}    \example{
}    \example{
}}}

\entry{toma mit}{\headword{toma mit}  \note{Parts of an insect}  \pronnote{toma mit}
  \pos{Noun}  \sense{
    \definition{metathorax}    \note{Parts of an insect}    \example{
}    \example{
}    \example{
}}}

\entry{tomäll}{\headword{tomäll}  \pronnote{toməɽ}
  \pos{Noun}  \sense{
    \definition{wart; fungal skin infection}}}

\entry{Tomas}{\headword{Tomas}
  \pos{Proper noun}  \sense{
    \definition{male personal name}}
  \variants{\SpellingVariant \vartext{Thomas}}}

\entry{tomato}{\headword{tomato}
  \pos{Noun}  \sense{
    \definition{tomato}}}

\entry{Tomato}{\headword{Tomato}  \pronnote{tomato}
  \pos{Proper noun}  \sense{
    \definition{female personal name}}}

\entry{tomoang}{\headword{tomoang}
  \pos{Noun}  \sense{
    \definition{tree liquid}}}

\entry{tomon}{\headword{tomon}  \pronnote{tomon}
  \pos{Transitive verb}  \sense{\textbf{1.} 
    \definition{to wait}    \example{
      \exsen{Ngäna kälakälae gontämon do yäbäd a angde goklowän dam ngäna ibi de gongkam.}      \extran{I waited a little until the sun set; then I started walking.}}}  \sense{\textbf{2.} 
    \definition{to wait for, await}    \example{
      \exsen{Ngäna mänmän e bantmonneg.}      \extran{I'll wait for the girls.}}    \example{
      \exsen{Bogo obom deyantämonän.}      \extran{She awaited him.}}}}

\entry{tomowang}{\headword{tomowang}  \pronnote{tomowaŋ}
  \pos{Modifier}  \sense{
    \definition{sour; bitter}    \example{
      \exsen{Ge wup bo moko da kälakälae tomowang dan.}      \extran{This banana's taste is a little sour.}}}}

\entry{Tomson}{\headword{Tomson}
  \pos{Proper noun}  \sense{
    \definition{male personal name}}}

\entry{tonang}{\headword{tonang}
  \pos{Modifier}  \sense{
    \definition{careful, cautious}    \example{
      \exsen{Bogo tonangae damllaemnegän.}      \extran{She held on carefully.}}    \example{
      \exsen{Abo tonangae agaebne ubira.}      \extran{You all must be wary of them.}}}}

\entry{tonangtonang}{\headword{tonangtonang}
  \pos{}  \sense{
    \definition{carefully, cautiously}    \example{
      \exsen{Ngäna mɨnyi yu de tonangtonang bikom.}      \extran{I'll bring the wood carefully.}}}}

\entry{tongg}{\headword{tongg}
  \pos{Transitive verb}  \sense{
    \definition{to point}    \example{
      \exsen{Liseng ttang dontogän Rose bom.}      \extran{Liseng pointed her finger at Rose.}}}}

\entry{Tonggar}{\headword{Tonggar}
  \pos{Proper noun}  \sense{
    \definition{ethnic group}}}

\entry{tonggo}{\headword{tonggo}  \lexeme{tonggo}  \pronnote{toŋɡo}
  \pos{Noun}  \sense{
    \definition{type of small bamboo that grows in the bush along creeks with a red interior; sharp when split and used as a cutting tool}}}

\entry{tongle}{\headword{tongle}  \note{?}  \pronnote{toŋle}
  \pos{Noun}  \sense{
    \definition{leech}    \note{?}    \example{
}    \example{
}    \example{
}}}

\entry{tongoe}{\headword{tongoe}  \note{Verbs of Motion}  \pronnote{toŋoe}
  \pos{Transitive verb}  \sense{\textbf{1.} 
    \definition{to play}    \note{Verbs of Motion}    \example{
      \exsen{Nga ubi dag tongoenen meae anggan.}      \extran{They are still playing.}}    \example{
      \exsen{Llɨg kälekäle da ebdo ulle atongoenegnan.}      \extran{The children were playing all day.}}    \example{
      \exsen{Ibi togotogol botongoe nalla.}      \extran{Let's play hide and seek.}}}  \sense{\textbf{2.} 
    \definition{to laugh (at)}    \note{Verbs of Motion}    \example{
      \exsen{Ubi ttongo lla de tongoenen erallo.}      \extran{They laugh at someone.}}    \example{
      \exsen{Lama käsre ngämira llɨtɨt dängkamän tongoetongoeang.}      \extran{Lama started to tell us (excl.) the story, laughing.}}}  \sense{\textbf{3.} 
    \definition{game; sport}    \note{Verbs of Motion}    \example{
      \exsen{Ngämo mokowang tongoe a basketbol dan.}      \extran{My favorite game is basketball.}}}}

\entry{tongoe ma ngätt}{\headword{tongoe ma ngätt}  \note{School}  \pronnote{toŋoe ma ŋəʈ͡ʂ}
  \pos{Noun}  \sense{
    \definition{playground}    \note{School}    \example{
}    \example{
}    \example{
}}}

\entry{tongoeang}{\headword{tongoeang}  \note{Oral Traditions}  \pronnote{toŋoeaŋ}
  \pos{Modifier}  \sense{
    \definition{funny, humorous}    \note{Oral Traditions}    \example{
      \exsen{tongoeang ttoen}      \extran{funny story}}    \example{
}    \example{
}}}

\entry{Toni}{\headword{Toni}
  \pos{Proper noun}  \sense{
    \definition{male personal name}}}

\entry{Tonny}{\headword{Tonny}  \pronnote{toni}
  \pos{Proper noun}  \sense{
    \definition{male personal name. Bearers: (1) Tonny Warama, son of Warama Kurupel and Wagiba Geser and member of Ende Language Committee.}}}

\entry{tontobabag}{\headword{tontobabag}  \note{good characteristics of people}  \pronnote{tontobabaɡ}
  \pos{Modifier}  \sense{
    \definition{neat, tidy}    \note{good characteristics of people}    \example{
}    \example{
}    \example{
}}}

\entry{tontobag}{\headword{tontobag}  \pronnote{tontobaɡ}
  \pos{Modifier}  \sense{
    \definition{wise}}}

\entry{tonton}{\headword{tonton}  \pronnote{tonton}
  \pos{Adverb}  \sense{
    \definition{directly}    \example{
      \exsen{Ewatta ke gänya täräp me lla bombllo da wasnen anggan ttowamang ttoen tonton ikop e ddapall att?}      \extran{Why does this generation beg to see miracles directly from heaven?}}}}

\entry{Tonzah}{\headword{Tonzah}
  \pos{Proper noun}  \sense{
    \definition{male personal name}}}

\entry{topoll}{\headword{topoll}  \note{Birds}  \pronnote{topoɽ}
  \pos{Noun}  \sense{
    \definition{flying fox}    \scientificname{Pteropus}    \note{Birds}    \example{
      \exsen{Abo ibi topoll gäddnan e dade angde bobäll.}      \extran{We can go kill flying foxes any time.}}    \example{
}    \example{
}}}

\entry{topotopoll}{\headword{topotopoll}  \note{Trees}  \pronnote{topotopoɽ}
  \pos{Noun}  \sense{
    \definition{type of tree that grows in the bush with white flowers; used as a yam stick}    \note{Trees}    \example{
}    \example{
}    \example{
}}}

\entry{torep}{\headword{torep}  \pronnote{torep}
  \pos{Noun}  \sense{
    \definition{brown quail}    \scientificname{Synoicus ypsilophorus dogwa}    \example{
}    \example{
}    \example{
}}}

\entry{toro}{\headword{toro}  \note{source: Jack Dipa}
  \pos{Noun}  \sense{
    \definition{a symbol of a slain animal, used to display and inform people of the type of animal killed. It was also a hunter's pride to compete or show other hunter' of his skill. Feathers, offcut tail, or animal fur were displayed on a small stick or pitpit or obäll tree stick or stem. In other instances pandanus leaves were symbols for pig, grass for wallaby.}    \note{source: Jack Dipa}}}

\entry{torok}{\headword{torok}  \lexeme{torok}  \pronnote{torok}
  \pos{Noun}  \sense{
    \definition{type of cane used for building houses, bows, and canoes}}}

\entry{Torok mittang}{\headword{Torok mittang}  \note{Place names around Limol}  \pronnote{torok miʈ͡ʂaŋ}
  \pos{Proper noun}  \sense{
    \definition{Torok mittang (toponym)}    \note{Place names around Limol}    \example{
}    \example{
}    \example{
}}}

\entry{toronggogo}{\headword{toronggogo}  \pronnote{toroŋɡoɡo}
  \pos{Noun}  \sense{
    \definition{bar-shouldered dove}    \scientificname{Geopelia humeralis}    \example{
}    \example{
}    \example{
}}}

\entry{torpoll}{\headword{torpoll}
  \variantof{SpellingVariant of \vartext{tärpoll}}}

\entry{torwam}{\headword{torwam}  \note{Verbs of Motion}  \pronnote{torwam}
  \pos{Intransitive verb}  \sense{\textbf{1.} 
    \definition{to lie down}    \note{Verbs of Motion}    \example{
      \exsen{tät me torwamang lla}      \extran{man lying down on a stretcher}}    \example{
      \exsen{Ngäna torwam allan.}      \extran{I am lying down.}}    \example{
}}  \sense{\textbf{2.} 
    \definition{to lay down}    \note{Verbs of Motion}    \example{
      \exsen{Ngäna obom buntruwam.}      \extran{I will lay him down.}}}}

\entry{torwamängg}{\headword{torwamängg}
  \pos{Transitive verb}  \sense{
    \definition{causative-applicative form of torwam}    \example{
      \exsen{Ngäna bam dantruwamängg.}      \extran{I made you lie down.}}}
  \variants{\Unspecified Variant \vartext{truwamängg}}}

\entry{tos}{\headword{tos}  \pronnote{tos}
  \pos{Noun}  \sense{
    \definition{(Commonwealth) torch, (US) flashlight}    \example{
      \exsen{Ngäna tos de dɨs.}      \extran{I turned off the flashlight.}}}}

\entry{tot}{\headword{tot}  \note{Trees}  \pronnote{tot}
  \pos{Noun}  \sense{
    \definition{piece}    \note{Trees}    \example{
      \exsen{Bäne aoli tot nyukukum dag?}      \extran{How many nyukukum bags do you have?}}}}

\entry{tot}{\headword{tot}  \note{Trees}  \pronnote{tot}
  \pos{Noun}  \sense{
    \definition{type of tree that gows along creeks with white and blue flowers and bark used to weave bags or sago baskets}    \note{Trees}    \example{
}    \example{
}    \example{
}}}

\entry{tot}{\headword{tot}  \note{Trees}  \pronnote{tot}
  \pos{Noun}  \sense{
    \definition{rubbish, trash, junk}    \note{Trees}    \example{
      \exsen{käm tot}      \extran{feces (lit. 'stomach trash')}}    \example{
      \exsen{emaemae tot}      \extran{all kinds of junk}}}}

\entry{totkoll}{\headword{totkoll}  \note{Water}  \pronnote{totkoɽ}
  \pos{Noun}  \sense{
    \definition{puddle}    \note{Water}    \example{
}    \example{
}    \example{
}}}

\entry{toto}{\headword{toto}  \pronnote{toto}
  \pos{Noun}  \sense{
    \definition{vertical house post}}}

\entry{toto}{\headword{toto}  \pronnote{toto}
  \pos{Noun}  \sense{
    \definition{afternoon; early evening (approx. 1 PM–5 PM)}}}

\entry{toto botta}{\headword{toto botta}  \note{Everything in the house}  \pronnote{toto boʈ͡ʂa}
  \pos{Noun}  \sense{
    \definition{horizontal house post}    \note{Everything in the house}    \example{
}    \example{
}    \example{
}}}

\entry{toto duwem}{\headword{toto duwem}  \note{Parts of the day}  \pronnote{toto duwem}
  \pos{Noun}  \sense{
    \definition{dinner}    \note{Parts of the day}    \example{
      \exsen{Ge ade ngämo toto duwem e dan.}      \extran{This is for my dinner.}}    \example{
}    \example{
}}}

\entry{toto nyäknyäk}{\headword{toto nyäknyäk}  \pronnote{toto ɲəkɲək}
  \pos{Noun}  \sense{
    \definition{Australian owlet-nightjar}    \scientificname{Aegotheles cristatus}    \example{
}    \example{
}    \example{
}}}

\entry{totoe}{\headword{totoe}
  \pos{Verb}  \sense{
}}

\entry{totrop}{\headword{totrop}  \pronnote{totrop}
  \pos{Transitive verb}  \sense{\textbf{1.} 
    \definition{to remove scales}}  \sense{\textbf{2.} 
    \definition{to clean up}}  \sense{\textbf{3.} 
    \definition{to scratch, like to scratch the tree with your fingers}}  \sense{\textbf{4.} 
    \definition{to scrape, e.g. when you shoot a deer and the arrow just scratches the animal}}  \sense{\textbf{5.} 
    \definition{to use the hoe or spade to remove small pieces of grass from a clear place}}}

\entry{towall}{\headword{towall}  \lexeme{towall}  \pronnote{towaɽ}
  \pos{Noun}  \sense{
    \definition{grass}    \example{
      \exsen{Ttall a towall wätätang dan.}      \extran{Wallabies are grass eaters.}}    \example{
}}}

\entry{towallang}{\headword{towallang}
  \pos{Modifier}  \sense{
    \definition{grassy}    \example{
      \exsen{Ge ttängäm a towallang dan.}      \extran{This village is grassy.}}}}

\entry{towallpipi}{\headword{towallpipi}  \note{Snakes}  \pronnote{towaɽpipi}
  \pos{Noun}  \sense{
    \definition{type of venomous snake}    \note{Snakes}    \example{
}    \example{
}    \example{
}}}

\entry{Towarwamang}{\headword{Towarwamang}
  \pos{Proper noun}  \sense{
    \definition{Towarwamang (toponym)}}}

\entry{traeb}{\headword{traeb}
  \pos{Noun}  \sense{
    \definition{tribe}}
  \variants{\SpellingVariant \vartext{traib}}}

\entry{traib}{\headword{traib}
  \variantof{SpellingVariant of \vartext{traeb}}}

\entry{training}{\headword{training}
  \variantof{Dialectal Variant of \vartext{English loanword}}}

\entry{trak}{\headword{trak}  \pronnote{trak}
  \pos{Noun}  \sense{
    \definition{truck}}
  \variants{\Unspecified Variant \vartext{träk}}}

\entry{träk}{\headword{träk}
  \variantof{Unspecified Variant of \vartext{trak}}}

\entry{trak}{\headword{trak}
  \variantof{Fast Speech Variant of \vartext{tärak}}}

\entry{träk}{\headword{träk}
  \variantof{Fast Speech Variant of \vartext{täräk}}  \pronnote{trək}}

\entry{träm}{\headword{träm}
  \variantof{Fast Speech Variant of \vartext{täräm}}  \pronnote{trəm}}

\entry{tram}{\headword{tram}
  \variantof{Fast Speech Variant of \vartext{täram}}}

\entry{trangg}{\headword{trangg}
  \variantof{Fast Speech Variant of \vartext{tärangg}}  \pronnote{traŋɡ}}

\entry{transfer}{\headword{transfer}
  \variantof{Dialectal Variant of \vartext{English loanword}}}

\entry{tratre}{\headword{tratre}
  \pos{Verb}  \sense{
}}

\entry{tri}{\headword{tri}
  \pos{Numeral}  \sense{
    \definition{three (English numeral)}}
  \variants{\SpellingVariant \vartext{three}}}

\entry{trik}{\headword{trik}  \pronnote{trik}
  \pos{Noun}  \sense{
    \definition{trick}}}

\entry{trimpeg}{\headword{trimpeg}
  \pos{Verb}  \sense{
}}

\entry{trongg}{\headword{trongg}
  \pos{Transitive verb}  \sense{
    \definition{to follow}    \example{
      \exsen{Ngäna däbe ddage dantrog do.}      \extran{I followed that stream there.}}}}

\entry{trongoe}{\headword{trongoe}  \pronnote{troŋoe}
  \pos{Transitive verb}  \sense{
    \definition{to check a bucket or container for water}}}

\entry{trungg}{\headword{trungg}
  \pos{Transitive verb}  \sense{
    \definition{to invite, call over, summon}    \example{
      \exsen{Ngänäm ngämo nag dantrugän Upiara we wayati kongkom e.}      \extran{My friend invited me to Upiara to collect watermelons.}}    \example{
      \exsen{Ngäna bam trungg allan, wiya!}      \extran{I am calling you over, come!}}}}

\entry{truwamängg}{\headword{truwamängg}
  \variantof{Unspecified Variant of \vartext{torwamängg}}  \pronnote{truwaməŋɡ}}

\entry{ttä~}{\headword{ttä~}
  \variantof{Infinitival reduplicant of \vartext{RED~}}}

\entry{ttäb}{\headword{ttäb}  \pronnote{ʈ͡ʂəb}
  \pos{Noun}  \sense{
    \definition{Pinon's imperial pigeon}    \scientificname{Ducula pinon pinon}    \example{
}    \example{
}    \example{
}}}

\entry{ttaba}{\headword{ttaba}
  \pos{Noun}  \sense{
}}

\entry{ttäbattäbe}{\headword{ttäbattäbe}
  \pos{Transitive verb}  \sense{
    \definition{to block}    \example{
      \exsen{Ttäbanen eran, ede ge nyongo da gagäll allan.}      \extran{It's blocked, so this road is no good.}}}}

\entry{ttäbe}{\headword{ttäbe}
  \pos{Noun}  \sense{
    \definition{a strong smelling plant whose bark, called ttäbe kollop, was traditionally worn around the neck to give fragrance and perfume smell. It was sometimes chewed and rubbed around the body and head to stop headache. The orignal purpose was also for protection from evil spirits.}}}

\entry{Ttäbe Ttäbe}{\headword{Ttäbe Ttäbe}  \note{Places around Limol}  \pronnote{ʈ͡ʂəbe ʈ͡ʂəbe}
  \pos{Proper noun}  \sense{
    \definition{Ttäbe Ttäbe (toponym)}    \note{Places around Limol}    \example{
}    \example{
}    \example{
}}}

\entry{ttäbottäbo}{\headword{ttäbottäbo}  \note{Parts of the Body}  \pronnote{ʈ͡ʂəboʈ͡ʂəbo}
  \pos{Noun}  \sense{
    \definition{rectum}    \note{Parts of the Body}    \example{
}    \example{
}    \example{
}}}

\entry{Ttae}{\headword{Ttae}  \pronnote{ʈ͡ʂae}
  \pos{Proper noun}  \sense{
    \definition{male personal name}}}

\entry{ttaem}{\headword{ttaem}  \pronnote{ʈ͡ʂaem}
  \pos{Transitive verb}  \sense{
    \definition{to pack}    \example{
      \exsen{Lama mängalae za de dättaemnegän.}      \extran{Lama quickly packed her things.}}}}

\entry{ttaempäg}{\headword{ttaempäg}  \pronnote{ʈ͡ʂaempəɡ}
  \pos{Transitive verb}  \sense{
    \definition{to seperate, divorce}    \example{
      \exsen{Ngämi ttaempägmeny dageyo.}      \extran{We two are inseparable.}}}}

\entry{ttaengän}{\headword{ttaengän}  \pronnote{ʈ͡ʂaeŋən}
  \pos{Transitive verb}  \sense{
    \definition{to pull (a plant sucker)}    \example{
}}}

\entry{ttaepnenttaepnen ma skul}{\headword{ttaepnenttaepnen ma skul}  \note{School}  \pronnote{ʈ͡ʂaepnenʈ͡ʂaepnen ma skul}
  \pos{Noun}  \sense{
    \definition{technical school}    \note{School}    \example{
}    \example{
}    \example{
}}}

\entry{ttägäll}{\headword{ttägäll}  \lexeme{ttägäll}  \pronnote{ʈ͡ʂəɡəɽ}
  \pos{Noun}  \sense{\textbf{1.} 
    \definition{termite mound, anthill (made by termites; ants may also live inside)}}  \sense{\textbf{2.} 
    \definition{mumu (oven made in the ground with fire, stones, leaves, and bark; the first mumus were made out of termite mounds)}    \example{
      \exsen{Bogo sɨmell sisiang de däbäddän ttägäll daogän.}      \extran{She killed the tame pig and made an earth oven.}}}}

\entry{ttägäll käp}{\headword{ttägäll käp}  \note{Conflict/war}  \pronnote{ʈ͡ʂəɡəɽ kəp}
  \pos{Noun}  \sense{\textbf{1.} 
    \definition{stone (esp. one used for cooking in mumus)}    \note{Conflict/war}    \example{
      \exsen{Dämetteya ttägäll käp de.}      \extran{We placed the oven stone.}}    \example{
      \exsen{Ngäna kumuddäga ttägäll käp tärpitärpi de nabllenegan.}      \extran{I found three smooth stones.}}    \example{
}}  \sense{\textbf{2.} 
    \definition{money}    \note{Conflict/war}    \example{
      \exsen{ttägäll käp mit me gäddgädd}      \extran{to fight about money}}}}

\entry{ttägäll ttoe}{\headword{ttägäll ttoe}  \note{Conflict/war}  \pronnote{ʈ͡ʂəɡəɽ ʈ͡ʂoe}
  \pos{Noun}  \sense{
    \definition{brown skin}    \note{Conflict/war}    \example{
}    \example{
}    \example{
}}}

\entry{Ttägällag kona}{\headword{Ttägällag kona}  \note{Places around Limol}  \pronnote{ʈ͡ʂəɡəɽaɡ kona}
  \pos{Proper noun}  \sense{
    \definition{Ttägälläg corner}    \note{Places around Limol}    \example{
}    \example{
}    \example{
}}}

\entry{Ttägällag pollon}{\headword{Ttägällag pollon}  \note{Places around Limol}  \pronnote{ʈ͡ʂəɡəɽaɡ poɽon}
  \pos{Proper noun}  \sense{
    \definition{Ttägällag pollon (toponym)}    \note{Places around Limol}    \example{
}    \example{
}    \example{
}}}

\entry{ttagbeag}{\headword{ttagbeag}  \note{Bad characteristics of people}  \pronnote{ʈ͡ʂaɡbeaɡ}
  \pos{Modifier}  \sense{
    \definition{disorganized, careless}    \note{Bad characteristics of people}    \example{
}    \example{
}    \example{
}}}

\entry{ttäk}{\headword{ttäk}  \note{Trees}  \pronnote{ʈ͡ʂək}
  \pos{Noun}  \sense{\textbf{1.} 
    \definition{type of tree that grows on floating grass (~3 m) with white or yellow flowers, green fruit, and a big trunk used to by children to float or for carfts}    \note{Trees}    \example{
}    \example{
}    \example{
}}  \sense{\textbf{2.} 
    \definition{soft wood}    \note{Trees}}}

\entry{ttäkäll}{\headword{ttäkäll}  \pronnote{ʈ͡ʂəkəɽ}
  \pos{Noun}  \sense{
    \definition{portion of yams mixed with coconut}}}

\entry{ttäkam}{\headword{ttäkam}  \pronnote{ʈ͡ʂəkam}
  \pos{Transitive verb}  \sense{
    \definition{to break, snap}    \example{
      \exsen{Obo ttɨle di dattkamän.}      \extran{She broke her leg.}}    \example{
      \exsen{Ubi ngäma sabi de ttäkam erallo.}      \extran{They are breaking our (excl.) law.}}}}

\entry{ttäkattäke}{\headword{ttäkattäke}
  \pos{Transitive verb}  \sense{
    \definition{to fold}}}

\entry{ttäkoe}{\headword{ttäkoe}  \note{Things a person can do}  \pronnote{ʈ͡ʂəkoe}
  \pos{Transitive verb}  \sense{
    \definition{to chop, cut down, mow; shave}    \note{Things a person can do}    \example{
      \exsen{Ngäna pos de ttokoenen de dängkamneg, tämamae naen dattkoeneg.}      \extran{I started to cut the posts; I chopped all nine.}}    \example{
      \exsen{Dade towall de battkoe.}      \extran{Maybe I'll mow the grass.}}    \example{
      \exsen{Ngäna obo ttatt kom de nattkoeyan.}      \extran{I shaved his beard.}}}
  \variants{\Dialectal Variant \vartext{ttokoe}}}

\entry{ttalam}{\headword{ttalam}  \pronnote{ʈ͡ʂalam}
  \pos{Intransitive verb}  \sense{
    \definition{to split, crack}    \example{
      \exsen{Käza bo ttatt a gottalamän.}      \extran{The crocodile's jaw cracked open.}}}}

\entry{ttalamttalam}{\headword{ttalamttalam}  \note{Trees}  \pronnote{ʈ͡ʂalamʈ͡ʂalam}
  \pos{Noun}  \sense{
    \definition{type of tree}    \note{Trees}    \example{
}    \example{
}    \example{
}}}

\entry{ttalamttalam}{\headword{ttalamttalam}  \pronnote{ʈ͡ʂalamʈ͡ʂalam}
  \pos{Adverb}  \sense{
    \definition{walking with one's legs spread far apart}}}

\entry{ttalängg}{\headword{ttalängg}  \pronnote{ʈ͡ʂaləŋɡ}
  \pos{Intransitive verb}  \sense{
    \definition{to aim}    \example{
      \exsen{Kwalde mäse pakos tupi di gonttalägän sɨmell zuwoe e.}      \extran{Kwalde tried aiming the long arrow to shoot the pig.}}}}

\entry{ttäle}{\headword{ttäle}  \pronnote{ʈ͡ʂəle}
  \pos{Noun}  \sense{
    \definition{type of tree}}}

\entry{ttäle}{\headword{ttäle}  \pronnote{ʈ͡ʂəle}
  \pos{Noun}  \sense{\textbf{1.} 
    \definition{leg}    \example{
      \exsen{Ngämo ttäle da kädkädag agallo.}      \extran{My legs got cold.}}}  \sense{\textbf{2.} 
    \definition{tendril}}
  \variants{\Dialectal Variant \vartext{ttɨle, ttälle}}}

\entry{Ttäle Bun}{\headword{Ttäle Bun}  \pronnote{ʈ͡ʂəle bun}
  \pos{Proper noun}  \sense{
    \definition{Ttäle Bun (toponym)}}}

\entry{ttäle koll}{\headword{ttäle koll}  \note{Parts of the body}  \pronnote{ʈ͡ʂəle koɽ}
  \pos{Noun}  \sense{
    \definition{instep}    \note{Parts of the body}    \example{
}    \example{
}    \example{
}}}

\entry{ttäle kum}{\headword{ttäle kum}
  \pos{Noun}  \sense{
}}

\entry{ttäle lläpät}{\headword{ttäle lläpät}  \note{Parts of an insect}  \pronnote{ʈ͡ʂəle ɽəpət}
  \pos{Noun}  \sense{
    \definition{claw}    \note{Parts of an insect}    \example{
}    \example{
}    \example{
}}}

\entry{Ttäle mitt}{\headword{Ttäle mitt}  \note{Places around Limol}  \pronnote{ʈ͡ʂəle miʈ͡ʂ}
  \pos{Proper noun}  \sense{
    \definition{place with a well in Limol}    \note{Places around Limol}    \example{
}    \example{
}    \example{
}}}

\entry{ttäle pätt}{\headword{ttäle pätt}  \note{Parts of a reptile}  \pronnote{ʈ͡ʂəle pəʈ͡ʂ}
  \pos{Noun}  \sense{
    \definition{hind leg, hind limb}    \note{Parts of a reptile}    \example{
}    \example{
}    \example{
}}}

\entry{ttäle pättkäp}{\headword{ttäle pättkäp}  \note{Parts of an insect}  \pronnote{ʈ͡ʂəle pəʈ͡ʂkəp}
  \pos{Noun}  \sense{
    \definition{tibia}    \note{Parts of an insect}    \example{
}    \example{
}    \example{
}}}

\entry{ttäle pe}{\headword{ttäle pe}  \note{Parts of a reptile}  \pronnote{ʈ͡ʂəle pe}
  \pos{Noun}  \sense{
    \definition{hind foot (of a reptile) with four toes}    \note{Parts of a reptile}    \example{
}    \example{
}    \example{
}}}

\entry{ttäle pitt}{\headword{ttäle pitt}  \note{Clothing}  \pronnote{ʈ͡ʂəle piʈ͡ʂ}
  \pos{Noun}  \sense{
    \definition{leg band}    \note{Clothing}    \example{
}    \example{
}    \example{
}}}

\entry{Ttälebun}{\headword{Ttälebun}  \note{Places around Limol}  \pronnote{ʈ͡ʂəlebun}
  \pos{Proper noun}  \sense{
    \definition{Ttalebun (toponym)}    \note{Places around Limol}    \example{
}    \example{
}    \example{
}}}

\entry{Ttall}{\headword{Ttall}
  \pos{Proper noun}  \sense{
    \definition{male personal name}}}

\entry{ttall}{\headword{ttall}  \lexeme{ttall}  \pronnote{ʈ͡ʂaɽ}
  \pos{Noun}  \sense{
    \definition{agile wallaby, sandy wallaby}    \scientificname{Notamacropus agilis papuanus}}}

\entry{ttall ip}{\headword{ttall ip}  \note{Trees}  \pronnote{ʈ͡ʂaɽ ip}
  \pos{Noun}  \sense{
    \definition{type of tree that grows in the bush, grassland, and along creeks with indigo flowers, bark that is chewed, and liquid used as glue for spears}    \note{Trees}    \example{
}    \example{
}    \example{
}}}

\entry{ttall källa}{\headword{ttall källa}  \pronnote{ʈ͡ʂaɽ kəɽa}
  \pos{Noun}  \sense{
    \definition{brahminy kite}    \scientificname{Haliastur indus girrenera}    \example{
}    \example{
}    \example{
}}}

\entry{ttall mätta}{\headword{ttall mätta}  \pronnote{ʈ͡ʂaɽ məʈ͡ʂa}
  \pos{Noun}  \sense{
    \definition{type of big yam with a white interior and few hairs}}}

\entry{ttall nge}{\headword{ttall nge}
  \pos{Noun}  \sense{
    \definition{type of palm with yellow leaves and coconuts with a yellow exocarp}}}

\entry{ttall ttoe}{\headword{ttall ttoe}  \note{Trees}  \pronnote{ʈ͡ʂaɽ ʈ͡ʂoe}
  \pos{Noun}  \sense{
    \definition{type of tree}    \note{Trees}    \example{
}    \example{
}    \example{
}}}

\entry{ttälla}{\headword{ttälla}  \pronnote{ʈ͡ʂəɽa}
  \pos{Verb}  \sense{
    \definition{to hurt}}}

\entry{ttällallem}{\headword{ttällallem}  \pronnote{ʈ͡ʂəɽaɽem}
  \pos{Intransitive coverb}  \sense{
    \definition{to make a light noise}}}

\entry{ttällam}{\headword{ttällam}  \pronnote{ʈ͡ʂəɽam}
  \pos{Transitive verb}  \sense{\textbf{1.} 
    \definition{to pass, hand}    \example{
      \exsen{Ngäna bäne pate neil de bänttllam.}      \extran{I will pass you the nail.}}}  \sense{\textbf{2.} 
    \definition{to extend, stretch out, reach out, put out}    \example{
      \exsen{Bäne ttang de ttättlemae nänttllam.}      \extran{Stretch your arm out straight.}}}}

\entry{ttällanenang}{\headword{ttällanenang}
  \pos{Modifier}  \sense{
    \example{
      \exsen{Ddob llayabira komo ddäddäg a malla ttälanenang dag.}      \extran{The centipede bite isn't painful for some people.}}}}

\entry{ttallängg}{\headword{ttallängg}
  \pos{Verb}  \sense{
}}

\entry{ttälläp}{\headword{ttälläp}  \note{Snakes}  \pronnote{ʈ͡ʂəɽəp}
  \pos{Noun}  \sense{
    \definition{type of venomous snake}    \note{Snakes}    \example{
}    \example{
}    \example{
}}}

\entry{ttällattälla}{\headword{ttällattälla}  \pronnote{ʈ͡ʂəɽaʈ͡ʂəɽa}
  \pos{Transitive coverb}  \sense{
    \definition{to hit lightly}}}

\entry{ttälle}{\headword{ttälle}
  \variantof{Dialectal Variant of \vartext{ttäle}}  \pronnote{ʈ͡ʂəɽe}}

\entry{ttällma tränymägäll}{\headword{ttällma tränymägäll}  \note{Trees}  \pronnote{ʈ͡ʂəɽma trəɲməɡəɽ}
  \pos{Noun}  \sense{
    \definition{tree type}    \note{Trees}    \example{
}    \example{
}    \example{
}}}

\entry{ttalme}{\headword{ttalme}  \note{Water}  \pronnote{ʈ͡ʂalme}
  \pos{Noun}  \sense{
    \definition{type of floating grass that grass, deer, and wallaby eat}    \note{Water}    \example{
}    \example{
}    \example{
}}}

\entry{ttam}{\headword{ttam}
  \pos{Intransitive verb}  \sense{
    \definition{to appear}    \example{
      \exsen{Yesu ngattong abal e Meri Magdalin pate gottamän.}      \extran{Jesus appeared to Mary Magdalene first.}}}}

\entry{ttäm}{\headword{ttäm}  \pronnote{ʈ͡ʂəm}
  \pos{Transitive verb}  \sense{
    \definition{to burn; heat on a fire}    \example{
      \exsen{Yu da obom dättämän.}      \extran{The fire burned him.}}    \example{
      \exsen{Därunggu de yu mi wandawandae dättämalle.}      \extran{The bamboo tube is heated rotating on the fire.}}}}

\entry{ttam}{\headword{ttam}  \pronnote{ʈ͡ʂam}
  \pos{Transitive verb}  \sense{\textbf{1.} 
    \definition{to call, name}    \example{
      \exsen{Da bongo mäse dowae me o utale me nänttam, ddone mɨnyi bongllaeyän adawatta obo llan a ttomoll dageyo.}      \extran{If you tried to call him, from near or far, he would not answer because he was deaf in his ears.}}    \example{
      \exsen{Kaya walle do Upiara ttängämang a ttaem nallo.}      \extran{Upiara villagers call me Kaya.}}    \example{
      \exsen{Kumddäga pemli de ngäna amim nättaemnegan, ngämi ubim mɨnyi umllang bägaebeya.}      \extran{The three families I named, we (excl.) will tell them.}}}  \sense{\textbf{2.} 
    \definition{to confess}    \example{
      \exsen{Angde bogo gonttamän, llamda da ddone ada mikutt gogon oba pate.}      \extran{When he confessed, the old man got very mad at them.}}}  \sense{\textbf{3.} 
    \definition{to mention}}}

\entry{ttam}{\headword{ttam}  \lexeme{ttam}  \pronnote{ʈ͡ʂam}
  \pos{Noun}  \sense{
    \definition{leaf}    \example{
      \exsen{Ngämlle kapalla ttam alle sana yuatt a ddobae moko dan.}      \extran{I love sago cooked with kapalla leaves.}}}}

\entry{ttam}{\headword{ttam}  \lexeme{ttam}  \pronnote{ʈ͡ʂam}
  \pos{Noun}  \sense{\textbf{1.} 
    \definition{life}    \example{
      \exsen{Mälla ause da llokttang ttam me dagirnän.}      \extran{The old woman lived a hard life.}}}  \sense{\textbf{2.} 
    \definition{alive}    \example{
      \exsen{Matthew kili allan adawatta bogo sisri ttam dan.}      \extran{Matthew is happy because he is alive now.}}    \example{
      \exsen{Bogo kuddäll atta ttam gogon.}      \extran{She survived death.}}}  \sense{\textbf{3.} 
    \definition{to save someone's life}    \example{
      \exsen{Bäne imomdae ttoen da bam ttam nagan.}      \extran{Your faith saved your life.}}}}

\entry{ttäm}{\headword{ttäm}
  \variantof{Fast Speech Variant of \vartext{ttängäm}}}

\entry{ttam giddoll}{\headword{ttam giddoll}
  \pos{Noun}  \sense{
    \definition{life}    \example{
      \exsen{Do ttam giddoll a mer dan.}      \extran{Life there is good.}}}}

\entry{ttaman}{\headword{ttaman}
  \variantof{Unspecified Variant of \vartext{ttamän}}  \pronnote{ʈ͡ʂaman}}

\entry{ttamän}{\headword{ttamän}  \pronnote{ʈ͡ʂamən}
  \pos{Transitive verb}  \sense{\textbf{1.} 
    \definition{to finish, end}    \example{
      \exsen{Soka tongoe bo ebdo da zäme wottamänan.}      \extran{The day of the soccer game is over.}}    \example{
      \exsen{Angde tine pittnen a gottamänän, dibaballe ddäganen de dängkam.}      \extran{After the sago weaving finished, I started on the roofing.}}}  \sense{\textbf{2.} 
    \definition{to finish, end, complete}    \example{
      \exsen{Ngäna greid siks de dättemän.}      \extran{I completed grade six.}}    \example{
      \exsen{Ngäna ngämo sana de dot dägadäga dättemän.}      \extran{I finished eating all my sago.}}}
  \variants{\Unspecified Variant \vartext{ttaman}}}

\entry{ttämattäme}{\headword{ttämattäme}  \pronnote{ʈ͡ʂəmaʈ͡ʂəme}
  \pos{Intransitive verb}  \sense{
    \definition{to be crowded}    \example{
      \exsen{Ngättma da lla walle gonttämawän.}      \extran{The place was crowded with people.}}}}

\entry{ttambällag}{\headword{ttambällag}  \pronnote{ʈ͡ʂambəɽaɡ}
  \pos{Noun}  \sense{
    \definition{the second stage of coconut growth in which it is planted}    \example{
      \exsen{Ngäna ngämo nge ttambällag de tätäm dibeny.}      \extran{I planted my ttambällag coconut yesterday.}}}}

\entry{ttämbe}{\headword{ttämbe}  \note{Trees}  \pronnote{ʈ͡ʂəmbe}
  \pos{Noun}  \sense{
    \definition{type of big tree that grows in the bush with blue, purple, and white flowers and red fruit}    \note{Trees}    \example{
}    \example{
}    \example{
}}}

\entry{ttämbe role}{\headword{ttämbe role}  \note{Trees}  \pronnote{ʈ͡ʂəmbe role}
  \pos{Noun}  \sense{
    \definition{type of tree}    \note{Trees}    \example{
}    \example{
}    \example{
}}}

\entry{ttamttem}{\headword{ttamttem}
  \variantof{SpellingVariant of \vartext{ttanttem}}}

\entry{ttän}{\headword{ttän}  \note{Trees}  \pronnote{ʈ͡ʂən}
  \pos{Noun}  \sense{
    \definition{type of tree that grows in the grassland (~60 m) with yellow flowers, small green fruit, and wood used for house posts}    \note{Trees}    \example{
}    \example{
}    \example{
}}}

\entry{ttän maigag}{\headword{ttän maigag}  \pronnote{ʈ͡ʂən maiɡaɡ}
  \pos{Noun}  \sense{
    \definition{type of bandicoot}    \example{
}    \example{
}    \example{
}}}

\entry{ttang}{\headword{ttang}  \pronnote{ʈ͡ʂaŋ}
  \pos{Noun}  \sense{\textbf{1.} 
    \definition{hand}    \example{
      \exsen{Kaemne de ttang alle mällam a mudan.}      \extran{Don't hold bees with your hands.}}}  \sense{\textbf{2.} 
    \definition{arm}    \example{
      \exsen{Kottllam obo ttang a wa ttäle da tubutubu dageyo.}      \extran{The turtle's arms and legs are short.}}}}

\entry{ttang käpän}{\headword{ttang käpän}  \pronnote{ʈ͡ʂaŋ kəpən}
  \pos{Intransitive compound verb}  \sense{
    \definition{to snap one's fingers}}}

\entry{ttang koll}{\headword{ttang koll}  \note{Parts of the body}  \pronnote{ʈ͡ʂaŋ koɽ}
  \pos{Noun}  \sense{
    \definition{palm}    \note{Parts of the body}    \example{
}    \example{
}    \example{
}}}

\entry{ttang kum}{\headword{ttang kum}  \note{Numbers}  \pronnote{ʈ͡ʂaŋ kum}
  \pos{Noun}  \sense{\textbf{1.} 
    \definition{elbow}    \note{Numbers}}  \sense{\textbf{2.} 
    \definition{seven (lit. elbow; body counting numeral)}    \note{Numbers}    \example{
}    \example{
}    \example{
}}}

\entry{ttang lläbe}{\headword{ttang lläbe}  \note{Parts of a reptile}  \pronnote{ʈ͡ʂaŋ ɽəbe}
  \pos{Noun}  \sense{
    \definition{talon, claw}    \note{Parts of a reptile}    \example{
}    \example{
}    \example{
}}}

\entry{ttang lläpät}{\headword{ttang lläpät}  \note{Parts of the Body}  \pronnote{ʈ͡ʂaŋ ɽəpət}
  \pos{Noun}  \sense{\textbf{1.} 
    \definition{finger}    \note{Parts of the Body}    \example{
      \exsen{Tupi ttang llɨpɨt de pop me nowanseg.}      \extran{Put your long finger in the hole.}}    \example{
}    \example{
}}  \sense{\textbf{2.} 
    \definition{unit of five (one hand)}    \note{Parts of the Body}    \example{
      \exsen{Obo mälla da apte ttang lläpät da.}      \extran{He has five wives.}}    \example{
      \exsen{Ngämo ttongo ttang lläpät a muminy dan.}      \extran{I owe five kina.}}    \example{
      \exsen{Ngämo kollba mu da komlla ttang lläpät dag.}      \extran{The price of my fish is 10 kina.}}}}

\entry{ttang päk}{\headword{ttang päk}
  \pos{Noun}  \sense{
}}

\entry{ttang pättkäp}{\headword{ttang pättkäp}  \note{check definition. shouldn't it be humerus, ulna, or radia? tibia is in leg}  \pronnote{ʈ͡ʂaŋ pəʈ͡ʂkəp}
  \pos{Noun}  \sense{
    \definition{tibia}    \note{check definition. shouldn't it be humerus, ulna, or radia? tibia is in leg}    \example{
}    \example{
}    \example{
}}}

\entry{ttang pe}{\headword{ttang pe}  \note{Parts of a reptile}  \pronnote{ʈ͡ʂaŋ pe}
  \pos{Noun}  \sense{
    \definition{upper part of foreleg (of a reptile)}    \note{Parts of a reptile}    \example{
}    \example{
}    \example{
}}}

\entry{ttang pitt}{\headword{ttang pitt}  \note{Clothing}  \pronnote{ʈ͡ʂaŋ piʈ͡ʂ}
  \pos{Noun}  \sense{
    \definition{bracelet}    \note{Clothing}    \example{
}    \example{
}    \example{
}}}

\entry{ttang pllallem}{\headword{ttang pllallem}  \note{Verbs of Motion}  \pronnote{ʈ͡ʂaŋ pɽaɽem}
  \pos{Intransitive coverb}  \sense{
    \definition{to clap}    \note{Verbs of Motion}    \example{
      \exsen{Bogo ttang pllallem allan.}      \extran{She is clapping.}}    \example{
}    \example{
}}}

\entry{ttang ttaempäg}{\headword{ttang ttaempäg}  \lexeme{ttaempäg}  \pronnote{ʈ͡ʂaempəɡ}
  \pos{Transitive verb}  \sense{
    \definition{to shake hands with}    \example{
      \exsen{Nagnag a ngämim ttang deyanttepmenyeyo.}      \extran{Friends shook hands with us (excl.).}}}}

\entry{ttang tupiang sod}{\headword{ttang tupiang sod}  \note{Clothing}  \pronnote{ʈ͡ʂaŋ tupiaŋ sod}
  \pos{Noun}  \sense{
    \definition{long-sleeved shirt}    \note{Clothing}    \example{
}    \example{
}    \example{
}}}

\entry{ttängäm}{\headword{ttängäm}  \pronnote{ʈ͡ʂəŋəm}
  \pos{Noun}  \sense{\textbf{1.} 
    \definition{village}    \example{
      \exsen{Obo ttängäm de enanae dowansegän.}      \extran{He left his village for good.}}}  \sense{\textbf{2.} 
    \definition{garden}    \example{
      \exsen{Ubi ttängäm kuddäll me polle de kame däkätteyo.}      \extran{They built a new fence around the abandoned garden.}}    \example{
}    \example{
}}  \sense{\textbf{3.} 
    \definition{place}    \example{
      \exsen{Ero bäne zegatt ttängäm a?}      \extran{Where is your birthplace?}}}
  \variants{\Fast Speech Variant \vartext{ttäm}}}

\entry{ttängämttängäm}{\headword{ttängämttängäm}  \pronnote{ʈ͡ʂəŋəmʈ͡ʂəŋəm}
  \pos{Noun}  \sense{
    \definition{small garden}}}

\entry{ttänganen~}{\headword{ttänganen~}
  \variantof{freevarlab of \vartext{RED~}}}

\entry{ttängattänge}{\headword{ttängattänge}  \note{School}  \pronnote{ʈ͡ʂəŋaʈ͡ʂəŋe}
  \pos{Transitive verb}  \sense{\textbf{1.} 
    \definition{to read, recite}    \note{School}    \example{
      \exsen{Ddob llaeyaba eka walle darbnen att eka da, ngäna mɨnyi bangttängeneg.}      \extran{Stories written in other people's languages, I will read them.}}    \example{
}    \example{
}}  \sense{\textbf{2.} 
    \definition{date}    \note{School}    \example{
      \exsen{Sisri ngasekäma ttängattänge da endan?}      \extran{What's the date today?}}}}

\entry{ttängkag}{\headword{ttängkag}  \pronnote{ʈ͡ʂəŋkaɡ}
  \pos{Transitive verb}  \sense{
    \definition{to break with force}    \example{
      \exsen{Pakos ulle walle ngäna ddia bo ddäg kutt deyanttäkeg.}      \extran{I broke the deer's backbone using a big spear.}}}}

\entry{ttängkamäll}{\headword{ttängkamäll}  \pronnote{ʈ͡ʂəŋkaməɽ}
  \pos{Transitive verb}  \sense{
    \definition{to meet, reach}    \example{
      \exsen{Ubi dinduag a duduli Awayang bom danttäkämälleyo.}      \extran{Running that way, they reached Awayang.}}}}

\entry{ttangkumttangkum}{\headword{ttangkumttangkum}  \pronnote{ʈ͡ʂaŋkumʈ͡ʂaŋkum}
  \pos{Adverb}  \sense{
    \definition{on all fours}}}

\entry{ttangkuttangkumang}{\headword{ttangkuttangkumang}  \note{Weaving}  \pronnote{ʈ͡ʂaŋkuʈ͡ʂaŋkumaŋ}
  \pos{Noun}  \sense{
    \definition{sideways chevron weaving pattern}    \note{Weaving}    \example{
}    \example{
}    \example{
}}}

\entry{ttangttang}{\headword{ttangttang}  \pronnote{ʈ͡ʂaŋʈ͡ʂaŋ}
  \pos{Adverb}  \sense{
    \definition{with one's hands full}    \example{
      \exsen{Ngäna ttangttang agan.}      \extran{My hands are full.}}}}

\entry{ttangttang}{\headword{ttangttang}  \note{Birds}  \pronnote{ʈ͡ʂaŋʈ͡ʂaŋ}
  \pos{Noun}  \sense{
    \definition{type of bird}    \note{Birds}    \example{
}    \example{
}    \example{
}}}

\entry{ttänttäm}{\headword{ttänttäm}  \note{Parts of a fire}  \pronnote{ʈ͡ʂənʈ͡ʂəm}
  \pos{Noun}  \sense{\textbf{1.} 
    \definition{heat}    \note{Parts of a fire}    \example{
      \exsen{Däbe lla da obo labalaba de dängkänän ttänttäm atta.}      \extran{That man took of his lap-lap because of the heat.}}}  \sense{\textbf{2.} 
    \definition{to burn}    \note{Parts of a fire}    \example{
      \exsen{Ikrol era yu ttänttämatt me pänpän dan.}      \extran{Ash is dust left after a fire burns.}}}}

\entry{ttänttämang}{\headword{ttänttämang}  \note{Adjectives}  \pronnote{ʈ͡ʂənʈ͡ʂəmaŋ}
  \pos{Modifier}  \sense{
    \definition{hot}    \note{Adjectives}    \example{
      \exsen{Ngämo moko da ttänttämang kollba ddäddäg e dan.}      \extran{I like to eat hot fish.}}    \example{
}    \example{
}}}

\entry{ttanttem}{\headword{ttanttem}
  \pos{Intransitive verb}  \sense{
    \definition{to be confused}    \example{
      \exsen{Obo eka me ubi gonttemnegän.}      \extran{They were confused by his story.}}}
  \variants{\SpellingVariant \vartext{ttamttem}}}

\entry{ttäp}{\headword{ttäp}
  \variantof{Unspecified Variant of \vartext{ttɨp}}}

\entry{ttape}{\headword{ttape}  \pronnote{ʈ͡ʂape}
  \pos{Noun}  \sense{
    \definition{type of small, flat fish}}}

\entry{ttäpen}{\headword{ttäpen}  \pronnote{ʈ͡ʂəpen}
  \pos{Transitive verb}  \sense{\textbf{1.} 
    \definition{to snap, break}    \example{
      \exsen{Kwalde mäse pakos de gonttalägän sämell zuwoe e, be bägäl wadär a gottäpenän.}      \extran{Kwalde tried to aim the arrow to shoot the pig, but the bow string snapped.}}}  \sense{\textbf{2.} 
    \definition{to break, tear (esp. something long)}    \example{
      \exsen{Ngäna ngämo bunkom de nättpenan.}      \extran{I broke a strand of my hair.}}    \example{
      \exsen{Obo tudi di kollba da dättpenän.}      \extran{The fish broke her fishing line.}}}  \sense{\textbf{3.} 
}}

\entry{ttapeyam}{\headword{ttapeyam}  \pronnote{ʈ͡ʂapejam}
  \pos{Transitive verb}  \sense{
    \definition{to open (something long, e.g. door, book)}}}

\entry{ttapeyamttapeyam}{\headword{ttapeyamttapeyam}  \pronnote{ʈ͡ʂapejamʈ͡ʂapejam}
  \pos{Adverb}  \sense{
    \definition{walking with one's legs spread far apart}}}

\entry{ttaputtapung}{\headword{ttaputtapung}  \pronnote{ʈ͡ʂapuʈ͡ʂapuŋ}
  \pos{Modifier}  \sense{
    \definition{joined together}    \example{
}    \example{
}    \example{
}}}

\entry{ttärtt}{\headword{ttärtt}
  \variantof{SpellingVariant of \vartext{sos}}}

\entry{ttätt}{\headword{ttätt}  \pronnote{ʈ͡ʂəʈ͡ʂ}
  \pos{Modifier}  \sense{
    \definition{right}    \example{
      \exsen{Ge ngämo ttätt ttang dan.}      \extran{This is my right hand.}}}}

\entry{ttatt}{\headword{ttatt}  \note{Parts of the Body}  \pronnote{ʈ͡ʂaʈ͡ʂ}
  \pos{Noun}  \sense{
    \definition{jaw, chin}    \note{Parts of the Body}    \example{
      \exsen{Simell ttatt a ulle dan.}      \extran{A pig's jaw is big.}}    \example{
}    \example{
}}}

\entry{ttätt käp}{\headword{ttätt käp}  \pronnote{ʈ͡ʂəʈ͡ʂ kəp}
  \pos{Noun}  \sense{\textbf{1.} 
    \definition{graceful honeyeater}    \scientificname{Microptilotis gracilis}    \example{
}    \example{
}    \example{
}}  \sense{\textbf{2.} 
    \definition{puff-backed honeyeater}    \scientificname{Meliphaga aruensis}}}

\entry{ttatt käpkäp}{\headword{ttatt käpkäp}  \pronnote{ʈ͡ʂaʈ͡ʂ kəpkəp}
  \pos{Noun}  \sense{
    \definition{chubby cheeks}}}

\entry{ttatt kom}{\headword{ttatt kom}  \note{Parts of the Body}  \pronnote{ʈ͡ʂaʈ͡ʂ kom}
  \pos{Noun}  \sense{
    \definition{beard}    \note{Parts of the Body}    \example{
}    \example{
}    \example{
}}}

\entry{ttatt kutt}{\headword{ttatt kutt}  \note{Musical instruments}  \pronnote{ʈ͡ʂaʈ͡ʂ kuʈ͡ʂ}
  \pos{Noun}  \sense{\textbf{1.} 
    \definition{mandible, jawbone}    \note{Musical instruments}}  \sense{\textbf{2.} 
    \definition{drum rim}    \note{Musical instruments}    \example{
}    \example{
}    \example{
}}}

\entry{ttätta}{\headword{ttätta}
  \pos{Nominal}  \sense{
}}

\entry{ttatta}{\headword{ttatta}  \note{Birth}  \pronnote{ʈ͡ʂaʈ͡ʂa}
  \pos{Noun}  \sense{
    \definition{lower back}    \note{Birth}    \example{
      \exsen{Obo ttatta pallall e deyazuwän.}      \extran{He shot towards his lower back.}}}}

\entry{ttatta}{\headword{ttatta}  \pronnote{ʈ͡ʂaʈ͡ʂa}
  \pos{Transitive verb}  \sense{
    \definition{to chop a tree}    \example{
      \exsen{Bogo llo de ttanen anggan.}      \extran{He's chopping a tree.}}}}

\entry{ttättäle}{\headword{ttättäle}
  \variantof{Unspecified Variant of \vartext{ttättle}}}

\entry{ttättälläg}{\headword{ttättälläg}  \pronnote{ʈ͡ʂəʈ͡ʂəɽəɡ}
  \pos{Transitive verb}  \sense{
    \definition{to tear, tear up}    \example{
      \exsen{Gudne iddpo ttättllägatt pop de mäzi abal ulle ubemang bägagän.}      \extran{It will make the torn holes in the old clothing even wider.}}}}

\entry{ttattang}{\headword{ttattang}  \note{Food}  \pronnote{ʈ͡ʂaʈ͡ʂaŋ}
  \pos{Noun}  \sense{
    \definition{type of big, round yam with a white interior}    \note{Food}    \example{
}    \example{
}    \example{
}}}

\entry{ttättäp}{\headword{ttättäp}
  \pos{Noun}  \sense{
    \definition{young leaf}}}

\entry{ttättawe}{\headword{ttättawe}  \note{Trees}  \pronnote{ʈ͡ʂəʈ͡ʂawe}
  \pos{Noun}  \sense{
    \definition{type of tree}    \note{Trees}    \example{
}    \example{
}    \example{
}}}

\entry{ttattel}{\headword{ttattel}  \lexeme{ttattel}  \pronnote{ʈ͡ʂaʈ͡ʂel}
  \pos{Noun}  \sense{
    \definition{type of thorny vine}}}

\entry{ttattep}{\headword{ttattep}
  \pos{Noun}  \sense{
    \definition{mature leaf}}}

\entry{ttattle}{\headword{ttattle}  \pronnote{ʈ͡ʂaʈ͡ʂle}
  \pos{Noun}  \sense{
    \definition{pain}    \example{
      \exsen{Bogo ttattle de umllang gogän.}      \extran{He felt pain.}}}
  \variants{\Unspecified Variant \vartext{ttattlle}}}

\entry{ttättle}{\headword{ttättle}  \note{Adjectives}  \pronnote{ʈ͡ʂəʈ͡ʂle}
  \pos{Modifier}  \sense{\textbf{1.} 
    \definition{correct, proper}    \note{Adjectives}    \example{
      \exsen{Ddob eka kutt panypeny a sisor llɨg aba ddone ttättle dag.}      \extran{Some of the words of young people are not pronounced correctly.}}}  \sense{\textbf{2.} 
    \definition{straight}    \note{Adjectives}    \example{
      \exsen{Malläm me nyongo da ddone ttättle dan.}      \extran{The road to Malam is not straight.}}    \example{
      \exsen{Ubi deyareyo ttättlemae do up wo kakab a erame daeya.}      \extran{They went straight to where the leftover ripe bananas were.}}    \example{
}}
  \variants{\Unspecified Variant \vartext{ttättäle}}}

\entry{ttattleang}{\headword{ttattleang}  \pronnote{ʈ͡ʂaʈ͡ʂleaŋ}
  \pos{Modifier}  \sense{
    \definition{sore, in pain; sick, ill}}
  \variants{\Fast Speech Variant \vartext{ttattllong, ttattlong}}}

\entry{ttattlläb}{\headword{ttattlläb}
  \pos{Intransitive verb}  \sense{
    \definition{to open}    \example{
      \exsen{Wäd a wättallban.}      \extran{The door opened.}}}}

\entry{ttattlle}{\headword{ttattlle}
  \variantof{Unspecified Variant of \vartext{ttattle}}}

\entry{ttattllong}{\headword{ttattllong}
  \variantof{Fast Speech Variant of \vartext{ttattleang}}  \pronnote{ʈ͡ʂaʈ͡ʂɽoŋ}}

\entry{ttattlong}{\headword{ttattlong}
  \variantof{Fast Speech Variant of \vartext{ttattleang}}  \pronnote{ʈ͡ʂaʈ͡ʂloŋ}}

\entry{ttätto}{\headword{ttätto}
  \variantof{SpellingVariant of \vartext{ttotto}}}

\entry{ttek}{\headword{ttek}  \note{Trees}  \pronnote{ʈ͡ʂek}
  \pos{Noun}  \sense{
    \definition{type of tree that grows in the grassland with white flowers and sap used as glue for spears}    \note{Trees}    \example{
}    \example{
}    \example{
}}}

\entry{ttette}{\headword{ttette}  \note{Everything in the house}  \pronnote{ʈ͡ʂeʈ͡ʂe}
  \pos{Noun}  \sense{
    \definition{rafter}    \note{Everything in the house}    \example{
}    \example{
}    \example{
}}}

\entry{ttima~}{\headword{ttima~}
  \variantof{Inflected Form of \vartext{RED~}}}

\entry{ttimattima}{\headword{ttimattima}
  \pos{Noun}  \sense{
    \definition{limp}    \example{
      \exsen{Bogo ttimattimang ibiag daeya.}      \extran{He walked with a limp.}}}}

\entry{ttɨle}{\headword{ttɨle}
  \variantof{Dialectal Variant of \vartext{ttäle}}}

\entry{ttɨp}{\headword{ttɨp}  \note{Sago palms}  \pronnote{ʈ͡ʂɪp}
  \pos{Noun}  \sense{\textbf{1.} 
    \definition{type of sago}    \note{Sago palms}    \example{
}    \example{
}    \example{
}}  \sense{\textbf{2.} 
    \definition{type of yam with a white interior}    \note{Sago palms}}
  \variants{\Unspecified Variant \vartext{ttäp}}}

\entry{ttoa}{\headword{ttoa}
  \variantof{SpellingVariant of \vartext{ttowa}}  \note{Birds}  \pronnote{ʈ͡ʂoa}}

\entry{ttoe}{\headword{ttoe}  \pronnote{ʈ͡ʂoe}
  \pos{Noun}  \sense{\textbf{1.} 
    \definition{skin (of a person or animal)}    \example{
      \exsen{Däbe käza ttoe de Sowati soltang dägagän.}      \extran{Sowati salted that crocodile skin.}}}  \sense{\textbf{2.} 
    \definition{bark}}  \sense{\textbf{3.} 
    \definition{to cover with skin}    \example{
      \exsen{Alläp de era gazibra ttoe alle ttoe amallo.}      \extran{Drums are covered with gazibra snakeskin.}}}}

\entry{ttoen}{\headword{ttoen}  \note{Oral Traditions}  \pronnote{ʈ͡ʂoen}
  \pos{Noun}  \sense{\textbf{1.} 
    \definition{story}    \note{Oral Traditions}    \example{
      \exsen{Ngämlle mer dan ttoenttoen därmall a.}      \extran{I like to listen to stories.}}    \example{
}}  \sense{\textbf{2.} 
    \definition{thing}    \note{Oral Traditions}    \example{
      \exsen{Oba tämamae ttoen bällam ngasnen a eragwaeya llame dagwaeya.}      \extran{They did everything together.}}}  \sense{\textbf{3.} 
    \definition{way, method}    \note{Oral Traditions}    \example{
      \exsen{Yu ikllo da ttongo eka ttoen ngangema dan.}      \extran{Smoke is one way of communicating.}}}}

\entry{ttoen~}{\headword{ttoen~}
  \variantof{Derivational reduplicant of \vartext{RED~}}}

\entry{ttoengg}{\headword{ttoengg}  \pronnote{ʈ͡ʂoeŋɡ}
  \pos{Intransitive verb}  \sense{
    \definition{to split up, scatter}    \example{
      \exsen{Ibi naeka peyang ttoengeny e dag.}      \extran{We will split up in tears.}}}}

\entry{ttoenglla}{\headword{ttoenglla}  \pronnote{ʈ͡ʂoeŋɽa}
  \pos{Noun}  \sense{
    \definition{unrelated to one's clan or tribe}    \example{
}    \example{
}    \example{
}}}

\entry{ttoenttoen}{\headword{ttoenttoen}  \pronnote{ʈ͡ʂoenʈ͡ʂoen}
  \pos{Noun}  \sense{
    \definition{small story}    \example{
      \exsen{Känazbag, ngäna mɨnyi bibra sisor ttoenttoen de bɨllɨt.}      \extran{Tomorrow, I will tell you all a new short story.}}}}

\entry{ttoep}{\headword{ttoep}  \note{Snakes}  \pronnote{ʈ͡ʂoep}
  \pos{Noun}  \sense{
    \definition{type of tree}    \note{Snakes}}}

\entry{ttoep}{\headword{ttoep}  \note{Snakes}  \pronnote{ʈ͡ʂoep}
  \pos{Noun}  \sense{
    \definition{type of snake}    \note{Snakes}    \example{
}    \example{
}    \example{
}}}

\entry{ttoettoe}{\headword{ttoettoe}  \pronnote{ʈ͡ʂoeʈ͡ʂoe}
  \pos{Noun}  \sense{
    \definition{blue-faced honeyeater}    \scientificname{Entomyzon cyanotis griseigularis}    \example{
}    \example{
}    \example{
}}}

\entry{ttogottogo}{\headword{ttogottogo}  \pronnote{ʈ͡ʂoɡoʈ͡ʂoɡo}
  \pos{Intransitive coverb}  \sense{
    \definition{to become blunt}}}

\entry{ttokoe}{\headword{ttokoe}
  \variantof{Dialectal Variant of \vartext{ttäkoe}}  \note{Water}  \pronnote{ʈ͡ʂokoe}}

\entry{ttomoll}{\headword{ttomoll}  \pronnote{ʈ͡ʂomoɽ}
  \pos{Modifier}  \sense{
    \definition{deaf}    \example{
      \exsen{Obo llan a ttomoll dageyo.}      \extran{His ears are deaf.}}}}

\entry{ttomttom}{\headword{ttomttom}
  \pos{Noun}  \sense{\textbf{1.} 
    \definition{yam heap}}  \sense{\textbf{2.} 
    \definition{yam cooked whole with skin}}}

\entry{ttonen~}{\headword{ttonen~}
  \variantof{freevarlab of \vartext{RED~}}}

\entry{ttongda}{\headword{ttongda}
  \variantof{Fast Speech Variant of \vartext{ttongdae}}}

\entry{ttongdae}{\headword{ttongdae}  \pronnote{ʈ͡ʂoŋdae}
  \pos{Numeral}  \sense{
    \definition{one (yam counting numeral; also general)}}
  \variants{\Fast Speech Variant \vartext{ttongda}}}

\entry{ttongg}{\headword{ttongg}  \note{Verbs}  \pronnote{ʈ͡ʂoŋɡ}
  \pos{Ditransitive verb}  \sense{
    \definition{to give}    \note{Verbs}    \example{
      \exsen{Ngämlle sawe alle giri de nanttog.}      \extran{Give me the knife with your left hand.}}    \example{
}    \example{
}}}

\entry{ttonggmeny}{\headword{ttonggmeny}  \pronnote{ʈ͡ʂoŋɡmeɲ}
  \pos{Modifier}  \sense{
    \definition{selfish, greedy}}}

\entry{ttonggttongg}{\headword{ttonggttongg}  \note{Trees}  \pronnote{ʈ͡ʂoŋɡʈ͡ʂoŋɡ}
  \pos{Modifier}  \sense{
    \definition{generous, giving}    \note{Trees}}}

\entry{ttongo}{\headword{ttongo}  \pronnote{ʈ͡ʂoŋo}
  \pos{Noun}  \sense{
    \definition{drum handle}}}

\entry{ttongo}{\headword{ttongo}  \pronnote{ʈ͡ʂoŋo}
  \pos{Numeral}  \sense{\textbf{1.} 
    \definition{one}    \example{
      \exsen{Ngämlle ttongo nanttog, a bäne ttongo.}      \extran{Give me one, and one for you.}}}  \sense{\textbf{2.} 
    \definition{a/an}    \example{
      \exsen{Abo bongo ttongo ma sisor de nongo.}      \extran{You must build a new house.}}}  \sense{\textbf{3.} 
    \definition{another}    \example{
      \exsen{Ttongo lla da agbänmällnan a ttongo lla da nindugan obo ingoll me.}      \extran{One man was jumping and another man ran in front of him.}}}  \sense{\textbf{4.} 
    \definition{next}    \example{
      \exsen{Ttongo ebdo me ako llo täräpnan e dalle.}      \extran{The next morning, I went to cut trees.}}}  \sense{\textbf{5.} 
    \definition{unique}    \example{
      \exsen{Bongo imomdae ttongo dan.}      \extran{You are truly unique.}}}}

\entry{ttongo iddob}{\headword{ttongo iddob}  \pronnote{ʈ͡ʂoŋo iɖ͡ʐob}
  \pos{Noun}  \sense{
    \definition{the day after tomorrow}}}

\entry{ttongo lla}{\headword{ttongo lla}  \pronnote{ʈ͡ʂoŋo ɽa}
  \pos{Personal pronoun}  \sense{
    \definition{someone; anyone}    \example{
      \exsen{Ttongo lla da ddone aya bäne peyang ine ma ibi allan?}      \extran{Is anyone going with you to get water?}}    \example{
}    \example{
}}}

\entry{ttongoalle ttongoalle}{\headword{ttongoalle ttongoalle}
  \variantof{Unspecified Variant of \vartext{ttongottongoalle}}}

\entry{ttongottongalle}{\headword{ttongottongalle}
  \variantof{Fast Speech Variant of \vartext{ttongottongoalle}}}

\entry{ttongottongoalle}{\headword{ttongottongoalle}  \pronnote{ʈ͡ʂoŋoʈ͡ʂoŋoaɽe}
  \pos{Adverb}  \sense{\textbf{1.} 
    \definition{each, one by one}    \example{
      \exsen{Tämamae mänmän a llɨg di ttongalle ttongalle ume nɨddɨlaeballo.}      \extran{Every girl kissed each boy.}}    \example{
      \exsen{Wayati de dikom llayabira dänye ttongottongo alle.}      \extran{I brought the watermelons and gave them out to people, one by one.}}}  \sense{\textbf{2.} 
    \definition{few}    \example{
      \exsen{Ddob ttongottongowalle lla da dade ami ddob ttängäm att a guddällmamän gänyme.}      \extran{There are a few others who arrived here from other villages.}}}
  \variants{\Fast Speech Variant \vartext{ttongottongalle}}
  \variants{\Unspecified Variant \vartext{ttongoalle ttongoalle}}
  \variants{\SpellingVariant \vartext{ttongottongowalle}}}

\entry{ttongottongowalle}{\headword{ttongottongowalle}
  \variantof{SpellingVariant of \vartext{ttongottongoalle}}}

\entry{ttongttong}{\headword{ttongttong}  \note{Trees}  \pronnote{ʈ͡ʂoŋʈ͡ʂoŋ}
  \pos{Noun}  \sense{
    \definition{type of tree}    \note{Trees}}}

\entry{ttope}{\headword{ttope}  \note{Water}  \pronnote{ʈ͡ʂope}
  \pos{Noun}  \sense{
    \definition{reed}    \note{Water}    \example{
}    \example{
}    \example{
}}}

\entry{ttotto}{\headword{ttotto}  \pronnote{ʈ͡ʂoʈ͡ʂo}
  \pos{Transitive verb}  \sense{\textbf{1.} 
    \definition{to tie}    \example{
      \exsen{Ngäna ngämo bägäl de ttotto dägneg toboll peyang.}      \extran{I tied my bow and arrows together.}}}  \sense{\textbf{2.} 
    \definition{to collect}    \example{
      \exsen{Bogo ekaklle me gontmonän up gälbänanatt ttonen e.}      \extran{He waited on the ground to collect the fallen bananas.}}}
  \variants{\SpellingVariant \vartext{ttätto}}}

\entry{ttottoem}{\headword{ttottoem}  \lexeme{ttottoem}  \pronnote{ʈ͡ʂoʈ͡ʂoem}
  \pos{Noun}  \sense{
    \definition{type of tree that grows inthe swamp and along creeks with a mango-like, inedible fruit that is yellow when ripe}}}

\entry{ttowa}{\headword{ttowa}  \pronnote{ʈ͡ʂowa}
  \pos{Noun}  \sense{
    \definition{Pacific koel}    \scientificname{Eudynamys orientalis}    \anthronote{its call is imitated to communicate with someone from a distance}}
  \variants{\SpellingVariant \vartext{ttoa}}}

\entry{ttowaemang}{\headword{ttowaemang}
  \pos{Modifier}  \sense{
    \definition{miraculous}}
  \variants{\Fast Speech Variant \vartext{ttowamang}}}

\entry{ttowaemang ttoen}{\headword{ttowaemang ttoen}
  \pos{Noun}  \sense{
    \definition{miracle}    \example{
      \exsen{Ewatta ke gänya täräp me lla bombllo da wasnen anggan ttowamang ttoen tonton ikop e ddapall att?}      \extran{Why does this generation beg to see miracles directly from heaven?}}}}

\entry{ttowamang}{\headword{ttowamang}
  \variantof{Fast Speech Variant of \vartext{ttowaemang}}}

\entry{ttu}{\headword{ttu}  \pronnote{ʈ͡ʂu}
  \pos{Noun}  \sense{
    \definition{deep}    \example{
      \exsen{Angde obom ttu wi dandämoeyän bogo ddone dängllawän, be gonddäwän.}      \extran{When she pushed him into the deep, he did not swim, but drowned.}}}}

\entry{ttuel}{\headword{ttuel}  \pronnote{ʈ͡ʂuel}
  \pos{Noun}  \sense{
    \definition{poor people, no one takes care of them.}}}

\entry{ttullong}{\headword{ttullong}  \pronnote{ʈ͡ʂuɽoŋ}
  \pos{Noun}  \sense{
    \definition{large-tailed nightjar}    \scientificname{Caprimulgus macrurus}    \example{
}    \example{
}    \example{
}}}

\entry{ttupe}{\headword{ttupe}  \pronnote{ʈ͡ʂupe}
  \pos{Modifier}  \sense{
    \definition{bad (of a coconut)}}}

\entry{tu}{\headword{tu}  \pronnote{tu}
  \pos{Numeral}  \sense{
    \definition{two (English numeral)}}
  \variants{\SpellingVariant \vartext{two}}}

\entry{tuba}{\headword{tuba}  \note{See O01-0234 in the EMKP collection}
  \pos{Noun}  \sense{
    \definition{coconut drink traditionally made to welcome people}    \note{See O01-0234 in the EMKP collection}}}

\entry{Tube}{\headword{Tube}  \note{pronounced ttube}  \pronnote{tube}
  \pos{Proper noun}  \sense{
    \definition{male personal name}    \note{pronounced ttube}}}

\entry{tubi}{\headword{tubi}  \pronnote{tubi}
  \pos{Noun}  \sense{
    \definition{type of spear made from sago leaf used to kill birds}}}

\entry{tubu}{\headword{tubu}  \pronnote{tubu}
  \pos{Modifier}  \sense{\textbf{1.} 
    \definition{short}    \example{
      \exsen{llo pallkoll tubu}      \extran{short piece of wood}}}  \sense{\textbf{2.} 
    \definition{end; stump}    \example{
      \exsen{llo tubu, rop tubu}      \extran{tree stump, end of rope}}}}

\entry{tubu~}{\headword{tubu~}
  \variantof{freevarlab of \vartext{RED~}}}

\entry{tubu}{\headword{tubu}  \pronnote{tubu}
  \pos{Noun}  \sense{
    \definition{knee}    \example{
      \exsen{Ngämo yae obo tubu da apte gagäll dan.}      \extran{One of my mother's knees is bad.}}}}

\entry{Tubu}{\headword{Tubu}
  \pos{Proper noun}  \sense{
    \definition{male personal name}}}

\entry{tubukätt}{\headword{tubukätt}  \note{Parts of the Body}  \pronnote{tubukəʈ͡ʂ}
  \pos{Noun}  \sense{
    \definition{kneecap}    \note{Parts of the Body}    \example{
}    \example{
}    \example{
}}}

\entry{tubutubu}{\headword{tubutubu}  \pronnote{tubutubu}
  \pos{Modifier}  \sense{
    \definition{a bit short}    \example{
      \exsen{Ngäna tubutubu dan be ngämo män a tupi dan.}      \extran{I am shorter than my sister.}}    \example{
}    \example{
}}}

\entry{tubutubu}{\headword{tubutubu}
  \pos{Modifier}  \sense{
    \definition{kneeling, on one's knees, on the ground; worshipping}    \example{
      \exsen{Bogo tubutubu gogon, obo ingoll alle ttäle dowae me.}      \extran{He kneeled, his face by her feet.}}    \example{
      \exsen{Zuu lla da Sabat ebdo me sipel gognegnän a Adi pate tubutubu gognegnän.}      \extran{The Jews rested and worshipped God on the Sabbath.}}}}

\entry{tubutubu}{\headword{tubutubu}
  \pos{Modifier}  \sense{
    \definition{nonsingular form of tubu}    \example{
      \exsen{Däräng bo llan a tubutubu dag.}      \extran{The dog's ears are short.}}}}

\entry{tudi}{\headword{tudi}  \pronnote{tudi}
  \pos{Noun}  \sense{\textbf{1.} 
    \definition{fishing}    \example{
      \exsen{Mälla da tudi ma dallän, ako bogo tudi gognän.}      \extran{The woman went to go fishing; then she was fishing.}}}  \sense{\textbf{2.} 
    \definition{fishing rod, fishing line}    \example{
      \exsen{Obo tudi di kollba ulle da danykoeyän.}      \extran{A big fish pulled on his line.}}}}

\entry{tudi ittaenen}{\headword{tudi ittaenen}  \note{Fishing}  \pronnote{tudi iʈ͡ʂaenen}
  \pos{Noun}  \sense{
    \definition{fishing style in which bait is put on fishing lines that are left in the river and checked later}    \note{Fishing}    \example{
}    \example{
}    \example{
}}}

\entry{tudi käp}{\headword{tudi käp}  \note{Diving}  \pronnote{tudi kəp}
  \pos{Noun}  \sense{
    \definition{fishing hook}    \note{Diving}    \example{
      \exsen{Tudi käp sapasapang alle mɨnyi tudi bognegnän.}      \extran{They will fish with various kinds of hooks.}}    \example{
}    \example{
}}}

\entry{tudi tär}{\headword{tudi tär}  \note{Tools}  \pronnote{tudi tər}
  \pos{Noun}  \sense{
    \definition{fishing line}    \note{Tools}    \example{
      \exsen{Tudi tär me guittän.}      \extran{It was caught on the fishing line.}}    \example{
}    \example{
}}}

\entry{tugul}{\headword{tugul}  \note{Trees}  \pronnote{tuɡul}
  \pos{Noun}  \sense{
    \definition{type of big tree that grows in the bush with green, leaflike flowers and straight wood used for house sticks}    \note{Trees}    \example{
}    \example{
}    \example{
}}}

\entry{tuk}{\headword{tuk}  \pronnote{tuk}
  \pos{Noun}  \sense{\textbf{1.} 
    \definition{air}    \example{
      \exsen{Ngäna tuk i pipllugag dan.}      \extran{I fly into the air.}}}  \sense{\textbf{2.} 
    \definition{top}    \example{
      \exsen{Gänya anykeanyke da gänyan tuk mi.}      \extran{This picture is on top.}}}}

\entry{tuk skul}{\headword{tuk skul}  \note{School}  \pronnote{tuk skul}
  \pos{Noun}  \sense{
    \definition{technical university, university}    \note{School}    \example{
}    \example{
}    \example{
}}
  \variants{\SpellingVariant \vartext{tuk sukul}}}

\entry{tuk sukul}{\headword{tuk sukul}
  \variantof{SpellingVariant of \vartext{tuk skul}}}

\entry{tukituki}{\headword{tukituki}  \pronnote{tukituki}
  \pos{Adverb}  \sense{\textbf{1.} 
    \definition{uphill}}  \sense{\textbf{2.} 
    \definition{more than, above, over}    \example{
      \exsen{ttongo taosen tukituki}      \extran{more than one thousand}}}}

\entry{tukpi}{\headword{tukpi}  \pronnote{tukpi}
  \pos{Transitive coverb}  \sense{
    \definition{to heap, pile; gather, collect}    \example{
      \exsen{Bogo mɨnyi bikomän a ttongdae ngättma me tukpi bägnegän.}      \extran{He will bring them over and pile them up in one spot.}}}}

\entry{tulgoe}{\headword{tulgoe}  \pronnote{tulɡoe}
  \pos{Transitive verb}  \sense{
    \definition{to gossip about}    \example{
      \exsen{Tätäm mälla da ngänäm era daitulgoeneyo duli.}      \extran{Yesterday, women were gossiping about me there.}}}}

\entry{tum}{\headword{tum}  \pronnote{tum}
  \pos{Modifier}  \sense{
    \definition{angry}    \example{
      \exsen{Pol bälle enda tum gogon.}      \extran{Paul got very angry.}}}}

\entry{tum}{\headword{tum}  \pronnote{tum}
  \pos{Noun}  \sense{
    \definition{heap}    \example{
      \exsen{Bogo käg me tum dägagän.}      \extran{She put it in a heap in the container.}}}}

\entry{tumang}{\headword{tumang}  \lexeme{tumang}
  \pos{Quantifier}  \sense{
    \definition{plenty, many}    \example{
      \exsen{Ngäna kollba tumang de naittnegan.}      \extran{I caught plenty of fish.}}    \example{
      \exsen{Tumang lla da Ende eka de panypeny erallo.}      \extran{Many people speak the Ende language.}}}}

\entry{tumku}{\headword{tumku}  \pronnote{tumku}
  \pos{Noun}  \sense{
    \definition{back of head}}}

\entry{tumtum}{\headword{tumtum}  \pronnote{tumtum}
  \pos{Transitive/Intransitive coverb}  \sense{
    \definition{to gather}    \example{
      \exsen{Tämamae ddäddäg de tumtum nägagallo.}      \extran{They gathered all the carcasses.}}    \example{
      \exsen{Ngämi tumtum gogmam.}      \extran{We (excl.) gathered.}}}}

\entry{tunen}{\headword{tunen}  \pronnote{tunen}
  \pos{Transitive verb}  \sense{
    \definition{to bump, obligatory plural object (even if singular object)}}}

\entry{tungg}{\headword{tungg}  \pronnote{tuŋɡ}
  \pos{Noun}  \sense{\textbf{1.} 
    \definition{support}}  \sense{\textbf{2.} 
    \definition{to support}}}

\entry{Tungnu}{\headword{Tungnu}  \note{Places around Limol}  \pronnote{tuŋnu}
  \pos{Proper noun}  \sense{
    \definition{Tungnu (toponym)}    \note{Places around Limol}    \example{
}    \example{
}    \example{
}}}

\entry{tunnel}{\headword{tunnel}
  \variantof{Dialectal Variant of \vartext{English loanword}}}

\entry{tupi}{\headword{tupi}  \pronnote{tupi}
  \pos{Modifier}  \sense{\textbf{1.} 
    \definition{tall}    \example{
      \exsen{Da bongo lla tupi di nälläd, bongo mɨnyi llɨg kamebiag ag.}      \extran{If you marry the tall man, you will have many children.}}}  \sense{\textbf{2.} 
    \definition{long}    \example{
      \exsen{Bituri walle da tupi abal dan.}      \extran{The Bituri river is very long.}}}}

\entry{tupi}{\headword{tupi}  \pronnote{tupi}
  \pos{Noun}  \sense{\textbf{1.} 
    \definition{pointer finger, index finger}}  \sense{\textbf{2.} 
    \definition{four (lit. pointer finger; body counting numeral)}    \example{
      \exsen{tupi ollong}      \extran{four times}}}}

\entry{tupitupi}{\headword{tupitupi}  \pronnote{tupitupi}
  \pos{Modifier}  \sense{
    \definition{nonsingular form of tupi}    \example{
      \exsen{Ngämo llan a tubutubu agnegän, obo da tupitupi.}      \extran{My ears are short; his are long.}}}
  \variants{\Fast Speech Variant \vartext{tutupi, tutpi}}}

\entry{tupol}{\headword{tupol}  \note{Traditional Arrows}  \pronnote{tupol}
  \pos{Noun}  \sense{
    \definition{type of spear}    \note{Traditional Arrows}    \example{
}    \example{
}    \example{
}}}

\entry{turik}{\headword{turik}  \note{Tools}  \pronnote{turik}
  \pos{Noun}  \sense{
    \definition{axe}    \note{Tools}    \example{
      \exsen{Bäne turik de abo iweny gall tɨt i.}      \extran{Bring your axe to hollow out the tree.}}    \example{
}    \example{
}}}

\entry{turku}{\headword{turku}
  \pos{Noun}  \sense{
    \definition{thinner piece used with bore to smoke tobacco}}}

\entry{turllo}{\headword{turllo}  \note{Insects}  \pronnote{turɽo}
  \pos{Noun}  \sense{
    \definition{type of lizard}    \note{Insects}    \example{
}    \example{
}    \example{
}}}

\entry{turwe}{\headword{turwe}  \note{Birds}  \pronnote{turwe}
  \pos{Noun}  \sense{
    \definition{shining bronze cuckoo}    \scientificname{Chrysococcyx lucidus}    \note{Birds}    \example{
}    \example{
}    \example{
}}}

\entry{tusde}{\headword{tusde}  \note{Part of the week}  \pronnote{tusde}
  \pos{Noun}  \sense{
    \definition{Tuesday}    \note{Part of the week}    \example{
}    \example{
}    \example{
}}}

\entry{tutpi}{\headword{tutpi}
  \variantof{Fast Speech Variant of \vartext{tupitupi}}  \pronnote{tutpi}}

\entry{tutu}{\headword{tutu}  \pronnote{tutu}
  \pos{Transitive verb}  \sense{
    \definition{to knock fruit continuously}}}

\entry{tutu}{\headword{tutu}
  \variantof{Baby Talk Variant of \vartext{susu}}  \pronnote{tutu}}

\entry{tutu}{\headword{tutu}  \pronnote{tutu}
  \pos{Noun}  \sense{\textbf{1.} 
    \definition{mountain, hill}    \example{
      \exsen{Ngängälätt dädär a mänyi tutu atta bolldaeyän kopek e.}      \extran{The round rock will roll from the hill into the valley.}}}  \sense{\textbf{2.} 
    \definition{land}    \example{
      \exsen{Tutu wi darullgoeya käza de gull peyang.}      \extran{We dragged the crocodile onto land with the net.}}}}

\entry{tutuang}{\headword{tutuang}
  \pos{Modifier}  \sense{
    \definition{steep}}}

\entry{tutuaram}{\headword{tutuaram}  \note{Types of Taro}  \pronnote{tutuaram}
  \pos{Noun}  \sense{
    \definition{type of taro}    \note{Types of Taro}    \example{
}    \example{
}    \example{
}}}

\entry{Tutuli}{\headword{Tutuli}  \pronnote{tutuli}
  \pos{Proper noun}  \sense{
    \definition{male personal name}}}

\entry{tutupi}{\headword{tutupi}
  \variantof{Fast Speech Variant of \vartext{tupitupi}}}

\entry{tutututu}{\headword{tutututu}
  \pos{Noun}  \sense{
    \definition{nonsingular form of tutu}    \example{
      \exsen{Tutututu da dag.}      \extran{There are mountains.}}}}

\entry{tuwenti}{\headword{tuwenti}
  \variantof{SpellingVariant of \vartext{twenti}}}

\entry{tuwetuwe}{\headword{tuwetuwe}  \note{Trees}  \pronnote{tuwetuwe}
  \pos{Noun}  \sense{
    \definition{type of small tree that grows in the bush with white flowers and edible red fruit}    \note{Trees}    \example{
}    \example{
}    \example{
}}}

\entry{tuwi}{\headword{tuwi}  \note{Trees}  \pronnote{tuwi}
  \pos{Noun}  \sense{
    \definition{type of tree with white flowers, yellow fruit, and wood used for posts}    \note{Trees}    \example{
}    \example{
}    \example{
}}}

\entry{tuwok}{\headword{tuwok}  \note{Trees}  \pronnote{tuwok}
  \pos{Noun}  \sense{
    \definition{type of tree}    \note{Trees}    \example{
}    \example{
}    \example{
}}}

\entry{tuyem}{\headword{tuyem}  \pronnote{tujem}
  \pos{Intransitive coverb}  \sense{
    \definition{to make a loud noise}}}

\entry{Tuyu}{\headword{Tuyu}
  \pos{Proper noun}  \sense{
    \definition{male personal name}}}

\entry{twelb}{\headword{twelb}
  \pos{Numeral}  \sense{
    \definition{twelve (English numeral)}}
  \variants{\SpellingVariant \vartext{twelve}}}

\entry{twelve}{\headword{twelve}
  \variantof{SpellingVariant of \vartext{twelb}}}

\entry{twenti}{\headword{twenti}
  \pos{Numeral}  \sense{
    \definition{twenty}}
  \variants{\SpellingVariant \vartext{tuwenti}}}

\entry{two}{\headword{two}
  \variantof{SpellingVariant of \vartext{tu}}}

\chapter{U}


\entry{Uba}{\headword{Uba}
  \pos{Proper noun}  \sense{
    \definition{female personal name}}}

\entry{Ubäd}{\headword{Ubäd}
  \pos{Proper noun}  \sense{
    \definition{male personal name}}}

\entry{ubemang}{\headword{ubemang}  \pronnote{ubemaŋ}
  \pos{Modifier}  \sense{
    \definition{wide}    \example{
      \exsen{Däbe abo obo pallall ngänttäg eka da ulle ubemang gogon.}      \extran{That story about his arrival spread wide and far.}}}}

\entry{ubi}{\headword{ubi}  \pronnote{ubi}
  \pos{Personal pronoun}  \sense{
    \definition{they (third person nonsingular pronoun, nominative form)}}
  \variants{\Dialectal Variant \vartext{obi}}}

\entry{ubim}{\headword{ubim}  \pronnote{ubim}
  \pos{Personal pronoun}  \sense{
    \definition{accusative form of ubi}    \example{
      \exsen{Ubim mɨnyi umllang bägaebeya.}      \extran{We will inform them.}}}}

\entry{ubira}{\headword{ubira}
  \pos{Personal pronoun}  \sense{
    \definition{dative form of ubi}    \example{
      \exsen{Bogo ddäddäg de dokonegän ubira.}      \extran{He cut the meat for them.}}}}

\entry{ubony}{\headword{ubony}  \lexeme{ubony}  \pronnote{uboɲ}
  \pos{Noun}  \sense{
    \definition{type of black bee}}}

\entry{Ubrag}{\headword{Ubrag}  \pronnote{ubraɡ}
  \pos{Proper noun}  \sense{
    \definition{male personal name}}}

\entry{ubrattäka}{\headword{ubrattäka}  \lexeme{ubrattäka}  \pronnote{ubraʈ͡ʂəka}
  \pos{Noun}  \sense{
    \definition{type of yam a red interior}}}

\entry{ud}{\headword{ud}  \pronnote{ud}
  \pos{Noun}  \sense{
    \definition{door; gate}    \example{
      \exsen{Ud a papekang dan.}      \extran{The door is closed.}}}
  \variants{\SpellingVariant \vartext{uud, wud, wäd}}}

\entry{udab}{\headword{udab}  \pronnote{udab}
  \pos{Transitive verb}  \sense{\textbf{1.} 
    \definition{to disappear, get lost, go missing}    \example{
      \exsen{Ngäna udab allan.}      \extran{I am getting lost.}}    \example{
      \exsen{Angde ubi gongoseyo, ubi ikop dägageyo mamam mätta da Meri erem de dibeyän audaban.}      \extran{When they came back, they saw that the red yam that Mary had planted was gone.}}}  \sense{\textbf{2.} 
    \definition{to lose}    \example{
      \exsen{Ubi obom daudabeyo.}      \extran{They lost him.}}}}

\entry{udaude}{\headword{udaude}  \pronnote{udaude}
  \pos{Transitive verb}  \sense{
    \definition{to light, start (a fire)}    \example{
      \exsen{Ubi yu de daudeaemeyo.}      \extran{Then started some fires.}}}
  \variants{\Dialectal Variant \vartext{odaode}}}

\entry{udu}{\headword{udu}  \pronnote{udu}
  \pos{Noun}  \sense{
    \definition{walking stick, cane, staff}    \example{
      \exsen{Bogo udu peyang ibiag dan.}      \extran{She walks with a cane.}}}
  \variants{\Dialectal Variant \vartext{wudu}}}

\entry{ugeuge}{\headword{ugeuge}  \note{Trees}  \pronnote{uɡeuɡe}
  \pos{Noun}  \sense{
    \definition{type of tree}    \note{Trees}    \example{
}    \example{
}    \example{
}}}

\entry{ugri}{\headword{ugri}
  \pos{Noun}  \sense{
    \definition{fever}}}

\entry{ugug}{\headword{ugug}  \pronnote{uɡuɡ}
  \pos{Transitive verb}  \sense{
    \definition{to make mumu}    \example{
      \exsen{Ngämi ttägäll daugeya.}      \extran{We made mumu.}}}}

\entry{uk}{\headword{uk}  \pronnote{uk}
  \pos{Intransitive coverb}  \sense{
    \definition{to shout, yell, cry out}    \example{
      \exsen{Ubi nga wukangae anggan.}      \extran{They are still shouting.}}    \example{
      \exsen{Mälla da uk allan tätäp atta.}      \extran{The woman is crying out in pain.}}}
  \variants{\Unspecified Variant \vartext{wäk}}
  \variants{\SpellingVariant \vartext{wuk}}}

\entry{ukär}{\headword{ukär}  \pronnote{ukər}
  \pos{Noun}  \sense{\textbf{1.} 
    \definition{glossy-mantled manucode}    \scientificname{Manucodia ater ater}}  \sense{\textbf{2.} 
    \definition{trumpet manucode}    \scientificname{Phonygammus keraudrenii}}}

\entry{ukeuke}{\headword{ukeuke}  \pronnote{u'keu'ke}
  \pos{Intransitive coverb}  \sense{\textbf{1.} 
    \definition{to flatten (grass)}    \example{
      \exsen{Bogo towall de ukeuke allan.}      \extran{He's flattening the grass.}}}  \sense{\textbf{2.} 
    \definition{to squeeze and tie grass above something to mark ownership}}}

\entry{ulle~}{\headword{ulle~}
  \variantof{Plural reduplicant of \vartext{RED~}}}

\entry{ulle}{\headword{ulle}  \note{Adjectives (Swadesh)}  \pronnote{uɽe}
  \pos{Modifier}  \sense{\textbf{1.} 
    \definition{big, large, great}    \note{Adjectives (Swadesh)}    \example{
      \exsen{Yu ttänttämatt a ai dan ute ulle we bogon.}      \extran{Burns can become big sores.}}}  \sense{\textbf{2.} 
    \definition{important, great}    \note{Adjectives (Swadesh)}    \example{
      \exsen{Tämamae Adi bo gwell makäp me, eran ke ngattong abal e ulle abal gwell a?}      \extran{Within all of God's commandments, which commandment is the most important?}}}  \sense{\textbf{3.} 
    \definition{long; tall}    \note{Adjectives (Swadesh)}    \example{
      \exsen{pinggudd ulle}      \extran{long skirt}}}  \sense{\textbf{4.} 
    \definition{a lot, abundant, plentiful (for uncountable nouns)}    \note{Adjectives (Swadesh)}    \example{
      \exsen{Ine da ulle dan.}      \extran{There is plenty of water.}}    \example{
      \exsen{Eso ulle.}      \extran{Many thanks.}}}  \sense{\textbf{5.} 
    \definition{entire, whole}    \note{Adjectives (Swadesh)}    \example{
      \exsen{ebdo ulle, iddob ulle}      \extran{all day, all night}}}  \sense{\textbf{6.} 
    \definition{to become big, grow up}    \note{Adjectives (Swadesh)}    \example{
      \exsen{Ngäna gänyme ulle gog.}      \extran{I grew up here.}}}  \sense{\textbf{7.} 
    \definition{to raise, rear}    \note{Adjectives (Swadesh)}    \example{
      \exsen{Aeya bam ulle dagän?}      \extran{Who raised you?}}}}

\entry{ulle binang}{\headword{ulle binang}
  \pos{Noun}  \sense{
    \definition{master, owner, ruler, important person}    \example{
      \exsen{Da ttongo bina makäp me ulle binang e moko bogon, ede ai dan bogo melemang bogon.}      \extran{If one of you wants to be a ruler, that person can be a servant instead.}}}}

\entry{ulle kottllam}{\headword{ulle kottllam}  \pronnote{uɽe koʈ͡ʂɽam}
  \pos{Noun}  \sense{
    \definition{type of big turtle}}}

\entry{ullegäll}{\headword{ullegäll}  \note{Trees}  \pronnote{uɽeɡəɽ}
  \pos{Noun}  \sense{
    \definition{type of tree that grows in the grassland with white flowers, black fruits, and red nuts}    \note{Trees}    \example{
}    \example{
}    \example{
}}}

\entry{ulleulle}{\headword{ulleulle}  \pronnote{uɽeuɽe}
  \pos{}  \sense{
    \definition{nonsingular form of ulle}    \example{
      \exsen{Obo mälläng ik a ulleulle abal dagaya.}      \extran{His nostrils were very big.}}    \example{
      \exsen{Ngämi ulleulle agmalla.}      \extran{We (excl.) grew up.}}}}

\entry{ulleulle binang}{\headword{ulleulle binang}
  \pos{Noun}  \sense{
    \definition{nonsingular form of ulle binang}}}

\entry{ullowae}{\headword{ullowae}  \pronnote{uɽowae}
  \pos{Adverb}  \sense{
    \definition{fast, quickly}    \example{
      \exsen{Däräng ullowae dinduag daeya.}      \extran{The dog was running quickly.}}}}

\entry{ullull}{\headword{ullull}  \pronnote{uɽuɽ}
  \pos{Transitive verb}  \sense{
    \definition{to cross over}    \example{
      \exsen{Ubi tätäm ttongo mo sisor de daulliyu.}      \extran{Yesterday, they crossed over a new bridge.}}}}

\entry{umaem}{\headword{umaem}  \pronnote{umaem}
  \pos{Intransitive verb}  \sense{\textbf{1.} 
    \definition{to gather}    \example{
      \exsen{Lla da gumaemalle a eka gontemenyalle täre pallall e.}      \extran{The people will gather and discuss the feast.}}}  \sense{\textbf{2.} 
    \definition{to gather, collect}    \example{
      \exsen{Ngäna towall de umaem anggan komlla tum.}      \extran{I gathered the grass into two heaps.}}}}

\entry{Umbuzag}{\headword{Umbuzag}
  \pos{Proper noun}  \sense{
    \definition{male personal name (name of the original Agob man)}}}

\entry{ume}{\headword{ume}
  \pos{Noun}  \sense{
}}

\entry{ume ddäddäl}{\headword{ume ddäddäl}  \pronnote{ume ɖ͡ʐəɖ͡ʐəl}
  \pos{Transitive compound verb}  \sense{
    \definition{to kiss}    \anthronote{No special word for romantic kiss since that is not custom. Only kiss children. Young people do kiss now though.}    \example{
      \exsen{Bina mokowang lla da gänyan ngäna erem ume bäddäl.}      \extran{The man you all want is the one that I will kiss.}}}}

\entry{ume nyanyu}{\headword{ume nyanyu}  \pronnote{ume ɲaɲu}
  \pos{Noun}  \sense{\textbf{1.} 
    \definition{type of string game, imitation game}    \example{
}    \example{
}    \example{
}}  \sense{\textbf{2.} 
    \definition{to imitate}}}

\entry{ume ttäp}{\headword{ume ttäp}
  \pos{Noun}  \sense{
    \definition{mouth}    \example{
      \exsen{Ngäna däbe giri tupi da ngämo ereya ada käza bo ume ttäp e däträkne.}      \extran{I had my long knife and pushed it into the crocodile's mouth.}}}
  \variants{\SpellingVariant \vartext{umettäp}}}

\entry{umettäp}{\headword{umettäp}
  \variantof{SpellingVariant of \vartext{ume ttäp}}  \pronnote{umeʈ͡ʂəp}}

\entry{umllang}{\headword{umllang}  \pronnote{umɽaŋ}
  \pos{Noun}  \sense{\textbf{1.} 
    \definition{knowledge, knowing, awareness}    \example{
      \exsen{Bäne umllang dan ada ngämo umllang dan pittpitt a.}      \extran{You know that I know how to sew (lit. your knowledge is that sewing is my knowledge).}}}  \sense{\textbf{2.} 
    \definition{to know; come to know, learn}    \example{
      \exsen{Ddone aya umllang gogon ada erowattäm da ge za da gogezän.}      \extran{No one knows from where these things came out.}}    \example{
      \exsen{Ngäna mäse umllang gogne orpmang eka de.}      \extran{I tried to learn the Wipi language.}}}  \sense{\textbf{3.} 
    \definition{to tell, inform}    \example{
      \exsen{Ngäna umllang däga, "Imanuel, bongo gall guwo me, ngäna walle we gäbän allan ddia llädäd e."}      \extran{I told him, "Imanuel, you stay in the canoe; I am jumping in the water to grab the deer."}}    \example{
      \exsen{Ngäna balle ngämo nag bom umllang bägag bablle ngämingg e.}      \extran{I will go tell my friend to help you.}}}}

\entry{umllang bällanen ma skul}{\headword{umllang bällanen ma skul}  \note{Seasons}  \pronnote{umɽaŋ bəɽanen ma skul}
  \pos{Noun}  \sense{
    \definition{vocational school}    \note{Seasons}    \example{
}    \example{
}    \example{
}}}

\entry{umllang eka}{\headword{umllang eka}
  \pos{Noun}  \sense{
    \definition{summary}}}

\entry{umull}{\headword{umull}
  \pos{Noun}  \sense{
    \definition{wisdom}}}

\entry{uniform}{\headword{uniform}
  \variantof{Dialectal Variant of \vartext{English loanword}}}

\entry{unkäm}{\headword{unkäm}
  \variantof{Unspecified Variant of \vartext{wänkäm}}}

\entry{until}{\headword{until}
  \variantof{Dialectal Variant of \vartext{English loanword}}}

\entry{up}{\headword{up}  \lexeme{up}  \pronnote{up}
  \pos{Noun}  \sense{
    \definition{banana}    \scientificname{Musa}    \example{
      \exsen{wup däg}      \extran{banana bunch}}    \example{
      \exsen{wup käp}      \extran{banana fruit}}    \example{
      \exsen{wup ttam}      \extran{banana leaf}}}
  \variants{\freevarlabs \vartext{wup, wäp}}}

\entry{upe}{\headword{upe}  \pronnote{upe}
  \pos{Locational}  \sense{
    \definition{outside, out}    \example{
      \exsen{Meri uziz allan upe we.}      \extran{Mary runs out.}}}
  \variants{\freevarlab \vartext{ope}}}

\entry{upeupe}{\headword{upeupe}  \pronnote{upeupe}
  \pos{Noun}  \sense{
    \definition{type of plant with a single stem and edible, tall fruit near the base of stem}    \example{
}    \example{
}    \example{
}}}

\entry{Upiara}{\headword{Upiara}  \pronnote{upiara}
  \pos{Proper noun}  \sense{
    \definition{Upiara (Bitur-speaking village in Oriomo-Bituri Rural LLG; GPS: -8.547170, 142.653008)}    \example{
}    \example{
}    \example{
}}}

\entry{upiye}{\headword{upiye}  \note{Source: Jack Dipa}
  \pos{Noun}  \sense{
    \definition{type of tree used to make kwib charcoal}    \note{Source: Jack Dipa}}}

\entry{upma}{\headword{upma}  \note{Friendship/social terms of endearment}  \pronnote{upma}
  \pos{Noun}  \sense{
    \definition{two friends who share a twin banana fruit}    \note{Friendship/social terms of endearment}    \example{
}    \example{
}    \example{
}}}

\entry{upoupoll}{\headword{upoupoll}  \note{Trees}  \pronnote{upoupoɽ}
  \pos{Noun}  \sense{
    \definition{type of tree}    \note{Trees}    \example{
}    \example{
}    \example{
}}}

\entry{upye}{\headword{upye}  \note{Trees}  \pronnote{upje}
  \pos{Noun}  \sense{
    \definition{type of tree with white flowers and black fruit that produces a black pigment}    \note{Trees}    \example{
}    \example{
}    \example{
}}}

\entry{Ur}{\headword{Ur}  \pronnote{ur}
  \pos{Proper noun}  \sense{
    \definition{Ur (toponym)}}
  \variants{\SpellingVariant \vartext{Wur}}}

\entry{ur yogoll}{\headword{ur yogoll}  \note{?}  \pronnote{ur joɡoɽ}
  \pos{Noun}  \sense{
    \definition{flood time, rainy season}    \note{?}    \example{
}    \example{
}    \example{
}}}

\entry{uriar}{\headword{uriar}  \pronnote{uriar}
  \pos{Noun}  \sense{
    \definition{type of palm with coconuts with a purple exocarp}}}

\entry{Urimba}{\headword{Urimba}
  \pos{Proper noun}  \sense{
    \definition{Urimba (toponym)}}}

\entry{Uroe}{\headword{Uroe}
  \pos{Proper noun}  \sense{
    \definition{Wuroi (in Oriomo-Bituri Rural LLG)}}
  \variants{\SpellingVariant \vartext{Wuroi}}}

\entry{uta}{\headword{uta}  \pronnote{uta}
  \pos{Intransitive verb}  \sense{
    \definition{go (imperative, singular form)}    \example{
}}}

\entry{utale}{\headword{utale}  \note{Adjectives (Swadesh)}  \pronnote{utale}
  \pos{Locational}  \sense{
    \definition{far}    \note{Adjectives (Swadesh)}    \example{
      \exsen{Alla utale agan ma da?}      \extran{How far is the village?}}    \example{
}    \example{
}}
  \variants{\SpellingVariant \vartext{utali}}}

\entry{utali}{\headword{utali}
  \variantof{SpellingVariant of \vartext{utale}}}

\entry{utamom}{\headword{utamom}
  \pos{Intransitive verb}  \sense{
    \definition{plural form of uta}    \example{
      \exsen{Bibi utamom ngattong.}      \extran{You all, go first.}}}}

\entry{ute}{\headword{ute}  \pronnote{ute}
  \pos{Noun}  \sense{
    \definition{sore, blister; wound}    \example{
      \exsen{Yu ttänttämatt a ai dan ute ulle we bogon.}      \extran{Burns can become big sores.}}}
  \variants{\SpellingVariant \vartext{wute}}
  \variants{\freevarlab \vartext{wäte}}}

\entry{ute mälanenang}{\headword{ute mälanenang}
  \pos{Noun}  \sense{
    \definition{healthcare worker}}}

\entry{uteute}{\headword{uteute}
  \pos{Modifier}  \sense{
    \definition{with sores}}}

\entry{uteute ma}{\headword{uteute ma}  \pronnote{uteute ma}
  \pos{Noun}  \sense{
    \definition{aid post}    \example{
}    \example{
}    \example{
}}}

\entry{utt}{\headword{utt}  \pronnote{uʈ͡ʂ}
  \pos{Noun}  \sense{
    \definition{shoot (of a plant)}    \example{
      \exsen{Duny wätätatt a nge wutt a gokenyamän.}      \extran{After the beetle eats the shoot, it bends.}}}
  \variants{\SpellingVariant \vartext{wutt}}}

\entry{utt}{\headword{utt}  \note{Dancing}  \pronnote{uʈ͡ʂ}
  \pos{Noun}  \sense{
    \definition{conch shell}    \note{Dancing}    \example{
      \exsen{Da Llimoll me lla da kuddäll bogon, komlla o kumuddäga lla da mɨnyi utt alle eka de bängaeyo.}      \extran{If someone in Limol dies, two or three men will deliver the news using a conch shell.}}    \example{
}    \example{
}}
  \variants{\SpellingVariant \vartext{wutt}}}

\entry{uttang ttatta}{\headword{uttang ttatta}  \note{Sago palms}  \pronnote{uʈ͡ʂaŋ ʈ͡ʂaʈ͡ʂa}
  \pos{Noun}  \sense{
    \definition{type of sago}    \note{Sago palms}    \example{
}    \example{
}    \example{
}}}

\entry{uttutt}{\headword{uttutt}  \pronnote{uʈ͡ʂuʈ͡ʂ}
  \pos{Noun}  \sense{
    \definition{small conch shell}}}

\entry{uwo}{\headword{uwo}  \pronnote{uwo}
  \pos{Noun}  \sense{
    \definition{magnificent riflebird}    \scientificname{Ptiloris magnificus}    \example{
}    \example{
}    \example{
}}}

\entry{uwo kottllam}{\headword{uwo kottllam}
  \pos{Noun}  \sense{
    \definition{type of turtle}}}

\entry{Uzaba}{\headword{Uzaba}  \pronnote{uzaba}
  \pos{Proper noun}  \sense{
    \definition{male personal name}}}

\entry{Uziag}{\headword{Uziag}  \pronnote{uziaɡ}
  \pos{Proper noun}  \sense{
    \definition{male personal name}}}

\entry{uziz}{\headword{uziz}
  \pos{Intransitive verb}  \sense{
    \definition{to run}}}

\chapter{UU}


\entry{uud}{\headword{uud}
  \variantof{SpellingVariant of \vartext{ud}}}

\chapter{V}


\entry{Vanessa}{\headword{Vanessa}  \pronnote{vanessa}
  \pos{Proper noun}  \sense{
    \definition{female personal name}}}

\entry{very}{\headword{very}
  \variantof{Dialectal Variant of \vartext{English loanword}}}

\entry{vice-chairman}{\headword{vice-chairman}
  \pos{Nominal}  \sense{
    \definition{vice-chairman}}}

\entry{Victoria}{\headword{Victoria}  \pronnote{viktoria}
  \pos{Proper noun}  \sense{
    \definition{female personal name}}}

\entry{Vincent}{\headword{Vincent}  \pronnote{vinsent}
  \pos{Proper noun}  \sense{
    \definition{male personal name}}}

\entry{vocational}{\headword{vocational}
  \variantof{Dialectal Variant of \vartext{English loanword}}}

\entry{volleyball}{\headword{volleyball}
  \pos{Noun}  \sense{
    \definition{volleyball}}}

\chapter{W}


\entry{wa}{\headword{wa}  \pronnote{wa}
  \pos{Noun}  \sense{
    \definition{penis}}}

\entry{=wa}{\headword{=wa}
  \pos{Nominal enclitic}  \sense{
    \definition{emphatic clitic}    \example{
      \exsen{Ngämo baba wa ngänäm dantemenyän.}      \extran{My dad, he taught me.}}}}

\entry{wa piye}{\headword{wa piye}
  \pos{Noun}  \sense{
    \definition{sperm}}}

\entry{wabeb}{\headword{wabeb}  \pronnote{wabeb}
  \pos{Transitive verb}  \sense{
    \definition{to beat, smash, pound}    \example{
      \exsen{Mälla da biye de duwebnegän.}      \extran{The woman pounded the taro.}}    \example{
      \exsen{Ttängäm mondrog lla da melemang de dällädeyo, gagagäll duwebeyo.}      \extran{The gardeners grabbed the servant and beat him badly.}}}}

\entry{wabeyawabeya}{\headword{wabeyawabeya}  \note{Trees}  \pronnote{'wabeya'wabeya}
  \pos{Noun}  \sense{
    \definition{type of tree}    \note{Trees}    \example{
}    \example{
}    \example{
}}}

\entry{wäd}{\headword{wäd}
  \variantof{SpellingVariant of \vartext{ud}}}

\entry{wadär}{\headword{wadär}  \pronnote{wadər}
  \pos{Noun}  \sense{\textbf{1.} 
    \definition{type of grass; cane}    \example{
      \exsen{Lla da da bogäddnän, mulldae dan wadär de bängällbänän, wadär alle bäbäddän.}      \extran{If a person is fighting, it is possible to get a cane and beat them with it.}}}  \sense{\textbf{2.} 
    \definition{bowstring}    \example{
      \exsen{Bägäl wadär a gottäpanän.}      \extran{The bowstring snapped.}}}}

\entry{wadär gullem}{\headword{wadär gullem}  \lexeme{wadär gullem}  \pronnote{wadər ɡuɽem}
  \pos{Noun}  \sense{
    \definition{Papuan python}    \scientificname{Apodora papuana}}}

\entry{wadär käp}{\headword{wadär käp}  \note{Types of Taro}  \pronnote{wadər kəp}
  \pos{Noun}  \sense{
    \definition{type of big taro}    \note{Types of Taro}    \example{
}    \example{
}    \example{
}}}

\entry{wadär mikel}{\headword{wadär mikel}  \note{Hunting}  \pronnote{wadər mikel}
  \pos{Noun}  \sense{
    \definition{extra bowstring}    \note{Hunting}    \example{
}    \example{
}    \example{
}}}

\entry{Wadär Mitang}{\headword{Wadär Mitang}  \note{Places around Limol}  \pronnote{wadər mitaŋ}
  \pos{Proper noun}  \sense{
    \definition{Wadär Mitang (toponym)}    \note{Places around Limol}    \example{
}    \example{
}    \example{
}}}

\entry{wadär nyongkoe}{\headword{wadär nyongkoe}  \note{Games}  \pronnote{wadər ɲoŋkoe}
  \pos{Noun}  \sense{
    \definition{type of game involving pulling cane}    \note{Games}    \example{
}    \example{
}    \example{
}}}

\entry{wadär pitt}{\headword{wadär pitt}  \note{Hunting}  \pronnote{wadər piʈ͡ʂ}
  \pos{Noun}  \sense{
    \definition{plant from cane, used on both sides of bow}    \note{Hunting}    \example{
}    \example{
}    \example{
}}}

\entry{wädɨwädɨg}{\headword{wädɨwädɨg}  \note{Trees}  \pronnote{wədɪwədɪɡ}
  \pos{Noun}  \sense{
    \definition{type of tree}    \note{Trees}    \example{
}    \example{
}    \example{
}}}

\entry{wädwäd}{\headword{wädwäd}  \note{Trees}  \pronnote{wədwəd}
  \pos{Noun}  \sense{
    \definition{type of tree}    \note{Trees}    \example{
}    \example{
}    \example{
}}}

\entry{Wae}{\headword{Wae}
  \pos{Proper noun}  \sense{
    \definition{female personal name}}}

\entry{Waedar}{\headword{Waedar}
  \pos{Proper noun}  \sense{
    \definition{female personal name}}}

\entry{Waenum}{\headword{Waenum}
  \pos{Proper noun}  \sense{
    \definition{female personal name}}}

\entry{waeo}{\headword{waeo}
  \variantof{SpellingVariant of \vartext{waeyo}}}

\entry{waetwaet}{\headword{waetwaet}  \note{Trees}  \pronnote{waetwaet}
  \pos{Noun}  \sense{
    \definition{type of tree}    \note{Trees}    \example{
}    \example{
}    \example{
}}}

\entry{waewae}{\headword{waewae}
  \pos{Noun}  \sense{
    \definition{song sung while beating sago}}}

\entry{waeya}{\headword{waeya}  \pronnote{waeja}
  \pos{Noun}  \sense{
    \definition{wire}    \example{
      \exsen{Ngäna nyäng tär ngasnges e nyäng de napoenan waeya kälsre alle.}      \extran{I used the small wire to make holes in the bag to put the new bag string.}}}}

\entry{waeyo}{\headword{waeyo}  \pronnote{waejo}
  \pos{Interjection}  \sense{
    \definition{shout made in distress}}
  \variants{\SpellingVariant \vartext{waeo}}}

\entry{Wagälla}{\headword{Wagälla}  \note{Places around Limol}  \pronnote{waɡəɽa}
  \pos{Proper noun}  \sense{
    \definition{Wagälla (toponym)}    \note{Places around Limol}    \example{
}    \example{
}    \example{
}}}

\entry{wägba}{\headword{wägba}  \note{Trees}  \pronnote{wəɡba}
  \pos{Noun}  \sense{
    \definition{type of tree that grows in the bush with white flowers, bark used as medicine, and strong wood used for posts; helps make dogs' noses more sensitive}    \note{Trees}    \example{
}    \example{
}    \example{
}}}

\entry{Wagiba}{\headword{Wagiba}  \pronnote{waɡiba}
  \pos{Proper noun}  \sense{
    \definition{female personal name}}}

\entry{Wagiya}{\headword{Wagiya}
  \pos{Proper noun}  \sense{
    \definition{personal name}}}

\entry{waglla}{\headword{waglla}  \pronnote{waɡɽa}
  \pos{Noun}  \sense{
    \definition{bullroarer}    \anthronote{sacred object, used in men's initiation ceremonies. Legend: It was a musical instrument that hung from beneath the skirts of the first woman and made a beautiful noise while she swept. The first man stole it from her, yanking it from her, and starting woman's first menstrual cycle. The bull roarer was kept secret from women and children; they were told the loud sound was giants.}    \example{
      \exsen{waglla eka}      \extran{bullroarer signal}}    \example{
}    \example{
}}}

\entry{Wai}{\headword{Wai}  \pronnote{wai}
  \pos{Proper noun}  \sense{
    \definition{female personal name}}}

\entry{Wainum}{\headword{Wainum}  \pronnote{wainum}
  \pos{Proper noun}  \sense{
    \definition{female personal name}}}

\entry{wäk}{\headword{wäk}
  \variantof{Unspecified Variant of \vartext{uk}}}

\entry{wak}{\headword{wak}  \pronnote{wak}
  \pos{Noun}  \sense{
    \definition{Papuan pitta}    \scientificname{Erythropitta macklotii}    \example{
}    \example{
}    \example{
}}}

\entry{Waka Källäm}{\headword{Waka Källäm}  \note{Places around Limol}  \pronnote{waka kəɽəm}
  \pos{Proper noun}  \sense{
    \definition{Waka Pond (in Limol)}    \note{Places around Limol}    \example{
}    \example{
}    \example{
}}}

\entry{wäkär}{\headword{wäkär}  \note{Birds}  \pronnote{wəkər}
  \pos{Noun}  \sense{
    \definition{type of bird}    \note{Birds}    \example{
}    \example{
}    \example{
}}}

\entry{wakata}{\headword{wakata}  \note{Names of bananas}  \pronnote{wakata}
  \pos{Noun}  \sense{
    \definition{type of introduced banana}    \note{Names of bananas}    \example{
}    \example{
}    \example{
}}}

\entry{wäkɨs}{\headword{wäkɨs}  \note{Birds}  \pronnote{wəkɪs}
  \pos{Noun}  \sense{
    \definition{type of bird}    \note{Birds}    \example{
}    \example{
}    \example{
}}}

\entry{wakllu}{\headword{wakllu}  \pronnote{wakɽu}
  \pos{Modifier}  \sense{
    \definition{rough}}}

\entry{wäl}{\headword{wäl}  \pronnote{wəl}
  \pos{Noun}  \sense{
    \definition{main surface-level stem of a plant with rhizomes (e.g. sweet potato, lemongrass)}    \example{
      \exsen{Kak bo nae a wäl peyang gognegän.}      \extran{Grandmother's sweet potato plants have stems.}}}
  \variants{\Unspecified Variant \vartext{ol}}}

\entry{Wala}{\headword{Wala}  \pronnote{wala}
  \pos{Proper noun}  \sense{
    \definition{personal name}}}

\entry{walap}{\headword{walap}  \note{Clothing}  \pronnote{walap}
  \pos{Noun}  \sense{
    \definition{cap}    \note{Clothing}    \example{
}    \example{
}    \example{
}}}

\entry{wälep}{\headword{wälep}  \note{Trees}  \pronnote{wəlep}
  \pos{Noun}  \sense{
    \definition{type of tree that grows in the bush with blue flowers and small blue fruit}    \note{Trees}    \example{
}    \example{
}    \example{
}}}

\entry{Wäli}{\headword{Wäli}
  \pos{Proper noun}  \sense{
    \definition{female personal name}}}

\entry{Waliyama}{\headword{Waliyama}  \pronnote{walijama}
  \pos{Proper noun}  \sense{
    \definition{Wariama (in Gogodala Rural LLG)}    \example{
}    \example{
}    \example{
}}}

\entry{wäll}{\headword{wäll}
  \variantof{Unspecified Variant of \vartext{oll}}  \note{Market}  \pronnote{wəɽ}}

\entry{walläg}{\headword{walläg}
  \pos{Verb}  \sense{
}}

\entry{wälländd}{\headword{wälländd}
  \variantof{Dialectal Variant of \vartext{ollondd}}  \pronnote{wəɽənɖ͡ʐ}}

\entry{wälläng}{\headword{wälläng}  \pronnote{wəɽəŋ}
  \pos{Noun}  \sense{
    \definition{backwoods, hinterland, (Australia) bush (any rural, undeveloped landscape)}    \example{
      \exsen{Lla da wälläng e dallän.}      \extran{The man goes to the bush.}}}
  \variants{\Dialectal Variant \vartext{wolläng, wollong}}}

\entry{wälläng gllall}{\headword{wälläng gllall}
  \pos{Noun}  \sense{
}}

\entry{wälläng ttäp}{\headword{wälläng ttäp}  \pronnote{wəɽəŋ ʈ͡ʂəp}
  \pos{Noun}  \sense{
    \definition{Papuan eagle}    \scientificname{Harpyopsis novaeguineae}    \example{
}    \example{
}    \example{
}}}

\entry{wällängakäbu}{\headword{wällängakäbu}  \pronnote{wəɽəŋakəbu}
  \pos{Noun}  \sense{
    \definition{wompoo fruit dove}    \scientificname{Ptilinopus magnificus}    \example{
}    \example{
}    \example{
}}}

\entry{wällawälla}{\headword{wällawälla}  \pronnote{wəɽawəɽa}
  \pos{Intransitive coverb}  \sense{\textbf{1.} 
    \definition{to intimidate}    \example{
      \exsen{De lla de ikop nägaeya, alla wällawälla allan?}      \extran{See that person, how he's showing off?}}}  \sense{\textbf{2.} 
    \definition{to wander, walk around aimlessly}}}

\entry{wälle}{\headword{wälle}
  \variantof{Unspecified Variant of \vartext{olle}}  \note{Verbs}  \pronnote{wəɽe}}

\entry{walle}{\headword{walle}  \note{Water}  \pronnote{waɽe}
  \pos{Noun}  \sense{
    \definition{body of water}    \note{Water}    \example{
      \exsen{Män a kädkäd dag walle we mɨnyi bogbaebän.}      \extran{The girl will jump into the cold water.}}}}

\entry{walle ddage}{\headword{walle ddage}  \lexeme{walle ddage}  \pronnote{waɽe ɖ͡ʐaɡe}
  \pos{Noun}  \sense{
    \definition{river, stream}}}

\entry{walle gutu}{\headword{walle gutu}  \pronnote{waɽe ɡutu}
  \pos{Noun}  \sense{
    \definition{dry riverbed}}}

\entry{walle mäg}{\headword{walle mäg}  \lexeme{walle mäg}  \pronnote{waɽe məɡ}
  \pos{Noun}  \sense{
    \definition{creek; river source}}}

\entry{walle menae}{\headword{walle menae}  \note{Water}  \pronnote{waɽe menae}
  \pos{Noun}  \sense{
    \definition{bank, shore, water's edge}    \note{Water}    \example{
}    \example{
}    \example{
}}}

\entry{wällegäll}{\headword{wällegäll}  \pronnote{wəɽeɡəɽ}
  \pos{Noun}  \sense{
    \definition{type of tree with fruit that are black and edible when ripe, leaves used to roll cigarettes, and roots used to treat toothache or asthma}}}

\entry{wallwall}{\headword{wallwall}
  \pos{Verb}  \sense{
}}

\entry{wällwäll}{\headword{wällwäll}  \note{Trees}  \pronnote{wəɽwəɽ}
  \pos{Noun}  \sense{
    \definition{type of tree}    \note{Trees}    \example{
}    \example{
}    \example{
}}}

\entry{wälsa}{\headword{wälsa}  \note{Trees}  \pronnote{wəlsa}
  \pos{Noun}  \sense{
    \definition{type of tree that grows in the bush}    \note{Trees}    \example{
}    \example{
}    \example{
}}}

\entry{Wama}{\headword{Wama}
  \pos{Proper noun}  \sense{
    \definition{female personal name}}}

\entry{wamän}{\headword{wamän}
  \pos{Intransitive verb}  \sense{
    \definition{to go out, dissipate, extinguish}    \example{
      \exsen{Yu da dowamänän.}      \extran{The fire went out.}}}}

\entry{Wamorong}{\headword{Wamorong}  \pronnote{wamoroŋ}
  \pos{Proper noun}  \sense{
    \definition{Wamorong (toponym)}    \example{
}    \example{
}    \example{
}}}

\entry{wän}{\headword{wän}  \pronnote{wən}
  \pos{Noun}  \sense{
    \definition{boil (on skin)}}}

\entry{wan}{\headword{wan}  \pronnote{wan}
  \pos{Numeral}  \sense{
    \definition{one (English numeral)}}
  \variants{\SpellingVariant \vartext{one}}}

\entry{Wän}{\headword{Wän}
  \pos{Proper noun}  \sense{
    \definition{female personal name}}}

\entry{=wän}{\headword{=wän}
  \variantof{Dialectal Variant of \vartext{=n}}}

\entry{wan pinga}{\headword{wan pinga}
  \pos{Noun}  \sense{
    \definition{metal fish spear}}}

\entry{wana}{\headword{wana}  \note{Names of bananas}  \pronnote{wana}
  \pos{Noun}  \sense{
    \definition{type of introduced banana}    \note{Names of bananas}    \example{
}    \example{
}    \example{
}}}

\entry{Wanadidi}{\headword{Wanadidi}
  \pos{Proper noun}  \sense{
    \definition{female personal name}}}

\entry{wänänang}{\headword{wänänang}  \pronnote{wənənaŋ}
  \pos{Modifier}  \sense{
    \definition{provider, bringing back food for one's family}    \example{
      \exsen{Bogo ddobae wänänang dan.}      \extran{She is a provider.}}}}

\entry{wandae}{\headword{wandae}
  \pos{Transitive verb}  \sense{
    \definition{to burn}}}

\entry{wändäg}{\headword{wändäg}  \pronnote{wəndəɡ}
  \pos{Transitive verb}  \sense{
    \definition{to crowd}    \example{
      \exsen{Yesu adame obo kollmällang de umllang dägnegän oblle gall kälsäre särämbae e, oba lla da ddone obom beyawändägallo.}      \extran{This is why Jesus told his disciples to prepare a small boat for him, so that people wouldn't come crowd him.}}}
  \variants{\SpellingVariant \vartext{wändɨg}}}

\entry{wandana}{\headword{wandana}  \note{Gardens}  \pronnote{wandana}
  \pos{Noun}  \sense{
    \definition{grass in the garden}    \note{Gardens}    \example{
}    \example{
}    \example{
}}}

\entry{wandawandae}{\headword{wandawandae}  \pronnote{wandawandae}
  \pos{Adverb}  \sense{
    \definition{rotating, spinning, rolling}}}

\entry{wändɨg}{\headword{wändɨg}
  \variantof{SpellingVariant of \vartext{wändäg}}}

\entry{Wane}{\headword{Wane}  \pronnote{wane}
  \pos{Proper noun}  \sense{
    \definition{male personal name}}}

\entry{=wane}{\headword{=wane}
  \variantof{Dialectal Variant of \vartext{=alle}}}

\entry{wangam}{\headword{wangam}  \pronnote{waŋam}
  \pos{Transitive verb}  \sense{
    \definition{to forget}    \example{
      \exsen{Ddone mänyi bam bawengameya.}      \extran{We will not forget you.}}}}

\entry{wängän}{\headword{wängän}  \note{Trees}  \pronnote{wəŋən}
  \pos{Noun}  \sense{
    \definition{type of tree}    \note{Trees}    \example{
}    \example{
}    \example{
}}}

\entry{wänkäm}{\headword{wänkäm}  \note{Parts of the Body}  \pronnote{wənkəm}
  \pos{Noun}  \sense{\textbf{1.} 
    \definition{belly button, navel}    \note{Parts of the Body}    \example{
}    \example{
}    \example{
}}  \sense{\textbf{2.} 
    \definition{anus}    \note{Parts of the Body}}
  \variants{\Unspecified Variant \vartext{wunkäm, unkäm}}}

\entry{wänkäm molle}{\headword{wänkäm molle}  \pronnote{wənkəm moɽe}
  \pos{Noun}  \sense{
    \definition{soft part of a shoot or sucker (e.g. of taro, banana, or sago)}}}

\entry{wänkäm tärpan}{\headword{wänkäm tärpan}  \note{Birth}  \pronnote{wənkəm tərpan}
  \pos{Noun}  \sense{
    \definition{umbilical cord}    \note{Birth}    \example{
}    \example{
}    \example{
}}}

\entry{wäno}{\headword{wäno}  \note{Trees}  \pronnote{wəno}
  \pos{Noun}  \sense{
    \definition{type of tree that grows in the grassland and along creeks with white flowers, small brown fruit, and bark used on rooves}    \note{Trees}    \example{
}    \example{
}    \example{
}}}

\entry{wanpadam}{\headword{wanpadam}  \note{clothing}  \pronnote{wanpadam}
  \pos{Noun}  \sense{
    \definition{lap-lap}    \note{clothing}    \example{
}    \example{
}    \example{
}}}

\entry{wanseg}{\headword{wanseg}  \pronnote{wanseɡ}
  \pos{Transitive verb}  \sense{\textbf{1.} 
    \definition{to put, place, set aside, leave}    \example{
      \exsen{Abo ddäddäg de nokomeyo a katrekatre toko me nowattälleyo.}      \extran{You all should take the pieces of meat and put them on the table.}}    \example{
      \exsen{Pa obom dowansegän a dallän nagda bom dingällbänän a kottllam pate gongttägeyo.}      \extran{The bird left him and went to get his friend and returned to the turtle.}}    \example{
      \exsen{Oba llɨg kälekäle tämamae ma me dowattälleyo.}      \extran{They left all their small children at home.}}}  \sense{\textbf{2.} 
    \definition{to be left (in a state)}    \example{
      \exsen{Lama wa Mana ubi obaoba ulle abal kandärmang me gowensegäneyo.}      \extran{Lama and Mana were left feeling very sorry for themselves.}}    \example{
      \exsen{Däbe källäm a yäbäd ulle atta dallän do kälae abal gowensegän ine da.}      \extran{Because of the strong sun, that pond kept getting smaller until there was very little water left.}}}
  \variants{\SpellingVariant \vartext{wensäg}}}

\entry{wanttawantta}{\headword{wanttawantta}  \note{Games}  \pronnote{wanʈ͡ʂawanʈ͡ʂa}
  \pos{Noun}  \sense{
    \definition{type of game like capture the flag but with a stick planted in the middle of a ring instead of a flag}    \note{Games}    \example{
}    \example{
}    \example{
}}}

\entry{wanwen}{\headword{wanwen}  \pronnote{wanwen}
  \pos{Transitive verb}  \sense{\textbf{1.} 
    \definition{to shake, swing}    \example{
      \exsen{Bogo bun gowenän ada, "Ddone."}      \extran{His head shook like this, "No."}}}  \sense{\textbf{2.} 
    \definition{to shake, swing}    \example{
      \exsen{Zämllallang a oba bun di dowanaemneyo mikutt me.}      \extran{Passers-by shook their heads in anger.}}}}

\entry{wäny}{\headword{wäny}
  \variantof{SpellingVariant of \vartext{ony}}}

\entry{wanyweny}{\headword{wanyweny}
  \pos{Transitive verb}  \sense{
    \definition{to burn}    \example{
      \exsen{Kullkull wanyweny e gotäbanegnän.}      \extran{They were planning to burn a grass fire.}}}}

\entry{waowaem}{\headword{waowaem}
  \variantof{SpellingVariant of \vartext{wawaem}}  \note{Swadesh}  \pronnote{waowaem}}

\entry{wap}{\headword{wap}  \pronnote{wap}
  \pos{Noun}  \sense{
    \definition{stick}    \example{
      \exsen{Ngäna wap alle pa de gäddnan dängkaemneg.}      \extran{I started killing birds with sticks.}}}}

\entry{wäp}{\headword{wäp}
  \variantof{freevarlab of \vartext{up}}}

\entry{wapo}{\headword{wapo}  \pronnote{wapo}
  \pos{Interjection}  \sense{
    \definition{oh no}}}

\entry{Wapotea}{\headword{Wapotea}  \note{Place names around Limol}  \pronnote{wapotea}
  \pos{Proper noun}  \sense{
    \definition{Wapotea (toponym)}    \note{Place names around Limol}    \example{
}    \example{
}    \example{
}}}

\entry{Wara}{\headword{Wara}  \note{Places around Limol}  \pronnote{wara}
  \pos{Proper noun}  \sense{
    \definition{Wara (toponym)}    \note{Places around Limol}    \example{
}    \example{
}    \example{
}}}

\entry{Warama}{\headword{Warama}  \pronnote{warama}
  \pos{Proper noun}  \sense{
    \definition{male personal name}}}

\entry{waramakae}{\headword{waramakae}  \note{Numbers}  \pronnote{waramakae}
  \pos{Numeral}  \sense{
    \definition{7776 (yam counting numeral; 6^5)}    \note{Numbers}    \example{
}    \example{
}    \example{
}}}

\entry{waramawarama}{\headword{waramawarama}  \note{Trees}  \pronnote{waramawarama}
  \pos{Noun}  \sense{
    \definition{type of tree that grows in the bush with white flowers and edible yellow fruit}    \note{Trees}    \example{
}    \example{
}    \example{
}}}

\entry{Warani}{\headword{Warani}  \pronnote{warani}
  \pos{Proper noun}  \sense{
    \definition{male personal name}}}

\entry{Wareka}{\headword{Wareka}  \pronnote{wareka}
  \pos{Proper noun}  \sense{
    \definition{female personal name}}}

\entry{wärenzbag}{\headword{wärenzbag}  \note{Types of Taro}  \pronnote{wərenzbaɡ}
  \pos{Noun}  \sense{
    \definition{type of taro}    \note{Types of Taro}    \example{
}    \example{
}    \example{
}}}

\entry{Wareya}{\headword{Wareya}
  \pos{Proper noun}  \sense{
    \definition{male personal name}}}

\entry{Wariabodolo}{\headword{Wariabodolo}  \note{Place names}  \pronnote{wariabodolo}
  \pos{Proper noun}  \sense{
    \definition{Variobadoro (in Kiwai Rural LLG)}    \note{Place names}    \example{
}    \example{
}    \example{
}}}

\entry{Warik}{\headword{Warik}
  \pos{Proper noun}  \sense{
    \definition{male personal name}}}

\entry{wariwari}{\headword{wariwari}
  \pos{Noun}  \sense{
    \definition{sago shoot}}}

\entry{waro}{\headword{waro}  \note{(picture in k-iphone)}  \pronnote{waro}
  \pos{Noun}  \sense{
    \definition{type of turtle}    \note{(picture in k-iphone)}}}

\entry{Warola}{\headword{Warola}  \note{Places around Limol}  \pronnote{warola}
  \pos{Proper noun}  \sense{
    \definition{Warola (camping place in Limol)}    \note{Places around Limol}    \example{
}    \example{
}    \example{
}}
  \variants{\Unspecified Variant \vartext{Warolla}}}

\entry{Warolla}{\headword{Warolla}
  \variantof{Unspecified Variant of \vartext{Warola}}}

\entry{wärpir}{\headword{wärpir}
  \pos{Noun}  \sense{
    \definition{slippery mud found by the river}}}

\entry{Wasang}{\headword{Wasang}
  \pos{Proper noun}  \sense{
    \definition{male personal name}}}

\entry{wasangwasang}{\headword{wasangwasang}
  \pos{Adverb}  \sense{
    \definition{always}    \example{
      \exsen{Angde Matthew llɨg kälsre me, dibe llo de wasangwasang dängkälne.}      \extran{When Matthiew was a little boy, he was always climbing that tree.}}}}

\entry{wasar}{\headword{wasar}  \note{Types of palms}  \pronnote{wasar}
  \pos{Noun}  \sense{
    \definition{type of edible palm}    \note{Types of palms}    \example{
}    \example{
}    \example{
}}}

\entry{waso}{\headword{waso}  \pronnote{waso}
  \pos{Noun}  \sense{
    \definition{eastern cattle egret}    \scientificname{Bubulcus coromandus}    \example{
}    \example{
}    \example{
}}}

\entry{Wasua}{\headword{Wasua}
  \pos{Proper noun}  \sense{
    \definition{Wasua (in Gogodala Rural LLG)}}
  \variants{\SpellingVariant \vartext{Wasuwa}}}

\entry{Wasuwa}{\headword{Wasuwa}
  \variantof{SpellingVariant of \vartext{Wasua}}  \pronnote{wasuwa}}

\entry{waswes}{\headword{waswes}  \note{Government}  \pronnote{waswes}
  \pos{Transitive verb}  \sense{\textbf{1.} 
    \definition{to ask, beg}    \note{Government}    \example{
      \exsen{Da ngämi e we bam baweseya, ngäma moko da bongo nangesneg ngämira.}      \extran{Whatever we (excl.) ask you for, we want you to do it for us.}}    \example{
      \exsen{Ge ngäna ewe de dawes ikopse alle.}      \extran{This is what I had asked for in prayer.}}}  \sense{\textbf{2.} 
    \definition{political group}    \note{Government}    \example{
}    \example{
}    \example{
}}}

\entry{wätaote}{\headword{wätaote}  \pronnote{wətaote}
  \pos{Noun}  \sense{\textbf{1.} 
    \definition{type of large vine that grows in bush; used to tie fence posts together}    \example{
}    \example{
}    \example{
}}  \sense{\textbf{2.} 
    \definition{type of taro}}}

\entry{wätät}{\headword{wätät}
  \variantof{Unspecified Variant of \vartext{otät}}}

\entry{wäte}{\headword{wäte}
  \variantof{freevarlab of \vartext{ute}}}

\entry{Wätt bo ma}{\headword{Wätt bo ma}  \note{Places around Limol}  \pronnote{wəʈ͡ʂ bo ma}
  \pos{Proper noun}  \sense{
    \definition{Wätt bo ma (toponym)}    \note{Places around Limol}    \example{
}    \example{
}    \example{
}}}

\entry{wawa}{\headword{wawa}  \note{Trees}  \pronnote{wawa}
  \pos{Noun}  \sense{
    \definition{type of tree that grows in the bush and along creeks white flowers and blue fruit}    \note{Trees}    \example{
}    \example{
}    \example{
}}}

\entry{wawaem}{\headword{wawaem}  \note{Animal verbs}  \pronnote{wawaem}
  \pos{Noun}  \sense{\textbf{1.} 
    \definition{hiss}    \note{Animal verbs}    \example{
      \exsen{Ngäna wawaem de dandär.}      \extran{I heard a hiss.}}}  \sense{\textbf{2.} 
    \definition{current}    \note{Animal verbs}    \example{
      \exsen{Bituri mängallang abal ine wawaem de dangalltallo mutt i.}      \extran{They paddled upstream against the strong current of the Bituri river.}}}  \sense{\textbf{3.} 
    \definition{to flow}    \note{Animal verbs}    \example{
      \exsen{Yogoll mamott me ine da bingkälän a wawaem bognän.}      \extran{During rainy times, the water rises and flows.}}}
  \variants{\SpellingVariant \vartext{waowaem}}}

\entry{Wawase}{\headword{Wawase}  \pronnote{wawase}
  \pos{Proper noun}  \sense{
    \definition{male personal name}}}

\entry{Waweba}{\headword{Waweba}  \pronnote{waweba}
  \pos{Proper noun}  \sense{
    \definition{male personal name}}}

\entry{Wawera}{\headword{Wawera}
  \pos{Proper noun}  \sense{
    \definition{male personal name}}}

\entry{wawonai}{\headword{wawonai}  \lexeme{wawonai}  \pronnote{wawonai}
  \pos{Noun}  \sense{
    \definition{type of long yam with a hooked end, white interior, and hairs}}}

\entry{waya}{\headword{waya}  \pronnote{waja}
  \pos{Noun}  \sense{
    \definition{type of pronged metal spear}}}

\entry{waya gullem}{\headword{waya gullem}  \note{Snakes}  \pronnote{waja ɡuɽem}
  \pos{Noun}  \sense{
    \definition{type of venomous snake}    \note{Snakes}    \example{
}    \example{
}    \example{
}}}

\entry{Wayapampe}{\headword{Wayapampe}
  \pos{Proper noun}  \sense{
    \definition{Wayapampe (toponym)}}}

\entry{wayati}{\headword{wayati}  \lexeme{wayati}  \pronnote{wajati}
  \pos{Noun}  \sense{
    \definition{watermelon}    \example{
      \exsen{Ngämo wayati kongkom e nalle.}      \extran{[You] go carry my watermelons.}}}}

\entry{Wäziag}{\headword{Wäziag}
  \pos{Proper noun}  \sense{
    \definition{personal name}}}

\entry{wẽ}{\headword{wẽ}  \pronnote{wẽ}
  \pos{Ideophone}  \sense{
    \definition{sound of crying}}}

\entry{Wed}{\headword{Wed}  \pronnote{wed}
  \pos{Proper noun}  \sense{\textbf{1.} 
    \definition{Wed (toponym)}}  \sense{\textbf{2.} 
    \definition{female personal name}}}

\entry{Wedereamo}{\headword{Wedereamo}  \note{Place names}  \pronnote{wedereamo}
  \pos{Proper noun}  \sense{
    \definition{Wederehiamo (on the south side of the mouth of the Fly River)}    \note{Place names}    \example{
}    \example{
}    \example{
}}}

\entry{wel}{\headword{wel}  \note{Weather}  \pronnote{wel}
  \pos{Noun}  \sense{
    \definition{wind}    \note{Weather}    \example{
      \exsen{Känazbag mänyi wel peyang bogon.}      \extran{Tomorrow, it will be windy.}}    \example{
}    \example{
}}}

\entry{wel ud}{\headword{wel ud}  \note{Everything in the house}  \pronnote{wel ud}
  \pos{Noun}  \sense{
    \definition{window}    \note{Everything in the house}    \example{
}    \example{
}    \example{
}}}

\entry{welwel}{\headword{welwel}  \note{Birds}  \pronnote{welwel}
  \pos{Noun}  \sense{
    \definition{type of bird}    \note{Birds}    \example{
}    \example{
}    \example{
}}}

\entry{welwele}{\headword{welwele}
  \pos{Noun}  \sense{
}}

\entry{Wendi}{\headword{Wendi}
  \pos{Noun}  \sense{
    \definition{female personal name}}
  \variants{\SpellingVariant \vartext{Wendy}}}

\entry{Wendy}{\headword{Wendy}
  \variantof{SpellingVariant of \vartext{Wendi}}  \pronnote{wɛndi}}

\entry{wensäg}{\headword{wensäg}
  \variantof{SpellingVariant of \vartext{wanseg}}  \pronnote{wensəɡ}}

\entry{wer}{\headword{wer}  \pronnote{wer}
  \pos{Noun}  \sense{
    \definition{type of tree with edible black fruit with one seed}}}

\entry{Wesley}{\headword{Wesley}
  \variantof{SpellingVariant of \vartext{Wesli}}}

\entry{Wesli}{\headword{Wesli}  \pronnote{wesli}
  \pos{Proper noun}  \sense{
    \definition{male personal name}}
  \variants{\SpellingVariant \vartext{Wesley}}}

\entry{=weya}{\headword{=weya}
  \variantof{freevarlab of \vartext{=aeya}}}

\entry{when}{\headword{when}
  \variantof{freevarlab of \vartext{English loanword}}}

\entry{wi}{\headword{wi}  \pronnote{wi}
  \pos{Intransitive verb}  \sense{
    \definition{to settle}    \example{
      \exsen{Särem a deyawinän.}      \extran{Darkness settled.}}    \example{
      \exsen{Ttängäm a zäme deyawinän toto me.}      \extran{The village had already settled down in the early evening.}}}}

\entry{wi}{\headword{wi}
  \pos{Interjection}  \sense{
    \definition{oh}}}

\entry{wia}{\headword{wia}
  \variantof{SpellingVariant of \vartext{wiya}}}

\entry{wiasara}{\headword{wiasara}
  \variantof{SpellingVariant of \vartext{wiyasara}}}

\entry{wib}{\headword{wib}  \note{Trees}  \pronnote{wib}
  \pos{Noun}  \sense{
    \definition{type of tree}    \note{Trees}    \example{
}    \example{
}    \example{
}}}

\entry{wibell}{\headword{wibell}  \note{Trees}  \pronnote{wibeɽ}
  \pos{Noun}  \sense{
    \definition{type of tree}    \note{Trees}    \example{
}    \example{
}    \example{
}}}

\entry{Wiben}{\headword{Wiben}  \pronnote{wiben}
  \pos{Proper noun}  \sense{
    \definition{male personal name}}}

\entry{Widama}{\headword{Widama}  \pronnote{widama}
  \pos{Proper noun}  \sense{
    \definition{Widama (toponym)}    \example{
}    \example{
}    \example{
}}}

\entry{widere}{\headword{widere}  \note{Tools}  \pronnote{widere}
  \pos{Noun}  \sense{
    \definition{paddle, oar}    \note{Tools}    \example{
      \exsen{Ngäna ngämlle ttongo widere sisor de popo e dan.}      \extran{I want to make a new paddle.}}    \example{
}    \example{
}}}

\entry{widwid}{\headword{widwid}  \pronnote{widwid}
  \pos{Noun}  \sense{
    \definition{type of plant with big leaves}    \example{
}    \example{
}    \example{
}}}

\entry{Wik}{\headword{Wik}
  \pos{Proper noun}  \sense{
    \definition{male personal name}}}

\entry{wik}{\headword{wik}  \pronnote{wik}
  \pos{Noun}  \sense{
    \definition{week}}}

\entry{Willie}{\headword{Willie}  \pronnote{wiɽie}
  \pos{Proper noun}  \sense{
    \definition{male personal name}}}

\entry{Wilma}{\headword{Wilma}  \pronnote{wilma}
  \pos{Proper noun}  \sense{
    \definition{female personal name}}}

\entry{wilwil}{\headword{wilwil}  \note{Trees}  \pronnote{wilwil}
  \pos{Noun}  \sense{
    \definition{type of tree that grows in the bush}    \note{Trees}    \example{
}    \example{
}    \example{
}}}

\entry{Wim}{\headword{Wim}  \pronnote{wim}
  \pos{Proper noun}  \sense{
    \definition{Wim (Kawam-speaking village in Oriomo-Bituri Rural LLG; GPS: 8.762144, 142.770164)}    \example{
}    \example{
}    \example{
}}}

\entry{win}{\headword{win}  \pronnote{win}
  \pos{Noun}  \sense{
    \definition{win}}}

\entry{Wingäm}{\headword{Wingäm}
  \pos{Proper noun}  \sense{
    \definition{female personal name}}}

\entry{Wini}{\headword{Wini}
  \variantof{SpellingVariant of \vartext{Winny}}}

\entry{winisde}{\headword{winisde}  \pronnote{winisde}
  \pos{Noun}  \sense{
    \definition{Wednesday}    \example{
}    \example{
}    \example{
}}}

\entry{Winny}{\headword{Winny}  \pronnote{wini}
  \pos{Proper noun}  \sense{
    \definition{female personal name}}
  \variants{\SpellingVariant \vartext{Wini}}}

\entry{Winsen}{\headword{Winsen}
  \variantof{SpellingVariant of \vartext{Winson}}}

\entry{Winson}{\headword{Winson}  \pronnote{winson}
  \pos{Proper noun}  \sense{
    \definition{male personal name}}
  \variants{\SpellingVariant \vartext{Winsen}}}

\entry{Winy}{\headword{Winy}  \note{Places around Limol}  \pronnote{wiɲ}
  \pos{Proper noun}  \sense{
    \definition{Winy (toponym)}    \note{Places around Limol}    \example{
}    \example{
}    \example{
}}}

\entry{winy}{\headword{winy}  \lexeme{winy}  \pronnote{wiɲ}
  \pos{Noun}  \sense{
    \definition{honey}    \anthronote{Used as medicine to treat diarhoea}    \example{
      \exsen{Däbe tätämatt winy a ka tomowangtomowang moko allan.}      \extran{That honey from yesterday tastes a little bitter.}}}}

\entry{winyteya}{\headword{winyteya}  \note{Bees}  \pronnote{wiɲteja}
  \pos{Noun}  \sense{
    \definition{honeycomb}    \note{Bees}    \example{
}    \example{
}    \example{
}}}

\entry{wiowa}{\headword{wiowa}
  \variantof{Dialectal Variant of \vartext{wiyo}}  \pronnote{wiowa}}

\entry{wipell}{\headword{wipell}  \note{Types of palms}  \pronnote{wipeɽ}
  \pos{Noun}  \sense{
    \definition{type of tall palm that grows along the riverside}    \note{Types of palms}    \example{
}    \example{
}    \example{
}}}

\entry{wipellgallagallab}{\headword{wipellgallagallab}  \pronnote{wipeɽɡaɽaɡaɽab}
  \pos{Noun}  \sense{
    \definition{chevron weaving pattern}    \example{
}    \example{
}    \example{
}}}

\entry{Wipi}{\headword{Wipi}  \note{Market}  \pronnote{wipi}
  \pos{Proper noun}  \sense{
    \definition{Wipi language}    \note{Market}    \example{
}    \example{
}    \example{
}}}

\entry{Wipim}{\headword{Wipim}  \pronnote{wipim}
  \pos{Proper noun}  \sense{
    \definition{Wipim (Wipi- and Kawam-speaking village in Oriomo-Bituri Rural LLG; GPS: -8.786616, 142.872201)}    \example{
}    \example{
}    \example{
}}}

\entry{wirog}{\headword{wirog}  \note{Names of bananas}  \pronnote{wiroɡ}
  \pos{Noun}  \sense{
    \definition{type of native banana}    \note{Names of bananas}    \example{
}    \example{
}    \example{
}}}

\entry{wiswis}{\headword{wiswis}  \note{Trees}  \pronnote{wiswis}
  \pos{Noun}  \sense{
    \definition{type of tree that grows in the bush with white flowers and edible orange fruit}    \note{Trees}    \example{
}    \example{
}    \example{
}}}

\entry{wit}{\headword{wit}  \pronnote{wit}
  \pos{Noun}  \sense{
    \definition{wheat}}}

\entry{witara}{\headword{witara}  \pronnote{witara}
  \pos{Noun}  \sense{
    \definition{type of garden}    \example{
      \exsen{Witara da karama me däm ibeny ma ttängäm dan.}}    \example{
      \exsen{Witara me wayati da, pamker a, biye da wa mompel a däbem de ibeb erallo.}      \extran{They are planting watermelon, pumpkin, taro, and aibika in the river garden.}}    \example{
      \exsen{Lla da witara me mondre anggan.}}}}

\entry{wiya}{\headword{wiya}  \pronnote{wija}
  \pos{Intransitive verb}  \sense{
    \definition{come (imperative, singular form)}    \example{
      \exsen{Wiya gänyerimae!}      \extran{[You] come over here!}}    \example{
}    \example{
}}
  \variants{\SpellingVariant \vartext{wia}}}

\entry{wiyamom}{\headword{wiyamom}  \pronnote{wijamom}
  \pos{Intransitive verb}  \sense{
    \definition{plural form of wiya}    \example{
      \exsen{Tämamae wiyamom!}      \extran{Come, everyone!}}}}

\entry{wiyasara}{\headword{wiyasara}  \pronnote{wijasara}
  \pos{Noun}  \sense{
    \definition{silver gull}    \scientificname{Chroicocephalus novaehollandiae}}
  \variants{\SpellingVariant \vartext{wiasara}}}

\entry{wiyebubug}{\headword{wiyebubug}  \note{Parts of the body}  \pronnote{wijebubuɡ}
  \pos{Noun}  \sense{
    \definition{body part}    \note{Parts of the body}    \example{
}    \example{
}    \example{
}}}

\entry{wiyo}{\headword{wiyo}  \pronnote{wijo}
  \pos{Interjection}  \sense{
    \definition{wow}}
  \variants{\Dialectal Variant \vartext{wiowa}}}

\entry{wiyowae}{\headword{wiyowae}  \pronnote{wijowae}
  \pos{Interjection}  \sense{
    \definition{wow; phew}}}

\entry{wiyowe}{\headword{wiyowe}  \note{Types of palms}  \pronnote{wijowe}
  \pos{Noun}  \sense{
    \definition{type of large palm that grows in the bush or along creeks with flowers that start from the top and spread downwards and white fruit that hangs like coconut}    \note{Types of palms}    \example{
}    \example{
}    \example{
}}}

\entry{wizarab}{\headword{wizarab}  \note{Types of palms}  \pronnote{wizarab}
  \pos{Noun}  \sense{
    \definition{type of pandanus with red fruit}    \note{Types of palms}    \example{
}    \example{
}    \example{
}}}

\entry{Wizing}{\headword{Wizing}  \pronnote{wiziŋ}
  \pos{Proper noun}  \sense{
    \definition{male personal name}}}

\entry{wɨndwɨnd}{\headword{wɨndwɨnd}  \pronnote{wɪndwɪnd}
  \pos{Transitive verb}  \sense{
    \definition{to cover, bury}    \example{
      \exsen{Pätt de ekaklle alle dawɨndeya.}      \extran{We covered the body with earth.}}}}

\entry{Wɨr}{\headword{Wɨr}
  \pos{Proper noun}  \sense{
    \definition{Wɨr (toponym)}}}

\entry{Wɨrbun}{\headword{Wɨrbun}
  \pos{Proper noun}  \sense{
    \definition{Wɨrbun (toponym)}}}

\entry{wɨtwɨt}{\headword{wɨtwɨt}
  \pos{Noun}  \sense{
    \definition{type of sago}}}

\entry{wo}{\headword{wo}
  \variantof{SpellingVariant of \vartext{o}}}

\entry{woboll}{\headword{woboll}  \pronnote{woboɽ}
  \pos{Noun}  \sense{
    \definition{type of plant}    \example{
}}}

\entry{wod}{\headword{wod}  \lexeme{wod}  \pronnote{wod}
  \pos{Noun}  \sense{\textbf{1.} 
    \definition{type of fatty fish}}  \sense{\textbf{2.} 
    \definition{type of long yam with a white interior and few hairs}}}

\entry{wodd memba}{\headword{wodd memba}  \pronnote{woɖ͡ʐ memba}
  \pos{Noun}  \sense{
    \definition{ward member}    \example{
}    \example{
}    \example{
}}}

\entry{woddowoddo}{\headword{woddowoddo}  \pronnote{woɖ͡ʐowoɖ͡ʐo}
  \pos{Noun}  \sense{\textbf{1.} 
    \definition{rusty pitohui}    \scientificname{Pseudorectes ferrugineus clarus}    \example{
}    \example{
}    \example{
}}  \sense{\textbf{2.} 
    \definition{Amboyna cuckoo-dove}    \scientificname{Macropygia amboinensis}}}

\entry{wolläb}{\headword{wolläb}  \pronnote{woɽəb}
  \pos{Noun}  \sense{
    \definition{plant term}    \example{
      \exsen{Bie wolläb de näddäg.}}}}

\entry{wolläng}{\headword{wolläng}
  \variantof{Dialectal Variant of \vartext{wälläng}}}

\entry{wolle}{\headword{wolle}
  \variantof{Unspecified Variant of \vartext{olle}}}

\entry{wollong}{\headword{wollong}
  \variantof{Dialectal Variant of \vartext{wälläng}}}

\entry{Wonie}{\headword{Wonie}  \pronnote{wonie}
  \pos{Proper noun}  \sense{
    \definition{Wonie (Wipi-speaking village in Oriomo-Bituri Rural LLG; GPS: -8.838939, 142.975007)}    \example{
}    \example{
}    \example{
}}}

\entry{worbam}{\headword{worbam}
  \variantof{SpellingVariant of \vartext{orbam}}  \pronnote{worbam}}

\entry{Wot}{\headword{Wot}  \note{Places around Limol}  \pronnote{wot}
  \pos{Proper noun}  \sense{
    \definition{Wot (toponym)}    \note{Places around Limol}    \example{
}    \example{
}    \example{
}}}

\entry{wowo}{\headword{wowo}  \pronnote{wowo}
  \pos{Transitive verb}  \sense{
    \definition{to clear floating grass by pushing through with a canoe}    \example{
      \exsen{Tätäm ibi tawa de duwonalla.}      \extran{Yesterday we cleared the floating grass.}}}}

\entry{wud}{\headword{wud}
  \variantof{SpellingVariant of \vartext{ud}}}

\entry{wudu}{\headword{wudu}
  \variantof{Dialectal Variant of \vartext{udu}}}

\entry{wuk}{\headword{wuk}
  \variantof{SpellingVariant of \vartext{uk}}}

\entry{Wun}{\headword{Wun}  \pronnote{wun}
  \pos{Proper noun}  \sense{
    \definition{male personal name}}}

\entry{wunkäm}{\headword{wunkäm}
  \variantof{Unspecified Variant of \vartext{wänkäm}}  \pronnote{wunkəm}}

\entry{wup}{\headword{wup}
  \variantof{freevarlab of \vartext{up}}}

\entry{wup ttämbällag}{\headword{wup ttämbällag}  \pronnote{wup ʈ͡ʂəmbəɽaɡ}
  \pos{Noun}  \sense{
    \definition{type of spear}    \example{
}    \example{
}    \example{
}}}

\entry{Wur}{\headword{Wur}
  \variantof{SpellingVariant of \vartext{Ur}}  \note{Place names around Limol}  \pronnote{wur}}

\entry{Wurlaimäll}{\headword{Wurlaimäll}  \note{Places around Limol}  \pronnote{wurlaiməɽ}
  \pos{Proper noun}  \sense{
    \definition{Wurlaimäll (toponym)}    \note{Places around Limol}    \example{
}    \example{
}    \example{
}}}

\entry{Wuroi}{\headword{Wuroi}
  \variantof{SpellingVariant of \vartext{Uroe}}}

\entry{wute}{\headword{wute}
  \variantof{SpellingVariant of \vartext{ute}}}

\entry{wutt}{\headword{wutt}
  \variantof{SpellingVariant of \vartext{utt}}}

\chapter{Y}


\entry{y-}{\headword{y-}
  \pos{Verb}  \sense{\textbf{1.} 
    \definition{du|1pl|2plP}}  \sense{\textbf{2.} 
    \definition{duS}}  \sense{\textbf{3.} 
}  \sense{\textbf{4.} 
}}

\entry{ya}{\headword{ya}  \pronnote{ja}
  \pos{Interjection}  \sense{
    \definition{scram, go away}}}

\entry{yäbäd}{\headword{yäbäd}  \lexeme{yäbäd}  \pronnote{jəbəd}
  \pos{Noun}  \sense{\textbf{1.} 
    \definition{sun}    \example{
      \exsen{Yäbäd de ddapall käkan a dakonewän.}      \extran{The clouds covered the sun.}}}  \sense{\textbf{2.} 
    \definition{season when the dry season starts and people go camping in the bush (eleventh season; corresponds to August)}}  \sense{\textbf{3.} 
    \definition{weather}    \example{
      \exsen{Mer abal yäbäd daeya.}      \extran{It was very good weather.}}}}

\entry{yäbäd bäng}{\headword{yäbäd bäng}  \lexeme{yäbäd bäng}  \pronnote{jəbəd bəŋ}
  \pos{Noun}  \sense{\textbf{1.} 
    \definition{dry, hot season characterized by burning grass (twelfth season; corresponds to September)}}  \sense{\textbf{2.} 
    \definition{drought}}}

\entry{yäbäd gazen}{\headword{yäbäd gazen}  \note{Parts of the day}  \pronnote{jəbəd ɡazen}
  \pos{Noun}  \sense{
    \definition{dawn}    \note{Parts of the day}    \example{
}    \example{
}    \example{
}}}

\entry{yäbäd källa}{\headword{yäbäd källa}  \pronnote{jəbəd kəɽa}
  \pos{Noun}  \sense{
    \definition{type of big taro}    \example{
}    \example{
}    \example{
}}}

\entry{yäbäd kuttakuttang}{\headword{yäbäd kuttakuttang}  \note{Weather}  \pronnote{jəbəd kuʈ͡ʂakuʈ͡ʂaŋ}
  \pos{Modifier}  \sense{
    \definition{really hot}    \note{Weather}    \example{
}    \example{
}    \example{
}}}

\entry{yäbäd ttänttämang}{\headword{yäbäd ttänttämang}  \pronnote{jəbəd ʈ͡ʂənʈ͡ʂəmaŋ}
  \pos{Noun}  \sense{
    \definition{hot season when new gardens are burnt (thirteenth month; corresponds to October)}    \example{
}    \example{
}    \example{
}}}

\entry{yäbäd tuktuk}{\headword{yäbäd tuktuk}  \note{Parts of the day}  \pronnote{jəbəd tuktuk}
  \pos{Noun}  \sense{
    \definition{solar noon (when the sun reaches its zenith)}    \note{Parts of the day}    \example{
}    \example{
}    \example{
}}}

\entry{yäbäd tuktuk duwem}{\headword{yäbäd tuktuk duwem}  \note{Parts of the day}  \pronnote{jəbəd tuktuk duwem}
  \pos{Noun}  \sense{
    \definition{lunch}    \note{Parts of the day}    \example{
}    \example{
}    \example{
}}}

\entry{yäbäyäbäd}{\headword{yäbäyäbäd}  \note{Trees}  \pronnote{jəbəjəbəd}
  \pos{Noun}  \sense{\textbf{1.} 
    \definition{type of tree that grows in the bush with white flowers and red fruit}    \note{Trees}    \example{
}    \example{
}    \example{
}}  \sense{\textbf{2.} 
    \definition{type of grass}    \note{Trees}}}

\entry{yäbäyäbäl}{\headword{yäbäyäbäl}  \note{Trees}  \pronnote{jəbəjəbəl}
  \pos{Noun}  \sense{
    \definition{type of tree}    \note{Trees}    \example{
}    \example{
}    \example{
}}}

\entry{yäbdo}{\headword{yäbdo}
  \variantof{freevarlab of \vartext{ebdo}}}

\entry{yaber}{\headword{yaber}  \note{Trees}  \pronnote{jaber}
  \pos{Noun}  \sense{
    \definition{type of tree that grows in the bush with white flowers and poisonous bark used to catch fish}    \note{Trees}    \example{
}    \example{
}    \example{
}}}

\entry{yäbik}{\headword{yäbik}  \lexeme{yäbik}  \pronnote{jəbik}
  \pos{Noun}  \sense{
    \definition{sharp gardening stick}}}

\entry{yad}{\headword{yad}  \pronnote{jad}
  \pos{Noun}  \sense{
    \definition{yard}}}

\entry{yae}{\headword{yae}  \lexeme{yae}  \pronnote{jae}
  \pos{Kinship noun}  \sense{
    \definition{mother}}}

\entry{yaedidib}{\headword{yaedidib}  \lexeme{yaedidib}  \pronnote{jaedidib}
  \pos{Noun}  \sense{
    \definition{type of long, narrow grass (~0.3 m)}}}

\entry{yag~}{\headword{yag~}
  \variantof{Infinitival reduplicant of \vartext{RED~}}}

\entry{yäg}{\headword{yäg}  \lexeme{yäg}  \pronnote{jəɡ}
  \pos{Noun}  \sense{
    \definition{bush rope}}}

\entry{yagäl}{\headword{yagäl}  \note{Trees}  \pronnote{jaɡəl}
  \pos{Noun}  \sense{\textbf{1.} 
    \definition{type of tree that grows in the grassland with leaves used to sand bows and spears}    \note{Trees}    \example{
}    \example{
}    \example{
}}  \sense{\textbf{2.} 
    \definition{to smooth with the yagäl leaf}    \note{Trees}}}

\entry{yagyag}{\headword{yagyag}
  \variantof{freevarlab of \vartext{yagyeg}}  \pronnote{jaɡjaɡ}}

\entry{yagyeg}{\headword{yagyeg}  \note{Verbs}  \pronnote{jaɡjeɡ}
  \pos{Transitive verb}  \sense{
    \definition{to search, look for}    \note{Verbs}    \example{
      \exsen{Män nagda bom yagnen dängkamän.}      \extran{‎‎The girl started to look for her friend.}}    \example{
      \exsen{Erodias nyongo de deyagnegnän Zon bälle gäz e.}      \extran{Herodias was looking for ways to kill John.}}}
  \variants{\freevarlab \vartext{yagyag}}}

\entry{yaiya}{\headword{yaiya}
  \variantof{Unspecified Variant of \vartext{yaya}}  \pronnote{jaija}}

\entry{yäkäl}{\headword{yäkäl}  \note{Family ties}  \pronnote{jəkəl}
  \pos{Kinship noun}  \sense{
    \definition{cousin}    \note{Family ties}    \example{
}    \example{
}    \example{
}}}

\entry{yäkälnda}{\headword{yäkälnda}
  \pos{Kinship noun}  \sense{
    \definition{uncle (one's mother's sister's husband)}}}

\entry{yäko}{\headword{yäko}
  \variantof{Unspecified Variant of \vartext{yoko}}}

\entry{yal}{\headword{yal}  \pronnote{jal}
  \pos{Noun}  \sense{
    \definition{yellow-billed kingfisher}    \scientificname{Syma torotoro}    \example{
}    \example{
}    \example{
}}}

\entry{Yam}{\headword{Yam}  \note{Place names around Limol}  \pronnote{jam}
  \pos{Proper noun}  \sense{
    \definition{Yam (toponym)}    \note{Place names around Limol}    \example{
}    \example{
}    \example{
}}}

\entry{yämak}{\headword{yämak}  \lexeme{yämak}  \pronnote{jəmak}
  \pos{Noun}  \sense{
    \definition{type of big tree found in the bush and by the river}}}

\entry{yämän}{\headword{yämän}  \note{Food}  \pronnote{jəmən}
  \pos{Noun}  \sense{
    \definition{type of big tuber with a reddish interior (not a yam)}    \note{Food}    \example{
}    \example{
}    \example{
}}}

\entry{Yamayama}{\headword{Yamayama}  \note{Places around Limol}  \pronnote{jamajama}
  \pos{Proper noun}  \sense{
    \definition{Yamayama (toponym)}    \note{Places around Limol}    \example{
}    \example{
}    \example{
}}}

\entry{yämbäg}{\headword{yämbäg}
  \pos{Transitive verb}  \sense{
    \definition{to disown, repudiate}    \example{
      \exsen{Ngäna mɨnyi ddone abal bam baembäg.}      \extran{I will never disown you.}}}}

\entry{Yamega}{\headword{Yamega}  \note{Place names}  \pronnote{jameɡa}
  \pos{Proper noun}  \sense{
    \definition{Iamega (Wipi-speaking village in Oriomo-Bituri Rural LLG)}    \note{Place names}    \example{
}    \example{
}    \example{
}}
  \variants{\SpellingVariant \vartext{Iamega}}}

\entry{Yamkong}{\headword{Yamkong}
  \pos{Noun}  \sense{
    \definition{name of clan}}}

\entry{yänddäna}{\headword{yänddäna}
  \variantof{Unspecified Variant of \vartext{enddäna}}  \pronnote{jənɖ͡ʐəna}}

\entry{yänkllollang}{\headword{yänkllollang}  \note{Bees}  \pronnote{jənkɽoɽaŋ}
  \pos{Modifier}  \sense{
    \definition{forgetful}    \note{Bees}    \example{
}    \example{
}}}

\entry{yante}{\headword{yante}  \pronnote{jante}
  \pos{Noun}  \sense{
    \definition{type of large tree that grows in the grassland with white flowers and wood used for house sticks}}}

\entry{Yao}{\headword{Yao}
  \pos{Proper noun}  \sense{
    \definition{Taeme language}}}

\entry{=yaralla}{\headword{=yaralla}
  \pos{Auxiliary verb}  \sense{
    \definition{aux.prs.1|2nsgA>duP}}}

\entry{=yaralle}{\headword{=yaralle}  \pronnote{jaraɽe}
  \pos{Auxiliary verb}  \sense{
    \definition{aux.prs.2sgA>duP}}}

\entry{=yarallo}{\headword{=yarallo}
  \pos{Auxiliary verb}  \sense{
    \definition{aux.prs.3nsgA>duP}}}

\entry{yärany}{\headword{yärany}
  \variantof{Unspecified Variant of \vartext{erany}}  \pronnote{jəraɲ}}

\entry{Yarbab}{\headword{Yarbab}  \pronnote{jarbab}
  \pos{Proper noun}  \sense{
    \definition{male personal name}}}

\entry{yärmuyärmu}{\headword{yärmuyärmu}  \note{Trees}  \pronnote{jərmujərmu}
  \pos{Noun}  \sense{
    \definition{type of tree}    \note{Trees}    \example{
}    \example{
}    \example{
}}}

\entry{yarte}{\headword{yarte}  \note{Trees}  \pronnote{jarte}
  \pos{Noun}  \sense{
    \definition{type of tree with young wood used for house sticks}    \note{Trees}    \example{
}    \example{
}    \example{
}}}

\entry{Yarte}{\headword{Yarte}  \note{Places around Limol}  \pronnote{jarte}
  \pos{Proper noun}  \sense{
    \definition{Yarte (camping place)}    \note{Places around Limol}    \example{
}    \example{
}    \example{
}}}

\entry{yäru}{\headword{yäru}  \lexeme{yäru}  \pronnote{jəru}
  \pos{Noun}  \sense{
    \definition{type of small tree with thorns}}}

\entry{yaru}{\headword{yaru}  \note{Food}  \pronnote{jaru}
  \pos{Noun}  \sense{
    \definition{type of yam with a purple interior, hairs, and thorns}    \note{Food}    \example{
}    \example{
}    \example{
}}}

\entry{yaryem}{\headword{yaryem}  \note{Food}  \pronnote{jarjem}
  \pos{Noun}  \sense{
    \definition{type of big yam with a white interior, hairs, and thorns}    \note{Food}    \example{
}    \example{
}    \example{
}}}

\entry{yatän}{\headword{yatän}
  \pos{Transitive verb}  \sense{
    \definition{to fetch (water)}    \example{
      \exsen{Oblle ine neyatän.}      \extran{[You] fetch water for him.}}}}

\entry{yätt}{\headword{yätt}  \pronnote{jəʈ͡ʂ}
  \pos{Noun}  \sense{
    \definition{forehead}    \example{
}    \example{
}    \example{
}}}

\entry{yattän}{\headword{yattän}
  \pos{Intransitive verb}  \sense{
    \definition{to disembark, get off, get out}    \example{
      \exsen{Gall atta goyattänän.}      \extran{She got out of the canoe.}}}}

\entry{Yau}{\headword{Yau}
  \variantof{SpellingVariant of \vartext{Yow}}}

\entry{yaul}{\headword{yaul}  \lexeme{yaul}  \pronnote{jaul}
  \pos{Noun}  \sense{
    \definition{type of long yam with a white interior and few hairs}}}

\entry{Yawani}{\headword{Yawani}  \pronnote{jawani}
  \pos{Proper noun}  \sense{
    \definition{male personal name}}}

\entry{Yawen}{\headword{Yawen}
  \pos{Proper noun}  \sense{
    \definition{female personal name}}
  \variants{\Unspecified Variant \vartext{Yawin}}}

\entry{Yawin}{\headword{Yawin}
  \variantof{Unspecified Variant of \vartext{Yawen}}  \pronnote{jawin}}

\entry{yaya}{\headword{yaya}  \note{Fishing}  \pronnote{jaja}
  \pos{Kinship noun}  \sense{
    \definition{father}    \note{Fishing}    \example{
}    \example{
}    \example{
}}
  \variants{\Unspecified Variant \vartext{yaiya}}}

\entry{yayo}{\headword{yayo}  \note{Animals/Plants}  \pronnote{jajo}
  \pos{Verb}  \sense{
    \definition{to slither}    \note{Animals/Plants}    \example{
      \exsen{yayoang dan}      \extran{he is slithering}}    \example{
}    \example{
}}}

\entry{=ye}{\headword{=ye}
  \variantof{Unspecified Variant of \vartext{=aeya}}}

\entry{year}{\headword{year}
  \variantof{Dialectal Variant of \vartext{English loanword}}}

\entry{years}{\headword{years}
  \variantof{Dialectal Variant of \vartext{English loanword}}}

\entry{yebdo}{\headword{yebdo}
  \variantof{SpellingVariant of \vartext{ebdo}}}

\entry{yerngän}{\headword{yerngän}
  \variantof{Unspecified Variant of \vartext{irängän}}  \pronnote{jerŋən}}

\entry{Yesu}{\headword{Yesu}  \pronnote{jesu}
  \pos{Proper noun}  \sense{
    \definition{Jesus}}}

\entry{yetedduem}{\headword{yetedduem}  \pronnote{jeteɖ͡ʐuem}
  \pos{Phrase}  \sense{
    \definition{raise eyebrows}}}

\entry{yewede}{\headword{yewede}
  \variantof{freevarlab of \vartext{ewede}}}

\entry{yid}{\headword{yid}
  \pos{Noun}  \sense{
    \definition{liquid extracted from a plant}    \example{
      \exsen{nge id}      \extran{coconut cream}}}
  \variants{\SpellingVariant \vartext{id, iid}}
  \variants{\Unspecified Variant \vartext{yɨd}}}

\entry{yimne}{\headword{yimne}
  \variantof{SpellingVariant of \vartext{imne}}}

\entry{Yina}{\headword{Yina}  \pronnote{jina}
  \pos{Proper noun}  \sense{
    \definition{female personal name}}}

\entry{yinbo}{\headword{yinbo}
  \variantof{SpellingVariant of \vartext{inbo}}}

\entry{yindrang}{\headword{yindrang}
  \variantof{SpellingVariant of \vartext{indrang}}  \pronnote{jindraŋ}}

\entry{yinu}{\headword{yinu}
  \variantof{SpellingVariant of \vartext{inu}}  \pronnote{jinu}}

\entry{yirɨm}{\headword{yirɨm}
  \variantof{SpellingVariant of \vartext{Iräm}}}

\entry{yɨb}{\headword{yɨb}  \note{Food}  \pronnote{jɨb}
  \pos{Noun}  \sense{
    \definition{type of yam}    \note{Food}    \example{
}    \example{
}    \example{
}}}

\entry{yɨd}{\headword{yɨd}
  \variantof{Unspecified Variant of \vartext{yid}}  \pronnote{jɪd}}

\entry{yɨdmeny}{\headword{yɨdmeny}  \pronnote{jɪdmeɲ}
  \pos{Noun}  \sense{
    \definition{the seventh and final stage of sago growth during which the pith will not yield any starch}    \example{
      \exsen{Tätäm ubi yɨdmeny sana de däkämeyo.}      \extran{Yesterday, they squeezed a dry sago palm.}}}}

\entry{yo}{\headword{yo}  \note{Parts of a reptile}  \pronnote{jo}
  \pos{Noun}  \sense{
    \definition{liver}    \note{Parts of a reptile}    \example{
}    \example{
}    \example{
}}}

\entry{yobeg}{\headword{yobeg}  \note{Trees}  \pronnote{jobeɡ}
  \pos{Noun}  \sense{
    \definition{type of cultivated shrub with white and yellow flowers and long leaves used to tie yam shoots to yam sticks}    \note{Trees}    \example{
}    \example{
}    \example{
}}}

\entry{yogoll}{\headword{yogoll}  \lexeme{yogoll}  \pronnote{joɡoɽ}
  \pos{Noun}  \sense{
    \definition{rain}    \example{
      \exsen{Sisri yogoll allan.}      \extran{It is raining now.}}}}

\entry{yogoll kutt}{\headword{yogoll kutt}  \note{Weather}  \pronnote{joɡoɽ kuʈʂ}
  \pos{Noun}  \sense{
    \definition{black cloud}    \note{Weather}    \example{
}    \example{
}    \example{
}}}

\entry{yogoll särem}{\headword{yogoll särem}  \lexeme{yogoll särem}  \pronnote{joɡoɽ sərem}
  \pos{Noun}  \sense{
    \definition{dark rain cloud}}}

\entry{Yokas}{\headword{Yokas}  \note{Places around Limol}  \pronnote{jokas}
  \pos{Proper noun}  \sense{
    \definition{Yokas (camping place in Limol)}    \note{Places around Limol}    \example{
}    \example{
}    \example{
}}}

\entry{yoko}{\headword{yoko}  \lexeme{yoko}  \pronnote{joko}
  \pos{Noun}  \sense{
    \definition{type of cane used for building houses, bows, and canoes}}
  \variants{\Unspecified Variant \vartext{yäko}}}

\entry{Yokon}{\headword{Yokon}
  \pos{Proper noun}  \sense{
    \definition{personal name}}}

\entry{Yokong}{\headword{Yokong}
  \pos{Proper noun}  \sense{
    \definition{language name}}}

\entry{yon}{\headword{yon}  \pronnote{jon}
  \pos{Noun}  \sense{
    \definition{dream}    \example{
      \exsen{Ngäna käza yon de wätaran.}      \extran{I dreamed about crocodiles.}}    \example{
}    \example{
}}}

\entry{Yop}{\headword{Yop}
  \pos{Proper noun}  \sense{
    \definition{Yop (toponym)}}}

\entry{yorko}{\headword{yorko}  \note{Trees}  \pronnote{jorko}
  \pos{Noun}  \sense{
    \definition{type of large cane found in the bush}    \note{Trees}    \example{
}    \example{
}    \example{
}}}

\entry{yorkoll}{\headword{yorkoll}  \pronnote{jorkoɽ}
  \pos{Noun}  \sense{
    \definition{dirt}}}

\entry{yorkollang}{\headword{yorkollang}  \pronnote{jorkoɽaŋ}
  \pos{Modifier}  \sense{
    \definition{dirty}    \example{
      \exsen{Bongo ddone yorkollang alle.}      \extran{You are not dirty.}}}}

\entry{Yoteang}{\headword{Yoteang}  \note{Places around Limol}  \pronnote{joteaŋ}
  \pos{Proper noun}  \sense{
    \definition{Yoteang (toponym)}    \note{Places around Limol}    \example{
}    \example{
}    \example{
}}}

\entry{yoto}{\headword{yoto}  \note{Birds}  \pronnote{joto}
  \pos{Noun}  \sense{
    \definition{type of biting bee found in trees}    \note{Birds}    \example{
}    \example{
}    \example{
}}}

\entry{you}{\headword{you}
  \variantof{freevarlab of \vartext{English loanword}}}

\entry{Yow}{\headword{Yow}  \pronnote{jou}
  \pos{Proper noun}  \sense{
    \definition{Yau (in Gogodala Rural LLG)}    \example{
}    \example{
}    \example{
}}
  \variants{\SpellingVariant \vartext{Yau}}}

\entry{yowa}{\headword{yowa}  \note{Parts of the body}  \pronnote{jowa}
  \pos{Noun}  \sense{
    \definition{vagina}    \note{Parts of the body}    \example{
      \exsen{Eiz itrell a lla bo wa piye me a mälla bo yowa käpät me dadeg.}      \extran{The HIV virus exists in men's sperm and in women's vaginal fluid.}}}}

\entry{yowede}{\headword{yowede}
  \variantof{Unspecified Variant of \vartext{ewede}}}

\entry{yu}{\headword{yu}  \pronnote{ju}
  \pos{Noun}  \sense{\textbf{1.} 
    \definition{fire}    \example{
      \exsen{Sali yu di mermerae lläklläk eran.}      \extran{Sali is spreading the fire nicely.}}}  \sense{\textbf{2.} 
    \definition{firewood}}  \sense{\textbf{3.} 
    \definition{to cook over fire}    \example{
      \exsen{Obo wätät de yu dägagän, duwem gogon.}      \extran{He cooked his food on the fire and ate it.}}    \example{
      \exsen{Pa de yu dägaebneyo.}      \extran{They cooked the birds on the fire.}}    \example{
      \exsen{Kollba de yu deagän.}      \extran{He cooked his two fish on the fire.}}}}

\entry{yu bägäl}{\headword{yu bägäl}  \pronnote{ju bəɡəl}
  \pos{Noun}  \sense{
    \definition{gun, firearm}    \example{
}    \example{
}    \example{
}}}

\entry{yu bäng}{\headword{yu bäng}  \pronnote{ju bəŋ}
  \pos{Noun}  \sense{
    \definition{firestick (to start a fire)}    \example{
      \exsen{Auma watta ako ngäna yu bäng de beyangäs.}      \extran{I'll bring the firestick back from the grave.}}    \example{
}    \example{
}}}

\entry{yu dumbrel}{\headword{yu dumbrel}  \note{Parts of a fire}  \pronnote{ju dumbrel}
  \pos{Noun}  \sense{
    \definition{flame}    \note{Parts of a fire}    \example{
}    \example{
}    \example{
}}}

\entry{yu kire}{\headword{yu kire}  \note{Everything in the house}  \pronnote{ju kire}
  \pos{Noun}  \sense{
    \definition{firewood}    \note{Everything in the house}    \example{
}    \example{
}    \example{
}}}

\entry{yu menanen}{\headword{yu menanen}  \pronnote{ju menanen}
  \pos{Noun}  \sense{
    \definition{heat of the fire}}}

\entry{yu molle täre}{\headword{yu molle täre}  \pronnote{ju moɽe təre}
  \pos{Noun}  \sense{
    \definition{feast type}}}

\entry{yu ngongom}{\headword{yu ngongom}  \pronnote{ju ŋoŋom}
  \pos{Noun}  \sense{
    \definition{small fire}}}

\entry{yu tärkmoll}{\headword{yu tärkmoll}
  \variantof{Fast Speech Variant of \vartext{yu torkomoll}}}

\entry{yu torkomoll}{\headword{yu torkomoll}  \note{Parts of a fire}  \pronnote{ju torkomoɽ}
  \pos{Noun}  \sense{
    \definition{charcoal}    \note{Parts of a fire}    \example{
}    \example{
}    \example{
}}
  \variants{\Fast Speech Variant \vartext{yu tärkmoll}}}

\entry{yu ttängäm}{\headword{yu ttängäm}  \note{Dialect names}  \pronnote{ju ʈ͡ʂəŋəm}
  \pos{Noun}  \sense{
    \definition{hell}    \note{Dialect names}    \example{
}    \example{
}    \example{
}}}

\entry{yu ttätta}{\headword{yu ttätta}  \pronnote{ju ʈ͡ʂəʈ͡ʂa}
  \pos{Noun}  \sense{
    \definition{burning wood, burnt wood}}}

\entry{yubud}{\headword{yubud}  \note{Trees}  \pronnote{jubud}
  \pos{Noun}  \sense{
    \definition{type of tree}    \note{Trees}    \example{
}    \example{
}    \example{
}}}

\entry{yuddädda}{\headword{yuddädda}  \note{Types of Taro}  \pronnote{juɖ͡ʐəɖ͡ʐa}
  \pos{Noun}  \sense{
    \definition{type of palm with branches used for armbands}    \note{Types of Taro}    \example{
}    \example{
}    \example{
}}}

\entry{Yuga}{\headword{Yuga}  \pronnote{juɡa}
  \pos{Proper noun}  \sense{
    \definition{male personal name}}}

\entry{Yugui}{\headword{Yugui}  \pronnote{juɡui}
  \pos{Proper noun}  \sense{
    \definition{female personal name}}}

\entry{Yun}{\headword{Yun}  \pronnote{jun}
  \pos{Proper noun}  \sense{
    \definition{female personal name}}}

\entry{yunipom}{\headword{yunipom}
  \pos{Noun}  \sense{
    \definition{uniform}}}

\entry{yunu}{\headword{yunu}
  \variantof{Unspecified Variant of \vartext{inu}}  \pronnote{junu}}

\entry{yure}{\headword{yure}  \note{School}  \pronnote{jure}
  \pos{Noun}  \sense{
    \definition{type of sago}    \note{School}    \example{
}    \example{
}    \example{
}}}

\entry{yuru}{\headword{yuru}  \pronnote{juru}
  \pos{Noun}  \sense{
    \definition{type of pandanus}    \example{
}    \example{
}    \example{
}}}

\entry{yurwe}{\headword{yurwe}
  \pos{Noun}  \sense{
    \definition{type of tree}    \example{
}}}

\entry{yus}{\headword{yus}
  \pos{Transitive coverb}  \sense{
    \definition{to use}    \example{
      \exsen{Kumuddäga eka de ngäna yus anggan.}      \extran{I use three languages.}}}}

\entry{yuwet~}{\headword{yuwet~}
  \variantof{freevarlab of \vartext{RED~}}}

\entry{yuwet}{\headword{yuwet}  \note{Market}  \pronnote{juwet}
  \pos{Noun}  \sense{
    \definition{short period of time}    \note{Market}    \example{
      \exsen{Ngäna bam mɨnyi yuwet e pen de banttog.}      \extran{I will lend a pen to you (lit. give you a pen for a short time).}}    \example{
}}}

\entry{yuwetyuwet}{\headword{yuwetyuwet}  \pronnote{juwetjuwet}
  \pos{Adverb}  \sense{
    \definition{temporarily, briefly}    \example{
      \exsen{Ubi siti we yuwetyuwet gobällnän.}      \extran{They went into the city briefly.}}}}

\entry{yuwog}{\headword{yuwog}  \note{Oral Traditions}  \pronnote{juwoɡ}
  \pos{Quantifier}  \sense{\textbf{1.} 
    \definition{many, a lot of (for countable nouns)}    \note{Oral Traditions}    \example{
      \exsen{yuwog abal mani}      \extran{a lot of money}}    \example{
      \exsen{Ngämo nagnag a yuwog abal dagän.}      \extran{I have many friends.}}    \example{
      \exsen{Yuwog abal kollba de ngämi daittaemnalla.}      \extran{We (excl.) were catching many fish.}}}  \sense{\textbf{2.} 
    \definition{for no reason, on a whim}    \note{Oral Traditions}    \example{
      \exsen{Ngäna yuwog dae ibi allan Bisuaka we.}      \extran{I'm going to Bisuaka on a whim.}}    \example{
      \exsen{Ngäna yuwog dae giddoll allan.}      \extran{I am living without any food.}}}}

\entry{yuwogyuwog}{\headword{yuwogyuwog}
  \variantof{Unspecified Variant of \vartext{yuwoyuwog}}}

\entry{yuwoyuwog}{\headword{yuwoyuwog}
  \pos{Adverb}  \sense{\textbf{1.} 
    \definition{needlessly, excessively}    \example{
      \exsen{Lla da o mälla da obaoba ddone yuwoyuwog bulignegnän.}      \extran{Men and women should not be having sex with each other needlessly.}}}  \sense{\textbf{2.} 
    \definition{not completely, not properly}    \example{
      \exsen{Sisor llɨg a ubi yuwoyuwog panypeny erallo Ende eka de.}      \extran{Young kids aren't speaking Ende properly.}}}
  \variants{\Unspecified Variant \vartext{yuwogyuwog}}}

\chapter{Z}


\entry{za~}{\headword{za~}
  \variantof{Infinitival reduplicant of \vartext{RED~}}}

\entry{za}{\headword{za}
  \pos{Noun}  \sense{
    \definition{thing}    \example{
      \exsen{Oba za da gänyag.}      \extran{Here are their things.}}}
  \variants{\SpellingVariant \vartext{zaa}}}

\entry{zäbo}{\headword{zäbo}  \pronnote{zəbo}
  \pos{Noun}  \sense{
    \definition{yellow-streaked lory}    \scientificname{Chalcopsitta scintillata}}}

\entry{zaga}{\headword{zaga}  \pronnote{zaɡa}
  \pos{Personal pronoun}  \sense{
    \definition{self (forms reflexive pronouns)}    \example{
      \exsen{Ngäna ngämo zaga butruwam.}      \extran{I will lay myself down.}}}}

\entry{zagu}{\headword{zagu}  \pronnote{zaɡu}
  \pos{Noun}  \sense{
    \definition{type of sugarcane-like plant that grows in the swamp}}}

\entry{Zakae}{\headword{Zakae}
  \pos{Proper noun}  \sense{
    \definition{male personal name}}
  \variants{\SpellingVariant \vartext{Zakai}}}

\entry{Zakai}{\headword{Zakai}
  \variantof{SpellingVariant of \vartext{Zakae}}}

\entry{zäm}{\headword{zäm}
  \pos{Intransitive verb}  \sense{
    \definition{to pass through}    \example{
      \exsen{Ubi tätäm gänyagaeya dinzämän.}      \extran{They passed through here yesterday.}}}}

\entry{zämae}{\headword{zämae}  \pronnote{zəmae}
  \pos{Transitive verb}  \sense{
    \definition{to pour, put, transfer}    \example{
      \exsen{Ubi ine de zämaenen erallo.}      \extran{They are pouring water.}}    \example{
      \exsen{Kollba dängälläbeya, nyäng e danzämaeya.}      \extran{We bought fish and put them in a bag.}}}}

\entry{zämängg}{\headword{zämängg}  \pronnote{zəməŋɡ}
  \pos{Intransitive verb}  \sense{
    \definition{to be in heaps}}}

\entry{zäme}{\headword{zäme}
  \variantof{Dialectal Variant of \vartext{zɨme}}  \pronnote{zəme}}

\entry{zämllall}{\headword{zämllall}
  \pos{Intransitive coverb}  \sense{
    \definition{to pass by}    \example{
      \exsen{Bogo zämllall allan.}      \extran{He is passing by.}}}}

\entry{zämllallang}{\headword{zämllallang}  \pronnote{zəmɽaɽaŋ}
  \pos{Noun}  \sense{
    \definition{passerby}}}

\entry{zämllazämllall}{\headword{zämllazämllall}  \note{Verbs of Motion}  \pronnote{zəmɽazəmɽaɽ}
  \pos{Adverb}  \sense{
    \definition{nonstop}    \note{Verbs of Motion}    \example{
      \exsen{Bogo zämllazämllall ibi allan.}      \extran{He walks without stopping.}}    \example{
}    \example{
}}}

\entry{zan}{\headword{zan}  \pronnote{zan}
  \pos{Intransitive verb}  \sense{\textbf{1.} 
    \definition{to enter, go in, go into}    \example{
      \exsen{lla zazer ma märäll me}      \extran{sized so that a person could enter}}    \example{
      \exsen{Ngäna mɨnyi bozen.}      \extran{I will enter.}}    \example{
      \exsen{Pollgo da täl dɨdɨr ik e gozenän.}      \extran{The frog went into the dry bamboo.}}    \example{
      \exsen{Bibi ngämo ma ik e zarnen amalla.}      \extran{You all are going in my house.}}}  \sense{\textbf{2.} 
    \definition{to put in}    \example{
      \exsen{Polis a obom säremang ma ik i däzaneyo.}      \extran{The police put him in prison.}}    \example{
      \exsen{Bogo up wo de nyäng e dazernän.}      \extran{He put the ripe bananas in the bag.}}}}

\entry{Zanger}{\headword{Zanger}
  \pos{Proper noun}  \sense{
    \definition{female personal name}}}

\entry{zanggae}{\headword{zanggae}  \pronnote{zaŋɡae}
  \pos{Intransitive verb}  \sense{
    \definition{to roam, go around}    \example{
      \exsen{Ngäna däba ngata me gonzagae.}      \extran{I wandered around in that area.}}}}

\entry{zangom}{\headword{zangom}  \pronnote{zaŋom}
  \pos{Temporal adverb}  \sense{
    \definition{now or today}}}

\entry{Zanor}{\headword{Zanor}  \pronnote{zanor}
  \pos{Proper noun}  \sense{
    \definition{Zanor (in Gogodala Rural LLG; GPS: -8.459817, 142.688025)}    \example{
}    \example{
}    \example{
}}}

\entry{zanzi}{\headword{zanzi}
  \pos{Intransitive verb}  \sense{
    \definition{to settle}    \example{
      \exsen{Ge llɨg a dedme gunziwän.}      \extran{This boy settled there.}}}}

\entry{Zarma}{\headword{Zarma}  \note{Places around Limol}  \pronnote{zarma}
  \pos{Proper noun}  \sense{
    \definition{Zarma (toponym)}    \note{Places around Limol}    \example{
}    \example{
}    \example{
}}}

\entry{zarmeny}{\headword{zarmeny}  \note{Yams}  \pronnote{zarmeɲ}
  \pos{Noun}  \sense{
    \definition{type of long yam with a white interior, hairs, and no thorns}    \note{Yams}    \example{
}    \example{
}    \example{
}}}

\entry{zawatt}{\headword{zawatt}
  \pos{Noun}  \sense{
    \definition{vagina}}}

\entry{zazaba}{\headword{zazaba}  \pronnote{zazaba}
  \pos{Noun}  \sense{
    \definition{type of bag}    \example{
      \exsen{Mänmän a sana de zazaba we dändäraebeyo.}      \extran{The girls filled their sago bags with sago.}}}}

\entry{zaze}{\headword{zaze}  \pronnote{zaze}
  \pos{Intransitive verb}  \sense{\textbf{1.} 
    \definition{to give birth, lay}    \example{
      \exsen{Ngäna käp tumang zazeag dan.}      \extran{I lay many eggs.}}}  \sense{\textbf{2.} 
    \definition{generation, siblings born of the same mother}    \example{
      \exsen{ngämo baba bakmall zaze}      \extran{my father's siblings (lit. those in the generation with my father)}}}  \sense{\textbf{3.} 
}}

\entry{zaa}{\headword{zaa}
  \variantof{SpellingVariant of \vartext{za}}}

\entry{Zedem}{\headword{Zedem}  \note{comes from ZM, "zealous missionary", Tonny Warama's band name.}  \pronnote{zedem}
  \pos{Proper noun}  \sense{
    \definition{male personal name}    \note{comes from ZM, "zealous missionary", Tonny Warama's band name.}}}

\entry{zeg}{\headword{zeg}  \pronnote{zeɡ}
  \pos{Intransitive verb}  \sense{
    \definition{to be born}    \example{
      \exsen{Bongo sisri ebdo ag me ozegalle.}      \extran{You were born this morning.}}    \example{
      \exsen{Ngäna do gozeg, Kibobma me.}      \extran{I was born there, in Kibobma.}}}}

\entry{Zegma}{\headword{Zegma}  \pronnote{zeɡma}
  \pos{Proper noun}  \sense{
    \definition{Zegma (camping place)}    \example{
}    \example{
}    \example{
}}}

\entry{Zeims}{\headword{Zeims}  \pronnote{zeims}
  \pos{Proper noun}  \sense{
    \definition{male personal name}}
  \variants{\SpellingVariant \vartext{James}}}

\entry{Zekraeya}{\headword{Zekraeya}
  \pos{Proper noun}  \sense{
    \definition{male personal name}}}

\entry{zel}{\headword{zel}
  \pos{Noun}  \sense{
    \definition{jail}    \example{
      \exsen{Obom zel ma ik i näzanallo.}      \extran{They put him in jail.}}}}

\entry{Zelma}{\headword{Zelma}  \pronnote{zelma}
  \pos{Proper noun}  \sense{
    \definition{female personal name}}}

\entry{zem}{\headword{zem}  \pronnote{zem}
  \pos{Noun}  \sense{
    \definition{germ}}}

\entry{Zemila}{\headword{Zemila}
  \variantof{SpellingVariant of \vartext{Jemila}}}

\entry{Zen}{\headword{Zen}
  \pos{Proper noun}  \sense{
    \definition{male personal name}}}

\entry{Zeri}{\headword{Zeri}
  \variantof{SpellingVariant of \vartext{Jerry}}}

\entry{Zeriko}{\headword{Zeriko}  \pronnote{zeriko}
  \pos{Proper noun}  \sense{
    \definition{Jericho}}}

\entry{zib}{\headword{zib}  \note{similar species: bunkombunkom}  \pronnote{zib}
  \pos{Noun}  \sense{
    \definition{type of big tree that grows in the bush}    \note{similar species: bunkombunkom}    \example{
}    \example{
}    \example{
}}}

\entry{zib mäka}{\headword{zib mäka}  \note{Nature}  \pronnote{zib məka}
  \pos{Noun}  \sense{
    \definition{type of introduced banana}    \note{Nature}    \example{
}    \example{
}    \example{
}}}

\entry{zigae}{\headword{zigae}  \pronnote{ziɡae}
  \pos{Transitive verb}  \sense{
    \definition{to wrap}    \example{
      \exsen{Ngäna mänyän de ttam alle zigae eran.}      \extran{I am wrapping the mänyän fish with leaves.}}    \example{
      \exsen{Ngämi sana mängalae dazgiaebeya.}      \extran{We (excl.) quickly wrapped the sago.}}}}

\entry{Zina}{\headword{Zina}
  \pos{Proper noun}  \sense{
    \definition{female personal name}}
  \variants{\SpellingVariant \vartext{Gina}}}

\entry{zire}{\headword{zire}  \pronnote{zire}
  \pos{Noun}  \sense{
    \definition{barramundi}    \scientificname{Lates calcarifer}    \example{
}    \example{
}    \example{
}}}

\entry{ziz}{\headword{ziz}  \pronnote{ziz}
  \pos{Noun}  \sense{
    \definition{insect}    \example{
}    \example{
}    \example{
}}}

\entry{zizag}{\headword{zizag}  \pronnote{zizaɡ}
  \pos{Noun}  \sense{
    \definition{owner, master, lord}    \example{
      \exsen{Zizag Adi}      \extran{the Lord God}}    \example{
      \exsen{ttängäm zizag}      \extran{landowner}}    \example{
      \exsen{Gänyaolle gall zizag a ngänawaenen.}      \extran{The owner of this canoe is me.}}}}

\entry{zizi}{\headword{zizi}  \pronnote{zizi}
  \pos{Transitive verb}  \sense{
    \definition{to uncover, lift}    \example{
      \exsen{Angde ngämi ttägäll de däziya, duwem gogmam.}      \extran{When we (excl.) uncovered the mumu, we ate.}}    \example{
      \exsen{Wel gullbe da gongttägän a mägda bo pite de däzinän.}      \extran{A strong wind came and lifted his mother's skirt.}}}}

\entry{zɨme}{\headword{zɨme}  \pronnote{zɨme}
  \pos{TAM particle}  \sense{
    \definition{already}    \example{
      \exsen{Ngäna bam zime ikop nagan.}      \extran{I already saw you.}}    \example{
      \exsen{Llɨg a zime ngänaeka gognegnän kallkäll atta.}      \extran{The children were already crying from cold.}}}
  \variants{\Dialectal Variant \vartext{zäme}}}

\entry{Zo}{\headword{Zo}
  \pos{Proper noun}  \sense{
    \definition{male personal name}}
  \variants{\SpellingVariant \vartext{Joe}}}

\entry{zo}{\headword{zo}  \pronnote{zo}
  \pos{Noun}  \sense{
    \definition{fawn-breasted bowerbird}    \scientificname{Chlamydera cerviniventris}    \example{
}    \example{
}    \example{
}}}

\entry{zobo ik}{\headword{zobo ik}  \pronnote{zobo ik}
  \pos{Noun}  \sense{
    \definition{walling (bark on house between sago walls and roof)}}}

\entry{zogam}{\headword{zogam}  \pronnote{zoɡam}
  \pos{Noun}  \sense{
    \definition{rat}}}

\entry{Zon}{\headword{Zon}
  \variantof{SpellingVariant of \vartext{John}}  \pronnote{zon}}

\entry{Zonas}{\headword{Zonas}  \pronnote{zonas}
  \pos{Proper noun}  \sense{
    \definition{male personal name}}
  \variants{\SpellingVariant \vartext{Jonas}}}

\entry{Zoni}{\headword{Zoni}
  \variantof{SpellingVariant of \vartext{Johnny}}}

\entry{zora}{\headword{zora}  \lexeme{zora}  \pronnote{zora}
  \pos{Noun}  \sense{
    \definition{sharp stick for peeling sago before beating it}}}

\entry{Zosep}{\headword{Zosep}  \pronnote{zosep}
  \pos{Proper noun}  \sense{
    \definition{male personal name}}
  \variants{\SpellingVariant \vartext{Joseph}}}

\entry{zowag}{\headword{zowag}  \pronnote{zowaɡ}
  \pos{Modifier}  \sense{
    \definition{hoarse}    \example{
      \exsen{Pitapo bo inkɨm a zowag dan.}      \extran{Pitapo's voice is hoarse.}}}}

\entry{zozo}{\headword{zozo}  \pronnote{zozo}
  \pos{Intransitive verb}  \sense{
    \definition{to rot, go bad}    \example{
      \exsen{Kumddäga ebdo ddɨgatt alle up a mɨnyi bozowän.}      \extran{After three days, the banana will rot.}}    \example{
      \exsen{Lläkäm a mamall zozoang dan.}      \extran{The mushroom goes bad quickly.}}}}

\entry{Zudas}{\headword{Zudas}  \pronnote{zudas}
  \pos{Proper noun}  \sense{
    \definition{male personal name}}
  \variants{\SpellingVariant \vartext{Judas}}}

\entry{Zudiya}{\headword{Zudiya}  \pronnote{zudija}
  \pos{Proper noun}  \sense{
    \definition{Judea}}}

\entry{Zugu}{\headword{Zugu}  \pronnote{zuɡu}
  \pos{Proper noun}  \sense{
    \definition{male personal name}}}

\entry{Zuli}{\headword{Zuli}
  \pos{Proper noun}  \sense{
    \definition{female personal name}}
  \variants{\SpellingVariant \vartext{Julie}}}

\entry{Zurusalem}{\headword{Zurusalem}  \pronnote{zurusalem}
  \pos{Proper noun}  \sense{
    \definition{Jerusalem}}}

\entry{zuwoe}{\headword{zuwoe}  \pronnote{zuwoe}
  \pos{Transitive verb}  \sense{\textbf{1.} 
    \definition{to shoot}    \example{
      \exsen{Bodog käbama dallän a sɨmell gullbe de dazuwän.}      \extran{Bodog went hunting and shot a boar.}}}  \sense{\textbf{2.} 
    \definition{to pierce; to inject, administer a shot}    \example{
      \exsen{Baba bom sana täkäll da nazunan nying me.}      \extran{The sago thorn pierced Dad in the leg.}}    \example{
      \exsen{Ngänäm doktor da pam alle nazuwan.}      \extran{The doctor gave me an injection.}}}}
